\documentclass[11pt,a4paper,oneside]{book}

% Load all styling and custom packages
% Preamble for IITM Machine Learning Foundations (BSCS2004)
% Common Preamble for IITM BS Textbook Project
\usepackage[utf8]{inputenc}
\usepackage[T1]{fontenc}
\usepackage{geometry}
\geometry{
    a4paper,
    left=0.9in,
    right=0.9in,
    top=1.0in,
    bottom=1.0in
}

% Typography
\usepackage{palatino} % Clean serif body font
\usepackage[scaled]{helvet} % Helvetica for sans-serif
\usepackage{microtype} % Better word spacing and justifications

% Math Packages
\usepackage{amsmath}
\usepackage{amssymb}
\usepackage{amsthm}
\usepackage{mathtools}
\usepackage{bm} % Bold math symbols
\usepackage{array}
\usepackage{booktabs}
\usepackage{pgfplots}
\pgfplotsset{compat=1.18}
\IfFileExists{venndiagram.sty}{%
  \usepackage{venndiagram}%
}{%
  \makeatletter
  \newenvironment{venndiagram2sets}[1][]{%
    \begin{tikzpicture}[scale=0.55, transform shape]%
      \draw[thick] (-2,-1.5) rectangle (3.5,1.5);
      \node at (3.1,1.1) {$S$};
      \draw[thick] (0,0) circle (1.2cm);
      \draw[thick] (1.5,0) circle (1.2cm);
      \node at (-0.5,0.7) {$A$};
      \node at (2.0,0.7) {$B$};
      \providecommand{\fillA}{\fill[gray!40,opacity=0.6] (0,0) circle (1.2cm);}%
      \providecommand{\fillB}{\fill[gray!40,opacity=0.6] (1.5,0) circle (1.2cm);}%
      \providecommand{\fillACapB}{%
        \begin{scope}%
          \clip (0,0) circle (1.2cm);%
          \fill[gray!60] (1.5,0) circle (1.2cm);%
        \end{scope}%
      }%
      \providecommand{\fillOnlyA}{%
        \begin{scope}%
          \clip (-2,-1.5) rectangle (3.5,1.5);%
          \fill[gray!40] (0,0) circle (1.2cm);%
          \fill[white] (1.5,0) circle (1.2cm);%
          \draw[thick] (0,0) circle (1.2cm);%
          \draw[thick] (1.5,0) circle (1.2cm);%
        \end{scope}%
      }%
      \providecommand{\fillNotA}{%
        \fill[gray!30] (-2,-1.5) rectangle (3.5,1.5);%
        \fill[white] (0,0) circle (1.2cm);%
        \draw[thick] (0,0) circle (1.2cm);%
        \draw[thick] (1.5,0) circle (1.2cm);%
      }%
  }{%
    \end{tikzpicture}%
  }
  \makeatother
}

% Headings styling
\usepackage{titlesec}
\titleformat{\chapter}[display]
  {\normalfont\sffamily\huge\bfseries\color{primarydark}}
  {\chaptername\ \thechapter}{20pt}{\Huge}
\titleformat{\section}
  {\normalfont\sffamily\Large\bfseries\color{primarydark}}
  {\thesection}{1em}{}
\titleformat{\subsection}
  {\normalfont\sffamily\large\bfseries\color{primarydark}}
  {\thesubsection}{1em}{}

% Colors
\usepackage{xcolor}
% Tailored Slate Blue/Navy for Theorems
\definecolor{primarydark}{HTML}{1e293b}
\definecolor{primarylight}{HTML}{eff6ff}
\definecolor{primaryborder}{HTML}{3b82f6}

% Tailored Teal for Definitions
\definecolor{tealdef}{HTML}{0f766e}
\definecolor{tealbg}{HTML}{f0fdfa}
\definecolor{tealborder}{HTML}{14b8a6}

% Tailored Purple for Examples
\definecolor{purpleex}{HTML}{6b21a8}
\definecolor{purplebg}{HTML}{faf5ff}
\definecolor{purpleborder}{HTML}{a855f7}

% Tailored Coral for Key Formulas
\definecolor{coralformula}{HTML}{c2410c}
\definecolor{coralbg}{HTML}{fff7ed}
\definecolor{coralborder}{HTML}{f97316}

% Tailored Amber for Warnings
\definecolor{amberwarning}{HTML}{b45309}
\definecolor{amberbg}{HTML}{fffbeb}
\definecolor{amberborder}{HTML}{f59e0b}

% Tailored Crimson/Ruby for Exam Gotchas & Misconceptions
\definecolor{traprad}{HTML}{991b1b}
\definecolor{trapbg}{HTML}{fef2f2}
\definecolor{trapborder}{HTML}{ef4444}

% Semantic Accent & Forest Colors for TikZ
\definecolor{forest}{HTML}{166534}
\definecolor{accent}{HTML}{2563eb}

% YouTube Dark Slate card
\definecolor{ytcardbg}{HTML}{f8fafc}
\definecolor{ytcardborder}{HTML}{cbd5e1}
\definecolor{yttext}{HTML}{ef4444} % YouTube Red for links

% Indexing, QR Codes, Listings & Lists
\usepackage{imakeidx}
\IfFileExists{qrcode.sty}{
    \usepackage{qrcode}
}{
    \newcommand{\qrcode}[2][]{}
}
\usepackage{listings}
\usepackage[inline]{enumitem}

% Hyperref configuration (loaded BEFORE custom tcolorbox environments using \href)
\usepackage{hyperref}
\hypersetup{
    colorlinks=true,
    linkcolor=blue,
    filecolor=magenta,      
    urlcolor=blue,
    pdftitle={IITM BS Textbook},
    pdfpagemode=FullScreen,
}

% Custom tcolorboxes
\usepackage[most]{tcolorbox}
\usetikzlibrary{shadows, calc, arrows.meta, shapes.geometric, positioning}

% Definition Box
\newtcolorbox{definitionbox}[1]{
    colback=tealbg,
    colframe=tealborder,
    coltitle=white,
    fonttitle=\bfseries\sffamily,
    title={Definition: #1},
    enhanced,
    attach boxed title to top left={yshift=-2mm, xshift=4mm},
    boxed title style={colback=tealdef, colframe=tealborder, arc=1mm},
    arc=2mm,
    boxrule=0.8mm,
    top=4mm,
    bottom=3mm,
    left=3mm,
    right=3mm,
    drop shadow={opacity=0.1, color=black}
}

% Theorem/Rule Box
\newtcolorbox{theorembox}[1]{
    colback=primarylight,
    colframe=primaryborder,
    coltitle=white,
    fonttitle=\bfseries\sffamily,
    title={Theorem/Rule: #1},
    enhanced,
    attach boxed title to top left={yshift=-2mm, xshift=4mm},
    boxed title style={colback=primarydark, colframe=primaryborder, arc=1mm},
    arc=2mm,
    boxrule=0.8mm,
    top=4mm,
    bottom=3mm,
    left=3mm,
    right=3mm,
    drop shadow={opacity=0.1, color=black}
}

% Example Box
\newtcolorbox{examplebox}[1]{
    colback=purplebg,
    colframe=purpleborder,
    coltitle=white,
    fonttitle=\bfseries\sffamily,
    title={Example: #1},
    enhanced,
    attach boxed title to top left={yshift=-2mm, xshift=4mm},
    boxed title style={colback=purpleex, colframe=purpleborder, arc=1mm},
    arc=2mm,
    boxrule=0.8mm,
    top=4mm,
    bottom=3mm,
    left=3mm,
    right=3mm,
    drop shadow={opacity=0.1, color=black}
}

% Intuition Box
\newtcolorbox{intuitionbox}[1]{
    colback=purple!5!white,
    colframe=purple!75!black,
    coltitle=white,
    fonttitle=\bfseries\sffamily,
    title={Core Intuition: #1},
    enhanced,
    attach boxed title to top left={yshift=-2mm, xshift=4mm},
    boxed title style={colback=purple!75!black, colframe=purple!90!black, arc=1mm},
    arc=2mm,
    boxrule=0.8mm,
    top=4mm,
    bottom=3mm,
    left=3mm,
    right=3mm,
    drop shadow={opacity=0.1, color=black}
}

% Key Formula Box
\newtcolorbox{formulabox}[1]{
    colback=coralbg,
    colframe=coralborder,
    coltitle=white,
    fonttitle=\bfseries\sffamily,
    title={Key Formula: #1},
    enhanced,
    attach boxed title to top left={yshift=-2mm, xshift=4mm},
    boxed title style={colback=coralformula, colframe=coralborder, arc=1mm},
    arc=2mm,
    boxrule=0.8mm,
    top=4mm,
    bottom=3mm,
    left=3mm,
    right=3mm,
    drop shadow={opacity=0.1, color=black}
}

% Warning Box (Common Mistakes)
\newtcolorbox{warningbox}[1]{
    colback=amberbg,
    colframe=amberborder,
    coltitle=white,
    fonttitle=\bfseries\sffamily,
    title={Warning: #1},
    enhanced,
    attach boxed title to top left={yshift=-2mm, xshift=4mm},
    boxed title style={colback=amberwarning, colframe=amberborder, arc=1mm},
    arc=2mm,
    boxrule=0.8mm,
    top=4mm,
    bottom=3mm,
    left=3mm,
    right=3mm,
    drop shadow={opacity=0.1, color=black}
}

% Exam Trap & Common Misconception Box (Ruby/Crimson)
\newtcolorbox{trapbox}[1]{
    colback=trapbg,
    colframe=trapborder,
    coltitle=white,
    fonttitle=\bfseries\sffamily,
    title={Exam Trap \& Common Misconception: #1},
    enhanced,
    attach boxed title to top left={yshift=-2mm, xshift=4mm},
    boxed title style={colback=traprad, colframe=trapborder, arc=1mm},
    arc=2mm,
    boxrule=0.8mm,
    top=4mm,
    bottom=3mm,
    left=3mm,
    right=3mm,
    drop shadow={opacity=0.1, color=black}
}

% YouTube Video Resource Card
\newtcolorbox{videocard}[2]{
    colback=ytcardbg,
    colframe=ytcardborder,
    enhanced,
    boxrule=0.5mm,
    arc=1.5mm,
    left=3mm,
    right=3mm,
    top=3mm,
    bottom=3mm,
    fontupper=\small\sffamily,
    title={\href{#2}{\textcolor{yttext}{\textbf{Video Resource:}} \textcolor{primarydark}{#1}} \hfill \raisebox{-1.5mm}{\qrcode[height=0.6cm]{#2}}},
    coltitle=primarydark,
    colbacktitle=ytcardbg,
    attach boxed title to top left={yshift=-2mm, xshift=2mm},
    boxed title style={colback=ytcardbg, colframe=ytcardbg},
    underlay={
        \path[draw=ytcardborder, line width=0.5mm] (frame.north west) -- (frame.south west);
    },
    title style={bottom=1mm}
}

% Python Code Snippets Styling
\definecolor{codegreen}{HTML}{10B981}
\definecolor{codegray}{HTML}{6B7280}
\definecolor{codepurple}{HTML}{8B5CF6}
\definecolor{codebg}{HTML}{F8FAFC}
\lstdefinestyle{pystyle}{
    language=Python,
    backgroundcolor=\color{codebg},   
    commentstyle=\color{codegray}\itshape,
    keywordstyle=\color{codepurple}\bfseries,
    numberstyle=\tiny\color{codegray},
    stringstyle=\color{codegreen},
    basicstyle=\ttfamily\footnotesize,
    breakatwhitespace=false,         
    breaklines=true,                 
    captionpos=b,                    
    keepspaces=true,                 
    numbers=left,                    
    numbersep=5pt,                  
    showspaces=false,                
    showstringspaces=false,
    showtabs=false,                  
    tabsize=4,
    frame=single,
    rulecolor=\color{ytcardborder},
    xleftmargin=10pt,
    xrightmargin=10pt
}
\lstset{style=pystyle}

% Custom Math Macros (Common across all courses)
\newcommand{\N}{\mathbb{N}}
\newcommand{\Z}{\mathbb{Z}}
\newcommand{\Q}{\mathbb{Q}}
\newcommand{\R}{\mathbb{R}}
\newcommand{\C}{\mathbb{C}}
\newcommand{\set}[1]{\left\{#1\right\}}
\newcommand{\abs}[1]{\left|#1\right|}
\newcommand{\func}[3]{#1 \colon #2 \to #3}
\newcommand{\limit}[2]{\lim_{#1 \to #2}}
\newcommand{\deriv}[2]{\frac{d#1}{d#2}}
\newcommand{\pderiv}[2]{\frac{\partial#1}{\partial#2}}
\newcommand{\dx}{\,dx}


% Hyperref options for clean PDF bookmarks and page labels
\hypersetup{
    plainpages=false,
    pdfpagelabels=true,
    pdftitle={Machine Learning Foundations - BSCS2004}
}

% Domain-specific packages for Machine Learning
\usepackage{tabularx}
\usepackage{longtable}

% Custom macros for Linear Algebra, Optimization, and Machine Learning
\DeclareMathOperator*{\argmin}{arg\,min}
\DeclareMathOperator*{\argmax}{arg\,max}
\DeclareMathOperator{\trace}{Tr}
\DeclareMathOperator{\rank}{rank}
\DeclareMathOperator{\diag}{diag}
\DeclareMathOperator{\proj}{proj}
\DeclareMathOperator{\Null}{Null}
\DeclareMathOperator{\Col}{Col}
\DeclareMathOperator{\Row}{Row}
\DeclareMathOperator{\Span}{span}
\DeclareMathOperator{\Cov}{Cov}
\DeclareMathOperator{\Var}{Var}
\DeclareMathOperator{\E}{\mathbb{E}}
\DeclareMathOperator{\sgn}{sgn}
\DeclareMathOperator{\prox}{prox}

% Theorem and Mathematical Environments
\theoremstyle{definition}
\newtheorem{theorem}{Theorem}[chapter]
\newtheorem{lemma}[theorem]{Lemma}
\newtheorem{corollary}[theorem]{Corollary}
\newtheorem{proposition}[theorem]{Proposition}

% Breakable Example Box for in-depth worked pedagogical examples
\renewtcolorbox{examplebox}[1]{
    colback=purplebg,
    colframe=purpleborder,
    coltitle=white,
    fonttitle=\bfseries\sffamily,
    title={Example: #1},
    enhanced,
    breakable,
    attach boxed title to top left={yshift=-2mm, xshift=4mm},
    boxed title style={colback=purpleex, colframe=purpleborder, arc=1mm},
    arc=2mm,
    boxrule=0.8mm,
    top=4mm,
    bottom=3mm,
    left=3mm,
    right=3mm
}



\title{
    \Huge \bfseries \sffamily \color{primarydark} Machine Learning Foundations \\
    \large \vspace{0.3cm} \color{tealdef} Mathematical Principles, Algorithmic Geometry \& Rigorous Foundations \\
    \vspace{0.2cm} \normalsize \color{primarydark} Course ID: BSCS2004
}
\author{
    \textbf{IIT Madras BS in Data Science and Applications} \\
    \small Based on Lectures by Dr.~Harish Guruprasad Ramaswamy, Dr.~Arun Rajkumar, and Dr.~Prashanth L.A. \\
    \small Prepared for Mathematical Mastery, Algorithmic Rigor, and Applied Machine Learning
}
\date{November 2026}

\makeindex

\begin{document}

\maketitle

\frontmatter
\tableofcontents
\newpage

\chapter*{Introduction to Machine Learning Foundations}
\addcontentsline{toc}{chapter}{Introduction to Machine Learning Foundations}

Welcome to the \textbf{Machine Learning Foundations (BSCS2004)} Comprehensive Textbook for the IIT Madras BS Degree Programme. Machine Learning Foundations sits at the vital juncture where the foundational mathematics of linear algebra, multivariable calculus, numerical optimization, and probability theory transform into the operative machinery of artificial intelligence.

Rather than treating modern machine learning algorithms as opaque heuristic routines, this textbook demystifies every model from first principles. By formalizing data spaces as vector spaces, objectives as convex energy landscapes, and learning as optimization or probabilistic inference, students develop the mathematical maturity required to diagnose failures, establish convergence rates, and innovate novel algorithmic architectures.

\section*{Course Organization \& Structural Flow}
The course is structured across \textbf{12 Weeks} systematically divided into four core pedagogical pillars:
\begin{itemize}
    \item \textbf{Weeks 1--4 (Linear Algebraic Foundations \& Linear Models):} Problem formulation, empirical risk minimization, univariate and multivariate calculus review, the four fundamental matrix subspaces, orthogonal projections, normal equations, Ordinary Least Squares (OLS), and matrix diagonalization.
    \item \textbf{Weeks 5--7 (Spectral Geometry, SVD \& Dimensionality Reduction):} Complex and Hermitian matrices, the Spectral Theorem, Singular Value Decomposition (SVD), low-rank matrix approximations (Eckart-Young-Mirsky Theorem), positive definite matrices, and Principal Component Analysis (PCA).
    \item \textbf{Weeks 8--10 (Convex Analysis \& Optimization in Machine Learning):} Unconstrained optimization, gradient descent dynamics and step-size bounds, multivariate Taylor expansions, constrained optimization, method of Lagrange multipliers, convex sets, convex functions, duality, and Karush-Kuhn-Tucker (KKT) optimality conditions.
    \item \textbf{Weeks 11--12 (Probability Foundations, Latent Models \& Concentration):} Joint, marginal, and conditional continuous distributions, expectation, exponential family, parameter estimation (MLE/MAP), Gaussian Mixture Models (GMMs), the Expectation-Maximization (EM) algorithm, and concentration inequalities (Markov, Chebyshev, Hoeffding, and Central Limit Theorem).
\end{itemize}

\section*{Pedagogical Anatomy \& Mastery Standard}
Every chapter in this textbook adheres strictly to an authoritative, self-contained nine-part pedagogical architecture:
\begin{itemize}
    \item \textcolor{primarydark}{\textbf{Part 1: Executive Roadmap \& Learning Objectives:}} Prerequisites, core competencies, and formal mathematical notation guide.
    \item \textcolor{tealdef}{\textbf{Part 2: Mathematical \& Computational Foundations:}} Rigorous definitions, governing theorems, lemmas, and mathematical proofs.
    \item \textcolor{primarydark}{\textbf{Part 3: Visual \& Geometric Intuition:}} Native TikZ architecture schematics, geometric subspace projections, and optimization contours.
    \item \textcolor{purpleex}{\textbf{Part 4: Deep Intuition \& Conceptual Derivations:}} First-principles mathematical derivations, mental models, and empirical reasoning.
    \item \textcolor{amberwarning}{\textbf{Part 5: Explicit Gotchas \& Diagnostic Misconceptions:}} Curated \texttt{trapbox} warnings exposing algebraic errors, geometric fallacies, and numerical pitfalls.
    \item \textcolor{tealdef}{\textbf{Part 6: Comparative Synthesis \& Summary Table:}} Multi-dimensional comparison matrices evaluating competing models, decomposition techniques, and optimization methods.
    \item \textcolor{primarydark}{\textbf{Part 7: Python / Scientific Computing Lab:}} Production-ready, runnable scientific Python implementations using \texttt{numpy}, \texttt{scipy}, and \texttt{matplotlib}.
    \item \textcolor{purpleex}{\textbf{Part 8: Comprehensive Practice Problems:}} Part A multi-part computational scenarios alongside Part B diagnostic multiple-choice questions.
    \item \textcolor{coralformula}{\textbf{Part 9: Detailed Solutions \& Distractor Rationales:}} Exhaustive mathematical derivations and thorough four-option distractor analyses in the Appendix.
\end{itemize}

\chapter*{Machine Learning Foundations --- Quick Reference Cheat Sheet}
\addcontentsline{toc}{chapter}{Machine Learning Foundations --- Quick Reference Cheat Sheet}

\begin{tcolorbox}[breakable, colback=primarylight!30, colframe=primarydark, title={\large \textbf{1. Linear Algebra, Subspaces \& Regularized Regression (Weeks 1, 3--4)}}]
Let $A \in \mathbb{R}^{m \times n}$ with rank $r = \rank(A) \le \min(m, n)$:
\begin{itemize}[leftmargin=*, noitemsep, topsep=2pt]
    \item \textbf{Strang's Four Fundamental Subspaces:}
    \begin{itemize}[leftmargin=*, noitemsep]
        \item Column space: $\Col(A) \subset \mathbb{R}^m$, dimension $\dim(\Col(A)) = r$.
        \item Row space: $\Row(A) \subset \mathbb{R}^n$, dimension $\dim(\Row(A)) = r$.
        \item Nullspace (Kernel): $\Null(A) \subset \mathbb{R}^n$, dimension $\dim(\Null(A)) = n - r$ (Rank-Nullity: $r + \dim(\Null(A)) = n$).
        \item Left nullspace: $\Null(A^T) \subset \mathbb{R}^m$, dimension $\dim(\Null(A^T)) = m - r$.
    \end{itemize}
    \item \textbf{Fundamental Orthogonal Decompositions:}
    \begin{align*}
    \Row(A) &\perp \Null(A) \quad (\mathbb{R}^n = \Row(A) \oplus \Null(A)) \\
    \Col(A) &\perp \Null(A^T) \quad (\mathbb{R}^m = \Col(A) \oplus \Null(A^T))
    \end{align*}
    \item \textbf{Orthogonal Projection Operators:}
    \begin{itemize}[leftmargin=*, noitemsep]
        \item Onto a unit vector $\mathbf{u}$: $P_{\mathbf{u}} = \mathbf{u}\mathbf{u}^T$.
        \item Onto $\Col(A)$ ($r = n$): $P = A(A^T A)^{-1}A^T$. Onto $\Null(A^T)$: $P_\perp = I_m - P$.
        \item Properties: Idempotent ($P^2 = P$), Symmetric ($P^T = P$), Eigenvalues $\lambda \in \{0, 1\}$, $\trace(P) = r$.
    \end{itemize}
    \item \textbf{Ordinary Least Squares (OLS) \& Normal Equations:}
    \begin{equation*}
    A^T A \mathbf{w}^* = A^T \mathbf{y} \implies \mathbf{w}^* = (A^T A)^{-1} A^T \mathbf{y}, \qquad \hat{\mathbf{y}} = P\mathbf{y}, \qquad \mathbf{e} = \mathbf{y} - \hat{\mathbf{y}} \in \Null(A^T)
    \end{equation*}
    \item \textbf{QR Factorization for Stable OLS:} $A = QR$ ($Q^T Q = I_n$, $R$ upper triangular). Solves $R \mathbf{w}^* = Q^T \mathbf{y}$ via back-substitution without forming $A^T A$, avoiding condition squaring.
    \item \textbf{Ridge Regression (Tikhonov / $\ell_2$ Regularization):}
    \begin{equation*}
    \mathbf{w}_{\text{ridge}}^* = (X^T X + \lambda I)^{-1} X^T \mathbf{y} = \sum_{i=1}^r \frac{\sigma_i^2}{\sigma_i^2 + \lambda} \frac{\mathbf{u}_i^T \mathbf{y}}{\sigma_i} \mathbf{v}_i
    \end{equation*}
    Guarantees invertibility because $(X^T X + \lambda I) \succ 0$ for $\lambda > 0$; strictly shrinks singular directions.
    \item \textbf{Lasso Regression ($\ell_1$ Regularization):} $\min_{\mathbf{w}} \frac{1}{2}\|X\mathbf{w} - \mathbf{y}\|_2^2 + \lambda \|\mathbf{w}\|_1$. Produces exact coordinate sparsity; 1D closed form via Soft-Thresholding:
    \begin{equation*}
    S_\lambda(z) \triangleq \operatorname{sign}(z)\max(|z| - \lambda, 0)
    \end{equation*}
    \item \textbf{Bias-Variance Decomposition:} $\operatorname{MSE}(\hat{\mathbf{w}}) = \operatorname{Bias}(\hat{\mathbf{w}})^2 + \operatorname{Var}(\hat{\mathbf{w}})$. Ridge introduces small bias to slash variance $\sigma^2 \trace\big((X^T X + \lambda I)^{-1} X^T X (X^T X + \lambda I)^{-1}\big)$.
    \item \textbf{Vector Norms \& Dualities:}
    \begin{itemize}[leftmargin=*, noitemsep]
        \item $\ell_1$: $\|\mathbf{x}\|_1 = \sum |x_i|$ (sparsity); $\ell_2$: $\|\mathbf{x}\|_2 = \sqrt{\sum x_i^2}$; $\ell_\infty$: $\|\mathbf{x}\|_\infty = \max |x_i|$.
        \item Cauchy-Schwarz: $|\mathbf{x}^T \mathbf{y}| \le \|\mathbf{x}\|_2 \|\mathbf{y}\|_2$; Hölder: $|\mathbf{x}^T \mathbf{y}| \le \|\mathbf{x}\|_p \|\mathbf{y}\|_q$ ($\frac{1}{p} + \frac{1}{q} = 1$).
    \end{itemize}
    \item \textbf{Matrix Norms \& Sensitivity:}
    \begin{itemize}[leftmargin=*, noitemsep]
        \item Spectral Norm: $\|A\|_2 \triangleq \sup_{\mathbf{x} \ne \mathbf{0}} \frac{\|A\mathbf{x}\|_2}{\|\mathbf{x}\|_2} = \sigma_{\max}(A) = \sigma_1$.
        \item Frobenius Norm: $\|A\|_F \triangleq \sqrt{\sum_{i,j} A_{ij}^2} = \sqrt{\trace(A^T A)} = \sqrt{\sum_{i=1}^r \sigma_i^2}$.
        \item Nuclear Norm: $\|A\|_* \triangleq \sum_{i=1}^r \sigma_i(A) = \trace(\sqrt{A^T A})$ (convex envelope of rank).
        \item Condition Number: $\kappa(A) \triangleq \|A\|_2 \|A^{-1}\|_2 = \frac{\sigma_{\max}(A)}{\sigma_{\min}(A)}$. Error: $\frac{\|\delta \mathbf{x}\|_2}{\|\mathbf{x}\|_2} \le \kappa(A) \frac{\|\delta \mathbf{b}\|_2}{\|\mathbf{b}\|_2}$.
        \item \textbf{Condition Squaring Disaster:} $\kappa(A^T A) = (\kappa(A))^2$. Forming normal equations squares condition number!
    \end{itemize}
\end{itemize}
\end{tcolorbox}

\begin{tcolorbox}[breakable, colback=tealbg!30, colframe=tealborder, title={\large \textbf{2. Spectral Theory, Complex Operators \& Rayleigh Quotient (Weeks 4--5)}}]
\begin{itemize}[leftmargin=*, noitemsep, topsep=2pt]
    \item \textbf{Eigenvalues \& Diagonalization:} $A \mathbf{v} = \lambda \mathbf{v}$ ($\mathbf{v} \ne \mathbf{0}$). Characteristic poly: $\det(A - \lambda I) = 0$.
    \begin{itemize}[leftmargin=*, noitemsep]
        \item Trace: $\trace(A) = \sum_{i=1}^n \lambda_i = \sum_{i=1}^n A_{ii}$, and Determinant: $\det(A) = \prod_{i=1}^n \lambda_i$.
        \item Algebraic vs. Geometric Multiplicity: $1 \le \text{g.m.}(\lambda) \le \text{a.m.}(\lambda)$. Diagonalizable $\iff \text{g.m.}(\lambda) = \text{a.m.}(\lambda)$ for all $\lambda$.
        \item Similarity Invariant: $B = P^{-1} A P \implies \det(B) = \det(A), \trace(B) = \trace(A)$, identical spectrum.
    \end{itemize}
    \item \textbf{Complex Inner Product Spaces:}
    \begin{itemize}[leftmargin=*, noitemsep]
        \item Inner product: $\langle \mathbf{u}, \mathbf{v} \rangle = \mathbf{u}^H \mathbf{v} = \sum \bar{u}_i v_i$. Conjugate transpose: $A^H = \overline{A^T}$.
        \item \textbf{Hermitian ($A^H = A$):} All eigenvalues are strictly real ($\lambda_i \in \mathbb{R}$); eigenvectors of distinct eigenvalues are mutually orthogonal.
        \item \textbf{Unitary ($U^H U = I$):} Preserves Euclidean length and angles ($\|U\mathbf{x}\|_2 = \|\mathbf{x}\|_2$); eigenvalues satisfy $|\lambda_i| = 1$.
        \item \textbf{Schur Decomposition:} For any square $A \in \mathbb{C}^{n \times n}$, there exists unitary $Q$ such that $A = Q T Q^H$ where $T$ is upper triangular.
    \end{itemize}
    \item \textbf{The Spectral Theorem (Real Symmetric $A = A^T$ or Hermitian $A = A^H$):}
    \begin{equation*}
    A = Q \Lambda Q^T = \sum_{i=1}^n \lambda_i \mathbf{q}_i \mathbf{q}_i^T, \qquad Q^T Q = I_n, \quad \lambda_i \in \mathbb{R}, \quad \mathbf{q}_i^T \mathbf{q}_j = \delta_{ij}
    \end{equation*}
    \item \textbf{Rayleigh Quotient \& Variational Bounds:}
    \begin{equation*}
    R_A(\mathbf{x}) \triangleq \frac{\mathbf{x}^T A \mathbf{x}}{\mathbf{x}^T \mathbf{x}} \implies \lambda_{\min}(A) \le R_A(\mathbf{x}) \le \lambda_{\max}(A) \quad (\forall \mathbf{x} \ne \mathbf{0})
    \end{equation*}
    $\argmax_{\|\mathbf{x}\|=1} R_A(\mathbf{x}) = \mathbf{q}_1$ (dominant eigenvector); $\argmin_{\|\mathbf{x}\|=1} R_A(\mathbf{x}) = \mathbf{q}_n$.
    \item \textbf{Courant-Fischer Min-Max Theorem:} Variational characterization of the $k$-th eigenvalue:
    \begin{equation*}
    \lambda_k(A) = \max_{\dim(V) = k} \min_{\mathbf{x} \in V, \mathbf{x} \ne \mathbf{0}} \frac{\mathbf{x}^T A \mathbf{x}}{\mathbf{x}^T \mathbf{x}} = \min_{\dim(V) = n - k + 1} \max_{\mathbf{x} \in V, \mathbf{x} \ne \mathbf{0}} \frac{\mathbf{x}^T A \mathbf{x}}{\mathbf{x}^T \mathbf{x}}
    \end{equation*}
    \item \textbf{Positive Definite Matrices ($A \succ 0$):} $\mathbf{x}^T A \mathbf{x} > 0 \; (\forall \mathbf{x} \ne \mathbf{0})$.
    \textbf{Four Equivalent Tests:} (1) All $\lambda_i > 0$; (2) Sylvester: all leading principal minors $\Delta_k > 0$; (3) Cholesky: unique $A = L L^T$ ($l_{ii} > 0$); (4) Energy landscape is a strictly convex bowl.
\end{itemize}
\end{tcolorbox}

\begin{tcolorbox}[breakable, colback=coralbg!30, colframe=coralborder, title={\large \textbf{3. SVD, Pseudoinverse \& Principal Component Analysis (Weeks 5--7)}}]
\begin{itemize}[leftmargin=*, noitemsep, topsep=2pt]
    \item \textbf{Singular Value Decomposition (SVD):}
    For any rectangular $A \in \mathbb{R}^{m \times n}$ with $\rank(A) = r$:
    \begin{equation*}
    A = U \Sigma V^T = \sum_{i=1}^r \sigma_i \mathbf{u}_i \mathbf{v}_i^T, \qquad U \in \mathbb{R}^{m \times m}, \; \Sigma \in \mathbb{R}^{m \times n}, \; V \in \mathbb{R}^{n \times n}
    \end{equation*}
    \begin{itemize}[leftmargin=*, noitemsep]
        \item Singular values: $\sigma_i = \sqrt{\lambda_i(A^T A)} \ge 0$, sorted: $\sigma_1 \ge \sigma_2 \ge \dots \ge \sigma_r > 0$.
        \item Right singular vectors $V$: Orthonormal eigenvectors of $A^T A$ ($A^T A \mathbf{v}_i = \sigma_i^2 \mathbf{v}_i$).
        \item Left singular vectors $U$: Orthonormal eigenvectors of $A A^T$ ($\mathbf{u}_i = \frac{1}{\sigma_i} A \mathbf{v}_i$).
    \end{itemize}
    \item \textbf{Eckart-Young-Mirsky Theorem (Optimal Low-Rank Approximation):}
    For $k < r$, the truncated SVD $A_k \triangleq \sum_{i=1}^k \sigma_i \mathbf{u}_i \mathbf{v}_i^T$ uniquely minimizes rank-$k$ approximation error:
    \begin{equation*}
    \min_{\rank(B) \le k} \|A - B\|_2 = \|A - A_k\|_2 = \sigma_{k+1}, \qquad \min_{\rank(B) \le k} \|A - B\|_F = \sqrt{\sum_{i=k+1}^r \sigma_i^2}
    \end{equation*}
    \item \textbf{Polar Decomposition:} $A = Q P$ where $Q = U V^T$ ($Q^T Q = I$) and $P = V \Sigma V^T = \sqrt{A^T A} \succeq 0$.
    \item \textbf{Moore-Penrose Pseudoinverse:} $A^+ = V \Sigma^+ U^T = \sum_{i=1}^r \frac{1}{\sigma_i} \mathbf{v}_i \mathbf{u}_i^T$.
    \begin{itemize}[leftmargin=*, noitemsep]
        \item Four Axioms: (1) $A A^+ A = A$, (2) $A^+ A A^+ = A^+$, (3) $(A A^+)^T = A A^+$, (4) $(A^+ A)^T = A^+ A$.
        \item Minimum-Norm OLS: $\mathbf{x}_{\text{LS}} = A^+ \mathbf{b}$ satisfies $\mathbf{x}_{\text{LS}} \in \Row(A)$ with $\|\mathbf{x}_{\text{LS}}\|_2 < \|\mathbf{x}\|_2$ for any other minimizer.
    \end{itemize}
    \item \textbf{Principal Component Analysis (PCA):}
    \begin{itemize}[leftmargin=*, noitemsep]
        \item Centering Data: $\tilde{X} = X - \mathbf{1}\bar{\mathbf{x}}^T$. Sample Covariance Matrix: $S = \frac{1}{N-1}\tilde{X}^T \tilde{X} \in \mathbb{R}^{d \times d}$.
        \item Covariance Eigendecomposition: $S \mathbf{u}_i = \lambda_i \mathbf{u}_i$ with $\Var(\tilde{X}\mathbf{u}_i) = \lambda_i$.
        \item SVD Connection: If $\tilde{X} = U \Sigma V^T$, then $S = \frac{1}{N-1} V \Sigma^2 V^T \implies \mathbf{u}_i = \mathbf{v}_i$ and $\lambda_i = \frac{\sigma_i^2}{N-1}$.
        \item Principal Component Scores: $Z = \tilde{X} V_k \in \mathbb{R}^{N \times k}$ with uncorrelated columns: $\frac{1}{N-1}Z^T Z = \Lambda_k$.
        \item Variance Explained \& Scree Ratio: $\eta_k = \frac{\sum_{i=1}^k \lambda_i}{\sum_{j=1}^d \lambda_j}$; Reconstruction Error: $\|X - \hat{X}_k\|_F^2 = \sum_{j=k+1}^d \lambda_j$.
        \item \textbf{Dual Gram Shortcut ($N \ll d$):} Compute $K = \tilde{X}\tilde{X}^T \in \mathbb{R}^{N \times N}$. If $K \mathbf{a}_i = \sigma_i^2 \mathbf{a}_i$, the high-dimensional loading vector is $\mathbf{v}_i = \frac{1}{\sigma_i} \tilde{X}^T \mathbf{a}_i$ (saving $\mathcal{O}(d^3)$ computations).
    \end{itemize}
\end{itemize}
\end{tcolorbox}

\begin{tcolorbox}[breakable, colback=purplebg!30, colframe=purpleborder, title={\large \textbf{4. Multivariate Calculus \& Numerical Optimization (Weeks 2 \& 8)}}]
\begin{itemize}[leftmargin=*, noitemsep, topsep=2pt]
    \item \textbf{Gradient, Hessian \& Directional Derivative:}
    \begin{equation*}
    \nabla f(\mathbf{x}) = \left[\frac{\partial f}{\partial x_1}, \dots, \frac{\partial f}{\partial x_d}\right]^T \in \mathbb{R}^d, \qquad H(\mathbf{x}) = \nabla^2 f(\mathbf{x}) \in \mathbb{R}^{d \times d}, \qquad D_{\mathbf{v}} f(\mathbf{x}) = \nabla f(\mathbf{x})^T \mathbf{v}
    \end{equation*}
    $\nabla f(\mathbf{x})$ points strictly in the direction of \textbf{steepest ascent} and is orthogonal to level curves.
    \item \textbf{Second-Order Taylor Expansion:}
    \begin{equation*}
    f(\mathbf{x} + \mathbf{v}) \approx f(\mathbf{x}) + \nabla f(\mathbf{x})^T \mathbf{v} + \frac{1}{2} \mathbf{v}^T \nabla^2 f(\mathbf{x}) \mathbf{v}
    \end{equation*}
    Critical point classification ($\nabla f(\mathbf{x}^*) = \mathbf{0}$): Minimum if $\nabla^2 f \succ 0$; Maximum if $\nabla^2 f \prec 0$; Saddle if indefinite.
    \item \textbf{Convexity Definitions \& Tests:}
    \begin{itemize}[leftmargin=*, noitemsep]
        \item Jensen's Inequality: $f(\theta \mathbf{x} + (1-\theta)\mathbf{y}) \le \theta f(\mathbf{x}) + (1-\theta) f(\mathbf{y})$ for $\theta \in [0, 1]$.
        \item First-Order: $f(\mathbf{y}) \ge f(\mathbf{x}) + \nabla f(\mathbf{x})^T(\mathbf{y} - \mathbf{x})$ (supporting hyperplanes).
        \item Second-Order: $f$ is convex $\iff \nabla^2 f(\mathbf{x}) \succeq 0$ (positive semidefinite Hessian).
    \end{itemize}
    \item \textbf{Descent Lemma ($L$-Smoothness):} If $\|\nabla f(\mathbf{x}) - \nabla f(\mathbf{y})\|_2 \le L \|\mathbf{x} - \mathbf{y}\|_2$:
    \begin{equation*}
    f(\mathbf{y}) \le f(\mathbf{x}) + \nabla f(\mathbf{x})^T(\mathbf{y} - \mathbf{x}) + \frac{L}{2}\|\mathbf{y} - \mathbf{x}\|_2^2
    \end{equation*}
    \item \textbf{Gradient Descent (GD) Dynamics \& Rates:} $\mathbf{w}_{k+1} = \mathbf{w}_k - \alpha \nabla f(\mathbf{w}_k)$.
    \begin{itemize}[leftmargin=*, noitemsep]
        \item Stability: Decreases monotonically if $\alpha < \frac{2}{L}$. Optimal step: $\alpha = \frac{1}{L}$.
        \item Convex $L$-smooth: Sublinear rate $f(\mathbf{w}_k) - f^* \le \frac{L \|\mathbf{w}_0 - \mathbf{w}^*\|^2}{2k} = \mathcal{O}(1/k)$.
        \item Strongly convex ($\mu I \preceq \nabla^2 f \preceq L I$): Linear rate $\|\mathbf{w}_k - \mathbf{w}^*\|_2 \le \left(\frac{\kappa - 1}{\kappa + 1}\right)^k \|\mathbf{w}_0 - \mathbf{w}^*\|_2$.
    \end{itemize}
    \item \textbf{Accelerated Momentum \& Modern Optimizers:}
    \begin{itemize}[leftmargin=*, noitemsep]
        \item Polyak Heavy-Ball: $\mathbf{v}_{k+1} = \beta \mathbf{v}_k + \alpha \nabla f(\mathbf{w}_k)$, $\mathbf{w}_{k+1} = \mathbf{w}_k - \mathbf{v}_{k+1}$.
        \item Nesterov Accelerated Gradient (NAG): Evaluates $\nabla f(\mathbf{w}_k - \beta \mathbf{v}_k)$, achieving optimal $\mathcal{O}(1/k^2)$ convex rate.
        \item Adam: Updates via $\mathbf{w}_{t+1} = \mathbf{w}_t - \frac{\eta}{\sqrt{\hat{v}_t} + \epsilon}\hat{m}_t$ with bias-corrected moments $\hat{m}_t, \hat{v}_t$.
        \item AdamW: Decouples weight decay: $\mathbf{w}_{t+1} = (1 - \eta \lambda)\mathbf{w}_t - \frac{\eta}{\sqrt{\hat{v}_t} + \epsilon}\hat{m}_t$.
        \item Newton-Raphson: $\mathbf{w}_{k+1} = \mathbf{w}_k - [H(\mathbf{w}_k)]^{-1}\nabla f(\mathbf{w}_k)$ (quadratic rate $\mathcal{O}(\epsilon^2)$ near minimum).
    \end{itemize}
\end{itemize}
\end{tcolorbox}

\begin{tcolorbox}[breakable, colback=primarylight!30, colframe=primarydark, title={\large \textbf{5. Constrained Optimization, KKT \& Support Vector Machines (Weeks 9--10)}}]
\begin{itemize}[leftmargin=*, noitemsep, topsep=2pt]
    \item \textbf{General Constrained Optimization Problem:}
    \begin{equation*}
    \min_{\mathbf{x}} f_0(\mathbf{x}) \quad \text{subject to} \quad f_i(\mathbf{x}) \le 0 \; (i = 1, \dots, m), \quad h_j(\mathbf{x}) = 0 \; (j = 1, \dots, p)
    \end{equation*}
    \item \textbf{Lagrangian Function \& Duality:}
    \begin{equation*}
    \mathcal{L}(\mathbf{x}, \boldsymbol{\lambda}, \boldsymbol{\nu}) \triangleq f_0(\mathbf{x}) + \sum_{i=1}^m \lambda_i f_i(\mathbf{x}) + \sum_{j=1}^p \nu_j h_j(\mathbf{x}), \qquad g(\boldsymbol{\lambda}, \boldsymbol{\nu}) \triangleq \inf_{\mathbf{x}} \mathcal{L}(\mathbf{x}, \boldsymbol{\lambda}, \boldsymbol{\nu})
    \end{equation*}
    \begin{itemize}[leftmargin=*, noitemsep]
        \item Dual Problem: $\max_{\boldsymbol{\lambda} \succeq \mathbf{0}, \boldsymbol{\nu}} g(\boldsymbol{\lambda}, \boldsymbol{\nu})$. The dual function $g$ is \textbf{always concave}.
        \item Weak Duality: $d^* \le p^*$ always. Strong Duality ($d^* = p^*$): Holds under Slater's Condition ($\exists \mathbf{x} : f_i(\mathbf{x}) < 0$).
    \end{itemize}
    \item \textbf{The Karush-Kuhn-Tucker (KKT) Optimality Conditions:}
    \begin{align*}
    \text{1. Primal Feasibility:} \quad & f_i(\mathbf{x}^*) \le 0 \quad (\forall i), \quad h_j(\mathbf{x}^*) = 0 \quad (\forall j) \\
    \text{2. Dual Feasibility:} \quad & \lambda_i^* \ge 0 \quad (\forall i) \\
    \text{3. Complementary Slackness:} \quad & \lambda_i^* f_i(\mathbf{x}^*) = 0 \quad (\forall i) \implies (\lambda_i^* > 0 \implies f_i(\mathbf{x}^*) = 0) \\
    \text{4. Stationarity:} \quad & \nabla f_0(\mathbf{x}^*) + \sum_{i=1}^m \lambda_i^* \nabla f_i(\mathbf{x}^*) + \sum_{j=1}^p \nu_j^* \nabla h_j(\mathbf{x}^*) = \mathbf{0}
    \end{align*}
    \item \textbf{Support Vector Machines (SVM):}
    \begin{itemize}[leftmargin=*, noitemsep]
        \item Hard-Margin SVM: $\min_{\mathbf{w}, b} \frac{1}{2}\|\mathbf{w}\|_2^2$ s.t. $y_i(\mathbf{w}^T \mathbf{x}_i + b) \ge 1 \; (\forall i)$. Margin: $\gamma = \frac{1}{\|\mathbf{w}\|_2}$.
        \item Soft-Margin SVM: $\min_{\mathbf{w}, b, \boldsymbol{\xi}} \frac{1}{2}\|\mathbf{w}\|_2^2 + C \sum_{i=1}^N \xi_i$ s.t. $y_i(\mathbf{w}^T \mathbf{x}_i + b) \ge 1 - \xi_i, \; \xi_i \ge 0$.
        \item Dual Soft-Margin Formulation:
        \begin{equation*}
        \max_{\boldsymbol{\alpha}} \sum_{i=1}^N \alpha_i - \frac{1}{2}\sum_{i=1}^N \sum_{j=1}^N \alpha_i \alpha_j y_i y_j K(\mathbf{x}_i, \mathbf{x}_j) \quad \text{s.t.} \quad 0 \le \alpha_i \le C, \quad \sum_{i=1}^N \alpha_i y_i = 0
        \end{equation*}
        \item Optimal Weights: $\mathbf{w}^* = \sum_{i=1}^N \alpha_i^* y_i \mathbf{x}_i$. Only Support Vectors ($\alpha_i^* > 0$) define the boundary!
        \item Complementary Slackness: $\alpha_i = 0 \implies$ correctly classified; $0 < \alpha_i < C \implies$ on margin; $\alpha_i = C \implies$ margin violation ($\xi_i > 0$).
    \end{itemize}
    \item \textbf{Kernel Theory \& The Representer Theorem:}
    \begin{itemize}[leftmargin=*, noitemsep]
        \item Mercer's Theorem: $K(\mathbf{x}, \mathbf{z})$ is a valid Mercer kernel $\iff$ Gram matrix $K_{ij} = K(\mathbf{x}_i, \mathbf{x}_j) \succeq 0$.
        \item Kernels: Linear $\mathbf{x}^T \mathbf{z}$; Polynomial $(\mathbf{x}^T \mathbf{z} + c)^d$; RBF/Gaussian $\exp(-\gamma \|\mathbf{x} - \mathbf{z}\|_2^2)$.
        \item Representer Theorem: Any regularized empirical risk minimizer admits $f^*(\mathbf{x}) = \sum_{i=1}^N \alpha_i K(\mathbf{x}_i, \mathbf{x})$.
        \item Kernel Ridge Regression: $\boldsymbol{\alpha}^* = (K + \lambda I)^{-1}\mathbf{y}$ with closed-form predictor $f(\mathbf{x}) = \mathbf{k}(\mathbf{x})^T (K + \lambda I)^{-1}\mathbf{y}$.
    \end{itemize}
\end{itemize}
\end{tcolorbox}

\begin{tcolorbox}[breakable, colback=tealbg!30, colframe=tealborder, title={\large \textbf{6. Probability Foundations, The Exponential Family \& GLMs (Weeks 10--11)}}]
\begin{itemize}[leftmargin=*, noitemsep, topsep=2pt]
    \item \textbf{Multivariate Density Transformations (Jacobian Rule):}
    For invertible $\mathbf{y} = \mathbf{g}(\mathbf{x})$:
    \begin{equation*}
    f_{\mathbf{Y}}(\mathbf{y}) = f_{\mathbf{X}}(\mathbf{g}^{-1}(\mathbf{y})) \cdot \big|\det J_{\mathbf{g}^{-1}}(\mathbf{y})\big| = \frac{f_{\mathbf{X}}(\mathbf{x})}{\big|\det J_{\mathbf{g}}(\mathbf{x})\big|}
    \end{equation*}
    Linear maps $\mathbf{Y} = A\mathbf{X} + \mathbf{b}$: $f_{\mathbf{Y}}(\mathbf{y}) = f_{\mathbf{X}}(A^{-1}(\mathbf{y} - \mathbf{b})) \cdot \frac{1}{|\det A|}$.
    \item \textbf{The Canonical Exponential Family:}
    \begin{equation*}
    p(x; \boldsymbol{\eta}) = h(x) \exp\left( \boldsymbol{\eta}^T \mathbf{T}(x) - A(\boldsymbol{\eta}) \right)
    \end{equation*}
    $\boldsymbol{\eta}$: natural parameter, $\mathbf{T}(x)$: sufficient statistics, $A(\boldsymbol{\eta})$: log-partition function.
    \item \textbf{Moment Properties of $A(\boldsymbol{\eta})$:}
    \begin{equation*}
    \nabla_{\boldsymbol{\eta}} A(\boldsymbol{\eta}) = \E[\mathbf{T}(X)], \qquad \nabla_{\boldsymbol{\eta}}^2 A(\boldsymbol{\eta}) = \Cov(\mathbf{T}(X)) \succeq 0 \implies A(\boldsymbol{\eta}) \text{ is strictly convex}
    \end{equation*}
    \item \textbf{MLE via Moment Matching:} Setting $\nabla \ell(\boldsymbol{\eta}) = \mathbf{0}$ yields:
    \begin{equation*}
    \nabla_{\boldsymbol{\eta}} A(\hat{\boldsymbol{\eta}}_{\text{MLE}}) = \frac{1}{N}\sum_{i=1}^N \mathbf{T}(x_i) \iff \E_{\hat{\boldsymbol{\eta}}}[\mathbf{T}(X)] = \bar{\mathbf{T}}_{\text{emp}}
    \end{equation*}
    \item \textbf{Fisher Information Matrix \& CRLB:}
    \begin{equation*}
    \mathcal{I}(\boldsymbol{\eta}) \triangleq \E\left[ (\nabla \log p)(\nabla \log p)^T \right] = \nabla_{\boldsymbol{\eta}}^2 A(\boldsymbol{\eta}) = \Cov(\mathbf{T}(X)), \qquad \Cov(\hat{\boldsymbol{\eta}}) \succeq \frac{1}{N}\mathcal{I}(\boldsymbol{\eta})^{-1}
    \end{equation*}
    \item \textbf{Universal GLM Gradient:} With $\eta_i = \mathbf{w}^T \mathbf{x}_i$ and $\mu_i = A'(\eta_i)$:
    \begin{equation*}
    \nabla_{\mathbf{w}} \ell(\mathbf{w}) = \sum_{i=1}^N (y_i - \mu_i)\mathbf{x}_i \quad \iff \quad \text{Prediction Error} \times \text{Input Feature}
    \end{equation*}
    \item \textbf{Information Theory \& Cross-Entropy Loss:}
    \begin{equation*}
    H(P) \triangleq -\sum P(x) \log P(x), \qquad H(P, Q) \triangleq -\sum P(x) \log Q(x) = H(P) + D_{\text{KL}}(P \parallel Q)
    \end{equation*}
    Minimizing cross-entropy directly minimizes KL divergence to Nature's empirical distribution.
    \item \textbf{KL Divergence \& Bregman Duality:}
    \begin{equation*}
    D_{\text{KL}}\left( p(x; \boldsymbol{\eta}_1) \parallel p(x; \boldsymbol{\eta}_2) \right) = A(\boldsymbol{\eta}_2) - A(\boldsymbol{\eta}_1) - (\boldsymbol{\eta}_2 - \boldsymbol{\eta}_1)^T \nabla A(\boldsymbol{\eta}_1) = B_A(\boldsymbol{\eta}_2 \parallel \boldsymbol{\eta}_1)
    \end{equation*}
    \item \textbf{Canonical Conjugate Prior:} $p(\boldsymbol{\eta}; \boldsymbol{\chi}, \nu) \propto \exp(\boldsymbol{\eta}^T \boldsymbol{\chi} - \nu A(\boldsymbol{\eta})) \implies \boldsymbol{\chi}_N = \boldsymbol{\chi}_0 + \sum \mathbf{T}(x_i), \; \nu_N = \nu_0 + N$.
\end{itemize}
\end{tcolorbox}

\begin{tcolorbox}[breakable, colback=coralbg!30, colframe=coralborder, title={\large \textbf{7. GMMs, Expectation-Maximization \& Concentration Inequalities (Week 12)}}]
\begin{itemize}[leftmargin=*, noitemsep, topsep=2pt]
    \item \textbf{Multivariate Gaussian Distribution:}
    \begin{equation*}
    \mathcal{N}(\mathbf{x}; \boldsymbol{\mu}, \Sigma) = \frac{1}{(2\pi)^{d/2} |\Sigma|^{1/2}} \exp\left( -\frac{1}{2}(\mathbf{x} - \boldsymbol{\mu})^T \Sigma^{-1} (\mathbf{x} - \boldsymbol{\mu}) \right)
    \end{equation*}
    \begin{itemize}[leftmargin=*, noitemsep]
        \item Mahalanobis Distance: $D_M(\mathbf{x}, \boldsymbol{\mu}) = \sqrt{(\mathbf{x} - \boldsymbol{\mu})^T \Sigma^{-1} (\mathbf{x} - \boldsymbol{\mu})}$. Contours align with eigenvectors $\mathbf{q}_i$.
        \item Sample Covariance Bias: $\hat{\Sigma}_{\text{MLE}} = \frac{1}{N}\sum (\mathbf{x}_i - \bar{\mathbf{x}})(\mathbf{x}_i - \bar{\mathbf{x}})^T \implies \E[\hat{\Sigma}_{\text{MLE}}] = \frac{N-1}{N}\Sigma$. Unbiased: $S = \frac{N}{N-1}\hat{\Sigma}_{\text{MLE}}$.
    \end{itemize}
    \item \textbf{Gaussian Mixture Models (GMM) \& The EM Algorithm:}
    $p(\mathbf{x}) = \sum_{k=1}^K \pi_k \mathcal{N}(\mathbf{x}; \boldsymbol{\mu}_k, \Sigma_k)$.
    \begin{itemize}[leftmargin=*, noitemsep]
        \item \textbf{E-Step:} $\gamma_{ik} \triangleq P(z_i = k \mid \mathbf{x}_i, \boldsymbol{\theta}^{(t)}) = \frac{\pi_k^{(t)} \mathcal{N}(\mathbf{x}_i; \boldsymbol{\mu}_k^{(t)}, \Sigma_k^{(t)})}{\sum_{j=1}^K \pi_j^{(t)} \mathcal{N}(\mathbf{x}_i; \boldsymbol{\mu}_j^{(t)}, \Sigma_j^{(t)})}$, and $N_k \triangleq \sum_{i=1}^N \gamma_{ik}$.
        \item \textbf{M-Step:} $\pi_k = \frac{N_k}{N}$, $\boldsymbol{\mu}_k = \frac{1}{N_k}\sum \gamma_{ik}\mathbf{x}_i$, $\Sigma_k = \frac{1}{N_k}\sum \gamma_{ik}(\mathbf{x}_i - \boldsymbol{\mu}_k)(\mathbf{x}_i - \boldsymbol{\mu}_k)^T$.
        \item \textbf{ELBO-KL Decomposition:} $\log p(\mathbf{X}; \boldsymbol{\theta}) = \text{ELBO}(q, \boldsymbol{\theta}) + D_{\text{KL}}\left( q(\mathbf{Z}) \parallel p(\mathbf{Z} \mid \mathbf{X}; \boldsymbol{\theta}) \right)$.
        \item \textbf{Monotonic Ascent Guarantee:} $\ell(\boldsymbol{\theta}^{(t+1)}) \ge \ell(\boldsymbol{\theta}^{(t)})$.
    \end{itemize}
    \item \textbf{High-Dimensional Gaussian Geometry (The Orange Peel Paradox):}
    For $\mathbf{X} \sim \mathcal{N}(\mathbf{0}, I_d)$, $R = \|\mathbf{X}\|_2 \approx \sqrt{d} \pm \frac{1}{\sqrt{2}}$. Mass concentrates in a thin spherical shell; the mode at the origin is an empty desert!
    \item \textbf{Hierarchy of Non-Asymptotic Concentration Inequalities:}
    \begin{itemize}[leftmargin=*, noitemsep]
        \item \textbf{Markov ($X \ge 0$):} $P(X \ge a) \le \frac{\E[X]}{a}$ for all $a > 0$ ($\mathcal{O}(1/a)$ decay).
        \item \textbf{Chebyshev:} $P(|X - \mu| \ge \epsilon) \le \frac{\sigma^2}{\epsilon^2}$ ($\mathcal{O}(1/\epsilon^2)$ decay; Sample size: $N \ge \frac{\sigma^2}{\epsilon^2 \delta}$).
        \item \textbf{Hoeffding (Bounded $X_i \in [a_i, b_i]$):}
        \begin{equation*}
        P\left( |\bar{X}_N - \mu| \ge \epsilon \right) \le 2 \exp\left( -\frac{2 N^2 \epsilon^2}{\sum_{i=1}^N (b_i - a_i)^2} \right) \implies N \ge \frac{(b-a)^2 \ln(2/\delta)}{2\epsilon^2}
        \end{equation*}
        \item \textbf{Bernstein (Variance-Aware):} For independent zero-mean $X_i$ with $|X_i| \le c$ and variance $\sigma^2$:
        \begin{equation*}
        P\left( \bar{X}_N - \mu \ge \epsilon \right) \le \exp\left( -\frac{N\epsilon^2}{2\sigma^2 + \frac{2}{3}c\epsilon} \right) \implies \text{substantially tighter when } \sigma^2 \ll \frac{(b-a)^2}{4}
        \end{equation*}
        \item \textbf{McDiarmid (Bounded Differences for Functions $f(X_1, \dots, X_N)$):}
        If $|f - f'_i| \le c_i$:
        \begin{equation*}
        P\left( f(\mathbf{X}) - \E[f(\mathbf{X})] \ge \epsilon \right) \le \exp\left( -\frac{2\epsilon^2}{\sum_{i=1}^N c_i^2} \right)
        \end{equation*}
    \end{itemize}
\end{itemize}
\end{tcolorbox}


\mainmatter

% Include all 12 weekly chapters
\chapter[Foundations of Machine Learning]{Foundations of Machine Learning \& Problem Formulation}
\label{chap:week01}

% ----------------------------------------------------------------------
% PART 1: EXECUTIVE ROADMAP & LEARNING OBJECTIVES
% ----------------------------------------------------------------------
\section{Executive Roadmap \& Learning Objectives}

\subsection{The Big Picture: Why Can't We Just Write the Rules?}
Every journey into machine learning begins with a simple, provocative question: \emph{Why do we need a new computational paradigm at all?} For more than half a century, traditional computer science solved the world's most demanding computational problems—from calculating spacecraft trajectories to managing global banking ledgers—using a single, elegant philosophical principle known as **deductive programming**:
\begin{equation}
\boxed{\text{Human Expert Writes Rules (Logic)} + \text{Computer Receives Input Data} \xrightarrow{\text{Deductive Execution}} \text{Answers (Output)}}
\label{eq:deductive_paradigm}
\end{equation}
In this classical paradigm, a human software engineer analyzes the problem, formalizes the complete set of operational rules, and writes deterministic code (using loops, conditionals, and mathematical formulas). If you want to compute income tax, you look up the tax brackets codified by law, translate them into a sequence of \texttt{if-else} statements, and run the program. The computer executes your rules faithfully, deterministically, and with microsecond precision.

\subsubsection{Where Traditional Logic Breaks Down: The Perception Paradox}
Now consider a deceptively modest task: write a computer program to determine whether a small $28 \times 28$ grayscale photograph contains the handwritten digit \textbf{3}.

A human child performs this recognition in a fraction of a second, effortlessly identifying the digit whether it is scribbled with a pencil, painted with a brush, slanted backward, or partially faded. But what does the computer actually see? The computer does not see a digit; it sees a grid of $28 \times 28 = 784$ numbers, where each number represents the brightness of a single pixel between $0$ (black) and $1$ (white). In mathematical terms, the input is a vector living in a 784-dimensional space:
\begin{equation}
\mathbf{x} = \begin{bmatrix} x_1 \\ x_2 \\ \vdots \\ x_{784} \end{bmatrix} \in [0, 1]^{784}
\end{equation}
How would you write classical code to recognize a ``3''?
\begin{itemize}
    \item You might try checking if certain pixels are dark: \emph{``If pixel 200 is dark and pixel 350 is dark, it might be a 3.''} But if the writer slants their handwriting slightly to the left, those exact pixels are now completely white!
    \item You might try writing geometric rules: \emph{``A 3 consists of two semi-circular loops open on the left, connected in the center.''} But how do you mathematically define a ``rough, hand-drawn semi-circular loop'' across a discrete pixel grid subject to ink splatters, varying stroke thickness, and broken strokes?
    \item Even if each pixel were strictly binary (black or white), the number of possible $28 \times 28$ images is:
    \begin{equation}
    2^{784} \approx 10^{236}
    \end{equation}
    To put this incomprehensible number in perspective, physicists estimate that the entire observable universe contains roughly $10^{80}$ subatomic particles. The configuration space of handwritten digits is larger than the physical universe by a factor of $10^{156}$!
\end{itemize}
No team of human engineers, given a century of effort, could write enough \texttt{if-else} statements to handle the combinatorial diversity of human handwriting. Deductive programming collapses whenever the rules governing the problem are too subtle, too complex, or too subconscious for human programmers to explicitly articulate.

\subsubsection{The Inductive Leap: Inverting the Computation}
**Machine Learning** inverts the computational pipeline into an **inductive paradigm**:
\begin{equation}
\boxed{\text{Input Observations} + \text{Desired Outputs (Targets)} \xrightarrow{\text{Inductive Learning}} \text{Discovered Rules (Hypothesis } h\text{)}}
\label{eq:inductive_paradigm}
\end{equation}
Instead of hand-coding the rules, we show the computer 60,000 examples of handwritten digits along with their correct labels. We define an optimization landscape (a mathematical surface measuring how far the computer's guesses are from the truth), and an algorithm systematically adjusts the parameters of a flexible mathematical function until the function automatically produces the correct labels.

In 1959, artificial intelligence pioneer Arthur Samuel captured this essence in words that remain the definitive layman's definition:
\begin{quote}
\emph{``Machine learning is the field of study that gives computers the ability to learn without being explicitly programmed.''}
\end{quote}
Four decades later, Tom Mitchell (1997) refined this into the standard engineering formulation that guides modern data science:
\begin{quote}
\emph{``A computer program is said to learn from experience $E$ with respect to some class of tasks $T$ and performance measure $P$, if its performance at tasks in $T$, as measured by $P$, improves with experience $E$.''}
\end{quote}
In our digit recognition problem:
\begin{itemize}
    \item \textbf{Task $T$:} Classify unseen $28 \times 28$ pixel images into categories $0$ through $9$.
    \item \textbf{Experience $E$:} A historical database of labeled handwritten images $S = \{(\mathbf{x}_i, y_i)\}_{i=1}^N$.
    \item \textbf{Performance Measure $P$:} The classification accuracy (percentage of correct predictions) evaluated on fresh, unseen images that the computer has never encountered before.
\end{itemize}

\subsection{Prerequisites \& Scaffolding}
To master this course, students do not need encyclopedic recall of obscure mathematics, but they do need intuitive fluency across three core collegiate subjects, mapped in Table~\ref{tab:w01_prereqs}.

\begin{table}[htbp]
\centering
\small
\begin{tabularx}{\textwidth}{lXX}
\toprule
\textbf{Subject} & \textbf{Key Mental Model Required} & \textbf{How Machine Learning Uses It} \\
\midrule
\textbf{Linear Algebra} & Vectors as points in space; dot products as directional similarity; matrices as space-transforming operators. & Every dataset is a matrix $X$; every observation is a vector $\mathbf{x} \in \mathbb{R}^d$; linear models draw flat hyperplanes through feature space. \\
\addlinespace
\textbf{Multivariable Calculus} & Derivatives as slopes; gradients $\nabla f$ as directions of steepest ascent; curvature captured by Hessians. & Training a model is an iterative mountain hike: gradients tell algorithms which way is ``downhill'' on the loss surface to find the optimal weights. \\
\addlinespace
\textbf{Probability \& Statistics} & Random variables as uncertain measurements; expectations $\E[X]$ as centers of mass; distributions capturing noise. & Real data is never clean; probability models the intrinsic noise of Nature and allows us to reason about unseen future data. \\
\bottomrule
\end{tabularx}
\caption{Prerequisite Scaffolding: Connecting Foundational Mathematics to Machine Learning Concepts.}
\label{tab:w01_prereqs}
\end{table}

\subsection{Core Competencies}
By the conclusion of this chapter, you will be able to:
\begin{enumerate}
    \item \textbf{Speak the Language of Statistical Learning:} Translate ambiguous business and scientific questions into the formal five-tuple $(\mathcal{X}, \mathcal{Y}, \mathcal{D}, \mathcal{H}, \ell)$.
    \item \textbf{Derive Nature's Theoretical Ceiling:} Prove why the conditional expectation $\E[Y \mid \mathbf{X}]$ is the mathematically best possible regressor under squared loss, and why majority voting is optimal for classification.
    \item \textbf{Dismantle the Overfitting Illusion:} Explain why an algorithm that achieves $100\%$ training accuracy can perform worse than random guessing on new data, and compute the exact Generalization Gap.
    \item \textbf{Diagnose Model Failures via Error Decomposition:} Partition predictive error into Approximation Error (model inflexibility / bias) and Estimation Error (finite data noise / variance).
    \item \textbf{Navigate High-Dimensional Geometry:} Mathematically demonstrate why Euclidean intuition fails in high dimensions, proving the Empty Hypercube theorem and the Distance Concentration phenomenon.
\end{enumerate}

\subsection{Notation \& Mathematical Conventions}
Throughout this textbook, we adhere to the standard mathematical notation detailed in Table~\ref{tab:w01_notation}.

\begin{table}[htbp]
\centering
\small
\begin{tabularx}{\textwidth}{llX}
\toprule
\textbf{Symbol} & \textbf{Formal Name} & \textbf{Intuitive Plain-English Meaning} \\
\midrule
$\mathbf{x} \in \mathcal{X} \subseteq \mathbb{R}^d$ & Feature Vector & A list of $d$ numerical measurements describing an object (e.g., patient age, blood pressure, cholesterol). \\
$y \in \mathcal{Y}$ & Ground-Truth Label & The outcome we want to predict (e.g., house price in $\mathbb{R}$, or whether a transaction is fraudulent in $\{0, 1\}$). \\
$\mathcal{D}$ & Data Distribution & Nature's secret probability distribution that generates real-world data pairs $(\mathbf{x}, y)$. \\
$S = \{(\mathbf{x}_i, y_i)\}_{i=1}^N$ & Training Sample & The limited collection of $N$ examples our algorithm is allowed to study. \\
$\mathcal{H}$ & Hypothesis Class & The ``menu'' of mathematical functions our algorithm is permitted to choose from (e.g., all straight lines). \\
$h(\mathbf{x}) \in \mathcal{H}$ & Candidate Hypothesis & A specific prediction rule chosen from the menu that maps features to predicted targets $\hat{y} = h(\mathbf{x})$. \\
$\ell(\hat{y}, y)$ & Loss Function & The price we pay for making a mistake when we predict $\hat{y}$ but the true answer was $y$. \\
$R(h)$ & True Risk & The average loss the model would suffer on all possible future data generated by Nature (the unseen final exam score). \\
$R_{\text{emp}}(h; S)$ & Empirical Risk & The average loss the model suffered on the observed training sample $S$ (the practice homework score). \\
$h^*$ & Bayes Optimal Rule & The theoretically best predictor possible if we possessed infinite data and knew Nature's distribution perfectly. \\
\bottomrule
\end{tabularx}
\caption{Mathematical Conventions and Learning Framework Notation.}
\label{tab:w01_notation}
\end{table}

% ----------------------------------------------------------------------
% PART 2: MATHEMATICAL & COMPUTATIONAL FOUNDATIONS
% ----------------------------------------------------------------------
\section{Mathematical \& Computational Foundations}

\subsection{How to Translate the World into Mathematics: The 5-Tuple}
To study machine learning mathematically, we must formalize how an algorithm interacts with reality. Every machine learning problem on Earth can be codified as a **five-tuple**:
\begin{equation}
\Big( \mathcal{X}, \; \mathcal{Y}, \; \mathcal{D}, \; \mathcal{H}, \; \ell \Big)
\end{equation}

\subsubsection{1. The Instance Space $\mathcal{X}$}
The instance space represents everything our sensors can measure about the world. An observation is converted into a vector of $d$ real numbers, $\mathbf{x} = [x_1, x_2, \dots, x_d]^T \in \mathbb{R}^d$.
\begin{itemize}
    \item If we are diagnosing diabetes, $\mathbf{x}$ might be a 5-dimensional vector: $[\text{Age}, \text{BMI}, \text{Glucose}, \text{Insulin}, \text{BloodPressure}]^T \in \mathbb{R}^5$.
    \item If we are classifying images, $\mathbf{x}$ is a vector of $784$ pixel intensities in $[0, 1]^{784}$.
\end{itemize}

\subsubsection{2. The Target Space $\mathcal{Y}$}
The target space is the set of all possible correct answers:
\begin{itemize}
    \item **Regression ($\mathcal{Y} \subseteq \mathbb{R}$):** The target is a continuous real number. Examples: predicting house price, vehicle fuel efficiency, or patient blood glucose level.
    \item **Binary Classification ($\mathcal{Y} \in \{0, 1\}$ or $\{-1, +1\}$):** The target is a binary decision. Examples: predicting whether an email is spam ($1$) or inbox ($0$), or whether a tumor is malignant ($+1$) or benign ($-1$).
    \item **Multiclass Classification ($\mathcal{Y} \in \{1, 2, \dots, K\}$):** The target is one of $K$ mutually exclusive categories. Examples: recognizing handwritten digits ($K = 10$) or identifying dog breeds ($K = 120$).
\end{itemize}

\subsubsection{3. Nature's Secret Law: The Data Distribution $\mathcal{D}$}
A common misconception among beginners is that there exists a clean mathematical function $y = f(\mathbf{x})$ that perfectly dictates the output from the inputs. In reality, the world is noisy, ambiguous, and filled with unobserved variables:
\begin{itemize}
    \item Two houses in the same neighborhood with identical square footage and bedrooms might sell for \$320,000 and \$360,000 because one has an updated kitchen that our feature vector did not record.
    \item Two patients with identical age, BMI, and glucose readings might have different diabetes outcomes due to genetic factors our medical test did not measure.
\end{itemize}
Therefore, we model the relationship between $\mathbf{x}$ and $y$ **probabilistically**. Nature possesses a joint probability distribution $\mathcal{D}$ over the product space $\mathcal{X} \times \mathcal{Y}$. When a data point is generated, Nature draws a pair $(\mathbf{x}, y) \sim \mathcal{D}$.
\begin{equation}
P(\mathbf{x}, y) = \underbrace{P(\mathbf{x})}_{\text{Marginal Feature Distribution}} \times \underbrace{P(y \mid \mathbf{x})}_{\text{Conditional Target Distribution}}
\end{equation}
The learner **never sees $\mathcal{D}$ directly**. Nature keeps $\mathcal{D}$ locked in a black box. The learner only sees a finite collection of random draws: the training sample $S = \{(\mathbf{x}_i, y_i)\}_{i=1}^N \stackrel{\text{i.i.d.}}{\sim} \mathcal{D}$.

\subsubsection{4. The Hypothesis Class $\mathcal{H}$ and Inductive Bias}
Why can't an algorithm simply consider \emph{every possible function} mapping $\mathcal{X} \to \mathcal{Y}$?

Consider Hume's classic philosophical riddle: \emph{If you see five black ravens in a row, what color is the sixth raven?} If you make zero assumptions about how the universe works, seeing five black ravens tells you strictly zero about the next raven. It could be white, purple, or an armchair. To conclude that the sixth raven is likely black, your brain must make an unprovable prior assumption: that Nature possesses continuity, smoothness, and patterns that generalize across instances.

In machine learning, this prior assumption is called **Inductive Bias**. An algorithm that considers every arbitrary mathematical function cannot learn; given $N$ points, it could fit a function that equals the correct answers on those $N$ points, but outputs completely random numbers everywhere else!

To make learning possible, we restrict our search to a specific family of functions $\mathcal{H}$, called the **Hypothesis Class**:
\begin{itemize}
    \item \textbf{Linear Models:} $\mathcal{H}_{\text{linear}} = \{ h(\mathbf{x}) = \mathbf{w}^T \mathbf{x} + b \mid \mathbf{w} \in \mathbb{R}^d, b \in \mathbb{R} \}$. The inductive bias is: \emph{``The relationship between features and target is planar.''}
    \item \textbf{Polynomial Models:} $\mathcal{H}_{\text{poly}, p} = \{ h(x) = \sum_{j=0}^p w_j x^j \}$. The inductive bias is: \emph{``The relationship is a smooth curve of degree at most $p$.''}
\end{itemize}

\subsubsection{5. The Loss Function $\ell(\hat{y}, y)$}
The loss function is the referee. It takes our model's guess $\hat{y} = h(\mathbf{x})$ and the true ground-truth $y$, and outputs a non-negative penalty score $\ell(\hat{y}, y) \ge 0$:
\begin{itemize}
    \item \textbf{Squared Error Loss (L2 Loss):}
    \begin{equation}
    \ell_{\text{sq}}(\hat{y}, y) \triangleq (\hat{y} - y)^2
    \label{eq:l2_loss_def}
    \end{equation}
    Penalizes large mistakes quadratically. If you miss by 2 units, your penalty is 4; if you miss by 10 units, your penalty is 100! This makes squared loss heavily focused on avoiding large outliers.
    \item \textbf{Absolute Error Loss (L1 Loss):}
    \begin{equation}
    \ell_{\text{abs}}(\hat{y}, y) \triangleq |\hat{y} - y|
    \label{eq:l1_loss_def}
    \end{equation}
    Penalizes mistakes linearly. Missing by 10 is penalized exactly 10 times as much as missing by 1. This makes absolute loss robust to outliers.
    \item \textbf{Zero-One Loss (Classification Loss):}
    \begin{equation}
    \ell_{0-1}(\hat{y}, y) \triangleq \mathbb{I}(\hat{y} \ne y) = \begin{cases} 0 & \text{if } \hat{y} = y \\ 1 & \text{if } \hat{y} \ne y \end{cases}
    \label{eq:zero_one_loss_def}
    \end{equation}
    Gives zero penalty for correct predictions and 1 penalty for errors, regardless of confidence.
\end{itemize}

\subsection{The Central Goal: Minimizing True Risk}
Now that we have our 5-tuple, we can state the ultimate goal of machine learning with mathematical precision:

\begin{definitionbox}{True Risk (Generalization Error)}
The \textbf{True Risk} $R(h)$ of a hypothesis $h \in \mathcal{H}$ under loss function $\ell$ is the expected loss over the entire, infinite data-generating distribution $\mathcal{D}$:
\begin{equation}
R(h) \triangleq \E_{(\mathbf{X}, Y) \sim \mathcal{D}} [\ell(h(\mathbf{X}), Y)] = \int_{\mathcal{X} \times \mathcal{Y}} \ell(h(\mathbf{x}), y) \, dP(\mathbf{x}, y)
\label{eq:true_risk_formal}
\end{equation}
\end{definitionbox}

Think of True Risk as a student's grade on the **unseen final examination**. It measures how well the hypothesis will perform on future data drawn from the identical environment.

\subsubsection{The Fundamental Dilemma of Machine Learning}
Look carefully at Equation~\eqref{eq:true_risk_formal}. To compute the integral, you need $dP(\mathbf{x}, y)$—Nature's secret joint probability distribution! But we do not know Nature's distribution! If we knew $\mathcal{D}$, we wouldn't need machine learning in the first place; we could just do probability calculations.

Because we cannot compute the true risk, we do the next best thing: we compute the average error on the historical homework assignments we actually have in our hands.

\begin{definitionbox}{Empirical Risk (Training Error)}
The \textbf{Empirical Risk} $R_{\text{emp}}(h; S)$ is the sample average loss realized by hypothesis $h$ on the concrete training sample $S = \{(\mathbf{x}_i, y_i)\}_{i=1}^N$:
\begin{equation}
R_{\text{emp}}(h; S) \triangleq \frac{1}{N} \sum_{i=1}^N \ell(h(\mathbf{x}_i), y_i)
\label{eq:emp_risk_formal}
\end{equation}
\end{definitionbox}

This leads directly to the core operating rule of modern artificial intelligence:

\begin{theorembox}{The Empirical Risk Minimization (ERM) Principle}
Because true risk $R(h)$ is unobservable, the learner searches through hypothesis class $\mathcal{H}$ and selects the hypothesis $\hat{h}_S$ that achieves the smallest possible empirical risk on the training set:
\begin{equation}
\hat{h}_S \triangleq \argmin_{h \in \mathcal{H}} R_{\text{emp}}(h; S) = \argmin_{h \in \mathcal{H}} \frac{1}{N} \sum_{i=1}^N \ell(h(\mathbf{x}_i), y_i)
\label{eq:erm_formal}
\end{equation}
\end{theorembox}

\subsubsection{Why ERM is Dangerous: The Downward Optimistic Bias}
Here lies the most critical conceptual trap in all of machine learning.

If you pick a \emph{fixed} function $h$ out of a hat before seeing any data, its empirical risk $R_{\text{emp}}(h; S)$ is an unbiased estimate of its true risk:
\begin{equation}
\E_{S \sim \mathcal{D}^N} [R_{\text{emp}}(h; S)] = R(h)
\end{equation}
By the Law of Large Numbers, as your sample size grows ($N \to \infty$), the homework score converges to the true exam score.

**However**, the hypothesis $\hat{h}_S$ output by the learning algorithm was **not chosen before seeing the data**! It was chosen specifically to minimize the error on that particular sample $S$. The algorithm examined every quirk, every outlier, and every random coincidence in $S$, and tailored $\hat{h}_S$ to fit them.

Consequently, the empirical risk of the chosen minimizer is **strictly downward-biased** (systematically optimistic):
\begin{equation}
\E_S [R_{\text{emp}}(\hat{h}_S; S)] \le \min_{h \in \mathcal{H}} R(h) \le R(\hat{h}_S)
\end{equation}
The model's score on its training homework is almost always better than its score on the unseen final exam.

\begin{definitionbox}{The Generalization Gap}
The difference between how a model performs in the real world versus how it performed on its training set is called the \textbf{Generalization Gap}:
\begin{equation}
\Delta_{\text{gen}}(\hat{h}_S) \triangleq R(\hat{h}_S) - R_{\text{emp}}(\hat{h}_S; S)
\label{eq:generalization_gap_def}
\end{equation}
\end{definitionbox}
When $\Delta_{\text{gen}}$ is large, the model has fallen victim to **overfitting**: it has memorized the training sample rather than learning the general physical law.

% ----------------------------------------------------------------------
% IN-CHAPTER WORKED EXAMPLES
% ----------------------------------------------------------------------
\subsection{In-Chapter Step-by-Step Pedagogical Examples}

Let us step away from abstract algebra and walk through concrete numbers that reveal these dynamics clearly.

\begin{examplebox}{Worked Example 1.1: A Toy Universe with 4 States}
\textbf{The Scenario:}
Imagine a tiny, transparent universe where the input feature $X$ can only take two values: $\{1, 2\}$, and the target label $Y$ can only take two values: $\{0, 1\}$.
Because this universe is simple, we can write down Nature's complete joint distribution $\mathcal{D}$ in a $2 \times 2$ table:
\begin{equation*}
\begin{array}{c|cc|c}
P(X = x, Y = y) & Y = 0 & Y = 1 & P(X = x) \\
\hline
X = 1 & 0.40 & 0.10 & 0.50 \\
X = 2 & 0.10 & 0.40 & 0.50 \\
\hline
P(Y = y) & 0.50 & 0.50 & 1.00
\end{array}
\end{equation*}

\textbf{Physical Meaning of Nature's Table:}
\begin{itemize}
    \item When $X = 1$, Nature produces label $Y = 0$ with $80\%$ probability ($0.40 / 0.50 = 0.80$) and label $Y = 1$ with $20\%$ probability ($0.10 / 0.50 = 0.20$). There is intrinsic noise!
    \item When $X = 2$, Nature produces label $Y = 1$ with $80\%$ probability ($0.40 / 0.50 = 0.80$) and label $Y = 0$ with $20\%$ probability.
\end{itemize}

Since an input can only be $1$ or $2$, and an output can only be $0$ or $1$, there are only $2^2 = 4$ possible deterministic prediction rules in the entire universe! Our hypothesis class $\mathcal{H} = \{h_1, h_2, h_3, h_4\}$ is:
\begin{equation*}
\begin{aligned}
h_1(x) &= [h(1)=0, \; h(2)=0] \quad \text{(Always guess 0)} \\
h_2(x) &= [h(1)=0, \; h(2)=1] \quad \text{(Guess 0 for 1, Guess 1 for 2)} \\
h_3(x) &= [h(1)=1, \; h(2)=0] \quad \text{(Guess 1 for 1, Guess 0 for 2)} \\
h_4(x) &= [h(1)=1, \; h(2)=1] \quad \text{(Always guess 1)}
\end{aligned}
\end{equation*}
We evaluate predictions using $0-1$ loss $\ell(\hat{y}, y) = \mathbb{I}(\hat{y} \ne y)$.

\textbf{Questions to Solve:}
\begin{enumerate}
    \item Calculate the exact True Risk $R(h)$ for each of the four rules. Which rule is Nature's true favorite?
    \item Now, suppose a machine learning student collects a tiny training dataset of $N = 4$ observations drawn from this universe:
    \begin{equation*}
    S = \big\{ (1, 0), \; (1, 1), \; (2, 0), \; (2, 0) \big\}
    \end{equation*}
    Calculate the Empirical Risk $R_{\text{emp}}(h; S)$ for each of the four hypotheses.
    \item Which rule does the ERM algorithm choose? What is its Generalization Gap?
\end{enumerate}

\tcblower
\textbf{Detailed Step-by-Step Solution:}

\textbf{Step 1: Calculating True Risk $R(h)$ across Nature's Entire Universe}
The True Risk is the probability that rule $h$ makes a mistake under Nature's joint distribution:
\begin{equation*}
R(h) = \sum_{x \in \{1, 2\}} \sum_{y \in \{0, 1\}} \mathbb{I}(h(x) \ne y) P(X = x, Y = y)
\end{equation*}
Let us compute this for each rule:
\begin{itemize}
    \item \textbf{Rule $h_1 = [0, 0]$:} Predicts $0$ everywhere. It is wrong whenever $Y = 1$.
    \begin{equation*}
    R(h_1) = P(X=1, Y=1) + P(X=2, Y=1) = 0.10 + 0.40 = \mathbf{0.50} \quad (50\% \text{ error})
    \end{equation*}
    \item \textbf{Rule $h_2 = [0, 1]$:} Predicts $0$ at $X=1$ and $1$ at $X=2$. It is wrong on $(1, 1)$ and $(2, 0)$.
    \begin{equation*}
    R(h_2) = P(X=1, Y=1) + P(X=2, Y=0) = 0.10 + 0.10 = \mathbf{0.20} \quad (20\% \text{ error})
    \end{equation*}
    \item \textbf{Rule $h_3 = [1, 0]$:} Predicts the exact opposite of $h_2$. It is wrong on $(1, 0)$ and $(2, 1)$.
    \begin{equation*}
    R(h_3) = P(X=1, Y=0) + P(X=2, Y=1) = 0.40 + 0.40 = \mathbf{0.80} \quad (80\% \text{ error})
    \end{equation*}
    \item \textbf{Rule $h_4 = [1, 1]$:} Predicts $1$ everywhere. It is wrong whenever $Y = 0$.
    \begin{equation*}
    R(h_4) = P(X=1, Y=0) + P(X=2, Y=0) = 0.40 + 0.10 = \mathbf{0.50} \quad (50\% \text{ error})
    \end{equation*}
\end{itemize}
\textbf{Conclusion 1:} Rule $h_2$ is Nature's true favorite, achieving $R(h_2) = 0.20$. Notice that even the best rule has a $20\%$ error rate! Why? Because Nature has intrinsic label noise: at $X=1$, $Y=1$ occurs $20\%$ of the time. This $0.20$ is the irreducible **Bayes Error**.

\textbf{Step 2: Calculating Empirical Risk on the Sample $S$}
The student's sample contains 4 points: $(1, 0), (1, 1), (2, 0), (2, 0)$. Notice what happened: by pure chance, the student drew two $(2, 0)$ samples, even though $P(Y=0 \mid X=2)$ is only $20\%$ in the real world!
\begin{align*}
R_{\text{emp}}(h_1; S) &= \frac{1}{4} \big[ \mathbb{I}(0 \ne 0) + \mathbb{I}(0 \ne 1) + \mathbb{I}(0 \ne 0) + \mathbb{I}(0 \ne 0) \big] = \frac{0 + 1 + 0 + 0}{4} = \mathbf{0.25} \\
R_{\text{emp}}(h_2; S) &= \frac{1}{4} \big[ \mathbb{I}(0 \ne 0) + \mathbb{I}(0 \ne 1) + \mathbb{I}(1 \ne 0) + \mathbb{I}(1 \ne 0) \big] = \frac{0 + 1 + 1 + 1}{4} = \mathbf{0.75} \\
R_{\text{emp}}(h_3; S) &= \frac{1}{4} \big[ \mathbb{I}(1 \ne 0) + \mathbb{I}(1 \ne 1) + \mathbb{I}(0 \ne 0) + \mathbb{I}(0 \ne 0) \big] = \frac{1 + 0 + 0 + 0}{4} = \mathbf{0.25} \\
R_{\text{emp}}(h_4; S) &= \frac{1}{4} \big[ \mathbb{I}(1 \ne 0) + \mathbb{I}(1 \ne 1) + \mathbb{I}(1 \ne 0) + \mathbb{I}(1 \ne 0) \big] = \frac{1 + 0 + 1 + 1}{4} = \mathbf{0.75}
\end{align*}

\textbf{Step 3: ERM Selection and the Trap}
The ERM algorithm inspects the empirical risks and seeks the minimum. The minimum is $0.25$, achieved by rules $h_1$ and $h_3$.
The algorithm picks $\hat{h}_S = h_1 = [0, 0]$.
Notice the catastrophe:
\begin{itemize}
    \item The algorithm **completely rejected the true optimal rule $h_2$** because on this sample, $h_2$ had a terrible score of $75\%$!
    \item The algorithm selected $h_1$, bragging that its training error is only $25\%$.
    \item But when we deploy $h_1$ into the real world, its true risk is $R(h_1) = 0.50$ ($50\%$ error!).
\end{itemize}
The Generalization Gap is:
\begin{equation*}
\Delta_{\text{gen}}(\hat{h}_S) = R(\hat{h}_S) - R_{\text{emp}}(\hat{h}_S; S) = 0.50 - 0.25 = \mathbf{+0.25}
\end{equation*}
The algorithm was fooled by the random sampling fluctuations of finite data.
\end{examplebox}

% ----------------------------------------------------------------------
% PART 3: VISUAL & GEOMETRIC INTUITION
% ----------------------------------------------------------------------
\section{Visual \& Geometric Intuition}

\subsection{The Geometry of the Learning Pipeline}
Figure~\ref{fig:ml_pipeline_geometry} visualizes how information flows from Nature's unseen generative process through a finite training sample, into an empirical optimization algorithm, and finally out into the real world where generalization error is evaluated.

\begin{figure}[htbp]
\centering
\begin{tikzpicture}[
    scale=0.92,
    font=\sffamily\footnotesize,
    node distance=1.1cm and 1.3cm,
    block/.style={rectangle, draw=primarydark, fill=primarylight, thick, rounded corners=3pt, minimum width=2.7cm, minimum height=1.1cm, align=center},
    model/.style={rectangle, draw=tealborder, fill=tealbg, thick, rounded corners=3pt, minimum width=2.9cm, minimum height=1.1cm, align=center},
    risk/.style={rectangle, draw=coralborder, fill=coralbg, thick, rounded corners=3pt, minimum width=2.7cm, minimum height=1.1cm, align=center},
    arr/.style={->, >=stealth, thick, color=primarydark}
]

% Top Row: Data Flow
\node[block] (world) {\textbf{Generative World}\\Joint Measure $\mathcal{D}$\\$(\mathbf{x}, y) \sim \mathcal{D}$};
\node[block, right=of world] (sample) {\textbf{Finite Training Set}\\$S = \{(\mathbf{x}_i, y_i)\}_{i=1}^N$\\$N$ i.i.d. draws};
\node[model, right=of sample] (erm) {\textbf{ERM Optimization}\\$\hat{h}_S = \argmin_{h \in \mathcal{H}} R_{\text{emp}}(h)$\\Inductive Bias $\mathcal{H}$};

% Bottom Row: Risk Assessment
\node[risk, below=1.4cm of erm] (remp) {\textbf{Empirical Risk}\\$R_{\text{emp}}(\hat{h}_S; S) = \frac{1}{N}\sum \ell_i$\\Observable Training Loss};
\node[risk, left=of remp] (rtrue) {\textbf{True Generalization Risk}\\$R(\hat{h}_S) = \E_{(\mathbf{x},y)}[\ell(\hat{h}_S(\mathbf{x}), y)]$\\Unobservable Expectation};
\node[block, left=of rtrue] (gap) {\textbf{Generalization Gap}\\$\Delta_{\text{gen}} = R(\hat{h}_S) - R_{\text{emp}}$\\Overfitting Magnitude};

% Arrows
\draw[arr] (world) -- node[above] {Sample $S \sim \mathcal{D}^N$} (sample);
\draw[arr] (sample) -- node[above] {Minimize} (erm);
\draw[arr] (erm) -- (remp);
\draw[arr] (remp) -- node[above] {Generalization Test} (rtrue);
\draw[arr] (rtrue) -- (gap);
\draw[dashed, ->, thick, color=coralborder] (world) |- node[below right, pos=0.25] {Evaluated Against Nature} (rtrue);

\end{tikzpicture}
\caption{The Statistical Learning Pipeline: Connecting unknown generative distribution $\mathcal{D}$, finite training sample $S$, empirical optimization within hypothesis class $\mathcal{H}$, and true out-of-sample risk.}
\label{fig:ml_pipeline_geometry}
\end{figure}

\subsection{Why Nature Cannot Be Beaten: Overlapping Distributions}
Why is zero generalization error fundamentally impossible in real-world problems? Consider Figure~\ref{fig:bayes_boundary_geometry}, which plots the probability densities of two competing classes along a single measurement axis $x$.

\begin{figure}[htbp]
\centering
\begin{tikzpicture}
\begin{axis}[
    width=0.88\textwidth,
    height=6.2cm,
    samples=120,
    domain=-3.5:6.5,
    xlabel={Observed Feature Value $x$},
    ylabel={Probability Density $p(x \mid y)$},
    ymin=0, ymax=0.45,
    legend pos=north east,
    axis lines=left,
    legend style={font=\footnotesize\sffamily, fill=white, draw=gray!30}
]

% Class 0: N(0, 1)
\addplot[thick, color=primarydark, fill=primarylight, fill opacity=0.3] {exp(-0.5*x^2)/sqrt(2*pi)};
\addlegendentry{$p(x \mid Y=0) \sim \mathcal{N}(0, 1)$ (Healthy)}

% Class 1: N(2.5, 1)
\addplot[thick, color=tealdef, fill=tealbg, fill opacity=0.3] {exp(-0.5*(x-2.5)^2)/sqrt(2*pi)};
\addlegendentry{$p(x \mid Y=1) \sim \mathcal{N}(2.5, 1)$ (Diseased)}

% Shaded Bayes Error Region
\addplot[domain=1.25:4.5, draw=none, fill=coralborder, fill opacity=0.4] {exp(-0.5*x^2)/sqrt(2*pi)} \closedcycle;
\addplot[domain=-2.0:1.25, draw=none, fill=coralborder, fill opacity=0.4] {exp(-0.5*(x-2.5)^2)/sqrt(2*pi)} \closedcycle;

% Optimal Boundary Line
\draw[dashed, thick, color=coralformula] (axis cs:1.25, 0) -- (axis cs:1.25, 0.42)
    node[above, font=\footnotesize\bfseries\sffamily] {Optimal Boundary $x^* = 1.25$};

\node[font=\footnotesize\sffamily, color=coralformula] at (axis cs:1.25, 0.08) {\textbf{Bayes Error Region}};
\node[font=\footnotesize\sffamily, color=primarydark] at (axis cs:-1.0, 0.28) {Predict $\hat{y} = 0$};
\node[font=\footnotesize\sffamily, color=tealdef] at (axis cs:3.8, 0.28) {Predict $\hat{y} = 1$};

\end{axis}
\end{tikzpicture}
\caption{The Geometry of Bayes Error: In regions where class densities overlap, no deterministic rule can achieve $100\%$ accuracy. The shaded area represents the irreducible ambiguity of the physical measurement.}
\label{fig:bayes_boundary_geometry}
\end{figure}

Notice what happens at $x = 1.25$. Both curves have positive height! If an observation arrives at $x = 1.5$, Class 1 is more likely than Class 0, so any rational model should guess $\hat{y} = 1$. But Nature \emph{still} produces Class 0 at $x = 1.5$ some fraction of the time. Every time that happens, our model is wrong! 

This shaded area is the **Bayes Error**. It represents Nature's irreducible ambiguity. No artificial intelligence, no deep neural network with a trillion parameters, and no supercomputer can ever achieve a lower error rate than this shaded area.

\subsection{The U-Shaped Generalization Curve}
Figure~\ref{fig:bias_variance_geometry} illustrates what happens as we make our hypothesis class $\mathcal{H}$ more flexible (e.g., increasing the degree of a polynomial from 1 to 20).

\begin{figure}[htbp]
\centering
\begin{tikzpicture}
\begin{axis}[
    width=0.88\textwidth,
    height=6.0cm,
    domain=0.5:5.5,
    xmin=0.5, xmax=5.5,
    ymin=0, ymax=10,
    ticks=none,
    xlabel={Model Complexity / Flexibility of $\mathcal{H}$},
    ylabel={Risk / Predictive Error},
    axis lines=left,
    legend pos=north east,
    legend style={font=\footnotesize\sffamily, fill=white, draw=gray!30}
]

% Training Error (Monotonically decreasing)
\addplot[thick, color=tealdef, domain=0.8:5.2] {7.5 * exp(-0.7*x) + 0.3};
\addlegendentry{Training Error (Empirical Risk)}

% True Risk (U-shaped)
\addplot[very thick, color=coralformula, domain=0.8:5.2] {7.5 * exp(-0.7*x) + 0.25*exp(0.85*x) + 1.2};
\addlegendentry{Real-World Error (True Risk)}

% Bayes Error line
\addplot[dashed, thick, color=gray, domain=0.5:5.5] {1.2};
\addlegendentry{Irreducible Bayes Noise $R(h^*)$}

% Optimal capacity vertical line
\draw[dotted, thick, color=primarydark] (axis cs:2.35, 0) -- (axis cs:2.35, 9.5)
    node[above, font=\footnotesize\bfseries\sffamily] {Sweet Spot};

% Annotations
\node[font=\footnotesize\sffamily, color=primarydark, align=center] at (axis cs:1.3, 7.8) {\textbf{Underfitting}\\Too Rigid\\(High Bias)};
\node[font=\footnotesize\sffamily, color=traprad, align=center] at (axis cs:4.4, 7.8) {\textbf{Overfitting}\\Too Flexible\\(High Variance)};

\end{axis}
\end{tikzpicture}
\caption{The Fundamental Dilemma of Machine Learning: Training error drops monotonically as models become more complex, but true test error forms a U-shape. The goal of machine learning is to find the sweet spot at the bottom of the U.}
\label{fig:bias_variance_geometry}
\end{figure}

% ----------------------------------------------------------------------
% PART 4: DEEP INTUITION & CONCEPTUAL DERIVATIONS
% ----------------------------------------------------------------------
\section{Deep Intuition \& Conceptual Derivations}

\subsection{Derivation of the Bayes Optimal Regressor: The Center of Mass}
Suppose for a moment that an omniscient oracle handed you Nature's true probability distribution $p(\mathbf{x}, y)$. If you are evaluated under squared error loss $\ell(\hat{y}, y) = (\hat{y} - y)^2$, what is the mathematically best prediction function $h^*(\mathbf{x})$ you could possibly build?

\subsubsection{The Dartboard Intuition Before the Calculus}
Before doing any calculus, consider an intuitive game. You must guess a single number $c$ to predict the height of a target. Nature generates three possible targets at heights:
\begin{equation*}
y_1 = 10, \quad y_2 = 20, \quad y_3 = 30
\end{equation*}
all occurring with equal probability ($1/3$ each). You are penalized by the sum of squared errors: $\sum_{i=1}^3 (c - y_i)^2$.
\begin{itemize}
    \item If you guess $c = 10$: Penalty is $(10 - 10)^2 + (10 - 20)^2 + (10 - 30)^2 = 0 + 100 + 400 = \mathbf{500}$.
    \item If you guess $c = 30$: Penalty is $(30 - 10)^2 + (30 - 20)^2 + (30 - 30)^2 = 400 + 100 + 0 = \mathbf{500}$.
    \item If you guess $c = 20$ (the arithmetic average!): Penalty is:
    \begin{equation*}
    (20 - 10)^2 + (20 - 20)^2 + (20 - 30)^2 = 100 + 0 + 100 = \mathbf{200}
    \end{equation*}
\end{itemize}
Under a quadratic penalty, the **average** always minimizes the total loss. In physics, the average is the **center of mass**. If you support a physical rod loaded with weights at $10, 20, 30$, it balances perfectly at $20$.

When we extend this from three numbers to a continuous probability distribution, the ``average of $Y$ given that $\mathbf{X} = \mathbf{x}$'' is simply the **conditional expectation**:
\begin{equation}
h^*(\mathbf{x}) = \E[Y \mid \mathbf{X} = \mathbf{x}]
\end{equation}
Now let us prove this rigorously.

\begin{theorem}[Bayes Optimal Regressor Under Squared Loss]
Under squared error loss $\ell(\hat{y}, y) = (\hat{y} - y)^2$, the predictor achieving the absolute minimum true risk across all measurable functions is:
\begin{equation}
h^*(\mathbf{x}) = \E[Y \mid \mathbf{X} = \mathbf{x}]
\label{eq:bayes_regressor_proof_target}
\end{equation}
\end{theorem}

\begin{proof}
By the Law of Total Expectation, we expand the True Risk:
\begin{align}
R(h) &= \E_{(\mathbf{X}, Y)} \left[ (h(\mathbf{X}) - Y)^2 \right] \nonumber \\
&= \E_{\mathbf{X}} \left[ \E_{Y \mid \mathbf{X}} \left[ (h(\mathbf{X}) - Y)^2 \;\middle|\; \mathbf{X} = \mathbf{x} \right] \right]
\end{align}
Because the outer expectation averages over $p(\mathbf{x}) \ge 0$, we can minimize the entire integral by minimizing the inner conditional expectation pointwise for each individual $\mathbf{x} \in \mathcal{X}$ separately.

Fix an arbitrary $\mathbf{x}$. Let $c \triangleq h(\mathbf{x}) \in \mathbb{R}$ be the scalar prediction we will output at this specific $\mathbf{x}$. Define the conditional expected loss function $g(c): \mathbb{R} \to \mathbb{R}$:
\begin{equation}
g(c) \triangleq \E_{Y \mid \mathbf{X}} \left[ (c - Y)^2 \;\middle|\; \mathbf{X} = \mathbf{x} \right]
\end{equation}
Expand the square inside the expectation:
\begin{align}
g(c) &= \E_{Y \mid \mathbf{X}} \left[ c^2 - 2cY + Y^2 \;\middle|\; \mathbf{X} = \mathbf{x} \right] \nonumber \\
&= c^2 - 2c \, \E[Y \mid \mathbf{X} = \mathbf{x}] + \E[Y^2 \mid \mathbf{X} = \mathbf{x}]
\end{align}
Notice that $g(c)$ is a simple, upward-opening parabola with respect to $c$. To find its minimum, take the derivative and set it to zero:
\begin{equation}
\frac{dg}{dc} = 2c - 2 \, \E[Y \mid \mathbf{X} = \mathbf{x}] = 0 \implies c^* = \E[Y \mid \mathbf{X} = \mathbf{x}]
\end{equation}
Confirming the minimum via the second derivative:
\begin{equation}
\frac{d^2 g}{dc^2} = 2 > 0
\end{equation}
Since the second derivative is positive, $c^* = \E[Y \mid \mathbf{X} = \mathbf{x}]$ is the unique global minimizer.
\end{proof}

\subsubsection{What is Nature's Irreducible Noise?}
If we plug Nature's optimal prediction $h^*(\mathbf{x}) = \E[Y \mid \mathbf{X} = \mathbf{x}]$ back into the True Risk equation, what error remains?
\begin{align}
R(h^*) &= \E_{\mathbf{X}} \left[ \E_{Y \mid \mathbf{X}} \left[ (Y - \E[Y \mid \mathbf{X}])^2 \;\middle|\; \mathbf{X} \right] \right] \nonumber \\
&= \E_{\mathbf{X}} \left[ \Var(Y \mid \mathbf{X}) \right]
\label{eq:irreducible_variance}
\end{align}
By definition, $\E_{Y \mid \mathbf{X}} [ (Y - \E[Y \mid \mathbf{X}])^2 ]$ is the **conditional variance** of $Y$ given $\mathbf{X}$. 

This equation carries profound philosophical significance: **even if you have infinite data, infinite computing power, and the perfect model, your prediction error cannot drop below $\E[\Var(Y \mid \mathbf{X})]$.** This is the variance of Nature itself—the unobserved quantum fluctuations, sensor thermal noise, and hidden variables that no algorithm can anticipate.

\subsection{Derivation of the Bayes Optimal Classifier: Majority Voting}
Now let us switch from continuous regression to discrete classification. Suppose $\mathcal{Y} \in \{0, 1\}$, and our penalty is $0-1$ loss: $\ell(\hat{y}, y) = \mathbb{I}(\hat{y} \ne y)$. What is the best possible decision rule?

\subsubsection{The Common Sense Intuition}
Suppose a patient has a specific symptom vector $\mathbf{x}$. You look up the medical records of everyone with identical symptoms:
\begin{itemize}
    \item $85\%$ of them tested positive ($Y = 1$).
    \item $15\%$ of them tested negative ($Y = 0$).
\end{itemize}
You must make a single call: guess $1$ or guess $0$.
\begin{itemize}
    \item If you guess $\hat{y} = 1$: You are wrong only when the patient is in the $15\%$ minority. Your expected error is **$0.15$**.
    \item If you guess $\hat{y} = 0$: You are wrong whenever the patient is in the $85\%$ majority. Your expected error is **$0.85$**.
\end{itemize}
To minimize mistakes, you must predict the label that has more than a $50\%$ probability of being true! This is simply the principle of **majority voting**.

\begin{theorem}[Bayes Optimal Binary Classifier]
Define the posterior probability of class 1 as:
\begin{equation}
\eta(\mathbf{x}) \triangleq P(Y = 1 \mid \mathbf{X} = \mathbf{x})
\end{equation}
Under symmetric $0-1$ loss, the Bayes optimal decision rule is:
\begin{equation}
h^*(\mathbf{x}) = \begin{cases} 
1 & \text{if } \eta(\mathbf{x}) \ge \frac{1}{2} \\
0 & \text{if } \eta(\mathbf{x}) < \frac{1}{2}
\end{cases}
\label{eq:bayes_classifier_rule_pedagogical}
\end{equation}
\end{theorem}

\begin{proof}
Condition on $\mathbf{X} = \mathbf{x}$. The pointwise expected $0-1$ loss for predicting label $\hat{y} \in \{0, 1\}$ is:
\begin{align}
\E_{Y \mid \mathbf{X}} [\mathbb{I}(\hat{y} \ne Y) \mid \mathbf{x}] &= P(Y \ne \hat{y} \mid \mathbf{x}) \nonumber \\
&= \mathbb{I}(\hat{y} = 0) P(Y = 1 \mid \mathbf{x}) + \mathbb{I}(\hat{y} = 1) P(Y = 0 \mid \mathbf{x}) \nonumber \\
&= \mathbb{I}(\hat{y} = 0) \eta(\mathbf{x}) + \mathbb{I}(\hat{y} = 1) \big(1 - \eta(\mathbf{x})\big)
\end{align}
The algorithm has only two choices:
\begin{itemize}
    \item Predict $\hat{y} = 1 \implies \text{Expected Error} = 1 - \eta(\mathbf{x})$.
    \item Predict $\hat{y} = 0 \implies \text{Expected Error} = \eta(\mathbf{x})$.
\end{itemize}
Predicting $\hat{y} = 1$ is superior whenever:
\begin{equation}
1 - \eta(\mathbf{x}) \le \eta(\mathbf{x}) \iff 1 \le 2\eta(\mathbf{x}) \iff \eta(\mathbf{x}) \ge \frac{1}{2}
\end{equation}
The minimal error rate achieved at this point is $\min\big(\eta(\mathbf{x}), 1 - \eta(\mathbf{x})\big)$. Averaging this over all possible inputs gives the irreducible **Bayes Error**:
\begin{equation}
R(h^*) = \E_{\mathbf{X}} \left[ \min\big(\eta(\mathbf{X}), 1 - \eta(\mathbf{X})\big) \right]
\end{equation}
\end{proof}

\begin{examplebox}{Worked Example 1.2: Two Overlapping Gaussians with Unequal Priors}
\textbf{The Scenario:}
A machine learning engineer at a semiconductor fabrication plant is designing an automated inspection system to classify computer chips as \textbf{Defective ($Y = 1$)} or \textbf{Functional ($Y = 0$)} based on a single sensor measurement $X \in \mathbb{R}$.
Extensive historical data establishes that:
\begin{itemize}
    \item Functional chips produce sensor readings distributed as $\mathcal{N}(\mu_0 = 0, \sigma^2 = 1)$.
    \item Defective chips produce sensor readings distributed as $\mathcal{N}(\mu_1 = 2, \sigma^2 = 1)$.
    \item Overall, $70\%$ of manufactured chips are Functional ($P(Y=0) = 0.70$) and $30\%$ are Defective ($P(Y=1) = 0.30$).
\end{itemize}

\textbf{Questions to Solve:}
\begin{enumerate}
    \item If you only looked at the means ($\mu_0 = 0$ and $\mu_1 = 2$), where is the naive geometric midpoint?
    \item Use Bayes' rule to derive the exact posterior probability $\eta(x) = P(Y = 1 \mid X = x)$.
    \item Compute the exact threshold $x^*$ where the model should switch from predicting Functional ($0$) to Defective ($1$). Why is $x^* > 1.0$?
\end{enumerate}

\tcblower
\textbf{Step-by-Step Analytical Solution:}

\textbf{Step 1: Naive Geometric Midpoint}
The center of the Functional distribution is $\mu_0 = 0$; the center of the Defective distribution is $\mu_1 = 2$.
The naive geometric midpoint between them is:
\begin{equation*}
x_{\text{midpoint}} = \frac{0 + 2}{2} = \mathbf{1.0}
\end{equation*}
If prior probabilities were equal ($50/50$), our decision boundary would sit exactly at $x = 1.0$.

\textbf{Step 2: Deriving the True Posterior $\eta(x)$ via Bayes' Rule}
By Bayes' Theorem:
\begin{equation*}
\eta(x) = P(Y=1 \mid x) = \frac{p(x \mid Y=1) P(Y=1)}{p(x \mid Y=1) P(Y=1) + p(x \mid Y=0) P(Y=0)}
\end{equation*}
Substituting the Gaussian densities $p(x \mid \mu, 1) = \frac{1}{\sqrt{2\pi}} e^{-(x - \mu)^2 / 2}$:
\begin{equation*}
p(x \mid Y=0) = \frac{1}{\sqrt{2\pi}} e^{-x^2 / 2}, \quad p(x \mid Y=1) = \frac{1}{\sqrt{2\pi}} e^{-(x - 2)^2 / 2}
\end{equation*}
Divide the numerator and denominator by $p(x \mid Y=1) P(Y=1)$:
\begin{equation*}
\eta(x) = \frac{1}{1 + \frac{P(Y=0)}{P(Y=1)} \cdot \frac{p(x \mid Y=0)}{p(x \mid Y=1)}}
\end{equation*}
Compute the density ratio:
\begin{align*}
\frac{p(x \mid Y=0)}{p(x \mid Y=1)} &= \frac{e^{-x^2 / 2}}{e^{-(x - 2)^2 / 2}} = \exp\left( -\frac{x^2}{2} + \frac{x^2 - 4x + 4}{2} \right) \\
&= \exp\left( \frac{-4x + 4}{2} \right) = \exp(2 - 2x) = e^{2(1 - x)}
\end{align*}
Now plug in the prior probabilities $P(Y=0) = 0.70$ and $P(Y=1) = 0.30$:
\begin{equation*}
\frac{P(Y=0)}{P(Y=1)} = \frac{0.70}{0.30} = \frac{7}{3}
\end{equation*}
This yields the exact formula for the probability that a chip is defective:
\begin{equation*}
\eta(x) = \frac{1}{1 + \frac{7}{3} e^{2(1 - x)}}
\end{equation*}

\textbf{Step 3: Finding the Bayes Optimal Threshold $x^*$}
By Theorem~\ref{eq:bayes_classifier_rule_pedagogical}, the decision boundary occurs when the probability of defect reaches exactly $50\%$ ($\eta(x^*) = 0.50$):
\begin{align*}
\frac{1}{1 + \frac{7}{3} e^{2(1 - x^*)}} = \frac{1}{2} &\iff 1 + \frac{7}{3} e^{2(1 - x^*)} = 2 \\
&\iff \frac{7}{3} e^{2(1 - x^*)} = 1 \\
&\iff e^{2(1 - x^*)} = \frac{3}{7}
\end{align*}
Take the natural logarithm of both sides:
\begin{align*}
2(1 - x^*) = \log\left(\frac{3}{7}\right) = -\log\left(\frac{7}{3}\right)
\end{align*}
Divide by 2 and solve for $x^*$:
\begin{equation*}
1 - x^* = -\frac{1}{2} \log\left(\frac{7}{3}\right) \implies x^* = 1 + \frac{1}{2} \log\left(\frac{7}{3}\right)
\end{equation*}
Since $\frac{7}{3} \approx 2.3333$ and $\log(2.3333) \approx 0.8473$:
\begin{equation*}
x^* = 1 + \frac{1}{2}(0.8473) = 1 + 0.4236 = \mathbf{1.4236}
\end{equation*}

\textbf{Pedagogical Intuition: Why Did the Boundary Shift Right?}
Because $70\%$ of all chips are functional, Nature has a strong prior bias toward functional chips. If you observe a sensor value of $x = 1.1$, even though it is slightly closer to the defective mean ($\mu_1 = 2$) than the functional mean ($\mu_0 = 0$), the chip is **still more likely to be functional** because functional chips are so much more common!
To justify calling a chip defective, the sensor reading must move far enough to the right ($x \ge 1.4236$) to overcome the strong baseline prior of the manufacturing line.
\end{examplebox}

\begin{examplebox}{Worked Example 1.3: Asymmetric Costs in Cancer Screening}
\textbf{The Dilemma:}
In medical diagnostics, treating all errors as equal ($\text{loss} = 1$) is a recipe for medical malpractice:
\begin{itemize}
    \item \textbf{False Positive ($\hat{y} = 1, y = 0$):} A healthy patient is flagged as having cancer. They undergo a follow-up ultrasound, experience a few days of anxiety, and are declared cancer-free. Cost: $C_{10} = \mathbf{1}$ unit of harm.
    \item \textbf{False Negative ($\hat{y} = 0, y = 1$):} A patient with an early-stage malignant tumor is told they are completely healthy. The tumor metastasizes over 6 months, causing irreparable harm. Cost: $C_{01} = \mathbf{100}$ units of harm.
    \item Correct decisions incur $0$ cost.
\end{itemize}
What should the model's posterior probability threshold $\tau^*$ be to trigger a biopsy? Should we wait until the model is $50\%$ sure?

\tcblower
\textbf{Step-by-Step Analytical Solution:}

Let $\eta(\mathbf{x}) = P(Y = 1 \mid \mathbf{x})$ be the model's calibrated estimate that the patient has cancer.
When a patient with symptoms $\mathbf{x}$ arrives, the doctor has two choices:
\begin{enumerate}
    \item \textbf{Choice 1: Order a Biopsy ($\hat{y} = 1$):}
    A cost is incurred only if the patient is actually healthy ($Y = 0$). That occurs with probability $1 - \eta(\mathbf{x})$.
    \begin{equation*}
    \text{Expected Harm}(\hat{y} = 1 \mid \mathbf{x}) = C_{10} \cdot P(Y = 0 \mid \mathbf{x}) = C_{10} \big(1 - \eta(\mathbf{x})\big)
    \end{equation*}
    \item \textbf{Choice 2: Send Patient Home ($\hat{y} = 0$):}
    A cost is incurred only if the patient actually has cancer ($Y = 1$). That occurs with probability $\eta(\mathbf{x})$.
    \begin{equation*}
    \text{Expected Harm}(\hat{y} = 0 \mid \mathbf{x}) = C_{01} \cdot P(Y = 1 \mid \mathbf{x}) = C_{01} \eta(\mathbf{x})
    \end{equation*}
\end{enumerate}

\textbf{Finding the Optimal Decision Threshold:}
The doctor should order a biopsy whenever the expected harm of testing is less than the expected harm of sending the patient home:
\begin{align*}
\text{Expected Harm}(\hat{y} = 1 \mid \mathbf{x}) &\le \text{Expected Harm}(\hat{y} = 0 \mid \mathbf{x}) \\
C_{10} \big(1 - \eta(\mathbf{x})\big) &\le C_{01} \eta(\mathbf{x}) \\
C_{10} - C_{10} \eta(\mathbf{x}) &\le C_{01} \eta(\mathbf{x}) \\
C_{10} &\le (C_{10} + C_{01}) \eta(\mathbf{x})
\end{align*}
Dividing both sides by $(C_{10} + C_{01})$:
\begin{equation*}
\eta(\mathbf{x}) \ge \frac{C_{10}}{C_{10} + C_{01}}
\end{equation*}
The Bayes optimal decision threshold is:
\begin{equation*}
\tau^* = \frac{C_{10}}{C_{10} + C_{01}}
\end{equation*}
Plugging in our clinical costs ($C_{10} = 1$ and $C_{01} = 100$):
\begin{equation*}
\tau^* = \frac{1}{1 + 100} = \frac{1}{101} \approx \mathbf{0.0099} \quad (\mathbf{0.99\%})
\end{equation*}

\textbf{The Life-Saving Lesson:}
The doctor must order a biopsy whenever the probability of malignancy exceeds **$1\%$**! 
If an engineer naively deployed a standard machine learning classifier with a default $50\%$ decision threshold ($\tau = 0.50$), the system would ignore patients with an $8\%, 20\%,$ or $45\%$ risk of cancer, killing hundreds of patients. The threshold of a classifier is not a mathematical constant; it is an economic and ethical choice dictated by the loss matrix.
\end{examplebox}

\subsection{Approximation Error vs. Estimation Error: The Archer Analogy}
To understand why choosing the right model family $\mathcal{H}$ is the most critical decision in machine learning, consider the classical analogy of an archer shooting arrows at a target bullseye:
\begin{itemize}
    \item \textbf{Bias (Approximation Error):} The archer's bow sights are misaligned by 5 degrees to the left. Even if the archer has perfectly steady hands, every arrow consistently clusters to the left of the bullseye. This is an intrinsic limitation of the equipment.
    \item \textbf{Variance (Estimation Error):} The archer's bow is perfectly calibrated, but the archer has shaky hands and the wind is gusting unpredictably. The arrows scatter wildly across the entire target.
\end{itemize}

In mathematical learning theory, we decompose the total excess risk $R(\hat{h}_S) - R(h^*)$ into these two exact components:
\begin{equation}
R(\hat{h}_S) - R(h^*) = \underbrace{\big( R(h^*_{\mathcal{H}}) - R(h^*) \big)}_{\text{Approximation Error (Bias)}} + \underbrace{\big( R(\hat{h}_S) - R(h^*_{\mathcal{H}}) \big)}_{\text{Estimation Error (Variance)}}
\label{eq:approx_estimation_decomp}
\end{equation}
where:
\begin{itemize}
    \item $h^*$ is Nature's unconstrained Bayes optimal predictor (Equation~\eqref{eq:bayes_regressor_proof_target}).
    \item $h^*_{\mathcal{H}} \triangleq \argmin_{h \in \mathcal{H}} R(h)$ is the best possible hypothesis inside our chosen class $\mathcal{H}$.
    \item $\hat{h}_S \triangleq \argmin_{h \in \mathcal{H}} R_{\text{emp}}(h; S)$ is the empirical hypothesis our algorithm actually picked from training data $S$.
\end{itemize}

\begin{figure}[htbp]
\centering
\begin{tikzpicture}[
    scale=0.95,
    font=\sffamily\footnotesize
]
% Space of all functions
\draw[thick, fill=gray!8, draw=gray!40] (-1, -1) rectangle (9, 4.5);
\node[color=gray!70, font=\bfseries\sffamily] at (7.2, 4.0) {All Measurable Functions};

% Hypothesis Space H
\draw[thick, fill=tealbg, draw=tealborder] (1.5, 0.5) ellipse (2.2cm and 1.4cm);
\node[color=tealdef, font=\bfseries\sffamily] at (1.5, 1.6) {Hypothesis Space $\mathcal{H}$};

% Points
\coordinate (bayes) at (7.0, 1.5);
\coordinate (best_h) at (3.2, 1.2);
\coordinate (erm_h) at (0.5, 0.8);

% Draw points
\filldraw[coralformula] (bayes) circle (2.5pt) node[right=3pt] {\textbf{Bayes Optimal } $h^*$};
\filldraw[primarydark] (best_h) circle (2.5pt) node[below right=2pt] {\textbf{Best in Class } $h^*_{\mathcal{H}}$};
\filldraw[tealdef] (erm_h) circle (2.5pt) node[left=3pt] {\textbf{ERM Solution } $\hat{h}_S$};

% Distance vectors
\draw[<->, >=stealth, very thick, color=coralformula] (bayes) -- node[above=2pt, font=\footnotesize\bfseries\sffamily] {Approximation Error $\epsilon_{\text{app}}$} (best_h);
\draw[<->, >=stealth, very thick, color=primarydark] (best_h) -- node[below=2pt, font=\footnotesize\bfseries\sffamily] {Estimation Error $\epsilon_{\text{est}}$} (erm_h);

\end{tikzpicture}
\caption{Geometric Decomposition of Generalization Error: The gap between Nature's truth $h^*$ and the best model in our class $h^*_{\mathcal{H}}$ is Approximation Error. The gap between the best model in class and the sample-trained model $\hat{h}_S$ is Estimation Error.}
\label{fig:approx_est_geometry}
\end{figure}

\begin{enumerate}
    \item \textbf{Approximation Error ($\epsilon_{\text{app}}$):} Measures the geometric distance between Nature's truth $h^*$ and our chosen model class $\mathcal{H}$. If Nature's law is a sine wave, but we restrict $\mathcal{H}$ to straight lines ($y = mx + b$), our model can never fit the sine wave, even if we collect ten billion data points! This error depends solely on the expressiveness of $\mathcal{H}$.
    \item \textbf{Estimation Error ($\epsilon_{\text{est}}$):} Measures how far our sample-trained model $\hat{h}_S$ drifts from the class optimum $h^*_{\mathcal{H}}$ due to finite sample noise. If we choose a massively flexible neural network with 10 million parameters, $\epsilon_{\text{app}} \approx 0$ (a neural network can fit almost any curve). But because the model has so many degrees of freedom, the small random noise in our training sample will cause the weights to oscillate wildly, driving estimation error $\epsilon_{\text{est}}$ through the roof!
\end{enumerate}

\subsection{The Curse of Dimensionality: Why 100 Dimensions Breaks Human Intuition}
Human beings evolved in a three-dimensional world. When we hear the word ``sphere'', we think of a basketball. When we hear ``cube'', we think of a cardboard box. In high-dimensional data science, these physical intuitions are catastrophic traps.

\subsubsection{The Orange Peel Paradox}
Imagine peeling an orange. In our everyday 3D world, the peel is a thin skin representing perhaps $10\%$ of the orange's total volume, while the juicy fruit in the center represents $90\%$.

What happens to an orange in 100-dimensional space?
Consider a $d$-dimensional hypercube of side length $2$: $\mathcal{C}_d = [-1, 1]^d$. Its volume is:
\begin{equation}
V(\mathcal{C}_d) = 2^d
\end{equation}
Now consider the interior core of the fruit, shaved down by a tiny boundary skin of thickness $\epsilon = 0.05$ (a $5\%$ skin thickness): $\mathcal{C}_{\text{core}} = [-(1 - \epsilon), 1 - \epsilon]^d$. Its volume is:
\begin{equation}
V(\mathcal{C}_{\text{core}}) = \big(2(1 - \epsilon)\big)^d = 2^d (1 - \epsilon)^d
\end{equation}
What fraction of the total volume lives in the outer $5\%$ peel?
\begin{equation}
\frac{V(\text{Peel})}{V(\mathcal{C}_d)} = \frac{V(\mathcal{C}_d) - V(\mathcal{C}_{\text{core}})}{V(\mathcal{C}_d)} = \frac{2^d - 2^d(1 - \epsilon)^d}{2^d} = 1 - (1 - \epsilon)^d
\end{equation}
Now, let us evaluate this simple formula as the dimension $d$ grows:
\begin{itemize}
    \item In $d = 1$ (a line segment): $1 - (1 - 0.05)^1 = 1 - 0.95 = \mathbf{0.05} \quad (5\% \text{ in the peel})$.
    \item In $d = 3$ (a physical box): $1 - (0.95)^3 = 1 - 0.857 = \mathbf{0.143} \quad (14.3\% \text{ in the peel})$.
    \item In $d = 50$: $1 - (0.95)^{50} = 1 - 0.0769 = \mathbf{0.923} \quad (\mathbf{92.3\%} \text{ in the peel!})$.
    \item In $d = 100$: $1 - (0.95)^{100} = 1 - 0.0059 = \mathbf{0.994} \quad (\mathbf{99.4\%} \text{ in the peel!})$.
    \item In $d = 500$: $1 - (0.95)^{500} = 1 - 7.2 \times 10^{-12} \approx \mathbf{0.99999999999} \quad (\mathbf{100\%})$.
\end{itemize}
\textbf{The Stunning Reality:} In 100 dimensions, a hypercube is entirely empty! There is no fruit in the center of the orange; **$99.4\%$ of all the volume resides in the outer $5\%$ crust.** If you draw random data points uniformly from a high-dimensional box, you will never find a point near the center; virtually every single data point lives on the extreme outer boundaries.

\subsubsection{The Distance Concentration Phenomenon}
Suppose you generate 100 random points in high-dimensional space and try to find the ``nearest neighbor'' to a query point.

Let $\mathbf{X}, \mathbf{Y}$ be two independent points drawn uniformly from $[0, 1]^d$. The squared Euclidean distance between them is:
\begin{equation}
D^2 = \sum_{j=1}^d (X_j - Y_j)^2
\end{equation}
Each coordinate difference $Z_j = (X_j - Y_j)^2$ is an independent random variable. By calculus integration:
\begin{align}
\E[Z_j] &= \int_0^1 \int_0^1 (x - y)^2 \, dx \, dy = \frac{1}{6} \\
\Var(Z_j) &= \frac{7}{180}
\end{align}
By the Central Limit Theorem, the total squared distance is the sum of $d$ independent random variables. Its expected value scales linearly with dimension:
\begin{equation}
\E[D^2] = \sum_{j=1}^d \E[Z_j] = \frac{d}{6} \implies D \approx \sqrt{\frac{d}{6}}
\end{equation}
The standard deviation of the distance, however, grows much slower, scaling only as $\mathcal{O}(\sqrt{d})$.
What happens to the relative contrast between the nearest neighbor $D_{\min}$ and the furthest neighbor $D_{\max}$?
\begin{equation}
\frac{D_{\max} - D_{\min}}{D_{\min}} \le \frac{\text{Spread}}{\text{Mean Distance}} \approx \frac{\mathcal{O}(\sqrt{d})}{\mathcal{O}(d)} = \mathcal{O}\left( \frac{1}{\sqrt{d}} \right) \xrightarrow{d \to \infty} \mathbf{0}
\end{equation}
\textbf{The Consequence for Machine Learning:} As dimension $d$ grows, the difference between your ``closest neighbor'' and your ``furthest neighbor'' shrinks to zero. In 1,000-dimensional space, every point is almost exactly the same distance away from you! 

Algorithms based on Euclidean distance (such as $k$-Nearest Neighbors or clustering) become completely useless in raw high-dimensional space unless we first apply **dimensionality reduction** (such as PCA or SVD, studied in Chapters 6 and 7).

% ----------------------------------------------------------------------
% PART 5: EXPLICIT GOTCHAS & DIAGNOSTIC MISCONCEPTIONS
% ----------------------------------------------------------------------
\section{Explicit Gotchas \& Diagnostic Misconceptions}

\begin{trapbox}{Gotcha 1: The 100\% Accuracy Delusion (Class Imbalance Trap)}
\textbf{The Scenario:} A junior data scientist trains an AI system to detect fraudulent credit card transactions. The training dataset contains 100,000 transactions, of which 100 are fraud ($0.1\%$) and 99,900 are legitimate ($99.9\%$). The model achieves $99.9\%$ accuracy! The company prepares to deploy it to production. \\
\textbf{The Catastrophe:} The model simply learned to predict:
\begin{equation*}
h(\mathbf{x}) = \text{``Legitimate''} \quad \text{for all inputs!}
\end{equation*}
It never caught a single fraudulent transaction. Its accuracy is $99.9\%$, but its utility is zero. \\
\textbf{The Lesson:} On imbalanced datasets, raw accuracy is meaningless. You must evaluate using precision, recall, cost-weighted loss, or Area Under the ROC Curve (AUC).
\end{trapbox}

\begin{trapbox}{Gotcha 2: The Discontinuous Cliff of 0-1 Loss}
\textbf{The Scenario:} A student defines a classification neural network with a step threshold function $\hat{y} = \mathbb{I}(\mathbf{w}^T \mathbf{x} \ge 0)$ and tries to train the weights $\mathbf{w}$ using gradient descent on the $0-1$ loss $\ell = \mathbb{I}(\hat{y} \ne y)$. \\
\textbf{The Catastrophe:} Gradient descent fails completely. The computer reports that the gradient is $0.0000$ at every iteration! \\
\textbf{Why it Happened:} The $0-1$ loss is a piecewise-constant step function. If you nudge the weights $\mathbf{w}$ by $0.0001$, the prediction does not change at all, so the derivative $\frac{\partial \ell}{\partial \mathbf{w}}$ is identically zero. The model is standing on a flat plateau; gradient descent has no slope to tell it which way is downhill. \\
\textbf{The Remedy:} We replace the step function with smooth, convex surrogate losses: Logistic Regression uses **Cross-Entropy Loss**, and Support Vector Machines use **Hinge Loss** (Chapter 10).
\end{trapbox}

\begin{trapbox}{Gotcha 3: The Contaminated Test Set (Data Leakage)}
\textbf{The Scenario:} An analyst wants to normalize features so they have mean 0 and variance 1. They compute the mean $\mu$ and standard deviation $\sigma$ across all 50,000 rows in their database, normalize the entire dataset, and \emph{then} split the data into $80\%$ training and $20\%$ test sets. \\
\textbf{The Catastrophe:} Information from the test set leaked into the training set through $\mu$ and $\sigma$! The test set is no longer an independent, pristine evaluation of the real world. The test accuracy will be artificially high, but the model will degrade when deployed to real users. \\
\textbf{The Golden Rule:} Split your data into Train and Test **first**. Compute $\mu$ and $\sigma$ strictly on the training set. Then use those exact frozen training values to scale the test set.
\end{trapbox}

\begin{trapbox}{Gotcha 4: Conflating Correlation with Causation in Feature Spaces}
\textbf{The Scenario:} A linear model trained on hospital data finds that patients who are prescribed heavy doses of morphine have much higher mortality rates. The algorithm concludes: \emph{``Morphine causes death; stop giving morphine to critical patients.''} \\
\textbf{The Fallacy:} Machine learning models discover statistical correlations in $P(y \mid \mathbf{x})$, not causal mechanisms. Critical illness causes both morphine administration and higher mortality (a classic confounding variable). Never confuse an associative predictor with a causal intervention.
\end{trapbox}

% ----------------------------------------------------------------------
% PART 6: COMPARATIVE SYNTHESIS & SUMMARY TABLE
% ----------------------------------------------------------------------
\section{Comparative Synthesis \& Summary Table}

Table~\ref{tab:ml_paradigms_synthesis} summarizes the primary problem formulations that you will master across the 12 weeks of this textbook.

\begin{table}[htbp]
\centering
\footnotesize
\setlength{\tabcolsep}{3.5pt}
\begin{tabularx}{\textwidth}{lp{2.4cm}p{2.6cm}p{2.8cm}p{3.8cm}}
\toprule
\textbf{Paradigm} & \textbf{Data Provided} & \textbf{Core Goal} & \textbf{Standard Loss Metric} & \textbf{Primary Engineering Danger} \\
\midrule
\textbf{Regression} & $(\mathbf{x}_i, y_i)$ with $y_i \in \mathbb{R}$ & Estimate conditional mean $\E[Y \mid \mathbf{X}]$ & Squared Error $(\hat{y} - y)^2$, Absolute Error $|\hat{y} - y|$ & Vulnerable to extreme outliers; high-degree polynomials oscillate wildly. \\
\addlinespace
\textbf{Classification} & $(\mathbf{x}_i, y_i)$ with $y_i \in \{0, 1\}$ & Find decision boundary separating classes & Cross-Entropy $-\log p(y \mid \mathbf{x})$, Hinge Loss $\max(0, 1 - y f)$ & Class imbalance collapse; non-differentiability of direct $0-1$ loss. \\
\addlinespace
\textbf{Dimensionality Reduction} & Unlabeled observations $\mathbf{x}_i \in \mathbb{R}^d$ & Compress $d$ dimensions to $k \ll d$ while saving variance & Reconstruction MSE $\|\mathbf{x} - \hat{\mathbf{x}}\|_2^2$ (SVD / PCA) & Linear subspaces fail to unfold non-linear manifolds (e.g., Swiss Roll). \\
\addlinespace
\textbf{Density Estimation} & Unlabeled observations $\mathbf{x}_i \in \mathbb{R}^d$ & Reconstruct full continuous probability density $p(\mathbf{x})$ & Negative Log-Likelihood $-\sum \log p_\theta(\mathbf{x}_i)$ & The curse of dimensionality: in high dimensions, empty space dominates. \\
\addlinespace
\textbf{Reinforcement Learning} & Actions, environment states, delayed rewards & Learn policy maximizing long-term discounted rewards & Bellman Error $(r + \gamma V(s') - V(s))^2$ & Credit assignment problem: reward arrives hundreds of steps after action. \\
\bottomrule
\end{tabularx}
\caption{Comparative Synthesis of Foundational Machine Learning Problem Paradigms.}
\label{tab:ml_paradigms_synthesis}
\end{table}

% ----------------------------------------------------------------------
% PART 7: PYTHON / SCIENTIFIC COMPUTING LAB
% ----------------------------------------------------------------------
\section{Python / Scientific Computing Lab}

The following runnable Python script empirically demonstrates the two core theoretical themes of this chapter:
\begin{enumerate}
    \item \textbf{The Bias-Variance Experiment:} We fit polynomials of degrees $1, 3, 9,$ and $15$ to noisy data. You will see training error drop monotonically to zero, while test error explodes at degree 15.
    \item \textbf{The High-Dimensional Experiment:} We generate random points in dimensions $d = 2, 10, 50, 200,$ and $1000$, computing the relative distance contrast $\frac{D_{\max} - D_{\min}}{D_{\min}}$ to prove that all points become equidistant.
\end{enumerate}

\begin{lstlisting}[language=Python, caption={Simulating Overfitting Dynamics and High-Dimensional Distance Concentration in Python.}]
import numpy as np

def generate_noisy_sine(n_samples: int, noise_std: float = 0.25, seed: int = 42):
    """Generates synthetic data from true law: y = sin(2*pi*x) + Gaussian noise."""
    rng = np.random.RandomState(seed)
    x = rng.uniform(0.0, 1.0, size=n_samples)
    y = np.sin(2 * np.pi * x) + rng.normal(0.0, noise_std, size=n_samples)
    return x, y

def fit_polynomial(x: np.ndarray, y: np.ndarray, degree: int) -> np.ndarray:
    """Fits polynomial weights w = (Phi^T Phi)^(-1) Phi^T y via Moore-Penrose pseudo-inverse."""
    Phi = np.vander(x, N=degree + 1, increasing=True)
    w = np.linalg.pinv(Phi) @ y
    return w

def compute_mse(x: np.ndarray, y: np.ndarray, w: np.ndarray) -> float:
    """Computes Mean Squared Error for fitted polynomial weights w."""
    Phi = np.vander(x, N=len(w), increasing=True)
    y_pred = Phi @ w
    return float(np.mean((y - y_pred) ** 2))

def run_bias_variance_lab():
    """Demonstrates how increasing polynomial degree causes training MSE to collapse while test MSE explodes."""
    print("=" * 70)
    print("EXPERIMENT 1: EMPIRICAL RISK MINIMIZATION (BIAS-VARIANCE TRADEOFF)")
    print("=" * 70)
    
    # 25 training points, 500 independent out-of-sample test points
    x_train, y_train = generate_noisy_sine(n_samples=25, seed=101)
    x_test, y_test = generate_noisy_sine(n_samples=500, seed=202)

    degrees = [1, 3, 5, 9, 15]
    print(f"{'Degree':<8}{'Train MSE (Homework)':<22}{'Test MSE (Final Exam)':<24}{'Diagnosis':<15}")
    print("-" * 70)

    for deg in degrees:
        w = fit_polynomial(x_train, y_train, degree=deg)
        train_mse = compute_mse(x_train, y_train, w)
        test_mse = compute_mse(x_test, y_test, w)
        
        if deg == 1:
            diagnosis = "Underfitting (High Bias)"
        elif deg == 3:
            diagnosis = "Optimal Sweet Spot"
        elif deg >= 9:
            diagnosis = "Catastrophic Overfitting"
        else:
            diagnosis = "Moderate Fit"
            
        print(f"{deg:<8}{train_mse:<22.5f}{test_mse:<24.5f}{diagnosis:<15}")

def run_curse_of_dimensionality_lab():
    """Demonstrates that in high dimensions, (D_max - D_min) / D_min collapses to 0."""
    print("\n" + "=" * 70)
    print("EXPERIMENT 2: THE CURSE OF DIMENSIONALITY (DISTANCE CONCENTRATION)")
    print("=" * 70)
    print(f"{'Dimension d':<14}{'Mean Distance':<18}{'D_min':<12}{'D_max':<12}{'Relative Spread':<15}")
    print("-" * 70)

    dimensions = [2, 10, 50, 200, 1000]
    n_points = 200
    rng = np.random.RandomState(42)

    for d in dimensions:
        # Draw 200 points uniformly from unit hypercube [0, 1]^d
        X = rng.uniform(0.0, 1.0, size=(n_points, d))
        # Compute all pairwise Euclidean distances
        diffs = X[:, np.newaxis, :] - X[np.newaxis, :, :]
        dists = np.sqrt(np.sum(diffs ** 2, axis=-1))
        # Extract off-diagonal pairs
        pairwise = dists[np.triu_indices(n_points, k=1)]
        
        d_mean = np.mean(pairwise)
        d_min = np.min(pairwise)
        d_max = np.max(pairwise)
        spread = (d_max - d_min) / d_min

        print(f"{d:<14}{d_mean:<18.4f}{d_min:<12.4f}{d_max:<12.4f}{spread:<15.4f}")

if __name__ == "__main__":
    run_bias_variance_lab()
    run_curse_of_dimensionality_lab()
\end{lstlisting}

% ----------------------------------------------------------------------
% PART 8: COMPREHENSIVE PRACTICE PROBLEMS
% ----------------------------------------------------------------------
\section{Comprehensive Practice Problems}

\subsection*{Part A: Analytical \& Mathematical Problems}

\begin{enumerate}[label=\textbf{P\arabic*.}]
    \item \textbf{Bayes Optimal Predictor Under Absolute Loss (MAE):}
    Let $Y \in \mathbb{R}$ be a continuous random variable with conditional probability density function $f(y \mid \mathbf{x})$. Suppose prediction error is penalized by the Mean Absolute Error (MAE) loss:
    \begin{equation*}
    \ell(\hat{y}, y) = |\hat{y} - y|
    \end{equation*}
    Using Leibniz's integral rule to differentiate the expected loss, prove that the Bayes optimal predictor $h^*(\mathbf{x}) = \argmin_{c \in \mathbb{R}} \E[|c - Y| \mid \mathbf{X} = \mathbf{x}]$ is the \textbf{conditional median}:
    \begin{equation*}
    h^*(\mathbf{x}) = \text{Median}(Y \mid \mathbf{X} = \mathbf{x})
    \end{equation*}

    \item \textbf{Optimal Threshold Under Asymmetric Fraud Costs:}
    A bank evaluates an automated fraud detection system. Misclassifying an unauthorized transaction ($Y = 1$) as legitimate ($Y = 0$) incurs a loss of $C_{01} = 50$, whereas an unnecessary security verification call for a legitimate transaction ($Y = 0$) incurs a nuisance cost of $C_{10} = 1$. Let $\eta(\mathbf{x}) = P(Y = 1 \mid \mathbf{X} = \mathbf{x})$. Derive the Bayes optimal decision threshold $\tau^*$ such that predicting $\hat{y} = 1$ when $\eta(\mathbf{x}) \ge \tau^*$ minimizes total expected loss.

    \item \textbf{Generalization Risk Under Stochastic Label Noise:}
    Let $Y \in \{-1, +1\}$ be the true ground-truth label. Suppose training annotations are collected under random worker error: each label is flipped independently with probability $\rho \in [0, 1/2)$, generating observed noisy labels $\tilde{Y}$. Express the true risk under noisy labels $R_{\text{noisy}}(h) \triangleq \E[\mathbb{I}(h(\mathbf{X}) \ne \tilde{Y})]$ strictly in terms of the clean risk $R(h) \triangleq \E[\mathbb{I}(h(\mathbf{X}) \ne Y)]$ and noise parameter $\rho$. Show that minimizing noisy risk preserves the optimal hypothesis ranking.

    \item \textbf{Analytical Proof of Distance Concentration in High Dimensions:}
    Let $\mathbf{X}, \mathbf{Y}$ be two independent random vectors drawn uniformly from the unit hypercube $[0, 1]^d$.
    \begin{enumerate}[label=(\alph*)]
        \item Compute the exact mean $\E[(X_j - Y_j)^2]$ and variance $\Var((X_j - Y_j)^2)$ for coordinate difference $j$.
        \item Using Chebyshev's inequality and the Central Limit Theorem, prove that as dimension $d \to \infty$, the relative distance spread satisfies:
        \begin{equation*}
        \frac{D_{\max} - D_{\min}}{D_{\min}} \xrightarrow{P} 0
        \end{equation*}
    \end{enumerate}

    \item \textbf{Uniform Generalization Bound for Finite Hypothesis Classes:}
    Let $\mathcal{H}$ be a finite hypothesis space containing $|\mathcal{H}| < \infty$ candidate models. Using Hoeffding's inequality for bounded loss functions $\ell \in [0, 1]$ and applying the union bound across all hypotheses, prove that with probability at least $1 - \delta$ over the random draw of training sample $S$ of size $N$, every hypothesis $h \in \mathcal{H}$ simultaneously satisfies:
    \begin{equation*}
    R(h) \le R_{\text{emp}}(h; S) + \sqrt{\frac{\log |\mathcal{H}| + \log(1/\delta)}{2N}}
    \end{equation*}
\end{enumerate}

\subsection*{Part B: Diagnostic Conceptual Multiple-Choice Questions}

\begin{enumerate}[label=\textbf{MCQ\arabic*.}]
    \item A machine learning engineer trains a binary classification model on a heavily imbalanced medical dataset where $98\%$ of instances belong to class $0$ (healthy). The trained model achieves an empirical accuracy of $98.0\%$. Which of the following conclusions is mathematically warranted?
    \begin{enumerate}[label=(\alph*)]
        \item The model has proven high generalization ability on the diseased class ($Y = 1$).
        \item The model could be trivially predicting the constant function $h(\mathbf{x}) = 0$, achieving zero empirical error on class $0$ and $100\%$ empirical error on class $1$.
        \item Minimizing unweighted $0-1$ loss guarantees that the precision of the minority class is maximized.
        \item The empirical risk under cross-entropy loss must evaluate to identically zero.
    \end{enumerate}

    \item Why is the empirical risk of the trained minimizer $R_{\text{emp}}(\hat{h}_S; S)$ a systematically downward-biased (optimistic) estimate of its true risk $R(\hat{h}_S)$?
    \begin{enumerate}[label=(\alph*)]
        \item Because the hypothesis $\hat{h}_S$ was selected specifically to exploit the idiosyncratic stochastic fluctuations of the specific sample $S$.
        \item Because by the Law of Large Numbers, $\frac{1}{N}\sum_{i=1}^N \ell(h(\mathbf{x}_i), y_i) < \E[\ell(h(\mathbf{X}), Y)]$ holds for every fixed $h$.
        \item Because the loss function is strictly concave with respect to the sample size $N$.
        \item Because Jensen's inequality guarantees that the minimum of expectations is strictly greater than the expectation of the minimum.
    \end{enumerate}
\end{enumerate}

\chapter[Differential Calculus for Learning]{Multivariate Differential Calculus \& Optimization Geometry}
\label{chap:week02}

% ----------------------------------------------------------------------
% PART 1: WHY CALCULUS IN ML? THE BRIDGE & THE BLINDFOLDED HIKER
% ----------------------------------------------------------------------
\section{Why Calculus in Machine Learning? The Paradigm, The Promise, and The Mountain}

\subsection{The Architectural Bridge: From Empirical Risk to Vector Calculus}
In Chapter~\ref{chap:week01}, we established the mathematical contract of machine learning through the formal 5-tuple:
\begin{equation}
(\mathcal{X}, \mathcal{Y}, \mathcal{D}, \mathcal{H}, \ell)
\end{equation}
We proved that Nature governs an unknown data-generating distribution $\mathcal{D}$, from which we observe a finite training sample $S = \{(\mathbf{x}_i, y_i)\}_{i=1}^N \sim \mathcal{D}^N$. Under the Empirical Risk Minimization (ERM) inductive principle, the learner selects an optimal hypothesis $h(\cdot; \mathbf{w}) \in \mathcal{H}$ parameterized by a weight vector $\mathbf{w} \in \mathbb{R}^d$ by minimizing the empirical risk functional:
\begin{equation}
\min_{\mathbf{w} \in \mathbb{R}^d} R_S(\mathbf{w}) \triangleq \frac{1}{N}\sum_{i=1}^N \ell\big(h(\mathbf{x}_i; \mathbf{w}), y_i\big)
\label{eq:w02_erm_bridge}
\end{equation}
Chapter~\ref{chap:week01} posed the foundational question: \emph{How does an algorithm actually find the parameter vector $\mathbf{w}^*$ that achieves this minimum?}

This is where Differential Calculus enters as the engine of learning. The loss functional $R_S(\mathbf{w})$ defines an elevation surface over the $d$-dimensional parameter space $\mathbb{R}^d$. Finding the best machine learning model is geometrically equivalent to navigating this high-dimensional terrain to locate its lowest valley floor.

\subsection{The Seductive Fantasy of High School Calculus}
In introductory single-variable calculus, finding the minimum of a function is presented as a clean, three-step algebraic exercise:
\begin{equation}
y = f(w) \xrightarrow{\quad\text{Step 1: Differentiate}\quad} f'(w) \xrightarrow{\quad\text{Step 2: Set to Zero}\quad} f'(w) = 0 \xrightarrow{\quad\text{Step 3: Solve for } w\quad} w^*
\end{equation}
This tidy schoolroom routine creates a seductive fantasy: \emph{to minimize any machine learning objective, simply set its vector derivative to zero and solve the resulting algebraic equation.}

In modern machine learning, this fantasy shatters against two insurmountable computational barriers:
\begin{enumerate}
    \item \textbf{The Dimensionality Explosion:} A modern machine learning model does not have one scalar variable $w$; it possesses millions, or even hundreds of billions, of adjustable parameters $\mathbf{w} = [w_1, w_2, \dots, w_d]^T \in \mathbb{R}^d$. Setting the vector derivative to zero produces a system of $d$ simultaneous coupled equations:
    \begin{equation}
    \nabla R_S(\mathbf{w}) = \begin{bmatrix} \frac{\partial R_S}{\partial w_1} \\ \frac{\partial R_S}{\partial w_2} \\ \vdots \\ \frac{\partial R_S}{\partial w_d} \end{bmatrix} = \begin{bmatrix} 0 \\ 0 \\ \vdots \\ 0 \end{bmatrix}
    \end{equation}
    Even for the simplest case of linear regression with Mean Squared Error, solving $\nabla R_S(\mathbf{w}) = \mathbf{0}$ yields the classic normal equations $(X^T X)\mathbf{w} = X^T \mathbf{y}$. Solving this linear system directly via Gaussian elimination or matrix inversion requires $\mathcal{O}(d^3)$ floating-point operations. If $d = 100{,}000$, direct matrix inversion would require roughly $10^{15}$ operations—demanding days of supercomputer time for a single model!
    
    \item \textbf{Transcendental and Non-Convex Equations:} In logistic regression, support vector machines with soft margins, and deep neural networks, the equations $\nabla R_S(\mathbf{w}) = \mathbf{0}$ involve transcendental functions (exponentials, logarithms, sigmoids) and nested non-linear compositions. No closed-form algebraic solution exists, even in principle.
\end{enumerate}

\subsection{Learning as Physical Navigation: The Blindfolded Hiker}
Because we cannot jump directly to the bottom of the valley algebraically, we are forced to \textbf{walk there iteratively}. 

Imagine you are a blindfolded hiker dropped onto an unfamiliar, fog-shrouded mountain range at twilight. You cannot see the landscape. You cannot see where the lowest valley lies. How do you find the bottom?
\begin{itemize}
    \item You reach down with your boots and feel the local slope of the ground beneath your feet.
    \item You find the direction where the earth slopes down most steeply.
    \item You take a careful, measured step in that downward direction.
    \item You stop, feel the slope anew at your new position, and take another step.
    \item You repeat this cycle until the ground beneath your boots feels completely flat in all directions.
\end{itemize}
This physical procedure is the essence of \textbf{Gradient Descent}. Multivariate differential calculus is the mathematical language that gives our blindfolded hiker an infallible local compass:
\begin{itemize}
    \item The \textbf{Gradient Vector} $\nabla f(\mathbf{w})$ computes the direction of steepest uphill ascent. Stepping in the opposite direction, $-\nabla f(\mathbf{w})$, guarantees local downhill progress.
    \item The \textbf{Hessian Matrix} $\nabla^2 f(\mathbf{w})$ quantifies the curvature of the terrain in every direction, revealing whether the hiker is standing inside a stable bowl, traversing an elongated ravine, perched on a knife-edge ridge, or stalled at a treacherous saddle point.
    \item The \textbf{Taylor Series} constructs a mathematical telescope that approximates any curving, complex loss landscape locally by a flat plane (first order) or a parabolic bowl (second order).
\end{itemize}

\subsection{The High-Dimensional Saddle Dilemma}
In single-variable calculus, when the first derivative vanishes ($f'(w) = 0$), the stationary point is almost always a local minimum or a local maximum (inflection points like $f(w) = w^3$ are rare exceptions).

In high-dimensional machine learning, this intuition is completely false:

\begin{ponderbox}{The High-Dimensional Saddle Dilemma: Why Minima are Vanishingly Rare}
Suppose you are training a machine learning model over a 100-dimensional parameter space ($d = 100$). An optimization algorithm reaches a stationary point $\mathbf{w}^*$ where the slope along every axis is zero: $\|\nabla R_S(\mathbf{w}^*)\|_2 = 0$.

Suppose that along each of the 100 independent coordinate directions, the ground has an equal $50\%$ chance of curving upward (like a valley) or curving downward (like a ridge). What is the probability that this stationary point is a true local minimum?
\tcblower
\textbf{Reflection:} For $\mathbf{w}^*$ to be a true local minimum, the terrain must curve upward along \textbf{every single one of the 100 independent directions simultaneously}!
\begin{equation}
P(\text{Local Minimum}) = \left(\frac{1}{2}\right)^{100} = \frac{1}{1{,}267{,}650{,}600{,}228{,}229{,}401{,}496{,}703{,}205{,}376} \approx 7.89 \times 10^{-31}
\end{equation}
The probability is virtually zero! If even a single independent direction curves downward, the hiker is perched on a \textbf{saddle point}: stepping along that downward direction allows the loss to keep decreasing. 

In high-dimensional spaces, almost all stationary points encountered by algorithms are saddle points, not local minima. Navigating high-dimensional terrain requires understanding the mathematics of multi-directional curvature.
\end{ponderbox}

\subsection{Prerequisites \& Notation Conventions}
Mastery of this chapter builds upon foundational concepts outlined in Table~\ref{tab:w02_prereqs} and Table~\ref{tab:w02_notation}.

\begin{table}[htbp]
\centering
\small
\begin{tabularx}{\textwidth}{lXX}
\toprule
\textbf{Foundational Area} & \textbf{Core Concepts Required} & \textbf{Operational Role in Differential Calculus} \\
\midrule
Single-Variable Calculus & Limits, power/chain rules, univariate Taylor polynomials & Baseline rates of change along 1D parameterized paths $\mathbf{x}(t)$. \\
\addlinespace
Vector Geometry & Dot products $\langle \mathbf{u}, \mathbf{v} \rangle = \mathbf{u}^T \mathbf{v}$, norms $\|\mathbf{v}\|_2$, unit directions & Projecting rates of change; defining directional derivatives and hyperplanes. \\
\addlinespace
Matrix Algebra & Matrix-vector products $A\mathbf{x}$, transpose $A^T$, symmetry $A = A^T$ & Representing Jacobians of vector transformations and quadratic forms $\mathbf{x}^T A \mathbf{x}$. \\
\bottomrule
\end{tabularx}
\caption{Prerequisite Scaffolding for Chapter 2.}
\label{tab:w02_prereqs}
\end{table}

\begin{table}[htbp]
\centering
\small
\begin{tabularx}{\textwidth}{llX}
\toprule
\textbf{Symbol} & \textbf{Formal Name} & \textbf{Intuitive Plain-English Meaning} \\
\midrule
$f: \mathbb{R}^d \to \mathbb{R}$ & Scalar Field / Loss Function & A real-valued landscape mapping $d$ model parameters to a single scalar loss penalty. \\
$\mathbf{f}: \mathbb{R}^n \to \mathbb{R}^m$ & Vector Field / Layer Map & A multi-output transformation mapping an $n$-dimensional input to an $m$-dimensional output. \\
$\nabla f(\mathbf{x}) \in \mathbb{R}^d$ & Gradient Vector & Column vector of all first partial derivatives; points in the direction of steepest uphill ascent. \\
$D_{\mathbf{u}} f(\mathbf{x}) \in \mathbb{R}$ & Directional Derivative & Scalar rate of change of $f$ when taking an infinitesimal step along unit vector $\mathbf{u}$. \\
$J_{\mathbf{f}}(\mathbf{x}) \in \mathbb{R}^{m \times n}$ & Jacobian Matrix & Matrix of all first-order partials; best local linear approximation of vector map $\mathbf{f}$. \\
$\nabla^2 f(\mathbf{x}) \in \mathbb{R}^{d \times d}$ & Hessian Matrix $H(\mathbf{x})$ & Symmetric matrix of second partial derivatives; quantifies local curvature in every direction. \\
$\mathbf{x}^T A \mathbf{x}$ & Quadratic Form & Energy/loss term scaling quadratically with vector $\mathbf{x}$; governed by matrix $A$. \\
\bottomrule
\end{tabularx}
\caption{Formal Mathematical Notation for Differential Calculus in Machine Learning.}
\label{tab:w02_notation}
\end{table}

% ----------------------------------------------------------------------
% PART 2: MATHEMATICAL & COMPUTATIONAL FOUNDATIONS
% ----------------------------------------------------------------------
\section{Mathematical \& Computational Foundations}

\subsection{Multivariate Differentiability: Beyond Partial Derivatives}
In single-variable calculus, differentiability simply means that the limit of the difference quotient exists. In multivariable space ($\mathbb{R}^d$), a function can be approached from infinitely many paths and directions.

\begin{definitionbox}{Multivariate Differentiability (Fréchet Form)}
A scalar function $f: \mathbb{R}^d \to \mathbb{R}$ is \textbf{differentiable} at a point $\mathbf{x} \in \mathbb{R}^d$ if there exists a vector $\mathbf{g} \in \mathbb{R}^d$ such that for any perturbation vector $\mathbf{v} \in \mathbb{R}^d$:
\begin{equation}
\lim_{\mathbf{v} \to \mathbf{0}} \frac{\big| f(\mathbf{x} + \mathbf{v}) - f(\mathbf{x}) - \mathbf{g}^T \mathbf{v} \big|}{\|\mathbf{v}\|_2} = 0
\label{eq:frechet_diff}
\end{equation}
When such a vector $\mathbf{g}$ exists, it is unique and defined as the \textbf{gradient} of $f$ at $\mathbf{x}$, denoted:
\begin{equation}
\nabla f(\mathbf{x}) \triangleq \mathbf{g} = \begin{bmatrix} \frac{\partial f}{\partial x_1} \\ \frac{\partial f}{\partial x_2} \\ \vdots \\ \frac{\partial f}{\partial x_d} \end{bmatrix} \in \mathbb{R}^d
\end{equation}
\end{definitionbox}

Equation~\eqref{eq:frechet_diff} carries an essential geometric meaning: it states that near $\mathbf{x}$, the curving function $f(\mathbf{x} + \mathbf{v})$ is extraordinarily well-approximated by the \textbf{affine tangent hyperplane}:
\begin{equation}
T_1(\mathbf{x} + \mathbf{v}) = f(\mathbf{x}) + \nabla f(\mathbf{x})^T \mathbf{v}
\end{equation}
The approximation error vanishes strictly faster than the step size $\|\mathbf{v}\|_2$ itself:
\begin{equation}
f(\mathbf{x} + \mathbf{v}) = f(\mathbf{x}) + \nabla f(\mathbf{x})^T \mathbf{v} + o(\|\mathbf{v}\|_2)
\end{equation}

\begin{ponderbox}{The Multi-Path Trap: Do Partial Derivatives Guarantee Differentiability?}
Suppose an objective function $f(x, y)$ has well-defined, finite partial derivatives $\frac{\partial f}{\partial x}(0, 0)$ and $\frac{\partial f}{\partial y}(0, 0)$ along both coordinate axes. Does this guarantee that $f$ is continuous or differentiable at the origin?
\tcblower
\textbf{Reflection:} \textbf{No!} Consider the classic pathological function:
\begin{equation*}
f(x, y) = \begin{cases} \frac{xy}{x^2 + y^2} & \text{if } (x, y) \ne (0, 0) \\ 0 & \text{if } (x, y) = (0, 0) \end{cases}
\end{equation*}
Along the $x$-axis ($y = 0$), $f(x, 0) = 0 \implies \frac{\partial f}{\partial x}(0, 0) = 0$. Along the $y$-axis ($x = 0$), $f(0, y) = 0 \implies \frac{\partial f}{\partial y}(0, 0) = 0$. Both coordinate partial derivatives exist and equal zero!

Yet, if you walk toward the origin along the diagonal line $y = x$, the function value is constant at $f(x, x) = \frac{x^2}{2x^2} = \frac{1}{2}$. The limit along the diagonal is $1/2$, while the limit along the coordinate axes is $0$. The function is not even continuous, let alone differentiable!

\textbf{Takeaway for ML:} For a loss function to be differentiable, its partial derivatives must not merely exist; they must be \textbf{continuous} in an open neighborhood around $\mathbf{x}$ ($f \in C^1$).
\end{ponderbox}

\subsection{Directional Derivatives and the Cauchy-Schwarz Steepest Descent Proof}
What is the rate of change of a loss function if we move along an arbitrary direction $\mathbf{u}$ that does not align with the coordinate axes?

\begin{definitionbox}{Directional Derivative}
Let $\mathbf{u} \in \mathbb{R}^d$ be a unit direction vector ($\|\mathbf{u}\|_2 = 1$). The \textbf{directional derivative} of $f$ at point $\mathbf{x}$ along direction $\mathbf{u}$ is defined as the 1D limit:
\begin{equation}
D_{\mathbf{u}} f(\mathbf{x}) \triangleq \lim_{t \to 0} \frac{f(\mathbf{x} + t\mathbf{u}) - f(\mathbf{x})}{t}
\label{eq:dir_deriv_def}
\end{equation}
For any differentiable function $f$, the directional derivative evaluates directly to the inner product:
\begin{equation}
D_{\mathbf{u}} f(\mathbf{x}) = \nabla f(\mathbf{x})^T \mathbf{u} = \langle \nabla f(\mathbf{x}), \mathbf{u} \rangle = \|\nabla f(\mathbf{x})\|_2 \cos(\theta)
\label{eq:dir_deriv_inner_prod}
\end{equation}
where $\theta$ is the angle between the gradient vector $\nabla f(\mathbf{x})$ and the search direction $\mathbf{u}$.
\end{definitionbox}

Equation~\eqref{eq:dir_deriv_inner_prod} immediately unlocks the mathematical justification for Gradient Descent:

\begin{theorembox}{The Steepest Descent Property}
Let $f: \mathbb{R}^d \to \mathbb{R}$ be differentiable at $\mathbf{x}$ with non-zero gradient ($\nabla f(\mathbf{x}) \ne \mathbf{0}$). Among all possible unit directions $\mathbf{u} \in \mathbb{R}^d$ ($\|\mathbf{u}\|_2 = 1$):
\begin{enumerate}
    \item The direction of \textbf{maximal rate of increase} (steepest ascent) is the normalized gradient:
    \begin{equation}
    \mathbf{u}^*_{\text{ascent}} = +\frac{\nabla f(\mathbf{x})}{\|\nabla f(\mathbf{x})\|_2}, \quad \text{achieving } D_{\mathbf{u}^*} f(\mathbf{x}) = +\|\nabla f(\mathbf{x})\|_2
    \end{equation}
    \item The direction of \textbf{maximal rate of decrease} (steepest descent) is the negative normalized gradient:
    \begin{equation}
    \mathbf{u}^*_{\text{descent}} = -\frac{\nabla f(\mathbf{x})}{\|\nabla f(\mathbf{x})\|_2}, \quad \text{achieving } D_{\mathbf{u}^*} f(\mathbf{x}) = -\|\nabla f(\mathbf{x})\|_2
    \end{equation}
    \item Any direction $\mathbf{u}_{\text{tangent}}$ orthogonal to $\nabla f(\mathbf{x})$ ($\langle \nabla f(\mathbf{x}), \mathbf{u} \rangle = 0$) has \textbf{zero instantaneous rate of change}.
\end{enumerate}
\end{theorembox}

\begin{proof}
By Equation~\eqref{eq:dir_deriv_inner_prod} and the Cauchy-Schwarz inequality ($|\langle \mathbf{a}, \mathbf{b} \rangle| \le \|\mathbf{a}\|_2 \|\mathbf{b}\|_2$):
\begin{equation}
-\|\nabla f(\mathbf{x})\|_2 \|\mathbf{u}\|_2 \le \langle \nabla f(\mathbf{x}), \mathbf{u} \rangle \le \|\nabla f(\mathbf{x})\|_2 \|\mathbf{u}\|_2
\end{equation}
Since $\mathbf{u}$ is a unit vector ($\|\mathbf{u}\|_2 = 1$):
\begin{equation}
-\|\nabla f(\mathbf{x})\|_2 \le D_{\mathbf{u}} f(\mathbf{x}) \le +\|\nabla f(\mathbf{x})\|_2
\end{equation}
Equality on the right-hand side holds if and only if $\mathbf{u}$ is positively collinear with $\nabla f(\mathbf{x})$ ($\theta = 0^\circ \implies \cos(0) = 1$). Since $\|\mathbf{u}\|_2 = 1$, $\mathbf{u}^*_{\text{ascent}} = \frac{\nabla f(\mathbf{x})}{\|\nabla f(\mathbf{x})\|_2}$.

Equality on the left-hand side holds if and only if $\mathbf{u}$ is anti-parallel to $\nabla f(\mathbf{x})$ ($\theta = 180^\circ \implies \cos(\pi) = -1$). Hence, $\mathbf{u}^*_{\text{descent}} = -\frac{\nabla f(\mathbf{x})}{\|\nabla f(\mathbf{x})\|_2}$.
\end{proof}

\begin{ponderbox}{Why is the Gradient Always Perpendicular to Level Sets (Contours)?}
A level curve of a loss function is defined by $f(\mathbf{x}) = c$ (a curve of constant elevation). Why must the gradient vector $\nabla f(\mathbf{x})$ point strictly perpendicular (normal) to this level curve, rather than at an angle?
\tcblower
\textbf{Reflection:} Imagine a path $\mathbf{x}(t)$ that travels strictly along a contour line of constant elevation $c$, so that $f(\mathbf{x}(t)) = c$ for all $t$. The velocity vector $\mathbf{x}'(t)$ represents the tangent direction of the contour.
Differentiating both sides with respect to $t$ using the multivariate chain rule:
\begin{equation*}
\frac{d}{dt} \big[ f(\mathbf{x}(t)) \big] = \frac{d}{dt}[c] \implies \nabla f(\mathbf{x}(t))^T \mathbf{x}'(t) = 0
\end{equation*}
Because their inner product is identically zero, the gradient vector $\nabla f(\mathbf{x})$ is mathematically \textbf{orthogonal} to the contour tangent $\mathbf{x}'(t)$. To change your elevation as rapidly as possible, you must step directly perpendicular to the lines of equal altitude!
\end{ponderbox}

\begin{ponderbox}{The Step-Size Dilemma: If $-\nabla f$ is Steepest Descent, Why Can't We Step Infinitely Far?}
If the negative gradient $-\nabla f(\mathbf{x})$ is proven to be the direction of maximal decrease, why can't we simply set our learning rate $\eta = 1{,}000$ and reach the global minimum in a single giant leap?
\tcblower
\textbf{Reflection:} The steepest descent property is an \textbf{infinitesimal local guarantee}. It holds in the limit as step size $t \to 0$.

When we take a finite, discrete step $\mathbf{v} = -\eta \nabla f(\mathbf{x})$, the second-order Taylor expansion reveals the true change in loss:
\begin{equation*}
f(\mathbf{x} - \eta \nabla f(\mathbf{x})) \approx f(\mathbf{x}) - \eta \|\nabla f(\mathbf{x})\|_2^2 + \frac{\eta^2}{2} \nabla f(\mathbf{x})^T \nabla^2 f(\mathbf{x}) \nabla f(\mathbf{x})
\end{equation*}
Notice the competition between the two terms:
\begin{itemize}
    \item The first-order linear term $-\eta \|\nabla f(\mathbf{x})\|_2^2$ is negative and pulls the loss \textbf{downward}.
    \item The second-order curvature term $+\frac{\eta^2}{2} \nabla f^T H \nabla f$ scales with $\eta^2$. If the ground curves upward ($H \succ 0$), this term pushes the loss \textbf{upward}!
\end{itemize}
If $\eta$ is chosen too large, the quadratic penalty $\mathcal{O}(\eta^2)$ overpowers the linear descent $\mathcal{O}(\eta)$. The hiker overshoots the valley floor completely and lands higher up the opposite cliff, causing explosive divergence! The curvature of the terrain dictates how far we dare step.
\end{ponderbox}

\subsection{The Matrix Calculus Reference Toolkit \& Layout Conventions}
In machine learning research and engineering, we rarely compute partial derivatives with respect to scalar coordinates $w_1, w_2, \dots, w_d$. Instead, we differentiate directly with respect to parameter vectors $\mathbf{w} \in \mathbb{R}^d$ and weight matrices $W \in \mathbb{R}^{m \times n}$.

\subsubsection{Layout Conventions: Numerator vs. Denominator}
In multivariable literature, two conventions exist for organizing derivatives:
\begin{itemize}
    \item \textbf{Numerator Layout (Jacobian formulation):} The derivative of a scalar with respect to a column vector $\mathbf{x} \in \mathbb{R}^d$ is formatted as a \textbf{row vector} of size $1 \times d$.
    \item \textbf{Denominator Layout (Machine Learning formulation):} The derivative of a scalar with respect to a column vector $\mathbf{x} \in \mathbb{R}^d$ is formatted as a \textbf{column vector} of size $d \times 1$.
\end{itemize}
\textbf{Standard Adoption:} In this textbook and across modern machine learning (including PyTorch, JAX, and optimization theory), we strictly adopt the \textbf{Denominator Layout}:
\begin{equation}
\nabla_{\mathbf{x}} f(\mathbf{x}) \triangleq \begin{bmatrix} \frac{\partial f}{\partial x_1} \\ \frac{\partial f}{\partial x_2} \\ \vdots \\ \frac{\partial f}{\partial x_d} \end{bmatrix} \in \mathbb{R}^{d \times 1} \quad \text{for } \mathbf{x} \in \mathbb{R}^{d \times 1}
\end{equation}
This ensures that the gradient $\nabla_{\mathbf{x}} f$ has the exact same shape as the parameter vector $\mathbf{x}$, allowing direct parameter update arithmetic: $\mathbf{x}_{t+1} = \mathbf{x}_t - \eta \nabla f(\mathbf{x}_t)$.

\subsubsection{Master Vector \& Matrix Derivative Identities}
Table~\ref{tab:matrix_calculus_toolkit} provides the indispensable matrix calculus identities required for all machine learning models.

\begin{table}[htbp]
\centering
\small
\begin{tabularx}{\textwidth}{llX}
\toprule
\textbf{Objective Expression} & \textbf{Gradient $\nabla_{\mathbf{x}} f(\mathbf{x})$} & \textbf{Mathematical Conditions / Derivation Note} \\
\midrule
$f(\mathbf{x}) = \mathbf{a}^T \mathbf{x} = \mathbf{x}^T \mathbf{a}$ & $\mathbf{a}$ & Linear form; constant slope along vector $\mathbf{a}$. \\
\addlinespace
$f(\mathbf{x}) = \mathbf{x}^T \mathbf{x} = \|\mathbf{x}\|_2^2$ & $2\mathbf{x}$ & Squared Euclidean norm ($L_2$ regularization penalty). \\
\addlinespace
$f(\mathbf{x}) = \mathbf{x}^T A \mathbf{x}$ & $(A + A^T)\mathbf{x}$ & General quadratic form; $A \in \mathbb{R}^{d \times d}$ arbitrary square matrix. \\
\addlinespace
$f(\mathbf{x}) = \mathbf{x}^T A \mathbf{x}$ & $2A\mathbf{x}$ & Symmetric quadratic form; holds if and only if $A = A^T$. \\
\addlinespace
$f(\mathbf{x}) = \|A\mathbf{x} - \mathbf{b}\|_2^2$ & $2A^T(A\mathbf{x} - \mathbf{b})$ & Ordinary least squares residual; chain rule on Euclidean norm. \\
\addlinespace
$f(\mathbf{x}) = \mathbf{a}^T g(\mathbf{x})$ & $J_g(\mathbf{x})^T \mathbf{a}$ & Linear combination of vector mapping $g: \mathbb{R}^d \to \mathbb{R}^m$. \\
\addlinespace
$f(X) = \operatorname{tr}(A X)$ & $A^T$ & Matrix trace derivative with respect to matrix $X \in \mathbb{R}^{m \times n}$. \\
\bottomrule
\end{tabularx}
\caption{The Master Matrix Calculus Toolkit for Machine Learning.}
\label{tab:matrix_calculus_toolkit}
\end{table}

\begin{theorembox}{Derivation of the Quadratic Form Gradient}
Let $A \in \mathbb{R}^{d \times d}$ be an arbitrary real square matrix (not necessarily symmetric). For the scalar quadratic form $f(\mathbf{x}) = \mathbf{x}^T A \mathbf{x}$, the gradient is:
\begin{equation}
\nabla_{\mathbf{x}} (\mathbf{x}^T A \mathbf{x}) = (A + A^T)\mathbf{x}
\label{eq:quad_form_deriv}
\end{equation}
\end{theorembox}

\begin{proof}
Expand the matrix multiplication into explicit scalar summation indices:
\begin{equation}
f(\mathbf{x}) = \mathbf{x}^T (A\mathbf{x}) = \sum_{i=1}^d x_i \left( \sum_{j=1}^d A_{ij} x_j \right) = \sum_{i=1}^d \sum_{j=1}^d A_{ij} x_i x_j
\end{equation}
To compute the $k$-th coordinate of the gradient, differentiate $f(\mathbf{x})$ with respect to a single variable $x_k$:
\begin{equation}
\frac{\partial f}{\partial x_k} = \frac{\partial}{\partial x_k} \left[ \sum_{i=1}^d \sum_{j=1}^d A_{ij} x_i x_j \right]
\end{equation}
The variable $x_k$ appears in this double sum in three distinct cases:
\begin{enumerate}
    \item When $i = k$ and $j \ne k$: terms of the form $A_{kj} x_k x_j$. Differentiating yields $\sum_{j \ne k} A_{kj} x_j$.
    \item When $j = k$ and $i \ne k$: terms of the form $A_{ik} x_i x_k$. Differentiating yields $\sum_{i \ne k} A_{ik} x_i$.
    \item When $i = k$ and $j = k$: the single diagonal term $A_{kk} x_k^2$. Differentiating yields $2A_{kk} x_k$.
\end{enumerate}
Combining all three cases:
\begin{equation}
\frac{\partial f}{\partial x_k} = \sum_{j=1}^d A_{kj} x_j + \sum_{i=1}^d A_{ik} x_i
\end{equation}
Notice the vector structure:
\begin{itemize}
    \item $\sum_{j=1}^d A_{kj} x_j$ is the $k$-th element of the matrix-vector product $A\mathbf{x}$.
    \item $\sum_{i=1}^d A_{ik} x_i$ is the $k$-th element of the matrix-vector product $A^T \mathbf{x}$.
\end{itemize}
Stacking these coordinate partial derivatives into a column vector:
\begin{equation}
\nabla_{\mathbf{x}} f(\mathbf{x}) = \begin{bmatrix} \frac{\partial f}{\partial x_1} \\ \vdots \\ \frac{\partial f}{\partial x_d} \end{bmatrix} = A\mathbf{x} + A^T \mathbf{x} = (A + A^T)\mathbf{x}
\end{equation}
If $A$ is symmetric ($A = A^T$), then $(A + A^T)\mathbf{x} = 2A\mathbf{x}$.
\end{proof}

\subsection{The Jacobian Matrix and the Multivariate Chain Rule}
When we move from scalar objective functions to multi-output transformations (such as hidden layers in a neural network, or coordinate changes), we evaluate vector-valued functions $\mathbf{f}: \mathbb{R}^n \to \mathbb{R}^m$:
\begin{equation}
\mathbf{f}(\mathbf{x}) = \begin{bmatrix} f_1(x_1, \dots, x_n) \\ f_2(x_1, \dots, x_n) \\ \vdots \\ f_m(x_1, \dots, x_n) \end{bmatrix}
\end{equation}

\begin{definitionbox}{The Jacobian Matrix}
The \textbf{Jacobian Matrix} $J_{\mathbf{f}}(\mathbf{x}) \in \mathbb{R}^{m \times n}$ collects all $m \times n$ first-order partial derivatives of vector mapping $\mathbf{f}$:
\begin{equation}
J_{\mathbf{f}}(\mathbf{x}) \triangleq \begin{bmatrix}
\frac{\partial f_1}{\partial x_1} & \frac{\partial f_1}{\partial x_2} & \dots & \frac{\partial f_1}{\partial x_n} \\
\frac{\partial f_2}{\partial x_1} & \frac{\partial f_2}{\partial x_2} & \dots & \frac{\partial f_2}{\partial x_n} \\
\vdots \\
\frac{\partial f_m}{\partial x_1} & \frac{\partial f_m}{\partial x_2} & \dots & \frac{\partial f_m}{\partial x_n}
\end{bmatrix} = \begin{bmatrix}
\nabla f_1(\mathbf{x})^T \\
\nabla f_2(\mathbf{x})^T \\
\vdots \\
\nabla f_m(\mathbf{x})^T
\end{bmatrix}
\label{eq:jacobian_def}
\end{equation}
Each \textbf{row} of the Jacobian is the transposed gradient of one output coordinate. Each \textbf{column} represents the sensitivity of the entire output vector to one input coordinate.
\end{definitionbox}

\begin{theorembox}{The Multivariate Chain Rule}
Let $\mathbf{g}: \mathbb{R}^k \to \mathbb{R}^n$ and $\mathbf{f}: \mathbb{R}^n \to \mathbb{R}^m$ be differentiable mappings. For the composite function $\mathbf{h} = \mathbf{f} \circ \mathbf{g}: \mathbb{R}^k \to \mathbb{R}^m$ defined by $\mathbf{h}(\mathbf{x}) = \mathbf{f}(\mathbf{g}(\mathbf{x}))$, the Jacobian is the clean matrix product:
\begin{equation}
J_{\mathbf{h}}(\mathbf{x}) = J_{\mathbf{f}}(\mathbf{g}(\mathbf{x})) \cdot J_{\mathbf{g}}(\mathbf{x})
\label{eq:chain_rule_matrix_prod}
\end{equation}
\end{theorembox}

Equation~\eqref{eq:chain_rule_matrix_prod} is the mathematical engine behind \textbf{Backpropagation} in deep neural networks. In a deep network with $L$ layers:
\begin{equation}
\mathbf{x} \xrightarrow{\quad\mathbf{f}_1\quad} \mathbf{h}_1 \xrightarrow{\quad\mathbf{f}_2\quad} \mathbf{h}_2 \dots \xrightarrow{\quad\mathbf{f}_L\quad} \hat{\mathbf{y}} \xrightarrow{\quad\ell\quad} \text{Loss}
\end{equation}
The total derivative of the scalar loss with respect to the input weights is simply a cascade of matrix-matrix multiplications of individual layer Jacobians!

\subsection{The Hessian Matrix and Second-Order Taylor Approximations}
First derivatives tell us the slope; second derivatives tell us the \textbf{curvature}.

\begin{definitionbox}{The Hessian Matrix}
For a scalar function $f: \mathbb{R}^d \to \mathbb{R}$ that is twice differentiable ($f \in C^2$), the \textbf{Hessian Matrix} $H(\mathbf{x}) = \nabla^2 f(\mathbf{x}) \in \mathbb{R}^{d \times d}$ collects all second-order partial derivatives:
\begin{equation}
H(\mathbf{x}) \triangleq \begin{bmatrix}
\frac{\partial^2 f}{\partial x_1^2} & \frac{\partial^2 f}{\partial x_1 \partial x_2} & \dots & \frac{\partial^2 f}{\partial x_1 \partial x_d} \\
\frac{\partial^2 f}{\partial x_2 \partial x_1} & \frac{\partial^2 f}{\partial x_2^2} & \dots & \frac{\partial^2 f}{\partial x_2 \partial x_d} \\
\vdots \\
\frac{\partial^2 f}{\partial x_d \partial x_1} & \frac{\partial^2 f}{\partial x_d \partial x_2} & \dots & \frac{\partial^2 f}{\partial x_d^2}
\end{bmatrix} \in \mathbb{R}^{d \times d}
\label{eq:hessian_def}
\end{equation}
\end{definitionbox}

\begin{theorembox}{Symmetry of the Hessian (Clairaut's / Schwarz's Theorem)}
If all second-order partial derivatives of $f: \mathbb{R}^d \to \mathbb{R}$ are continuous in an open neighborhood around $\mathbf{x}$, the order of differentiation does not matter:
\begin{equation}
\frac{\partial^2 f}{\partial x_i \partial x_j} = \frac{\partial^2 f}{\partial x_j \partial x_i} \quad \text{for all } i, j
\end{equation}
Consequently, the Hessian matrix is strictly \textbf{symmetric}: $H(\mathbf{x}) = H(\mathbf{x})^T$.
\end{theorembox}

\subsubsection{Multivariate Taylor Polynomials}
Using the gradient and Hessian, we construct the local polynomial approximations of a loss function $f$ around a base point $\mathbf{x}_0$:
\begin{enumerate}
    \item \textbf{0th-Order Approximation (Constant):} $T_0(\mathbf{x}_0 + \mathbf{v}) = f(\mathbf{x}_0)$.
    \item \textbf{1st-Order Approximation (Tangent Hyperplane):}
    \begin{equation}
    T_1(\mathbf{x}_0 + \mathbf{v}) = f(\mathbf{x}_0) + \nabla f(\mathbf{x}_0)^T \mathbf{v}
    \end{equation}
    \item \textbf{2nd-Order Approximation (Quadratic Bowl):}
    \begin{equation}
    T_2(\mathbf{x}_0 + \mathbf{v}) = f(\mathbf{x}_0) + \nabla f(\mathbf{x}_0)^T \mathbf{v} + \frac{1}{2} \mathbf{v}^T \nabla^2 f(\mathbf{x}_0) \mathbf{v}
    \label{eq:second_order_taylor}
    \end{equation}
\end{enumerate}

\subsection{Classification of Critical Points via Hessian Eigenvalues}
At a \textbf{stationary point} (or critical point) $\mathbf{x}^*$, the first-order slope vanishes: $\nabla f(\mathbf{x}^*) = \mathbf{0}$. 

Plugging $\nabla f(\mathbf{x}^*) = \mathbf{0}$ into the Taylor expansion~\eqref{eq:second_order_taylor}, the change in loss for a small displacement $\mathbf{v}$ is governed entirely by the quadratic form:
\begin{equation}
f(\mathbf{x}^* + \mathbf{v}) - f(\mathbf{x}^*) \approx \frac{1}{2} \mathbf{v}^T H(\mathbf{x}^*) \mathbf{v}
\end{equation}
Because $H(\mathbf{x}^*)$ is real and symmetric, the Spectral Theorem (which we prove in Week 6) guarantees that $H$ possesses $d$ real eigenvalues $\lambda_1, \lambda_2, \dots, \lambda_d$ with an orthonormal basis of eigenvectors $\mathbf{q}_1, \dots, \mathbf{q}_d$.

\begin{theorembox}{The Second-Order Critical Point Classification}
Let $\mathbf{x}^*$ be a stationary point ($\nabla f(\mathbf{x}^*) = \mathbf{0}$) of a $C^2$ function $f$, and let $\lambda_1, \dots, \lambda_d$ be the eigenvalues of its Hessian $H(\mathbf{x}^*)$:
\begin{enumerate}
    \item \textbf{Strict Local Minimum:} $H(\mathbf{x}^*) \succ 0$ (Strictly Positive Definite). All eigenvalues are strictly positive:
    \begin{equation}
    \lambda_i > 0 \quad \text{for all } i \in \{1, \dots, d\}
    \end{equation}
    Every non-zero perturbation increases the loss: $\mathbf{v}^T H \mathbf{v} > 0 \implies f(\mathbf{x}^* + \mathbf{v}) > f(\mathbf{x}^*)$.
    
    \item \textbf{Strict Local Maximum:} $H(\mathbf{x}^*) \prec 0$ (Strictly Negative Definite). All eigenvalues are strictly negative:
    \begin{equation}
    \lambda_i < 0 \quad \text{for all } i \in \{1, \dots, d\}
    \end{equation}
    Every non-zero perturbation decreases the loss: $\mathbf{v}^T H \mathbf{v} < 0 \implies f(\mathbf{x}^* + \mathbf{v}) < f(\mathbf{x}^*)$.

    \item \textbf{Saddle Point:} $H(\mathbf{x}^*)$ is \textbf{indefinite}. The Hessian has at least one strictly positive eigenvalue and at least one strictly negative eigenvalue:
    \begin{equation}
    \exists i, j \quad \text{such that } \lambda_i > 0 \quad \text{and} \quad \lambda_j < 0
    \end{equation}
    Perturbing along eigenvector $\mathbf{q}_i$ increases the loss, while perturbing along eigenvector $\mathbf{q}_j$ decreases it.

    \item \textbf{Inconclusive / Valley Floor:} At least one eigenvalue is zero ($\lambda_k = 0$), and all other eigenvalues have identical signs ($H \succeq 0$ or $H \preceq 0$). The second-order test is inconclusive; third- or fourth-order terms dictate stability.
\end{enumerate}
\end{theorembox}

% ----------------------------------------------------------------------
% PART 3: IN-CHAPTER STEP-BY-STEP WORKED EXAMPLES
% ----------------------------------------------------------------------
\section{In-Chapter Step-by-Step Pedagogical Examples}

\begin{examplebox}{Worked Example 2.1: Directional Derivative \& Angular Sweep}
\label{ex:w02_directional_derivative}
\textbf{The Scenario:}
Consider a two-dimensional loss function representing the cost surface of an industrial heating system:
\begin{equation*}
f(x_1, x_2) = 3x_1^2 + 2x_1 x_2 + x_2^2
\end{equation*}
An engineer is standing at the parameter point $\mathbf{x}_0 = (1, 2)$.

\textbf{Questions to Solve:}
\begin{enumerate}
    \item Calculate the exact gradient vector $\nabla f(1, 2)$.
    \item Let $\mathbf{u}(\theta) = [\cos\theta, \sin\theta]^T$ be a unit direction vector parameterized by angle $\theta \in [0, 2\pi)$. Express the directional derivative $D_{\mathbf{u}(\theta)} f(1, 2)$ as a scalar function of $\theta$.
    \item At what exact angle $\theta^*$ does the function increase most rapidly? What is this maximal rate?
    \item At what angle is the rate of change exactly zero? Confirm that this direction is tangent to the contour curve passing through $(1, 2)$.
\end{enumerate}

\tcblower
\textbf{Step-by-Step Analytical Solution:}

\textbf{Step 1: Computing the Gradient Vector}
Compute the partial derivatives:
\begin{align*}
\frac{\partial f}{\partial x_1} &= \frac{\partial}{\partial x_1}(3x_1^2 + 2x_1 x_2 + x_2^2) = 6x_1 + 2x_2 \\
\frac{\partial f}{\partial x_2} &= \frac{\partial}{\partial x_2}(3x_1^2 + 2x_1 x_2 + x_2^2) = 2x_1 + 2x_2
\end{align*}
Evaluate at $\mathbf{x}_0 = (1, 2)$:
\begin{equation*}
\nabla f(1, 2) = \begin{bmatrix} 6(1) + 2(2) \\ 2(1) + 2(2) \end{bmatrix} = \begin{bmatrix} 6 + 4 \\ 2 + 4 \end{bmatrix} = \begin{bmatrix} 10 \\ 6 \end{bmatrix}
\end{equation*}

\textbf{Step 2: Directional Derivative as a Function of $\theta$}
Using the inner product formula $D_{\mathbf{u}} f(\mathbf{x}) = \nabla f(\mathbf{x})^T \mathbf{u}$:
\begin{equation*}
D_{\mathbf{u}(\theta)} f(1, 2) = \begin{bmatrix} 10 & 6 \end{bmatrix} \begin{bmatrix} \cos\theta \\ \sin\theta \end{bmatrix} = 10\cos\theta + 6\sin\theta
\end{equation*}

\textbf{Step 3: Finding the Direction of Maximal Ascent}
To maximize $g(\theta) = 10\cos\theta + 6\sin\theta$, differentiate with respect to $\theta$ and set to zero:
\begin{equation*}
\frac{dg}{d\theta} = -10\sin\theta + 6\cos\theta = 0 \implies 10\sin\theta = 6\cos\theta \implies \tan\theta = \frac{6}{10} = 0.60
\end{equation*}
Solving for $\theta^*$:
\begin{equation*}
\theta^* = \arctan(0.60) \approx 0.5404 \text{ radians} \approx \mathbf{30.96^\circ}
\end{equation*}
Notice that $\theta^* = 30.96^\circ$ is precisely the geometric angle of the gradient vector $\nabla f = [10, 6]^T$ itself!
The maximal rate of increase achieved is:
\begin{equation*}
D_{\max} = \|\nabla f(1, 2)\|_2 = \sqrt{10^2 + 6^2} = \sqrt{100 + 36} = \sqrt{136} \approx \mathbf{11.6619}
\end{equation*}

\textbf{Step 4: Finding the Zero-Change Direction (Contour Tangent)}
The directional derivative is zero when:
\begin{equation*}
10\cos\theta + 6\sin\theta = 0 \implies \tan\theta = -\frac{10}{6} = -\frac{5}{3} \approx -1.6667
\end{equation*}
Solving for $\theta$:
\begin{equation*}
\theta_{\text{zero}} = \arctan(-5/3) \approx -59.04^\circ \quad (\text{or } 120.96^\circ)
\end{equation*}
Notice the profound geometric relationship:
\begin{equation*}
\theta^* - \theta_{\text{zero}} = 30.96^\circ - (-59.04^\circ) = \mathbf{90.00^\circ}
\end{equation*}
The zero-change direction is \textbf{strictly perpendicular} to the gradient, precisely tangent to the contour ellipse $f(x_1, x_2) = f(1, 2) = 3(1)^2 + 2(1)(2) + 2^2 = 11$.
\end{examplebox}

\begin{examplebox}{Worked Example 2.2: Matrix Calculus of Linear Regression (OLS)}
\label{ex:w02_ols_matrix_calculus}
\textbf{The Scenario:}
In Ordinary Least Squares (OLS) regression, we are given a design matrix $X \in \mathbb{R}^{N \times d}$ and a target vector $\mathbf{y} \in \mathbb{R}^N$. We wish to find weights $\mathbf{w} \in \mathbb{R}^d$ minimizing the Mean Squared Error (MSE) empirical risk:
\begin{equation*}
\mathcal{L}(\mathbf{w}) = \frac{1}{2N} \|\mathbf{y} - X\mathbf{w}\|_2^2
\end{equation*}

\textbf{Questions to Solve:}
\begin{enumerate}
    \item Expand the squared Euclidean norm using vector transpose identities.
    \item Derive the exact gradient $\nabla_{\mathbf{w}} \mathcal{L}(\mathbf{w})$ step-by-step using matrix calculus identities. Verify dimensions.
    \item Derive the exact Hessian matrix $H(\mathbf{w}) = \nabla^2_{\mathbf{w}} \mathcal{L}(\mathbf{w})$.
    \item Prove that the loss function is convex everywhere in $\mathbb{R}^d$.
\end{enumerate}

\tcblower
\textbf{Step-by-Step Analytical Solution:}

\textbf{Step 1: Expanding the Norm}
Recall that for any vector $\mathbf{v}$, $\|\mathbf{v}\|_2^2 = \mathbf{v}^T \mathbf{v}$. Let $\mathbf{e} = \mathbf{y} - X\mathbf{w}$:
\begin{align*}
\mathcal{L}(\mathbf{w}) &= \frac{1}{2N} (\mathbf{y} - X\mathbf{w})^T (\mathbf{y} - X\mathbf{w}) \\
&= \frac{1}{2N} \left( \mathbf{y}^T - \mathbf{w}^T X^T \right) (\mathbf{y} - X\mathbf{w}) \\
&= \frac{1}{2N} \left[ \mathbf{y}^T \mathbf{y} - \mathbf{y}^T X\mathbf{w} - \mathbf{w}^T X^T \mathbf{y} + \mathbf{w}^T X^T X \mathbf{w} \right]
\end{align*}
Notice that $\mathbf{y}^T X\mathbf{w}$ is a scalar ($1 \times 1$). The transpose of a scalar is itself:
\begin{equation*}
(\mathbf{y}^T X\mathbf{w})^T = \mathbf{w}^T X^T \mathbf{y}
\end{equation*}
Thus, the two middle cross-terms combine:
\begin{equation*}
\mathcal{L}(\mathbf{w}) = \frac{1}{2N} \mathbf{y}^T \mathbf{y} - \frac{1}{N} \mathbf{w}^T X^T \mathbf{y} + \frac{1}{2N} \mathbf{w}^T (X^T X) \mathbf{w}
\end{equation*}

\textbf{Step 2: Differentiating Term by Term}
We differentiate with respect to vector $\mathbf{w}$ using our Matrix Calculus Toolkit (Table~\ref{tab:matrix_calculus_toolkit}):
\begin{itemize}
    \item Term 1 ($\frac{1}{2N}\mathbf{y}^T \mathbf{y}$): Independent of $\mathbf{w} \implies \nabla_{\mathbf{w}} = \mathbf{0}$.
    \item Term 2 ($-\frac{1}{N}\mathbf{w}^T (X^T \mathbf{y})$): Linear form $\mathbf{w}^T \mathbf{a}$ where $\mathbf{a} = X^T \mathbf{y} \implies \nabla_{\mathbf{w}} = -\frac{1}{N} X^T \mathbf{y}$.
    \item Term 3 ($\frac{1}{2N}\mathbf{w}^T (X^T X)\mathbf{w}$): Quadratic form $\mathbf{w}^T A \mathbf{w}$ where $A = X^T X$. Since $A^T = (X^T X)^T = X^T X = A$, the matrix is symmetric:
    \begin{equation*}
    \nabla_{\mathbf{w}} \left( \frac{1}{2N} \mathbf{w}^T (X^T X)\mathbf{w} \right) = \frac{1}{2N} (2 X^T X \mathbf{w}) = \frac{1}{N} X^T X \mathbf{w}
    \end{equation*}
\end{itemize}
Combining terms gives the famous OLS gradient:
\begin{equation*}
\nabla_{\mathbf{w}} \mathcal{L}(\mathbf{w}) = \frac{1}{N} X^T X \mathbf{w} - \frac{1}{N} X^T \mathbf{y} = \mathbf{\frac{1}{N} X^T (X\mathbf{w} - \mathbf{y})}
\end{equation*}
\textbf{Dimension Check:} $X \in \mathbb{R}^{N \times d}$, so $X^T \in \mathbb{R}^{d \times N}$. The residual $(X\mathbf{w} - \mathbf{y}) \in \mathbb{R}^N$.
Multiplying: $(d \times N) \times (N \times 1) = d \times 1$. The gradient is indeed a $d \times 1$ vector matching $\mathbf{w}$!

\textbf{Step 3: Deriving the Hessian Matrix}
Differentiating the gradient $\nabla_{\mathbf{w}} \mathcal{L}(\mathbf{w}) = \frac{1}{N}X^T X\mathbf{w} - \frac{1}{N}X^T \mathbf{y}$ with respect to $\mathbf{w}$:
\begin{equation*}
H(\mathbf{w}) = \nabla^2_{\mathbf{w}} \mathcal{L}(\mathbf{w}) = \mathbf{\frac{1}{N} X^T X}
\end{equation*}
Notice that the Hessian is \textbf{constant} everywhere—it does not depend on $\mathbf{w}$!

\textbf{Step 4: Proving Global Convexity}
For any non-zero vector $\mathbf{v} \in \mathbb{R}^d \setminus \{\mathbf{0}\}$:
\begin{equation*}
\mathbf{v}^T H \mathbf{v} = \mathbf{v}^T \left( \frac{1}{N} X^T X \right) \mathbf{v} = \frac{1}{N} (X\mathbf{v})^T (X\mathbf{v}) = \frac{1}{N} \|X\mathbf{v}\|_2^2 \ge 0
\end{equation*}
Because the squared Euclidean norm $\|X\mathbf{v}\|_2^2$ is non-negative for any vector, $\mathbf{v}^T H \mathbf{v} \ge 0$ for all $\mathbf{v}$.
Hence, the Hessian is \textbf{positive semi-definite} ($H \succeq 0$) globally. The MSE loss surface is a convex bowl with no local traps; setting the gradient to zero yields the global minimum $(X^T X)\mathbf{w}^* = X^T \mathbf{y}$!
\end{examplebox}

\begin{examplebox}{Worked Example 2.3: Concrete Numerical OLS Step \& Inversion on Toy Data}
\label{ex:w02_numerical_ols_step}
\textbf{The Scenario:}
To demystify vector calculus and observe the numbers in action, consider a concrete toy dataset of $N = 3$ samples with a single input feature $x$, augmented with a bias column ($d = 2$ parameters $\mathbf{w} = [w_0, w_1]^T$ where $\hat{y} = w_0 + w_1 x$):
\begin{equation*}
S = \big\{ (1, 2), \; (2, 3), \; (3, 5) \big\}
\end{equation*}
The design matrix $X \in \mathbb{R}^{3 \times 2}$ and target vector $\mathbf{y} \in \mathbb{R}^3$ are:
\begin{equation*}
X = \begin{bmatrix} 1 & 1 \\ 1 & 2 \\ 1 & 3 \end{bmatrix}, \quad \mathbf{y} = \begin{bmatrix} 2 \\ 3 \\ 5 \end{bmatrix}
\end{equation*}

\textbf{Questions to Solve:}
\begin{enumerate}
    \item Compute the matrix product $X^T X$ and vector $X^T \mathbf{y}$ by hand.
    \item Invert $X^T X$ and compute the exact closed-form analytical OLS solution $\mathbf{w}^* = (X^T X)^{-1} X^T \mathbf{y}$.
    \item Initialize the weights at the origin: $\mathbf{w}_0 = [0, 0]^T$. Compute the initial loss $\mathcal{L}(\mathbf{w}_0) = \frac{1}{2N}\|\mathbf{y} - X\mathbf{w}_0\|_2^2$.
    \item Compute the exact numerical gradient $\mathbf{g}_0 = \nabla \mathcal{L}(\mathbf{w}_0)$.
    \item Take one manual Gradient Descent step with learning rate $\eta = 0.1$ to produce $\mathbf{w}_1 = \mathbf{w}_0 - \eta \mathbf{g}_0$.
    \item Compute the new loss $\mathcal{L}(\mathbf{w}_1)$ and verify that $\mathbf{w}_1$ has progressed directly toward $\mathbf{w}^*$.
\end{enumerate}

\tcblower
\textbf{Step-by-Step Analytical Solution:}

\textbf{Step 1: Computing $X^T X$ and $X^T \mathbf{y}$}
\begin{align*}
X^T X &= \begin{bmatrix} 1 & 1 & 1 \\ 1 & 2 & 3 \end{bmatrix} \begin{bmatrix} 1 & 1 \\ 1 & 2 \\ 1 & 3 \end{bmatrix} = \begin{bmatrix} 1+1+1 & 1+2+3 \\ 1+2+3 & 1+4+9 \end{bmatrix} = \begin{bmatrix} 3 & 6 \\ 6 & 14 \end{bmatrix} \\
X^T \mathbf{y} &= \begin{bmatrix} 1 & 1 & 1 \\ 1 & 2 & 3 \end{bmatrix} \begin{bmatrix} 2 \\ 3 \\ 5 \end{bmatrix} = \begin{bmatrix} 2 + 3 + 5 \\ 2(1) + 3(2) + 5(3) \end{bmatrix} = \begin{bmatrix} 10 \\ 2 + 6 + 15 \end{bmatrix} = \begin{bmatrix} 10 \\ 23 \end{bmatrix}
\end{align*}

\textbf{Step 2: Inverting $X^T X$ to Find $\mathbf{w}^*$}
The determinant of $X^T X$ is:
\begin{equation*}
\det(X^T X) = (3)(14) - (6)(6) = 42 - 36 = 6
\end{equation*}
Using the standard $2 \times 2$ matrix inversion formula:
\begin{equation*}
(X^T X)^{-1} = \frac{1}{6} \begin{bmatrix} 14 & -6 \\ -6 & 3 \end{bmatrix} = \begin{bmatrix} 7/3 & -1 \\ -1 & 1/2 \end{bmatrix}
\end{equation*}
Multiplying by $X^T \mathbf{y}$ gives the exact optimal OLS weights:
\begin{equation*}
\mathbf{w}^* = \begin{bmatrix} 7/3 & -1 \\ -1 & 1/2 \end{bmatrix} \begin{bmatrix} 10 \\ 23 \end{bmatrix} = \begin{bmatrix} \frac{70}{3} - 23 \\ -10 + \frac{23}{2} \end{bmatrix} = \begin{bmatrix} \frac{70 - 69}{3} \\ \frac{-20 + 23}{2} \end{bmatrix} = \begin{bmatrix} 1/3 \\ 3/2 \end{bmatrix} = \mathbf{\begin{bmatrix} 0.3333 \\ 1.5000 \end{bmatrix}}
\end{equation*}
The optimal fitted line is $\hat{y} = \frac{1}{3} + 1.5 x$.

\textbf{Step 3: Initial Loss at $\mathbf{w}_0 = [0, 0]^T$}
When $\mathbf{w}_0 = \mathbf{0}$, the prediction is $\hat{\mathbf{y}}_0 = X\mathbf{w}_0 = [0, 0, 0]^T$. The initial residual vector is $\mathbf{e}_0 = X\mathbf{w}_0 - \mathbf{y} = -\mathbf{y} = [-2, -3, -5]^T$.
\begin{equation*}
\mathcal{L}(\mathbf{w}_0) = \frac{1}{2(3)} \left( (-2)^2 + (-3)^2 + (-5)^2 \right) = \frac{1}{6} (4 + 9 + 25) = \frac{38}{6} \approx \mathbf{6.3333}
\end{equation*}

\textbf{Step 4: Computing the Initial Gradient $\mathbf{g}_0$}
Using the derived OLS gradient identity $\nabla \mathcal{L}(\mathbf{w}) = \frac{1}{N} X^T (X\mathbf{w} - \mathbf{y})$:
\begin{equation*}
\mathbf{g}_0 = \frac{1}{3} X^T \mathbf{e}_0 = \frac{1}{3} \begin{bmatrix} 1 & 1 & 1 \\ 1 & 2 & 3 \end{bmatrix} \begin{bmatrix} -2 \\ -3 \\ -5 \end{bmatrix} = \frac{1}{3} \begin{bmatrix} -10 \\ -23 \end{bmatrix} = \mathbf{\begin{bmatrix} -3.3333 \\ -7.6667 \end{bmatrix}}
\end{equation*}

\textbf{Step 5: Taking One Gradient Descent Step}
With learning rate $\eta = 0.1$:
\begin{equation*}
\mathbf{w}_1 = \mathbf{w}_0 - \eta \mathbf{g}_0 = \begin{bmatrix} 0 \\ 0 \end{bmatrix} - 0.1 \begin{bmatrix} -3.3333 \\ -7.6667 \end{bmatrix} = \mathbf{\begin{bmatrix} 0.3333 \\ 0.7667 \end{bmatrix}}
\end{equation*}

\textbf{Step 6: Loss Verification and Convergence Analysis}
At $\mathbf{w}_1 = [0.3333, 0.7667]^T$, the new model predictions are:
\begin{equation*}
\hat{\mathbf{y}}_1 = X\mathbf{w}_1 = \begin{bmatrix} 1 & 1 \\ 1 & 2 \\ 1 & 3 \end{bmatrix} \begin{bmatrix} 0.3333 \\ 0.7667 \end{bmatrix} = \begin{bmatrix} 1.1000 \\ 1.8667 \\ 2.6333 \end{bmatrix}
\end{equation*}
The new residuals $\mathbf{e}_1 = \hat{\mathbf{y}}_1 - \mathbf{y}$ are $[-0.9000, -1.1333, -2.3667]^T$.
The updated loss evaluates to:
\begin{equation*}
\mathcal{L}(\mathbf{w}_1) = \frac{1}{6} \left( (-0.9000)^2 + (-1.1333)^2 + (-2.3667)^2 \right) = \frac{1}{6}(0.8100 + 1.2844 + 5.6013) = \frac{7.6957}{6} \approx \mathbf{1.2826}
\end{equation*}
\textbf{The Verdict:} In a single step of gradient descent, the loss plummeted from $\mathbf{6.3333}$ to $\mathbf{1.2826}$ (a \textbf{79.7\% reduction} in error)! Furthermore, the weights moved from $[0, 0]^T$ to $[0.3333, 0.7667]^T$, progressing directly along the trajectory toward the analytical minimum $\mathbf{w}^* = [0.3333, 1.5000]^T$.
\end{examplebox}

\begin{examplebox}{Worked Example 2.4: Critical Point Classification on a Non-Convex Landscape}
\label{ex:w02_critical_point_classification}
\textbf{The Scenario:}
A machine learning researcher evaluates a toy non-convex loss function with two features $(x, y)$:
\begin{equation*}
f(x, y) = x^4 - 2x^2 + y^2
\end{equation*}

\textbf{Questions to Solve:}
\begin{enumerate}
    \item Find all stationary points of $f(x, y)$ in $\mathbb{R}^2$ where $\nabla f(x, y) = \mathbf{0}$.
    \item Compute the general Hessian matrix $H(x, y) = \nabla^2 f(x, y)$.
    \item Evaluate the eigenvalues of $H$ at each stationary point and classify each point as a local minimum, local maximum, or saddle point.
\end{enumerate}

\tcblower
\textbf{Step-by-Step Analytical Solution:}

\textbf{Step 1: Finding Stationary Points}
Compute the gradient vector:
\begin{equation*}
\nabla f(x, y) = \begin{bmatrix} \frac{\partial f}{\partial x} \\ \frac{\partial f}{\partial y} \end{bmatrix} = \begin{bmatrix} 4x^3 - 4x \\ 2y \end{bmatrix}
\end{equation*}
Set the gradient equal to the zero vector:
\begin{align*}
4x^3 - 4x = 0 &\implies 4x(x^2 - 1) = 0 \implies x \in \{0, +1, -1\} \\
2y = 0 &\implies y = 0
\end{align*}
Pairing the roots yields exactly \textbf{three stationary points}:
\begin{equation*}
P_1 = (0, 0), \quad P_2 = (1, 0), \quad P_3 = (-1, 0)
\end{equation*}

\textbf{Step 2: Computing the Hessian Matrix}
Compute the second partial derivatives:
\begin{equation*}
\frac{\partial^2 f}{\partial x^2} = 12x^2 - 4, \quad \frac{\partial^2 f}{\partial y^2} = 2, \quad \frac{\partial^2 f}{\partial x \partial y} = \frac{\partial^2 f}{\partial y \partial x} = 0
\end{equation*}
The Hessian matrix is:
\begin{equation*}
H(x, y) = \begin{bmatrix} 12x^2 - 4 & 0 \\ 0 & 2 \end{bmatrix}
\end{equation*}
Because the off-diagonal entries are zero, the diagonal entries are the eigenvalues!

\textbf{Step 3: Classifying Each Critical Point}
\begin{itemize}
    \item \textbf{Point $P_1 = (0, 0)$:}
    \begin{equation*}
    H(0, 0) = \begin{bmatrix} -4 & 0 \\ 0 & 2 \end{bmatrix} \implies \lambda_1 = -4, \; \lambda_2 = +2
    \end{equation*}
    Since $\lambda_1 < 0$ and $\lambda_2 > 0$, the Hessian is \textbf{indefinite}.
    \textbf{Classification:} $P_1 = (0, 0)$ is a \textbf{Saddle Point}. Moving along the $x$-axis decreases the loss, while moving along the $y$-axis increases it.
    
    \item \textbf{Point $P_2 = (1, 0)$:}
    \begin{equation*}
    H(1, 0) = \begin{bmatrix} 12(1)^2 - 4 & 0 \\ 0 & 2 \end{bmatrix} = \begin{bmatrix} 8 & 0 \\ 0 & 2 \end{bmatrix} \implies \lambda_1 = 8, \; \lambda_2 = 2
    \end{equation*}
    Both eigenvalues are strictly positive ($\lambda_1, \lambda_2 > 0$). $H(1, 0) \succ 0$ is strictly positive definite.
    \textbf{Classification:} $P_2 = (1, 0)$ is a strict \textbf{Local Minimum} with value $f(1, 0) = 1 - 2 + 0 = -1$.

    \item \textbf{Point $P_3 = (-1, 0)$:}
    \begin{equation*}
    H(-1, 0) = \begin{bmatrix} 12(-1)^2 - 4 & 0 \\ 0 & 2 \end{bmatrix} = \begin{bmatrix} 8 & 0 \\ 0 & 2 \end{bmatrix} \implies \lambda_1 = 8, \; \lambda_2 = 2
    \end{equation*}
    Identically, both eigenvalues are positive ($\lambda_1, \lambda_2 > 0$). $H(-1, 0) \succ 0$ is strictly positive definite.
    \textbf{Classification:} $P_3 = (-1, 0)$ is a strict \textbf{Local Minimum} with value $f(-1, 0) = 1 - 2 + 0 = -1$.
\end{itemize}
This example perfectly demonstrates the double-well potential: two symmetrical valleys separated by a central mountain saddle pass at the origin.
\end{examplebox}

% ----------------------------------------------------------------------
% PART 4: VISUAL & GEOMETRIC INTUITION
% ----------------------------------------------------------------------
\section{Visual \& Geometric Intuition: Contours, Trajectories, and Saddles}

\subsection{The Geometry of Contours and Tangents}
Figure~\ref{fig:contour_geometry} visualizes the fundamental geometric relationship between level sets (contours), the gradient vector, and the tangent hyperplane.

\begin{figure}[htbp]
\centering
\begin{tikzpicture}[
    scale=0.95,
    font=\sffamily\footnotesize
]

% Elliptical contour curves (level sets of quadratic function)
\draw[thick, color=gray!50] (0, 0) ellipse (4.8cm and 2.6cm);
\node[right, color=gray!80!black] at (4.8, 0) {$f(\mathbf{x}) = c_3$};

\draw[thick, color=primarydark!60] (0, 0) ellipse (3.4cm and 1.8cm);
\node[right, color=primarydark!80!black] at (3.4, 0) {$f(\mathbf{x}) = c_2$};

\draw[thick, color=primarydark] (0, 0) ellipse (1.8cm and 1.0cm);
\node[right, color=primarydark] at (1.8, 0) {$f(\mathbf{x}) = c_1 < c_2 < c_3$};

% Optimum center
\filldraw[color=tealdef] (0, 0) circle (2.5pt) node[below=3pt] {$\mathbf{x}^*$ (Global Minimum)};

% Evaluation point on contour c2
\coordinate (P) at (2.40, 1.27);
\filldraw[color=traprad] (P) circle (2.5pt) node[above right=2pt] {$\mathbf{x}_0$};

% Tangent line at P
\draw[dashed, line width=1.2pt, color=tealdef] ($(P) + (-2.2, 1.5)$) -- ($(P) + (2.2, -1.5)$);
\node[above, color=tealdef] at ($(P) + (1.6, -1.5)$) {Tangent Plane: $\nabla f(\mathbf{x}_0)^T(\mathbf{x} - \mathbf{x}_0) = 0$};

% Gradient vector (orthogonal to tangent)
\draw[->, >=stealth, line width=2.0pt, color=traprad] (P) -- ($(P) + (1.2, 1.7)$) node[right=2pt] {$\nabla f(\mathbf{x}_0)$ (Steepest Ascent)};

% Negative gradient vector (towards minimum)
\draw[->, >=stealth, line width=2.0pt, color=primarydark] (P) -- ($(P) + (-1.2, -1.7)$) node[below left=2pt] {$-\nabla f(\mathbf{x}_0)$ (Steepest Descent)};

\end{tikzpicture}
\caption{The Geometry of Contours: The gradient vector $\nabla f(\mathbf{x}_0)$ is strictly perpendicular (orthogonal) to the contour tangent line, pointing in the direction of maximal rate of increase. The negative gradient points directly downhill toward lower level sets.}
\label{fig:contour_geometry}
\end{figure}

\subsection{Optimization Dynamics: Trajectories in Anisotropic Ravines}
Figure~\ref{fig:gd_trajectories} illustrates the dynamic stepping behavior of Gradient Descent across elliptical loss contours. In high-dimensional machine learning, different parameters often have vastly different curvatures (governed by the eigenvalues $\lambda_{\max}$ and $\lambda_{\min}$ of the Hessian).

\begin{figure}[htbp]
\centering
\begin{tikzpicture}[
    scale=0.92,
    font=\sffamily\footnotesize
]

% Elliptical contour curves (elongated valley)
\draw[thick, color=gray!40] (0, 0) ellipse (5.5cm and 1.8cm);
\draw[thick, color=primarydark!50] (0, 0) ellipse (4.0cm and 1.2cm);
\draw[thick, color=primarydark!80] (0, 0) ellipse (2.5cm and 0.7cm);
\draw[thick, color=primarydark] (0, 0) ellipse (1.0cm and 0.3cm);

\filldraw[color=tealdef] (0, 0) circle (2.5pt) node[below=4pt] {$\mathbf{w}^*$ (Optimum)};

% Trajectory points showing zig-zagging in ill-conditioned valley
\coordinate (W0) at (4.8, 1.4);
\coordinate (W1) at (3.5, -0.9);
\coordinate (W2) at (2.4, 0.6);
\coordinate (W3) at (1.4, -0.4);
\coordinate (W4) at (0.6, 0.2);

\filldraw[color=traprad] (W0) circle (2.2pt) node[above right] {$\mathbf{w}_0$};
\filldraw[color=primarydark] (W1) circle (2.2pt) node[below right] {$\mathbf{w}_1$};
\filldraw[color=primarydark] (W2) circle (2.2pt) node[above] {$\mathbf{w}_2$};
\filldraw[color=primarydark] (W3) circle (2.2pt) node[below] {$\mathbf{w}_3$};
\filldraw[color=primarydark] (W4) circle (2.2pt) node[above] {$\mathbf{w}_4$};

% Arrows connecting trajectory
\draw[->, >=stealth, line width=1.5pt, color=traprad] (W0) -- (W1);
\draw[->, >=stealth, line width=1.5pt, color=traprad] (W1) -- (W2);
\draw[->, >=stealth, line width=1.5pt, color=traprad] (W2) -- (W3);
\draw[->, >=stealth, line width=1.5pt, color=traprad] (W3) -- (W4);
\draw[->, >=stealth, line width=1.5pt, dashed, color=tealdef] (W4) -- (0, 0);

\node[align=left, fill=white, draw=gray!40, rounded corners=3pt, inner sep=4pt] at (3.2, 2.2) {
\textbf{The Ravine Zig-Zag:}\\
Because $\nabla f$ is orthogonal to contours,\\
descent steps bounce violently across the steep walls\\
while creeping painfully slowly along the gentle valley floor.
};

\end{tikzpicture}
\caption{The Trajectory of Gradient Descent in an Anisotropic Ravine: When the Hessian has high condition number ($\kappa = \lambda_{\max}/\lambda_{\min} \gg 1$), the negative gradient points almost perpendicular to the long axis of the valley, causing severe oscillatory zig-zagging.}
\label{fig:gd_trajectories}
\end{figure}

\subsection{Saddle Point Geometry vs. Local Minima}
Figure~\ref{fig:saddle_vs_minimum} contrasts the geometry of a local minimum (where all directions curve upward) with a saddle point (where curvature changes sign depending on the direction of travel).

\begin{figure}[htbp]
\centering
\begin{tikzpicture}[
    scale=0.92,
    font=\sffamily\footnotesize
]

% Left: Local Minimum (Paraboloid bowl)
\begin{scope}[xshift=-4.5cm]
    \draw[thick, fill=tealbg, draw=tealborder] (0, 0) ellipse (2.8cm and 1.2cm);
    \draw[thick, fill=tealdef!15, draw=tealborder] (0, 0.8) ellipse (2.0cm and 0.8cm);
    \draw[thick, fill=tealdef!30, draw=tealborder] (0, 1.5) ellipse (1.1cm and 0.4cm);
    \filldraw[tealdef] (0, 0) circle (2.5pt) node[below=2pt] {$\mathbf{x}^*$ (Bowl Bottom)};
    
    % Curvature arrows pointing up
    \draw[->, >=stealth, thick, color=tealdef] (0, 0) -- (1.5, 0.7) node[above right] {$\lambda_1 > 0$};
    \draw[->, >=stealth, thick, color=tealdef] (0, 0) -- (-1.2, 0.6) node[above left] {$\lambda_2 > 0$};
    
    \node[align=center, font=\bfseries\sffamily, color=primarydark] at (0, -1.8) {Local Minimum\\$H \succ 0$ (Positive Definite)};
\end{scope}

% Right: Saddle Point (Hyperbolic Paraboloid)
\begin{scope}[xshift=4.5cm]
    % Saddle curve
    \draw[domain=-2.2:2.2, smooth, variable=\x, thick, color=primarydark] plot ({\x}, {0.35*\x*\x - 0.5});
    \node[color=primarydark, right] at (2.2, 1.2) {Uphill Axis ($\lambda_1 > 0$)};
    
    \draw[domain=-1.8:1.8, smooth, variable=\y, very thick, dashed, color=traprad] plot ({\y}, {-0.4*\y*\y + 0.5});
    \node[color=traprad, below] at (1.8, -0.9) {Downhill Axis ($\lambda_2 < 0$)};
    
    \filldraw[traprad] (0, 0) circle (2.5pt) node[above=4pt] {$\mathbf{x}^*$ (Saddle Point)};
    
    \node[align=center, font=\bfseries\sffamily, color=traprad] at (0, -1.8) {Saddle Point\\$H$ Indefinite ($\lambda_1 > 0, \lambda_2 < 0$)};
\end{scope}

\end{tikzpicture}
\caption{Geometric Anatomy of Critical Points: At a local minimum (left), all eigenvalues of the Hessian are positive ($\lambda_i > 0$), creating a stable bowl. At a saddle point (right), the surface curves upward along one axis ($\lambda_1 > 0$) but plunges downward along another ($\lambda_2 < 0$).}
\label{fig:saddle_vs_minimum}
\end{figure}

% ----------------------------------------------------------------------
% PART 5: EXPLICIT GOTCHAS & DIAGNOSTIC MISCONCEPTIONS
% ----------------------------------------------------------------------
\section{Explicit Gotchas \& Diagnostic Misconceptions}

\begin{trapbox}{Gotcha 1: The Vanishing Gradient Trap: Conflating Zero Gradient with Convergence}
\textbf{The Scenario:} A practitioner trains a deep neural network. After 10 epochs, the gradient norm drops to $\|\nabla \mathcal{L}(\mathbf{w})\|_2 = 10^{-6}$. The engineer halts training, declaring that the algorithm has found the optimal minimum. \\
\textbf{The Trap:} A vanishing gradient does \textbf{not} imply that the model has reached a low loss! In high-dimensional non-convex landscapes, vanishing gradients occur on vast barren plateaus, near saddle points, or in saturated activation regimes (e.g., sigmoids saturated at $0$ or $1$, or dead ReLUs). The loss may be catastrophically high while the gradient is zero. Always monitor the empirical risk $R_S(\mathbf{w})$ alongside the gradient norm.
\end{trapbox}

\begin{trapbox}{Gotcha 2: Assuming Hessian Symmetry for Non-$C^2$ Functions}
\textbf{The Scenario:} A programmer assumes that for any twice-differentiable loss function, $\frac{\partial^2 f}{\partial x \partial y} = \frac{\partial^2 f}{\partial y \partial x}$ always holds. \\
\textbf{The Trap:} Mixed partial derivatives commute only if the second partial derivatives are continuous in a neighborhood of the point (Schwarz's Theorem). On non-smooth loss functions (e.g., those involving piecewise linear ReLUs, hinge losses, or discontinuous thresholds), mixed partials can disagree, producing an asymmetric Hessian that crashes Newton-Raphson solvers.
\end{trapbox}

\begin{trapbox}{Gotcha 3: Conflating Stationary Points with Local Minima in High Dimensions}
\textbf{The Scenario:} In training a 100-parameter deep network, an optimization algorithm terminates because $\|\nabla f(\mathbf{w})\|_2 < 10^{-7}$. The engineer reports that the global or local minimum has been successfully found. \\
\textbf{The Trap:} In high-dimensional optimization surfaces, stationary points ($\nabla f = \mathbf{0}$) are overwhelmingly \textbf{saddle points}, not local minima. If the Hessian has 99 positive eigenvalues and just 1 negative eigenvalue, the point is a saddle. The optimizer has not minimized the loss; it is merely stalled on a high-dimensional plateau.
\end{trapbox}

\begin{trapbox}{Gotcha 4: Erroneous Matrix Transposition in Gradient Identities (Dimension Check)}
\textbf{The Scenario:} A student computes the gradient of the MSE loss $f(\mathbf{w}) = \|\mathbf{y} - X\mathbf{w}\|_2^2$ and writes $\nabla_{\mathbf{w}} f(\mathbf{w}) = -2X(\mathbf{y} - X\mathbf{w})$. \\
\textbf{The Trap:} Check matrix dimensions! If $X \in \mathbb{R}^{N \times d}$ and $\mathbf{w} \in \mathbb{R}^d$, then the residual $(\mathbf{y} - X\mathbf{w})$ is an $N \times 1$ vector. Multiplying $X$ ($N \times d$) by an $N \times 1$ vector is algebraically undefined! The correct gradient is $-2X^T (\mathbf{y} - X\mathbf{w}) \in \mathbb{R}^d$. Always verify that the gradient has the identical shape as the parameter vector.
\end{trapbox}

\begin{trapbox}{Gotcha 5: The Exploding Step-Size Trap ($\eta > 2/\lambda_{\max}$)}
\textbf{The Scenario:} An engineer observes that gradient descent is converging slowly and decides to increase the learning rate from $\eta = 0.01$ to $\eta = 1.0$. Immediately, the loss explodes to \texttt{NaN}. \\
\textbf{The Trap:} On a quadratic surface with Hessian $H$, gradient descent converges if and only if the step size satisfies $\eta < \frac{2}{\lambda_{\max}(H)}$, where $\lambda_{\max}$ is the largest eigenvalue (maximum curvature). If $\eta > \frac{2}{\lambda_{\max}}$, each step overshoots the valley floor so violently that the distance to the minimum increases exponentially at every iteration!
\end{trapbox}

% ----------------------------------------------------------------------
% PART 6: COMPARATIVE SYNTHESIS & SUMMARY TABLE
% ----------------------------------------------------------------------
\section{Comparative Synthesis \& Summary Table}

Table~\ref{tab:calculus_operators} synthesizes the principal differential operators employed throughout machine learning foundations.

\begin{table}[htbp]
\centering
\footnotesize
\setlength{\tabcolsep}{3pt}
\begin{tabularx}{\textwidth}{lp{2.3cm}p{2.9cm}XX}
\toprule
\textbf{Operator} & \textbf{Input / Output} & \textbf{Mathematical Definition} & \textbf{Geometric Meaning} & \textbf{Primary Use in ML} \\
\midrule
\textbf{Derivative} & $f: \mathbb{R} \to \mathbb{R}$ & $f'(x) = \lim_{h\to 0}\frac{f(x+h)-f(x)}{h}$ & Tangent slope of 1D curve & Univariate line search; learning rate schedules. \\
\addlinespace
\textbf{Directional Deriv.} & $f: \mathbb{R}^d \to \mathbb{R}, \mathbf{u} \in \mathbb{R}^d$ & $D_{\mathbf{u}} f(\mathbf{x}) = \nabla f(\mathbf{x})^T \mathbf{u}$ & Slope along an arbitrary 3D/nD line of travel & Testing descent directions; line search along momentum. \\
\addlinespace
\textbf{Gradient Vector} & $f: \mathbb{R}^d \to \mathbb{R} \implies \mathbb{R}^d$ & $[\nabla f]_i = \frac{\partial f}{\partial x_i}$ & Orthogonal to contours; steepest ascent direction & First-order optimizers (SGD, Adam, RMSprop). \\
\addlinespace
\textbf{Jacobian Matrix} & $\mathbf{f}: \mathbb{R}^n \to \mathbb{R}^m \implies \mathbb{R}^{m \times n}$ & $J_{ij} = \frac{\partial f_i}{\partial x_j}$ & Best local linear mapping $d\mathbf{y} = J d\mathbf{x}$ & Layer-to-layer chain rule; Backpropagation in deep networks. \\
\addlinespace
\textbf{Hessian Matrix} & $f: \mathbb{R}^d \to \mathbb{R} \implies \mathbb{R}^{d \times d}$ & $H_{ij} = \frac{\partial^2 f}{\partial x_i \partial x_j}$ & Local curvature tensor; ellipsoid axes of bowl & Second-order optimizers (Newton-Raphson, BFGS, Laplace approximation). \\
\bottomrule
\end{tabularx}
\caption{Comparative Synthesis of Foundational Differential Operators in Machine Learning.}
\label{tab:calculus_operators}
\end{table}

% ----------------------------------------------------------------------
% PART 7: PYTHON / SCIENTIFIC COMPUTING LAB
% ----------------------------------------------------------------------
\section{Python / Scientific Computing Lab}

The following production script implements two core numerical verification procedures:
\begin{enumerate}
    \item \textbf{Numerical Gradient Checking:} Compares the analytical gradient $\frac{1}{N}X^T(X\mathbf{w} - \mathbf{y})$ against central finite differences to guarantee implementation accuracy.
    \item \textbf{Hessian Curvature Evaluation:} Analyzes the eigenvalues of the Hessian to classify critical points and detect saddle point geometries.
\end{enumerate}

\begin{lstlisting}[language=Python, caption={Numerical Gradient Checking and Hessian Curvature Analysis in Python.}]
import numpy as np

def loss_func(w: np.ndarray, X: np.ndarray, y: np.ndarray) -> float:
    # Computes Mean Squared Error loss: f(w) = (1 / (2*N)) * ||Xw - y||^2
    N = X.shape[0]
    residuals = X @ w - y
    return float((1.0 / (2.0 * N)) * np.sum(residuals ** 2))

def analytical_gradient(w: np.ndarray, X: np.ndarray, y: np.ndarray) -> np.ndarray:
    # Computes exact analytical gradient: grad = (1 / N) * X^T (Xw - y)
    N = X.shape[0]
    return (1.0 / N) * (X.T @ (X @ w - y))

def numerical_gradient_central(func, w: np.ndarray, eps: float = 1e-6, *args) -> np.ndarray:
    # Computes central difference numerical gradient: O(eps^2) truncation error
    grad = np.zeros_like(w)
    for i in range(len(w)):
        w_plus = w.copy()
        w_minus = w.copy()
        w_plus[i] += eps
        w_minus[i] -= eps
        grad[i] = (func(w_plus, *args) - func(w_minus, *args)) / (2.0 * eps)
    return grad

def classify_critical_point(H: np.ndarray) -> str:
    # Classifies a stationary point using the eigenvalues of its Hessian
    eigvals = np.linalg.eigvalsh(H)
    tol = 1e-7
    if np.all(eigvals > tol):
        return f"Local Minimum (Strictly Positive Definite, min eig: {eigvals.min():.4f})"
    elif np.all(eigvals < -tol):
        return f"Local Maximum (Strictly Negative Definite, max eig: {eigvals.max():.4f})"
    elif np.any(eigvals > tol) and np.any(eigvals < -tol):
        return f"Saddle Point (Indefinite, eig range: [{eigvals.min():.4f}, {eigvals.max():.4f}])"
    else:
        return f"Degenerate / Inconclusive (Zero eigenvalue detected)"

def run_gradient_check_lab():
    print("=" * 70)
    print("EXPERIMENT 1: ANALYTICAL VS NUMERICAL GRADIENT VERIFICATION")
    print("=" * 70)
    rng = np.random.RandomState(42)
    N, d = 100, 5
    X = rng.randn(N, d)
    w_true = rng.randn(d)
    y = X @ w_true + 0.1 * rng.randn(N)

    w_test = rng.randn(d)
    grad_analytic = analytical_gradient(w_test, X, y)
    grad_numeric = numerical_gradient_central(loss_func, w_test, 1e-6, X, y)

    rel_error = np.linalg.norm(grad_analytic - grad_numeric) / (np.linalg.norm(grad_analytic) + 1e-12)
    print(f"Analytical Gradient: {grad_analytic.round(4)}")
    print(f"Numerical Gradient:  {grad_numeric.round(4)}")
    print(f"Relative Error:      {rel_error:.2e} (Pass threshold: < 1e-5)")

    H = (1.0 / N) * (X.T @ X)
    print(f"\nHessian Eigenvalue Classification: {classify_critical_point(H)}")

if __name__ == "__main__":
    run_gradient_check_lab()
\end{lstlisting}

% ----------------------------------------------------------------------
% PART 8: COMPREHENSIVE PRACTICE PROBLEMS
% ----------------------------------------------------------------------
\section{Comprehensive Practice Problems}

\subsection*{Part A: Analytical \& Mathematical Problems}

\begin{enumerate}[label=\textbf{P\arabic*.}]
    \item \textbf{Gradient of Regularized Ridge Objective:}
    Consider the Ridge Regression loss functional $\mathcal{L}(\mathbf{w}) = \frac{1}{2}\|\mathbf{y} - X\mathbf{w}\|_2^2 + \frac{\lambda}{2}\|\mathbf{w}\|_2^2$, where $X \in \mathbb{R}^{N \times d}$, $\mathbf{y} \in \mathbb{R}^N$, and $\lambda > 0$.
    \begin{enumerate}[label=(\alph*)]
        \item Derive the exact gradient $\nabla_{\mathbf{w}} \mathcal{L}(\mathbf{w})$ using matrix calculus.
        \item Derive the exact Hessian matrix $H(\mathbf{w}) = \nabla^2_{\mathbf{w}} \mathcal{L}(\mathbf{w})$.
        \item Prove that if $\lambda > 0$, $H(\mathbf{w})$ is strictly positive definite regardless of whether $X$ has full column rank.
    \end{enumerate}

    \item \textbf{Jacobian of the Softmax Transformation:}
    Let $\mathbf{z} \in \mathbb{R}^K$. The softmax function $\boldsymbol{\sigma}: \mathbb{R}^K \to \mathbb{R}^K$ is defined by $\sigma_i(\mathbf{z}) = \frac{\exp(z_i)}{\sum_{j=1}^K \exp(z_j)}$ for $i = 1, \dots, K$.
    Prove that the entries of the Jacobian matrix $J_{\boldsymbol{\sigma}}(\mathbf{z}) \in \mathbb{R}^{K \times K}$ are given by:
    \begin{equation}
    \frac{\partial \sigma_i}{\partial z_j} = \sigma_i(\mathbf{z}) \left( \delta_{ij} - \sigma_j(\mathbf{z}) \right)
    \end{equation}
    where $\delta_{ij}$ is the Kronecker delta.

    \item \textbf{Directional Derivative Along Non-Gradient Paths:}
    Let $f(x_1, x_2) = 3x_1^2 + 2x_1 x_2 + x_2^2$.
    \begin{enumerate}[label=(\alph*)]
        \item Compute $\nabla f(1, 2)$.
        \item Find the rate of change of $f$ at $(1, 2)$ in the direction toward the point $(4, 6)$.
        \item Determine the unit vector $\mathbf{v}$ along which the directional derivative at $(1, 2)$ is exactly zero.
    \end{enumerate}

    \item \textbf{Second-Order Taylor Approximation Error Bound:}
    Let $f: \mathbb{R}^d \to \mathbb{R}$ have an $M$-Lipschitz continuous Hessian, meaning $\|\nabla^2 f(\mathbf{x}) - \nabla^2 f(\mathbf{y})\|_2 \le M\|\mathbf{x} - \mathbf{y}\|_2$ for all $\mathbf{x}, \mathbf{y}$. Prove that the quadratic approximation error satisfies:
    \begin{equation}
    \left| f(\mathbf{x} + \mathbf{v}) - \left[ f(\mathbf{x}) + \nabla f(\mathbf{x})^T \mathbf{v} + \frac{1}{2}\mathbf{v}^T \nabla^2 f(\mathbf{x})\mathbf{v} \right] \right| \le \frac{M}{6}\|\mathbf{v}\|_2^3
    \end{equation}

    \item \textbf{Saddle Point vs. Local Extrema Classification:}
    Consider the non-convex scalar function $f(x, y) = x^3 - 3x - y^2 + 4y$.
    \begin{enumerate}[label=(\alph*)]
        \item Find all critical points of $f(x, y)$ in $\mathbb{R}^2$ where $\nabla f(x, y) = \mathbf{0}$.
        \item Compute the Hessian matrix $H(x, y)$.
        \item Classify each critical point as a local minimum, local maximum, or saddle point via eigenvalue evaluation.
    \end{enumerate}
\end{enumerate}

\subsection*{Part B: Diagnostic Multiple-Choice Questions}

\begin{enumerate}[label=\textbf{MCQ\arabic*.}]
    \item Let $f(\mathbf{x}) = \mathbf{x}^T A \mathbf{x} + \mathbf{b}^T \mathbf{x} + c$, where $A \in \mathbb{R}^{d \times d}$ is not necessarily symmetric. What is the exact gradient $\nabla_{\mathbf{x}} f(\mathbf{x})$?
    \begin{enumerate}[label=(\alph*)]
        \item $2A\mathbf{x} + \mathbf{b}$
        \item $(A + A^T)\mathbf{x} + \mathbf{b}$
        \item $A^T \mathbf{x} + \mathbf{b}$
        \item $2A^T \mathbf{x} + \mathbf{b}$
    \end{enumerate}

    \item At a critical point $\mathbf{x}^*$ of a $C^2$ function $f: \mathbb{R}^d \to \mathbb{R}$, the Hessian matrix has eigenvalues $\lambda_1 = 4.2$, $\lambda_2 = 1.8$, and $\lambda_3 = -0.5$. Which statement is definitively true?
    \begin{enumerate}[label=(\alph*)]
        \item The point $\mathbf{x}^*$ is a local minimum because the trace $\sum \lambda_i = 5.5 > 0$.
        \item The point $\mathbf{x}^*$ is a local maximum because the determinant is negative.
        \item The point $\mathbf{x}^*$ is a saddle point because the Hessian is indefinite.
        \item Gradient descent with a small enough learning rate will always converge to $\mathbf{x}^*$.
    \end{enumerate}
\end{enumerate}

\chapter[Subspaces, Projections \& Least Squares]{Four Fundamental Subspaces, Orthogonal Projections \& Ordinary Least Squares}
\label{chap:week03}

% ----------------------------------------------------------------------
% PART 1: EXECUTIVE ROADMAP & LEARNING OBJECTIVES
% ----------------------------------------------------------------------
\section{Executive Roadmap \& Learning Objectives}

The foundational geometry of machine learning resides in the algebra of vector spaces and orthogonal projections. In classical high school algebra, students learn to solve systems of linear equations $A\mathbf{x} = \mathbf{b}$ where the number of equations equals the number of unknowns ($m = n$). However, modern machine learning operates almost exclusively in the regime where empirical observations vastly outnumber parameters: $m \gg n$. 

When ten thousand data points are collected to estimate five feature weights, the system $A\mathbf{x} = \mathbf{b}$ becomes heavily \textbf{overdetermined}. Because real-world measurements are corrupted by physical noise, unobserved latent variables, and stochastic disturbances, the target vector $\mathbf{b}$ almost never resides in the subspace spanned by the feature columns of $A$. High school algebra declares that \emph{no exact solution exists}. 

Ordinary Least Squares (OLS) does not give up in the face of this impossibility. Instead of seeking an unattainable exact solution, OLS reinterprets the problem through linear algebra: it seeks the vector $\hat{\mathbf{x}}$ whose image $A\hat{\mathbf{x}} = \hat{\mathbf{b}}$ is the \textbf{orthogonal projection} of $\mathbf{b}$ onto the column space $\Col(A)$. This chapter uncovers the profound geometric framework behind this projection: Gilbert Strang's Four Fundamental Subspaces, the Rank-Nullity Theorem, the algebraic properties of projectors, and the numerical stability of solving the normal equations.

\subsection{Prerequisites \& Scaffolding}
Mastery of this chapter requires fluency in vector space axioms, linear independence, and matrix rank as outlined in Table~\ref{tab:w03_prereqs}.

\begin{table}[htbp]
\centering
\small
\begin{tabularx}{\textwidth}{lXX}
\toprule
\textbf{Foundational Area} & \textbf{Core Concepts Required} & \textbf{Role in Machine Learning Foundations} \\
\midrule
Linear Algebra & Vector subspaces, span, linear independence, rank & Defining the structural domains where feature vectors and data reside \\
Inner Product Spaces & Dot products, orthogonality $\mathbf{u}^T \mathbf{v} = 0$, Pythagoras theorem & Decomposing arbitrary observation vectors into signal and noise \\
Matrix Operations & Transpose, inverse, reduced row echelon form (RREF) & Solving normal equations and identifying fundamental subspace bases \\
\bottomrule
\end{tabularx}
\caption{Prerequisite Scaffolding for Chapter 3.}
\label{tab:w03_prereqs}
\end{table}

\subsection{Core Competencies}
By the conclusion of this chapter, students will be able to:
\begin{itemize}
    \item Define, construct bases for, and compute the dimensions of the four fundamental subspaces ($\Col(A)$, $\Null(A)$, $\Row(A)$, and $\Null(A^T)$) associated with any matrix $A \in \mathbb{R}^{m \times n}$.
    \item State and prove the Fundamental Theorem of Linear Algebra and orthogonal complement relationships ($\Col(A)^\perp = \Null(A^T)$ and $\Row(A)^\perp = \Null(A)$).
    \item Derive the algebraic projection matrix $P = A(A^T A)^{-1}A^T$ from geometric first principles and verify its algebraic properties ($P^2 = P, P^T = P$).
    \item Formulate and solve the Normal Equations $A^T A \hat{\mathbf{x}} = A^T \mathbf{b}$ through both geometric orthogonality and matrix calculus.
    \item Diagnose multicollinearity, rank deficiency, and ill-conditioned Gram matrices $A^T A$, and explain why QR decomposition is numerically preferred in production software.
\end{itemize}

\subsection{Notation \& Mathematical Conventions}
\begin{itemize}
    \item $A \in \mathbb{R}^{m \times n}$: A design matrix with $m$ observations (rows) and $n$ features (columns), having rank $r = \rank(A) \le \min(m, n)$.
    \item $\mathbf{a}_1, \dots, \mathbf{a}_n \in \mathbb{R}^m$: Column vectors of $A$, representing feature vectors across all samples.
    \item $\Col(A) \triangleq \{\mathbf{y} \in \mathbb{R}^m \mid \exists \mathbf{x} \in \mathbb{R}^n, \mathbf{y} = A\mathbf{x}\}$: Column space (range) of $A$, a subspace of $\mathbb{R}^m$.
    \item $\Null(A) \triangleq \{\mathbf{x} \in \mathbb{R}^n \mid A\mathbf{x} = \mathbf{0}\}$: Nullspace (kernel) of $A$, a subspace of $\mathbb{R}^n$.
    \item $\Row(A) \triangleq \Col(A^T) \subseteq \mathbb{R}^n$: Row space of $A$, spanned by the rows of $A$.
    \item $\Null(A^T) \subseteq \mathbb{R}^m$: Left nullspace of $A$, orthogonal to the column space.
    \item $V^\perp$: The orthogonal complement of subspace $V$.
    \item $P_V$: The unique orthogonal projection matrix onto subspace $V$.
    \item $\hat{\mathbf{b}} \triangleq P_{\Col(A)}\mathbf{b}$: The orthogonal projection of $\mathbf{b}$ onto $\Col(A)$ (the fitted signal).
    \item $\mathbf{e} \triangleq \mathbf{b} - \hat{\mathbf{b}} = (I - P)\mathbf{b}$: The residual error vector (the orthogonal noise).
\end{itemize}

% ----------------------------------------------------------------------
% PART 2: THE GRAND PARADOX & INTUITIVE PUZZLE
% ----------------------------------------------------------------------
\section{The Grand Paradox \& Intuitive Motivation}

\subsection{The Inconsistent Measurement Paradox}
Consider a high-school student attempting to determine a person's height $h$ and weight $w$. If the student takes two independent measurements, they get a system of two equations in two unknowns:
\begin{equation}
\begin{aligned}
2h + 3w &= 160 \\
h + 2w &= 95
\end{aligned}
\end{equation}
Assuming the equations are independent, Gaussian elimination delivers a unique, exact solution. 

Now suppose the student attempts to be more scientifically rigorous by taking a \emph{third} measurement with a slightly different instrument:
\begin{equation}
3h + 4w = 230
\end{equation}
Substituting the solution from the first two equations ($h = 35, w = 30$) into the third equation yields $3(35) + 4(30) = 105 + 120 = 225 \ne 230$. The system is now contradictory! 

\begin{quote}
\textit{\textbf{The Paradox:}} Why should collecting \textbf{more} data destroy our ability to find a mathematical solution? If more data yields more information, why does classical algebra declare that 3 equations in 2 unknowns has \textbf{zero} solutions?
\end{quote}

The resolution requires moving from the \textbf{Row Picture} of linear equations to the \textbf{Column Picture}:
\begin{itemize}
    \item \textbf{The Row Picture:} Each equation represents a line in $\mathbb{R}^2$. In an inconsistent system, the three lines do not intersect at a single point; instead, they form a triangle. There is no single $(h, w)$ that lies on all three lines simultaneously.
    \item \textbf{The Column Picture:} We rewrite the system as a single vector equation:
    \begin{equation}
    h \begin{bmatrix} 2 \\ 1 \\ 3 \end{bmatrix} + w \begin{bmatrix} 3 \\ 2 \\ 4 \end{bmatrix} = \begin{bmatrix} 160 \\ 95 \\ 230 \end{bmatrix} \iff A\mathbf{x} = \mathbf{b}
    \end{equation}
    The two column vectors $\mathbf{a}_1 = [2, 1, 3]^T$ and $\mathbf{a}_2 = [3, 2, 4]^T$ span a 2-dimensional plane embedded in 3-dimensional observation space $\mathbb{R}^3$. The measurement vector $\mathbf{b} = [160, 95, 230]^T$ does not lie on this 2D plane! Because $\mathbf{b} \notin \Col(A)$, no linear combination of $\mathbf{a}_1$ and $\mathbf{a}_2$ can ever reach $\mathbf{b}$.
\end{itemize}

Dropping the third equation would discard real information and subject the model to the arbitrary noise of the remaining two measurements. The genius of Ordinary Least Squares is to accept that while $\mathbf{b}$ cannot be reached, there is a \emph{unique closest point} on the plane $\Col(A)$. That closest point is the \textbf{orthogonal shadow} of $\mathbf{b}$ cast onto the plane.

\begin{ponderbox}[The Overdetermination Dilemma]
Suppose you are training a neural network or linear model with 1,000,000 training examples and 50 weights. The system $X\mathbf{w} = \mathbf{y}$ has 1,000,000 equations and 50 unknowns. 

If someone suggests: \textit{"To make the system solvable, let's randomly delete 999,950 rows so that we have 50 equations and 50 unknowns,"} why is this mathematically catastrophic from the perspective of variance and information geometry? What does orthogonal projection do with those 999,950 equations instead?
\end{ponderbox}

% ----------------------------------------------------------------------
% PART 3: THE FOUR FUNDAMENTAL SUBSPACES
% ----------------------------------------------------------------------
\section{The Four Fundamental Subspaces}

Every linear transformation represented by a real matrix $A \in \mathbb{R}^{m \times n}$ partitions the input space $\mathbb{R}^n$ and output space $\mathbb{R}^m$ into two pairs of mutually orthogonal vector subspaces. Understanding these four subspaces is the key to mastering both regression and dimensionality reduction.

\begin{definitionbox}{The Four Fundamental Subspaces}
For any matrix $A \in \mathbb{R}^{m \times n}$ of rank $r = \rank(A)$:
\begin{enumerate}
    \item \textbf{Column Space} $\Col(A) \subseteq \mathbb{R}^m$: The span of the columns of $A$. It represents all possible model predictions $A\mathbf{x}$. $\dim(\Col(A)) = r$.
    \item \textbf{Nullspace} $\Null(A) \subseteq \mathbb{R}^n$: The set of all input vectors mapped to zero: $\{\mathbf{x} \in \mathbb{R}^n \mid A\mathbf{x} = \mathbf{0}\}$. $\dim(\Null(A)) = n - r$.
    \item \textbf{Row Space} $\Row(A) \triangleq \Col(A^T) \subseteq \mathbb{R}^n$: The span of the rows of $A$ (columns of $A^T$). $\dim(\Row(A)) = r$.
    \item \textbf{Left Nullspace} $\Null(A^T) \subseteq \mathbb{R}^m$: The set of all vectors $\mathbf{y} \in \mathbb{R}^m$ such that $A^T \mathbf{y} = \mathbf{0}$. $\dim(\Null(A^T)) = m - r$.
\end{enumerate}
\end{definitionbox}

\begin{figure}[htbp]
\centering
\begin{tikzpicture}[scale=0.9, font=\sffamily\footnotesize]
% Input space R^n
\draw[thick, fill=primarylight!30] (0, 0) rectangle (4.5, 5);
\node[above, font=\bfseries] at (2.25, 5) {Input Space $\mathbb{R}^n$};

% Row space
\draw[thick, fill=primarydark!20] (0.3, 2.6) rectangle (4.2, 4.7);
\node[font=\bfseries, color=primarydark] at (2.25, 3.8) {Row Space $\Row(A)$};
\node at (2.25, 3.2) {$\dim = r$};
\node at (2.25, 2.8) {$\Col(A^T)$};

% Nullspace
\draw[thick, fill=tealdef!20] (0.3, 0.3) rectangle (4.2, 2.4);
\node[font=\bfseries, color=tealdef] at (2.25, 1.5) {Nullspace $\Null(A)$};
\node at (2.25, 0.9) {$\dim = n - r$};

% Orthogonal symbol in R^n
\draw[line width=1.5pt, draw=black] (2.25, 2.4) -- (2.25, 2.6);
\node at (2.25, 2.5) {\footnotesize $\perp$};

% Output space R^m
\draw[thick, fill=primarylight!30] (8.5, 0) rectangle (13, 5);
\node[above, font=\bfseries] at (10.75, 5) {Output Space $\mathbb{R}^m$};

% Column space
\draw[thick, fill=primarydark!20] (8.8, 2.6) rectangle (12.7, 4.7);
\node[font=\bfseries, color=primarydark] at (10.75, 3.8) {Column Space $\Col(A)$};
\node at (10.75, 3.2) {$\dim = r$};

% Left Nullspace
\draw[thick, fill=tealdef!20] (8.8, 0.3) rectangle (12.7, 2.4);
\node[font=\bfseries, color=tealdef] at (10.75, 1.5) {Left Nullspace $\Null(A^T)$};
\node at (10.75, 0.9) {$\dim = m - r$};

% Orthogonal symbol in R^m
\draw[line width=1.5pt, draw=black] (10.75, 2.4) -- (10.75, 2.6);
\node at (10.75, 2.5) {\footnotesize $\perp$};

% Transformation arrows
\draw[->, >=stealth, line width=1.8pt, color=primarydark] (4.2, 3.65) -- (8.8, 3.65) node[midway, above] {$A\mathbf{x}_r \in \Col(A)$ (invertible)};
\draw[->, >=stealth, line width=1.8pt, dashed, color=tealdef] (4.2, 1.35) -- (8.8, 1.35) node[midway, below] {$A\mathbf{x}_n = \mathbf{0}$};

\end{tikzpicture}
\caption{The Fundamental Theorem of Linear Algebra (Gilbert Strang's Big Picture). The row space is mapped invertibly onto the column space, while the nullspace is crushed to $\mathbf{0}$.}
\label{fig:strang_subspaces}
\end{figure}

\begin{theorem}[Fundamental Theorem of Linear Algebra]
Let $A \in \mathbb{R}^{m \times n}$ have rank $r$. Then:
\begin{enumerate}
    \item The row space and nullspace are orthogonal complements in $\mathbb{R}^n$:
    \begin{equation}
    \Row(A)^\perp = \Null(A) \quad \text{and} \quad \Row(A) \oplus \Null(A) = \mathbb{R}^n
    \label{eq:ftla_rn}
    \end{equation}
    \item The column space and left nullspace are orthogonal complements in $\mathbb{R}^m$:
    \begin{equation}
    \Col(A)^\perp = \Null(A^T) \quad \text{and} \quad \Col(A) \oplus \Null(A^T) = \mathbb{R}^m
    \label{eq:ftla_rm}
    \end{equation}
\end{enumerate}
\end{theorem}

\begin{proof}
We first prove that $\Col(A)^\perp = \Null(A^T)$.
\begin{enumerate}
    \item \textbf{Direction 1: $\Null(A^T) \subseteq \Col(A)^\perp$.}
    Let $\mathbf{y} \in \Null(A^T)$, which means $A^T \mathbf{y} = \mathbf{0}$.
    Let $\mathbf{v} \in \Col(A)$ be an arbitrary vector in the column space. By definition of the column space, there exists some $\mathbf{x} \in \mathbb{R}^n$ such that $\mathbf{v} = A\mathbf{x}$.
    We compute the Euclidean inner product between $\mathbf{v}$ and $\mathbf{y}$:
    \begin{equation}
    \langle \mathbf{v}, \mathbf{y} \rangle = \mathbf{v}^T \mathbf{y} = (A\mathbf{x})^T \mathbf{y} = \mathbf{x}^T (A^T \mathbf{y}) = \mathbf{x}^T \mathbf{0} = 0
    \end{equation}
    Since $\mathbf{v}^T \mathbf{y} = 0$ for \emph{every} $\mathbf{v} \in \Col(A)$, the vector $\mathbf{y}$ is orthogonal to the entire column space: $\mathbf{y} \in \Col(A)^\perp$. Thus, $\Null(A^T) \subseteq \Col(A)^\perp$.

    \item \textbf{Direction 2: Dimension match.}
    By the Rank-Nullity Theorem applied to $A^T \in \mathbb{R}^{n \times m}$, the rank of $A^T$ is $r$, and its domain is $\mathbb{R}^m$. Therefore:
    \begin{equation}
    \dim(\Null(A^T)) = m - \rank(A^T) = m - r
    \end{equation}
    Because $\Col(A)$ is an $r$-dimensional subspace of the $m$-dimensional space $\mathbb{R}^m$, its orthogonal complement has dimension:
    \begin{equation}
    \dim(\Col(A)^\perp) = m - \dim(\Col(A)) = m - r
    \end{equation}
    Since $\Null(A^T) \subseteq \Col(A)^\perp$ and their dimensions are identical, we conclude that $\Col(A)^\perp = \Null(A^T)$.
\end{enumerate}
The proof that $\Row(A)^\perp = \Null(A)$ is obtained by substituting $A^T$ in place of $A$.
\end{proof}

\begin{examplebox}{Finding the Four Fundamental Subspaces of a Concrete Matrix}
Let matrix $A \in \mathbb{R}^{3 \times 4}$ be defined as:
\begin{equation}
A = \begin{bmatrix} 1 & 2 & 0 & 1 \\ 0 & 1 & 1 & 0 \\ 1 & 3 & 1 & 1 \end{bmatrix}
\end{equation}
Find the rank, bases for all four fundamental subspaces, and verify their orthogonal complement relationships.

\textbf{Step 1: Row Reduction to RREF.}
Row 3 minus Row 1 ($R_3 \leftarrow R_3 - R_1$):
\begin{equation}
\begin{bmatrix} 1 & 2 & 0 & 1 \\ 0 & 1 & 1 & 0 \\ 0 & 1 & 1 & 0 \end{bmatrix}
\end{equation}
Row 3 minus Row 2 ($R_3 \leftarrow R_3 - R_2$):
\begin{equation}
\begin{bmatrix} 1 & 2 & 0 & 1 \\ 0 & 1 & 1 & 0 \\ 0 & 0 & 0 & 0 \end{bmatrix}
\end{equation}
Row 1 minus 2 times Row 2 ($R_1 \leftarrow R_1 - 2R_2$):
\begin{equation}
R = \text{rref}(A) = \begin{bmatrix} 1 & 0 & -2 & 1 \\ 0 & 1 & 1 & 0 \\ 0 & 0 & 0 & 0 \end{bmatrix}
\end{equation}

\textbf{Step 2: Identify Pivot and Free Variables.}
The pivot columns are columns 1 and 2. Columns 3 and 4 are free variables ($x_3, x_4$).
Therefore, the matrix rank is $r = \rank(A) = 2$.

\textbf{Step 3: Basis for Column Space $\Col(A) \subseteq \mathbb{R}^3$.}
The basis is formed by the pivot columns of the \textbf{original} matrix $A$ (not $R$!):
\begin{equation}
\mathcal{B}_{\Col(A)} = \left\{ \begin{bmatrix} 1 \\ 0 \\ 1 \end{bmatrix}, \begin{bmatrix} 2 \\ 1 \\ 3 \end{bmatrix} \right\}, \quad \dim(\Col(A)) = 2
\end{equation}

\textbf{Step 4: Basis for Row Space $\Row(A) \subseteq \mathbb{R}^4$.}
The basis is formed by the non-zero rows of the reduced matrix $R$:
\begin{equation}
\mathcal{B}_{\Row(A)} = \left\{ \begin{bmatrix} 1 \\ 0 \\ -2 \\ 1 \end{bmatrix}, \begin{bmatrix} 0 \\ 1 \\ 1 \\ 0 \end{bmatrix} \right\}, \quad \dim(\Row(A)) = 2
\end{equation}

\textbf{Step 5: Basis for Nullspace $\Null(A) \subseteq \mathbb{R}^4$.}
From $R\mathbf{x} = \mathbf{0}$, we have:
\begin{align}
x_1 - 2x_3 + x_4 = 0 &\implies x_1 = 2x_3 - x_4 \\
x_2 + x_3 = 0 &\implies x_2 = -x_3
\end{align}
Setting free variables $(x_3, x_4) = (1, 0)$ and $(0, 1)$ yields the special solutions:
\begin{equation}
\mathbf{s}_1 = \begin{bmatrix} 2 \\ -1 \\ 1 \\ 0 \end{bmatrix}, \quad \mathbf{s}_2 = \begin{bmatrix} -1 \\ 0 \\ 0 \\ 1 \end{bmatrix} \implies \mathcal{B}_{\Null(A)} = \left\{ \begin{bmatrix} 2 \\ -1 \\ 1 \\ 0 \end{bmatrix}, \begin{bmatrix} -1 \\ 0 \\ 0 \\ 1 \end{bmatrix} \right\}
\end{equation}
Notice $\dim(\Null(A)) = n - r = 4 - 2 = 2$.

\textbf{Step 6: Basis for Left Nullspace $\Null(A^T) \subseteq \mathbb{R}^3$.}
Vectors $\mathbf{y} = [y_1, y_2, y_3]^T$ satisfying $A^T \mathbf{y} = \mathbf{0}$. We observe Row 3 of $A$ was Row 1 + Row 2. Thus:
\begin{equation}
1 \cdot \text{Row}_1 + 1 \cdot \text{Row}_2 - 1 \cdot \text{Row}_3 = \mathbf{0}^T \implies \mathbf{y} = \begin{bmatrix} 1 \\ 1 \\ -1 \end{bmatrix}
\end{equation}
Indeed, $\dim(\Null(A^T)) = m - r = 3 - 2 = 1$, and:
\begin{equation}
A^T \begin{bmatrix} 1 \\ 1 \\ -1 \end{bmatrix} = \begin{bmatrix} 1 & 0 & 1 \\ 2 & 1 & 3 \\ 0 & 1 & 1 \\ 1 & 0 & 1 \end{bmatrix} \begin{bmatrix} 1 \\ 1 \\ -1 \end{bmatrix} = \begin{bmatrix} 1 - 1 \\ 2 + 1 - 3 \\ 1 - 1 \\ 1 - 1 \end{bmatrix} = \begin{bmatrix} 0 \\ 0 \\ 0 \\ 0 \end{bmatrix}
\end{equation}

\textbf{Step 7: Orthogonality Verification.}
Testing $\Row(A) \perp \Null(A)$:
\begin{equation}
[1, 0, -2, 1] \begin{bmatrix} 2 \\ -1 \\ 1 \\ 0 \end{bmatrix} = 1(2) + 0(-1) - 2(1) + 1(0) = 2 - 2 = 0 \quad \checkmark
\end{equation}
Testing $\Col(A) \perp \Null(A^T)$:
\begin{equation}
\mathbf{y}^T \mathbf{a}_1 = [1, 1, -1] \begin{bmatrix} 1 \\ 0 \\ 1 \end{bmatrix} = 1 + 0 - 1 = 0 \quad \checkmark, \quad \mathbf{y}^T \mathbf{a}_2 = [1, 1, -1] \begin{bmatrix} 2 \\ 1 \\ 3 \end{bmatrix} = 2 + 1 - 3 = 0 \quad \checkmark
\end{equation}
\end{examplebox}

\begin{ponderbox}[Why $X^T X$ Preserves the Kernel]
Let $X \in \mathbb{R}^{m \times n}$. Why is $\Null(X^T X)$ always \emph{exactly} identical to $\Null(X)$?

It is obvious that if $X\mathbf{v} = \mathbf{0}$, then $X^T (X\mathbf{v}) = X^T \mathbf{0} = \mathbf{0}$. But why can't there exist some magical non-zero vector $\mathbf{u} \in \mathbb{R}^n$ such that $X\mathbf{u} \ne \mathbf{0}$, but when multiplied by $X^T$, $X^T (X\mathbf{u})$ becomes $\mathbf{0}$? 

\textit{Hint: If $X^T(X\mathbf{u}) = \mathbf{0}$, where must the vector $X\mathbf{u}$ lie? Can a vector simultaneously live in the column space of $X$ and the left nullspace of $X$ without being zero?}
\end{ponderbox}

% ----------------------------------------------------------------------
% PART 4: ORTHOGONAL PROJECTIONS & DERIVATION
% ----------------------------------------------------------------------
\section{Orthogonal Projections onto Subspaces}

\subsection{Projection onto a 1-Dimensional Line}
Before deriving the general multi-dimensional projection matrix, we begin with the simplest geometric case: projecting an observation vector $\mathbf{b} \in \mathbb{R}^m$ onto the line spanned by a single non-zero vector $\mathbf{a} \in \mathbb{R}^m$.

The projected point $\hat{\mathbf{b}}$ must lie on the line, meaning it is a scalar multiple of $\mathbf{a}$:
\begin{equation}
\hat{\mathbf{b}} = \hat{x}\mathbf{a} \quad \text{for some scalar } \hat{x} \in \mathbb{R}
\end{equation}
The error vector connecting $\mathbf{b}$ to its projection is:
\begin{equation}
\mathbf{e} = \mathbf{b} - \hat{\mathbf{b}} = \mathbf{b} - \hat{x}\mathbf{a}
\end{equation}
To minimize the Euclidean distance $\|\mathbf{e}\|_2 = \|\mathbf{b} - \hat{x}\mathbf{a}\|_2$, the error vector must be perpendicular to the line vector $\mathbf{a}$:
\begin{equation}
\mathbf{a}^T \mathbf{e} = 0 \implies \mathbf{a}^T (\mathbf{b} - \hat{x}\mathbf{a}) = 0 \implies \mathbf{a}^T \mathbf{b} - \hat{x}(\mathbf{a}^T \mathbf{a}) = 0
\end{equation}
Solving for the optimal scalar coefficient $\hat{x}$:
\begin{equation}
\hat{x} = \frac{\mathbf{a}^T \mathbf{b}}{\mathbf{a}^T \mathbf{a}}
\label{eq:scalar_proj_coeff}
\end{equation}
Substituting $\hat{x}$ back into $\hat{\mathbf{b}} = \hat{x}\mathbf{a}$:
\begin{equation}
\hat{\mathbf{b}} = \mathbf{a} \hat{x} = \mathbf{a} \left( \frac{\mathbf{a}^T \mathbf{b}}{\mathbf{a}^T \mathbf{a}} \right) = \left( \frac{\mathbf{a}\mathbf{a}^T}{\mathbf{a}^T \mathbf{a}} \right) \mathbf{b} = P_{\mathbf{a}} \mathbf{b}
\end{equation}
The matrix $P_{\mathbf{a}} = \frac{\mathbf{a}\mathbf{a}^T}{\mathbf{a}^T \mathbf{a}} \in \mathbb{R}^{m \times m}$ is the \textbf{rank-1 orthogonal projection matrix} onto the line spanned by $\mathbf{a}$. Notice the outer product $\mathbf{a}\mathbf{a}^T$ produces an $m \times m$ matrix, while the inner product $\mathbf{a}^T \mathbf{a} = \|\mathbf{a}\|_2^2$ is a scalar normalizing factor.

\begin{examplebox}{1D Orthogonal Projection and Vector Decomposition}
Let $\mathbf{a} = \begin{bmatrix} 2 \\ 1 \\ 2 \end{bmatrix}$ and target vector $\mathbf{b} = \begin{bmatrix} 1 \\ 2 \\ 3 \end{bmatrix}$. 
\begin{enumerate}
    \item Compute the projection matrix $P_{\mathbf{a}}$.
    \item Compute the projection $\hat{\mathbf{b}}$ and error vector $\mathbf{e}$.
    \item Verify orthogonality $\mathbf{a}^T \mathbf{e} = 0$ and the Pythagorean theorem $\|\mathbf{b}\|_2^2 = \|\hat{\mathbf{b}}\|_2^2 + \|\mathbf{e}\|_2^2$.
\end{enumerate}

\textbf{Step 1: Compute Outer and Inner Products.}
\begin{equation}
\mathbf{a}^T \mathbf{a} = 2^2 + 1^2 + 2^2 = 4 + 1 + 4 = 9
\end{equation}
The outer product matrix is:
\begin{equation}
\mathbf{a}\mathbf{a}^T = \begin{bmatrix} 2 \\ 1 \\ 2 \end{bmatrix} \begin{bmatrix} 2 & 1 & 2 \end{bmatrix} = \begin{bmatrix} 4 & 2 & 4 \\ 2 & 1 & 2 \\ 4 & 2 & 4 \end{bmatrix}
\end{equation}
Thus, the projection matrix is:
\begin{equation}
P_{\mathbf{a}} = \frac{1}{9} \begin{bmatrix} 4 & 2 & 4 \\ 2 & 1 & 2 \\ 4 & 2 & 4 \end{bmatrix}
\end{equation}

\textbf{Step 2: Compute Projection and Error.}
Scalar coefficient:
\begin{equation}
\hat{x} = \frac{\mathbf{a}^T \mathbf{b}}{\mathbf{a}^T \mathbf{a}} = \frac{2(1) + 1(2) + 2(3)}{9} = \frac{2 + 2 + 6}{9} = \frac{10}{9}
\end{equation}
The projection vector is:
\begin{equation}
\hat{\mathbf{b}} = \hat{x}\mathbf{a} = \frac{10}{9} \begin{bmatrix} 2 \\ 1 \\ 2 \end{bmatrix} = \begin{bmatrix} 20/9 \\ 10/9 \\ 20/9 \end{bmatrix}
\end{equation}
The error vector is:
\begin{equation}
\mathbf{e} = \mathbf{b} - \hat{\mathbf{b}} = \begin{bmatrix} 1 - 20/9 \\ 2 - 10/9 \\ 3 - 20/9 \end{bmatrix} = \begin{bmatrix} -11/9 \\ 8/9 \\ 7/9 \end{bmatrix} = \frac{1}{9} \begin{bmatrix} -11 \\ 8 \\ 7 \end{bmatrix}
\end{equation}

\textbf{Step 3: Verification of Orthogonality and Pythagoras.}
\begin{equation}
\mathbf{a}^T \mathbf{e} = 2(-11/9) + 1(8/9) + 2(7/9) = \frac{-22 + 8 + 14}{9} = \frac{0}{9} = 0 \quad \checkmark
\end{equation}
Computing squared Euclidean norms:
\begin{align}
\|\mathbf{b}\|_2^2 &= 1^2 + 2^2 + 3^2 = 1 + 4 + 9 = 14 \\
\|\hat{\mathbf{b}}\|_2^2 &= \left(\frac{10}{9}\right)^2 (2^2 + 1^2 + 2^2) = \frac{100}{81} \cdot 9 = \frac{100}{9} \approx 11.111 \\
\|\mathbf{e}\|_2^2 &= \frac{(-11)^2 + 8^2 + 7^2}{81} = \frac{121 + 64 + 49}{81} = \frac{234}{81} = \frac{26}{9} \approx 2.889
\end{align}
Summing the components:
\begin{equation}
\|\hat{\mathbf{b}}\|_2^2 + \|\mathbf{e}\|_2^2 = \frac{100}{9} + \frac{26}{9} = \frac{126}{9} = 14 = \|\mathbf{b}\|_2^2 \quad \checkmark
\end{equation}
\end{examplebox}

\subsection{Projection onto General Subspaces (\texorpdfstring{$n$}{n}-Dimensions)}

\begin{theorem}[Invertibility of the Gram Matrix]
\label{thm:gram_invertible}
Let $A \in \mathbb{R}^{m \times n}$. The square symmetric matrix $A^T A \in \mathbb{R}^{n \times n}$ is invertible if and only if $A$ has full column rank ($\rank(A) = n$).
\end{theorem}

\begin{proof}
It suffices to prove that $A^T A$ and $A$ share the exact same nullspace: $\Null(A^T A) = \Null(A)$.
\begin{enumerate}
    \item Let $\mathbf{x} \in \Null(A)$. Then $A\mathbf{x} = \mathbf{0} \implies A^T (A\mathbf{x}) = A^T \mathbf{0} = \mathbf{0}$, hence $\mathbf{x} \in \Null(A^T A)$. This establishes $\Null(A) \subseteq \Null(A^T A)$.
    \item Conversely, let $\mathbf{x} \in \Null(A^T A)$. Then $(A^T A)\mathbf{x} = \mathbf{0}$. Pre-multiplying by $\mathbf{x}^T$:
    \begin{equation}
    \mathbf{x}^T (A^T A \mathbf{x}) = \mathbf{x}^T \mathbf{0} = 0 \implies (A\mathbf{x})^T (A\mathbf{x}) = 0 \implies \|A\mathbf{x}\|_2^2 = 0 \implies A\mathbf{x} = \mathbf{0}
    \end{equation}
    Thus $\mathbf{x} \in \Null(A)$, which proves $\Null(A^T A) \subseteq \Null(A)$.
\end{enumerate}
Since $\Null(A^T A) = \Null(A)$, the square matrix $A^T A \in \mathbb{R}^{n \times n}$ is invertible if and only if its nullspace is trivial ($\Null(A^T A) = \{\mathbf{0}\}$), which occurs if and only if $\Null(A) = \{\mathbf{0}\}$, meaning $\rank(A) = n$.
\end{proof}

When projecting onto an $n$-dimensional subspace spanned by the linearly independent columns of $A = [\mathbf{a}_1, \dots, \mathbf{a}_n] \in \mathbb{R}^{m \times n}$ ($m \ge n$), the target vector $\hat{\mathbf{b}}$ is a linear combination of all $n$ columns:
\begin{equation}
\hat{\mathbf{b}} = \hat{x}_1 \mathbf{a}_1 + \hat{x}_2 \mathbf{a}_2 + \dots + \hat{x}_n \mathbf{a}_n = A\hat{\mathbf{x}} \quad \text{for some } \hat{\mathbf{x}} \in \mathbb{R}^n
\end{equation}
The error vector is:
\begin{equation}
\mathbf{e} = \mathbf{b} - \hat{\mathbf{b}} = \mathbf{b} - A\hat{\mathbf{x}}
\end{equation}
For $\hat{\mathbf{b}}$ to be the closest point in $\Col(A)$ to $\mathbf{b}$, the error vector $\mathbf{e}$ must be orthogonal to \emph{every single column} of $A$:
\begin{equation}
\mathbf{a}_1^T \mathbf{e} = 0, \quad \mathbf{a}_2^T \mathbf{e} = 0, \quad \dots, \quad \mathbf{a}_n^T \mathbf{e} = 0
\end{equation}
Stacking these $n$ separate scalar equations into a single compact matrix product:
\begin{equation}
\begin{bmatrix} \mathbf{a}_1^T \\ \mathbf{a}_2^T \\ \vdots \\ \mathbf{a}_n^T \end{bmatrix} \mathbf{e} = \mathbf{0} \iff A^T \mathbf{e} = \mathbf{0}
\end{equation}
This reveals the beautiful connection to the Fundamental Theorem of Linear Algebra:
\begin{equation}
\mathbf{e} \in \Null(A^T) = \Col(A)^\perp
\end{equation}
Substituting $\mathbf{e} = \mathbf{b} - A\hat{\mathbf{x}}$ into $A^T \mathbf{e} = \mathbf{0}$:
\begin{equation}
A^T (\mathbf{b} - A\hat{\mathbf{x}}) = \mathbf{0} \implies A^T \mathbf{b} - A^T A \hat{\mathbf{x}} = \mathbf{0} \implies A^T A \hat{\mathbf{x}} = A^T \mathbf{b}
\label{eq:normal_equations_geometry}
\end{equation}
Equation~\eqref{eq:normal_equations_geometry} is the celebrated \textbf{Normal Equations}. Because $A$ has full column rank ($\rank(A) = n$), the $n \times n$ symmetric Gram matrix $A^T A$ is invertible (Theorem~\ref{thm:gram_invertible}). 

Pre-multiplying by $(A^T A)^{-1}$:
\begin{equation}
\hat{\mathbf{x}} = (A^T A)^{-1} A^T \mathbf{b}
\end{equation}
Substituting the optimal parameter vector $\hat{\mathbf{x}}$ back into $\hat{\mathbf{b}} = A\hat{\mathbf{x}}$:
\begin{equation}
\hat{\mathbf{b}} = A\hat{\mathbf{x}} = A (A^T A)^{-1} A^T \mathbf{b} = P \mathbf{b}
\end{equation}
where the matrix $P \in \mathbb{R}^{m \times m}$ is defined as:
\begin{equation}
P \triangleq A(A^T A)^{-1}A^T
\label{eq:proj_matrix_formula}
\end{equation}

\begin{theorem}[Fundamental Properties of the Orthogonal Projector]
The matrix $P = A(A^T A)^{-1}A^T$ satisfies:
\begin{enumerate}
    \item \textbf{Symmetry:} $P^T = P$.
    \item \textbf{Idempotence:} $P^2 = P$.
    \item \textbf{Complementary Projector:} $I - P$ is the orthogonal projector onto $\Null(A^T) = \Col(A)^\perp$.
    \item \textbf{Trace-Rank Identity:} $\trace(P) = \rank(A) = n$.
\end{enumerate}
\end{theorem}

\begin{proof}
\begin{enumerate}
    \item \textbf{Symmetry:} Using $(BC)^T = C^T B^T$ and $((A^T A)^{-1})^T = ((A^T A)^T)^{-1} = (A^T A)^{-1}$:
    \begin{equation}
    P^T = \left( A(A^T A)^{-1}A^T \right)^T = (A^T)^T \left( (A^T A)^{-1} \right)^T A^T = A (A^T A)^{-1} A^T = P
    \end{equation}
    \item \textbf{Idempotence:} Multiplying $P$ by itself:
    \begin{align}
    P^2 &= \left( A(A^T A)^{-1}A^T \right) \left( A(A^T A)^{-1}A^T \right) \\
    &= A (A^T A)^{-1} \underbrace{\left( A^T A \right) (A^T A)^{-1}}_{= I_n} A^T \\
    &= A (A^T A)^{-1} A^T = P
    \end{align}
    \item \textbf{Complementary Projector:} Let $Q = I - P$.
    \begin{align}
    Q^T &= (I - P)^T = I - P^T = I - P = Q \quad (\text{symmetric}) \\
    Q^2 &= (I - P)(I - P) = I - 2P + P^2 = I - 2P + P = I - P = Q \quad (\text{idempotent})
    \end{align}
    Furthermore, for any vector $\mathbf{v} \in \Col(A)$, $P\mathbf{v} = \mathbf{v} \implies Q\mathbf{v} = (I - P)\mathbf{v} = \mathbf{0}$. For any vector $\mathbf{u} \in \Col(A)^\perp = \Null(A^T)$, $A^T \mathbf{u} = \mathbf{0} \implies P\mathbf{u} = A(A^T A)^{-1}A^T \mathbf{u} = \mathbf{0} \implies Q\mathbf{u} = \mathbf{u}$. Thus, $Q$ projects orthogonally onto $\Col(A)^\perp$.
    \item \textbf{Trace:} Using the cyclic trace identity $\trace(BC) = \trace(CB)$ with $B = A$ and $C = (A^T A)^{-1}A^T$:
    \begin{equation}
    \trace(P) = \trace\left( A \cdot (A^T A)^{-1}A^T \right) = \trace\left( (A^T A)^{-1}A^T \cdot A \right) = \trace(I_n) = n = \rank(A)
    \end{equation}
\end{enumerate}
\end{proof}

\begin{ponderbox}[The Idempotence Test]
Imagine you are standing on a sunny field. The sun casts your shadow onto the 2D ground.
\begin{enumerate}
    \item If you take the shadow on the ground and project it onto the ground a second time, does it change? Why does this explain $P^2 = P$?
    \item If $P^2 = P$ is true for \emph{any} shadow (even a slanted, oblique shadow cast by a streetlight at an angle), what extra condition is required to ensure that the shadow is the \emph{closest} point to the original object? Why must $P = P^T$?
\end{enumerate}
\end{ponderbox}

% ----------------------------------------------------------------------
% PART 5: WORKED NUMERICAL REGRESSION EXAMPLE
% ----------------------------------------------------------------------
\section{In-Depth Worked Numerical Example: OLS Line Fitting}

\begin{examplebox}{Full Ordinary Least Squares Line Fitting via Projection Geometry}
Fit an affine line $y = \beta_0 + \beta_1 x$ to the three empirical observations:
\begin{equation}
(x_1, y_1) = (1, 2), \quad (x_2, y_2) = (2, 2), \quad (x_3, y_3) = (3, 4)
\end{equation}
Derive the design matrix, the Gram matrix, the projection matrix, the fitted values, the residual vector, and verify the zero-mean residual property.

\textbf{Step 1: Formulate the Overdetermined System $X\boldsymbol{\beta} = \mathbf{y}$.}
With intercept feature $x_0 = 1$ and slope feature $x_1$:
\begin{equation}
X = \begin{bmatrix} 1 & 1 \\ 1 & 2 \\ 1 & 3 \end{bmatrix} \in \mathbb{R}^{3 \times 2}, \quad \boldsymbol{\beta} = \begin{bmatrix} \beta_0 \\ \beta_1 \end{bmatrix} \in \mathbb{R}^2, \quad \mathbf{y} = \begin{bmatrix} 2 \\ 2 \\ 4 \end{bmatrix} \in \mathbb{R}^3
\end{equation}
Notice $m = 3$ equations and $n = 2$ unknowns. The columns $\mathbf{x}_0 = [1, 1, 1]^T$ and $\mathbf{x}_1 = [1, 2, 3]^T$ are clearly linearly independent, so $\rank(X) = 2$.

\textbf{Step 2: Compute the Gram Matrix $X^T X$ and $X^T \mathbf{y}$.}
\begin{equation}
X^T X = \begin{bmatrix} 1 & 1 & 1 \\ 1 & 2 & 3 \end{bmatrix} \begin{bmatrix} 1 & 1 \\ 1 & 2 \\ 1 & 3 \end{bmatrix} = \begin{bmatrix} 1+1+1 & 1+2+3 \\ 1+2+3 & 1+4+9 \end{bmatrix} = \begin{bmatrix} 3 & 6 \\ 6 & 14 \end{bmatrix}
\end{equation}
The determinant is $\det(X^T X) = 3(14) - 6(6) = 42 - 36 = 6 \ne 0$.
The inverse Gram matrix is:
\begin{equation}
(X^T X)^{-1} = \frac{1}{6} \begin{bmatrix} 14 & -6 \\ -6 & 3 \end{bmatrix} = \begin{bmatrix} 7/3 & -1 \\ -1 & 1/2 \end{bmatrix}
\end{equation}
The right-hand side vector is:
\begin{equation}
X^T \mathbf{y} = \begin{bmatrix} 1 & 1 & 1 \\ 1 & 2 & 3 \end{bmatrix} \begin{bmatrix} 2 \\ 2 \\ 4 \end{bmatrix} = \begin{bmatrix} 2 + 2 + 4 \\ 2(1) + 2(2) + 4(3) \end{bmatrix} = \begin{bmatrix} 8 \\ 18 \end{bmatrix}
\end{equation}

\textbf{Step 3: Solve the Normal Equations for $\hat{\boldsymbol{\beta}}$.}
\begin{equation}
\hat{\boldsymbol{\beta}} = (X^T X)^{-1} X^T \mathbf{y} = \frac{1}{6} \begin{bmatrix} 14 & -6 \\ -6 & 3 \end{bmatrix} \begin{bmatrix} 8 \\ 18 \end{bmatrix} = \frac{1}{6} \begin{bmatrix} 14(8) - 6(18) \\ -6(8) + 3(18) \end{bmatrix} = \frac{1}{6} \begin{bmatrix} 112 - 108 \\ -48 + 54 \end{bmatrix} = \frac{1}{6} \begin{bmatrix} 4 \\ 6 \end{bmatrix} = \begin{bmatrix} 2/3 \\ 1 \end{bmatrix}
\end{equation}
Thus, the fitted Ordinary Least Squares line is:
\begin{equation}
\hat{y} = \frac{2}{3} + 1 \cdot x
\end{equation}

\textbf{Step 4: Compute the Orthogonal Projection Matrix $P$.}
\begin{align}
P &= X (X^T X)^{-1} X^T = \begin{bmatrix} 1 & 1 \\ 1 & 2 \\ 1 & 3 \end{bmatrix} \frac{1}{6} \begin{bmatrix} 14 & -6 \\ -6 & 3 \end{bmatrix} \begin{bmatrix} 1 & 1 & 1 \\ 1 & 2 & 3 \end{bmatrix} \\
&= \frac{1}{6} \begin{bmatrix} 1 & 1 \\ 1 & 2 \\ 1 & 3 \end{bmatrix} \begin{bmatrix} 8 & 2 & -4 \\ -3 & 0 & 3 \end{bmatrix} = \frac{1}{6} \begin{bmatrix} 5 & 2 & -1 \\ 2 & 2 & 2 \\ -1 & 2 & 5 \end{bmatrix}
\end{align}
Notice $P$ is symmetric ($P^T = P$) and $\trace(P) = \frac{5 + 2 + 5}{6} = \frac{12}{6} = 2 = \rank(X)$.

\textbf{Step 5: Compute Fitted Values $\hat{\mathbf{y}}$ and Residuals $\mathbf{e}$.}
\begin{equation}
\hat{\mathbf{y}} = X\hat{\boldsymbol{\beta}} = \begin{bmatrix} 1 & 1 \\ 1 & 2 \\ 1 & 3 \end{bmatrix} \begin{bmatrix} 2/3 \\ 1 \end{bmatrix} = \begin{bmatrix} 2/3 + 1 \\ 2/3 + 2 \\ 2/3 + 3 \end{bmatrix} = \begin{bmatrix} 5/3 \\ 8/3 \\ 11/3 \end{bmatrix}
\end{equation}
The residual error vector is:
\begin{equation}
\mathbf{e} = \mathbf{y} - \hat{\mathbf{y}} = \begin{bmatrix} 2 - 5/3 \\ 2 - 8/3 \\ 4 - 11/3 \end{bmatrix} = \begin{bmatrix} 1/3 \\ -2/3 \\ 1/3 \end{bmatrix} = \frac{1}{3} \begin{bmatrix} 1 \\ -2 \\ 1 \end{bmatrix}
\end{equation}

\textbf{Step 6: Confirm Error Orthogonality and Zero-Mean Residuals.}
We test $X^T \mathbf{e}$:
\begin{equation}
X^T \mathbf{e} = \begin{bmatrix} 1 & 1 & 1 \\ 1 & 2 & 3 \end{bmatrix} \left( \frac{1}{3} \begin{bmatrix} 1 \\ -2 \\ 1 \end{bmatrix} \right) = \frac{1}{3} \begin{bmatrix} 1 - 2 + 1 \\ 1(1) + 2(-2) + 3(1) \end{bmatrix} = \frac{1}{3} \begin{bmatrix} 0 \\ 1 - 4 + 3 \end{bmatrix} = \begin{bmatrix} 0 \\ 0 \end{bmatrix} \quad \checkmark
\end{equation}
Because the first column of $X$ is the constant intercept column of ones $\mathbf{1}$, the first equation of $X^T \mathbf{e} = \mathbf{0}$ states:
\begin{equation}
\mathbf{1}^T \mathbf{e} = \sum_{i=1}^m e_i = \frac{1}{3} - \frac{2}{3} + \frac{1}{3} = 0
\end{equation}
This proves mathematically why \textbf{any linear regression model that contains an intercept term always has residuals that sum to exactly zero}!
\end{examplebox}

% ----------------------------------------------------------------------
% PART 6: EXAM TRAPS & DIAGNOSTIC MISCONCEPTIONS
% ----------------------------------------------------------------------
\section{Exam Traps \& Diagnostic Misconceptions}

\begin{trapbox}{Gotcha 1: The Rectangular Matrix Pseudo-Cancellation Fallacy}
A student attempts to simplify the projection matrix $P = A(A^T A)^{-1}A^T$ by expanding the matrix inverse $(BC)^{-1} = C^{-1} B^{-1}$:
\begin{equation}
P = A (A^T A)^{-1} A^T \stackrel{?}{=} A A^{-1} (A^T)^{-1} A^T = I_m \cdot I_m = I_m
\end{equation}
\textbf{The Trap:} The matrix $A \in \mathbb{R}^{m \times n}$ with $m > n$ is \textbf{rectangular}. Rectangular matrices do not possess a two-sided matrix inverse $A^{-1}$! The algebraic identity $(BC)^{-1} = C^{-1}B^{-1}$ holds \emph{strictly} when both $B$ and $C$ are square invertible matrices. 

If $P = I_m$, every vector in $\mathbb{R}^m$ would be projected onto itself, meaning $\Col(A) = \mathbb{R}^m$. But $\Col(A)$ has dimension at most $n < m$, which is a proper subspace!
\end{trapbox}

\begin{trapbox}{Gotcha 2: Confusing General (Oblique) Projectors with Orthogonal Projectors}
A student designs a projection operator that satisfies $P^2 = P$, but discovers that the residual $\|\mathbf{b} - P\mathbf{b}\|_2$ is strictly larger than the distance to the subspace:
\begin{equation}
\|\mathbf{b} - P\mathbf{b}\|_2 > \min_{\mathbf{v} \in V} \|\mathbf{b} - \mathbf{v}\|_2
\end{equation}
\textbf{The Trap:} Idempotence ($P^2 = P$) is the defining property of \textbf{any} projection (including skewed, oblique projections). To guarantee that $P\mathbf{b}$ is the \textbf{closest point in Euclidean distance}, the projection must be \textbf{orthogonal}, which requires the matrix to be \textbf{symmetric}: $P = P^T$. 

Without symmetry, the error vector $\mathbf{b} - P\mathbf{b}$ drops along an oblique angle rather than perpendicular to the subspace, violating least-squares optimality.
\end{trapbox}

\begin{trapbox}{Gotcha 3: Exact Collinearity Causing Inversion Failure}
A data scientist adds a redundant feature (e.g., $x_3 = 2x_1 - 5x_2$) into a design matrix $X$, and attempts to solve for regression coefficients using \texttt{np.linalg.inv(X.T @ X)}. The script crashes with a \texttt{LinAlgError: Singular matrix}.
\textbf{The Trap:} When columns are linearly dependent, $\rank(X) < n$, which means $\Null(X) \ne \{\mathbf{0}\}$. By Theorem~\ref{thm:gram_invertible}, $\Null(X^T X) = \Null(X) \ne \{\mathbf{0}\}$, so $\det(X^T X) = 0$. The Gram matrix is strictly singular and has no inverse. 

The normal equations still have solutions, but they are \textbf{infinitely many} (an $(n-r)$-dimensional affine space). To resolve this in practice, one must either eliminate redundant features, apply Ridge regularization ($X^T X + \lambda I$), or use the Moore-Penrose pseudoinverse $X^\dagger$.
\end{trapbox}

\begin{trapbox}{Gotcha 4: The Gram Condition Number Squaring Curse}
A numerical programmer decides to solve $X^T X \hat{\boldsymbol{\beta}} = X^T \mathbf{y}$ directly using Cholesky or LU decomposition when $X$ has condition number $\kappa(X) = 10^8$. The resulting parameter estimates are completely contaminated by garbage values.
\textbf{The Trap:} The condition number of a matrix measures how much relative input perturbations are amplified in the output. For the Gram matrix:
\begin{equation}
\kappa(X^T X) = (\kappa(X))^2
\end{equation}
If $\kappa(X) = 10^8$, then $\kappa(X^T X) = 10^{16}$. Standard IEEE-754 double-precision floating-point arithmetic provides approximately $15$ to $17$ decimal digits of precision ($53$ mantissa bits). A condition number of $10^{16}$ wipes out \emph{all} digits of precision! 

For this reason, production scientific computing libraries (\texttt{scipy.linalg.lstsq}, PyTorch, LAPACK) \textbf{never} form the product $X^T X$ directly. Instead, they factor $X$ using QR decomposition ($X = QR$) or SVD, solving $R\hat{\boldsymbol{\beta}} = Q^T \mathbf{y}$ which operates directly on $\kappa(X)$ without squaring it!
\end{trapbox}

\begin{ponderbox}[The Condition Number Squaring Curse]
Suppose you have an ill-conditioned design matrix $X \in \mathbb{R}^{1000 \times 10}$ with condition number $\kappa(X) = 10^5$. 
\begin{enumerate}
    \item What is the condition number of the normal equations system matrix $X^T X$?
    \item How many decimal digits of precision do you lose when solving $(X^T X)\mathbf{w} = X^T \mathbf{y}$ compared to solving via QR decomposition $R\mathbf{w} = Q^T \mathbf{y}$?
\end{enumerate}
\end{ponderbox}

% ----------------------------------------------------------------------
% PART 7: COMPARATIVE SYNTHESIS TABLE
% ----------------------------------------------------------------------
\section{Comparative Synthesis Table}

Table~\ref{tab:subspaces_synthesis} provides an exhaustive structural comparison of the Four Fundamental Subspaces and their operational role in machine learning models.

\begin{table}[htbp]
\centering
\small
\begin{tabularx}{\textwidth}{lXXXXX}
\toprule
\textbf{Subspace} & \textbf{Notation} & \textbf{Ambient Space} & \textbf{Dimension} & \textbf{Orthogonal Complement} & \textbf{Operational Role in ML} \\
\midrule
Column Space & $\Col(A)$ & $\mathbb{R}^m$ & $r = \rank(A)$ & Left Nullspace $\Null(A^T)$ & The space of all realizable model predictions $A\mathbf{x}$ \\
\addlinespace
Nullspace & $\Null(A)$ & $\mathbb{R}^n$ & $n - r$ & Row Space $\Row(A)$ & Directions of parameter unidentifiability (zero-gradient perturbations) \\
\addlinespace
Row Space & $\Row(A)$ & $\mathbb{R}^n$ & $r = \rank(A)$ & Nullspace $\Null(A)$ & The effective parameter space; minimum-norm solutions reside here \\
\addlinespace
Left Nullspace & $\Null(A^T)$ & $\mathbb{R}^m$ & $m - r$ & Column Space $\Col(A)$ & The space of irrecoverable empirical residuals and model errors $\mathbf{e}$ \\
\bottomrule
\end{tabularx}
\caption{Comparative Structural Synthesis of the Four Fundamental Subspaces.}
\label{tab:subspaces_synthesis}
\end{table}

% ----------------------------------------------------------------------
% PART 8: PYTHON SCIENTIFIC COMPUTING LAB
% ----------------------------------------------------------------------
\section{Python / Scientific Computing Lab}

The production script below demonstrates subspace projector construction, verifies the spectral trace-rank identity, checks residual orthogonality, and numerically validates why QR decomposition beats normal equations on collinear data.

\begin{lstlisting}[language=Python, caption={Subspace Projections and Numerical Stability Comparison in Python.}]
import numpy as np

def build_projection_matrix(A: np.ndarray) -> np.ndarray:
    """Computes orthogonal projection matrix P = A (A^T A)^{-1} A^T."""
    AtA = A.T @ A
    AtA_inv = np.linalg.inv(AtA)
    return A @ AtA_inv @ A.T

def solve_ols_normal_equations(A: np.ndarray, b: np.ndarray) -> np.ndarray:
    """Solves normal equations directly: (A^T A) x = A^T b."""
    return np.linalg.solve(A.T @ A, A.T @ b)

def solve_ols_qr(A: np.ndarray, b: np.ndarray) -> np.ndarray:
    """Solves OLS stably using QR decomposition: A = QR => R x = Q^T b."""
    Q, R = np.linalg.qr(A)
    return np.linalg.solve(R, Q.T @ b)

def main():
    rng = np.random.RandomState(42)
    m, n = 8, 3
    A = rng.randn(m, n)
    b = rng.randn(m)

    # 1. Compute projection matrix
    P = build_projection_matrix(A)

    # 2. Verify algebraic properties
    is_idempotent = np.allclose(P @ P, P)
    is_symmetric = np.allclose(P, P.T)
    trace_val = np.trace(P)

    print(f"P is Idempotent (P^2 == P): {is_idempotent}")
    print(f"P is Symmetric (P == P^T):  {is_symmetric}")
    print(f"Trace of P: {trace_val:.4f} (Expected rank = {n})")

    # 3. Verify error vector orthogonality: A^T e = 0
    b_hat = P @ b
    e = b - b_hat
    orthogonality_check = np.linalg.norm(A.T @ e)
    print(f"Orthogonality ||A^T e||: {orthogonality_check:.2e} (Expected ~ 0)")

    # 4. Compare Normal Equations vs. QR on an ill-conditioned design matrix
    A_ill = A.copy()
    # Create near-collinearity: column 2 is almost a linear combination of cols 0 and 1
    A_ill[:, 2] = A_ill[:, 0] + 1e-7 * A_ill[:, 1]
    
    cond_A = np.linalg.cond(A_ill)
    cond_AtA = np.linalg.cond(A_ill.T @ A_ill)
    print(f"\nCondition number kappa(A):     {cond_A:.2e}")
    print(f"Condition number kappa(A^T A): {cond_AtA:.2e} (Squared!)")

    x_qr = solve_ols_qr(A_ill, b)
    x_ne = solve_ols_normal_equations(A_ill, b)
    print(f"QR Solution:               {x_qr}")
    print(f"Normal Equations Solution: {x_ne}")
    print(f"Relative discrepancy:      {np.linalg.norm(x_qr - x_ne) / np.linalg.norm(x_qr):.2e}")

if __name__ == "__main__":
    main()
\end{lstlisting}

% ----------------------------------------------------------------------
% PART 9: COMPREHENSIVE PRACTICE PROBLEMS
% ----------------------------------------------------------------------
\section{Comprehensive Practice Problems}

\subsection*{Part A: Analytical \& Mathematical Problems}

\begin{enumerate}[label=\textbf{P\arabic*.}]
    \item \textbf{Explicit Projection Matrix Computation:}
    Let matrix $A = \begin{bmatrix} 1 & 0 \\ 1 & 1 \\ 0 & 1 \end{bmatrix}$ and vector $\mathbf{b} = \begin{bmatrix} 3 \\ 1 \\ 2 \end{bmatrix}$.
    \begin{enumerate}[label=(\alph*)]
        \item Compute the Gram matrix $A^T A$ and its inverse $(A^T A)^{-1}$.
        \item Compute the $3 \times 3$ projection matrix $P$ onto $\Col(A)$.
        \item Find the projected point $\hat{\mathbf{b}} = P\mathbf{b}$ and the error vector $\mathbf{e} = \mathbf{b} - \hat{\mathbf{b}}$.
        \item Confirm by direct calculation that $A^T \mathbf{e} = \mathbf{0}$.
    \end{enumerate}

    \item \textbf{Properties of Complementary Projectors:}
    Let $P$ be an orthogonal projector onto subspace $V \subseteq \mathbb{R}^m$. Define $Q \triangleq I - P$.
    \begin{enumerate}[label=(\alph*)]
        \item Prove that $Q$ is idempotent ($Q^2 = Q$) and symmetric ($Q^T = Q$).
        \item Show that $\Col(Q) = \Null(P) = V^\perp$.
        \item Prove the generalized Pythagorean identity: for any $\mathbf{x} \in \mathbb{R}^m$, $\|\mathbf{x}\|_2^2 = \|P\mathbf{x}\|_2^2 + \|Q\mathbf{x}\|_2^2$.
    \end{enumerate}

    \item \textbf{Rank-1 Projection Matrix:}
    Let $\mathbf{a} \in \mathbb{R}^m$ be a non-zero column vector.
    \begin{enumerate}[label=(\alph*)]
        \item Derive the projection matrix $P_{\mathbf{a}}$ onto the 1-dimensional subspace spanned by $\mathbf{a}$.
        \item Show that $\trace(P_{\mathbf{a}}) = 1$.
        \item Show that for any vector $\mathbf{b}$, the scalar multiple $c = \frac{\mathbf{a}^T \mathbf{b}}{\mathbf{a}^T \mathbf{a}}$ minimizes $\|c\mathbf{a} - \mathbf{b}\|_2$.
    \end{enumerate}

    \item \textbf{Affine Fit with Offset:}
    We wish to fit a line $y = mx + c$ through three points $(x_1, y_1) = (0, 1)$, $(x_2, y_2) = (1, 3)$, and $(x_3, y_3) = (2, 4)$.
    \begin{enumerate}[label=(\alph*)]
        \item Set up the overdetermined linear system $A\mathbf{w} = \mathbf{y}$, identifying $A$, $\mathbf{w} = [c, m]^T$, and $\mathbf{y}$.
        \item Formulate and solve the Normal Equations $A^T A \mathbf{w} = A^T \mathbf{y}$.
        \item Compute the residual sum of squares $\|\mathbf{y} - A\hat{\mathbf{w}}\|_2^2$.
    \end{enumerate}

    \item \textbf{Subspace Intersection and Projection Commutativity:}
    Let $P_1$ and $P_2$ be orthogonal projection matrices onto subspaces $V_1$ and $V_2$, respectively.
    Prove that the composite operator $P_1 P_2$ is an orthogonal projection matrix if and only if $P_1$ and $P_2$ commute ($P_1 P_2 = P_2 P_1$), in which case it projects onto $V_1 \cap V_2$.
\end{enumerate}

\subsection*{Part B: Diagnostic Multiple-Choice Questions}

\begin{enumerate}[label=\textbf{MCQ\arabic*.}]
    \item Let $A \in \mathbb{R}^{m \times n}$ with $\rank(A) = r < n \le m$. Which statement regarding the linear system $A^T A \mathbf{x} = A^T \mathbf{b}$ is definitively true?
    \begin{enumerate}[label=(\alph*)]
        \item The system is inconsistent and possesses no solution.
        \item The system is always consistent, but has infinitely many solutions.
        \item The matrix $A^T A$ is strictly positive definite.
        \item The projection matrix $P = A(A^T A)^{-1}A^T$ is well-defined.
    \end{enumerate}

    \item If $P \in \mathbb{R}^{5 \times 5}$ is an orthogonal projection matrix onto a 3-dimensional subspace $V \subset \mathbb{R}^5$, what are the eigenvalues of $P$?
    \begin{enumerate}[label=(\alph*)]
        \item $1, 1, 1, 1, 1$
        \item $3, 3, 3, 0, 0$
        \item $1, 1, 1, 0, 0$
        \item $1, 1, 1, -1, -1$
    \end{enumerate}
\end{enumerate}

\chapter[Eigenvalues \& Diagonalization]{Polynomial Regression, Eigenvalues, Eigenspaces \& Diagonalization}
\label{chap:week04}

% ----------------------------------------------------------------------
% PART 1: EXECUTIVE ROADMAP & LEARNING OBJECTIVES
% ----------------------------------------------------------------------
\section{Executive Roadmap \& Learning Objectives}

At first glance, linear models appear fundamentally limited: they can only construct flat hyperplanes and linear decision boundaries. Yet, in modern data science, linear algorithms are routinely deployed to capture highly non-linear physical phenomena, aerodynamic drag curves, and financial volatility smiles. This remarkable versatility is unlocked through \textbf{non-linear feature expansions}: by projecting low-dimensional inputs into higher-dimensional feature spaces via polynomial basis functions, a linear model learns complex curves while retaining the convex optimization guarantees of Ordinary Least Squares.

However, once we operate in expanded feature spaces or analyze iterated transformations (such as recurrence relations in deep networks or transitions in Markov models), we confront matrices acting repeatedly on vector states: $\mathbf{x}_{k} = A^k \mathbf{x}_0$. Computing $A^k$ directly via brute-force matrix multiplication costs $\mathcal{O}(k n^3)$ operations and tangles coordinate axes uncontrollably. 

The mathematical remedy is the \textbf{spectral theory of linear operators}. By identifying the invariant directions—the \textbf{eigenvectors}—and their corresponding scale factors—the \textbf{eigenvalues}—we decouple a coupled matrix transformation into $n$ independent scalar scalings. This chapter builds the theoretical and practical bridge from polynomial feature representations to the spectral decomposition of linear operators.

\subsection{Prerequisites \& Scaffolding}
Mastery of this chapter builds upon vector spaces, determinant calculations, and ordinary least squares, as outlined in Table~\ref{tab:w04_prereqs}.

\begin{table}[htbp]
\centering
\small
\begin{tabularx}{\textwidth}{lXX}
\toprule
\textbf{Foundational Area} & \textbf{Core Concepts Required} & \textbf{Role in Machine Learning Foundations} \\
\midrule
Linear Algebra & Determinants $\det(A)$, matrix invertibility, linear independence & Solving characteristic polynomial $\det(A - \lambda I) = 0$ \\
Subspace Theory & Nullspace $\Null(A - \lambda I)$, rank-nullity theorem & Computing geometric eigenspaces and diagnosing deficiency \\
Least Squares & Design matrix $X$, normal equations $(X^T X)\mathbf{w} = X^T \mathbf{y}$ & Fitting polynomial curves via Vandermonde expansions \\
\bottomrule
\end{tabularx}
\caption{Prerequisite Scaffolding for Chapter 4.}
\label{tab:w04_prereqs}
\end{table}

\subsection{Core Competencies}
By the conclusion of this chapter, students will be able to:
\begin{itemize}
    \item Formulate polynomial regression problems via the non-linear feature map $\boldsymbol{\phi}(x)$ and construct Vandermonde design matrices.
    \item Compute eigenvalues, eigenvectors, and geometric eigenspaces $E_\lambda = \Null(A - \lambda I)$ for arbitrary square matrices.
    \item Distinguish rigorously between algebraic multiplicity $am(\lambda)$ and geometric multiplicity $gm(\lambda)$ to diagnose defective matrices.
    \item Diagonalize square matrices into $A = PDP^{-1}$ and compute iterated matrix powers $A^k = P D^k P^{-1}$ in closed form.
    \item Formulate discrete linear dynamical systems (including recurrence relations and Markov chains) and determine their asymptotic stability.
\end{itemize}

\subsection{Notation \& Mathematical Conventions}
\begin{itemize}
    \item $A \in \mathbb{R}^{n \times n}$: A real square matrix representing a linear transformation from $\mathbb{R}^n$ to $\mathbb{R}^n$.
    \item $\lambda \in \mathbb{C}$: An eigenvalue of $A$, satisfying $A\mathbf{v} = \lambda \mathbf{v}$.
    \item $\mathbf{v} \in \mathbb{C}^n \setminus \{\mathbf{0}\}$: A non-zero eigenvector associated with eigenvalue $\lambda$.
    \item $E_\lambda \triangleq \Null(A - \lambda I)$: The eigenspace corresponding to eigenvalue $\lambda$.
    \item $p_A(\lambda) \triangleq \det(A - \lambda I)$: The degree-$n$ characteristic polynomial of $A$.
    \item $am(\lambda)$: Algebraic multiplicity (the multiplicity of $\lambda$ as a root of $p_A(\lambda)$).
    \item $gm(\lambda) \triangleq \dim(E_\lambda) = n - \rank(A - \lambda I)$: Geometric multiplicity (the number of linearly independent eigenvectors for $\lambda$).
    \item $P = [\mathbf{v}_1, \dots, \mathbf{v}_n] \in \mathbb{C}^{n \times n}$: The modal matrix whose columns are linearly independent eigenvectors.
    \item $D = \diag(\lambda_1, \dots, \lambda_n)$: The diagonal spectral matrix of eigenvalues.
    \item $\Phi \in \mathbb{R}^{N \times (p+1)}$: The Vandermonde design matrix for degree-$p$ polynomial regression.
\end{itemize}

% ----------------------------------------------------------------------
% PART 2: THE GRAND PARADOX & MOTIVATION
% ----------------------------------------------------------------------
\section{The Grand Paradox: The Linear Model That Curves}

\subsection{The Linearity Illusion}
When newcomers encounter the phrase \textbf{Linear Regression}, they almost universally imagine fitting a straight line $y = mx + c$ through data points on a 2D plane. When presented with a parabolic dataset—such as projectile motion under gravity $y(t) = -\frac{1}{2}gt^2 + v_0 t + y_0$—the novice assumes that linear regression has failed, and that one must abandon linear models in favor of "non-linear algorithms."

\begin{quote}
\textit{\textbf{The Paradox:}} How can a strictly \textbf{linear} model fit a parabolas, cubics, or oscillatory curves? If an algorithm can fit an arbitrary polynomial curve, why is it still classified as "linear"?
\end{quote}

The resolution hinges on the distinction between \textbf{linearity in variables} versus \textbf{linearity in parameters}:
\begin{equation}
f(x; \mathbf{w}) = w_0 + w_1 x + w_2 x^2 + \dots + w_p x^p = \sum_{j=0}^p w_j \phi_j(x) = \mathbf{w}^T \boldsymbol{\phi}(x)
\end{equation}
While $f(x)$ is non-linear with respect to the input variable $x$, it is \textbf{strictly linear with respect to the parameter vector $\mathbf{w}$}:
\begin{equation}
f(x; \alpha \mathbf{u} + \beta \mathbf{v}) = \alpha f(x; \mathbf{u}) + \beta f(x; \mathbf{v})
\end{equation}
By defining a feature expansion $\boldsymbol{\phi}: \mathbb{R} \to \mathbb{R}^{p+1}$ where $\boldsymbol{\phi}(x) = [1, x, x^2, \dots, x^p]^T$, we transform the 1D non-linear curve into a \textbf{flat $(p+1)$-dimensional hyperplane} in feature space! Because the loss function $\|\mathbf{y} - \Phi \mathbf{w}\|_2^2$ remains a convex quadratic with respect to $\mathbf{w}$, all the guarantees of Ordinary Least Squares derived in Chapter~\ref{chap:week03} apply without alteration.

\begin{ponderbox}[The Non-Linear Illusion]
Consider the model $y = w_0 + w_1 \sin(x) + w_2 \exp(x) + w_3 \log(x)$. 
\begin{enumerate}
    \item Is this model linear regression or non-linear regression?
    \item What is the feature mapping $\boldsymbol{\phi}(x)$ that allows this model to be solved using the standard Normal Equations $\Phi^T \Phi \mathbf{w} = \Phi^T \mathbf{y}$?
    \item In contrast, consider the model $y = \frac{w_0}{1 + \exp(-w_1 x)}$. Can this model be solved via standard linear least squares? Why or why not?
\end{enumerate}
\end{ponderbox}

\subsection{The Coordinate Untangling Paradox}
Consider a general $n \times n$ matrix $A$. If you multiply an arbitrary vector $\mathbf{x} \in \mathbb{R}^n$ by $A$, the resulting vector $A\mathbf{x}$ typically rotates, skews, and stretches simultaneously. If we iterate this operation $k$ times:
\begin{equation}
\mathbf{x}_k = A \mathbf{x}_{k-1} = A^2 \mathbf{x}_{k-2} = \dots = A^k \mathbf{x}_0
\end{equation}
Every coordinate of $\mathbf{x}_k$ becomes an entangled combination of all previous coordinates.

Yet, in almost every linear system, there exist special directions where matrix multiplication does \textbf{not} rotate the vector at all! Along these privileged directions, matrix multiplication reduces to \textbf{pure scalar multiplication}:
\begin{equation}
A\mathbf{v} = \lambda \mathbf{v}
\end{equation}
These vectors $\mathbf{v}$ are the \textbf{eigenvectors}, and the scalars $\lambda$ are the \textbf{eigenvalues}. If we change coordinates to a basis of eigenvectors, the matrix $A$ becomes a simple diagonal matrix $D = \diag(\lambda_1, \dots, \lambda_n)$, and computing $A^k$ collapses to raising individual scalar numbers to the power $k$: $\lambda_i^k$.

% ----------------------------------------------------------------------
% PART 3: POLYNOMIAL REGRESSION & VANDERMONDE MATRICES
% ----------------------------------------------------------------------
\section{Polynomial Regression \& The Vandermonde Matrix}

Let $\{ (x_i, y_i) \}_{i=1}^N$ be a dataset of $N$ scalar observations. We wish to fit a degree-$p$ polynomial:
\begin{equation}
\hat{y}(x) = w_0 + w_1 x + w_2 x^2 + \dots + w_p x^p
\end{equation}
We define the polynomial feature mapping $\boldsymbol{\phi}: \mathbb{R} \to \mathbb{R}^{p+1}$:
\begin{equation}
\boldsymbol{\phi}(x) \triangleq \begin{bmatrix} 1 & x & x^2 & \dots & x^p \end{bmatrix}^T
\end{equation}
Stacking the feature vectors for all $N$ data points into rows forms the \textbf{Vandermonde Design Matrix}:
\begin{equation}
\Phi = \begin{bmatrix}
\boldsymbol{\phi}(x_1)^T \\
\boldsymbol{\phi}(x_2)^T \\
\vdots \\
\boldsymbol{\phi}(x_N)^T
\end{bmatrix} = \begin{bmatrix}
1 & x_1 & x_1^2 & \dots & x_1^p \\
1 & x_2 & x_2^2 & \dots & x_2^p \\
\vdots & \vdots & \vdots & \ddots & \vdots \\
1 & x_N & x_N^2 & \dots & x_N^p
\end{bmatrix} \in \mathbb{R}^{N \times (p+1)}
\label{eq:vandermonde_matrix}
\end{equation}
The polynomial regression problem is now formulated as the standard linear system:
\begin{equation}
\Phi \mathbf{w} \approx \mathbf{y} \quad \text{where } \mathbf{w} = \begin{bmatrix} w_0 \\ w_1 \\ \vdots \\ w_p \end{bmatrix} \in \mathbb{R}^{p+1}, \quad \mathbf{y} = \begin{bmatrix} y_1 \\ y_2 \\ \vdots \\ y_N \end{bmatrix} \in \mathbb{R}^N
\end{equation}

\begin{theorem}[Full Rank of Vandermonde Matrices]
Let $x_1, x_2, \dots, x_N$ be $N$ distinct real numbers. If $N \ge p + 1$, the Vandermonde matrix $\Phi \in \mathbb{R}^{N \times (p+1)}$ has full column rank: $\rank(\Phi) = p + 1$.
\end{theorem}

\begin{proof}
Suppose $\Phi \mathbf{c} = \mathbf{0}$ for some coefficient vector $\mathbf{c} = [c_0, c_1, \dots, c_p]^T \in \mathbb{R}^{p+1}$.
Each row $i$ of the product $\Phi \mathbf{c} = \mathbf{0}$ represents the evaluation of the polynomial $P(x) = \sum_{j=0}^p c_j x^j$ at point $x_i$:
\begin{equation}
P(x_i) = c_0 + c_1 x_i + c_2 x_i^2 + \dots + c_p x_i^p = 0 \quad \text{for all } i \in \{1, 2, \dots, N\}
\end{equation}
Thus, $P(x)$ is a polynomial of degree at most $p$ that has $N$ distinct roots $x_1, \dots, x_N$.
By the Fundamental Theorem of Algebra, a non-zero polynomial of degree $p$ can have at most $p$ distinct roots. Since $N \ge p + 1 > p$, the polynomial must be identically zero: $c_0 = c_1 = \dots = c_p = 0$.
Hence, the only vector in $\Null(\Phi)$ is $\mathbf{c} = \mathbf{0}$, proving $\rank(\Phi) = p + 1$.
\end{proof}

\begin{examplebox}{Polynomial Feature Expansion \& Exact Parabolic Interpolation}
Given three data points:
\begin{equation}
(x_1, y_1) = (-1, 1), \quad (x_2, y_2) = (0, 0), \quad (x_3, y_3) = (1, 2)
\end{equation}
Fit a quadratic polynomial $y = w_0 + w_1 x + w_2 x^2$ ($p = 2, N = 3$).

\textbf{Step 1: Construct the Square Vandermonde Matrix $\Phi$.}
\begin{equation}
\Phi = \begin{bmatrix} 1 & x_1 & x_1^2 \\ 1 & x_2 & x_2^2 \\ 1 & x_3 & x_3^2 \end{bmatrix} = \begin{bmatrix} 1 & -1 & (-1)^2 \\ 1 & 0 & 0^2 \\ 1 & 1 & 1^2 \end{bmatrix} = \begin{bmatrix} 1 & -1 & 1 \\ 1 & 0 & 0 \\ 1 & 1 & 1 \end{bmatrix}
\end{equation}
The target vector is $\mathbf{y} = [1, 0, 2]^T$.

\textbf{Step 2: Solve the Linear System $\Phi \mathbf{w} = \mathbf{y}$.}
Because $N = p + 1 = 3$, the system is exactly determined. We can solve via row operations:
From row 2: $1 \cdot w_0 + 0 \cdot w_1 + 0 \cdot w_2 = 0 \implies w_0 = 0$.
Substitute $w_0 = 0$ into rows 1 and 3:
\begin{align}
-w_1 + w_2 &= 1 \\
w_1 + w_2 &= 2
\end{align}
Adding the two equations: $2w_2 = 3 \implies w_2 = \frac{3}{2} = 1.5$.
Subtracting the equations: $2w_1 = 1 \implies w_1 = \frac{1}{2} = 0.5$.

\textbf{Step 3: Verification and Out-of-Sample Evaluation.}
The fitted polynomial is:
\begin{equation}
\hat{y}(x) = 0 + \frac{1}{2}x + \frac{3}{2}x^2
\end{equation}
Checking training nodes:
\begin{align}
\hat{y}(-1) &= 0.5(-1) + 1.5(-1)^2 = -0.5 + 1.5 = 1.0 \quad \checkmark \\
\hat{y}(0) &= 0 + 0 = 0.0 \quad \checkmark \\
\hat{y}(1) &= 0.5(1) + 1.5(1)^2 = 0.5 + 1.5 = 2.0 \quad \checkmark
\end{align}
Evaluating out-of-sample at $x = 2$:
\begin{equation}
\hat{y}(2) = 0.5(2) + 1.5(4) = 1 + 6 = 7.0
\end{equation}
\end{examplebox}

% ----------------------------------------------------------------------
% PART 4: SPECTRAL THEORY: EIGENVALUES & EIGENSPACES
% ----------------------------------------------------------------------
\section{Spectral Theory: Eigenvalues, Eigenspaces \& Multiplicity}

\subsection{The Characteristic Polynomial}
To find the invariant directions of an $n \times n$ matrix $A$, we seek non-zero vectors $\mathbf{v}$ and scalars $\lambda$ satisfying:
\begin{equation}
A\mathbf{v} = \lambda \mathbf{v} \iff (A - \lambda I)\mathbf{v} = \mathbf{0} \quad (\mathbf{v} \ne \mathbf{0})
\label{eq:eigen_fundamental}
\end{equation}
Equation~\eqref{eq:eigen_fundamental} states that the matrix $(A - \lambda I)$ maps a non-zero vector to zero. Therefore, $(A - \lambda I)$ must have a non-trivial nullspace, which means it cannot be invertible:
\begin{equation}
\det(A - \lambda I) = 0
\label{eq:char_eq}
\end{equation}
Expanding this determinant yields the \textbf{characteristic polynomial} $p_A(\lambda)$ of degree $n$:
\begin{equation}
p_A(\lambda) = (-1)^n \lambda^n + (-1)^{n-1} \trace(A) \lambda^{n-1} + \dots + \det(A) = 0
\end{equation}

\subsection{Algebraic vs. Geometric Multiplicity}

\begin{definitionbox}{Algebraic and Geometric Multiplicity}
Let $\lambda_0$ be an eigenvalue of $A \in \mathbb{R}^{n \times n}$:
\begin{enumerate}
    \item \textbf{Algebraic Multiplicity} $am(\lambda_0)$: The power of $(\lambda - \lambda_0)$ in the factorization of the characteristic polynomial:
    \begin{equation}
    p_A(\lambda) = (-1)^n (\lambda - \lambda_0)^{am(\lambda_0)} \prod_{j} (\lambda - \lambda_j)^{am(\lambda_j)}
    \end{equation}
    \item \textbf{Geometric Multiplicity} $gm(\lambda_0)$: The dimension of the eigenspace $E_{\lambda_0} = \Null(A - \lambda_0 I)$:
    \begin{equation}
    gm(\lambda_0) \triangleq \dim(E_{\lambda_0}) = n - \rank(A - \lambda_0 I)
    \end{equation}
\end{enumerate}
\end{definitionbox}

\begin{theorem}[Fundamental Multiplicity Inequality]
For every eigenvalue $\lambda$ of an $n \times n$ matrix $A$:
\begin{equation}
1 \le gm(\lambda) \le am(\lambda) \le n
\label{eq:mult_inequality}
\end{equation}
\end{theorem}

\begin{proof}
Because $\lambda$ is an eigenvalue, there exists at least one non-zero eigenvector, so $\dim(\Null(A - \lambda I)) \ge 1$, proving $gm(\lambda) \ge 1$.

To prove $gm(\lambda) \le am(\lambda)$, let $k = gm(\lambda)$, and choose an orthonormal basis $\{\mathbf{v}_1, \dots, \mathbf{v}_k\}$ for $E_\lambda$. Extend this to an orthonormal basis $\{\mathbf{v}_1, \dots, \mathbf{v}_k, \mathbf{u}_1, \dots, \mathbf{u}_{n-k}\}$ for all of $\mathbb{C}^n$. Assemble these vectors as columns of a unitary matrix $Q = [V_1 \mid U]$.

Because $A\mathbf{v}_i = \lambda \mathbf{v}_i$ for each $i \in \{1, \dots, k\}$, we have $A V_1 = \lambda V_1$.
Transforming $A$ by $Q^H (\cdot) Q$:
\begin{equation}
Q^H A Q = \begin{bmatrix} V_1^H \\ U^H \end{bmatrix} A \begin{bmatrix} V_1 & U \end{bmatrix} = \begin{bmatrix} V_1^H A V_1 & V_1^H A U \\ U^H A V_1 & U^H A U \end{bmatrix} = \begin{bmatrix} \lambda I_k & B \\ 0 & C \end{bmatrix}
\end{equation}
The characteristic polynomial of $Q^H A Q$ (which equals $p_A(t)$ because similar matrices share characteristic polynomials) is:
\begin{equation}
\det(Q^H A Q - t I_n) = \det\begin{bmatrix} (\lambda - t)I_k & B \\ 0 & C - t I_{n-k} \end{bmatrix} = (\lambda - t)^k \det(C - t I_{n-k})
\end{equation}
Thus, $(\lambda - t)^k$ is a factor of $p_A(t)$, which proves that the algebraic multiplicity of $\lambda$ must be at least $k = gm(\lambda)$.
\end{proof}

\begin{ponderbox}[The Geometric Multiplicity Bound]
Suppose a $4 \times 4$ matrix has characteristic polynomial $p(\lambda) = (\lambda - 3)^3 (\lambda + 1)$.
\begin{enumerate}
    \item What is the algebraic multiplicity of $\lambda = 3$?
    \item What are the possible values for the geometric multiplicity $gm(3)$?
    \item What is the geometric multiplicity of $\lambda = -1$? Can it ever equal 2? Why or why not?
\end{enumerate}
\end{ponderbox}

\begin{examplebox}{Diagnosing a Defective Shear Matrix}
Let matrix $B \in \mathbb{R}^{2 \times 2}$ be the horizontal shear operator:
\begin{equation}
B = \begin{bmatrix} 3 & 1 \\ 0 & 3 \end{bmatrix}
\end{equation}
Compute its characteristic polynomial, eigenvalues, algebraic and geometric multiplicities, and explain why $B$ is not diagonalizable.

\textbf{Step 1: Compute Characteristic Polynomial.}
\begin{equation}
p_B(\lambda) = \det(B - \lambda I) = \det\begin{bmatrix} 3 - \lambda & 1 \\ 0 & 3 - \lambda \end{bmatrix} = (3 - \lambda)^2 = 0
\end{equation}
The only eigenvalue is $\lambda = 3$, with algebraic multiplicity $am(3) = 2$.

\textbf{Step 2: Compute Eigenspace $E_3 = \Null(B - 3I)$.}
\begin{equation}
B - 3I = \begin{bmatrix} 3 - 3 & 1 \\ 0 & 3 - 3 \end{bmatrix} = \begin{bmatrix} 0 & 1 \\ 0 & 0 \end{bmatrix}
\end{equation}
Setting $(B - 3I)\mathbf{v} = \mathbf{0}$:
\begin{equation}
\begin{bmatrix} 0 & 1 \\ 0 & 0 \end{bmatrix} \begin{bmatrix} v_1 \\ v_2 \end{bmatrix} = \begin{bmatrix} 0 \\ 0 \end{bmatrix} \implies 0 \cdot v_1 + 1 \cdot v_2 = 0 \implies v_2 = 0
\end{equation}
The first coordinate $v_1$ is completely free. Thus, any eigenvector must be of the form $\mathbf{v} = [v_1, 0]^T = v_1 [1, 0]^T$.
The eigenspace is:
\begin{equation}
E_3 = \Span\left\{ \begin{bmatrix} 1 \\ 0 \end{bmatrix} \right\} \implies gm(3) = \dim(E_3) = 1
\end{equation}

\textbf{Step 3: Diagnosis of Defectiveness.}
Because $gm(3) = 1 < am(3) = 2$, matrix $B$ is \textbf{defective}. It possesses only one linearly independent eigenvector in $\mathbb{R}^2$. It is physically impossible to construct an invertible $2 \times 2$ modal matrix $P = [\mathbf{v}_1, \mathbf{v}_2]$, so $B$ \textbf{cannot be diagonalized}!
\end{examplebox}

% ----------------------------------------------------------------------
% PART 5: MATRIX DIAGONALIZATION & DYNAMICAL SYSTEMS
% ----------------------------------------------------------------------
\section{Matrix Diagonalization \& Iterated Powers}

\begin{theorem}[The Diagonalization Theorem]
An $n \times n$ matrix $A$ is diagonalizable if and only if the sum of geometric multiplicities equals $n$:
\begin{equation}
\sum_{i} gm(\lambda_i) = n \iff gm(\lambda_i) = am(\lambda_i) \quad \text{for all } i
\end{equation}
When diagonalizable, $A = P D P^{-1}$ where:
\begin{equation}
P = \begin{bmatrix} \mathbf{v}_1 & \mathbf{v}_2 & \dots & \mathbf{v}_n \end{bmatrix}, \quad D = \begin{bmatrix} \lambda_1 & & \\ & \ddots & \\ & & \lambda_n \end{bmatrix}
\label{eq:diagonalization_decomp}
\end{equation}
\end{theorem}

\begin{figure}[htbp]
\centering
\begin{tikzpicture}[
    scale=0.9,
    font=\sffamily\footnotesize
]

% Left: Original Unit Circle
\draw[thick, fill=primarylight!30] (-4.5, 0) circle (1.8cm);
\draw[->, >=stealth, color=gray] (-6.8, 0) -- (-2.2, 0) node[right] {$x_1$};
\draw[->, >=stealth, color=gray] (-4.5, -2.3) -- (-4.5, 2.3) node[above] {$x_2$};

% Eigenvectors on unit circle
\draw[->, >=stealth, line width=1.6pt, color=tealborder] (-4.5, 0) -- (-3.0, 0.9) node[above right] {$\mathbf{v}_1$};
\draw[->, >=stealth, line width=1.6pt, color=coralborder] (-4.5, 0) -- (-5.4, 1.5) node[above left] {$\mathbf{v}_2$};
\node[below] at (-4.5, -2.4) {\textbf{Input Space: } Unit Sphere $\|\mathbf{x}\| = 1$};

% Transformation arrow
\draw[->, >=stealth, line width=2pt, color=primarydark] (-1.8, 0) -- (0.8, 0) node[midway, above] {$\mathbf{x} \mapsto A\mathbf{x}$};

% Right: Transformed Ellipse
\begin{scope}[shift={(4.0, 0)}, rotate=30]
    \draw[thick, fill=tealbg!40] (0, 0) ellipse (3.2cm and 1.2cm);
    % Stretched eigenvectors
    \draw[->, >=stealth, line width=1.8pt, color=tealborder] (0, 0) -- (3.2, 0) node[above right] {$A\mathbf{v}_1 = \lambda_1 \mathbf{v}_1$};
    \draw[->, >=stealth, line width=1.8pt, color=coralborder] (0, 0) -- (0, 1.2) node[above left] {$A\mathbf{v}_2 = \lambda_2 \mathbf{v}_2$};
\end{scope}
\node[below] at (4.0, -2.4) {\textbf{Output Space: } Ellipsoid with axes $\lambda_i \mathbf{v}_i$};

\end{tikzpicture}
\caption{The Geometry of Eigendecomposition: The linear transformation $A$ deforms the unit circle into an ellipse whose principal axes align with the eigenvectors $\mathbf{v}_i$ scaled by eigenvalues $\lambda_i$.}
\label{fig:eigen_geometry}
\end{figure}

\begin{examplebox}{Step-by-Step Matrix Eigendecomposition and Power $A^k$}
Let $A = \begin{bmatrix} 4 & 2 \\ 1 & 3 \end{bmatrix}$. 
\begin{enumerate}
    \item Find eigenvalues $\lambda_1, \lambda_2$.
    \item Find eigenvectors $\mathbf{v}_1, \mathbf{v}_2$ and construct $P$ and $P^{-1}$.
    \item Compute the closed-form formula for $A^k$ and evaluate $A^5$.
\end{enumerate}

\textbf{Step 1: Compute Eigenvalues.}
\begin{equation}
\det(A - \lambda I) = \det\begin{bmatrix} 4 - \lambda & 2 \\ 1 & 3 - \lambda \end{bmatrix} = (4 - \lambda)(3 - \lambda) - 2 = \lambda^2 - 7\lambda + 10 = (\lambda - 5)(\lambda - 2) = 0
\end{equation}
The eigenvalues are $\lambda_1 = 5$ and $\lambda_2 = 2$. Both are distinct, guaranteeing that $A$ is diagonalizable.

\textbf{Step 2: Compute Eigenvectors.}
For $\lambda_1 = 5$:
\begin{equation}
(A - 5I)\mathbf{v}_1 = \begin{bmatrix} -1 & 2 \\ 1 & -2 \end{bmatrix} \begin{bmatrix} v_{11} \\ v_{12} \end{bmatrix} = \begin{bmatrix} 0 \\ 0 \end{bmatrix} \implies -v_{11} + 2v_{12} = 0 \implies \mathbf{v}_1 = \begin{bmatrix} 2 \\ 1 \end{bmatrix}
\end{equation}
For $\lambda_2 = 2$:
\begin{equation}
(A - 2I)\mathbf{v}_2 = \begin{bmatrix} 2 & 2 \\ 1 & 1 \end{bmatrix} \begin{bmatrix} v_{21} \\ v_{22} \end{bmatrix} = \begin{bmatrix} 0 \\ 0 \end{bmatrix} \implies v_{21} + v_{22} = 0 \implies \mathbf{v}_2 = \begin{bmatrix} 1 \\ -1 \end{bmatrix}
\end{equation}

\textbf{Step 3: Construct Modal Matrix $P$ and Invert.}
\begin{equation}
P = \begin{bmatrix} \mathbf{v}_1 & \mathbf{v}_2 \end{bmatrix} = \begin{bmatrix} 2 & 1 \\ 1 & -1 \end{bmatrix}, \quad \det(P) = 2(-1) - 1(1) = -3
\end{equation}
\begin{equation}
P^{-1} = \frac{1}{-3} \begin{bmatrix} -1 & -1 \\ -1 & 2 \end{bmatrix} = \frac{1}{3} \begin{bmatrix} 1 & 1 \\ 1 & -2 \end{bmatrix}
\end{equation}

\textbf{Step 4: Compute General Matrix Power $A^k = P D^k P^{-1}$.}
\begin{align}
A^k &= \begin{bmatrix} 2 & 1 \\ 1 & -1 \end{bmatrix} \begin{bmatrix} 5^k & 0 \\ 0 & 2^k \end{bmatrix} \left( \frac{1}{3} \begin{bmatrix} 1 & 1 \\ 1 & -2 \end{bmatrix} \right) \\
&= \frac{1}{3} \begin{bmatrix} 2 \cdot 5^k & 2^k \\ 5^k & -2^k \end{bmatrix} \begin{bmatrix} 1 & 1 \\ 1 & -2 \end{bmatrix} \\
&= \frac{1}{3} \begin{bmatrix} 2 \cdot 5^k + 2^k & 2 \cdot 5^k - 2 \cdot 2^k \\ 5^k - 2^k & 5^k + 2 \cdot 2^k \end{bmatrix}
\end{align}

\textbf{Step 5: Specific Evaluation for $k = 5$.}
Using $5^5 = 3125$ and $2^5 = 32$:
\begin{align}
(A^5)_{11} &= \frac{2(3125) + 32}{3} = \frac{6250 + 32}{3} = \frac{6282}{3} = 2094 \\
(A^5)_{12} &= \frac{2(3125) - 2(32)}{3} = \frac{6250 - 64}{3} = \frac{6186}{3} = 2062 \\
(A^5)_{21} &= \frac{3125 - 32}{3} = \frac{3093}{3} = 1031 \\
(A^5)_{22} &= \frac{3125 + 2(32)}{3} = \frac{3125 + 64}{3} = \frac{3189}{3} = 1063
\end{align}
Thus:
\begin{equation}
A^5 = \begin{bmatrix} 2094 & 2062 \\ 1031 & 1063 \end{bmatrix}
\end{equation}
Notice how the dominant eigenvalue $\lambda_1 = 5$ completely dictates the asymptotic scale of $A^k$ as $k \to \infty$!
\end{examplebox}

\begin{ponderbox}[The Spectral Radius and Dynamical Stability]
Let $\mathbf{x}_{t+1} = W \mathbf{x}_t$ describe the hidden state evolution of a linear recurrent neural network (RNN). The \textbf{spectral radius} is $\rho(W) \triangleq \max_i |\lambda_i|$.
\begin{enumerate}
    \item If $\rho(W) > 1$, what happens to $\|\mathbf{x}_t\|_2$ as $t \to \infty$? Why does this cause \textbf{exploding activations}?
    \item If $\rho(W) < 1$, what happens to $\|\mathbf{x}_t\|_2$ as $t \to \infty$? Why does this cause \textbf{vanishing gradients}?
    \item Why do modern recurrent architectures (e.g. orthogonal RNNs, unitary RNNs) initialize weights so that $\rho(W) \approx 1$?
\end{enumerate}
\end{ponderbox}

\subsection{Solving Discrete Dynamical Systems: Derivation of Binet's Fibonacci Formula}
The Fibonacci recurrence is defined as $F_{k+2} = F_{k+1} + F_k$ with initial conditions $F_0 = 0, F_1 = 1$.
We can express this second-order recurrence as a first-order two-dimensional linear dynamical system:
\begin{equation}
\mathbf{u}_{k+1} \triangleq \begin{bmatrix} F_{k+1} \\ F_k \end{bmatrix} = \begin{bmatrix} 1 & 1 \\ 1 & 0 \end{bmatrix} \begin{bmatrix} F_k \\ F_{k-1} \end{bmatrix} = A \mathbf{u}_k
\label{eq:fib_dyn}
\end{equation}
Iterating $k$ times from initial state $\mathbf{u}_0 = \begin{bmatrix} 1 \\ 0 \end{bmatrix}$:
\begin{equation}
\mathbf{u}_k = A^k \mathbf{u}_0 = A^k \begin{bmatrix} 1 \\ 0 \end{bmatrix}
\end{equation}
To compute $A^k$, we diagonalize $A = \begin{bmatrix} 1 & 1 \\ 1 & 0 \end{bmatrix}$:
\begin{equation}
\det(A - \lambda I) = \det\begin{bmatrix} 1 - \lambda & 1 \\ 1 & -\lambda \end{bmatrix} = \lambda^2 - \lambda - 1 = 0
\end{equation}
The eigenvalues are the Golden Ratio $\phi$ and its conjugate:
\begin{equation}
\lambda_1 = \frac{1 + \sqrt{5}}{2} = \phi \approx 1.618, \quad \lambda_2 = \frac{1 - \sqrt{5}}{2} = 1 - \phi = -\frac{1}{\phi} \approx -0.618
\end{equation}
Finding eigenvectors: $(A - \lambda_i I)\mathbf{v}_i = \mathbf{0} \implies \mathbf{v}_i = \begin{bmatrix} \lambda_i \\ 1 \end{bmatrix}$.
Thus:
\begin{equation}
P = \begin{bmatrix} \lambda_1 & \lambda_2 \\ 1 & 1 \end{bmatrix}, \quad P^{-1} = \frac{1}{\lambda_1 - \lambda_2} \begin{bmatrix} 1 & -\lambda_2 \\ -1 & \lambda_1 \end{bmatrix} = \frac{1}{\sqrt{5}} \begin{bmatrix} 1 & -\lambda_2 \\ -1 & \lambda_1 \end{bmatrix}
\end{equation}
Now, expanding $\mathbf{u}_k = P D^k P^{-1} \mathbf{u}_0$:
\begin{align}
\begin{bmatrix} F_{k+1} \\ F_k \end{bmatrix} &= \frac{1}{\sqrt{5}} \begin{bmatrix} \lambda_1 & \lambda_2 \\ 1 & 1 \end{bmatrix} \begin{bmatrix} \lambda_1^k & 0 \\ 0 & \lambda_2^k \end{bmatrix} \begin{bmatrix} 1 & -\lambda_2 \\ -1 & \lambda_1 \end{bmatrix} \begin{bmatrix} 1 \\ 0 \end{bmatrix} \nonumber \\
&= \frac{1}{\sqrt{5}} \begin{bmatrix} \lambda_1 & \lambda_2 \\ 1 & 1 \end{bmatrix} \begin{bmatrix} \lambda_1^k \\ -\lambda_2^k \end{bmatrix} = \frac{1}{\sqrt{5}} \begin{bmatrix} \lambda_1^{k+1} - \lambda_2^{k+1} \\ \lambda_1^k - \lambda_2^k \end{bmatrix}
\end{align}
Reading the second component yields Binet's celebrated formula:
\begin{equation}
F_k = \frac{1}{\sqrt{5}} \left( \left(\frac{1 + \sqrt{5}}{2}\right)^k - \left(\frac{1 - \sqrt{5}}{2}\right)^k \right)
\label{eq:binet}
\end{equation}

% ----------------------------------------------------------------------
% PART 6: EXPLICIT GOTCHAS & DIAGNOSTIC MISCONCEPTIONS
% ----------------------------------------------------------------------
\section{Explicit Gotchas \& Diagnostic Misconceptions}

\begin{trapbox}{Gotcha 1: Assuming Every Square Matrix is Diagonalizable}
A practitioner encounters matrix $J = \begin{bmatrix} 3 & 1 \\ 0 & 3 \end{bmatrix}$ and attempts to compute $J^{100} = P D^{100} P^{-1}$.
\textbf{The Trap:} The characteristic polynomial is $(\lambda - 3)^2 = 0$, giving eigenvalue $\lambda = 3$ with algebraic multiplicity $am(3) = 2$. However, the eigenspace is $\Null(J - 3I) = \Null\begin{bmatrix} 0 & 1 \\ 0 & 0 \end{bmatrix} = \Span\left\{\begin{bmatrix} 1 \\ 0 \end{bmatrix}\right\}$, which has dimension $gm(3) = 1$. Because $gm(3) < am(3)$, the matrix is \textbf{defective} (non-diagonalizable) and lacks a complete basis of eigenvectors. To compute powers of defective matrices, one must use the Jordan Normal Form.
\end{trapbox}

\begin{trapbox}{Gotcha 2: High-Degree Polynomial Regression and Runge's Phenomenon}
To achieve zero training error on $N = 15$ data points, a student fits a degree-$14$ polynomial, observing wild oscillations near the boundary edges of the interval $[-1, 1]$.
\textbf{The Trap:} This is \textbf{Runge's phenomenon}. In equispaced polynomial interpolation, the condition number of the Vandermonde matrix grows exponentially: $\kappa(\Phi) \sim \mathcal{O}(e^p)$. The resulting high-degree polynomial exhibits catastrophic out-of-sample test error. In production machine learning, high-degree polynomial expansions must be paired with Ridge regularization ($\lambda \|\mathbf{w}\|_2^2$), Chebyshev-spaced nodes, or replaced by local splines.
\end{trapbox}

\begin{trapbox}{Gotcha 3: Believing Real Matrices Always Have Real Eigenvalues}
A programmer computes the eigenvalues of a $2 \times 2$ 90-degree counterclockwise rotation matrix $R = \begin{bmatrix} 0 & -1 \\ 1 & 0 \end{bmatrix}$ and expects real scale factors.
\textbf{The Trap:} The characteristic equation is $\det(R - \lambda I) = \lambda^2 + 1 = 0 \implies \lambda = \pm i$. A pure rotation leaves no 1-dimensional real line invariant. Eigenvalues of real matrices are generally \textbf{complex numbers}. Real eigenvalues are guaranteed if and only if the matrix is symmetric ($A = A^T$), as established by the Spectral Theorem (Chapter~\ref{chap:week06}).
\end{trapbox}

\begin{trapbox}{Gotcha 4: Illegally Adding Eigenvalues for Matrix Sums}
An engineer knows the eigenvalues of $A$ are $\{2, 5\}$ and the eigenvalues of $B$ are $\{1, 3\}$. They conclude the eigenvalues of $A + B$ are $\{3, 8\}$.
\textbf{The Trap:} In general, $\lambda_i(A + B) \ne \lambda_i(A) + \lambda_i(B)$. The eigenvalues of a sum equal the sum of eigenvalues if and only if $A$ and $B$ share a simultaneous basis of eigenvectors (which occurs $\iff AB = BA$, meaning the matrices commute). For non-commuting matrices, the eigenvectors mix non-linearly.
\end{trapbox}

% ----------------------------------------------------------------------
% PART 7: COMPARATIVE SYNTHESIS & SUMMARY TABLE
% ----------------------------------------------------------------------
\section{Comparative Synthesis Table}

Table~\ref{tab:matrix_classes} compares the spectral properties across canonical matrix classes.

\begin{table}[htbp]
\centering
\small
\begin{tabularx}{\textwidth}{lXXXX}
\toprule
\textbf{Matrix Class} & \textbf{Symmetry Condition} & \textbf{Eigenvalue Field} & \textbf{Eigenvectors} & \textbf{Diagonalizability} \\
\midrule
General Square & $A \in \mathbb{R}^{n \times n}$ & $\lambda \in \mathbb{C}$ & Linearly independent if distinct & Diagonalizable $\iff gm(\lambda) = am(\lambda)$ \\
\addlinespace
Symmetric & $A = A^T$ & $\lambda \in \mathbb{R}$ strictly & Mutually orthogonal ($Q^T Q = I$) & Always orthogonally diagonalizable ($Q \Lambda Q^T$) \\
\addlinespace
Orthogonal & $Q^T Q = I$ & $|\lambda| = 1$ ($\lambda = e^{i\theta}$) & Mutually orthogonal in $\mathbb{C}^n$ & Always diagonalizable over $\mathbb{C}$ \\
\addlinespace
Projection & $P^2 = P = P^T$ & $\lambda \in \{0, 1\}$ strictly & Orthogonal bases for $V$ and $V^\perp$ & Always diagonalizable; $\rank(P) = \trace(P)$ \\
\addlinespace
Defective & $gm(\lambda) < am(\lambda)$ & $\lambda \in \mathbb{C}$ & Incomplete spanning set ($< n$) & Not diagonalizable; requires Jordan Normal Form \\
\bottomrule
\end{tabularx}
\caption{Comparative Synthesis of Spectral Properties Across Matrix Classes.}
\label{tab:matrix_classes}
\end{table}

% ----------------------------------------------------------------------
% PART 8: PYTHON / SCIENTIFIC COMPUTING LAB
% ----------------------------------------------------------------------
\section{Python / Scientific Computing Lab}

The production script below implements polynomial regression, evaluates the condition number explosion of Vandermonde matrices, and verifies Binet's Fibonacci formula via matrix diagonalization.

\begin{lstlisting}[language=Python, caption={Polynomial Regression and Spectral Diagonalization in Python.}]
import numpy as np

def fit_polynomial(x: np.ndarray, y: np.ndarray, degree: int):
    """Fits polynomial regression via Vandermonde expansion and evaluates condition number."""
    Phi = np.vander(x, N=degree + 1, increasing=True)
    cond_num = np.linalg.cond(Phi)
    # Solve via QR for numerical stability
    Q, R = np.linalg.qr(Phi)
    w = np.linalg.solve(R, Q.T @ y)
    return w, cond_num

def fibonacci_via_diagonalization(k: int) -> int:
    """Computes k-th Fibonacci number via spectral eigendecomposition of transition matrix."""
    A = np.array([[1.0, 1.0], [1.0, 0.0]])
    eigvals, P = np.linalg.eig(A)

    # Sort eigenvalues descending so lambda_1 is the Golden Ratio
    idx = np.argsort(eigvals)[::-1]
    eigvals = eigvals[idx]
    P = P[:, idx]

    # Compute A^k = P @ diag(lambda^k) @ P^{-1}
    D_k = np.diag(eigvals ** k)
    P_inv = np.linalg.inv(P)
    A_k = P @ D_k @ P_inv

    u0 = np.array([1.0, 0.0])
    u_k = A_k @ u0
    return int(np.round(u_k[1]))

def main():
    # 1. Vandermonde condition number explosion
    x_sample = np.linspace(-1.0, 1.0, 20)
    print("Vandermonde Matrix Condition Number vs. Degree:")
    for deg in [2, 5, 10, 15]:
        Phi = np.vander(x_sample, N=deg + 1, increasing=True)
        print(f"  Degree {deg:2d}: kappa(Phi) = {np.linalg.cond(Phi):.2e}")

    # 2. Fibonacci closed-form verification
    print("\nFibonacci Sequence via Matrix Powers:")
    for k in range(1, 11):
        f_k = fibonacci_via_diagonalization(k)
        print(f"  F_{k:02d} = {f_k}")

if __name__ == "__main__":
    main()
\end{lstlisting}

% ----------------------------------------------------------------------
% PART 9: COMPREHENSIVE PRACTICE PROBLEMS
% ----------------------------------------------------------------------
\section{Comprehensive Practice Problems}

\subsection*{Part A: Analytical \& Mathematical Problems}

\begin{enumerate}[label=\textbf{P\arabic*.}]
    \item \textbf{Eigendecomposition of a $2 \times 2$ Matrix:}
    Let $A = \begin{bmatrix} 4 & 2 \\ 1 & 3 \end{bmatrix}$.
    \begin{enumerate}[label=(\alph*)]
        \item Compute the characteristic equation $p_A(\lambda) = \det(A - \lambda I) = 0$ and find all eigenvalues.
        \item For each eigenvalue, find a corresponding eigenvector.
        \item Construct the modal matrix $P$ and diagonal matrix $D$ such that $A = PDP^{-1}$.
        \item Use this diagonalization to compute the exact matrix $A^5$.
    \end{enumerate}

    \item \textbf{Defective Matrix and Algebraic vs. Geometric Multiplicity:}
    Let $B = \begin{bmatrix} 2 & 1 & 0 \\ 0 & 2 & 0 \\ 0 & 0 & 5 \end{bmatrix}$.
    \begin{enumerate}[label=(\alph*)]
        \item Find the eigenvalues of $B$ and their algebraic multiplicities.
        \item Compute the eigenspace $E_\lambda = \Null(B - \lambda I)$ for each eigenvalue.
        \item Determine the geometric multiplicity of each eigenvalue.
        \item State whether $B$ is diagonalizable, justifying your answer rigorously.
    \end{enumerate}

    \item \textbf{Spectral Theorem for Symmetric Matrices:}
    Let $S \in \mathbb{R}^{n \times n}$ be a real symmetric matrix ($S = S^T$).
    \begin{enumerate}[label=(\alph*)]
        \item Prove that all eigenvalues of $S$ are real ($\lambda \in \mathbb{R}$).
        \item Prove that eigenvectors corresponding to distinct eigenvalues are mutually orthogonal: if $\lambda_1 \ne \lambda_2$, then $\mathbf{v}_1^T \mathbf{v}_2 = 0$.
    \end{enumerate}

    \item \textbf{Cayley-Hamilton Theorem Application:}
    The Cayley-Hamilton Theorem states that every square matrix satisfies its own characteristic equation: $p_A(A) = 0$.
    For matrix $A = \begin{bmatrix} 2 & 1 \\ 3 & 4 \end{bmatrix}$:
    \begin{enumerate}[label=(\alph*)]
        \item Find the characteristic polynomial $p_A(\lambda)$.
        \item Verify directly that $A^2 - \trace(A)A + \det(A)I = 0$.
        \item Use the Cayley-Hamilton equation to express $A^{-1}$ as a linear combination of $A$ and $I$.
    \end{enumerate}

    \item \textbf{Markov Transition Matrix Dynamics:}
    Consider a 2-state Markov transition matrix $T = \begin{bmatrix} 0.8 & 0.3 \\ 0.2 & 0.7 \end{bmatrix}$.
    \begin{enumerate}[label=(\alph*)]
        \item Prove that $\lambda_1 = 1$ is always an eigenvalue for any column-stochastic matrix.
        \item Find the second eigenvalue $\lambda_2$.
        \item Diagonalize $T = P D P^{-1}$ and compute the steady-state distribution $\boldsymbol{\pi} = \lim_{k \to \infty} T^k \mathbf{p}_0$ for an arbitrary initial probability vector $\mathbf{p}_0$.
    \end{enumerate}
\end{enumerate}

\subsection*{Part B: Diagnostic Multiple-Choice Questions}

\begin{enumerate}[label=\textbf{MCQ\arabic*.}]
    \item A matrix $A \in \mathbb{R}^{3 \times 3}$ has trace $\trace(A) = 6$ and determinant $\det(A) = 8$. Two of its eigenvalues are $\lambda_1 = 2$ and $\lambda_2 = 4$. What is the third eigenvalue $\lambda_3$, and is $A$ invertible?
    \begin{enumerate}[label=(\alph*)]
        \item $\lambda_3 = 0$; $A$ is singular.
        \item $\lambda_3 = 1$; $A$ is invertible.
        \item $\lambda_3 = 0$; $A$ is invertible.
        \item $\lambda_3 = -2$; $A$ is invertible.
    \end{enumerate}

    \item Why does fitting a degree-$p$ polynomial with $p = 15$ to $N = 16$ data points via Ordinary Least Squares typically perform poorly on new test data, even when training error is zero?
    \begin{enumerate}[label=(\alph*)]
        \item Because the Vandermonde matrix is non-invertible for $p < N$.
        \item Because high-degree polynomials have high bias and cannot fit the training points accurately.
        \item Because the model has excessive variance and overfits to training noise, exhibiting Runge's phenomenon at interval boundaries.
        \item Because polynomial regression requires the response variable $y$ to follow a Gaussian distribution.
    \end{enumerate}
\end{enumerate}

\chapter{Complex Vector Spaces, Hermitian \& Unitary Matrices, and the Spectral Theorem}
\label{chap:week05}

% ----------------------------------------------------------------------
% PART 1: EXECUTIVE ROADMAP & LEARNING OBJECTIVES
% ----------------------------------------------------------------------
\section{Executive Roadmap \& Learning Objectives}

While data vectors in machine learning frequently consist of real measurements, the algebraic closure of polynomials guarantees that real matrices can possess complex eigenvalues and complex eigenvectors. Extending linear algebra to complex vector spaces $\mathbb{C}^n$ provides the complete mathematical framework needed to understand spectral theory. In this domain, the familiar concepts of matrix transpose and symmetry generalize into the \textbf{conjugate transpose} and \textbf{Hermitian matrices}.

\subsection{Prerequisites \& Scaffolding}
Mastery of this chapter builds upon vector spaces, inner products, and single-variable complex numbers ($z = a + ib$), as shown in Table~\ref{tab:w05_prereqs}.

\begin{table}[htbp]
\centering
\small
\begin{tabularx}{\textwidth}{lXX}
\toprule
\textbf{Foundational Area} & \textbf{Core Concepts Required} & \textbf{Role in Machine Learning Foundations} \\
\midrule
Complex Numbers & Complex conjugate $\bar{z}$, modulus $|z|^2 = z\bar{z}$, Euler's formula $e^{i\theta}$ & Defining norms and scalar fields over $\mathbb{C}^n$ \\
Linear Algebra & Eigenvalue characteristic equations, orthogonality, matrix rank & Formulating unitary transformations and spectral projections \\
Matrix Decompositions & Diagonalization $A = PDP^{-1}$, orthogonal matrices $Q^T Q = I$ & Foundation for the complex Spectral Theorem and SVD \\
\bottomrule
\end{tabularx}
\caption{Prerequisite Scaffolding for Chapter 5.}
\label{tab:w05_prereqs}
\end{table}

\subsection{Core Competencies}
By the conclusion of this chapter, students will be able to:
\begin{itemize}
    \item Perform matrix algebra with the conjugate transpose (Hermitian transpose) $A^H$.
    \item Define and characterize Hermitian, Skew-Hermitian, and Unitary matrices.
    \item Rigorously prove that all eigenvalues of a Hermitian matrix are strictly real.
    \item State and prove the Spectral Theorem for Hermitian matrices, expressing operators as sums of rank-1 projectors.
    \item Formulate the Rayleigh quotient and derive the Courant-Fischer min-max characterization of eigenvalues.
\end{itemize}

\subsection{Notation \& Mathematical Conventions}
\begin{itemize}
    \item $\mathbb{C}^n$: The $n$-dimensional complex vector space.
    \item $A^H \triangleq \overline{A}^T$: The conjugate transpose (Hermitian adjoint), where $[A^H]_{ij} = \overline{A_{ji}}$.
    \item $\langle \mathbf{u}, \mathbf{v} \rangle \triangleq \mathbf{u}^H \mathbf{v} = \sum_{i=1}^n \bar{u}_i v_i$: The standard complex inner product.
    \item $\|\mathbf{z}\|_2 \triangleq \sqrt{\mathbf{z}^H \mathbf{z}} = \sqrt{\sum_{i=1}^n |z_i|^2}$: The Euclidean norm on $\mathbb{C}^n$.
    \item $U(n) \triangleq \{U \in \mathbb{C}^{n \times n} \mid U^H U = I\}$: The unitary group of order $n$.
    \item $R_A(\mathbf{x}) \triangleq \frac{\mathbf{x}^H A \mathbf{x}}{\mathbf{x}^H \mathbf{x}}$: The Rayleigh quotient of matrix $A$.
\end{itemize}

% ----------------------------------------------------------------------
% PART 2: MATHEMATICAL & COMPUTATIONAL FOUNDATIONS
% ----------------------------------------------------------------------
\section{Mathematical \& Computational Foundations}

\subsection{Complex Inner Products and Adjoints}

\begin{definitionbox}{Complex Inner Product}
An inner product on a complex vector space $V$ is a mapping $\langle \cdot, \cdot \rangle: V \times V \to \mathbb{C}$ satisfying for all $\mathbf{u}, \mathbf{v}, \mathbf{w} \in V$ and $\alpha, \beta \in \mathbb{C}$:
\begin{enumerate}
    \item \textbf{Conjugate Symmetry:} $\langle \mathbf{u}, \mathbf{v} \rangle = \overline{\langle \mathbf{v}, \mathbf{u} \rangle}$.
    \item \textbf{Linearity in Second Argument:} $\langle \mathbf{u}, \alpha \mathbf{v} + \beta \mathbf{w} \rangle = \alpha \langle \mathbf{u}, \mathbf{v} \rangle + \beta \langle \mathbf{u}, \mathbf{w} \rangle$.
    \item \textbf{Positive Definiteness:} $\langle \mathbf{u}, \mathbf{u} \rangle \ge 0$, and $\langle \mathbf{u}, \mathbf{u} \rangle = 0 \iff \mathbf{u} = \mathbf{0}$.
\end{enumerate}
\end{definitionbox}

\begin{definitionbox}{Hermitian and Unitary Matrices}
Let $A \in \mathbb{C}^{n \times n}$ and $U \in \mathbb{C}^{n \times n}$:
\begin{itemize}
    \item $A$ is \textbf{Hermitian} (self-adjoint) if and only if $A^H = A$.
    \item $A$ is \textbf{Skew-Hermitian} if and only if $A^H = -A$.
    \item $U$ is \textbf{Unitary} if and only if $U^H U = U U^H = I$.
\end{itemize}
\end{definitionbox}

\subsection{Spectral Theorems for Hermitian Matrices}

\begin{theorem}[Real Eigenvalues of Hermitian Matrices]
All eigenvalues of a Hermitian matrix $A \in \mathbb{C}^{n \times n}$ are strictly real numbers: $\lambda \in \mathbb{R}$.
\end{theorem}

\begin{proof}
Let $(\lambda, \mathbf{v})$ be an eigenpair of $A$, with $A\mathbf{v} = \lambda \mathbf{v}$ and $\mathbf{v} \ne \mathbf{0}$.
Pre-multiply by $\mathbf{v}^H$:
\begin{equation}
\mathbf{v}^H A \mathbf{v} = \mathbf{v}^H (\lambda \mathbf{v}) = \lambda (\mathbf{v}^H \mathbf{v}) = \lambda \|\mathbf{v}\|_2^2
\label{eq:herm_proof1}
\end{equation}
Taking the conjugate transpose of the scalar $\mathbf{v}^H A \mathbf{v}$:
\begin{equation}
(\mathbf{v}^H A \mathbf{v})^H = \mathbf{v}^H A^H (\mathbf{v}^H)^H = \mathbf{v}^H A \mathbf{v}
\end{equation}
Thus, the complex scalar $\mathbf{v}^H A \mathbf{v}$ is equal to its own complex conjugate, meaning it is strictly real: $\mathbf{v}^H A \mathbf{v} \in \mathbb{R}$.
From equation~\eqref{eq:herm_proof1}, we solve for $\lambda$:
\begin{equation}
\lambda = \frac{\mathbf{v}^H A \mathbf{v}}{\|\mathbf{v}\|_2^2}
\end{equation}
Since both the numerator $\mathbf{v}^H A \mathbf{v}$ and denominator $\|\mathbf{v}\|_2^2 > 0$ are real, $\lambda \in \mathbb{R}$.
\end{proof}

\begin{theorem}[Orthogonality of Eigenvectors]
Let $A$ be Hermitian. Eigenvectors corresponding to distinct eigenvalues $\lambda_1 \ne \lambda_2$ are mutually orthogonal under the complex inner product: $\mathbf{v}_1^H \mathbf{v}_2 = 0$.
\end{theorem}

\begin{proof}
Let $A\mathbf{v}_1 = \lambda_1 \mathbf{v}_1$ and $A\mathbf{v}_2 = \lambda_2 \mathbf{v}_2$. Consider the quantity $\mathbf{v}_1^H A \mathbf{v}_2$:
\begin{equation}
\mathbf{v}_1^H A \mathbf{v}_2 = \mathbf{v}_1^H (\lambda_2 \mathbf{v}_2) = \lambda_2 (\mathbf{v}_1^H \mathbf{v}_2)
\end{equation}
Applying the adjoint property:
\begin{equation}
\mathbf{v}_1^H A \mathbf{v}_2 = (A^H \mathbf{v}_1)^H \mathbf{v}_2 = (A\mathbf{v}_1)^H \mathbf{v}_2 = (\lambda_1 \mathbf{v}_1)^H \mathbf{v}_2 = \bar{\lambda}_1 (\mathbf{v}_1^H \mathbf{v}_2)
\end{equation}
Because $\lambda_1 \in \mathbb{R}$, $\bar{\lambda}_1 = \lambda_1$. Equating the two expressions:
\begin{equation}
\lambda_2 (\mathbf{v}_1^H \mathbf{v}_2) = \lambda_1 (\mathbf{v}_1^H \mathbf{v}_2) \implies (\lambda_1 - \lambda_2)(\mathbf{v}_1^H \mathbf{v}_2) = 0
\end{equation}
Since $\lambda_1 \ne \lambda_2$, we must have $\mathbf{v}_1^H \mathbf{v}_2 = 0$, establishing orthogonality.
\end{proof}

\begin{theorem}[The Spectral Theorem for Hermitian Matrices]
Every Hermitian matrix $A \in \mathbb{C}^{n \times n}$ is unitarily diagonalizable:
\begin{equation}
A = U \Lambda U^H = \sum_{i=1}^n \lambda_i \mathbf{u}_i \mathbf{u}_i^H
\label{eq:spectral_theorem}
\end{equation}
where $U = [\mathbf{u}_1, \dots, \mathbf{u}_n]$ is a unitary matrix ($U^H U = I$), $\Lambda = \diag(\lambda_1, \dots, \lambda_n)$ contains real eigenvalues, and each $P_i = \mathbf{u}_i \mathbf{u}_i^H$ is an orthogonal rank-1 projection operator.
\end{theorem}

% ----------------------------------------------------------------------
% PART 3: VISUAL & GEOMETRIC INTUITION
% ----------------------------------------------------------------------
\section{Visual \& Geometric Intuition: Unitary Invariance and Spectral Decompositions}

Figure~\ref{fig:spectral_geometry} visualizes the Spectral Theorem: a Hermitian operator decomposes space into a collection of mutually orthogonal 1-dimensional eigenspaces, scaling each component independently by its real eigenvalue $\lambda_i$.

\begin{figure}[htbp]
\centering
\begin{tikzpicture}[
    scale=0.95,
    font=\sffamily\footnotesize
]

% Coordinate axes
\draw[->, >=stealth, color=gray] (-3.5, 0) -- (3.5, 0) node[right] {$\text{Re}(z_1)$};
\draw[->, >=stealth, color=gray] (0, -3.0) -- (0, 3.0) node[above] {$\text{Re}(z_2)$};

% Eigenspace lines
\draw[dashed, color=tealborder, line width=1pt] (-3.0, -1.5) -- (3.0, 1.5) node[above right] {$E_{\lambda_1} = \Span\{\mathbf{u}_1\}$};
\draw[dashed, color=coralborder, line width=1pt] (-1.5, 3.0) -- (1.5, -3.0) node[below right] {$E_{\lambda_2} = \Span\{\mathbf{u}_2\} \perp E_{\lambda_1}$};

% Eigenvectors
\draw[->, >=stealth, line width=1.8pt, color=tealborder] (0, 0) -- (2.0, 1.0) node[above left] {$\mathbf{u}_1$ ($\|\mathbf{u}_1\|=1$)};
\draw[->, >=stealth, line width=1.8pt, color=coralborder] (0, 0) -- (-1.0, 2.0) node[above right] {$\mathbf{u}_2$ ($\|\mathbf{u}_2\|=1$)};

% Arbitrary vector v
\coordinate (V) at (1.0, 3.0);
\draw[->, >=stealth, line width=1.8pt, color=primarydark] (0, 0) -- (V) node[above] {$\mathbf{v} = c_1 \mathbf{u}_1 + c_2 \mathbf{u}_2$};

% Projections
\coordinate (P1) at (1.6, 0.8);
\coordinate (P2) at (-0.6, 1.2);
\draw[dotted, thick, color=gray] (V) -- (P1);
\draw[dotted, thick, color=gray] (V) -- (P2);

% Transformed vector Av
\coordinate (AV) at (2.4, -0.6);
\draw[->, >=stealth, line width=2pt, color=red] (0, 0) -- (3.2, 1.6) node[below right] {$A\mathbf{v} = \lambda_1 c_1 \mathbf{u}_1 + \lambda_2 c_2 \mathbf{u}_2$};

\end{tikzpicture}
\caption{The Spectral Theorem Decomposition: Any vector $\mathbf{v}$ is resolved onto the orthogonal principal axes $\mathbf{u}_i$, where $A$ scales each projection independently by $\lambda_i \in \mathbb{R}$.}
\label{fig:spectral_geometry}
\end{figure}

% ----------------------------------------------------------------------
% PART 4: DEEP INTUITION & CONCEPTUAL DERIVATIONS
% ----------------------------------------------------------------------
\section{Deep Intuition \& Conceptual Derivations}

\subsection{The Rayleigh Quotient and Optimization on the Unit Sphere}
In machine learning (e.g., Principal Component Analysis and Fisher's Linear Discriminant), we seek directions $\mathbf{x}$ that maximize or minimize quadratic forms. The \textbf{Rayleigh Quotient} formalizes this optimization:

\begin{theorem}[Rayleigh-Ritz Theorem]
Let $A \in \mathbb{C}^{n \times n}$ be Hermitian with eigenvalues ordered as $\lambda_{\min} = \lambda_1 \le \lambda_2 \le \dots \le \lambda_n = \lambda_{\max}$. For all non-zero $\mathbf{x} \in \mathbb{C}^n$:
\begin{equation}
\lambda_{\min} \le R_A(\mathbf{x}) \triangleq \frac{\mathbf{x}^H A \mathbf{x}}{\mathbf{x}^H \mathbf{x}} \le \lambda_{\max}
\label{eq:rayleigh_bounds}
\end{equation}
Furthermore, the maximum is attained if and only if $\mathbf{x}$ is an eigenvector associated with $\lambda_{\max}$, and the minimum is attained if and only if $\mathbf{x}$ is an eigenvector associated with $\lambda_{\min}$.
\end{theorem}

\begin{proof}
By the Spectral Theorem, we can write $A = U \Lambda U^H$, where $U = [\mathbf{u}_1, \dots, \mathbf{u}_n]$ forms an orthonormal basis for $\mathbb{C}^n$.
Any vector $\mathbf{x} \in \mathbb{C}^n$ can be expanded in this basis: $\mathbf{x} = \sum_{i=1}^n c_i \mathbf{u}_i = U\mathbf{c}$, where $\mathbf{c} = U^H \mathbf{x}$.
Because $U$ is unitary, inner products are preserved:
\begin{equation}
\mathbf{x}^H \mathbf{x} = (U\mathbf{c})^H (U\mathbf{c}) = \mathbf{c}^H U^H U \mathbf{c} = \mathbf{c}^H \mathbf{c} = \sum_{i=1}^n |c_i|^2
\end{equation}
Similarly, the quadratic form evaluates to:
\begin{equation}
\mathbf{x}^H A \mathbf{x} = (U\mathbf{c})^H (U \Lambda U^H) (U\mathbf{c}) = \mathbf{c}^H \Lambda \mathbf{c} = \sum_{i=1}^n \lambda_i |c_i|^2
\end{equation}
The Rayleigh quotient becomes a weighted average of the eigenvalues:
\begin{equation}
R_A(\mathbf{x}) = \frac{\sum_{i=1}^n \lambda_i |c_i|^2}{\sum_{i=1}^n |c_i|^2}
\end{equation}
Let $w_i = \frac{|c_i|^2}{\sum_{j=1}^n |c_j|^2} \ge 0$, with $\sum_{i=1}^n w_i = 1$. Then:
\begin{equation}
R_A(\mathbf{x}) = \sum_{i=1}^n w_i \lambda_i
\end{equation}
Since $\lambda_1 \le \lambda_i \le \lambda_n$ for every $i$:
\begin{equation}
\lambda_1 = \sum_{i=1}^n w_i \lambda_1 \le \sum_{i=1}^n w_i \lambda_i \le \sum_{i=1}^n w_i \lambda_n = \lambda_n
\end{equation}
Choosing $\mathbf{x} = \mathbf{u}_n \implies w_n = 1, w_i = 0 \; (\forall i < n)$, giving $R_A(\mathbf{u}_n) = \lambda_n = \lambda_{\max}$.
Similarly, setting $\mathbf{x} = \mathbf{u}_1$ achieves the lower bound $\lambda_1 = \lambda_{\min}$.
\end{proof}

% ----------------------------------------------------------------------
% PART 5: EXPLICIT GOTCHAS & DIAGNOSTIC MISCONCEPTIONS
% ----------------------------------------------------------------------
\section{Explicit Gotchas \& Diagnostic Misconceptions}

\begin{trapbox}{Gotcha 1: Neglecting Conjugation in Complex Inner Products}
A student writes the complex inner product as $\langle \mathbf{u}, \mathbf{v} \rangle = \mathbf{u}^T \mathbf{v}$.
\textbf{The Trap:} If $\mathbf{z} = [1, i]^T$, then $\mathbf{z}^T \mathbf{z} = 1^2 + i^2 = 1 - 1 = 0$. A non-zero vector would have a ``norm'' of zero! In complex spaces, the inner product requires the conjugate transpose: $\langle \mathbf{u}, \mathbf{v} \rangle = \mathbf{u}^H \mathbf{v}$. Under this definition, $\|\mathbf{z}\|_2^2 = \mathbf{z}^H \mathbf{z} = |1|^2 + |i|^2 = 1 + 1 = 2 > 0$, preserving the positive-definiteness axiom.
\end{trapbox}

\begin{trapbox}{Gotcha 2: Assuming Real Eigenvalues Guarantee a Hermitian Matrix}
An analyst encounters matrix $M = \begin{bmatrix} 2 & 7 \\ 0 & 5 \end{bmatrix}$ with eigenvalues $\{2, 5\} \subset \mathbb{R}$. They conclude that $M$ is Hermitian.
\textbf{The Trap:} The implication only runs in one direction: Hermitian $\implies$ real eigenvalues. The converse is false. Triangular matrices with real diagonal entries always have real eigenvalues, but are not Hermitian unless they are diagonal. Hermitian matrices must satisfy $M = M^H$ and possess mutually orthogonal eigenvectors.
\end{trapbox}

\begin{trapbox}{Gotcha 3: Confusing Unitary Transformations with Hermitian Matrices}
A programmer assumes that because unitary matrices preserve norms, they must be Hermitian ($U = U^H$).
\textbf{The Trap:} Unitary matrices satisfy $U^H U = I$, meaning $U^H = U^{-1}$, NOT $U^H = U$. The eigenvalues of a unitary matrix lie on the unit circle in the complex plane ($|\lambda| = 1$), whereas the eigenvalues of a Hermitian matrix are real ($\lambda \in \mathbb{R}$). A matrix is simultaneously Hermitian and Unitary if and only if $U^2 = I$ (a reflection operator).
\end{trapbox}

\begin{trapbox}{Gotcha 4: Conflating Quadratic Forms Over $\mathbb{R}$ vs. $\mathbb{C}$}
A student takes a real symmetric matrix $A \in \mathbb{R}^{n \times n}$ and calculates $\mathbf{z}^T A \mathbf{z}$ for complex $\mathbf{z} \in \mathbb{C}^n$, expecting a real scalar.
\textbf{The Trap:} For complex vectors, $\mathbf{z}^T A \mathbf{z}$ is generally a complex number. To guarantee a real result, one must evaluate the sesquilinear quadratic form $\mathbf{z}^H A \mathbf{z}$. Because $(\mathbf{z}^H A \mathbf{z})^H = \mathbf{z}^H A^H \mathbf{z} = \mathbf{z}^H A \mathbf{z}$, the quantity is guaranteed real if and only if the conjugate transpose is used.
\end{trapbox}

% ----------------------------------------------------------------------
% PART 6: COMPARATIVE SYNTHESIS & SUMMARY TABLE
% ----------------------------------------------------------------------
\section{Comparative Synthesis \& Summary Table}

Table~\ref{tab:real_vs_complex} presents the exact mathematical dictionary mapping real linear algebraic concepts to their complex generalizations.

\begin{table}[htbp]
\centering
\small
\begin{tabularx}{\textwidth}{lXX}
\toprule
\textbf{Linear Algebra Concept} & \textbf{Real Formulation ($\mathbb{R}^n$)} & \textbf{Complex Formulation ($\mathbb{C}^n$)} \\
\midrule
Vector Transpose & $A^T$ ($A_{ji}$) & $A^H = \overline{A}^T$ (Conjugate Transpose) \\
\addlinespace
Inner Product & $\langle \mathbf{u}, \mathbf{v} \rangle = \mathbf{u}^T \mathbf{v}$ & $\langle \mathbf{u}, \mathbf{v} \rangle = \mathbf{u}^H \mathbf{v}$ (Sesquilinear) \\
\addlinespace
Vector Norm & $\|\mathbf{x}\|_2^2 = \mathbf{x}^T \mathbf{x} = \sum x_i^2$ & $\|\mathbf{z}\|_2^2 = \mathbf{z}^H \mathbf{z} = \sum |z_i|^2$ \\
\addlinespace
Self-Adjoint Matrix & Symmetric: $A^T = A$ & Hermitian: $A^H = A$ \\
\addlinespace
Isometric Transformation & Orthogonal: $Q^T Q = I$ & Unitary: $U^H U = I$ \\
\addlinespace
Eigenvalue Domain & $\lambda \in \mathbb{C}$ in general; $\lambda \in \mathbb{R}$ if symmetric & $\lambda \in \mathbb{R}$ strictly if Hermitian; $|\lambda|=1$ if Unitary \\
\addlinespace
Spectral Decomposition & $A = Q \Lambda Q^T$ ($Q \in O(n)$) & $A = U \Lambda U^H$ ($U \in U(n)$) \\
\bottomrule
\end{tabularx}
\caption{Systematic Dictionary: Real vs. Complex Linear Algebraic Structures.}
\label{tab:real_vs_complex}
\end{table}

% ----------------------------------------------------------------------
% PART 7: PYTHON / SCIENTIFIC COMPUTING LAB
% ----------------------------------------------------------------------
\section{Python / Scientific Computing Lab}

The following production script generates complex Hermitian matrices, computes their unitary spectral decomposition, and numerically validates the Rayleigh quotient bounds.

\begin{lstlisting}[language=Python, caption={Hermitian Matrix Decomposition and Rayleigh Quotient Verification in Python.}]
import numpy as np

def generate_random_hermitian(n: int, seed: int = 42) -> np.ndarray:
    """Generates an n x n complex Hermitian matrix: A = (M + M^H) / 2."""
    rng = np.random.RandomState(seed)
    real_part = rng.randn(n, n)
    imag_part = rng.randn(n, n)
    M = real_part + 1j * imag_part
    return (M + M.conj().T) / 2.0

def verify_spectral_theorem(A: np.ndarray):
    """Diagonalizes Hermitian matrix A and verifies unitary properties."""
    # np.linalg.eigh is optimized for Hermitian / symmetric matrices
    eigvals, U = np.linalg.eigh(A)

    # 1. Check realness of eigenvalues
    is_real = np.allclose(eigvals.imag, 0.0)

    # 2. Check unitary property: U^H @ U == I
    is_unitary = np.allclose(U.conj().T @ U, np.eye(A.shape[0]))

    # 3. Reconstruct A = U @ diag(eigvals) @ U^H
    A_reconstructed = U @ np.diag(eigvals) @ U.conj().T
    recon_error = np.linalg.norm(A - A_reconstructed)

    print(f"Eigenvalues are strictly real: {is_real}")
    print(f"Eigenvector matrix is Unitary: {is_unitary}")
    print(f"Spectral Reconstruction Error: {recon_error:.2e}")
    return eigvals, U

def test_rayleigh_quotient(A: np.ndarray, eigvals: np.ndarray, n_trials: int = 1000):
    """Numerically verifies that lambda_min <= R_A(x) <= lambda_max for random complex vectors."""
    lambda_min = eigvals.min()
    lambda_max = eigvals.max()

    rng = np.random.RandomState(101)
    n = A.shape[0]

    all_in_bounds = True
    for _ in range(n_trials):
        x = rng.randn(n) + 1j * rng.randn(n)
        R_x = (x.conj().T @ A @ x) / (x.conj().T @ x)
        val = R_x.real
        if val < lambda_min - 1e-9 or val > lambda_max + 1e-9:
            all_in_bounds = False
            break

    print(f"Rayleigh Bounds [{lambda_min:.4f}, {lambda_max:.4f}] respected in all trials: {all_in_bounds}")

def main():
    n = 4
    A = generate_random_hermitian(n)
    eigvals, U = verify_spectral_theorem(A)
    test_rayleigh_quotient(A, eigvals)

if __name__ == "__main__":
    main()
\end{lstlisting}

% ----------------------------------------------------------------------
% PART 8: COMPREHENSIVE PRACTICE PROBLEMS
% ----------------------------------------------------------------------
\section{Comprehensive Practice Problems}

\subsection*{Part A: Analytical \& Mathematical Problems}

\begin{enumerate}[label=\textbf{P\arabic*.}]
    \item \textbf{Verifying Hermitian Properties:}
    Consider the complex matrix $H = \begin{bmatrix} 2 & 1 - i \\ 1 + i & 3 \end{bmatrix}$.
    \begin{enumerate}[label=(\alph*)]
        \item Compute $H^H$ and confirm that $H$ is Hermitian.
        \item Compute the characteristic polynomial $\det(H - \lambda I) = 0$ and find its eigenvalues.
        \item Find the normalized eigenvectors $\mathbf{u}_1, \mathbf{u}_2 \in \mathbb{C}^2$ and verify $\mathbf{u}_1^H \mathbf{u}_2 = 0$.
        \item Verify the projection formula $H = \lambda_1 \mathbf{u}_1 \mathbf{u}_1^H + \lambda_2 \mathbf{u}_2 \mathbf{u}_2^H$.
    \end{enumerate}

    \item \textbf{Spectral Properties of Unitary Transformations:}
    Let $U \in \mathbb{C}^{n \times n}$ be a unitary matrix ($U^H U = I$).
    \begin{enumerate}[label=(\alph*)]
        \item Prove that for any $\mathbf{x} \in \mathbb{C}^n$, $\|U\mathbf{x}\|_2 = \|\mathbf{x}\|_2$.
        \item Prove that every eigenvalue $\lambda$ of $U$ satisfies $|\lambda| = 1$ (i.e., $\lambda = e^{i\theta}$ for some $\theta \in \mathbb{R}$).
        \item Prove that if $\lambda_1 \ne \lambda_2$, their corresponding eigenvectors are orthogonal: $\mathbf{v}_1^H \mathbf{v}_2 = 0$.
    \end{enumerate}

    \item \textbf{Skew-Hermitian Matrices:}
    A matrix $S \in \mathbb{C}^{n \times n}$ is skew-Hermitian if $S^H = -S$.
    \begin{enumerate}[label=(\alph*)]
        \item Prove that the eigenvalues of $S$ are purely imaginary (of the form $\lambda = i\mu$, where $\mu \in \mathbb{R}$).
        \item Prove that matrix exponential $U = \exp(S)$ is a unitary matrix.
    \end{enumerate}

    \item \textbf{Decomposition into Hermitian and Skew-Hermitian Components:}
    Let $M \in \mathbb{C}^{n \times n}$ be an arbitrary square matrix.
    \begin{enumerate}[label=(\alph*)]
        \item Show that $M$ can be uniquely written as $M = H + S$, where $H$ is Hermitian and $S$ is skew-Hermitian.
        \item Express $H$ and $S$ explicitly in terms of $M$ and $M^H$.
        \item Apply this decomposition to $M = \begin{bmatrix} 1 + i & 4 \\ 2i & 3 \end{bmatrix}$.
    \end{enumerate}

    \item \textbf{Courant-Fischer Min-Max Theorem:}
    Let $A$ be an $n \times n$ Hermitian matrix with eigenvalues $\lambda_1 \le \lambda_2 \le \dots \le \lambda_n$.
    The Courant-Fischer theorem states that the $k$-th eigenvalue satisfies:
    \begin{equation}
    \lambda_k = \min_{S: \dim(S) = k} \max_{\mathbf{x} \in S, \mathbf{x} \ne \mathbf{0}} \frac{\mathbf{x}^H A \mathbf{x}}{\mathbf{x}^H \mathbf{x}}
    \end{equation}
    Using this variational characterization, prove that if $A$ and $B$ are Hermitian and $B \succeq 0$ (positive semidefinite), then $\lambda_k(A + B) \ge \lambda_k(A)$ for all $k \in \{1, \dots, n\}$.
\end{enumerate}

\subsection*{Part B: Diagnostic Multiple-Choice Questions}

\begin{enumerate}[label=\textbf{MCQ\arabic*.}]
    \item Let $A \in \mathbb{C}^{n \times n}$ be a Hermitian matrix. Which of the following conditions is NOT guaranteed to hold?
    \begin{enumerate}[label=(\alph*)]
        \item All eigenvalues of $A$ are real numbers.
        \item $A$ possesses an orthonormal basis of eigenvectors in $\mathbb{C}^n$.
        \item The determinant $\det(A)$ is a real number.
        \item All entries on the main diagonal of $A$ are strictly positive.
    \end{enumerate}

    \item What is the geometric action of a unitary operator $U \in \mathbb{C}^{n \times n}$ on the complex vector space $\mathbb{C}^n$?
    \begin{enumerate}[label=(\alph*)]
        \item It projects vectors orthogonally onto a lower-dimensional subspace.
        \item It scales all vectors by an arbitrary real constant $c > 0$.
        \item It is a rigid isometry that preserves Euclidean norms and inner products between all vectors.
        \item It maps all eigenvalues to zero.
    \end{enumerate}
\end{enumerate}


\chapter[SVD \& Positive Definite Matrices]{Singular Value Decomposition (SVD), Low-Rank Approximations \& Positive Definite Matrices}
\label{chap:week06}

% ----------------------------------------------------------------------
% PART 1: EXECUTIVE ROADMAP & LEARNING OBJECTIVES
% ----------------------------------------------------------------------
\section{Executive Roadmap \& Learning Objectives}

While eigendecomposition is strictly confined to square matrices and can fail entirely if a matrix is defective, the \textbf{Singular Value Decomposition (SVD)} stands as the undisputed crown jewel of linear algebra: it exists for \emph{every} matrix $A \in \mathbb{R}^{m \times n}$, whether square or rectangular, full-rank or rank-deficient, well-conditioned or completely singular. 

SVD unlocks the deepest geometric truth of linear transformations: \textbf{every linear map transforms a unit sphere into a hyper-ellipsoid}. SVD dissects this geometric transformation into three elementary, decoupled physical actions: a rigid rotation in the domain, a coordinate scaling along perpendicular axes by non-negative \textbf{singular values}, and a final rigid rotation in the codomain.

In data science and modern machine learning, SVD is the foundation of dimensionality reduction (Principal Component Analysis), matrix compression, latent semantic analysis, recommender systems (collaborative filtering), and regularized pseudo-inversion. This chapter develops the complete mathematical theory of SVD, the Eckart-Young-Mirsky low-rank approximation theorem, and the characterization of positive definite matrices.

\subsection{Prerequisites \& Scaffolding}
Mastery of this chapter builds directly upon orthogonal matrices, fundamental subspaces, and the Spectral Theorem as mapped in Table~\ref{tab:w06_prereqs}.

\begin{table}[htbp]
\centering
\small
\begin{tabularx}{\textwidth}{lXX}
\toprule
\textbf{Foundational Area} & \textbf{Core Concepts Required} & \textbf{Role in Machine Learning Foundations} \\
\midrule
Spectral Theorem & Orthogonal diagonalization $S = Q \Lambda Q^T$ of symmetric matrices & Foundation for decomposing Gram matrices $A^T A$ and $A A^T$ \\
Fundamental Subspaces & $\Col(A), \Row(A), \Null(A), \Null(A^T)$, dimensions and orthogonality & Explicitly mapped to the singular vectors $U$ and $V$ \\
Matrix Norms & Frobenius norm $\|A\|_F$, spectral operator norm $\|A\|_2$ & Quantifying error in low-rank data compression \\
\bottomrule
\end{tabularx}
\caption{Prerequisite Scaffolding for Chapter 6.}
\label{tab:w06_prereqs}
\end{table}

\subsection{Core Competencies}
By the conclusion of this chapter, students will be able to:
\begin{itemize}
    \item Derive and compute both the full SVD ($A = U \Sigma V^T$) and compact SVD ($A = U_r \Sigma_r V_r^T$) for arbitrary rectangular matrices.
    \item Relate singular values and singular vectors directly to the eigendecompositions of the Gram matrices $A^T A$ and $A A^T$.
    \item State and apply the four equivalent characterizations of Positive Definite and Semidefinite matrices (eigenvalues, quadratic forms, Sylvester's minors, and Cholesky decomposition).
    \item State and apply the Eckart-Young-Mirsky Theorem to compute the provably optimal rank-$k$ matrix approximation under both Spectral and Frobenius norms.
    \item Construct the Moore-Penrose pseudoinverse $A^+$ via SVD to find the unique minimum-norm solution to arbitrary linear systems.
\end{itemize}

\subsection{Notation \& Mathematical Conventions}
\begin{itemize}
    \item $A \in \mathbb{R}^{m \times n}$: An arbitrary real matrix of rank $r \le \min(m, n)$.
    \item $\sigma_1 \ge \sigma_2 \ge \dots \ge \sigma_r > 0$: The strictly positive singular values of $A$, ordered decreasingly.
    \item $U = [\mathbf{u}_1, \dots, \mathbf{u}_m] \in \mathbb{R}^{m \times m}$: Orthogonal matrix of left singular vectors (eigenvectors of $A A^T$).
    \item $V = [\mathbf{v}_1, \dots, \mathbf{v}_n] \in \mathbb{R}^{n \times n}$: Orthogonal matrix of right singular vectors (eigenvectors of $A^T A$).
    \item $\Sigma \in \mathbb{R}^{m \times n}$: Rectangular pseudo-diagonal matrix with singular values $\Sigma_{ii} = \sigma_i$.
    \item $A_k \triangleq \sum_{i=1}^k \sigma_i \mathbf{u}_i \mathbf{v}_i^T$: The optimal truncated rank-$k$ approximation of $A$.
    \item $A^+ \triangleq V \Sigma^+ U^T = \sum_{i=1}^r \frac{1}{\sigma_i} \mathbf{v}_i \mathbf{u}_i^T$: The Moore-Penrose pseudoinverse.
    \item $M \succ 0$ ($M \succeq 0$): Denotes a symmetric positive definite (positive semidefinite) matrix.
\end{itemize}

% ----------------------------------------------------------------------
% PART 2: THE GRAND PARADOX & MOTIVATION
% ----------------------------------------------------------------------
\section{The Grand Paradox: The Sphere-to-Ellipsoid Transformation}

\subsection{The Universal Geometric Action}
Consider an arbitrary $2 \times 2$ matrix with messy, non-symmetric entries:
\begin{equation}
A = \begin{bmatrix} 3 & 0 \\ 4 & 5 \end{bmatrix}
\end{equation}
If you feed a circle of unit vectors $\mathbf{x} \in \mathbb{R}^2$ ($\|\mathbf{x}\|_2 = 1$) into this matrix, what does the output set of vectors $\{A\mathbf{x} \mid \|\mathbf{x}\|_2 = 1\}$ look like? Does it warp into a peanut shape? Does it fold into a polygon?

\begin{quote}
\textit{\textbf{The Geometric Discovery:}} \textbf{Every} linear transformation—without exception—maps a unit sphere into a hyper-ellipsoid. 
\end{quote}

Why? Because if $\mathbf{y} = A\mathbf{x}$ and $A$ is invertible, then $\mathbf{x} = A^{-1} \mathbf{y}$. The unit sphere condition $\|\mathbf{x}\|_2^2 = 1$ becomes:
\begin{equation}
\mathbf{x}^T \mathbf{x} = 1 \implies (A^{-1}\mathbf{y})^T (A^{-1}\mathbf{y}) = 1 \implies \mathbf{y}^T (A A^T)^{-1} \mathbf{y} = 1
\end{equation}
The matrix $(A A^T)^{-1}$ is strictly symmetric positive definite! By the Spectral Theorem, any quadratic equation $\mathbf{y}^T M \mathbf{y} = 1$ with $M \succ 0$ defines an \textbf{ellipsoid}. 

The principal axes of this output ellipsoid are the left singular vectors $\mathbf{u}_i$, the half-lengths of the axes are the singular values $\sigma_i$, and the original orthogonal directions on the unit sphere that get stretched into these axes are the right singular vectors $\mathbf{v}_i$!

\begin{ponderbox}[Singular Values vs. Eigenvalues]
Consider the matrix $A = \begin{bmatrix} 0 & 3 \\ -3 & 0 \end{bmatrix}$.
\begin{enumerate}
    \item Compute the eigenvalues of $A$. Are they real or complex?
    \item Compute the Gram matrix $A^T A$ and find its eigenvalues.
    \item What are the singular values of $A$? Can a singular value ever be negative or complex?
    \item Why does SVD succeed in extracting real scale factors when eigendecomposition yields purely imaginary numbers?
\end{enumerate}
\end{ponderbox}

% ----------------------------------------------------------------------
% PART 3: MATHEMATICAL THEORY OF SVD
% ----------------------------------------------------------------------
\section{Mathematical Theory of SVD}

\begin{theorem}[Singular Value Decomposition (SVD)]
Let $A \in \mathbb{R}^{m \times n}$ have rank $r \le \min(m, n)$. There exist orthogonal matrices $U \in \mathbb{R}^{m \times m}$ ($U^T U = I_m$) and $V \in \mathbb{R}^{n \times n}$ ($V^T V = I_n$), and a unique diagonal matrix $\Sigma \in \mathbb{R}^{m \times n}$ such that:
\begin{equation}
A = U \Sigma V^T = \sum_{i=1}^r \sigma_i \mathbf{u}_i \mathbf{v}_i^T
\label{eq:svd_outer}
\end{equation}
where the singular values are strictly positive and ordered:
\begin{equation}
\sigma_1 \ge \sigma_2 \ge \dots \ge \sigma_r > 0
\end{equation}
\end{theorem}

\begin{proof}
Consider the Gram matrix $A^T A \in \mathbb{R}^{n \times n}$.
\begin{enumerate}
    \item \textbf{Symmetry and Semidefiniteness:}
    $(A^T A)^T = A^T (A^T)^T = A^T A$, so $A^T A$ is symmetric.
    For any $\mathbf{x} \in \mathbb{R}^n$, $\mathbf{x}^T (A^T A)\mathbf{x} = (A\mathbf{x})^T (A\mathbf{x}) = \|A\mathbf{x}\|_2^2 \ge 0$. Thus $A^T A$ is positive semidefinite.
    \item \textbf{Orthonormal Basis $V$:}
    By the Spectral Theorem for real symmetric matrices, $A^T A$ possesses an orthonormal basis of eigenvectors $\mathbf{v}_1, \dots, \mathbf{v}_n \in \mathbb{R}^n$ with non-negative real eigenvalues:
    \begin{equation}
    A^T A \mathbf{v}_i = \lambda_i \mathbf{v}_i, \quad \lambda_1 \ge \lambda_2 \ge \dots \ge \lambda_n \ge 0
    \end{equation}
    Because $\rank(A^T A) = \rank(A) = r$, exactly $r$ eigenvalues are strictly positive: $\lambda_i > 0$ for $i \le r$, and $\lambda_i = 0$ for $i > r$.
    \item \textbf{Definition of Singular Values:}
    Define $\sigma_i \triangleq \sqrt{\lambda_i} > 0$ for $i \in \{1, \dots, r\}$.
    \item \textbf{Construction of Orthonormal Vectors $U$:}
    For each $i \in \{1, \dots, r\}$, define:
    \begin{equation}
    \mathbf{u}_i \triangleq \frac{1}{\sigma_i} A \mathbf{v}_i \in \mathbb{R}^m
    \label{eq:def_ui}
    \end{equation}
    Testing the inner product between $\mathbf{u}_i$ and $\mathbf{u}_j$:
    \begin{align}
    \mathbf{u}_i^T \mathbf{u}_j &= \left( \frac{1}{\sigma_i} A \mathbf{v}_i \right)^T \left( \frac{1}{\sigma_j} A \mathbf{v}_j \right) = \frac{1}{\sigma_i \sigma_j} \mathbf{v}_i^T (A^T A \mathbf{v}_j) \\
    &= \frac{1}{\sigma_i \sigma_j} \mathbf{v}_i^T (\lambda_j \mathbf{v}_j) = \frac{\lambda_j}{\sigma_i \sigma_j} \mathbf{v}_i^T \mathbf{v}_j = \frac{\sigma_j^2}{\sigma_i \sigma_j} \delta_{ij} = \delta_{ij}
    \end{align}
    Thus $\{\mathbf{u}_1, \dots, \mathbf{u}_r\}$ forms an orthonormal set in $\mathbb{R}^m$. Extend this set via Gram-Schmidt to an orthonormal basis $\{\mathbf{u}_1, \dots, \mathbf{u}_m\}$ for all of $\mathbb{R}^m$.
    \item \textbf{Matrix Assembly:}
    From Equation~\eqref{eq:def_ui}, $A\mathbf{v}_i = \sigma_i \mathbf{u}_i$ for $i \le r$. For $i > r$, $\lambda_i = 0 \implies \|A\mathbf{v}_i\|_2^2 = \mathbf{v}_i^T A^T A \mathbf{v}_i = 0 \implies A\mathbf{v}_i = \mathbf{0}$.
    Stacking into matrices:
    \begin{equation}
    A V = U \Sigma \implies A = U \Sigma V^T
    \end{equation}
\end{enumerate}
\end{proof}

\begin{figure}[htbp]
\centering
\begin{tikzpicture}[
    scale=0.9,
    font=\sffamily\footnotesize
]

% 1. Input Space: Unit Circle
\draw[thick, fill=primarylight!30] (-5.0, 0) circle (1.5cm);
\draw[->, >=stealth, color=gray] (-6.8, 0) -- (-3.2, 0) node[right] {$x_1$};
\draw[->, >=stealth, color=gray] (-5.0, -1.8) -- (-5.0, 1.8) node[above] {$x_2$};
\draw[->, >=stealth, line width=1.5pt, color=primarydark] (-5.0, 0) -- (-3.8, 0.9) node[above right] {$\mathbf{v}_1$};
\draw[->, >=stealth, line width=1.5pt, color=primarydark] (-5.0, 0) -- (-5.9, 1.2) node[above left] {$\mathbf{v}_2$};
\node[below] at (-5.0, -2.0) {\textbf{1. Domain: } Basis $V$};

% First arrow
\draw[->, >=stealth, line width=1.5pt, color=primarydark] (-3.0, 0) -- (-1.5, 0) node[midway, above] {$V^T$};

% 2. Intermediate Space: Scaled Ellipse aligned with axes
\draw[thick, fill=tealbg!30] (0.5, 0) ellipse (2.0cm and 0.9cm);
\draw[->, >=stealth, color=gray] (-1.8, 0) -- (2.8, 0);
\draw[->, >=stealth, color=gray] (0.5, -1.5) -- (0.5, 1.5);
\draw[->, >=stealth, line width=1.5pt, color=tealborder] (0.5, 0) -- (2.5, 0) node[below] {$\sigma_1 \mathbf{e}_1$};
\draw[->, >=stealth, line width=1.5pt, color=tealborder] (0.5, 0) -- (0.5, 0.9) node[left] {$\sigma_2 \mathbf{e}_2$};
\node[below] at (0.5, -2.0) {\textbf{2. Stretched: } $\Sigma$};

% Second arrow
\draw[->, >=stealth, line width=1.5pt, color=primarydark] (3.0, 0) -- (4.5, 0) node[midway, above] {$U$};

% 3. Output Space: Rotated Ellipse
\begin{scope}[shift={(6.5, 0)}, rotate=40]
    \draw[thick, fill=coralbg!40] (0, 0) ellipse (2.0cm and 0.9cm);
    \draw[->, >=stealth, line width=1.5pt, color=coralborder] (0, 0) -- (2.0, 0) node[above right] {$\sigma_1 \mathbf{u}_1$};
    \draw[->, >=stealth, line width=1.5pt, color=coralborder] (0, 0) -- (0, 0.9) node[above left] {$\sigma_2 \mathbf{u}_2$};
\end{scope}
\node[below] at (6.5, -2.0) {\textbf{3. Codomain: } Basis $U$};

\end{tikzpicture}
\caption{The Geometric Action of SVD: Any matrix transformation factorizes into a domain rotation $V^T$, axis stretching $\Sigma$, and codomain rotation $U$.}
\label{fig:svd_geometry}
\end{figure}

% ----------------------------------------------------------------------
% PART 4: POSITIVE DEFINITE MATRICES
% ----------------------------------------------------------------------
\section{Positive Definite \& Semidefinite Matrices}

\begin{definitionbox}{Positive Definite and Semidefinite Matrices}
A real symmetric matrix $M \in \mathbb{R}^{n \times n}$ ($M = M^T$) is:
\begin{itemize}
    \item \textbf{Positive Definite} ($M \succ 0$) if for all $\mathbf{x} \in \mathbb{R}^n \setminus \{\mathbf{0}\}$: $\mathbf{x}^T M \mathbf{x} > 0$.
    \item \textbf{Positive Semidefinite} ($M \succeq 0$) if for all $\mathbf{x} \in \mathbb{R}^n$: $\mathbf{x}^T M \mathbf{x} \ge 0$.
\end{itemize}
\end{definitionbox}

\begin{theorem}[Equivalent Tests for Positive Definiteness]
For any real symmetric matrix $M \in \mathbb{R}^{n \times n}$, the following four conditions are mathematically equivalent:
\begin{enumerate}
    \item \textbf{Eigenvalue Test:} All eigenvalues are strictly positive: $\lambda_i > 0$ for all $i \in \{1, \dots, n\}$.
    \item \textbf{Energy Test:} The quadratic form $\mathbf{x}^T M \mathbf{x} > 0$ for all non-zero $\mathbf{x} \in \mathbb{R}^n$.
    \item \textbf{Sylvester's Leading Minors Test:} All $n$ leading principal minors are strictly positive:
    \begin{equation}
    \det(M_{1:k, 1:k}) > 0 \quad \text{for } k = 1, 2, \dots, n
    \end{equation}
    \item \textbf{Cholesky Decomposition Test:} There exists a unique lower-triangular matrix $L$ with strictly positive diagonal entries ($L_{ii} > 0$) such that:
    \begin{equation}
    M = L L^T
    \label{eq:cholesky}
    \end{equation}
\end{enumerate}
\end{theorem}

\begin{proof}[Proof Sketch]
$(1 \iff 2)$: By the Spectral Theorem, $M = Q \Lambda Q^T$. Let $\mathbf{y} = Q^T \mathbf{x}$. Then $\mathbf{x}^T M \mathbf{x} = \mathbf{y}^T \Lambda \mathbf{y} = \sum_{i=1}^n \lambda_i y_i^2$. Because $Q$ is invertible, $\mathbf{x} \ne \mathbf{0} \iff \mathbf{y} \ne \mathbf{0}$. The sum $\sum \lambda_i y_i^2$ is strictly positive for all $\mathbf{y} \ne \mathbf{0}$ if and only if every $\lambda_i > 0$.

$(1 \iff 4)$: If $M = L L^T$, then for $\mathbf{x} \ne \mathbf{0}$: $\mathbf{x}^T M \mathbf{x} = \mathbf{x}^T L L^T \mathbf{x} = \|L^T \mathbf{x}\|_2^2$. Since $L$ has positive diagonals, $\det(L) = \prod L_{ii} > 0$, so $L$ is invertible. Hence $L^T \mathbf{x} \ne \mathbf{0}$ whenever $\mathbf{x} \ne \mathbf{0}$, ensuring $\|L^T \mathbf{x}\|_2^2 > 0$.
\end{proof}

% ----------------------------------------------------------------------
% PART 5: LOW-RANK APPROXIMATION & ECKART-YOUNG
% ----------------------------------------------------------------------
\section{Low-Rank Matrix Approximation \& Eckart-Young Theorem}

In massive machine learning datasets (user-item rating matrices in Netflix, document-word matrices in NLP, or high-resolution images), storing a full matrix $A \in \mathbb{R}^{m \times n}$ requires $\mathcal{O}(mn)$ numbers. Can we compress $A$ into a matrix $B$ of low rank $k \ll \min(m, n)$ while discarding as little information as possible?

\begin{theorem}[Eckart-Young-Mirsky Theorem]
Let $A = \sum_{i=1}^r \sigma_i \mathbf{u}_i \mathbf{v}_i^T$ be the SVD of $A \in \mathbb{R}^{m \times n}$ with $\sigma_1 \ge \sigma_2 \ge \dots \ge \sigma_r > 0$. For any target rank $k < r$, the truncated SVD:
\begin{equation}
A_k \triangleq \sum_{i=1}^k \sigma_i \mathbf{u}_i \mathbf{v}_i^T
\label{eq:truncated_svd_formula}
\end{equation}
is the \textbf{globally optimal} rank-$k$ approximation of $A$ under both the Spectral (operator) norm and Frobenius norm:
\begin{align}
\min_{\rank(B) \le k} \|A - B\|_2 &= \|A - A_k\|_2 = \sigma_{k+1} \label{eq:ey_spectral_error} \\
\min_{\rank(B) \le k} \|A - B\|_F &= \|A - A_k\|_F = \sqrt{\sum_{i=k+1}^r \sigma_i^2} \label{eq:ey_frobenius_error}
\end{align}
\end{theorem}

\begin{ponderbox}[The Low-Rank Compression Miracle]
Consider a $1000 \times 1000$ image ($1,000,000$ pixels).
\begin{enumerate}
    \item How many numbers are required to store the full image $A$?
    \item If you approximate the image with a rank-$k = 20$ truncated SVD $A_{20} = U_{20} \Sigma_{20} V_{20}^T$, how many numbers do you need to store (columns of $U$, singular values, columns of $V$)?
    \item What is the compression ratio?
    \item Why is this compression provably superior to any other rank-20 matrix that could ever be created?
\end{enumerate}
\end{ponderbox}

% ----------------------------------------------------------------------
% PART 6: IN-DEPTH WORKED EXAMPLES
% ----------------------------------------------------------------------
\section{In-Depth Worked Numerical Examples}

\begin{examplebox}{Manual SVD of a Rectangular $3 \times 2$ Matrix}
Let matrix $A \in \mathbb{R}^{3 \times 2}$ be defined as:
\begin{equation}
A = \begin{bmatrix} 1 & 1 \\ 0 & 1 \\ -1 & 1 \end{bmatrix}
\end{equation}
Compute the full SVD $A = U \Sigma V^T$ step-by-step by hand.

\textbf{Step 1: Compute Gram Matrix $A^T A$ and Eigendecomposition.}
\begin{equation}
A^T A = \begin{bmatrix} 1 & 0 & -1 \\ 1 & 1 & 1 \end{bmatrix} \begin{bmatrix} 1 & 1 \\ 0 & 1 \\ -1 & 1 \end{bmatrix} = \begin{bmatrix} 1+0+1 & 1+0-1 \\ 1+0-1 & 1+1+1 \end{bmatrix} = \begin{bmatrix} 2 & 0 \\ 0 & 3 \end{bmatrix}
\end{equation}
Notice $A^T A$ is already diagonal!
The eigenvalues are $\lambda_1 = 3$ and $\lambda_2 = 2$ (ordered decreasingly).
The singular values are:
\begin{equation}
\sigma_1 = \sqrt{3}, \quad \sigma_2 = \sqrt{2}
\end{equation}

\textbf{Step 2: Determine Orthonormal Right Singular Vectors $V$.}
Because $A^T A = \begin{bmatrix} 2 & 0 \\ 0 & 3 \end{bmatrix}$:
\begin{itemize}
    \item For $\lambda_1 = 3$: $(A^T A - 3I)\mathbf{v}_1 = \begin{bmatrix} -1 & 0 \\ 0 & 0 \end{bmatrix}\mathbf{v}_1 = \mathbf{0} \implies \mathbf{v}_1 = \begin{bmatrix} 0 \\ 1 \end{bmatrix}$.
    \item For $\lambda_2 = 2$: $(A^T A - 2I)\mathbf{v}_2 = \begin{bmatrix} 0 & 0 \\ 0 & 1 \end{bmatrix}\mathbf{v}_2 = \mathbf{0} \implies \mathbf{v}_2 = \begin{bmatrix} 1 \\ 0 \end{bmatrix}$.
\end{itemize}
Thus:
\begin{equation}
V = \begin{bmatrix} \mathbf{v}_1 & \mathbf{v}_2 \end{bmatrix} = \begin{bmatrix} 0 & 1 \\ 1 & 0 \end{bmatrix}
\end{equation}

\textbf{Step 3: Compute Left Singular Vectors $\mathbf{u}_1, \mathbf{u}_2 \in \mathbb{R}^3$.}
Using $\mathbf{u}_i = \frac{1}{\sigma_i} A \mathbf{v}_i$:
\begin{align}
\mathbf{u}_1 &= \frac{1}{\sqrt{3}} \begin{bmatrix} 1 & 1 \\ 0 & 1 \\ -1 & 1 \end{bmatrix} \begin{bmatrix} 0 \\ 1 \end{bmatrix} = \frac{1}{\sqrt{3}} \begin{bmatrix} 1 \\ 1 \\ 1 \end{bmatrix} \\
\mathbf{u}_2 &= \frac{1}{\sqrt{2}} \begin{bmatrix} 1 & 1 \\ 0 & 1 \\ -1 & 1 \end{bmatrix} \begin{bmatrix} 1 \\ 0 \end{bmatrix} = \frac{1}{\sqrt{2}} \begin{bmatrix} 1 \\ 0 \\ -1 \end{bmatrix}
\end{align}
Testing orthogonality: $\mathbf{u}_1^T \mathbf{u}_2 = \frac{1}{\sqrt{6}}(1(1) + 1(0) + 1(-1)) = 0 \quad \checkmark$.

\textbf{Step 4: Find the Third Basis Vector $\mathbf{u}_3 \in \Null(A^T)$.}
To construct full $U \in \mathbb{R}^{3 \times 3}$, we need $\mathbf{u}_3$ orthogonal to both $\mathbf{u}_1$ and $\mathbf{u}_2$. Taking cross product or solving $A^T \mathbf{u}_3 = \mathbf{0}$:
\begin{equation}
\begin{bmatrix} 1 & 0 & -1 \\ 1 & 1 & 1 \end{bmatrix} \begin{bmatrix} z_1 \\ z_2 \\ z_3 \end{bmatrix} = \begin{bmatrix} 0 \\ 0 \end{bmatrix} \implies z_1 = z_3, \quad z_2 = -2z_3
\end{equation}
Choosing $z_3 = 1 \implies [1, -2, 1]^T$. Normalizing:
\begin{equation}
\mathbf{u}_3 = \frac{1}{\sqrt{6}} \begin{bmatrix} 1 \\ -2 \\ 1 \end{bmatrix}
\end{equation}

\textbf{Step 5: Assemble Full SVD and Reconstruct.}
\begin{equation}
U = \begin{bmatrix} 1/\sqrt{3} & 1/\sqrt{2} & 1/\sqrt{6} \\ 1/\sqrt{3} & 0 & -2/\sqrt{6} \\ 1/\sqrt{3} & -1/\sqrt{2} & 1/\sqrt{6} \end{bmatrix}, \quad \Sigma = \begin{bmatrix} \sqrt{3} & 0 \\ 0 & \sqrt{2} \\ 0 & 0 \end{bmatrix}, \quad V^T = \begin{bmatrix} 0 & 1 \\ 1 & 0 \end{bmatrix}
\end{equation}
Reconstructing $U \Sigma V^T$:
\begin{align}
U \Sigma &= \begin{bmatrix} 1 & 1 \\ 1 & 0 \\ 1 & -1 \end{bmatrix} \\
U \Sigma V^T &= \begin{bmatrix} 1 & 1 \\ 1 & 0 \\ 1 & -1 \end{bmatrix} \begin{bmatrix} 0 & 1 \\ 1 & 0 \end{bmatrix} = \begin{bmatrix} 1 & 1 \\ 0 & 1 \\ -1 & 1 \end{bmatrix} = A \quad \checkmark
\end{align}
\end{examplebox}

\begin{examplebox}{Positive Definiteness Diagnostic: Sylvester vs. Cholesky}
Let matrix $M = \begin{bmatrix} 2 & -1 & 0 \\ -1 & 2 & -1 \\ 0 & -1 & 2 \end{bmatrix}$.
\begin{enumerate}
    \item Compute all leading principal minors (Sylvester's criterion).
    \item Construct the lower-triangular Cholesky factor $L$ ($M = L L^T$).
    \item Expand the quadratic form $\mathbf{x}^T M \mathbf{x}$ as a sum of squares.
\end{enumerate}

\textbf{Step 1: Compute Leading Principal Minors.}
\begin{align}
\Delta_1 &= \det[2] = 2 > 0 \\
\Delta_2 &= \det\begin{bmatrix} 2 & -1 \\ -1 & 2 \end{bmatrix} = 4 - 1 = 3 > 0 \\
\Delta_3 &= \det(M) = 2(4 - 1) - (-1)(-2 - 0) + 0 = 6 - 2 = 4 > 0
\end{align}
Since $\Delta_1 = 2, \Delta_2 = 3, \Delta_3 = 4 > 0$, Sylvester's criterion guarantees $M \succ 0$.

\textbf{Step 2: Cholesky Factorization $M = L L^T$.}
Let $L = \begin{bmatrix} l_{11} & 0 & 0 \\ l_{21} & l_{22} & 0 \\ l_{31} & l_{32} & l_{33} \end{bmatrix}$. Matching entries of $L L^T = M$:
\begin{align}
l_{11}^2 = 2 &\implies l_{11} = \sqrt{2} \\
l_{21} l_{11} = -1 &\implies l_{21} = -\frac{1}{\sqrt{2}} \\
l_{21}^2 + l_{22}^2 = 2 &\implies \frac{1}{2} + l_{22}^2 = 2 \implies l_{22} = \sqrt{\frac{3}{2}} \\
l_{31} l_{11} = 0 &\implies l_{31} = 0 \\
l_{31} l_{21} + l_{32} l_{22} = -1 &\implies 0 + l_{32} \sqrt{\frac{3}{2}} = -1 \implies l_{32} = -\sqrt{\frac{2}{3}} \\
l_{31}^2 + l_{32}^2 + l_{33}^2 = 2 &\implies 0 + \frac{2}{3} + l_{33}^2 = 2 \implies l_{33} = \sqrt{\frac{4}{3}} = \frac{2}{\sqrt{3}}
\end{align}
The Cholesky factor is:
\begin{equation}
L = \begin{bmatrix} \sqrt{2} & 0 & 0 \\ -1/\sqrt{2} & \sqrt{3/2} & 0 \\ 0 & -\sqrt{2/3} & 2/\sqrt{3} \end{bmatrix}
\end{equation}

\textbf{Step 3: Quadratic Form as Sum of Squares.}
\begin{align}
\mathbf{x}^T M \mathbf{x} &= 2x_1^2 + 2x_2^2 + 2x_3^2 - 2x_1 x_2 - 2x_2 x_3 \\
&= x_1^2 + (x_1^2 - 2x_1 x_2 + x_2^2) + (x_2^2 - 2x_2 x_3 + x_3^2) + x_3^2 \\
&= x_1^2 + (x_1 - x_2)^2 + (x_2 - x_3)^2 + x_3^2
\end{align}
Because this is a sum of four squares, it is strictly non-negative. It equals zero if and only if $x_1 = 0$, $x_1 = x_2$, $x_2 = x_3$, and $x_3 = 0 \implies \mathbf{x} = \mathbf{0}$. 

Hence $\mathbf{x}^T M \mathbf{x} > 0$ for all $\mathbf{x} \ne \mathbf{0}$, providing independent visual proof that $M \succ 0$!
\end{examplebox}

% ----------------------------------------------------------------------
% PART 7: EXAM TRAPS & DIAGNOSTIC MISCONCEPTIONS
% ----------------------------------------------------------------------
\section{Exam Traps \& Diagnostic Misconceptions}

\begin{trapbox}{Gotcha 1: Conflating Singular Values with Eigenvalues}
A student claims that because a matrix has eigenvalues $\lambda_1 = 5, \lambda_2 = -3$, its singular values must be $5$ and $-3$.
\textbf{The Trap:} Singular values are the square roots of the eigenvalues of $A^T A$: $\sigma_i = \sqrt{\lambda_i(A^T A)}$. Because $A^T A \succeq 0$, singular values are strictly \textbf{non-negative real numbers} ($\sigma_i \ge 0$). They never take negative or complex values. For a real symmetric matrix with negative eigenvalues, the singular values are the absolute values $|\lambda_i|$.
\end{trapbox}

\begin{trapbox}{Gotcha 2: Believing All-Positive Matrix Entries Imply Positive Definiteness}
A programmer checks a covariance candidate matrix $C = \begin{bmatrix} 1 & 3 \\ 3 & 1 \end{bmatrix}$ and assumes it is positive definite because all entries are positive ($C_{ij} > 0$).
\textbf{The Trap:} Positive matrix entries have \textbf{no relation} to positive definiteness. Evaluating the determinant: $\det(C) = 1(1) - 3(3) = -8 < 0$. The eigenvalues are $\lambda_1 = 4$ and $\lambda_2 = -2$. For $\mathbf{x} = [1, -1]^T$, the quadratic form is $\mathbf{x}^T C \mathbf{x} = 1 - 3 - 3 + 1 = -4 < 0$. The matrix is indefinite!
\end{trapbox}

\begin{trapbox}{Gotcha 3: Assuming $U$ and $V$ are Interchangeable}
An analyst switches $U$ and $V$, writing $A = V \Sigma U^T$.
\textbf{The Trap:} $U \in \mathbb{R}^{m \times m}$ contains the eigenvectors of $A A^T$ (spanning $\mathbb{R}^m$), while $V \in \mathbb{R}^{n \times n}$ contains the eigenvectors of $A^T A$ (spanning $\mathbb{R}^n$). If $A$ is rectangular ($m \ne n$), swapping $U$ and $V$ yields a matrix of dimension $n \times m$, which is the transpose $A^T$, not $A$.
\end{trapbox}

\begin{trapbox}{Gotcha 4: Measuring Variance Retained by Linear Sum of Singular Values}
In truncated SVD, an engineer claims that rank-$k$ approximation retains $\frac{\sum_{i=1}^k \sigma_i}{\sum_{i=1}^r \sigma_i}$ of the total dataset variance.
\textbf{The Trap:} Total variance (Frobenius energy) corresponds to the sum of \textbf{squared} singular values: $\|A\|_F^2 = \sum_{i=1}^r \sigma_i^2$. The true proportion of retained energy is $\frac{\sum_{i=1}^k \sigma_i^2}{\sum_{i=1}^r \sigma_i^2}$. Using linear sums severely overestimates the information preserved in weak singular components.
\end{trapbox}

% ----------------------------------------------------------------------
% PART 8: COMPARATIVE SYNTHESIS TABLE
% ----------------------------------------------------------------------
\section{Comparative Synthesis Table}

Table~\ref{tab:eigen_vs_svd} provides a comprehensive comparison between Eigendecomposition and Singular Value Decomposition.

\begin{table}[htbp]
\centering
\small
\begin{tabularx}{\textwidth}{lXX}
\toprule
\textbf{Property} & \textbf{Eigendecomposition ($A = P D P^{-1}$)} & \textbf{Singular Value Decomposition ($A = U \Sigma V^T$)} \\
\midrule
Matrix Applicability & Square matrices only ($n \times n$) & Any rectangular matrix ($m \times n$) \\
\addlinespace
Existence Guarantee & Only if matrix is non-defective & Always exists for every matrix \\
\addlinespace
Basis Vectors & Columns of $P$ (generally non-orthogonal) & Columns of $U, V$ (strictly orthonormal) \\
\addlinespace
Diagonal Elements & Eigenvalues $\lambda_i \in \mathbb{C}$ (can be negative/complex) & Singular values $\sigma_i \in \mathbb{R}_{\ge 0}$ (sorted decreasingly) \\
\addlinespace
Geometric Action & Stretches along invariant eigenspaces & Rotates domain ($V^T$), scales axes ($\Sigma$), rotates codomain ($U$) \\
\addlinespace
Low-Rank Truncation & Generally not optimal for low-rank approximation & Provably optimal under both $L_2$ and Frobenius norms \\
\bottomrule
\end{tabularx}
\caption{Comparative Synthesis: Eigendecomposition vs. Singular Value Decomposition.}
\label{tab:eigen_vs_svd}
\end{table}

% ----------------------------------------------------------------------
% PART 9: PYTHON / SCIENTIFIC COMPUTING LAB
% ----------------------------------------------------------------------
\section{Python / Scientific Computing Lab}

The production script below implements truncated SVD for matrix compression, verifies the Eckart-Young-Mirsky theorem, and validates positive definiteness via Cholesky decomposition.

\begin{lstlisting}[language=Python, caption={SVD Low-Rank Approximation and Positive Definiteness in Python.}]
import numpy as np

def compute_svd_approximation(A: np.ndarray, k: int):
    """Computes truncated rank-k SVD approximation: A_k = sum_{i=1}^k sigma_i u_i v_i^T."""
    U, s, Vt = np.linalg.svd(A, full_matrices=False)
    # Truncate to top k components
    U_k = U[:, :k]
    s_k = s[:k]
    Vt_k = Vt[:k, :]
    A_k = U_k @ np.diag(s_k) @ Vt_k

    # Compute theoretical vs empirical Frobenius error
    emp_frob_err = np.linalg.norm(A - A_k, 'fro')
    theory_frob_err = np.sqrt(np.sum(s[k:] ** 2))
    return A_k, emp_frob_err, theory_frob_err

def check_positive_definite(M: np.ndarray) -> bool:
    """Checks positive definiteness using Cholesky factorization."""
    # Symmetrize to guard against numerical float errors
    M_sym = (M + M.T) / 2.0
    try:
        _ = np.linalg.cholesky(M_sym)
        return True
    except np.linalg.LinAlgError:
        return False

def main():
    rng = np.random.RandomState(42)
    m, n = 50, 20
    # Generate low-rank matrix with noise
    rank_true = 5
    A_clean = rng.randn(m, rank_true) @ rng.randn(rank_true, n)
    noise = 0.05 * rng.randn(m, n)
    A = A_clean + noise

    print("Verifying Eckart-Young-Mirsky Theorem:")
    for k in [1, 3, 5, 10]:
        _, emp_err, theory_err = compute_svd_approximation(A, k)
        print(f"  Rank-{k:2d}: Empirical ||A - A_k||_F = {emp_err:.4f}, Theoretical = {theory_err:.4f}")

    # Positive Definite Checks
    M_pos = A.T @ A
    M_indef = np.array([[1.0, 3.0], [3.0, 1.0]])

    print(f"\nA^T A is Positive Definite: {check_positive_definite(M_pos)}")
    print(f"Matrix [[1, 3], [3, 1]] is Positive Definite: {check_positive_definite(M_indef)}")

if __name__ == "__main__":
    main()
\end{lstlisting}

% ----------------------------------------------------------------------
% PART 10: COMPREHENSIVE PRACTICE PROBLEMS
% ----------------------------------------------------------------------
\section{Comprehensive Practice Problems}

\subsection*{Part A: Analytical \& Mathematical Problems}

\begin{enumerate}[label=\textbf{P\arabic*.}]
    \item \textbf{Explicit Full SVD Computation by Hand:}
    Let $A = \begin{bmatrix} 3 & 0 \\ 4 & 5 \end{bmatrix}$.
    \begin{enumerate}[label=(\alph*)]
        \item Compute the Gram matrix $A^T A$ and find its eigenvalues $\lambda_1, \lambda_2$.
        \item Determine the singular values $\sigma_1, \sigma_2$.
        \item Find the orthonormal right singular vectors $\mathbf{v}_1, \mathbf{v}_2$.
        \item Compute left singular vectors $\mathbf{u}_i = \frac{1}{\sigma_i} A \mathbf{v}_i$ and verify $U \Sigma V^T = A$.
    \end{enumerate}

    \item \textbf{Sylvester's Criterion vs. Cholesky Factorization:}
    Let $M = \begin{bmatrix} 2 & -1 & 0 \\ -1 & 2 & -1 \\ 0 & -1 & 2 \end{bmatrix}$.
    \begin{enumerate}[label=(\alph*)]
        \item Compute all three leading principal minors: $\det(M_{1:1})$, $\det(M_{1:2, 1:2})$, and $\det(M)$.
        \item Use Sylvester's Criterion to determine if $M$ is strictly positive definite.
        \item Construct the lower-triangular Cholesky factor $L$ such that $M = L L^T$.
    \end{enumerate}

    \item \textbf{Frobenius Norm Representation in Terms of Singular Values:}
    Prove that for any matrix $A \in \mathbb{R}^{m \times n}$ with singular values $\sigma_1, \dots, \sigma_r$:
    \begin{equation}
    \|A\|_F^2 = \trace(A^T A) = \sum_{i=1}^r \sigma_i^2
    \end{equation}

    \item \textbf{Condition Number and Sensitivity in Linear Inversion:}
    The condition number of a matrix $A$ is defined as $\kappa(A) = \frac{\sigma_{\max}}{\sigma_{\min}}$.
    Suppose $A\mathbf{x} = \mathbf{b}$ and $A(\mathbf{x} + \delta \mathbf{x}) = \mathbf{b} + \delta \mathbf{b}$. Prove that:
    \begin{equation}
    \frac{\|\delta \mathbf{x}\|_2}{\|\mathbf{x}\|_2} \le \kappa(A) \frac{\|\delta \mathbf{b}\|_2}{\|\mathbf{b}\|_2}
    \end{equation}

    \item \textbf{Pseudoinverse of Rank-Deficient Matrix:}
    Let $A = \begin{bmatrix} 1 & 1 \\ 2 & 2 \end{bmatrix}$.
    \begin{enumerate}[label=(\alph*)]
        \item Determine $\rank(A)$ and compute its compact SVD $A = \sigma_1 \mathbf{u}_1 \mathbf{v}_1^T$.
        \item Compute the Moore-Penrose pseudoinverse $A^+ = \frac{1}{\sigma_1} \mathbf{v}_1 \mathbf{u}_1^T$.
        \item Verify that $A A^+ A = A$ and $A^+ A A^+ = A^+$.
    \end{enumerate}
\end{enumerate}

\subsection*{Part B: Diagnostic Multiple-Choice Questions}

\begin{enumerate}[label=\textbf{MCQ\arabic*.}]
    \item Let $A \in \mathbb{R}^{20 \times 10}$ have singular values $\sigma_1 = 12, \sigma_2 = 8, \sigma_3 = 5, \sigma_4 = 2$, with all other singular values equal to zero. What is the minimum Frobenius norm error $\min_{\rank(B) \le 2} \|A - B\|_F$?
    \begin{enumerate}[label=(\alph*)]
        \item $\sqrt{5^2 + 2^2} = \sqrt{29}$
        \item $5$
        \item $5 + 2 = 7$
        \item $\sqrt{12^2 + 8^2} = \sqrt{208}$
    \end{enumerate}

    \item Which of the following conditions is necessary and sufficient for a real symmetric matrix $A$ to be positive definite?
    \begin{enumerate}[label=(\alph*)]
        \item All entries $A_{ij} > 0$.
        \item The trace $\trace(A) > 0$ and determinant $\det(A) > 0$.
        \item All eigenvalues of $A$ are strictly positive.
        \item The diagonal entries are strictly greater than off-diagonal entries.
    \end{enumerate}
\end{enumerate}

\chapter[Principal Component Analysis]{Principal Component Analysis (PCA) and Dimensionality Reduction}
\label{chap:week07}

% ----------------------------------------------------------------------
% PART 1: EXECUTIVE ROADMAP & LEARNING OBJECTIVES
% ----------------------------------------------------------------------
\section{Executive Roadmap \& Learning Objectives}

Modern data science routinely handles empirical observations embedded in extraordinarily high-dimensional feature spaces ($d \sim 10^3$ to $10^6$), such as genomic expression profiles, facial images, consumer purchase histories, and neural embedding vectors. In these high-dimensional regimes, models confront the \textbf{curse of dimensionality}: volume grows exponentially, sample density vanishes into near-vacuum, pairwise distances between all points become equidistant, and severe multicollinearity renders linear models numerically unstable.

Yet, real-world data does not fill high-dimensional space uniformly. A set of facial photographs in $\mathbb{R}^{4096}$ (a $64 \times 64$ pixel grid) does not scatter across all 4,096 dimensions; rather, anatomical geometry, lighting, and head poses force the points to cluster tightly along a smooth, low-dimensional manifold.

\textbf{Principal Component Analysis (PCA)} is the foundational unsupervised statistical technique designed to discover these low-dimensional linear manifolds embedded within high-dimensional observations. PCA achieves this by projecting data onto an orthogonal coordinate system whose axes sequentially maximize empirical variance. In doing so, PCA:
\begin{enumerate}
    \item \textbf{Decorrelates Features:} Rotates the coordinate system so that all cross-feature covariances vanish ($\Cov(z_i, z_j) = 0$ for $i \ne j$).
    \item \textbf{Compresses Information:} Identifies the provably optimal $k$-dimensional linear subspace ($k \ll d$) that minimizes the mean squared Euclidean reconstruction error.
    \item \textbf{Unifies Statistics \& Geometry:} Bridges the statistical theory of covariance tensors with the geometric machinery of the Singular Value Decomposition (SVD).
\end{enumerate}

This chapter develops the rigorous mathematical foundations of PCA, establishes the equivalence between variance maximization and reconstruction error minimization, derives the high-dimensional dual Gram matrix shortcut for $d \gg N$, and explores critical diagnostic failure modes in machine learning pipelines.

\subsection{Prerequisites \& Scaffolding}
Mastery of this chapter synthesizes multivariate statistics, the Spectral Theorem, and Singular Value Decomposition as mapped in Table~\ref{tab:w07_prereqs}.

\begin{table}[htbp]
\centering
\small
\begin{tabularx}{\textwidth}{lXX}
\toprule
\textbf{Foundational Area} & \textbf{Core Concepts Required} & \textbf{Role in Machine Learning Foundations} \\
\midrule
Multivariate Statistics & Sample mean vector $\boldsymbol{\mu}$, sample covariance matrix $S$, variance & Defining the empirical dispersion tensor of feature spaces \\
\addlinespace
Spectral Theory & Eigendecomposition $S = Q \Lambda Q^T$, Rayleigh quotient & Identifying principal directions maximizing projected variance \\
\addlinespace
Matrix Factorizations & SVD $\tilde{X} = U \Sigma V^T$, Eckart-Young-Mirsky Theorem & Deriving stable and dual computational algorithms for PCA \\
\addlinespace
Constrained Optimization & Lagrange multipliers, stationary point conditions & Formulating the constrained quadratic optimization for principal axes \\
\bottomrule
\end{tabularx}
\caption{Prerequisite Scaffolding for Chapter 7.}
\label{tab:w07_prereqs}
\end{table}

\subsection{Core Competencies}
By the conclusion of this chapter, students will be able to:
\begin{itemize}
    \item Formulate PCA under two dual mathematical perspectives: Maximum Variance and Minimum Reconstruction Error, and prove their equivalence via the Pythagorean theorem.
    \item Rigorously derive using Lagrange multipliers why principal component loading vectors are the eigenvectors of the empirical sample covariance matrix.
    \item Establish the exact mathematical connection between SVD on centered data ($\frac{1}{\sqrt{N}}\tilde{X} = U \Sigma V^T$) and the spectral decomposition of the covariance matrix ($S = V \Lambda V^T$).
    \item Derive and implement the dual ($N \times N$) Gram matrix shortcut to compute exact PCA in high-dimensional settings where feature dimension vastly exceeds sample size ($d \gg N$).
    \item Diagnose scale sensitivity, select the optimal subspace dimension $k$ via scree plots and cumulative variance thresholds, and eliminate data leakage in validation pipelines.
\end{itemize}

\subsection{Notation \& Mathematical Conventions}
\begin{itemize}
    \item $X \in \mathbb{R}^{N \times d}$: The raw data matrix with $N$ rows (samples) and $d$ columns (features).
    \item $\boldsymbol{\mu} \triangleq \frac{1}{N} \sum_{i=1}^N \mathbf{x}_i \in \mathbb{R}^d$: The empirical sample mean vector.
    \item $\tilde{X} \triangleq X - \mathbf{1}_N \boldsymbol{\mu}^T \in \mathbb{R}^{N \times d}$: The mean-centered data matrix.
    \item $S \triangleq \frac{1}{N} \tilde{X}^T \tilde{X} \in \mathbb{R}^{d \times d}$: The empirical sample covariance matrix ($S \succeq 0$).
    \item $\mathbf{w}_1, \dots, \mathbf{w}_k \in \mathbb{R}^d$: Orthonormal principal component loading vectors (eigenvectors of $S$).
    \item $\mathbf{z}_i \triangleq W_k^T \tilde{\mathbf{x}}_i \in \mathbb{R}^k$: The low-dimensional representation (scores) of sample $\mathbf{x}_i$.
    \item $W_k = [\mathbf{w}_1, \dots, \mathbf{w}_k] \in \mathbb{R}^{d \times k}$: The orthogonal projection matrix ($W_k^T W_k = I_k$).
    \item $\hat{\mathbf{x}}_i \triangleq \boldsymbol{\mu} + W_k \mathbf{z}_i$: The reconstructed approximation of sample $\mathbf{x}_i$ in the original $d$-dimensional space.
    \item $K \triangleq \frac{1}{N}\tilde{X}\tilde{X}^T \in \mathbb{R}^{N \times N}$: The dual sample Gram matrix.
\end{itemize}

% ----------------------------------------------------------------------
% PART 2: THE GRAND PARADOX & INTUITIVE MOTIVATION
% ----------------------------------------------------------------------
\section{The Shadow That Tells the Whole Story}

\subsection{The Information Compression Paradox}
Imagine holding a complex 3D wireframe sculpture in front of a flat white wall. If you illuminate it with a single lamp, the sculpture casts a flat 2D shadow. 

If you shine the light from an arbitrary angle, the shadow is an incomprehensible, overlapping mess: intricate structural curves collapse onto each other, depth is scrambled, and vital details are obliterated. But if you slowly rotate the sculpture until you discover its widest, most expansive orientation, a miraculous transformation occurs: suddenly, the flat 2D shadow reveals almost the entire architecture of the 3D sculpture! Every curve, contour, and branch becomes instantly recognizable from the 2D silhouette alone.

\begin{quote}
\textit{\textbf{The Central Paradox:}} When projecting data from $d$ dimensions down to $k \ll d$ dimensions, we permanently discard $(d - k)$ dimensions of geometric information. How can a linear algorithm discard $99\%$ of its coordinates while preserving nearly $100\%$ of its meaningful informational content?
\end{quote}

The answer rests upon the geometry of \textbf{covariance}. High-dimensional data points in machine learning do not scatter randomly throughout $\mathbb{R}^d$. Due to underlying physical laws, biological regulatory networks, or economic correlations, the coordinates are intensely coupled. The data occupies a low-dimensional hyper-ellipse embedded inside $\mathbb{R}^d$.

Furthermore, linear projections are governed by an unyielding conservation law rooted in the Pythagorean theorem:
\begin{equation}
\text{Total Dataset Variance} = \text{Captured Shadow Variance} + \text{Lost Perpendicular Depth}
\end{equation}
Because the total variance of the centered points is a fixed physical constant, \textbf{maximizing the spread of the projected points} is mathematically identical to \textbf{minimizing the perpendicular drop distance from the points to the projection subspace}!

\begin{ponderbox}{The Shadow Geometry \& The Pythagorean Conservation Law}
Consider a zero-mean cloud of $N$ data points in $\mathbb{R}^d$. Let $\mathbf{w} \in \mathbb{R}^d$ be any candidate unit vector ($\|\mathbf{w}\|_2 = 1$).
\begin{enumerate}
    \item For each centered observation $\tilde{\mathbf{x}}_i$, decompose the vector into its parallel component along $\mathbf{w}$ and its perpendicular component orthogonal to $\mathbf{w}$.
    \item Apply the Pythagorean theorem to relate $\|\tilde{\mathbf{x}}_i\|_2^2$, the projected variance along $\mathbf{w}$, and the orthogonal reconstruction error.
    \item Sum this relation across all $N$ observations. Why does this prove that the direction maximizing projected variance is identically the direction minimizing reconstruction error?
\end{enumerate}
\tcblower
\textbf{Reflections and Mathematical Rationale:}
\begin{enumerate}
    \item The orthogonal projection of $\tilde{\mathbf{x}}_i$ onto unit vector $\mathbf{w}$ is:
    \begin{equation*}
    \mathbf{p}_i = (\mathbf{w}^T \tilde{\mathbf{x}}_i)\mathbf{w}
    \end{equation*}
    The perpendicular residual vector is:
    \begin{equation*}
    \mathbf{e}_i = \tilde{\mathbf{x}}_i - \mathbf{p}_i = \tilde{\mathbf{x}}_i - (\mathbf{w}^T \tilde{\mathbf{x}}_i)\mathbf{w}
    \end{equation*}
    By construction, $\mathbf{p}_i \perp \mathbf{e}_i$, because:
    \begin{equation*}
    \mathbf{w}^T \mathbf{e}_i = \mathbf{w}^T \tilde{\mathbf{x}}_i - (\mathbf{w}^T \tilde{\mathbf{x}}_i)(\mathbf{w}^T \mathbf{w}) = \mathbf{w}^T \tilde{\mathbf{x}}_i - \mathbf{w}^T \tilde{\mathbf{x}}_i(1) = 0
    \end{equation*}
    \item By the Pythagorean theorem in Euclidean space:
    \begin{equation*}
    \|\tilde{\mathbf{x}}_i\|_2^2 = \|\mathbf{p}_i\|_2^2 + \|\mathbf{e}_i\|_2^2 = (\mathbf{w}^T \tilde{\mathbf{x}}_i)^2 + \|\tilde{\mathbf{x}}_i - (\mathbf{w}^T \tilde{\mathbf{x}}_i)\mathbf{w}\|_2^2
    \end{equation*}
    \item Summing over all $N$ data points and dividing by $N$:
    \begin{equation*}
    \frac{1}{N}\sum_{i=1}^N \|\tilde{\mathbf{x}}_i\|_2^2 = \frac{1}{N}\sum_{i=1}^N (\mathbf{w}^T \tilde{\mathbf{x}}_i)^2 + \frac{1}{N}\sum_{i=1}^N \|\tilde{\mathbf{x}}_i - (\mathbf{w}^T \tilde{\mathbf{x}}_i)\mathbf{w}\|_2^2
    \end{equation*}
    Notice that the left-hand side is the total sample variance of the dataset, which is a fixed constant:
    \begin{equation*}
    \frac{1}{N}\sum_{i=1}^N \|\tilde{\mathbf{x}}_i\|_2^2 = \trace(S) = \text{Constant}
    \end{equation*}
    The first term on the right-hand side is the \textbf{projected variance} along $\mathbf{w}$. The second term is the \textbf{mean squared reconstruction error}. Since their sum is strictly constant, any choice of $\mathbf{w}$ that increases projected variance must decrease reconstruction error by the exact same amount. Maximizing spread and minimizing error are two sides of the same coin!
\end{enumerate}
\end{ponderbox}

\begin{figure}[htbp]
\centering
\begin{tikzpicture}[scale=1.0, font=\sffamily\footnotesize]

% Covariance Ellipse
\draw[dashed, color=gray!60, line width=1pt] (0, 0) [rotate=32] ellipse (4.0cm and 1.5cm);

% Scatter points
\foreach \pt in {
    (-2.8, -1.6), (-2.2, -1.1), (-1.9, -1.5), (-1.4, -0.6), (-1.1, -1.0),
    (-0.6, -0.2), (-0.4, -0.6), (0.1, 0.2), (0.5, -0.2), (0.8, 0.6),
    (1.3, 0.7), (1.6, 1.2), (2.0, 1.0), (2.4, 1.7), (2.9, 1.6)
}{
    \filldraw[color=primarydark] \pt circle (2.2pt);
}

% Origin / Mean
\filldraw[color=traprad] (0, 0) circle (3pt) node[below left] {$\boldsymbol{\mu} = \mathbf{0}$};

% First Principal Component (PC1)
\draw[->, >=stealth, line width=2.2pt, color=tealborder] (0, 0) -- (3.39, 2.12) node[above right] {\textbf{PC1 } $\mathbf{w}_1$ ($\lambda_1 = \sigma_1^2$)};
\draw[line width=1.1pt, color=tealborder, dashed] (-4.0, -2.5) -- (4.0, 2.5);

% Second Principal Component (PC2)
\draw[->, >=stealth, line width=2.2pt, color=coralborder] (0, 0) -- (-1.06, 1.70) node[above left] {\textbf{PC2 } $\mathbf{w}_2 \perp \mathbf{w}_1$ ($\lambda_2$)};
\draw[line width=1.1pt, color=coralborder, dashed] (1.3, -2.08) -- (-1.3, 2.08);

% Specific projection demonstration
\coordinate (Xpt) at (1.3, 0.7);
% Projection onto line y = tan(32) x = 0.6248 x
% w1 = [cos 32, sin 32] = [0.848, 0.530]
% <w1, Xpt> = 1.3*0.848 + 0.7*0.530 = 1.1024 + 0.371 = 1.473
% Proj = 1.473 * [0.848, 0.530] = [1.25, 0.78]
\coordinate (ProjPt) at (1.45, 0.91);
\draw[dotted, line width=1.4pt, color=traprad] (Xpt) -- (ProjPt);
\filldraw[color=tealborder] (ProjPt) circle (2.5pt);

\draw[<->, >=stealth, line width=0.8pt, color=traprad] (1.35, 0.73) -- (1.43, 0.88);
\node[right, color=traprad, font=\scriptsize] at (1.5, 0.75) {Orthogonal Error $\|\mathbf{e}_i\|_2$};

\end{tikzpicture}
\caption{The Dual Geometry of Principal Component Analysis. The first principal component $\mathbf{w}_1$ aligns with the major axis of the covariance dispersion ellipse. Projecting onto $\mathbf{w}_1$ maximizes empirical variance while simultaneously minimizing the orthogonal Euclidean drop distance $\|\mathbf{e}_i\|_2$.}
\label{fig:pca_geometry}
\end{figure}

% ----------------------------------------------------------------------
% PART 3: MATHEMATICAL FORMULATIONS OF PCA
% ----------------------------------------------------------------------
\section{Mathematical Formulations of PCA}

\subsection{Formulation 1: Maximizing Projected Variance}
Let $\tilde{\mathbf{x}} \in \mathbb{R}^d$ be a zero-mean random vector with sample covariance matrix:
\begin{equation}
S \triangleq \frac{1}{N}\sum_{i=1}^N \tilde{\mathbf{x}}_i \tilde{\mathbf{x}}_i^T = \frac{1}{N}\tilde{X}^T \tilde{X} \in \mathbb{R}^{d \times d}
\end{equation}
We seek a unit direction $\mathbf{w} \in \mathbb{R}^d$ ($\|\mathbf{w}\|_2 = 1$) such that projecting the centered points onto $\mathbf{w}$ achieves maximal empirical variance:
\begin{equation}
\Var(\mathbf{w}^T \tilde{\mathbf{x}}) = \frac{1}{N} \sum_{i=1}^N (\mathbf{w}^T \tilde{\mathbf{x}}_i)^2 = \mathbf{w}^T \left( \frac{1}{N}\sum_{i=1}^N \tilde{\mathbf{x}}_i \tilde{\mathbf{x}}_i^T \right) \mathbf{w} = \mathbf{w}^T S \mathbf{w}
\end{equation}

\begin{theorembox}{First Principal Component Optimization}
The unit direction $\mathbf{w}_1$ that maximizes the projected sample variance satisfies the constrained optimization problem:
\begin{equation}
\max_{\mathbf{w} \in \mathbb{R}^d} \mathbf{w}^T S \mathbf{w} \quad \text{subject to } \mathbf{w}^T \mathbf{w} = 1
\label{eq:pca_var_opt}
\end{equation}
The optimal direction $\mathbf{w}_1$ is the eigenvector of the covariance matrix $S$ associated with its largest eigenvalue $\lambda_1 = \lambda_{\max}(S)$:
\begin{equation}
S \mathbf{w}_1 = \lambda_1 \mathbf{w}_1
\end{equation}
Furthermore, the maximum projected variance achieved along this axis is identically $\lambda_1$:
\begin{equation}
\Var(\mathbf{w}_1^T \tilde{\mathbf{x}}) = \mathbf{w}_1^T S \mathbf{w}_1 = \lambda_1
\end{equation}
\end{theorembox}

\begin{proof}
We formulate the Lagrangian function introducing scalar multiplier $\lambda \in \mathbb{R}$ to enforce the unit-norm constraint:
\begin{equation}
\mathcal{L}(\mathbf{w}, \lambda) = \mathbf{w}^T S \mathbf{w} - \lambda (\mathbf{w}^T \mathbf{w} - 1)
\end{equation}
Differentiating with respect to $\mathbf{w}$ using denominator layout matrix calculus:
\begin{equation}
\nabla_{\mathbf{w}} \mathcal{L} = 2S\mathbf{w} - 2\lambda \mathbf{w} = \mathbf{0} \implies S\mathbf{w} = \lambda \mathbf{w}
\label{eq:pca_lagrange_condition}
\end{equation}
Thus, every stationary point of the Lagrangian is an eigenvector of $S$ with eigenvalue $\lambda$. Pre-multiplying both sides by $\mathbf{w}^T$:
\begin{equation}
\mathbf{w}^T S \mathbf{w} = \lambda \mathbf{w}^T \mathbf{w} = \lambda (1) = \lambda
\end{equation}
To achieve the absolute global maximum of the projected variance $\mathbf{w}^T S \mathbf{w}$, we must select the eigenvector associated with the largest eigenvalue $\lambda_1 = \lambda_{\max}(S)$.

For the second principal component $\mathbf{w}_2$, we require $\mathbf{w}_2$ to maximize $\mathbf{w}^T S \mathbf{w}$ subject to $\mathbf{w}^T \mathbf{w} = 1$ and mutual orthogonality $\mathbf{w}_2^T \mathbf{w}_1 = 0$. Formulating the Lagrangian with two multipliers $\lambda, \phi$:
\begin{equation}
\mathcal{L}(\mathbf{w}_2, \lambda, \phi) = \mathbf{w}_2^T S \mathbf{w}_2 - \lambda (\mathbf{w}_2^T \mathbf{w}_2 - 1) - \phi (\mathbf{w}_2^T \mathbf{w}_1)
\end{equation}
Differentiating:
\begin{equation}
2S\mathbf{w}_2 - 2\lambda \mathbf{w}_2 - \phi \mathbf{w}_1 = \mathbf{0}
\end{equation}
Pre-multiplying by $\mathbf{w}_1^T$:
\begin{equation}
2\mathbf{w}_1^T S \mathbf{w}_2 - 2\lambda \mathbf{w}_1^T \mathbf{w}_2 - \phi \mathbf{w}_1^T \mathbf{w}_1 = 0
\end{equation}
Since $\mathbf{w}_1^T \mathbf{w}_2 = 0$ and $S\mathbf{w}_1 = \lambda_1 \mathbf{w}_1$, we have $\mathbf{w}_1^T S \mathbf{w}_2 = (S\mathbf{w}_1)^T \mathbf{w}_2 = \lambda_1 \mathbf{w}_1^T \mathbf{w}_2 = 0$. Thus $\phi (1) = 0 \implies \phi = 0$.
The condition reduces to $S\mathbf{w}_2 = \lambda \mathbf{w}_2$. To maximize variance while remaining orthogonal to $\mathbf{w}_1$, we select the second largest eigenvalue $\lambda_2$. By mathematical induction, the $k$-th principal component is the $k$-th eigenvector of $S$.
\end{proof}

\subsection{Formulation 2: Minimizing Reconstruction Error}
Now consider the dual perspective. We seek an orthonormal basis $W_k = [\mathbf{w}_1, \dots, \mathbf{w}_k] \in \mathbb{R}^{d \times k}$ ($W_k^T W_k = I_k$) spanning a $k$-dimensional linear subspace such that projecting each centered sample $\tilde{\mathbf{x}}_i$ onto this subspace minimizes the mean squared Euclidean reconstruction error:
\begin{equation}
\min_{W_k \in \mathbb{R}^{d \times k}, W_k^T W_k = I_k} \mathcal{J}(W_k) \triangleq \frac{1}{N} \sum_{i=1}^N \|\tilde{\mathbf{x}}_i - W_k W_k^T \tilde{\mathbf{x}}_i\|_2^2
\label{eq:pca_recon_opt}
\end{equation}

\begin{theorembox}{Equivalence of PCA Formulations}
Maximizing empirical projected variance is mathematically identical to minimizing the mean squared reconstruction error. For any target dimension $k \le d$, the optimal projection matrix $W_k$ consists of the top $k$ eigenvectors of the sample covariance matrix $S$.

Furthermore, the minimum reconstruction error achieved is strictly the sum of the discarded eigenvalues:
\begin{equation}
\min_{W_k} \frac{1}{N}\sum_{i=1}^N \|\tilde{\mathbf{x}}_i - W_k W_k^T \tilde{\mathbf{x}}_i\|_2^2 = \sum_{j=k+1}^d \lambda_j
\label{eq:pca_error_sum}
\end{equation}
\end{theorembox}

\begin{proof}
Let $P_k = W_k W_k^T \in \mathbb{R}^{d \times d}$ be the rank-$k$ orthogonal projection operator onto $\Col(W_k)$. The orthogonal complement projector is $P_\perp = I_d - P_k$.
Expanding the squared reconstruction error for sample $\tilde{\mathbf{x}}_i$:
\begin{equation}
\|\tilde{\mathbf{x}}_i - P_k \tilde{\mathbf{x}}_i\|_2^2 = \|(I_d - P_k)\tilde{\mathbf{x}}_i\|_2^2 = \tilde{\mathbf{x}}_i^T (I_d - P_k)^T (I_d - P_k) \tilde{\mathbf{x}}_i
\end{equation}
Because orthogonal projectors are symmetric and idempotent ($(I_d - P_k)^2 = I_d - P_k$):
\begin{equation}
\|\tilde{\mathbf{x}}_i - P_k \tilde{\mathbf{x}}_i\|_2^2 = \tilde{\mathbf{x}}_i^T (I_d - P_k) \tilde{\mathbf{x}}_i = \tilde{\mathbf{x}}_i^T \tilde{\mathbf{x}}_i - \tilde{\mathbf{x}}_i^T P_k \tilde{\mathbf{x}}_i
\end{equation}
Summing over all $N$ observations and dividing by $N$:
\begin{align}
\mathcal{J}(W_k) &= \frac{1}{N}\sum_{i=1}^N \tilde{\mathbf{x}}_i^T \tilde{\mathbf{x}}_i - \frac{1}{N}\sum_{i=1}^N \tilde{\mathbf{x}}_i^T (W_k W_k^T) \tilde{\mathbf{x}}_i \notag \\
&= \trace\left( \frac{1}{N}\sum_{i=1}^N \tilde{\mathbf{x}}_i \tilde{\mathbf{x}}_i^T \right) - \trace\left( W_k^T \left( \frac{1}{N}\sum_{i=1}^N \tilde{\mathbf{x}}_i \tilde{\mathbf{x}}_i^T \right) W_k \right) \notag \\
&= \trace(S) - \trace(W_k^T S W_k)
\end{align}
Notice that $\trace(S) = \sum_{j=1}^d \lambda_j$ is the total variance of the data, which is completely independent of $W_k$.
Therefore:
\begin{equation}
\min_{W_k} \mathcal{J}(W_k) \iff \max_{W_k} \trace(W_k^T S W_k) = \max_{W_k} \sum_{i=1}^k \mathbf{w}_i^T S \mathbf{w}_i
\end{equation}
This establishes exact mathematical equivalence. Maximizing $\sum_{i=1}^k \mathbf{w}_i^T S \mathbf{w}_i$ is achieved by choosing the top $k$ eigenvectors of $S$, giving captured variance $\sum_{i=1}^k \lambda_i$.
The residual error is:
\begin{equation}
\mathcal{J}(W_k) = \trace(S) - \sum_{i=1}^k \lambda_i = \sum_{j=1}^d \lambda_j - \sum_{i=1}^k \lambda_i = \sum_{j=k+1}^d \lambda_j
\end{equation}
\end{proof}

\begin{figure}[htbp]
\centering
\begin{tikzpicture}[scale=0.95, font=\sffamily\footnotesize]

% Scree plot axes
\draw[->, >=stealth, color=gray!60] (0, 0) -- (7.5, 0) node[right] {Principal Component Index $k$};
\draw[->, >=stealth, color=gray!60] (0, 0) -- (0, 4.5) node[above] {Eigenvalue / Variance $\lambda_k$};

% Bar chart representing eigenvalues
\draw[fill=tealborder!70, draw=tealborder] (0.6, 0) rectangle (1.2, 3.8);
\node[above] at (0.9, 3.8) {$\lambda_1=38\%$};

\draw[fill=tealborder!70, draw=tealborder] (1.8, 0) rectangle (2.4, 2.7);
\node[above] at (2.1, 2.7) {$\lambda_2=27\%$};

\draw[fill=tealborder!70, draw=tealborder] (3.0, 0) rectangle (3.6, 1.6);
\node[above] at (3.3, 1.6) {$\lambda_3=16\%$};

\draw[fill=tealborder!40, draw=tealborder] (4.2, 0) rectangle (4.8, 0.6);
\node[above] at (4.5, 0.6) {$\lambda_4=6\%$};

\draw[fill=tealborder!40, draw=tealborder] (5.4, 0) rectangle (6.0, 0.4);
\node[above] at (5.7, 0.4) {$\lambda_5=4\%$};

\draw[fill=tealborder!40, draw=tealborder] (6.6, 0) rectangle (7.2, 0.2);
\node[above] at (6.9, 0.2) {$\lambda_6=2\%$};

% Scree line
\draw[thick, color=traprad, mark=*] plot coordinates {
    (0.9, 3.8) (2.1, 2.7) (3.3, 1.6) (4.5, 0.6) (5.7, 0.4) (6.9, 0.2)
};

% Elbow annotation
\draw[->, >=stealth, line width=1.5pt, color=traprad] (5.0, 2.2) -- (4.6, 0.9);
\node[align=center, fill=white, draw=traprad, rounded corners=3pt, inner sep=3pt] at (5.8, 2.6) {
    \textbf{The "Elbow" Cutoff} \\ Signal ends; noise floor begins!
};

\end{tikzpicture}
\caption{The Scree Plot Diagnostic for Subspace Dimensionality Selection. The individual eigenvalues $\lambda_k$ drop steeply across the signal subspace before flattening out along the "scree" (noise floor). Truncating at the elbow captures the majority of variance while discarding high-frequency noise.}
\label{fig:scree_plot}
\end{figure}

% ----------------------------------------------------------------------
% PART 4: PCA COMPUTATION VIA SVD & DUAL SHORTCUT
% ----------------------------------------------------------------------
\section{Algorithmic Computation: SVD \& The Dual Shortcut}

\subsection{Computing PCA via Direct SVD}
In production scientific computing, computing the sample covariance matrix $S = \frac{1}{N}\tilde{X}^T \tilde{X}$ explicitly is an \textbf{anti-pattern}. Forming $\tilde{X}^T \tilde{X}$ requires matrix multiplication that squares the condition number:
\begin{equation}
\kappa(S) = \kappa(\tilde{X})^2
\end{equation}
If $\tilde{X}$ has condition number $\kappa = 10^8$, forming $S$ squares it to $10^{16}$, completely destroying numerical precision in 64-bit IEEE floating-point arithmetic!

Instead, standard software (`scikit-learn`, MATLAB, LAPACK) computes PCA directly via the thin SVD of the scaled, centered data matrix:
\begin{equation}
\frac{1}{\sqrt{N}}\tilde{X} = U \Sigma V^T
\label{eq:pca_svd_direct}
\end{equation}
Let us evaluate the sample covariance matrix under this factorization:
\begin{equation}
S = \left(\frac{1}{\sqrt{N}}\tilde{X}\right)^T \left(\frac{1}{\sqrt{N}}\tilde{X}\right) = (U \Sigma V^T)^T (U \Sigma V^T) = V \Sigma^T (U^T U) \Sigma V^T = V \Sigma^2 V^T
\end{equation}
Because $V$ is orthogonal ($V^T V = I_d$) and $\Sigma^2$ is diagonal, this is \textbf{identically the spectral decomposition of $S$}:
\begin{itemize}
    \item The right singular vectors $V = [\mathbf{v}_1, \dots, \mathbf{v}_d]$ are the \textbf{principal component loading directions} $\mathbf{w}_i$.
    \item The singular values $\sigma_i$ are directly proportional to the standard deviation of each principal component:
    \begin{equation}
    \lambda_i = \sigma_i^2 \implies \sigma_i = \sqrt{\lambda_i} = \sqrt{\Var(\mathbf{w}_i^T \tilde{\mathbf{x}})}
    \end{equation}
    \item The projected low-dimensional scores $Z \in \mathbb{R}^{N \times k}$ are computed in one line:
    \begin{equation}
    Z = \tilde{X} V_k = \sqrt{N} U_k \Sigma_k
    \end{equation}
\end{itemize}

\begin{ponderbox}{The Center of Mass Fallacy: Why Mean Centering is Mandatory}
Suppose a data scientist bypasses mean centering and applies SVD directly to the raw, uncentered data matrix $X \in \mathbb{R}^{N \times d}$, where the global mean vector $\boldsymbol{\mu} \ne \mathbf{0}$.
\begin{enumerate}
    \item Expand the uncentered second-moment matrix $\frac{1}{N} X^T X$ in terms of the centered covariance matrix $S = \frac{1}{N}\tilde{X}^T \tilde{X}$ and mean vector $\boldsymbol{\mu}$.
    \item When $\|\boldsymbol{\mu}\|_2^2 \gg \lambda_1(S)$, what direction will the first principal component $\mathbf{v}_1$ point toward?
    \item Why does this completely destroy the dimensionality reduction objective of PCA?
\end{enumerate}
\tcblower
\textbf{Reflections and Mathematical Rationale:}
\begin{enumerate}
    \item Decomposing each raw observation as $\mathbf{x}_i = \tilde{\mathbf{x}}_i + \boldsymbol{\mu}$, the uncentered second-moment matrix expands as:
    \begin{equation*}
    \frac{1}{N} X^T X = \frac{1}{N}\sum_{i=1}^N (\tilde{\mathbf{x}}_i + \boldsymbol{\mu})(\tilde{\mathbf{x}}_i + \boldsymbol{\mu})^T = \frac{1}{N}\sum_{i=1}^N \tilde{\mathbf{x}}_i \tilde{\mathbf{x}}_i^T + \frac{1}{N}\sum_{i=1}^N \tilde{\mathbf{x}}_i \boldsymbol{\mu}^T + \frac{1}{N}\sum_{i=1}^N \boldsymbol{\mu}\tilde{\mathbf{x}}_i^T + \boldsymbol{\mu}\boldsymbol{\mu}^T
    \end{equation*}
    Because $\sum \tilde{\mathbf{x}}_i = \mathbf{0}$ by definition of the mean, the two cross-terms vanish identically:
    \begin{equation*}
    \frac{1}{N} X^T X = S + \boldsymbol{\mu}\boldsymbol{\mu}^T
    \end{equation*}
    \item Notice that $\boldsymbol{\mu}\boldsymbol{\mu}^T$ is a rank-1 symmetric matrix whose sole non-zero eigenvalue is $\|\boldsymbol{\mu}\|_2^2$, with eigenvector $\frac{\boldsymbol{\mu}}{\|\boldsymbol{\mu}\|_2}$. When the data cloud is shifted far from the origin ($\|\boldsymbol{\mu}\|_2^2 \gg \lambda_1(S)$), Weyl's eigenvalue perturbation theorem guarantees that the dominant eigenvector of $\frac{1}{N}X^T X$ will align almost perfectly with the mean vector:
    \begin{equation*}
    \mathbf{v}_1 \approx \frac{\boldsymbol{\mu}}{\|\boldsymbol{\mu}\|_2}
    \end{equation*}
    \item This completely subverts PCA! Instead of identifying the internal axis of maximal variance across the data points, the first principal component is entirely consumed by describing the static displacement from the arbitrary origin to the center of mass. Mean centering strips out the rank-1 component $\boldsymbol{\mu}\boldsymbol{\mu}^T$, ensuring that every singular vector isolates true geometric variance.
\end{enumerate}
\end{ponderbox}

\subsection{The Dual Gram Matrix Shortcut When \texorpdfstring{$d \gg N$}{d >> N}}
In genomics, microarrays, text embeddings, and computer vision, we frequently encounter datasets where the feature dimension vastly exceeds the number of observations:
\begin{equation}
d \gg N \quad (\text{e.g., } N = 50 \text{ patients}, \; d = 100{,}000 \text{ genes})
\end{equation}
Forming the standard covariance matrix $S = \frac{1}{N}\tilde{X}^T \tilde{X}$ requires constructing a gigantic $100{,}000 \times 100{,}000$ matrix, consuming $40$ Gigabytes of RAM and $\mathcal{O}(d^3) \approx 10^{15}$ operations for eigendecomposition!

\begin{theorembox}{The Dual PCA Gram Shortcut}
Let $\tilde{X} \in \mathbb{R}^{N \times d}$ with $d \gg N$. The non-zero eigenvalues of the enormous $d \times d$ covariance matrix $S = \frac{1}{N}\tilde{X}^T \tilde{X}$ are identical to the non-zero eigenvalues of the compact $N \times N$ sample Gram matrix:
\begin{equation}
K \triangleq \frac{1}{N}\tilde{X}\tilde{X}^T \in \mathbb{R}^{N \times N}
\label{eq:dual_gram_matrix}
\end{equation}
Furthermore, if $(\lambda, \mathbf{u})$ is an eigenpair of $K$ ($K\mathbf{u} = \lambda \mathbf{u}$ with $\lambda > 0$ and $\|\mathbf{u}\|_2 = 1$), then:
\begin{equation}
\mathbf{v} = \frac{1}{\sqrt{N\lambda}} \tilde{X}^T \mathbf{u} \in \mathbb{R}^d
\label{eq:dual_v_reconstruction}
\end{equation}
is a strictly normalized unit eigenvector of the full $d \times d$ covariance matrix $S$ with eigenvalue $\lambda$.
\end{theorembox}

\begin{proof}
Let $K\mathbf{u} = \lambda \mathbf{u}$ with $\lambda > 0$. Expanding the definition of $K$:
\begin{equation}
\left( \frac{1}{N}\tilde{X}\tilde{X}^T \right) \mathbf{u} = \lambda \mathbf{u}
\end{equation}
Pre-multiply both sides of this equation by the transposed data matrix $\tilde{X}^T$:
\begin{equation}
\tilde{X}^T \left( \frac{1}{N}\tilde{X}\tilde{X}^T \mathbf{u} \right) = \lambda (\tilde{X}^T \mathbf{u}) \implies \left( \frac{1}{N}\tilde{X}^T \tilde{X} \right) (\tilde{X}^T \mathbf{u}) = \lambda (\tilde{X}^T \mathbf{u})
\end{equation}
Recognizing the sample covariance matrix $S = \frac{1}{N}\tilde{X}^T \tilde{X}$:
\begin{equation}
S (\tilde{X}^T \mathbf{u}) = \lambda (\tilde{X}^T \mathbf{u})
\end{equation}
Thus, the high-dimensional vector $\mathbf{v} \propto \tilde{X}^T \mathbf{u}$ is an exact eigenvector of $S$ with the exact same eigenvalue $\lambda$!

Now, let us calculate the squared Euclidean norm of $\tilde{X}^T \mathbf{u}$:
\begin{equation}
\|\tilde{X}^T \mathbf{u}\|_2^2 = (\tilde{X}^T \mathbf{u})^T (\tilde{X}^T \mathbf{u}) = \mathbf{u}^T (\tilde{X}\tilde{X}^T) \mathbf{u} = \mathbf{u}^T (N K \mathbf{u}) = N \lambda \mathbf{u}^T \mathbf{u} = N \lambda
\end{equation}
where we have used $\|\mathbf{u}\|_2 = 1$. Taking the positive square root gives $\|\tilde{X}^T \mathbf{u}\|_2 = \sqrt{N\lambda}$.
Dividing by this normalization constant yields the unit eigenvector:
\begin{equation}
\mathbf{v} = \frac{1}{\sqrt{N\lambda}} \tilde{X}^T \mathbf{u}
\end{equation}
\end{proof}

\begin{ponderbox}{The Dual Gram Matrix Shortcut}
Suppose you are analyzing $N = 50$ cancer patient tissue samples, where each sample is profiled across $d = 100{,}000$ gene expression biomarkers.
\begin{enumerate}
    \item What is the matrix dimension and RAM memory required to store the standard covariance matrix $S = \frac{1}{N}\tilde{X}^T \tilde{X}$ in 32-bit floating-point numbers?
    \item What is the matrix dimension and RAM memory required to store the dual Gram matrix $K = \frac{1}{N}\tilde{X}\tilde{X}^T$?
    \item Compare the computational FLOPs required to compute the eigendecomposition of $S$ vs. $K$. What is the speedup factor?
\end{enumerate}
\tcblower
\textbf{Reflections and Mathematical Rationale:}
\begin{enumerate}
    \item The standard covariance matrix $S$ has dimension $d \times d = 100{,}000 \times 100{,}000 = 10^{10}$ entries. In standard 32-bit floats (4 bytes per number):
    \begin{equation*}
    \text{RAM} = 10^{10} \times 4 \text{ bytes} = 40{,}000{,}000{,}000 \text{ bytes} = \mathbf{40 \text{ Gigabytes!}}
    \end{equation*}
    Most standard laptops and cloud worker instances would instantly crash with an Out-Of-Memory (OOM) error.
    \item The dual Gram matrix $K$ has dimension $N \times N = 50 \times 50 = 2{,}500$ entries:
    \begin{equation*}
    \text{RAM} = 2{,}500 \times 4 \text{ bytes} = 10{,}000 \text{ bytes} = \mathbf{10 \text{ Kilobytes!}}
    \end{equation*}
    The dual matrix fits comfortably inside the L1 CPU cache!
    \item Eigendecomposition of a symmetric matrix scales cubically $\mathcal{O}(n^3)$:
    \begin{itemize}
        \item For $S$: $\mathcal{O}(d^3) = (10^5)^3 = 10^{15}$ floating-point operations.
        \item For $K$: $\mathcal{O}(N^3) = (50)^3 = 1.25 \times 10^5$ operations.
    \end{itemize}
    The computational speedup is:
    \begin{equation*}
    \text{Speedup} = \frac{10^{15}}{1.25 \times 10^5} = \mathbf{8{,}000{,}000{,}000 \times}
    \end{equation*}
    The dual shortcut runs in less than one millisecond, producing mathematically identical eigenvalues and eigenvectors!
\end{enumerate}
\end{ponderbox}

% ----------------------------------------------------------------------
% PART 5: WORKED NUMERICAL EXAMPLES
% ----------------------------------------------------------------------
\section{In-Depth Worked Numerical Examples}

\begin{examplebox}{Manual PCA on a 2D Dataset from Scratch}
Consider a dataset with $N = 4$ observations in $d = 2$ dimensions:
\begin{equation}
\mathbf{x}_1 = \begin{bmatrix} 1 \\ 1 \end{bmatrix}, \quad \mathbf{x}_2 = \begin{bmatrix} 2 \\ 3 \end{bmatrix}, \quad \mathbf{x}_3 = \begin{bmatrix} 3 \\ 2 \end{bmatrix}, \quad \mathbf{x}_4 = \begin{bmatrix} 4 \\ 4 \end{bmatrix}
\end{equation}
\begin{enumerate}
    \item Compute the empirical sample mean $\boldsymbol{\mu}$ and centered data matrix $\tilde{X}$.
    \item Compute the sample covariance matrix $S = \frac{1}{N}\tilde{X}^T \tilde{X}$.
    \item Compute the eigenvalues and orthonormal eigenvectors of $S$.
    \item Find the percentage of variance explained by the first principal component.
    \item Project sample $\mathbf{x}_4$ onto the first principal component and compute its reconstructed point $\hat{\mathbf{x}}_4$.
\end{enumerate}

\textbf{Step 1: Compute Sample Mean and Centering.}
\begin{equation}
\boldsymbol{\mu} = \frac{1}{4} \left( \begin{bmatrix} 1 \\ 1 \end{bmatrix} + \begin{bmatrix} 2 \\ 3 \end{bmatrix} + \begin{bmatrix} 3 \\ 2 \end{bmatrix} + \begin{bmatrix} 4 \\ 4 \end{bmatrix} \right) = \frac{1}{4} \begin{bmatrix} 10 \\ 10 \end{bmatrix} = \begin{bmatrix} 2.5 \\ 2.5 \end{bmatrix}
\end{equation}
Subtracting $\boldsymbol{\mu}$ from each row yields $\tilde{X} \in \mathbb{R}^{4 \times 2}$:
\begin{equation}
\tilde{X} = \begin{bmatrix} 1 - 2.5 & 1 - 2.5 \\ 2 - 2.5 & 3 - 2.5 \\ 3 - 2.5 & 2 - 2.5 \\ 4 - 2.5 & 4 - 2.5 \end{bmatrix} = \begin{bmatrix} -1.5 & -1.5 \\ -0.5 & 0.5 \\ 0.5 & -0.5 \\ 1.5 & 1.5 \end{bmatrix}
\end{equation}
Quick verification: $\sum \tilde{\mathbf{x}}_i = [-1.5 - 0.5 + 0.5 + 1.5, -1.5 + 0.5 - 0.5 + 1.5]^T = [0, 0]^T \quad \checkmark$.

\textbf{Step 2: Compute Sample Covariance Matrix $S$.}
\begin{align}
S &= \frac{1}{4} \tilde{X}^T \tilde{X} = \frac{1}{4} \begin{bmatrix} -1.5 & -0.5 & 0.5 & 1.5 \\ -1.5 & 0.5 & -0.5 & 1.5 \end{bmatrix} \begin{bmatrix} -1.5 & -1.5 \\ -0.5 & 0.5 \\ 0.5 & -0.5 \\ 1.5 & 1.5 \end{bmatrix} \notag \\
&= \frac{1}{4} \begin{bmatrix} 2.25 + 0.25 + 0.25 + 2.25 & 2.25 - 0.25 - 0.25 + 2.25 \\ 2.25 - 0.25 - 0.25 + 2.25 & 2.25 + 0.25 + 0.25 + 2.25 \end{bmatrix} = \frac{1}{4} \begin{bmatrix} 5.0 & 4.0 \\ 4.0 & 5.0 \end{bmatrix} = \begin{bmatrix} 1.25 & 1.00 \\ 1.00 & 1.25 \end{bmatrix}
\end{align}

\textbf{Step 3: Compute Eigenvalues and Eigenvectors.}
\begin{align}
\det(S - \lambda I) &= \det\begin{bmatrix} 1.25 - \lambda & 1.00 \\ 1.00 & 1.25 - \lambda \end{bmatrix} = (1.25 - \lambda)^2 - 1.00^2 \notag \\
&= \lambda^2 - 2.5\lambda + 1.5625 - 1.0 = \lambda^2 - 2.5\lambda + 0.5625 = (\lambda - 2.25)(\lambda - 0.25) = 0
\end{align}
The eigenvalues are $\lambda_1 = 2.25$ and $\lambda_2 = 0.25$.
\begin{itemize}
    \item \textbf{For $\lambda_1 = 2.25$:}
    \begin{equation}
    (S - 2.25I)\mathbf{w}_1 = \begin{bmatrix} -1.0 & 1.0 \\ 1.0 & -1.0 \end{bmatrix} \begin{bmatrix} w_{11} \\ w_{12} \end{bmatrix} = \mathbf{0} \implies w_{11} = w_{12} \implies \mathbf{w}_1 = \frac{1}{\sqrt{2}}\begin{bmatrix} 1 \\ 1 \end{bmatrix}
    \end{equation}
    \item \textbf{For $\lambda_2 = 0.25$:}
    \begin{equation}
    (S - 0.25I)\mathbf{w}_2 = \begin{bmatrix} 1.0 & 1.0 \\ 1.0 & 1.0 \end{bmatrix} \begin{bmatrix} w_{21} \\ w_{22} \end{bmatrix} = \mathbf{0} \implies w_{21} = -w_{22} \implies \mathbf{w}_2 = \frac{1}{\sqrt{2}}\begin{bmatrix} -1 \\ 1 \end{bmatrix}
    \end{equation}
\end{itemize}

\textbf{Step 4: Proportion of Explained Variance.}
Total variance is $\trace(S) = 1.25 + 1.25 = 2.50 = \lambda_1 + \lambda_2$.
\begin{equation}
\frac{\lambda_1}{\lambda_1 + \lambda_2} = \frac{2.25}{2.50} = 0.90 \implies \mathbf{90\% \text{ of variance is captured by PC1 alone!}}
\end{equation}

\textbf{Step 5: Projection and Reconstruction of $\mathbf{x}_4$.}
Centered sample is $\tilde{\mathbf{x}}_4 = [4 - 2.5, 4 - 2.5]^T = [1.5, 1.5]^T$.
The 1D scalar score is:
\begin{equation}
z_1 = \mathbf{w}_1^T \tilde{\mathbf{x}}_4 = \frac{1}{\sqrt{2}} (1.5(1) + 1.5(1)) = \frac{3.0}{\sqrt{2}} = 1.5\sqrt{2} \approx 2.121
\end{equation}
The reconstructed observation $\hat{\mathbf{x}}_4$ in original coordinates is:
\begin{equation}
\hat{\mathbf{x}}_4 = \boldsymbol{\mu} + z_1 \mathbf{w}_1 = \begin{bmatrix} 2.5 \\ 2.5 \end{bmatrix} + \left( \frac{3.0}{\sqrt{2}} \right) \left( \frac{1}{\sqrt{2}}\begin{bmatrix} 1 \\ 1 \end{bmatrix} \right) = \begin{bmatrix} 2.5 \\ 2.5 \end{bmatrix} + \begin{bmatrix} 1.5 \\ 1.5 \end{bmatrix} = \begin{bmatrix} 4.0 \\ 4.0 \end{bmatrix} = \mathbf{x}_4
\end{equation}
The reconstruction error is exactly zero ($\|\mathbf{x}_4 - \hat{\mathbf{x}}_4\|_2 = 0$) because $\tilde{\mathbf{x}}_4$ lies directly along the principal axis $\mathbf{w}_1$.
\end{examplebox}

\begin{examplebox}{The Collapse of Scale Invariance in PCA}
Consider a dataset where feature 1 is measured in meters and feature 2 is measured in kilometers, giving sample covariance:
\begin{equation}
S = \begin{bmatrix} 10 & 0 \\ 0 & 1 \end{bmatrix}
\end{equation}
An analyst changes feature 2 from kilometers to meters, multiplying all values by $c = 10$.
\begin{enumerate}
    \item Find the eigenvalues and principal component direction $\mathbf{w}_1$ before scaling.
    \item Compute the new covariance matrix $S'$ after scaling.
    \item Find the new principal component direction $\mathbf{w}_1'$.
    \item Contrast the results to show why feature standardization is mandatory.
\end{enumerate}

\textbf{Step 1: Before Scaling.}
$S = \diag(10, 1)$. The eigenvalues are $\lambda_1 = 10$ and $\lambda_2 = 1$.
The first principal component is:
\begin{equation}
\mathbf{w}_1 = \begin{bmatrix} 1 \\ 0 \end{bmatrix}
\end{equation}
PC1 aligns 100\% with feature 1.

\textbf{Step 2: Rescaled Covariance Matrix.}
Feature 2 is transformed by $x_2' = 10 x_2$. In matrix form, $X' = X \begin{bmatrix} 1 & 0 \\ 0 & 10 \end{bmatrix}$. The new covariance matrix is:
\begin{equation}
S' = \begin{bmatrix} 1 & 0 \\ 0 & 10 \end{bmatrix} \begin{bmatrix} 10 & 0 \\ 0 & 1 \end{bmatrix} \begin{bmatrix} 1 & 0 \\ 0 & 10 \end{bmatrix} = \begin{bmatrix} 10 & 0 \\ 0 & 100 \end{bmatrix}
\end{equation}

\textbf{Step 3: After Scaling.}
The eigenvalues of $S'$ are $\lambda_1' = 100$ and $\lambda_2' = 10$.
The first principal component is now:
\begin{equation}
\mathbf{w}_1' = \begin{bmatrix} 0 \\ 1 \end{bmatrix}
\end{equation}
PC1 now aligns 100\% with feature 2!

\textbf{Step 4: Pedagogical Takeaway.}
Merely converting kilometers to meters rotated the principal axis of the dataset by a full $90^\circ$!
This proves that \textbf{PCA is strictly not scale-invariant}. Features with large numerical scales will violently dominate the principal components, regardless of their actual informative value. To prevent artificial unit dominance, features must be standardized to unit variance ($\sigma_j = 1$), which is mathematically equivalent to performing PCA on the \textbf{correlation matrix} rather than the covariance matrix.
\end{examplebox}

\begin{examplebox}{The Dual High-Dimensional PCA Shortcut by Hand}
Consider a dataset with $N = 3$ observations and $d = 4$ features whose mean-centered data matrix is:
\begin{equation}
\tilde{X} = \begin{bmatrix} 1 & 1 & 0 & 0 \\ 0 & 1 & 1 & 0 \\ -1 & -2 & -1 & 0 \end{bmatrix} \in \mathbb{R}^{3 \times 4}
\end{equation}
Compute the dominant eigenvalue and 4-dimensional principal component $\mathbf{v}_1$ using the dual Gram shortcut.

\textbf{Step 1: Construct $3 \times 3$ Dual Gram Matrix $K$.}
Instead of computing the $4 \times 4$ matrix $\frac{1}{3}\tilde{X}^T \tilde{X}$, we compute $K = \frac{1}{3}\tilde{X}\tilde{X}^T \in \mathbb{R}^{3 \times 3}$:
\begin{equation}
\tilde{X}\tilde{X}^T = \begin{bmatrix} 1 & 1 & 0 & 0 \\ 0 & 1 & 1 & 0 \\ -1 & -2 & -1 & 0 \end{bmatrix} \begin{bmatrix} 1 & 0 & -1 \\ 1 & 1 & -2 \\ 0 & 1 & -1 \\ 0 & 0 & 0 \end{bmatrix} = \begin{bmatrix} 2 & 1 & -3 \\ 1 & 2 & -3 \\ -3 & -3 & 6 \end{bmatrix}
\end{equation}
Thus:
\begin{equation}
K = \frac{1}{3} \begin{bmatrix} 2 & 1 & -3 \\ 1 & 2 & -3 \\ -3 & -3 & 6 \end{bmatrix}
\end{equation}

\textbf{Step 2: Find Dominant Eigenpair of $K$.}
Notice the row sums are all zero (since the columns of $\tilde{X}$ are zero-mean), confirming $\lambda_3 = 0$ with eigenvector $\mathbf{u}_3 = \frac{1}{\sqrt{3}}[1, 1, 1]^T$.
Solving the characteristic equation for the non-zero roots:
\begin{equation}
\det(3K - \mu I) = \det\begin{bmatrix} 2 - \mu & 1 & -3 \\ 1 & 2 - \mu & -3 \\ -3 & -3 & 6 - \mu \end{bmatrix} = 0 \implies \mu^2(\mu - 9)(\mu - 1) = 0
\end{equation}
The eigenvalues of $\tilde{X}\tilde{X}^T$ are $\mu_1 = 9, \mu_2 = 1, \mu_3 = 0$.
Dividing by $N = 3$, the eigenvalues of $K$ are:
\begin{equation}
\lambda_1 = \frac{9}{3} = 3, \quad \lambda_2 = \frac{1}{3}
\end{equation}
For $\mu_1 = 9$:
\begin{equation}
(\tilde{X}\tilde{X}^T - 9I)\mathbf{u}_1 = \begin{bmatrix} -7 & 1 & -3 \\ 1 & -7 & -3 \\ -3 & -3 & -3 \end{bmatrix}\mathbf{u}_1 = \mathbf{0} \implies \mathbf{u}_1 = \frac{1}{\sqrt{6}}\begin{bmatrix} 1 \\ 1 \\ -2 \end{bmatrix}
\end{equation}

\textbf{Step 3: Reconstruct 4-Dimensional Principal Loading $\mathbf{v}_1$.}
Using the dual reconstruction formula $\mathbf{v}_1 = \frac{1}{\sqrt{N\lambda_1}} \tilde{X}^T \mathbf{u}_1$, where $\sqrt{N\lambda_1} = \sqrt{3 \cdot 3} = 3$:
\begin{equation}
\tilde{X}^T \mathbf{u}_1 = \begin{bmatrix} 1 & 0 & -1 \\ 1 & 1 & -2 \\ 0 & 1 & -1 \\ 0 & 0 & 0 \end{bmatrix} \left( \frac{1}{\sqrt{6}}\begin{bmatrix} 1 \\ 1 \\ -2 \end{bmatrix} \right) = \frac{1}{\sqrt{6}} \begin{bmatrix} 1(1) + 0 - 1(-2) \\ 1(1) + 1(1) - 2(-2) \\ 0(1) + 1(1) - 1(-2) \\ 0 \end{bmatrix} = \frac{1}{\sqrt{6}} \begin{bmatrix} 3 \\ 6 \\ 3 \\ 0 \end{bmatrix}
\end{equation}
Dividing by $3$:
\begin{equation}
\mathbf{v}_1 = \frac{1}{3} \left( \frac{1}{\sqrt{6}}\begin{bmatrix} 3 \\ 6 \\ 3 \\ 0 \end{bmatrix} \right) = \frac{1}{\sqrt{6}} \begin{bmatrix} 1 \\ 2 \\ 1 \\ 0 \end{bmatrix}
\end{equation}
Checking normalization: $\|\mathbf{v}_1\|_2^2 = \frac{1}{6}(1^2 + 2^2 + 1^2 + 0^2) = \frac{6}{6} = 1 \quad \checkmark$.
We have computed the 4-dimensional principal component strictly through a $3 \times 3$ matrix!
\end{examplebox}

\begin{ponderbox}{The Decorrelation Paradox: Uncorrelated vs. Independent}
In PCA, let $Z = \tilde{X} W \in \mathbb{R}^{N \times d}$ be the matrix of transformed principal component scores.
\begin{enumerate}
    \item Prove that the sample covariance matrix of the scores $S_Z = \frac{1}{N} Z^T Z$ is strictly diagonal:
    \begin{equation*}
    S_Z = \diag(\lambda_1, \dots, \lambda_d)
    \end{equation*}
    \item Does this diagonal structure prove that the features are \textbf{statistically independent}?
    \item Construct a simple 2D probability distribution where two features are strictly uncorrelated ($\Cov(z_1, z_2) = 0$), yet are completely dependent. What non-linear algorithm is required when true independence is needed?
\end{enumerate}
\tcblower
\textbf{Reflections and Mathematical Rationale:}
\begin{enumerate}
    \item Substituting $Z = \tilde{X} W$ and using $S = \frac{1}{N}\tilde{X}^T \tilde{X}$:
    \begin{equation*}
    S_Z = \frac{1}{N} Z^T Z = \frac{1}{N} (W^T \tilde{X}^T) (\tilde{X} W) = W^T \left(\frac{1}{N}\tilde{X}^T \tilde{X}\right) W = W^T S W
    \end{equation*}
    Since $W$ contains the orthonormal eigenvectors of $S$, $S W = W \Lambda$. Therefore:
    \begin{equation*}
    S_Z = W^T (W \Lambda) = (W^T W) \Lambda = I_d \Lambda = \diag(\lambda_1, \dots, \lambda_d)
    \end{equation*}
    \item \textbf{NO!} Diagonal covariance strictly guarantees that the features are \textbf{uncorrelated} (zero linear correlation: $\Cov(z_i, z_j) = 0$). Uncorrelatedness is merely a second-order linear property. Statistical independence requires the joint probability density to factorize completely across all higher moments:
    \begin{equation*}
    P(z_1, z_2) = P(z_1) P(z_2)
    \end{equation*}
    Uncorrelatedness implies independence \emph{only} when data follows a multivariate Gaussian distribution.
    \item \textbf{Counterexample:} Let $z_1 \sim \text{Uniform}[-1, 1]$, and let $z_2 = z_1^2$.
    Here, $z_2$ is fully determined by $z_1$; they are 100\% dependent.
    However, by symmetry:
    \begin{equation*}
    \E[z_1] = 0, \quad \E[z_1 z_2] = \E[z_1^3] = 0 \implies \Cov(z_1, z_2) = \E[z_1 z_2] - \E[z_1]\E[z_2] = 0 - 0 = 0
    \end{equation*}
    Linear PCA is completely blind to this non-linear dependency. When true statistical independence across higher moments is required, one must employ \textbf{Independent Component Analysis (ICA)} or non-linear Kernel PCA.
\end{enumerate}
\end{ponderbox}

\begin{figure}[htbp]
\centering
\begin{tikzpicture}[scale=0.95, font=\sffamily\footnotesize]

% Class 0 (Blue/Teal Ribbon)
\draw[fill=tealbg!80, draw=tealborder, rounded corners=4pt] (-3.5, 0.4) rectangle (3.5, 1.4);
\node[color=tealdef, font=\bfseries] at (-2.4, 0.9) {Class $Y = 0$};

% Class 1 (Red/Coral Ribbon)
\draw[fill=trapbg!80, draw=traprad, rounded corners=4pt] (-3.5, -1.4) rectangle (3.5, -0.4);
\node[color=traprad, font=\bfseries] at (-2.4, -0.9) {Class $Y = 1$};

% PCA Projection Axis (Horizontal)
\draw[->, >=stealth, line width=2.4pt, color=tealborder] (-4.2, 0) -- (4.2, 0) node[right] {\textbf{PCA Axis } $\mathbf{w}_{\text{PCA}}$};
\node[below, color=tealdef, align=center] at (2.2, -0.1) {Captures Maximum Variance \\ \textcolor{traprad}{\textbf{Destroys Class Separation!}}};

% LDA Projection Axis (Vertical)
\draw[->, >=stealth, line width=2.4pt, color=primarydark] (0, -2.2) -- (0, 2.2) node[above] {\textbf{LDA Axis } $\mathbf{w}_{\text{LDA}}$};
\node[right, color=primarydark, align=left] at (0.2, 1.8) {Low Variance, \\ \textbf{Optimal Class Separation}};

\end{tikzpicture}
\caption{The Supervised Failure Mode: PCA vs. Linear Discriminant Analysis (LDA). Because PCA is strictly unsupervised, it projects along the horizontal axis of maximum variance, collapsing the two distinct classes into an overlapping mixture. Supervised LDA identifies the vertical axis, preserving perfect class separability.}
\label{fig:pca_vs_lda}
\end{figure}

% ----------------------------------------------------------------------
% PART 6: EXAM TRAPS & DIAGNOSTIC MISCONCEPTIONS
% ----------------------------------------------------------------------
\section{Exam Traps \& Diagnostic Misconceptions}

\begin{trapbox}{Gotcha 1: Omitting Data Centering Before SVD/Covariance}
An analyst feeds raw uncentered data $X$ directly into SVD or computes $\frac{1}{N}X^T X$.
\textbf{The Trap:} If data is not mean-centered, the first principal component simply captures the displacement vector from the origin to the global center of mass $\boldsymbol{\mu}$, rather than the internal axis of maximum variation. Always verify that $\frac{1}{N}\sum_{i=1}^N \tilde{\mathbf{x}}_i = \mathbf{0}$ before performing PCA.
\end{trapbox}

\begin{trapbox}{Gotcha 2: Failure to Standardize Features with Incommensurate Units}
A data scientist runs PCA on a dataset containing age (range $20$--$80$) and annual salary in rupees (range $3{,}00{,}000$--$50{,}00{,}000$).
\textbf{The Trap:} PCA is strictly scale-dependent. Because salary has a numerical variance billions of times larger than age, the first principal component will align $99.999\%$ with salary, completely ignoring age. When features have different units, features must be standardized to zero mean and unit variance: $\mathbf{x}_{\text{std}} = \frac{\mathbf{x} - \boldsymbol{\mu}}{\boldsymbol{\sigma}}$ (performing PCA on the correlation matrix instead of covariance).
\end{trapbox}

\begin{trapbox}{Gotcha 3: Assuming PCA is Optimal for Supervised Classification}
A practitioner uses PCA to reduce 100 features to 2 dimensions before training a classifier, but discovers that classification accuracy collapses.
\textbf{The Trap:} PCA is strictly \textbf{unsupervised}. It identifies directions of maximum variance within $\mathcal{X}$, with zero knowledge of target labels $\mathcal{Y}$. If the class separation lies along a low-variance direction (e.g., two parallel dense ribbons as shown in Figure~\ref{fig:pca_vs_lda}), PCA will discard the discriminative direction as "noise." In supervised settings, Linear Discriminant Analysis (LDA) should be used.
\end{trapbox}

\begin{trapbox}{Gotcha 4: Data Leakage by Fitting PCA on the Full Dataset}
A programmer applies PCA to the full combined dataset of $10{,}000$ points before performing an $80/20$ train-test split.
\textbf{The Trap:} The principal directions $W$ and sample mean $\boldsymbol{\mu}$ encode information about the global test distribution. Fitting PCA globally induces data leakage, producing over-optimistic validation metrics. PCA must be fitted strictly on the training set: $\boldsymbol{\mu}_{\text{train}}$ and $W_{\text{train}}$ are then used to project the test partition out-of-sample: $Z_{\text{test}} = (X_{\text{test}} - \mathbf{1}\boldsymbol{\mu}_{\text{train}}^T) W_{\text{train}}$.
\end{trapbox}

\begin{trapbox}{Gotcha 5: Conflating PCA with Factor Analysis}
A student treats PCA and Factor Analysis (FA) as identical algorithms with different names.
\textbf{The Trap:} PCA is a deterministic orthogonal transformation that decomposes \textbf{total observed variance}. Factor Analysis is a probabilistic generative model that separates shared variance (latent common factors) from feature-specific independent noise:
\begin{equation*}
\Sigma = \Lambda \Lambda^T + \Psi
\end{equation*}
where $\Psi$ is a diagonal uniqueness matrix. PCA makes zero distinction between common variance and idiosyncratic noise, whereas Factor Analysis explicitly estimates and subtracts the noise floor $\Psi$.
\end{trapbox}

% ----------------------------------------------------------------------
% PART 7: COMPARATIVE SYNTHESIS TABLE
% ----------------------------------------------------------------------
\section{Comparative Synthesis Table}

Table~\ref{tab:dim_reduction_comp} compares principal dimensionality reduction paradigms across machine learning.

\begin{table}[htbp]
\centering
\small
\begin{tabularx}{\textwidth}{lXXXXX}
\toprule
\textbf{Method} & \textbf{Paradigm} & \textbf{Linearity} & \textbf{Optimization Objective} & \textbf{Computational Cost} & \textbf{Primary Use Case} \\
\midrule
PCA & Unsupervised & Linear & Maximize variance / minimize reconstruction error & $\mathcal{O}(\min(N d^2, N^2 d))$ & Compression, denoising, feature decorrelation \\
\addlinespace
Kernel PCA & Unsupervised & Non-linear & Maximize variance in Reproducing Kernel Hilbert Space & $\mathcal{O}(N^3)$ & Non-linear manifold unrolling (e.g., Swiss roll) \\
\addlinespace
LDA & Supervised & Linear & Maximize between-class variance / minimize within-class & $\mathcal{O}(d^3 + N d^2)$ & Class-separable feature projection for classification \\
\addlinespace
t-SNE & Unsupervised & Non-linear & Minimize KL divergence of pairwise Gaussian/Student-t similarities & $\mathcal{O}(N^2)$ or $\mathcal{O}(N \log N)$ & 2D/3D cluster visualization (no out-of-sample mapping) \\
\addlinespace
UMAP & Unsupervised & Non-linear & Preserve local fuzzy simplicial set geometry & $\mathcal{O}(N \log N)$ & Fast high-dimensional non-linear visualization \\
\bottomrule
\end{tabularx}
\caption{Comparative Synthesis: Dimensionality Reduction Techniques.}
\label{tab:dim_reduction_comp}
\end{table}

% ----------------------------------------------------------------------
% PART 8: PYTHON / SCIENTIFIC COMPUTING LAB
% ----------------------------------------------------------------------
\section{Python / Scientific Computing Lab}

The production script below implements PCA from scratch via covariance eigendecomposition, direct SVD, and the dual Gram matrix shortcut for $d \gg N$.

\begin{lstlisting}[language=Python, caption={PCA Implementations and Dual High-Dimensional Shortcut in Python.}]
import numpy as np

class RobustPCA:
    """
    Production-grade PCA implementation supporting:
    1. Standard direct SVD for d <= N.
    2. Dual Gram Matrix Shortcut for d >> N (O(N^3) complexity).
    3. Out-of-sample projection without data leakage.
    """
    def __init__(self, n_components: int):
        self.k = n_components
        self.mean_ = None
        self.components_ = None
        self.explained_variance_ratio_ = None

    def fit(self, X: np.ndarray):
        N, d = X.shape
        self.mean_ = np.mean(X, axis=0)
        X_centered = X - self.mean_

        if d <= N:
            # Standard direct SVD on centered data
            U, s, Vt = np.linalg.svd(X_centered, full_matrices=False)
            self.components_ = Vt[:self.k]
            eigenvalues = (s ** 2) / N
        else:
            # Dual High-Dimensional Shortcut for d >> N
            # Gram matrix K = (1/N) * X_c @ X_c^T: shape (N, N)
            K = (1.0 / N) * (X_centered @ X_centered.T)
            eigvals, U = np.linalg.eigh(K)
            # Sort descending
            idx = np.argsort(eigvals)[::-1]
            eigvals = eigvals[idx]
            U = U[:, idx]

            # Reconstruct V = (1 / sqrt(N * lambda)) * X_c^T @ u
            V_list = []
            for i in range(self.k):
                lam_i = max(eigvals[i], 1e-12)
                u_i = U[:, i]
                v_i = (1.0 / np.sqrt(N * lam_i)) * (X_centered.T @ u_i)
                # Normalize to ensure unit norm
                v_i = v_i / np.linalg.norm(v_i)
                V_list.append(v_i)
            self.components_ = np.array(V_list)
            eigenvalues = eigvals

        total_variance = np.sum(eigenvalues)
        self.explained_variance_ratio_ = eigenvalues[:self.k] / total_variance
        return self

    def transform(self, X: np.ndarray) -> np.ndarray:
        """Projects out-of-sample data onto learned principal components."""
        return (X - self.mean_) @ self.components_.T

    def inverse_transform(self, Z: np.ndarray) -> np.ndarray:
        """Reconstructs low-dimensional scores back to original feature space."""
        return Z @ self.components_ + self.mean_

def main():
    rng = np.random.RandomState(42)
    # Test on high-dimensional setting: N = 50 samples, d = 500 features (d >> N)
    N, d = 50, 500
    X_train = rng.randn(N, d) + 3.0
    X_test = rng.randn(10, d) + 3.0

    pca = RobustPCA(n_components=5)
    pca.fit(X_train)
    Z_train = pca.transform(X_train)
    Z_test = pca.transform(X_test)
    X_recon = pca.inverse_transform(Z_test)

    print("=== HIGH-DIMENSIONAL DUAL PCA (d >> N) ===")
    print(f"Original Feature Dimension: {d}, Reduced Dimension: {pca.k}")
    print(f"Transformed Training Shape: {Z_train.shape}")
    print(f"Transformed Test Shape:     {Z_test.shape}")
    print(f"Explained Variance per Component: {np.round(pca.explained_variance_ratio_, 4)}")
    print(f"Cumulative Variance Captured:     {np.sum(pca.explained_variance_ratio_):.4f}")
    print(f"Test Reconstruction MSE:          {np.mean((X_test - X_recon)**2):.4f}")

if __name__ == "__main__":
    main()
\end{lstlisting}

% ----------------------------------------------------------------------
% PART 9: COMPREHENSIVE PRACTICE PROBLEMS
% ----------------------------------------------------------------------
\section{Comprehensive Practice Problems}

\subsection*{Part A: Analytical \& Mathematical Problems}

\begin{enumerate}[label=\textbf{P\arabic*.}]
    \item \textbf{Analytical PCA on a Small $2 \times 2$ Covariance Matrix:}
    Suppose a 2-dimensional zero-mean dataset yields sample covariance matrix $S = \begin{bmatrix} 5 & 2 \\ 2 & 2 \end{bmatrix}$.
    \begin{enumerate}[label=(\alph*)]
        \item Compute the eigenvalues $\lambda_1, \lambda_2$ of $S$.
        \item Find the normalized eigenvectors $\mathbf{w}_1, \mathbf{w}_2$.
        \item Compute the proportion of total variance explained by the first principal component.
        \item For a centered observation $\tilde{\mathbf{x}} = [2, 1]^T$, compute its scalar score $z_1 = \mathbf{w}_1^T \tilde{\mathbf{x}}$ and reconstructed point $\hat{\mathbf{x}} = z_1 \mathbf{w}_1$.
    \end{enumerate}

    \item \textbf{Scale Invariance Counterexample:}
    Let $X \in \mathbb{R}^{N \times 2}$ have covariance matrix $S = \begin{bmatrix} 10 & 0 \\ 0 & 1 \end{bmatrix}$.
    Suppose the second feature is multiplied by scale factor $c = 10$, producing scaled data matrix $X'$.
    \begin{enumerate}[label=(\alph*)]
        \item Compute the new covariance matrix $S'$.
        \item Find the first principal component $\mathbf{w}_1$ before and after scaling.
        \item Conclude why standardizing features to unit variance is essential when features represent different physical units.
    \end{enumerate}

    \item \textbf{Orthogonality of Principal Component Scores:}
    Let $Z = \tilde{X} W \in \mathbb{R}^{N \times k}$ be the matrix of principal component scores, where $W = [\mathbf{w}_1, \dots, \mathbf{w}_k]$ contains the top $k$ eigenvectors of $S = \frac{1}{N}\tilde{X}^T \tilde{X}$.
    Prove that the sample covariance matrix of the scores $S_Z = \frac{1}{N} Z^T Z$ is strictly diagonal:
    \begin{equation}
    S_Z = \diag(\lambda_1, \lambda_2, \dots, \lambda_k)
    \end{equation}
    What does this prove regarding the correlation between distinct principal components?

    \item \textbf{Reconstruction Error Bounds:}
    Let $\lambda_1 \ge \lambda_2 \ge \dots \ge \lambda_d$ be the eigenvalues of $S$.
    Prove that the minimum mean squared reconstruction error for a $k$-dimensional PCA approximation satisfies:
    \begin{equation}
    \frac{1}{N} \sum_{i=1}^N \|\tilde{\mathbf{x}}_i - W_k W_k^T \tilde{\mathbf{x}}_i\|_2^2 = \sum_{j=k+1}^d \lambda_j
    \end{equation}

    \item \textbf{High-Dimensional Dual PCA Rank Bound:}
    Let $\tilde{X} \in \mathbb{R}^{N \times d}$ with $N < d$. Prove that the covariance matrix $S = \frac{1}{N}\tilde{X}^T \tilde{X}$ has at most $N - 1$ strictly positive eigenvalues, and deduce that at least $d - N + 1$ eigenvalues must be identically zero.
\end{enumerate}

\subsection*{Part B: Diagnostic Multiple-Choice Questions}

\begin{enumerate}[label=\textbf{MCQ\arabic*.}]
    \item A covariance matrix $S \in \mathbb{R}^{4 \times 4}$ has eigenvalues $\lambda_1 = 12, \lambda_2 = 5, \lambda_3 = 2, \lambda_4 = 1$. What proportion of the total dataset variance is explained by retaining the top 2 principal components?
    \begin{enumerate}[label=(\alph*)]
        \item $\frac{12 + 5}{12 + 5 + 2 + 1} = \frac{17}{20} = 85\%$
        \item $\frac{12^2 + 5^2}{12^2 + 5^2 + 2^2 + 1^2} = \frac{169}{174} \approx 97.1\%$
        \item $\frac{12}{20} = 60\%$
        \item $\frac{\sqrt{17}}{\sqrt{20}} \approx 92.2\%$
    \end{enumerate}

    \item In what scenario can Principal Component Analysis (PCA) directly hurt downstream classification accuracy when used for feature reduction?
    \begin{enumerate}[label=(\alph*)]
        \item When the features are completely uncorrelated.
        \item When the direction that best separates the target classes has low variance compared to background noise directions.
        \item When the number of observations $N$ exceeds the feature dimension $d$.
        \item When the data is centered to have zero mean.
    \end{enumerate}
\end{enumerate}

\chapter[Unconstrained Optimization \& Gradient Descent]{Unconstrained Optimization, Taylor Expansions and Gradient Descent Convergence}
\label{chap:week08}

% ----------------------------------------------------------------------
% PART 1: EXECUTIVE ROADMAP & THE NARROW RAVINE PARADOX
% ----------------------------------------------------------------------
\section{Executive Roadmap \& Learning Objectives}

The formulation of a machine learning problem as empirical risk minimization culminates in a mathematical optimization problem: finding the parameter vector $\mathbf{w}^* \in \mathbb{R}^d$ that minimizes an unconstrained objective functional $f(\mathbf{w})$. While closed-form solutions exist for Ordinary Least Squares (Chapter~\ref{chap:week03}), modern machine learning models---including logistic regression, support vector machines, and deep neural networks---lack closed-form stationary points. They necessitate iterative numerical optimization. This chapter develops the foundational mathematical theory governing gradient-based and curvature-based algorithms.

\begin{figure}[htbp]
\centering
\begin{tcolorbox}[
    colback=coralbg,
    colframe=coralborder,
    title={\large \textbf{The Narrow Ravine Paradox: Why Steepest Descent Fails to Head for the Minimum}},
    fonttitle=\bfseries\sffamily,
    arc=3mm,
    boxrule=1pt,
    width=\textwidth
]
\textbf{The Intuitive Belief:} Gradient descent follows the direction of \emph{steepest descent} ($-\nabla f(\mathbf{w})$). Intuitively, if you take steps in the direction where the loss drops most steeply, each step must point directly toward the bottom of the bowl (the global minimizer $\mathbf{w}^*$).

\textbf{The Mathematical Reality:} The gradient $\nabla f(\mathbf{w})$ is defined strictly by \emph{infinitesimal local geometry}: it is mathematically orthogonal to the tangent plane of the level set (contour line) passing through $\mathbf{w}$. Unless the loss surface has perfectly spherical contours (equal curvature in all directions), the negative gradient \textbf{almost never points toward the minimum}!

\textbf{The Catastrophe in a Narrow Ravine:} Consider an elongated canyon or ravine where the walls are steep ($L = 100$) but the valley floor slopes gently toward the minimum ($\mu = 1$). If an algorithm is positioned on the canyon wall, the steepest direction does not point along the canyon floor toward the exit; it points directly across the ravine into the opposite wall. As a result, standard gradient descent spends $99.9\%$ of its energy violently ping-ponging back and forth between the cliffs, creeping along the valley floor at a glacial pace of $\mathcal{O}(\kappa)$ steps, where $\kappa = L/\mu$ is the condition number.
\end{tcolorbox}
\caption{The Narrow Ravine Paradox: The fundamental pathology of first-order optimization.}
\label{fig:ravine_paradox}
\end{figure}

\subsection{Prerequisites \& Scaffolding}
Mastery of this chapter builds upon vector calculus, matrix norms, and quadratic form curvature as mapped in Table~\ref{tab:w08_prereqs}.

\begin{table}[htbp]
\centering
\small
\begin{tabularx}{\textwidth}{lXX}
\toprule
\textbf{Foundational Area} & \textbf{Core Concepts Required} & \textbf{Role in Machine Learning Foundations} \\
\midrule
Vector Calculus & Gradient vectors $\nabla f(\mathbf{w})$, Hessian matrices $\nabla^2 f(\mathbf{w})$, Taylor series & Determining local descent directions and second-order curvature \\
Linear Algebra & Spectral radius $\rho(A)$, matrix powers $A^k$, condition number $\kappa$ & Analyzing convergence rates and stability bounds of update operators \\
Mathematical Analysis & Lipschitz continuity, Cauchy sequences, contraction mappings & Establishing formal convergence guarantees for iterative algorithms \\
\bottomrule
\end{tabularx}
\caption{Prerequisite Scaffolding for Chapter 8.}
\label{tab:w08_prereqs}
\end{table}

\subsection{Core Competencies}
By the conclusion of this chapter, students will be able to:
\begin{itemize}
    \item Rigorously state and verify First- and Second-Order Necessary and Sufficient Optimality Conditions.
    \item Formulate $L$-Lipschitz smoothness and prove the fundamental Descent Lemma via integral calculus.
    \item Prove the sublinear $\mathcal{O}(1/k)$ convergence rate of Gradient Descent for convex $L$-smooth functions.
    \item Prove the linear $\mathcal{O}\left(\left(\frac{\kappa - 1}{\kappa + 1}\right)^k\right)$ convergence rate for $\mu$-strongly convex quadratics and determine the optimal step size.
    \item Formulate Newton-Raphson optimization, prove its affine invariance, and analyze the $\mathcal{O}(d^3)$ computational barrier.
    \item Analyze momentum acceleration (Polyak heavy-ball) via discrete dynamical systems.
\end{itemize}

\subsection{Notation \& Dimension Contracts}
Throughout this chapter, we adhere to the standard dimension contracts:
\begin{itemize}
    \item Parameter vector: $\mathbf{w} \in \mathbb{R}^d$ (column vector).
    \item Objective function: $f: \mathbb{R}^d \to \mathbb{R}$.
    \item Gradient vector: $\nabla f(\mathbf{w}) = \left[ \frac{\partial f}{\partial w_1}, \dots, \frac{\partial f}{\partial w_d} \right]^T \in \mathbb{R}^{d \times 1}$.
    \item Hessian matrix: $\nabla^2 f(\mathbf{w}) \in \mathbb{R}^{d \times d}$ where $[\nabla^2 f(\mathbf{w})]_{ij} = \frac{\partial^2 f}{\partial w_i \partial w_j}$.
    \item Learning rate (step size): $\eta > 0$.
    \item Lipschitz smoothness constant: $L = \sup_{\mathbf{w}} \|\nabla^2 f(\mathbf{w})\|_2 = \lambda_{\max}(\nabla^2 f)$.
    \item Strong convexity parameter: $\mu = \inf_{\mathbf{w}} \lambda_{\min}(\nabla^2 f(\mathbf{w})) > 0$.
    \item Condition number: $\kappa \triangleq \frac{L}{\mu} \ge 1$.
\end{itemize}

% ----------------------------------------------------------------------
% PART 2: MATHEMATICAL & COMPUTATIONAL FOUNDATIONS
% ----------------------------------------------------------------------
\section{Mathematical \& Computational Foundations}

\subsection{Multivariable Taylor Approximations}

Let $f: \mathbb{R}^d \to \mathbb{R}$ be twice continuously differentiable ($\mathcal{C}^2$). For any displacement vector $\mathbf{v} \in \mathbb{R}^d$, the scalar function $g(t) = f(\mathbf{w} + t\mathbf{v})$ for $t \in [0, 1]$ allows us to transfer one-dimensional Taylor calculus to $\mathbb{R}^d$:
\begin{align}
g'(t) &= \nabla f(\mathbf{w} + t\mathbf{v})^T \mathbf{v} \\
g''(t) &= \mathbf{v}^T \nabla^2 f(\mathbf{w} + t\mathbf{v}) \mathbf{v}
\end{align}
By Taylor's Theorem with Lagrange remainder on $g(t)$:
\begin{equation}
f(\mathbf{w} + \mathbf{v}) = f(\mathbf{w}) + \nabla f(\mathbf{w})^T \mathbf{v} + \frac{1}{2}\mathbf{v}^T \nabla^2 f(\mathbf{w} + \xi \mathbf{v}) \mathbf{v} \quad \text{for some } \xi \in (0, 1)
\label{eq:taylor_expansion}
\end{equation}
The linear approximation models the function as a hyper-plane, while the quadratic approximation captures local parabolic curvature.

\subsection{Optimality Conditions for Unconstrained Optimization}

\begin{theorem}[First-Order Necessary Condition (FONC)]
Let $f: \mathbb{R}^d \to \mathbb{R}$ be continuously differentiable. If $\mathbf{w}^*$ is a local minimum of $f$, then:
\begin{equation}
\nabla f(\mathbf{w}^*) = \mathbf{0}
\label{eq:fonc}
\end{equation}
A point satisfying $\nabla f(\mathbf{w}) = \mathbf{0}$ is called a \textbf{stationary (critical) point}.
\end{theorem}

\begin{proof}
Suppose $\nabla f(\mathbf{w}^*) \ne \mathbf{0}$. Choose direction $\mathbf{v} = -\nabla f(\mathbf{w}^*)$. By Taylor expansion:
\begin{equation}
f(\mathbf{w}^* + t\mathbf{v}) = f(\mathbf{w}^*) + t \nabla f(\mathbf{w}^*)^T (-\nabla f(\mathbf{w}^*)) + o(t) = f(\mathbf{w}^*) - t \|\nabla f(\mathbf{w}^*)\|_2^2 + o(t)
\end{equation}
For sufficiently small $t > 0$, $-t \|\nabla f(\mathbf{w}^*)\|_2^2 + o(t) < 0$, which implies $f(\mathbf{w}^* + t\mathbf{v}) < f(\mathbf{w}^*)$, directly contradicting that $\mathbf{w}^*$ is a local minimum. Hence $\nabla f(\mathbf{w}^*) = \mathbf{0}$.
\end{proof}

\begin{theorem}[Second-Order Optimality Conditions]
Let $f: \mathbb{R}^d \to \mathbb{R}$ be twice continuously differentiable:
\begin{enumerate}
    \item \textbf{Second-Order Necessary Condition (SONC):} If $\mathbf{w}^*$ is a local minimum, then $\nabla f(\mathbf{w}^*) = \mathbf{0}$ and the Hessian is positive semidefinite:
    \begin{equation}
    \nabla^2 f(\mathbf{w}^*) \succeq 0 \quad (\forall \mathbf{v} \in \mathbb{R}^d, \; \mathbf{v}^T \nabla^2 f(\mathbf{w}^*) \mathbf{v} \ge 0)
    \end{equation}
    \item \textbf{Second-Order Sufficient Condition (SOSC):} If $\nabla f(\mathbf{w}^*) = \mathbf{0}$ and the Hessian is strictly positive definite:
    \begin{equation}
    \nabla^2 f(\mathbf{w}^*) \succ 0 \quad (\forall \mathbf{v} \ne \mathbf{0}, \; \mathbf{v}^T \nabla^2 f(\mathbf{w}^*) \mathbf{v} > 0)
    \end{equation}
    then $\mathbf{w}^*$ is a strict (isolated) local minimum.
\end{enumerate}
\end{theorem}

\begin{ponderbox}[Stationary Points vs.\ Local Minima]
If $\nabla f(\mathbf{w}^*) = \mathbf{0}$ and $\nabla^2 f(\mathbf{w}^*) \succeq 0$, does this guarantee that $\mathbf{w}^*$ is a local minimum?

\textbf{Reflection:} No! Consider the scalar function $f(w) = w^3$. At $w^* = 0$, $f'(0) = 0$ and $f''(0) = 0 \ge 0$. The necessary conditions are fully satisfied. Yet $w^* = 0$ is an inflection point (saddle point), neither a minimum nor a maximum ($f(\epsilon) = \epsilon^3 > 0$ and $f(-\epsilon) = -\epsilon^3 < 0$). Strict positive definiteness ($\nabla^2 f \succ 0$) is required for sufficiency.
\end{ponderbox}

\subsection{Lipschitz Smoothness and the Descent Lemma}

\begin{definitionbox}{$L$-Lipschitz Smoothness}
A continuously differentiable function $f: \mathbb{R}^d \to \mathbb{R}$ is \textbf{$L$-smooth} (its gradient is $L$-Lipschitz continuous) if there exists a constant $L > 0$ such that for all $\mathbf{u}, \mathbf{v} \in \mathbb{R}^d$:
\begin{equation}
\|\nabla f(\mathbf{u}) - \nabla f(\mathbf{v})\|_2 \le L \|\mathbf{u} - \mathbf{v}\|_2
\label{eq:lipschitz_grad}
\end{equation}
For twice-differentiable functions, this condition is mathematically equivalent to bounding the spectral norm of the Hessian everywhere:
\begin{equation}
\|\nabla^2 f(\mathbf{w})\|_2 = \sigma_{\max}(\nabla^2 f(\mathbf{w})) \le L \quad \forall \mathbf{w} \in \mathbb{R}^d
\end{equation}
Geometrically, $L$-smoothness means the curvature of the loss surface cannot exceed $L$ in any direction.
\end{definitionbox}

\begin{theorem}[The Descent Lemma / Quadratic Upper Bound]
If $f$ is $L$-smooth, then for all $\mathbf{u}, \mathbf{v} \in \mathbb{R}^d$:
\begin{equation}
f(\mathbf{v}) \le f(\mathbf{u}) + \nabla f(\mathbf{u})^T (\mathbf{v} - \mathbf{u}) + \frac{L}{2} \|\mathbf{v} - \mathbf{u}\|_2^2
\label{eq:descent_lemma}
\end{equation}
\end{theorem}

\begin{proof}
By the Fundamental Theorem of Calculus along the segment $\mathbf{w}(t) = \mathbf{u} + t(\mathbf{v} - \mathbf{u})$:
\begin{equation}
f(\mathbf{v}) - f(\mathbf{u}) = \int_0^1 \nabla f(\mathbf{u} + t(\mathbf{v} - \mathbf{u}))^T (\mathbf{v} - \mathbf{u}) \, dt
\end{equation}
Add and subtract $\nabla f(\mathbf{u})^T (\mathbf{v} - \mathbf{u})$ inside the integral:
\begin{equation}
f(\mathbf{v}) - f(\mathbf{u}) = \nabla f(\mathbf{u})^T (\mathbf{v} - \mathbf{u}) + \int_0^1 \left( \nabla f(\mathbf{u} + t(\mathbf{v} - \mathbf{u})) - \nabla f(\mathbf{u}) \right)^T (\mathbf{v} - \mathbf{u}) \, dt
\end{equation}
Applying the Cauchy-Schwarz inequality and the $L$-Lipschitz gradient property:
\begin{align}
\left| \int_0^1 \left( \nabla f(\mathbf{u} + t(\mathbf{v} - \mathbf{u})) - \nabla f(\mathbf{u}) \right)^T (\mathbf{v} - \mathbf{u}) \, dt \right|
&\le \int_0^1 \|\nabla f(\mathbf{u} + t(\mathbf{v} - \mathbf{u})) - \nabla f(\mathbf{u})\|_2 \|\mathbf{v} - \mathbf{u}\|_2 \, dt \nonumber \\
&\le \int_0^1 L t \|\mathbf{v} - \mathbf{u}\|_2^2 \, dt = \frac{L}{2} \|\mathbf{v} - \mathbf{u}\|_2^2
\end{align}
Substituting this back establishes the global quadratic upper bound~\eqref{eq:descent_lemma}.
\end{proof}

\begin{theorem}[Sufficient Decrease per Iteration]
Let $f$ be $L$-smooth. Under the Gradient Descent update $\mathbf{w}_{k+1} = \mathbf{w}_k - \eta \nabla f(\mathbf{w}_k)$ with constant step size $\eta \in \left(0, \frac{2}{L}\right)$:
\begin{equation}
f(\mathbf{w}_{k+1}) \le f(\mathbf{w}_k) - \eta \left( 1 - \frac{L\eta}{2} \right) \|\nabla f(\mathbf{w}_k)\|_2^2
\label{eq:sufficient_decrease}
\end{equation}
In particular, for the canonical step size $\eta = \frac{1}{L}$:
\begin{equation}
f(\mathbf{w}_{k+1}) \le f(\mathbf{w}_k) - \frac{1}{2L} \|\nabla f(\mathbf{w}_k)\|_2^2
\label{eq:canonical_decrease}
\end{equation}
\end{theorem}

\begin{proof}
Substitute $\mathbf{u} = \mathbf{w}_k$ and $\mathbf{v} = \mathbf{w}_{k+1} = \mathbf{w}_k - \eta \nabla f(\mathbf{w}_k)$ into the Descent Lemma:
\begin{align}
f(\mathbf{w}_{k+1}) &\le f(\mathbf{w}_k) + \nabla f(\mathbf{w}_k)^T (-\eta \nabla f(\mathbf{w}_k)) + \frac{L}{2} \|-\eta \nabla f(\mathbf{w}_k)\|_2^2 \nonumber \\
&= f(\mathbf{w}_k) - \eta \|\nabla f(\mathbf{w}_k)\|_2^2 + \frac{L\eta^2}{2} \|\nabla f(\mathbf{w}_k)\|_2^2 \nonumber \\
&= f(\mathbf{w}_k) - \eta \left( 1 - \frac{L\eta}{2} \right) \|\nabla f(\mathbf{w}_k)\|_2^2
\end{align}
The coefficient $\eta\left(1 - \frac{L\eta}{2}\right) > 0$ if and only if $0 < \eta < \frac{2}{L}$. Setting the derivative with respect to $\eta$ to zero:
\begin{equation}
\frac{d}{d\eta} \left( \eta - \frac{L\eta^2}{2} \right) = 1 - L\eta = 0 \implies \eta^*_{\text{descent}} = \frac{1}{L}
\end{equation}
Thus, $\eta = \frac{1}{L}$ maximizes the guaranteed one-step decrease: $\Delta f \ge \frac{1}{2L}\|\nabla f(\mathbf{w}_k)\|_2^2$.
\end{proof}

\begin{ponderbox}[The Two Step Sizes: $1/L$ vs.\ $2/(L+\mu)$]
Why do textbooks emphasize two different "optimal" step sizes for Gradient Descent? When should you use $\eta = 1/L$, and when should you use $\eta = \frac{2}{L+\mu}$?

\textbf{Reflection:}
\begin{itemize}
    \item $\eta = 1/L$ is the \textbf{conservative descent step size}. It guarantees that the objective function strictly decreases on \emph{every single step} ($f(\mathbf{w}_{k+1}) \le f(\mathbf{w}_k)$) for \emph{any} $L$-smooth function (even non-convex ones!).
    \item $\eta = \frac{2}{L+\mu}$ is the \textbf{asymptotic contraction step size}. It is designed specifically for \emph{strongly convex quadratics} to minimize the spectral radius of the error transition operator. It yields the fastest possible asymptotic convergence rate $\frac{\kappa-1}{\kappa+1}$, but early steps along the high-curvature axis can overshoot the valley floor (though the overall error in parameter space contractively shrinks).
\end{itemize}
\end{ponderbox}

\subsection{Convergence Rates: Convex vs.\ Strongly Convex}

\begin{theorem}[$\mathcal{O}(1/k)$ Sublinear Convergence for Convex $L$-Smooth Functions]
Let $f: \mathbb{R}^d \to \mathbb{R}$ be convex and $L$-smooth. With constant step size $\eta = \frac{1}{L}$, after $T$ iterations of Gradient Descent:
\begin{equation}
f(\mathbf{w}_T) - f^* \le \frac{2L \|\mathbf{w}_0 - \mathbf{w}^*\|_2^2}{T} = \mathcal{O}\left(\frac{1}{T}\right)
\label{eq:convex_rate}
\end{equation}
To achieve an $\epsilon$-accurate solution ($f(\mathbf{w}_T) - f^* \le \epsilon$), Gradient Descent requires $\mathcal{O}(1/\epsilon)$ iterations.
\end{theorem}

\begin{definitionbox}{Strong Convexity}
A continuously differentiable function $f: \mathbb{R}^d \to \mathbb{R}$ is \textbf{$\mu$-strongly convex} ($\mu > 0$) if for all $\mathbf{u}, \mathbf{v} \in \mathbb{R}^d$:
\begin{equation}
f(\mathbf{v}) \ge f(\mathbf{u}) + \nabla f(\mathbf{u})^T (\mathbf{v} - \mathbf{u}) + \frac{\mu}{2} \|\mathbf{v} - \mathbf{u}\|_2^2
\label{eq:strong_convexity}
\end{equation}
For twice-differentiable functions, this is equivalent to $\nabla^2 f(\mathbf{w}) \succeq \mu I$ for all $\mathbf{w}$. Strong convexity provides a quadratic \emph{lower} bound, guaranteeing that the objective grows quadratically away from the minimum $\mathbf{w}^*$.
\end{definitionbox}

Combining $L$-smoothness (quadratic upper bound) and $\mu$-strong convexity (quadratic lower bound), we obtain the sandwich inequality for all $\mathbf{w}$:
\begin{equation}
\frac{\mu}{2} \|\mathbf{w} - \mathbf{w}^*\|_2^2 \le f(\mathbf{w}) - f^* \le \frac{L}{2} \|\mathbf{w} - \mathbf{w}^*\|_2^2
\label{eq:sandwich_bound}
\end{equation}

% ----------------------------------------------------------------------
% PART 3: VISUAL & GEOMETRIC INTUITION
% ----------------------------------------------------------------------
\section{Visual \& Geometric Intuition: The Geometry of Gradient Descent Trajectories}

Figure~\ref{fig:gd_trajectories} contrasts the convergence trajectory of Gradient Descent on well-conditioned isotropic surfaces against ill-conditioned anisotropic surfaces.

\begin{figure}[htbp]
\centering
\begin{tikzpicture}[
    scale=0.9,
    font=\sffamily\footnotesize
]

% Left: Well-conditioned (circular contours)
\draw[thick, color=primarydark!40] (-4.5, 0) circle (2.2cm);
\draw[thick, color=primarydark!65] (-4.5, 0) circle (1.4cm);
\draw[thick, color=primarydark!90] (-4.5, 0) circle (0.6cm);
\filldraw[color=tealborder] (-4.5, 0) circle (2.5pt) node[below] {$\mathbf{w}^*$};

% Direct descent trajectory
\draw[->, >=stealth, line width=1.4pt, color=tealborder] (-6.3, 1.3) -- (-5.3, 0.6) -- (-4.8, 0.2) -- (-4.5, 0);
\filldraw[color=coralborder] (-6.3, 1.3) circle (2.5pt) node[above left] {$\mathbf{w}_0$};
\node[below, align=center] at (-4.5, -2.5) {\textbf{Isotropic Surface: } $\kappa = 1$ \\ Direct radial descent to $\mathbf{w}^*$};

% Right: Ill-conditioned (stretched elliptical contours)
\begin{scope}[shift={(4.0, 0)}]
    \draw[thick, color=primarydark!40] (0, 0) ellipse (3.6cm and 0.9cm);
    \draw[thick, color=primarydark!65] (0, 0) ellipse (2.4cm and 0.6cm);
    \draw[thick, color=primarydark!90] (0, 0) ellipse (1.2cm and 0.3cm);
    \filldraw[color=tealborder] (0, 0) circle (2.5pt) node[below] {$\mathbf{w}^*$};

    % Zig-zagging trajectory
    \coordinate (W0) at (-3.2, 1.4);
    \coordinate (W1) at (-2.2, -0.7);
    \coordinate (W2) at (-1.3, 0.5);
    \coordinate (W3) at (-0.6, -0.3);
    \coordinate (W4) at (-0.2, 0.1);

    \draw[->, >=stealth, line width=1.4pt, color=traprad] (W0) -- (W1);
    \draw[->, >=stealth, line width=1.4pt, color=traprad] (W1) -- (W2);
    \draw[->, >=stealth, line width=1.4pt, color=traprad] (W2) -- (W3);
    \draw[->, >=stealth, line width=1.4pt, color=traprad] (W3) -- (W4);
    \draw[->, >=stealth, line width=1.4pt, color=traprad] (W4) -- (0, 0);

    % Newton single step arrow
    \draw[->, >=stealth, line width=1.6pt, dashed, color=tealborder] (W0) to[bend left=15] node[above, sloped] {Newton Step (1 iteration)} (0, 0);

    \filldraw[color=coralborder] (W0) circle (2.5pt) node[above] {$\mathbf{w}_0$};
    \node[below, align=center] at (0, -2.5) {\textbf{Ill-Conditioned Surface: } $\kappa \gg 1$ \\ Severe cross-ravine zig-zagging};
\end{scope}

\end{tikzpicture}
\caption{The Geometry of Optimization Trajectories: On isotropic surfaces ($\kappa = 1$), gradient arrows point straight to the minimizer $\mathbf{w}^*$. On ill-conditioned surfaces ($\kappa \gg 1$), gradients point nearly perpendicular to the valley axis, causing severe oscillation. In contrast, Newton's method (dashed line) jumps directly to $\mathbf{w}^*$ in a single iteration.}
\label{fig:gd_trajectories}
\end{figure}

% ----------------------------------------------------------------------
% PART 4: DEEP INTUITION & CONCEPTUAL DERIVATIONS
% ----------------------------------------------------------------------
\section{Deep Intuition \& Conceptual Derivations}

\subsection{Exact Dynamics on Strictly Convex Quadratics}
Consider minimizing the quadratic objective:
\begin{equation}
f(\mathbf{w}) = \frac{1}{2}\mathbf{w}^T A \mathbf{w} - \mathbf{b}^T \mathbf{w}
\label{eq:quadratic_obj}
\end{equation}
where $A \in \mathbb{R}^{d \times d}$ is symmetric positive definite ($A \succ 0$) with eigenvalues:
\begin{equation}
0 < \mu = \lambda_{\min}(A) = \lambda_1 \le \lambda_2 \le \dots \le \lambda_d = \lambda_{\max}(A) = L
\end{equation}
The gradient is $\nabla f(\mathbf{w}) = A\mathbf{w} - \mathbf{b}$, and the unique global minimizer is $\mathbf{w}^* = A^{-1}\mathbf{b}$.

The Gradient Descent recurrence with step size $\eta$ is:
\begin{equation}
\mathbf{w}_{k+1} = \mathbf{w}_k - \eta (A\mathbf{w}_k - \mathbf{b})
\end{equation}
Subtracting $\mathbf{w}^* = A^{-1}\mathbf{b}$ from both sides:
\begin{align}
\mathbf{w}_{k+1} - \mathbf{w}^* &= \mathbf{w}_k - \mathbf{w}^* - \eta (A\mathbf{w}_k - A\mathbf{w}^*) \nonumber \\
&= (I - \eta A)(\mathbf{w}_k - \mathbf{w}^*)
\end{align}
Iterating this recurrence $k$ times:
\begin{equation}
\mathbf{w}_k - \mathbf{w}^* = (I - \eta A)^k (\mathbf{w}_0 - \mathbf{w}^*)
\label{eq:gd_matrix_power}
\end{equation}

\subsubsection{Eigen-Decomposition of the Error Vector}
Because $A$ is symmetric, it admits an orthonormal spectral decomposition $A = Q \Lambda Q^T = \sum_{i=1}^d \lambda_i \mathbf{q}_i \mathbf{q}_i^T$.
Expressing the initial error in the eigenbasis of $A$:
\begin{equation}
\mathbf{w}_0 - \mathbf{w}^* = \sum_{i=1}^d \alpha_i \mathbf{q}_i, \quad \text{where } \alpha_i = \mathbf{q}_i^T (\mathbf{w}_0 - \mathbf{w}^*)
\end{equation}
Since $(I - \eta A) \mathbf{q}_i = (1 - \eta \lambda_i) \mathbf{q}_i$, the error after $k$ steps uncouples into $d$ independent scalar progressions:
\begin{equation}
\mathbf{w}_k - \mathbf{w}^* = \sum_{i=1}^d (1 - \eta \lambda_i)^k \alpha_i \mathbf{q}_i
\label{eq:decoupled_error}
\end{equation}

\subsubsection{Stability Condition}
Convergence requires that every modal error shrinks to zero:
\begin{equation}
\lim_{k \to \infty} (1 - \eta \lambda_i)^k = 0 \iff |1 - \eta \lambda_i| < 1 \quad \forall i \in \{1, \dots, d\}
\end{equation}
This holds if and only if $-1 < 1 - \eta \lambda_i < 1$, which yields:
\begin{equation}
0 < \eta < \frac{2}{\lambda_{\max}} = \frac{2}{L}
\label{eq:stability_threshold}
\end{equation}
If $\eta \ge \frac{2}{L}$, the error along the eigenvector $\mathbf{q}_d$ of maximum eigenvalue $\lambda_{\max} = L$ oscillates with factor $|1 - \eta L| \ge 1$ and fails to converge (exploding if $\eta > 2/L$).

\subsubsection{Optimal Learning Rate and Convergence Factor}
The asymptotic convergence rate is governed by the spectral radius $\rho(I - \eta A) = \max_i |1 - \eta \lambda_i| = \max(|1 - \eta \mu|, |1 - \eta L|)$.
To minimize this worst-case contraction factor:
\begin{equation}
1 - \eta^* \mu = -(1 - \eta^* L) \implies \eta^* = \frac{2}{L + \mu}
\label{eq:optimal_eta}
\end{equation}
Substituting $\eta^*$ back into the contraction factor:
\begin{equation}
\rho^* = 1 - \eta^* \mu = 1 - \frac{2\mu}{L + \mu} = \frac{L - \mu}{L + \mu} = \frac{\frac{L}{\mu} - 1}{\frac{L}{\mu} + 1} = \frac{\kappa - 1}{\kappa + 1}
\label{eq:condition_contraction}
\end{equation}
Taking the Euclidean norm of~\eqref{eq:decoupled_error}:
\begin{equation}
\|\mathbf{w}_k - \mathbf{w}^*\|_2 \le \left( \frac{\kappa - 1}{\kappa + 1} \right)^k \|\mathbf{w}_0 - \mathbf{w}^*\|_2
\label{eq:linear_conv_rate}
\end{equation}
This is \textbf{linear convergence} (geometric decay in error). When the condition number is large ($\kappa \gg 1$), a first-order Taylor expansion reveals:
\begin{equation}
\frac{\kappa - 1}{\kappa + 1} = \frac{1 - 1/\kappa}{1 + 1/\kappa} \approx \left(1 - \frac{1}{\kappa}\right)\left(1 - \frac{1}{\kappa}\right) \approx 1 - \frac{2}{\kappa}
\end{equation}
To reduce the initial error by a factor of $\epsilon$, the number of iterations required is:
\begin{equation}
\left( 1 - \frac{2}{\kappa} \right)^k \le \epsilon \implies k \approx \frac{\kappa}{2} \ln\left(\frac{1}{\epsilon}\right) = \mathcal{O}(\kappa \ln(1/\epsilon))
\end{equation}
If $\kappa = 10{,}000$, gradient descent requires over $50{,}000$ iterations to achieve moderate precision!

\subsection{Newton-Raphson Method: Curvature Inversion}
Newton's method addresses the condition number curse by modeling the loss surface with its full second-order Taylor expansion around $\mathbf{w}_k$:
\begin{equation}
q(\mathbf{w}) = f(\mathbf{w}_k) + \nabla f(\mathbf{w}_k)^T (\mathbf{w} - \mathbf{w}_k) + \frac{1}{2}(\mathbf{w} - \mathbf{w}_k)^T \nabla^2 f(\mathbf{w}_k) (\mathbf{w} - \mathbf{w}_k)
\end{equation}
Minimizing this quadratic model with respect to $\mathbf{w}$ yields:
\begin{equation}
\nabla q(\mathbf{w}) = \nabla f(\mathbf{w}_k) + \nabla^2 f(\mathbf{w}_k)(\mathbf{w} - \mathbf{w}_k) = \mathbf{0}
\end{equation}
Assuming $\nabla^2 f(\mathbf{w}_k) \succ 0$, the exact minimizer of the local quadratic model defines the Newton update:
\begin{equation}
\mathbf{w}_{k+1} = \mathbf{w}_k - [\nabla^2 f(\mathbf{w}_k)]^{-1} \nabla f(\mathbf{w}_k)
\label{eq:newton_update}
\end{equation}

\subsubsection{One-Step Convergence on Quadratics}
For any quadratic $f(\mathbf{w}) = \frac{1}{2}\mathbf{w}^T A \mathbf{w} - \mathbf{b}^T \mathbf{w}$ with $A \succ 0$:
The gradient is $\nabla f(\mathbf{w}) = A\mathbf{w} - \mathbf{b}$ and the Hessian is $\nabla^2 f(\mathbf{w}) = A$.
Applying Newton's update from \emph{any} initial point $\mathbf{w}_0$:
\begin{equation}
\mathbf{w}_1 = \mathbf{w}_0 - A^{-1}(A\mathbf{w}_0 - \mathbf{b}) = \mathbf{w}_0 - \mathbf{w}_0 + A^{-1}\mathbf{b} = A^{-1}\mathbf{b} = \mathbf{w}^*
\end{equation}
Newton's method finds the exact global minimum in \textbf{precisely one iteration}, completely independent of the condition number $\kappa$!

\subsubsection{Affine Invariance of Newton's Method}
Unlike Gradient Descent, Newton's method is completely invariant to non-singular linear changes of coordinates.
Let $\mathbf{y} = B \mathbf{w}$ where $B \in \mathbb{R}^{d \times d}$ is invertible. Define $\tilde{f}(\mathbf{y}) = f(B^{-1}\mathbf{y})$.
By the chain rule:
\begin{align}
\nabla \tilde{f}(\mathbf{y}) &= B^{-T} \nabla f(\mathbf{w}) \\
\nabla^2 \tilde{f}(\mathbf{y}) &= B^{-T} \nabla^2 f(\mathbf{w}) B^{-1}
\end{align}
The Newton step in $\mathbf{y}$-coordinates is:
\begin{align}
\Delta \mathbf{y} &= -[\nabla^2 \tilde{f}(\mathbf{y})]^{-1} \nabla \tilde{f}(\mathbf{y}) \nonumber \\
&= -\left[ B^{-T} \nabla^2 f(\mathbf{w}) B^{-1} \right]^{-1} \left[ B^{-T} \nabla f(\mathbf{w}) \right] \nonumber \\
&= -B [\nabla^2 f(\mathbf{w})]^{-1} B^T B^{-T} \nabla f(\mathbf{w}) \nonumber \\
&= B \left( -[\nabla^2 f(\mathbf{w})]^{-1} \nabla f(\mathbf{w}) \right) = B \Delta \mathbf{w}
\end{align}
The update in the transformed space is identical to the transformation of the update in the original space: $\mathbf{y}_{k+1} = B \mathbf{w}_{k+1}$. Newton's method effectively rescales the coordinate axes so that the contours appear perfectly spherical.

\begin{ponderbox}[Why isn't Newton's Method used for Deep Learning?]
If Newton's method converges quadratically ($\| \mathbf{w}_{k+1} - \mathbf{w}^* \| \le C \| \mathbf{w}_k - \mathbf{w}^* \|^2$) and is invariant to ill-conditioning, why do modern deep learning frameworks (PyTorch, TensorFlow) almost universally rely on first-order variants (SGD, Adam)?

\textbf{Reflection:} Three severe obstacles prevent direct application of Newton's method in deep learning:
\begin{enumerate}
    \item \textbf{Memory Storage $\mathcal{O}(d^2)$:} A modest modern neural network has $d = 10^8$ parameters. Storing the Hessian requires $d^2 = 10^{16}$ floats ($40$ Petabytes of RAM).
    \item \textbf{Computational Inversion $\mathcal{O}(d^3)$:} Inverting a $10^8 \times 10^8$ matrix takes $\approx 10^{24}$ FLOPs, requiring centuries of computation per gradient step.
    \item \textbf{Non-Convexity and Saddle Points:} In high dimensions, random stationary points are overwhelmingly saddle points rather than local minima. Standard Newton's method is attracted to local \emph{maxima} and saddle points because it only seeks $\nabla f = \mathbf{0}$, jumping violently in directions of negative curvature ($\lambda_i < 0$).
\end{enumerate}
\end{ponderbox}

\subsection{Polyak Momentum: Heavy-Ball Acceleration}
The Heavy-Ball momentum method introduces a physical inertia term:
\begin{equation}
\mathbf{w}_{k+1} = \mathbf{w}_k - \alpha \nabla f(\mathbf{w}_k) + \beta (\mathbf{w}_k - \mathbf{w}_{k-1})
\label{eq:momentum_update}
\end{equation}
where $\beta \in [0, 1)$ is the momentum parameter.

\subsubsection{Discrete Dynamical System Formulation}
On a one-dimensional quadratic $f(w) = \frac{1}{2} a w^2$ with $a > 0$:
\begin{equation}
w_{k+1} = (1 - \alpha a + \beta) w_k - \beta w_{k-1}
\end{equation}
Stacking the state vector $\mathbf{u}_k = [w_k, w_{k-1}]^T \in \mathbb{R}^2$:
\begin{equation}
\begin{bmatrix} w_{k+1} \\ w_k \end{bmatrix} = \begin{bmatrix} 1 - \alpha a + \beta & -\beta \\ 1 & 0 \end{bmatrix} \begin{bmatrix} w_k \\ w_{k-1} \end{bmatrix} \implies \mathbf{u}_{k+1} = T \mathbf{u}_k
\label{eq:momentum_system}
\end{equation}
The characteristic polynomial of the transition matrix $T$ is:
\begin{equation}
\det(T - \lambda I) = \lambda^2 - (1 - \alpha a + \beta)\lambda + \beta = 0
\end{equation}
By Vieta's formulas, the product of the eigenvalues is $\lambda_1 \lambda_2 = \det(T) = \beta$.
When the roots are complex conjugates, their magnitudes are identically $|\lambda_1| = |\lambda_2| = \sqrt{\beta}$!
By optimally setting:
\begin{equation}
\alpha^* = \frac{4}{(\sqrt{L} + \sqrt{\mu})^2}, \quad \beta^* = \left( \frac{\sqrt{L} - \sqrt{\mu}}{\sqrt{L} + \sqrt{\mu}} \right)^2 = \left( \frac{\sqrt{\kappa} - 1}{\sqrt{\kappa} + 1} \right)^2
\label{eq:optimal_momentum}
\end{equation}
the spectral radius of the transition matrix becomes:
\begin{equation}
\rho(T) = \sqrt{\beta^*} = \frac{\sqrt{\kappa} - 1}{\sqrt{\kappa} + 1}
\end{equation}
This improves the condition number dependence from $\mathcal{O}(\kappa)$ to $\mathcal{O}(\sqrt{\kappa})$. For $\kappa = 10{,}000$, $\kappa$ drops from $10{,}000$ to $\sqrt{\kappa} = 100$, accelerating convergence by a factor of $100\times$!

% ----------------------------------------------------------------------
% PART 5: MULTI-PAGE WORKED EXAMPLES
% ----------------------------------------------------------------------
\section{Multi-Page Breakable Numerical Worked Examples}

\begin{examplebox}[breakable]{Step-by-Step Gradient Descent Dynamics on an Ill-Conditioned Quadratic}
\textbf{Problem Statement:}
Consider the quadratic objective function:
\begin{equation}
f(x_1, x_2) = \frac{1}{2} x_1^2 + 5 x_2^2 = \frac{1}{2} \mathbf{x}^T \begin{bmatrix} 1 & 0 \\ 0 & 10 \end{bmatrix} \mathbf{x}
\end{equation}
The unique global minimum is at $\mathbf{x}^* = [0, 0]^T$ with $f^* = 0$.
The starting point is $\mathbf{x}_0 = [10.0, 1.0]^T$.

\textbf{Task Requirements:}
\begin{enumerate}
    \item Identify the eigenvalues $\lambda_1, \lambda_2$, Lipschitz constant $L$, strong convexity parameter $\mu$, and condition number $\kappa$.
    \item Determine the maximum non-divergent step size $\eta_{\max}$ and the theoretical optimal step size $\eta^*$.
    \item Manually calculate the exact values of iterates $\mathbf{x}_1, \mathbf{x}_2$ and objective losses $f(\mathbf{x}_0), f(\mathbf{x}_1), f(\mathbf{x}_2)$ under the optimal learning rate $\eta^* = \frac{2}{L+\mu}$.
    \item Contrast this with the divergent behavior under step size $\eta = 0.22 > \eta_{\max}$.
\end{enumerate}

\tcblower

\textbf{Step 1: Curvature Parameters and Condition Number}
The Hessian is constant:
\begin{equation}
\nabla^2 f(\mathbf{x}) = A = \begin{bmatrix} 1 & 0 \\ 0 & 10 \end{bmatrix}
\end{equation}
The eigenvalues are:
\begin{equation}
\lambda_1 = \mu = 1.0, \quad \lambda_2 = L = 10.0
\end{equation}
The condition number is:
\begin{equation}
\kappa = \frac{L}{\mu} = \frac{10.0}{1.0} = 10
\end{equation}

\textbf{Step 2: Step Size Boundaries}
\begin{itemize}
    \item Maximum stable step size:
    \begin{equation}
    \eta_{\max} = \frac{2}{L} = \frac{2}{10.0} = 0.200
    \end{equation}
    \item Optimal asymptotic step size:
    \begin{equation}
    \eta^* = \frac{2}{L + \mu} = \frac{2}{10 + 1} = \frac{2}{11} \approx 0.181818
    \end{equation}
\end{itemize}

\textbf{Step 3: Exact Iterate Computations Under $\eta^* = \frac{2}{11}$}
The gradient is:
\begin{equation}
\nabla f(\mathbf{x}) = \begin{bmatrix} x_1 \\ 10 x_2 \end{bmatrix}
\end{equation}
The initial objective value at $\mathbf{x}_0 = [10, 1]^T$ is:
\begin{equation}
f(\mathbf{x}_0) = \frac{1}{2}(10)^2 + 5(1)^2 = 50 + 5 = 55.000
\end{equation}
The gradient at $\mathbf{x}_0$ is:
\begin{equation}
\nabla f(\mathbf{x}_0) = \begin{bmatrix} 10.0 \\ 10(1.0) \end{bmatrix} = \begin{bmatrix} 10.0 \\ 10.0 \end{bmatrix}
\end{equation}

\textbf{Iteration 1 ($\mathbf{x}_1$):}
\begin{align}
\mathbf{x}_1 &= \mathbf{x}_0 - \eta^* \nabla f(\mathbf{x}_0) \nonumber \\
&= \begin{bmatrix} 10 \\ 1 \end{bmatrix} - \frac{2}{11} \begin{bmatrix} 10 \\ 10 \end{bmatrix}
= \begin{bmatrix} 10 - \frac{20}{11} \\ 1 - \frac{20}{11} \end{bmatrix}
= \begin{bmatrix} \frac{90}{11} \\ -\frac{9}{11} \end{bmatrix} \approx \begin{bmatrix} 8.1818 \\ -0.8182 \end{bmatrix}
\end{align}
Notice the critical diagnostic event: $x_2$ has flipped sign from $+1.0$ to $-0.8182$! It overshot the bottom of the ravine and landed on the opposite wall.
The objective value at $\mathbf{x}_1$ is:
\begin{align}
f(\mathbf{x}_1) &= \frac{1}{2}\left(\frac{90}{11}\right)^2 + 5\left(-\frac{9}{11}\right)^2 \nonumber \\
&= \frac{1}{2} \cdot \frac{8100}{121} + 5 \cdot \frac{81}{121} = \frac{4050 + 405}{121} = \frac{4455}{121} \approx 36.8182
\end{align}

\textbf{Iteration 2 ($\mathbf{x}_2$):}
The gradient at $\mathbf{x}_1$ is:
\begin{equation}
\nabla f(\mathbf{x}_1) = \begin{bmatrix} \frac{90}{11} \\ 10\left(-\frac{9}{11}\right) \end{bmatrix} = \begin{bmatrix} \frac{90}{11} \\ -\frac{90}{11} \end{bmatrix}
\end{equation}
Now update:
\begin{align}
\mathbf{x}_2 &= \mathbf{x}_1 - \eta^* \nabla f(\mathbf{x}_1) \nonumber \\
&= \begin{bmatrix} \frac{90}{11} \\ -\frac{9}{11} \end{bmatrix} - \frac{2}{11} \begin{bmatrix} \frac{90}{11} \\ -\frac{90}{11} \end{bmatrix}
= \begin{bmatrix} \frac{90}{11}\left(1 - \frac{2}{11}\right) \\ -\frac{9}{11}\left(1 - \frac{20}{11}\right) \end{bmatrix}
= \begin{bmatrix} \frac{90}{11} \cdot \frac{9}{11} \\ -\frac{9}{11} \left(-\frac{9}{11}\right) \end{bmatrix}
= \begin{bmatrix} \frac{810}{121} \\ \frac{81}{121} \end{bmatrix} \approx \begin{bmatrix} 6.6942 \\ +0.6694 \end{bmatrix}
\end{align}
The sign of $x_2$ flips back to positive ($+0.6694$).
The coordinate progression for any step $k$ is given analytically by:
\begin{align}
x_1^{(k)} &= \left(1 - \eta^* \lambda_1\right)^k x_1^{(0)} = \left(1 - \frac{2}{11}\right)^k (10) = 10 \left(\frac{9}{11}\right)^k \\
x_2^{(k)} &= \left(1 - \eta^* \lambda_2\right)^k x_2^{(0)} = \left(1 - \frac{20}{11}\right)^k (1) = 1 \left(-\frac{9}{11}\right)^k
\end{align}
Observe that both coordinate errors contract at the exact same rate:
\begin{equation}
|x_1^{(k)}| = 10 \left(\frac{9}{11}\right)^k, \quad |x_2^{(k)}| = 1 \left(\frac{9}{11}\right)^k
\end{equation}
The contraction factor is precisely $\rho^* = \frac{\kappa - 1}{\kappa + 1} = \frac{10 - 1}{10 + 1} = \frac{9}{11} \approx 0.8182$.
The objective value at $\mathbf{x}_2$ is:
\begin{equation}
f(\mathbf{x}_2) = \left(\frac{9}{11}\right)^2 f(\mathbf{x}_1) = \frac{81}{121} \times \frac{4455}{121} \approx 24.646
\end{equation}

\textbf{Step 4: Divergence with $\eta = 0.22 > \eta_{\max}$}
Suppose a user chooses step size $\eta = 0.22 > 0.20$.
The contraction factors along the two axes are:
\begin{align}
1 - \eta \lambda_1 &= 1 - 0.22(1) = +0.78 \quad (\text{stable, contracts}) \\
1 - \eta \lambda_2 &= 1 - 0.22(10) = 1 - 2.20 = -1.20 \quad (\text{UNSTABLE!})
\end{align}
Along the $x_2$ axis, the error expands exponentially:
\begin{equation}
x_2^{(k)} = (-1.20)^k \cdot (1.0)
\end{equation}
\begin{itemize}
    \item $x_2^{(0)} = +1.0$
    \item $x_2^{(1)} = -1.20$
    \item $x_2^{(2)} = +1.44$
    \item $x_2^{(5)} = -2.488$
    \item $x_2^{(20)} \approx +38.33$
    \item $x_2^{(100)} \approx 8.28 \times 10^7$
\end{itemize}
Even though $x_1 \to 0$, $x_2$ explodes violently to $\pm \infty$, triggering numerical overflow (\texttt{NaN} in Python).
\end{examplebox}

\begin{examplebox}[breakable]{Newton's Method vs.\ Gradient Descent on a Multi-Variable Objective}
\textbf{Problem Statement:}
Consider the non-quadratic function:
\begin{equation}
f(x_1, x_2) = x_1^4 + 2x_1 x_2 + 2x_2^2
\end{equation}
Starting at the initial test point $\mathbf{x}_0 = [1.0, 1.0]^T$:
\begin{enumerate}
    \item Compute the gradient $\nabla f(\mathbf{x}_0)$.
    \item Compute the Hessian matrix $\nabla^2 f(\mathbf{x}_0)$ and verify that it is strictly positive definite.
    \item Compute the exact inverse Hessian $[\nabla^2 f(\mathbf{x}_0)]^{-1}$.
    \item Compute the first Newton iterate $\mathbf{x}_1^{\text{Newton}}$ and compare it with the Gradient Descent iterate $\mathbf{x}_1^{\text{GD}}$ using $\eta = 0.1$.
\end{enumerate}

\tcblower

\textbf{Step 1: Gradient Vector}
The partial derivatives are:
\begin{align}
\frac{\partial f}{\partial x_1} &= 4x_1^3 + 2x_2 \\
\frac{\partial f}{\partial x_2} &= 2x_1 + 4x_2
\end{align}
Evaluating at $\mathbf{x}_0 = [1, 1]^T$:
\begin{equation}
\nabla f(\mathbf{x}_0) = \begin{bmatrix} 4(1)^3 + 2(1) \\ 2(1) + 4(1) \end{bmatrix} = \begin{bmatrix} 6 \\ 6 \end{bmatrix}
\end{equation}

\textbf{Step 2: Hessian Matrix Evaluation}
The second partial derivatives are:
\begin{align}
\frac{\partial^2 f}{\partial x_1^2} &= 12x_1^2, \quad \frac{\partial^2 f}{\partial x_1 \partial x_2} = 2 \\
\frac{\partial^2 f}{\partial x_2 \partial x_1} &= 2, \quad \frac{\partial^2 f}{\partial x_2^2} = 4
\end{align}
Evaluating at $\mathbf{x}_0 = [1, 1]^T$:
\begin{equation}
H_0 = \nabla^2 f(\mathbf{x}_0) = \begin{bmatrix} 12 & 2 \\ 2 & 4 \end{bmatrix}
\end{equation}
Check positive definiteness:
\begin{itemize}
    \item Leading principal minor 1: $\Delta_1 = 12 > 0$.
    \item Leading principal minor 2: $\Delta_2 = \det(H_0) = (12)(4) - (2)(2) = 48 - 4 = 44 > 0$.
\end{itemize}
By Sylvester's Criterion, $H_0 \succ 0$ strictly.

\textbf{Step 3: Inverse Hessian Matrix}
Using the $2 \times 2$ matrix inverse formula:
\begin{equation}
H_0^{-1} = \frac{1}{\det(H_0)} \begin{bmatrix} 4 & -2 \\ -2 & 12 \end{bmatrix} = \frac{1}{44} \begin{bmatrix} 4 & -2 \\ -2 & 12 \end{bmatrix} = \begin{bmatrix} \frac{1}{11} & -\frac{1}{22} \\ -\frac{1}{22} & \frac{3}{11} \end{bmatrix}
\end{equation}

\textbf{Step 4: Comparison of Iterates}
\begin{itemize}
    \item \textbf{Newton Iteration:}
    \begin{align}
    \Delta \mathbf{x}^{\text{Newton}} &= -H_0^{-1} \nabla f(\mathbf{x}_0) = -\frac{1}{44} \begin{bmatrix} 4 & -2 \\ -2 & 12 \end{bmatrix} \begin{bmatrix} 6 \\ 6 \end{bmatrix} \nonumber \\
    &= -\frac{1}{44} \begin{bmatrix} 24 - 12 \\ -12 + 72 \end{bmatrix} = -\frac{1}{44} \begin{bmatrix} 12 \\ 60 \end{bmatrix} = \begin{bmatrix} -\frac{3}{11} \\ -\frac{15}{11} \end{bmatrix} \approx \begin{bmatrix} -0.2727 \\ -1.3636 \end{bmatrix}
    \end{align}
    The updated parameter vector is:
    \begin{equation}
    \mathbf{x}_1^{\text{Newton}} = \mathbf{x}_0 + \Delta \mathbf{x}^{\text{Newton}} = \begin{bmatrix} 1 \\ 1 \end{bmatrix} + \begin{bmatrix} -\frac{3}{11} \\ -\frac{15}{11} \end{bmatrix} = \begin{bmatrix} \frac{8}{11} \\ -\frac{4}{11} \end{bmatrix} \approx \begin{bmatrix} 0.7273 \\ -0.3636 \end{bmatrix}
    \end{equation}

    \item \textbf{Gradient Descent Iteration ($\eta = 0.1$):}
    \begin{equation}
    \mathbf{x}_1^{\text{GD}} = \mathbf{x}_0 - \eta \nabla f(\mathbf{x}_0) = \begin{bmatrix} 1 \\ 1 \end{bmatrix} - 0.1 \begin{bmatrix} 6 \\ 6 \end{bmatrix} = \begin{bmatrix} 0.4 \\ 0.4 \end{bmatrix}
    \end{equation}
\end{itemize}
Notice that Gradient Descent moves equally along both coordinates ($0.4, 0.4$) because the gradient vector had identical components ($6, 6$). However, Newton's method recognized that the curvature along $x_1$ ($12$) is three times greater than along $x_2$ ($4$), so it took a modest step along $x_1$ and a much larger, corrective step along $x_2$, adapting automatically to the local geometry.
\end{examplebox}

% ----------------------------------------------------------------------
% PART 6: EXPLICIT GOTCHAS & DIAGNOSTIC MISCONCEPTIONS
% ----------------------------------------------------------------------
\section{Explicit Gotchas \& Diagnostic Misconceptions}

\begin{trapbox}{Gotcha 1: Selecting Learning Rate Above the Theoretical Stability Bound ($2/L$)}
A student sets learning rate $\eta = 0.25$ on a problem where $L = 10$. They observe the loss oscillating wildly and overflowing to \texttt{NaN} within 15 iterations.
\textbf{The Trap:} The divergence boundary is strictly $\eta_{\max} = \frac{2}{L}$. For $L = 10$, $\eta_{\max} = \frac{2}{10} = 0.20$. Any step size $\eta \ge 0.20$ creates an error amplification factor $|1 - \eta L| \ge 1$. The error grows exponentially as $|1 - \eta L|^k \to \infty$. Step sizes must satisfy $\eta < \frac{2}{L}$. In machine learning practice, setting $\eta \le \frac{1}{L}$ guarantees monotonic objective decrease.
\end{trapbox}

\begin{trapbox}{Gotcha 2: Assuming Gradient Descent Always Finds the Global Minimum}
A practitioner executes gradient descent on a complex deep neural network and asserts the final parameter vector is the global minimizer.
\textbf{The Trap:} Gradient descent is guaranteed to find the global minimum only if the objective function is \textbf{convex}. On non-convex loss surfaces, gradient descent can be trapped in suboptimal local minima or stall indefinitely near saddle points and plateaus where $\nabla f(\mathbf{w}) \approx \mathbf{0}$.
\end{trapbox}

\begin{trapbox}{Gotcha 3: Conflating Stationary Points with Local Minima}
An optimization script terminates when $\|\nabla f(\mathbf{w})\|_2 < 10^{-5}$. The engineer logs that a local minimum has been discovered.
\textbf{The Trap:} A vanishing gradient ($\nabla f = \mathbf{0}$) is a \textbf{necessary} condition, not a sufficient condition. On non-convex surfaces, stationary points can be local maxima or saddle points. Verifying a local minimum requires evaluating the Hessian eigenvalues: $\lambda_i(\nabla^2 f) > 0$ for all $i$.
\end{trapbox}

\begin{trapbox}{Gotcha 4: Inverting Hessians in High Dimensions (The $\mathcal{O}(d^3)$ Barrier)}
An engineer decides to replace gradient descent with Newton's method on a modern computer vision model with $d = 10^7$ parameters.
\textbf{The Trap:} Storing the Hessian requires storing $d^2 = 10^{14}$ floating-point values ($800$ Terabytes of RAM). Inverting it via Gaussian elimination requires $\mathcal{O}(d^3) = 10^{21}$ FLOPs, which would take centuries of compute time. Second-order methods in modern ML require quasi-Newton approximations (L-BFGS) or Hessian-free conjugate gradient methods.
\end{trapbox}

% ----------------------------------------------------------------------
% PART 7: COMPARATIVE SYNTHESIS & SUMMARY TABLE
% ----------------------------------------------------------------------
\section{Comparative Synthesis \& Summary Table}

Table~\ref{tab:optim_algorithms} compares the primary numerical optimization algorithms used in machine learning.

\begin{table}[htbp]
\centering
\small
\begin{tabularx}{\textwidth}{lXXXXX}
\toprule
\textbf{Algorithm} & \textbf{Order} & \textbf{Per-Iter Cost} & \textbf{Convergence Rate (Convex)} & \textbf{Memory Requirement} & \textbf{Sensitivity to $\kappa$} \\
\midrule
Gradient Descent (GD) & 1st & $\mathcal{O}(d)$ & Sublinear $\mathcal{O}(1/k)$ & $\mathcal{O}(d)$ & Extreme zig-zagging for large $\kappa$ \\
\addlinespace
GD with Momentum & 1st & $\mathcal{O}(d)$ & Accelerated $\mathcal{O}(1/k^2)$ & $\mathcal{O}(d)$ & Damps high-frequency oscillations \\
\addlinespace
Newton-Raphson & 2nd & $\mathcal{O}(d^3)$ & Quadratic $\mathcal{O}(e^{-2^k})$ & $\mathcal{O}(d^2)$ & Completely invariant to condition number \\
\addlinespace
Quasi-Newton (BFGS) & 1st + approx 2nd & $\mathcal{O}(d^2)$ & Superlinear & $\mathcal{O}(d^2)$ & Robust; builds Hessian estimate iteratively \\
\addlinespace
L-BFGS & Limited memory & $\mathcal{O}(m d)$ & Superlinear & $\mathcal{O}(m d)$ ($m \approx 10$) & Industrial standard for medium-scale convex ML \\
\bottomrule
\end{tabularx}
\caption{Comparative Synthesis: Numerical Optimization Algorithms in Machine Learning.}
\label{tab:optim_algorithms}
\end{table}

% ----------------------------------------------------------------------
% PART 8: PYTHON / SCIENTIFIC COMPUTING LAB
% ----------------------------------------------------------------------
\section{Python / Scientific Computing Lab}

The following production script implements Gradient Descent, empirically demonstrates divergence when $\eta > 2/L$, and contrasts isotropic vs.\ ill-conditioned convergence dynamics.

\begin{lstlisting}[language=Python, caption={Gradient Descent Dynamics and Stability Threshold Simulation in Python.}]
import numpy as np

def run_gradient_descent(A: np.ndarray, b: np.ndarray, w0: np.ndarray, lr: float, n_iters: int):
    """Simulates GD on quadratic objective f(w) = 0.5 * w^T A w - b^T w."""
    w = w0.copy()
    history = []
    for k in range(n_iters):
        grad = A @ w - b
        loss = 0.5 * float(w.T @ A @ w) - float(b.T @ w)
        history.append(loss)
        w = w - lr * grad
        if np.isnan(loss) or np.isinf(loss) or loss > 1e12:
            break
    return w, history

def main():
    # Construct an ill-conditioned 2D quadratic: lambda_min = 1.0, lambda_max = 50.0 (kappa = 50)
    L = 50.0
    mu = 1.0
    A = np.array([[L, 0.0], [0.0, mu]])
    b = np.array([0.0, 0.0])  # Minimum at w* = [0, 0]^T
    w0 = np.array([2.0, 2.0])

    stable_lr_bound = 2.0 / L
    optimal_lr = 2.0 / (L + mu)

    print(f"Theoretical Stability Bound (2/L): {stable_lr_bound:.4f}")
    print(f"Optimal Learning Rate 2/(L+mu):   {optimal_lr:.4f}")

    # Case 1: Optimal Learning Rate
    _, hist_opt = run_gradient_descent(A, b, w0, lr=optimal_lr, n_iters=25)
    # Case 2: Sub-optimal Small Learning Rate
    _, hist_small = run_gradient_descent(A, b, w0, lr=0.005, n_iters=25)
    # Case 3: Divergent Learning Rate (eta > 2/L)
    _, hist_div = run_gradient_descent(A, b, w0, lr=0.045, n_iters=10)

    print("\nIteration Loss Comparison:")
    print(f"{'Iter':<6}{'Optimal LR':<18}{'Small LR':<18}{'Divergent LR (> 2/L)':<25}")
    print("-" * 65)
    for i in range(min(10, len(hist_opt))):
        div_val = f"{hist_div[i]:.2e}" if i < len(hist_div) else "OVERFLOW"
        print(f"{i:<6}{hist_opt[i]:<18.4e}{hist_small[i]:<18.4e}{div_val:<25}")

if __name__ == "__main__":
    main()
\end{lstlisting}

% ----------------------------------------------------------------------
% PART 9: COMPREHENSIVE PRACTICE PROBLEMS
% ----------------------------------------------------------------------
\section{Comprehensive Practice Problems}

\subsection*{Part A: Analytical \& Mathematical Problems}

\begin{enumerate}[label=\textbf{P\arabic*.}]
    \item \textbf{Analytical Step-Size Bound on a 2D Quadratic:}
    Let $f(x_1, x_2) = 10 x_1^2 + x_2^2 - 4x_1 - 2x_2$.
    \begin{enumerate}[label=(\alph*)]
        \item Compute the gradient $\nabla f(\mathbf{x})$ and Hessian $\nabla^2 f(\mathbf{x})$.
        \item Find the unique global minimizer $\mathbf{x}^*$.
        \item Find the Lipschitz constant $L$ and strong convexity constant $\mu$.
        \item Calculate the maximum stable step size $\eta_{\max}$ and the optimal step size $\eta^*$.
    \end{enumerate}

    \item \textbf{Derivation of Convergence Rate for Convex $L$-Smooth Functions:}
    Let $f$ be convex and $L$-smooth. Suppose we run Gradient Descent with $\eta = \frac{1}{L}$ for $T$ iterations.
    Using the descent lemma and convexity inequality $f(\mathbf{w}_k) - f^* \le \nabla f(\mathbf{w}_k)^T (\mathbf{w}_k - \mathbf{w}^*)$, prove that:
    \begin{equation}
    f(\mathbf{w}_T) - f^* \le \frac{L \|\mathbf{w}_0 - \mathbf{w}^*\|_2^2}{2T}
    \end{equation}

    \item \textbf{Exact Newton Step on Non-Quadratic Function:}
    Let $f(x) = x^4 + 2x^2 + 4x$.
    \begin{enumerate}[label=(\alph*)]
        \item Compute $f'(x)$ and $f''(x)$.
        \item Starting from initial guess $x_0 = 1$, compute the first two Newton-Raphson iterates $x_1$ and $x_2$.
        \item Contrast the rapid convergence of Newton's method against Gradient Descent with $\eta = 0.05$.
    \end{enumerate}

    \item \textbf{Momentum Acceleration Dynamics:}
    The Heavy-Ball Momentum method updates parameters via:
    \begin{equation}
    \mathbf{w}_{k+1} = \mathbf{w}_k - \alpha \nabla f(\mathbf{w}_k) + \beta (\mathbf{w}_k - \mathbf{w}_{k-1})
    \end{equation}
    Show that for a 1D quadratic $f(w) = \frac{1}{2} a w^2$, this recurrence can be written as a $2 \times 2$ discrete dynamical system. Determine the eigenvalues of the transition matrix and state the conditions on $\alpha, \beta$ for asymptotic stability.

    \item \textbf{Strongly Convex Linear Convergence Rate:}
    Let $f$ be $\mu$-strongly convex and $L$-smooth. Prove that with step size $\eta = \frac{2}{L + \mu}$:
    \begin{equation}
    \|\mathbf{w}_k - \mathbf{w}^*\|_2 \le \left( \frac{\kappa - 1}{\kappa + 1} \right)^k \|\mathbf{w}_0 - \mathbf{w}^*\|_2
    \end{equation}
    where $\kappa = L / \mu$.
\end{enumerate}

\subsection*{Part B: Diagnostic Multiple-Choice Questions}

\begin{enumerate}[label=\textbf{MCQ\arabic*.}]
    \item A loss function $f(\mathbf{w})$ is $L$-smooth with $L = 10$. What is the largest constant step size $\eta$ that guarantees the gradient descent sequence does NOT diverge on an arbitrary convex quadratic loss?
    \begin{enumerate}[label=(\alph*)]
        \item $\eta = 0.1$
        \item $\eta = 0.2$
        \item $\eta = 0.05$
        \item $\eta = 1.0$
    \end{enumerate}

    \item Why does Gradient Descent exhibit slow, zig-zagging convergence on an objective with condition number $\kappa = 10{,}000$?
    \begin{enumerate}[label=(\alph*)]
        \item Because the gradient vector is dominated by the direction of highest curvature, forcing very small step sizes to prevent divergence along that axis, while progress along low-curvature directions is sluggish.
        \item Because the objective function is non-convex.
        \item Because the gradient norm evaluates to zero at all points along the trajectory.
        \item Because the Hessian matrix is singular and cannot be inverted.
    \end{enumerate}
\end{enumerate}

\chapter[Constrained Optimization \& Convex Sets]{Constrained Optimization, Lagrange Multipliers, Convex Sets \& Convex Functions}
\label{chap:week09}

% ----------------------------------------------------------------------
% PART 1: EXECUTIVE ROADMAP & ARCHITECTURAL FOUNDATIONS
% ----------------------------------------------------------------------
\section{Executive Roadmap \& Learning Objectives}

\subsection{Opening Hook: The Boundary Fence and the Rolling Marble}
Imagine a vast, smooth hillside forming an elliptical bowl. At the very bottom of the bowl lies a deep depression: the unconstrained global minimum where a marble placed anywhere on the hillside would naturally roll and come to rest. 

Now imagine that a surveyor erects a rigid, straight wooden fence line cutting across the hillside high above the valley floor. If you release the marble high on the upper slope, where will it come to rest against the fence?

A common intuitive instinct is to look at where the bottom of the bowl lies, draw a straight perpendicular line from the valley floor to the fence, and assume the marble stops at that closest geometric point (the orthogonal Euclidean projection). But physical dynamics and mathematical optimization defy this naive projection! 

The marble does \emph{not} halt at the closest point to the valley floor. Gravity pulls the marble down the local slope; as it strikes the fence, the component of gravity parallel to the fence rail drives the marble rolling sideways along the wood. The marble continues to roll until the lateral downhill pull vanishes completely---that is, until the downward gravitational force is oriented purely perpendicular to the fence rail, exactly opposed by the normal contact force of the wood. 

The marble comes to rest at the precise point of \textbf{tangency}, where an elevation contour of the hill runs parallel to the fence line. 

In continuous mathematical optimization, this physical reality resolves the \textbf{Boundary Curse Paradox}. Finding a constrained minimum is fundamentally different from projecting an unconstrained minimum onto a constraint boundary. The contact force exerted by the fence is the physical embodiment of the \textbf{Lagrange Multiplier} $\nu^*$. It provides the exact normal counter-force required to balance the gradient of the objective functional.

\subsection{Architectural Bridge from Previous Chapters}
In Chapter~\ref{chap:week08}, we explored unconstrained optimization algorithms where parameter vectors $\mathbf{w}$ were permitted to roam freely across the entire Euclidean space $\mathbb{R}^d$. We analyzed gradient descent, Polyak momentum, and Newton-Raphson updates governed purely by local Taylor expansions and Hessian curvature.

However, in modern statistical machine learning, parameters are almost never unconstrained:
\begin{itemize}
    \item In multinomial classification, attention mechanisms, and Gaussian Mixture Models (Chapter~\ref{chap:week12}), parameter vectors represent probabilities and must reside on the \textbf{probability simplex}: $\sum_{i=1}^d w_i = 1$ with $w_i \ge 0$.
    \item In sparse feature selection via Lasso regression, parameters are bounded within an $L_1$ norm ball: $\|\mathbf{w}\|_1 \le s$.
    \item In hard-margin and soft-margin Support Vector Machines (Chapter~\ref{chap:week10}), classification margins enforce strict linear inequality boundaries: $y_i(\mathbf{w}^T \mathbf{x}_i + b) \ge 1$.
\end{itemize}
Chapter 9 bridges the gap between unconstrained vector calculus and constrained mathematical programming. We establish the geometric theory of Lagrange multipliers for manifolds, characterize convex sets and convex functions, prove the fundamental hyperplane separation theorems, and develop the mechanics of the Euclidean projection operator. This provides the mathematical launchpad for Karush-Kuhn-Tucker (KKT) duality in Chapter~\ref{chap:week10}.

\subsection{Prerequisites \& Scaffolding}
Mastery of this chapter builds upon vector calculus, linear geometry, and matrix analysis as organized in Table~\ref{tab:w09_prereqs}.

\begin{table}[htbp]
\centering
\small
\begin{tabularx}{\textwidth}{lXX}
\toprule
\textbf{Foundational Area} & \textbf{Core Concepts Required} & \textbf{Role in Machine Learning Foundations} \\
\midrule
Vector Calculus & Gradient vectors $\nabla f$, level surfaces $f(\mathbf{x}) = c$, tangent hyperplanes & Formulating orthogonality between objective gradients and constraint manifolds \\
\addlinespace
Linear Algebra & Hyperplanes $\mathbf{a}^T \mathbf{x} = b$, halfspaces $\mathbf{a}^T \mathbf{x} \le b$, orthogonal projections & Defining polyhedral feasible regions, normal cones, and projection operators \\
\addlinespace
Convex Analysis & Jensen's inequality, epigraphs, positive semidefinite cones $S^d_+$ & Establishing global optimality guarantees and non-expansiveness bounds \\
\bottomrule
\end{tabularx}
\caption{Prerequisite Scaffolding for Chapter 9.}
\label{tab:w09_prereqs}
\end{table}

\subsection{Core Competencies}
By the conclusion of this chapter, students will be able to:
\begin{itemize}
    \item Formulate equality-constrained optimization problems and apply the Method of Lagrange Multipliers.
    \item Prove that at a constrained local extremum, the objective gradient is collinear with the normal space of the constraint manifold.
    \item Define convex sets and verify domain convexity using intersection, affine map, and Cartesian product closure.
    \item State and apply the Separating Hyperplane Theorem and Supporting Hyperplane Theorem.
    \item Characterize convex functions via Jensen's inequality, first-order supporting hyperplanes, and second-order Hessian conditions.
    \item Prove the non-expansiveness of the Euclidean projection operator onto closed convex sets.
    \item Implement Projected Gradient Descent (PGD) with sparse $L_1$ ball projection.
\end{itemize}

\subsection{Notation \& Mathematical Conventions}
Throughout this chapter, we adhere to the standard mathematical conventions:
\begin{itemize}
    \item Feasible constraint set: $C \subseteq \mathbb{R}^d$.
    \item Euclidean projection operator: $\Pi_C(\mathbf{x}) \triangleq \argmin_{\mathbf{z} \in C} \|\mathbf{x} - \mathbf{z}\|_2$.
    \item Epigraph of a functional: $\operatorname{epi}(f) \triangleq \{(\mathbf{x}, t) \in \mathbb{R}^d \times \mathbb{R} \mid f(\mathbf{x}) \le t\}$.
    \item Sublevel set: $\mathcal{S}_\alpha(f) \triangleq \{\mathbf{x} \in \mathbb{R}^d \mid f(\mathbf{x}) \le \alpha\}$.
    \item Lagrangian functional: $\mathcal{L}(\mathbf{w}, \boldsymbol{\nu}) \triangleq f(\mathbf{w}) + \sum_{j=1}^p \nu_j h_j(\mathbf{w})$.
    \item Positive semidefinite cone: $S^d_+ \triangleq \{A \in \mathbb{R}^{d \times d} \mid A = A^T, \, \mathbf{v}^T A \mathbf{v} \ge 0 \; \forall \mathbf{v} \in \mathbb{R}^d\}$.
\end{itemize}

% ----------------------------------------------------------------------
% PART 2: MATHEMATICAL & COMPUTATIONAL FOUNDATIONS
% ----------------------------------------------------------------------
\section{Mathematical \& Computational Foundations}

\subsection{Equality Constraints and the Method of Lagrange Multipliers}
Consider the equality-constrained continuous optimization problem:
\begin{equation}
\min_{\mathbf{w} \in \mathbb{R}^d} f(\mathbf{w}) \quad \text{subject to } h_j(\mathbf{w}) = 0 \quad (j = 1, \dots, p)
\label{eq:eq_constrained_prob}
\end{equation}
where $f, h_1, \dots, h_p: \mathbb{R}^d \to \mathbb{R}$ are continuously differentiable functions with $p < d$. The feasible domain forms a $(d - p)$-dimensional geometric manifold:
\begin{equation}
\mathcal{M} = \{\mathbf{w} \in \mathbb{R}^d \mid h_j(\mathbf{w}) = 0, \; j = 1, \dots, p\}
\end{equation}

\begin{theorembox}{Geometric Condition of Optimality (Lagrange Multipliers)}
Let $\mathbf{w}^* \in \mathcal{M}$ be a local minimizer of $f$ on the feasible manifold $\mathcal{M}$. Assume the constraint gradients $\{\nabla h_1(\mathbf{w}^*), \dots, \nabla h_p(\mathbf{w}^*)\}$ are linearly independent (the \textbf{Linear Independence Constraint Qualification}, LICQ).

Then the gradient of the objective functional must reside strictly within the normal space spanned by the constraint gradients:
\begin{equation}
\nabla f(\mathbf{w}^*) + \sum_{j=1}^p \nu_j^* \nabla h_j(\mathbf{w}^*) = \mathbf{0}
\label{eq:lagrange_stationarity}
\end{equation}
for unique real scalars $\nu_1^*, \dots, \nu_p^* \in \mathbb{R}$ designated as \textbf{Lagrange Multipliers}.
\end{theorembox}

\begin{proof}
Let $\mathbf{w}^*$ be a local minimizer on $\mathcal{M}$. Consider any smooth curve $\boldsymbol{\gamma}: (-\epsilon, \epsilon) \to \mathcal{M}$ lying entirely within the feasible manifold such that $\boldsymbol{\gamma}(0) = \mathbf{w}^*$ and $\boldsymbol{\gamma}'(0) = \mathbf{v}$.
Since $\boldsymbol{\gamma}(t) \in \mathcal{M}$ for all $t$, every constraint satisfies $h_j(\boldsymbol{\gamma}(t)) = 0$ identically. Differentiating with respect to $t$ at $t = 0$ via the multivariable chain rule:
\begin{equation}
\frac{d}{dt} h_j(\boldsymbol{\gamma}(t))\Big|_{t=0} = \nabla h_j(\mathbf{w}^*)^T \boldsymbol{\gamma}'(0) = \nabla h_j(\mathbf{w}^*)^T \mathbf{v} = 0 \quad \forall j \in \{1, \dots, p\}
\end{equation}
Thus, the tangent space to the feasible manifold at $\mathbf{w}^*$ is precisely the orthogonal complement of the span of the constraint gradients:
\begin{equation}
T_{\mathcal{M}}(\mathbf{w}^*) = \left\{ \mathbf{v} \in \mathbb{R}^d \;\middle|\; \nabla h_j(\mathbf{w}^*)^T \mathbf{v} = 0, \, \forall j \right\} = \left( \Span\{\nabla h_1(\mathbf{w}^*), \dots, \nabla h_p(\mathbf{w}^*)\} \right)^\perp
\end{equation}
Now evaluate the rate of change of the objective functional along the feasible curve $\boldsymbol{\gamma}(t)$. Because $\mathbf{w}^*$ is a local minimum, the scalar function $\phi(t) = f(\boldsymbol{\gamma}(t))$ attains a local minimum at $t = 0$. By elementary calculus, $\phi'(0) = 0$:
\begin{equation}
\phi'(0) = \nabla f(\mathbf{w}^*)^T \boldsymbol{\gamma}'(0) = \nabla f(\mathbf{w}^*)^T \mathbf{v} = 0
\end{equation}
Because this orthogonality holds for \emph{every} feasible tangent vector $\mathbf{v} \in T_{\mathcal{M}}(\mathbf{w}^*)$:
\begin{equation}
\nabla f(\mathbf{w}^*) \in \left( T_{\mathcal{M}}(\mathbf{w}^*) \right)^\perp
\end{equation}
By the Fundamental Theorem of Linear Algebra (Chapter~\ref{chap:week02}), $\left( (V)^\perp \right)^\perp = V$. Therefore:
\begin{equation}
\nabla f(\mathbf{w}^*) \in \Span\{\nabla h_1(\mathbf{w}^*), \dots, \nabla h_p(\mathbf{w}^*)\}
\end{equation}
Expressing $\nabla f(\mathbf{w}^*)$ as a linear combination of $\{\nabla h_j(\mathbf{w}^*)\}$ with coefficients $-\nu_j^*$ yields~\eqref{eq:lagrange_stationarity}. Linear independence (LICQ) guarantees that these multiplier coefficients $\nu_j^*$ are uniquely determined.
\end{proof}

\begin{theorembox}{Lagrange Multiplier Sensitivity \& Shadow Price Theorem}
Let $f, h_1, \dots, h_p: \mathbb{R}^d \to \mathbb{R}$ be twice continuously differentiable. Consider the perturbed family of equality-constrained optimization problems:
\begin{equation}
f^*(\mathbf{c}) \triangleq \min_{\mathbf{w} \in \mathbb{R}^d} f(\mathbf{w}) \quad \text{subject to } h_j(\mathbf{w}) = c_j \quad (j = 1, \dots, p)
\label{eq:perturbed_opt_prob}
\end{equation}
where $\mathbf{c} = [c_1, \dots, c_p]^T \in \mathbb{R}^p$ represents a vector of constraint shifts.

Assume that at $\mathbf{c} = \mathbf{0}$, $\mathbf{w}^* = \mathbf{w}^*(\mathbf{0})$ is a local minimizer satisfying LICQ and the second-order sufficient condition ($\mathbf{v}^T \nabla_{\mathbf{w}}^2 \mathcal{L}(\mathbf{w}^*, \boldsymbol{\nu}^*) \mathbf{v} > 0$ for all non-zero $\mathbf{v} \in T_{\mathcal{M}}(\mathbf{w}^*)$).

Then the optimal value function $f^*(\mathbf{c})$ is continuously differentiable in an open neighborhood of $\mathbf{c} = \mathbf{0}$, and its gradient with respect to constraint perturbations is given exactly by the negative Lagrange multipliers:
\begin{equation}
\nabla_{\mathbf{c}} f^*(\mathbf{0}) = -\boldsymbol{\nu}^* \iff \frac{\partial f^*(\mathbf{0})}{\partial c_j} = -\nu_j^* \quad (j = 1, \dots, p)
\label{eq:lagrange_sensitivity_grad}
\end{equation}
\end{theorembox}

\begin{proof}
By the Implicit Function Theorem applied to the stationary Karush-Kuhn-Tucker system:
\begin{equation}
F(\mathbf{w}, \boldsymbol{\nu}, \mathbf{c}) = \begin{bmatrix} \nabla f(\mathbf{w}) + \sum_{j=1}^p \nu_j \nabla h_j(\mathbf{w}) \\ h_1(\mathbf{w}) - c_1 \\ \vdots \\ h_p(\mathbf{w}) - c_p \end{bmatrix} = \mathbf{0}
\end{equation}
the Jacobian matrix with respect to $(\mathbf{w}, \boldsymbol{\nu})$ is:
\begin{equation}
J_{(\mathbf{w}, \boldsymbol{\nu})} = \begin{bmatrix} \nabla_{\mathbf{w}}^2 \mathcal{L}(\mathbf{w}, \boldsymbol{\nu}) & \nabla \mathbf{h}(\mathbf{w}) \\ \nabla \mathbf{h}(\mathbf{w})^T & O \end{bmatrix}
\end{equation}
Under LICQ and the second-order sufficient condition, this bordered Hessian matrix is non-singular. Consequently, there exist unique, continuously differentiable local solution trajectories $\mathbf{w}^*(\mathbf{c})$ and $\boldsymbol{\nu}^*(\mathbf{c})$ satisfying $\mathbf{w}^*(\mathbf{0}) = \mathbf{w}^*$, $\boldsymbol{\nu}^*(\mathbf{0}) = \boldsymbol{\nu}^*$, and $h_j(\mathbf{w}^*(\mathbf{c})) = c_j$ identically.

Now evaluate the derivative of the optimal value function $f^*(\mathbf{c}) = f(\mathbf{w}^*(\mathbf{c}))$ with respect to coordinate $c_k$ via the multivariable chain rule:
\begin{equation}
\frac{\partial f^*(\mathbf{c})}{\partial c_k} = \sum_{i=1}^d \frac{\partial f(\mathbf{w}^*(\mathbf{c}))}{\partial w_i} \frac{\partial w_i^*(\mathbf{c})}{\partial c_k} = \nabla f(\mathbf{w}^*(\mathbf{c}))^T \frac{\partial \mathbf{w}^*(\mathbf{c})}{\partial c_k}
\end{equation}
Substitute the stationarity condition $\nabla f(\mathbf{w}^*(\mathbf{c})) = -\sum_{j=1}^p \nu_j^*(\mathbf{c}) \nabla h_j(\mathbf{w}^*(\mathbf{c}))$:
\begin{equation}
\frac{\partial f^*(\mathbf{c})}{\partial c_k} = -\sum_{j=1}^p \nu_j^*(\mathbf{c}) \left( \nabla h_j(\mathbf{w}^*(\mathbf{c}))^T \frac{\partial \mathbf{w}^*(\mathbf{c})}{\partial c_k} \right)
\label{eq:chain_rule_sub}
\end{equation}
Next, differentiate the identity $h_j(\mathbf{w}^*(\mathbf{c})) = c_j$ with respect to $c_k$ on both sides:
\begin{equation}
\frac{\partial}{\partial c_k} h_j(\mathbf{w}^*(\mathbf{c})) = \nabla h_j(\mathbf{w}^*(\mathbf{c}))^T \frac{\partial \mathbf{w}^*(\mathbf{c})}{\partial c_k} = \frac{\partial c_j}{\partial c_k} = \delta_{jk}
\end{equation}
where $\delta_{jk}$ is the Kronecker delta ($\delta_{jk} = 1$ if $j = k$, and $0$ otherwise).

Substituting this identity into~\eqref{eq:chain_rule_sub}:
\begin{equation}
\frac{\partial f^*(\mathbf{c})}{\partial c_k} = -\sum_{j=1}^p \nu_j^*(\mathbf{c}) \delta_{jk} = -\nu_k^*(\mathbf{c})
\end{equation}
Evaluating at $\mathbf{c} = \mathbf{0}$ yields the fundamental shadow price identity $\frac{\partial f^*(\mathbf{0})}{\partial c_k} = -\nu_k^*$.
\end{proof}

\begin{ponderbox}{What Do Lagrange Multipliers Actually Measure?}
What is the physical and economic significance of the optimal Lagrange multiplier $\nu_j^*$?
\tcblower
\textbf{The Socratic Resolution:} The Lagrange multiplier measures the \textbf{marginal sensitivity (shadow price)} of the optimal objective value with respect to perturbations of the constraint boundary. 

Suppose the constraint is perturbed by a small scalar shift: $h_j(\mathbf{w}) = \epsilon$. Let $\mathbf{w}^*(\epsilon)$ denote the perturbed optimal solution, and let $f^*(\epsilon) = f(\mathbf{w}^*(\epsilon))$ denote the perturbed optimal loss. Applying the chain rule:
\begin{equation}
\frac{d f^*(\epsilon)}{d\epsilon}\Big|_{\epsilon = 0} = \nabla f(\mathbf{w}^*)^T \frac{d\mathbf{w}^*}{d\epsilon}
\end{equation}
Differentiating the perturbed constraint $h_j(\mathbf{w}^*(\epsilon)) = \epsilon$ yields $\nabla h_j(\mathbf{w}^*)^T \frac{d\mathbf{w}^*}{d\epsilon} = 1$. Substituting the Lagrange stationarity condition $\nabla f(\mathbf{w}^*) = -\sum_k \nu_k^* \nabla h_k(\mathbf{w}^*)$:
\begin{equation}
\frac{d f^*(\epsilon)}{d\epsilon}\Big|_{\epsilon = 0} = -\sum_{k=1}^p \nu_k^* \nabla h_k(\mathbf{w}^*)^T \frac{d\mathbf{w}^*}{d\epsilon} = -\nu_j^*(1) = -\nu_j^*
\end{equation}
Thus, $-\nu_j^*$ is the exact marginal derivative of optimal loss with respect to constraint tightening. If $|\nu_j^*|$ is massive, shifting the constraint boundary slightly will severely degrade performance. If $\nu_j^* = 0$, the constraint has zero marginal impact on the optimal objective value.
\end{ponderbox}

\subsection{Convex Sets and Hyperplane Separation}

\begin{definitionbox}{Convex Set}
A subset $C \subseteq \mathbb{R}^d$ is \textbf{convex} if the line segment connecting any two arbitrary points in $C$ lies entirely within $C$:
\begin{equation}
\forall \mathbf{x}, \mathbf{y} \in C, \; \forall \theta \in [0, 1] \implies \theta \mathbf{x} + (1 - \theta)\mathbf{y} \in C
\label{eq:convex_set_def}
\end{equation}
A point of the form $\sum_{i=1}^k \theta_i \mathbf{x}_i$ with $\theta_i \ge 0$ and $\sum_{i=1}^k \theta_i = 1$ is a \textbf{convex combination}. A set is convex if and only if it contains all convex combinations of its elements.
\end{definitionbox}

\begin{theorembox}{Operations Preserving Set Convexity}
Convexity of sets is invariant under the following fundamental mathematical operations:
\begin{enumerate}
    \item \textbf{Arbitrary Intersections:} If $C_\alpha$ is a convex set for every index $\alpha \in \mathcal{A}$ (including uncountably infinite index collections), then their intersection is convex:
    \begin{equation}
    C = \bigcap_{\alpha \in \mathcal{A}} C_\alpha \quad \text{is convex.}
    \end{equation}
    Consequently, polyhedra defined as intersections of finitely many halfspaces $\{\mathbf{x} \mid A\mathbf{x} \le \mathbf{b}\}$ are always convex sets.
    \item \textbf{Affine Transformations:} If $C \subseteq \mathbb{R}^n$ is convex and $f: \mathbb{R}^n \to \mathbb{R}^m$ is an affine map $f(\mathbf{x}) = A\mathbf{x} + \mathbf{b}$, then the image $f(C) = \{A\mathbf{x} + \mathbf{b} \mid \mathbf{x} \in C\}$ is convex. Furthermore, the inverse image (pre-image) $f^{-1}(C) = \{\mathbf{x} \mid A\mathbf{x} + \mathbf{b} \in C\}$ is convex.
    \item \textbf{Cartesian Products:} If $C_1 \subseteq \mathbb{R}^n$ and $C_2 \subseteq \mathbb{R}^m$ are convex sets, their Cartesian product:
    \begin{equation}
    C_1 \times C_2 = \{(\mathbf{x}_1, \mathbf{x}_2) \mid \mathbf{x}_1 \in C_1, \; \mathbf{x}_2 \in C_2\} \subseteq \mathbb{R}^{n+m} \quad \text{is convex.}
    \end{equation}
\end{enumerate}
\end{theorembox}

\begin{theorembox}{Separating Hyperplane Theorem}
Let $C_1, C_2 \subset \mathbb{R}^d$ be two non-empty, disjoint convex sets ($C_1 \cap C_2 = \emptyset$). 
Then there exists a non-zero normal vector $\mathbf{a} \in \mathbb{R}^d \setminus \{\mathbf{0}\}$ and a scalar threshold $b \in \mathbb{R}$ such that:
\begin{equation}
\mathbf{a}^T \mathbf{x} \le b \le \mathbf{a}^T \mathbf{y} \quad \forall \mathbf{x} \in C_1, \; \forall \mathbf{y} \in C_2
\label{eq:separating_hyperplane}
\end{equation}
The affine hyperplane $\mathcal{H} = \{\mathbf{z} \in \mathbb{R}^d \mid \mathbf{a}^T \mathbf{z} = b\}$ separates $C_1$ and $C_2$. If $C_1$ is closed and $C_2$ is compact (or a singleton outside $C_1$), the separation is \textbf{strict}: $\mathbf{a}^T \mathbf{x} < b < \mathbf{a}^T \mathbf{y}$.
\end{theorembox}

\begin{ponderbox}{Why Intersection Preserves Convexity While Union Destroys It}
Why does an intersection of arbitrarily many convex sets always remain convex, while the union of even two convex sets almost never remains convex?
\tcblower
\textbf{The Socratic Resolution:} Convexity is an \textbf{internal closure property} across lines. 
If $\mathbf{x}, \mathbf{y} \in \bigcap_{\alpha} C_\alpha$, then $\mathbf{x}$ and $\mathbf{y}$ belong to every single set $C_\alpha$. Because each individual $C_\alpha$ is convex, the entire line segment $[\mathbf{x}, \mathbf{y}] = \{\theta \mathbf{x} + (1-\theta)\mathbf{y} \mid \theta \in [0, 1]\}$ lies inside $C_\alpha$ for every $\alpha$. Therefore, the entire segment lies in the intersection: $[\mathbf{x}, \mathbf{y}] \subseteq \bigcap_\alpha C_\alpha$.

Conversely, consider the union of two disjoint convex sets $C_1 \cup C_2$. Pick $\mathbf{x} \in C_1$ and $\mathbf{y} \in C_2$. The line segment connecting $\mathbf{x}$ to $\mathbf{y}$ must physically traverse the empty void separating the two sets. Points along the segment belong to neither $C_1$ nor $C_2$, instantly violating the convexity condition.
\end{ponderbox}

\subsection{Convex Functions}

\begin{definitionbox}{Convex Function}
A function $f: C \to \mathbb{R}$ defined on a convex domain $C \subseteq \mathbb{R}^d$ is \textbf{convex} if for all $\mathbf{x}, \mathbf{y} \in C$ and all $\theta \in [0, 1]$:
\begin{equation}
f(\theta \mathbf{x} + (1 - \theta)\mathbf{y}) \le \theta f(\mathbf{x}) + (1 - \theta) f(\mathbf{y})
\label{eq:convex_func_def}
\end{equation}
A function is \textbf{strictly convex} if the inequality is strict for all $\mathbf{x} \ne \mathbf{y}$ and $\theta \in (0, 1)$. Geometrically, convexity states that the secant chord connecting $( \mathbf{x}, f(\mathbf{x}) )$ and $( \mathbf{y}, f(\mathbf{y}) )$ lies everywhere on or above the graph of the function.
\end{definitionbox}

\begin{definitionbox}{Epigraph Characterization}
The \textbf{epigraph} of a function $f: C \to \mathbb{R}$ is the subset of $\mathbb{R}^{d+1}$ lying on or above its graph:
\begin{equation}
\operatorname{epi}(f) \triangleq \left\{ (\mathbf{x}, t) \in \mathbb{R}^d \times \mathbb{R} \;\middle|\; \mathbf{x} \in C, \; f(\mathbf{x}) \le t \right\}
\label{eq:epigraph_def}
\end{equation}
\textbf{The Fundamental Bridge:} A function $f$ is convex if and only if its epigraph $\operatorname{epi}(f)$ is a convex set in $\mathbb{R}^{d+1}$.
\end{definitionbox}

\begin{theorembox}{First-Order Characterization of Convexity}
Let $f: C \to \mathbb{R}$ be continuously differentiable on an open convex set $C$. Then $f$ is convex if and only if for all $\mathbf{x}, \mathbf{y} \in C$:
\begin{equation}
f(\mathbf{y}) \ge f(\mathbf{x}) + \nabla f(\mathbf{x})^T (\mathbf{y} - \mathbf{x})
\label{eq:first_order_convex}
\end{equation}
Geometrically, the first-order Taylor approximation forms a global supporting hyperplane that lies everywhere below or touches the graph of $f$.
\end{theorembox}

\begin{proof}
Assume $f$ is convex. For any $\theta \in (0, 1]$:
\begin{equation}
f(\mathbf{x} + \theta(\mathbf{y} - \mathbf{x})) = f((1 - \theta)\mathbf{x} + \theta \mathbf{y}) \le (1 - \theta)f(\mathbf{x}) + \theta f(\mathbf{y})
\end{equation}
Subtracting $f(\mathbf{x})$ and dividing by $\theta$:
\begin{equation}
\frac{f(\mathbf{x} + \theta(\mathbf{y} - \mathbf{x})) - f(\mathbf{x})}{\theta} \le f(\mathbf{y}) - f(\mathbf{x})
\end{equation}
Taking the limit as $\theta \to 0^+$, the left-hand side converges to the directional derivative $\nabla f(\mathbf{x})^T (\mathbf{y} - \mathbf{x})$, yielding $f(\mathbf{y}) \ge f(\mathbf{x}) + \nabla f(\mathbf{x})^T (\mathbf{y} - \mathbf{x})$.

Conversely, suppose~\eqref{eq:first_order_convex} holds for all points. Let $\mathbf{z} = \theta \mathbf{x} + (1 - \theta)\mathbf{y}$ for $\theta \in [0, 1]$. Applying~\eqref{eq:first_order_convex} at $\mathbf{z}$ with respect to $\mathbf{x}$ and $\mathbf{y}$:
\begin{align}
f(\mathbf{x}) &\ge f(\mathbf{z}) + \nabla f(\mathbf{z})^T (\mathbf{x} - \mathbf{z}) \\
f(\mathbf{y}) &\ge f(\mathbf{z}) + \nabla f(\mathbf{z})^T (\mathbf{y} - \mathbf{z})
\end{align}
Multiplying the first inequality by $\theta$, the second by $(1 - \theta)$, and adding:
\begin{align}
\theta f(\mathbf{x}) + (1 - \theta)f(\mathbf{y}) &\ge f(\mathbf{z}) + \nabla f(\mathbf{z})^T \left( \theta(\mathbf{x} - \mathbf{z}) + (1 - \theta)(\mathbf{y} - \mathbf{z}) \right) \nonumber \\
&= f(\mathbf{z}) + \nabla f(\mathbf{z})^T (\mathbf{0}) = f(\theta \mathbf{x} + (1 - \theta)\mathbf{y})
\end{align}
which proves that $f$ is convex.
\end{proof}

\begin{theorembox}{Second-Order Characterization of Convexity}
Let $f: C \to \mathbb{R}$ be twice continuously differentiable on an open convex set $C \subseteq \mathbb{R}^d$. Then:
\begin{itemize}
    \item $f$ is convex if and only if its Hessian matrix is positive semidefinite everywhere:
    \begin{equation}
    \nabla^2 f(\mathbf{x}) \succeq 0 \quad \forall \mathbf{x} \in C
    \end{equation}
    \item If the Hessian is strictly positive definite everywhere ($\nabla^2 f(\mathbf{x}) \succ 0$ for all $\mathbf{x} \in C$), then $f$ is \textbf{strictly convex}.
\end{itemize}
\end{theorembox}

\begin{theorembox}{The Strong Convexity \& Smoothness Quadratic Sandwich}
Let $f: \mathbb{R}^d \to \mathbb{R}$ be twice continuously differentiable.
\begin{enumerate}
    \item $f$ is \textbf{$m$-strongly convex} ($m > 0$) if and only if $\nabla^2 f(\mathbf{x}) \succeq m I$ for all $\mathbf{x} \in \mathbb{R}^d$. This is equivalent to the global quadratic under-estimator:
    \begin{equation}
    f(\mathbf{y}) \ge f(\mathbf{x}) + \nabla f(\mathbf{x})^T(\mathbf{y} - \mathbf{x}) + \frac{m}{2}\|\mathbf{y} - \mathbf{x}\|_2^2 \quad \forall \mathbf{x}, \mathbf{y} \in \mathbb{R}^d
    \label{eq:strong_convexity_lower}
    \end{equation}
    \item $f$ is \textbf{$L$-smooth} ($L > 0$) if and only if its gradient is $L$-Lipschitz continuous ($\|\nabla f(\mathbf{x}) - \nabla f(\mathbf{y})\|_2 \le L \|\mathbf{x} - \mathbf{y}\|_2$), which is equivalent to $\nabla^2 f(\mathbf{x}) \preceq L I$ for all $\mathbf{x} \in \mathbb{R}^d$. This implies the global quadratic over-estimator (Descent Lemma):
    \begin{equation}
    f(\mathbf{y}) \le f(\mathbf{x}) + \nabla f(\mathbf{x})^T(\mathbf{y} - \mathbf{x}) + \frac{L}{2}\|\mathbf{y} - \mathbf{x}\|_2^2 \quad \forall \mathbf{x}, \mathbf{y} \in \mathbb{R}^d
    \label{eq:smoothness_upper}
    \end{equation}
\end{enumerate}
Combining both bounds, any $m$-strongly convex and $L$-smooth functional is strictly sandwiched between two quadratic paraboloids of curvature $m$ and $L$:
\begin{equation}
f(\mathbf{x}) + \nabla f(\mathbf{x})^T(\mathbf{y} - \mathbf{x}) + \frac{m}{2}\|\mathbf{y} - \mathbf{x}\|_2^2 \le f(\mathbf{y}) \le f(\mathbf{x}) + \nabla f(\mathbf{x})^T(\mathbf{y} - \mathbf{x}) + \frac{L}{2}\|\mathbf{y} - \mathbf{x}\|_2^2
\label{eq:quadratic_sandwich}
\end{equation}
The ratio $\kappa \triangleq \frac{L}{m} \ge 1$ is the \textbf{condition number} of the objective functional, dictating the geometric eccentricity of all sublevel sets and governing the linear convergence rate $(1 - 1/\kappa)$ of both Gradient Descent and Projected Gradient Descent.
\end{theorembox}

\begin{proof}
By Taylor's theorem with integral remainder applied to the scalar function $g(t) = f(\mathbf{x} + t(\mathbf{y} - \mathbf{x}))$ for $t \in [0, 1]$:
\begin{equation}
f(\mathbf{y}) = f(\mathbf{x}) + \nabla f(\mathbf{x})^T(\mathbf{y} - \mathbf{x}) + \int_0^1 (1 - t) (\mathbf{y} - \mathbf{x})^T \nabla^2 f(\mathbf{x} + t(\mathbf{y} - \mathbf{x})) (\mathbf{y} - \mathbf{x}) \, dt
\end{equation}
If $\nabla^2 f(\mathbf{z}) \succeq m I$ for all $\mathbf{z}$, then $(\mathbf{y} - \mathbf{x})^T \nabla^2 f(\mathbf{z}) (\mathbf{y} - \mathbf{x}) \ge m \|\mathbf{y} - \mathbf{x}\|_2^2$. Integrating:
\begin{equation}
\int_0^1 (1 - t) m \|\mathbf{y} - \mathbf{x}\|_2^2 \, dt = m \|\mathbf{y} - \mathbf{x}\|_2^2 \left[ t - \frac{t^2}{2} \right]_0^1 = \frac{m}{2} \|\mathbf{y} - \mathbf{x}\|_2^2
\end{equation}
which proves~\eqref{eq:strong_convexity_lower}. Applying the identical integration to the upper bound $(\mathbf{y} - \mathbf{x})^T \nabla^2 f(\mathbf{z}) (\mathbf{y} - \mathbf{x}) \le L \|\mathbf{y} - \mathbf{x}\|_2^2$ yields~\eqref{eq:smoothness_upper}.
\end{proof}

\begin{ponderbox}{Convex Sets vs.\ Convex Functions: The Quasiconvexity Ambiguity}
Every convex function has convex sublevel sets $\mathcal{S}_\alpha = \{\mathbf{x} \mid f(\mathbf{x}) \le \alpha\}$. Does having convex sublevel sets guarantee that a function is convex?
\tcblower
\textbf{The Socratic Resolution:} No! Functions whose sublevel sets are all convex are called \textbf{quasiconvex}, and quasiconvexity is strictly weaker than convexity. 

Consider the one-dimensional function $f(x) = \sqrt{|x|}$. For any $\alpha \ge 0$, the sublevel set is $\mathcal{S}_\alpha = [-\alpha^2, \alpha^2]$, which is an interval and therefore a convex set. Yet $f(x)$ is blatantly non-convex: between $x = 0$ and $x = 4$, the secant chord lies \emph{below} the function curve ($f(2) = \sqrt{2} \approx 1.414$, while the chord midpoint is $\frac{1}{2}(0 + 2) = 1.0 < 1.414$).

Similarly, the Gaussian probability density function $f(x) = e^{-x^2}$ has convex upper level sets (it is quasiconcave), but fails standard concavity. Convexity requires the epigraph in $\mathbb{R}^{d+1}$ to be convex; having convex horizontal cross-sections (sublevel sets) does not ensure that the vertical profile satisfies the secant inequality.
\end{ponderbox}

% ----------------------------------------------------------------------
% PART 3: VISUAL & GEOMETRIC ANCHORING
% ----------------------------------------------------------------------
\section{Visual \& Geometric Anchoring}

Optimization geometry clarifies mathematical mechanics. In this section, we anchor the theoretical conditions in three native TikZ visualizations.

\subsection{The Geometry of Lagrange Multipliers and Tangency}
Figure~\ref{fig:constrained_geometry} contrasts the true constrained minimum with the naive orthogonal projection of the unconstrained minimum, demonstrating why constrained extrema are points of contour tangency.

\begin{figure}[htbp]
\centering
\begin{tikzpicture}[scale=1.0, font=\sffamily\footnotesize]
    % Axes
    \draw[->, thick, color=primarydark!50] (-0.8, 0) -- (4.2, 0) node[right] {$x$};
    \draw[->, thick, color=primarydark!50] (0, -0.8) -- (0, 4.2) node[above] {$y$};

    % Objective Contours: f(x, y) = 2x^2 + y^2 = c => x^2 / (c/2) + y^2 / c = 1
    % c = 2: rx = 1.0, ry = 1.414
    \draw[thick, color=primarydark!35] (0, 0) ellipse (1.0cm and 1.414cm);
    \node[right, color=primarydark!40] at (1.0, 0.4) {\scriptsize $f = 2$};

    % c = 6: rx = 1.732, ry = 2.449 (TANGENT CONTOUR!)
    \draw[very thick, color=primarydark!85] (0, 0) ellipse (1.732cm and 2.449cm);
    \node[right, color=primarydark!90] at (1.75, 1.2) {\textbf{Optimal Contour} $f = 6.0$};

    % c = 6.75: rx = 1.837, ry = 2.598 (SUBOPTIMAL PROJECTION CONTOUR)
    \draw[thick, dashed, color=traprad!65] (0, 0) ellipse (1.837cm and 2.598cm);
    \node[right, color=traprad!80] at (1.85, 2.1) {\scriptsize Naive Contour $f = 6.75$};

    % Constraint Line: x + y = 3
    \draw[line width=1.8pt, color=tealborder] (-0.3, 3.3) -- (3.3, -0.3);
    \node[above right, color=tealborder] at (2.4, 0.7) {\textbf{Feasible Constraint Line } $x + y = 3$};

    % Unconstrained minimum
    \filldraw[color=primarydark] (0, 0) circle (2.5pt) node[below left] {$\mathbf{w}_{\text{unc}}^* = (0, 0)$};

    % Naive Euclidean Projection: (1.5, 1.5)
    \coordinate (Pnaive) at (1.5, 1.5);
    \filldraw[color=traprad] (Pnaive) circle (3pt);
    \draw[dashed, color=traprad, thick] (0, 0) -- (Pnaive);
    \node[above right, color=traprad, align=left] at (1.55, 1.55) {
        \textbf{Naive Projection } $\Pi_C(\mathbf{0}) = (1.5, 1.5)$ \\
        \scriptsize $f = 6.75$ ($+12.5\%$ suboptimality)
    };

    % True Constrained Minimum: (1.0, 2.0)
    \coordinate (Wstar) at (1.0, 2.0);
    \filldraw[color=forest] (Wstar) circle (3.5pt);
    \node[above left, color=forest, align=right] at (0.95, 2.05) {
        \textbf{True Constrained Min } $\mathbf{w}_C^* = (1, 2)$ \\
        \scriptsize $f^* = 6.00$ (\textbf{Exact Tangency})
    };

    % Collinear gradient arrows at (1, 2)
    % grad f(1, 2) = [4, 4]^T => scale down by 0.25 => [1.0, 1.0]
    \draw[->, >=stealth, line width=2pt, color=coralborder] (Wstar) -- ($(Wstar) + (0.8, 0.8)$) 
        node[right] {$\nabla f = [4, 4]^T$};
    % grad h(1, 2) = [1, 1]^T => scale down by 0.6 => [0.6, 0.6]
    \draw[->, >=stealth, line width=1.6pt, color=tealborder] (Wstar) -- ($(Wstar) + (-0.6, -0.6)$) 
        node[below left] {$\nu^* \nabla h = -4 [1, 1]^T$};
\end{tikzpicture}
\caption{The Geometry of Lagrange Multipliers: Curvature dictates that constrained optima are points of contour tangency, not orthogonal projections. The unconstrained minimum is at $(0, 0)$. The naive Euclidean projection of $(0, 0)$ onto the constraint line $x + y = 3$ is $(1.5, 1.5)$, yielding a suboptimal loss of $f = 6.75$. The true constrained minimum occurs at $(1, 2)$, where the objective ellipse $f = 6.00$ is tangent to the line, aligning the gradient vectors $\nabla f$ and $\nabla h$ co-linearly.}
\label{fig:constrained_geometry}
\end{figure}

\subsection{Convex Sets and the Separating Hyperplane Theorem}
Figure~\ref{fig:convex_sets_separation} diagrams the internal chord test for set convexity and illustrates the geometric separation of disjoint convex bodies by an affine hyperplane.

\begin{figure}[htbp]
\centering
\begin{tikzpicture}[scale=0.9, font=\sffamily\footnotesize]
    % Panel 1: Convex Set
    \begin{scope}[shift={(-5.0, 0)}]
        \draw[thick, fill=tealborder!15, color=tealborder] 
            (-1.5, -1.0) to[out=90, in=180] (0, 1.5)
                         to[out=0, in=90] (1.8, 0.2)
                         to[out=-90, in=0] (0.8, -1.5)
                         to[out=180, in=-90] cycle;
        \coordinate (X1) at (-0.8, -0.3);
        \coordinate (Y1) at (1.0, 0.6);
        \draw[very thick, color=forest] (X1) -- (Y1);
        \filldraw[color=forest] (X1) circle (2.5pt) node[below left] {$\mathbf{x}$};
        \filldraw[color=forest] (Y1) circle (2.5pt) node[above right] {$\mathbf{y}$};
        \node[above, color=forest] at (0.1, 0.2) {\scriptsize $[\mathbf{x}, \mathbf{y}] \subset C$};
        \node[below, align=center] at (0, -2.0) {\textbf{(a) Convex Set} \\ All secant segments contained};
    \end{scope}

    % Panel 2: Non-Convex Set (Crescent)
    \begin{scope}[shift={(0, 0)}]
        \draw[thick, fill=traprad!15, color=traprad] 
            (-1.5, -1.2) to[out=80, in=180] (0, 1.5)
                         to[out=0, in=100] (1.5, -1.2)
                         to[out=120, in=0] (0, 0.2)
                         to[out=180, in=60] cycle;
        \coordinate (X2) at (-1.0, 0.2);
        \coordinate (Y2) at (1.0, 0.2);
        \draw[very thick, dashed, color=traprad] (X2) -- (Y2);
        \filldraw[color=traprad] (X2) circle (2.5pt) node[left] {$\mathbf{x}$};
        \filldraw[color=traprad] (Y2) circle (2.5pt) node[right] {$\mathbf{y}$};
        \node[above, color=traprad] at (0, 0.3) {\scriptsize Segment escapes!};
        \node[below, align=center] at (0, -2.0) {\textbf{(b) Non-Convex Set} \\ Segment exits the set};
    \end{scope}

    % Panel 3: Separating Hyperplane
    \begin{scope}[shift={(5.2, 0)}]
        % Set C1
        \draw[thick, fill=primarydark!15, color=primarydark] 
            (-1.8, -0.8) circle (0.9cm);
        \node at (-1.8, -0.8) {$C_1$};

        % Set C2
        \draw[thick, fill=forest!15, color=forest] 
            (1.5, 0.5) ellipse (1.1cm and 0.7cm);
        \node at (1.5, 0.5) {$C_2$};

        % Separating Hyperplane line
        \draw[very thick, color=coralborder] (-0.3, -2.0) -- (0.5, 2.0) 
            node[above] {\textbf{Hyperplane } $\mathbf{a}^T\mathbf{x} = b$};
        
        % Normal vector a
        \draw[->, >=stealth, line width=1.5pt, color=coralborder] (0.1, 0) -- (1.1, -0.2) 
            node[below right] {$\mathbf{a}$};

        \node[below, align=center] at (0, -2.0) {\textbf{(c) Separating Hyperplane} \\ Disjoint convex sets};
    \end{scope}
\end{tikzpicture}
\caption{Convex Geometry and the Separating Hyperplane Theorem: (a) In a convex set, the straight segment between any two points $\mathbf{x}, \mathbf{y}$ remains entirely enclosed within the set. (b) In a non-convex set (e.g., an indented crescent), chords cross through the external void. (c) The Separating Hyperplane Theorem guarantees that any two non-empty disjoint convex sets $C_1, C_2$ can be completely separated by an affine hyperplane $\mathbf{a}^T\mathbf{x} = b$.}
\label{fig:convex_sets_separation}
\end{figure}

\subsection{First-Order Supporting Hyperplanes and Epigraph Geometry}
Figure~\ref{fig:epigraph_supporting} illustrates the equivalence between function convexity, epigraph convexity, and the global supporting hyperplane property.

\begin{figure}[htbp]
\centering
\begin{tikzpicture}[scale=1.1, font=\sffamily\footnotesize]
    % Axes
    \draw[->, thick, color=primarydark!60] (-3.2, 0) -- (3.2, 0) node[right] {$x$};
    \draw[->, thick, color=primarydark!60] (0, -0.5) -- (0, 4.5) node[above] {$t$};

    % Shaded Epigraph region: t >= 0.4*x^2 + 0.8
    \fill[tealborder!15, domain=-2.8:2.8] 
        (-2.8, 4.2) -- plot (\x, {0.4*\x*\x + 0.8}) -- (2.8, 4.2) -- cycle;
    \node[color=tealborder!90!black] at (1.4, 3.2) {\large \textbf{Epigraph } $\operatorname{epi}(f)$};

    % Convex function curve f(x) = 0.4*x^2 + 0.8
    \draw[line width=1.8pt, color=primarydark, domain=-2.8:2.8] 
        plot (\x, {0.4*\x*\x + 0.8}) node[right] {$t = f(x)$};

    % Base point x0 = -1.0 => f(-1.0) = 0.4(1) + 0.8 = 1.2
    % f'(x) = 0.8*x => f'(-1.0) = -0.8
    % Tangent line: t = 1.2 - 0.8*(x - (-1.0)) = 1.2 - 0.8(x + 1) = 0.4 - 0.8*x
    \coordinate (X0) at (-1.0, 1.2);
    \filldraw[color=forest] (X0) circle (2.5pt) node[above left=2pt] {$(x_0, f(x_0))$};

    % Supporting Hyperplane (Tangent line)
    \draw[very thick, color=coralborder, domain=-3.0:2.5] 
        plot (\x, {1.2 - 0.8*(\x + 1.0)}) 
        node[below right] {\textbf{Supporting Hyperplane} $t = f(x_0) + f'(x_0)(x - x_0)$};

    % Normal vector to supporting hyperplane: [f'(x0), -1]^T = [-0.8, -1]^T pointing downward
    % or [-f'(x0), 1]^T pointing into epigraph
    \draw[->, >=stealth, line width=1.5pt, color=forest] (X0) -- ($(X0) + (0.8, 1.0)$) 
        node[above] {\scriptsize Inward normal $(-\nabla f(x_0), 1)$};

    % Annotation demonstrating global under-estimation
    \draw[<->, >=stealth, thick, color=traprad] (1.5, {0.4*(1.5)*(1.5) + 0.8}) -- (1.5, {1.2 - 0.8*(1.5 + 1.0)});
    \node[right, color=traprad, align=left] at (1.55, 0.4) {
        \scriptsize Gap $\ge 0$ \\
        \scriptsize $f(y) \ge f(x_0) + \nabla f(x_0)^T(y - x_0)$
    };
\end{tikzpicture}
\caption{First-Order Characterization and Epigraph Geometry: A functional $f$ is convex if and only if its epigraph $\operatorname{epi}(f)$ is a convex set in $\mathbb{R}^{d+1}$. Consequently, at every point $x_0$, the affine first-order Taylor expansion defines a global supporting hyperplane that under-estimates the functional everywhere across its domain.}
\label{fig:epigraph_supporting}
\end{figure}

\subsection{The Geometry of Water-Filling and Resource Allocation}
Figure~\ref{fig:waterfilling_geometry} visualizes the fundamental \textbf{Water-Filling Principle} originating in Shannon information theory and convex resource allocation. When allocating a finite total resource (such as transmit power $P_{\text{total}}$ or portfolio capital) across independent channels subjected to varying noise levels $\sigma_i^2$, the optimal Lagrange multiplier establishes a global equilibrium threshold $\mu^* \triangleq 1/\nu^*$. 

\begin{figure}[htbp]
\centering
\begin{tikzpicture}[scale=1.0, font=\sffamily\footnotesize]
    % Channel 1 basin: x in [0.5, 2.5], floor at y = 0.4
    \draw[thick, fill=gray!20] (0.5, 0.4) rectangle (2.5, 0);
    \draw[thick] (0.5, 3.5) -- (0.5, 0.4) -- (2.5, 0.4) -- (2.5, 3.5);
    \node at (1.5, 0.2) {\textbf{Noise } $\sigma_1^2 = 0.2$};
    
    % Channel 2 basin: x in [3.5, 5.5], floor at y = 1.2
    \draw[thick, fill=gray!20] (3.5, 1.2) rectangle (5.5, 0);
    \draw[thick] (3.5, 3.5) -- (3.5, 1.2) -- (5.5, 1.2) -- (5.5, 3.5);
    \node at (4.5, 0.6) {\textbf{Noise } $\sigma_2^2 = 0.6$};
    
    % Channel 3 basin: x in [6.5, 8.5], floor at y = 3.0
    \draw[thick, fill=gray!20] (6.5, 3.0) rectangle (8.5, 0);
    \draw[thick] (6.5, 3.5) -- (6.5, 3.0) -- (8.5, 3.0) -- (8.5, 3.5);
    \node at (7.5, 1.5) {\textbf{Noise } $\sigma_3^2 = 1.5$};

    % Water level mu* = 1.40 => y = 2.8
    \draw[very thick, dashed, color=coralborder] (-0.2, 2.8) -- (9.0, 2.8) 
        node[right, color=coralborder] {\textbf{Water Level } $\mu^* = \frac{1}{\nu^*} = 1.40$};

    % Water in Channel 1 (height 2.4, y from 0.4 to 2.8)
    \fill[tealborder!35] (0.52, 0.42) rectangle (2.48, 2.78);
    \node[color=primarydark, font=\bfseries] at (1.5, 1.6) {$P_1^* = 1.20$};

    % Water in Channel 2 (height 1.6, y from 1.2 to 2.8)
    \fill[tealborder!35] (3.52, 1.22) rectangle (5.48, 2.78);
    \node[color=primarydark, font=\bfseries] at (4.5, 2.0) {$P_2^* = 0.80$};

    % Channel 3 has floor at 3.0 > 2.8 => DRY!
    \node[color=traprad, font=\bfseries, align=center] at (7.5, 3.25) {\textsc{Dry Basin} \\ $P_3^* = 0.0$};

    % Labels
    \node[below=4pt] at (1.5, 0) {\textbf{Channel 1} (Clean)};
    \node[below=4pt] at (4.5, 0) {\textbf{Channel 2} (Moderate)};
    \node[below=4pt] at (7.5, 0) {\textbf{Channel 3} (Severe Noise)};
\end{tikzpicture}
\caption{The Water-Filling Principle for Parallel Gaussian Channels: Transmit power is poured into independent communication basins whose floor elevations represent noise power spectral densities $\sigma_i^2$. Channels with noise strictly below the equilibrium water level $\mu^*$ receive active transmit power $P_i^* = \mu^* - \sigma_i^2$, while channels with excessive noise ($\sigma_3^2 \ge \mu^*$) remain completely dry ($P_3^* = 0$), mathematically demonstrating how inequality constraints prune suboptimal dimensions.}
\label{fig:waterfilling_geometry}
\end{figure}

% ----------------------------------------------------------------------
% PART 4: DEEP INTUITION & CONCEPTUAL DERIVATIONS
% ----------------------------------------------------------------------
\section{Deep Intuition \& Conceptual Derivations}

\subsection{Projected Gradient Descent (PGD) and Non-Expansiveness}
When minimizing a differentiable convex loss functional $f: \mathbb{R}^d \to \mathbb{R}$ over a closed convex constraint set $C \subset \mathbb{R}^d$:
\begin{equation}
\min_{\mathbf{w} \in C} f(\mathbf{w})
\end{equation}
the \textbf{Projected Gradient Descent (PGD)} algorithm interleaves standard unconstrained gradient steps with orthogonal Euclidean projections back onto $C$:
\begin{equation}
\mathbf{w}_{k+1} = \Pi_C\left( \mathbf{w}_k - \eta \nabla f(\mathbf{w}_k) \right)
\label{eq:pgd_update}
\end{equation}

\begin{theorembox}{Variational Characterization of Euclidean Projection}
Let $C \subseteq \mathbb{R}^d$ be a non-empty closed convex set, and let $\mathbf{w} \in \mathbb{R}^d$. A point $\mathbf{z}^* \in C$ is the unique Euclidean projection of $\mathbf{w}$ onto $C$ ($\mathbf{z}^* = \Pi_C(\mathbf{w}) = \argmin_{\mathbf{z} \in C} \frac{1}{2}\|\mathbf{z} - \mathbf{w}\|_2^2$) if and only if:
\begin{equation}
\langle \mathbf{w} - \mathbf{z}^*, \mathbf{x} - \mathbf{z}^* \rangle \le 0 \quad \forall \mathbf{x} \in C
\label{eq:proj_variational_char}
\end{equation}
Geometrically, this obtuse angle condition ensures that moving from the projection $\mathbf{z}^*$ toward any other feasible point $\mathbf{x} \in C$ forms an angle $\ge 90^\circ$ with the displacement vector $\mathbf{w} - \mathbf{z}^*$.
\end{theorembox}

\begin{proof}
\textbf{Necessity ($\implies$):} Suppose $\mathbf{z}^* = \Pi_C(\mathbf{w})$. For any arbitrary point $\mathbf{x} \in C$, consider the convex combination $\mathbf{z}(t) = (1 - t)\mathbf{z}^* + t \mathbf{x} = \mathbf{z}^* + t(\mathbf{x} - \mathbf{z}^*)$ for $t \in [0, 1]$. Because $C$ is convex, $\mathbf{z}(t) \in C$ for all $t \in [0, 1]$. Define the scalar quadratic objective:
\begin{equation}
\phi(t) \triangleq \frac{1}{2}\|\mathbf{z}(t) - \mathbf{w}\|_2^2 = \frac{1}{2}\|\mathbf{z}^* + t(\mathbf{x} - \mathbf{z}^*) - \mathbf{w}\|_2^2
\end{equation}
Since $\mathbf{z}^* = \mathbf{z}(0)$ minimizes the distance over $C$, the one-sided directional derivative at $t = 0$ must be non-negative: $\phi'(0) \ge 0$. Differentiating with respect to $t$:
\begin{equation}
\phi'(t) = \langle \mathbf{z}(t) - \mathbf{w}, \mathbf{x} - \mathbf{z}^* \rangle \implies \phi'(0) = \langle \mathbf{z}^* - \mathbf{w}, \mathbf{x} - \mathbf{z}^* \rangle \ge 0
\end{equation}
Multiplying by $-1$ reverses the inequality:
\begin{equation}
\langle \mathbf{w} - \mathbf{z}^*, \mathbf{x} - \mathbf{z}^* \rangle \le 0
\end{equation}
\textbf{Sufficiency ($\impliedby$):} Suppose $\mathbf{z}^* \in C$ satisfies~\eqref{eq:proj_variational_char} for all $\mathbf{x} \in C$. For any feasible point $\mathbf{x} \in C$ with $\mathbf{x} \ne \mathbf{z}^*$, expand the squared distance:
\begin{align}
\|\mathbf{x} - \mathbf{w}\|_2^2 &= \|(\mathbf{x} - \mathbf{z}^*) + (\mathbf{z}^* - \mathbf{w})\|_2^2 \nonumber \\
&= \|\mathbf{x} - \mathbf{z}^*\|_2^2 + \|\mathbf{z}^* - \mathbf{w}\|_2^2 + 2\langle \mathbf{z}^* - \mathbf{w}, \mathbf{x} - \mathbf{z}^* \rangle \nonumber \\
&= \|\mathbf{x} - \mathbf{z}^*\|_2^2 + \|\mathbf{z}^* - \mathbf{w}\|_2^2 - 2\langle \mathbf{w} - \mathbf{z}^*, \mathbf{x} - \mathbf{z}^* \rangle
\end{align}
Because $-2\langle \mathbf{w} - \mathbf{z}^*, \mathbf{x} - \mathbf{z}^* \rangle \ge 0$ by~\eqref{eq:proj_variational_char} and $\|\mathbf{x} - \mathbf{z}^*\|_2^2 > 0$ for $\mathbf{x} \ne \mathbf{z}^*$:
\begin{equation}
\|\mathbf{x} - \mathbf{w}\|_2^2 > \|\mathbf{z}^* - \mathbf{w}\|_2^2
\end{equation}
Thus, $\mathbf{z}^*$ is the unique global minimizer of Euclidean distance over $C$.
\end{proof}

\begin{theorembox}{Firm Non-Expansiveness of the Convex Projection Operator}
Let $C \subseteq \mathbb{R}^d$ be a closed convex set. The Euclidean projection operator $\Pi_C: \mathbb{R}^d \to C$ is \textbf{firmly non-expansive}:
\begin{equation}
\|\Pi_C(\mathbf{u}) - \Pi_C(\mathbf{v})\|_2^2 + \|(I - \Pi_C)(\mathbf{u}) - (I - \Pi_C)(\mathbf{v})\|_2^2 \le \|\mathbf{u} - \mathbf{v}\|_2^2 \quad \forall \mathbf{u}, \mathbf{v} \in \mathbb{R}^d
\label{eq:firm_non_expansive}
\end{equation}
Consequently, $\Pi_C$ is strictly non-expansive ($\|\Pi_C(\mathbf{u}) - \Pi_C(\mathbf{v})\|_2 \le \|\mathbf{u} - \mathbf{v}\|_2$).
\end{theorembox}

\begin{proof}
Let $\hat{\mathbf{u}} = \Pi_C(\mathbf{u})$ and $\hat{\mathbf{v}} = \Pi_C(\mathbf{v})$. Applying the variational characterization~\eqref{eq:proj_variational_char}:
\begin{align}
\langle \mathbf{u} - \hat{\mathbf{u}}, \mathbf{x} - \hat{\mathbf{u}} \rangle &\le 0 \quad \forall \mathbf{x} \in C \label{eq:var_u} \\
\langle \mathbf{v} - \hat{\mathbf{v}}, \mathbf{x}' - \hat{\mathbf{v}} \rangle &\le 0 \quad \forall \mathbf{x}' \in C \label{eq:var_v}
\end{align}
Since $\hat{\mathbf{v}} \in C$, substitute $\mathbf{x} = \hat{\mathbf{v}}$ into~\eqref{eq:var_u}. Similarly, substitute $\mathbf{x}' = \hat{\mathbf{u}}$ into~\eqref{eq:var_v}:
\begin{align}
\langle \mathbf{u} - \hat{\mathbf{u}}, \hat{\mathbf{v}} - \hat{\mathbf{u}} \rangle &\le 0 \implies \langle \hat{\mathbf{u}} - \mathbf{u}, \hat{\mathbf{u}} - \hat{\mathbf{v}} \rangle \le 0 \\
\langle \mathbf{v} - \hat{\mathbf{v}}, \hat{\mathbf{u}} - \hat{\mathbf{v}} \rangle &\le 0
\end{align}
Summing these two inequalities:
\begin{equation}
\langle (\hat{\mathbf{u}} - \hat{\mathbf{v}}) - (\mathbf{u} - \mathbf{v}), \hat{\mathbf{u}} - \hat{\mathbf{v}} \rangle \le 0 \implies \|\hat{\mathbf{u}} - \hat{\mathbf{v}}\|_2^2 \le \langle \mathbf{u} - \mathbf{v}, \hat{\mathbf{u}} - \hat{\mathbf{v}} \rangle
\label{eq:firm_step}
\end{equation}
Now consider the vector identity for any vectors $\mathbf{a}, \mathbf{b} \in \mathbb{R}^d$: $\|\mathbf{a} + \mathbf{b}\|_2^2 = \|\mathbf{a}\|_2^2 + \|\mathbf{b}\|_2^2 + 2\langle \mathbf{a}, \mathbf{b} \rangle$. Setting $\mathbf{a} = \hat{\mathbf{u}} - \hat{\mathbf{v}}$ and $\mathbf{b} = (\mathbf{u} - \hat{\mathbf{u}}) - (\mathbf{v} - \hat{\mathbf{v}})$ so that $\mathbf{a} + \mathbf{b} = \mathbf{u} - \mathbf{v}$:
\begin{align}
\|\mathbf{u} - \mathbf{v}\|_2^2 &= \|\hat{\mathbf{u}} - \hat{\mathbf{v}}\|_2^2 + \|(\mathbf{u} - \hat{\mathbf{u}}) - (\mathbf{v} - \hat{\mathbf{v}})\|_2^2 + 2\langle \hat{\mathbf{u}} - \hat{\mathbf{v}}, (\mathbf{u} - \hat{\mathbf{u}}) - (\mathbf{v} - \hat{\mathbf{v}}) \rangle \nonumber \\
&= \|\Pi_C(\mathbf{u}) - \Pi_C(\mathbf{v})\|_2^2 + \|(I - \Pi_C)(\mathbf{u}) - (I - \Pi_C)(\mathbf{v})\|_2^2 + 2\big(\langle \hat{\mathbf{u}} - \hat{\mathbf{v}}, \mathbf{u} - \mathbf{v} \rangle - \|\hat{\mathbf{u}} - \hat{\mathbf{v}}\|_2^2\big)
\end{align}
By~\eqref{eq:firm_step}, the cross-term $\langle \hat{\mathbf{u}} - \hat{\mathbf{v}}, \mathbf{u} - \mathbf{v} \rangle - \|\hat{\mathbf{u}} - \hat{\mathbf{v}}\|_2^2 \ge 0$. Dropping this non-negative term yields~\eqref{eq:firm_non_expansive}. Dropping the non-negative residual norm $\|(I - \Pi_C)(\mathbf{u}) - (I - \Pi_C)(\mathbf{v})\|_2^2 \ge 0$ immediately gives $\|\Pi_C(\mathbf{u}) - \Pi_C(\mathbf{v})\|_2 \le \|\mathbf{u} - \mathbf{v}\|_2$.
\end{proof}

\begin{ponderbox}{Does Projection Accelerate or Retard Convergence?}
Since the projection operator satisfies $\|\Pi_C(\mathbf{u}) - \Pi_C(\mathbf{v})\|_2 \le \|\mathbf{u} - \mathbf{v}\|_2$, does the projection step shrink the distance to the constrained global minimizer $\mathbf{w}^* \in C$?
\tcblower
\textbf{The Socratic Resolution:} Yes! Notice that because the optimal solution $\mathbf{w}^*$ resides in $C$, projecting it leaves it completely unaltered: $\Pi_C(\mathbf{w}^*) = \mathbf{w}^*$. 

Setting $\mathbf{u} = \mathbf{w}_k - \eta \nabla f(\mathbf{w}_k)$ and $\mathbf{v} = \mathbf{w}^*$ in the non-expansiveness inequality:
\begin{align}
\|\mathbf{w}_{k+1} - \mathbf{w}^*\|_2 &= \|\Pi_C(\mathbf{w}_k - \eta \nabla f(\mathbf{w}_k)) - \Pi_C(\mathbf{w}^*)\|_2 \nonumber \\
&\le \|(\mathbf{w}_k - \eta \nabla f(\mathbf{w}_k)) - \mathbf{w}^*\|_2
\end{align}
The projection operator \textbf{never} pushes the parameter iterate farther away from the target $\mathbf{w}^*$ than the intermediate unconstrained step did. If the unconstrained gradient step overshoots the feasible domain, the orthogonal projection strictly contracts the distance to the optimum:
\begin{equation}
\|\mathbf{w}_{k+1} - \mathbf{w}^*\|_2^2 \le \|(\mathbf{w}_k - \eta \nabla f(\mathbf{w}_k)) - \mathbf{w}^*\|_2^2 - \|\mathbf{w}_{k+1} - (\mathbf{w}_k - \eta \nabla f(\mathbf{w}_k))\|_2^2
\end{equation}
This Pythagorean contraction ensures that Projected Gradient Descent achieves the identical $\mathcal{O}(1/k)$ sublinear convergence rate on convex sets as unconstrained Gradient Descent.
\end{ponderbox}

\subsection{Jensen's Inequality: The Engine of Probabilistic Learning}
For any convex functional $g: \mathbb{R} \to \mathbb{R}$ and any random variable $X$ defined on a probability space:
\begin{equation}
g(\E[X]) \le \E[g(X)]
\label{eq:jensen}
\end{equation}
Geometrically, Jensen's inequality is the statistical realization of secant chord convexity. If we take any probability mass distribution $p_1, \dots, p_n$ across values $x_1, \dots, x_n$, their center of mass $\sum p_i x_i = \E[X]$ evaluates to a point $g(\E[X])$ on the curve. In contrast, the expectation $\E[g(X)] = \sum p_i g(x_i)$ is the vertical center of mass of the corresponding points on the curve, which must float on or above the convex bowl.

In Machine Learning Foundations, Jensen's inequality serves as the primary theoretical engine for:
\begin{enumerate}
    \item \textbf{The Evidence Lower Bound (ELBO):} Derived in Chapter~\ref{chap:week12} for the Expectation-Maximization (EM) algorithm. Applying Jensen's inequality to the concave function $\log(\cdot)$:
    \begin{equation}
    \log p(\mathbf{x}) = \log \sum_{\mathbf{z}} q(\mathbf{z}) \frac{p(\mathbf{x}, \mathbf{z})}{q(\mathbf{z})} \ge \sum_{\mathbf{z}} q(\mathbf{z}) \log \frac{p(\mathbf{x}, \mathbf{z})}{q(\mathbf{z})} \triangleq \operatorname{ELBO}(q)
    \end{equation}
    \item \textbf{Non-Negativity of Kullback-Leibler Divergence:}
    \begin{equation}
    D_{\mathrm{KL}}(p \parallel q) = \E_p\left[ -\log \frac{q(X)}{p(X)} \right] \ge -\log \E_p\left[ \frac{q(X)}{p(X)} \right] = -\log(1) = 0
    \end{equation}
    guaranteeing that relative entropy is strictly non-negative, with equality if and only if $p = q$ almost everywhere.
\end{enumerate}

\begin{ponderbox}{The Water-Filling Paradox: Pruning Corrupted Channels}
In parallel Gaussian communication or multi-task resource allocation, common intuition suggests that every available channel should receive at least a small sliver of power ($\epsilon > 0$) to harvest its information throughput. Why does the optimal KKT solution ruthlessly assign $P_k^* = 0$ to noisy channels, and why would providing $\epsilon > 0$ power to those channels strictly degrade total transmission capacity?
\tcblower
\textbf{The Socratic Resolution:} This phenomenon is governed by the \textbf{Equal Marginal Returns Principle} of constrained convex optimization.

The marginal capacity gain obtained from injecting an incremental watt into channel $i$ is:
\begin{equation}
\frac{\partial C}{\partial P_i} = \frac{1}{\ln 2} \cdot \frac{1}{P_i + \sigma_i^2}
\end{equation}
Notice that when power is zero ($P_i = 0$), the initial marginal return is bounded above by the inverse noise floor: $\left.\frac{\partial C}{\partial P_i}\right|_{P_i = 0} = \frac{1}{\sigma_i^2 \ln 2}$. 

Now suppose the system reaches equilibrium across the clean channels at water level $\mu^* = 1/\nu^*$. Every active channel satisfies $P_j^* + \sigma_j^2 = \mu^*$, yielding identical marginal capacity returns: $\frac{\partial C}{\partial P_j} = \frac{1}{\mu^* \ln 2}$.

If a severely corrupted channel has noise variance $\sigma_k^2 > \mu^*$, its maximum possible marginal return (at $P_k = 0$) is:
\begin{equation}
\left.\frac{\partial C}{\partial P_k}\right|_{P_k = 0} = \frac{1}{\sigma_k^2 \ln 2} < \frac{1}{\mu^* \ln 2}
\end{equation}
Because the total power budget $\sum P_i = P_{\text{total}}$ is strictly conserved, granting $\epsilon > 0$ power to the noisy channel requires stealing $\epsilon$ power from clean channels. 
The capacity loss incurred by depriving clean channels is $\approx \frac{\epsilon}{\mu^* \ln 2}$, whereas the capacity gained from the noisy channel is only $\approx \frac{\epsilon}{\sigma_k^2 \ln 2}$. Because $\sigma_k^2 > \mu^*$:
\begin{equation}
\Delta C_{\text{total}} \approx \epsilon \left( \frac{1}{\sigma_k^2 \ln 2} - \frac{1}{\mu^* \ln 2} \right) < 0
\end{equation}
Every watt diverted to the noisy channel produces less capacity than it destroys in the clean channels. In constrained optimization, the boundary condition $P_k^* = 0$ is not a defect; it is nature's mechanism for pruning suboptimal dimensions.
\end{ponderbox}

% ----------------------------------------------------------------------
% PART 5: MULTI-PAGE WORKED EXAMPLES
% ----------------------------------------------------------------------
\section{Worked Numerical Examples}

\begin{examplebox}[breakable]{Step-by-Step Lagrange Multiplier Optimization on an Anisotropic Objective}
\textbf{Problem Statement:}
Consider minimizing the anisotropic quadratic objective functional:
\begin{equation}
f(x, y) = 2x^2 + y^2
\end{equation}
subject to the linear equality constraint:
\begin{equation}
h(x, y) = x + y - 3 = 0
\end{equation}

\textbf{Task Requirements:}
\begin{enumerate}
    \item Formulate the Lagrangian functional $\mathcal{L}(x, y, \nu)$.
    \item Compute the partial derivatives and set up the stationarity linear system.
    \item Solve for the exact optimal coordinates $(x^*, y^*)$ and the optimal multiplier $\nu^*$.
    \item Verify that the objective gradient $\nabla f(x^*, y^*)$ and constraint gradient $\nabla h(x^*, y^*)$ satisfy the collinearity condition.
    \item Compute the optimal objective value $f(x^*, y^*)$ and rigorously contrast it with the value obtained by projecting the unconstrained minimizer $(0, 0)$ onto $x + y = 3$.
\end{enumerate}

\tcblower

\textbf{Step 1: Formulation of the Lagrangian}
The Lagrangian introduces multiplier $\nu \in \mathbb{R}$:
\begin{equation}
\mathcal{L}(x, y, \nu) = f(x, y) + \nu \, h(x, y) = 2x^2 + y^2 + \nu(x + y - 3)
\end{equation}

\textbf{Step 2: Stationarity Conditions}
Setting all partial derivatives to zero:
\begin{align}
\frac{\partial \mathcal{L}}{\partial x} &= 4x + \nu = 0 \label{eq:ex_stat_x} \\
\frac{\partial \mathcal{L}}{\partial y} &= 2y + \nu = 0 \label{eq:ex_stat_y} \\
\frac{\partial \mathcal{L}}{\partial \nu} &= x + y - 3 = 0 \label{eq:ex_stat_nu}
\end{align}

\textbf{Step 3: Analytical Solution}
From~\eqref{eq:ex_stat_x} and~\eqref{eq:ex_stat_y}, express $x$ and $y$ as functions of $\nu$:
\begin{align}
x &= -\frac{\nu}{4} \\
y &= -\frac{\nu}{2}
\end{align}
Substitute these expressions into the primal feasibility equation~\eqref{eq:ex_stat_nu}:
\begin{equation}
-\frac{\nu}{4} - \frac{\nu}{2} - 3 = 0 \implies -\frac{3\nu}{4} = 3 \implies \nu^* = -4
\end{equation}
Now back-substitute $\nu^* = -4$:
\begin{align}
x^* &= -\frac{-4}{4} = 1.0 \\
y^* &= -\frac{-4}{2} = 2.0
\end{align}
The optimal constrained solution is $\mathbf{w}^* = (x^*, y^*) = (1, 2)$.

\textbf{Step 4: Verification of Gradient Collinearity}
Compute the gradient of the objective function at $(1, 2)$:
\begin{equation}
\nabla f(1, 2) = \begin{bmatrix} 4x \\ 2y \end{bmatrix}\Bigg|_{(1, 2)} = \begin{bmatrix} 4(1) \\ 2(2) \end{bmatrix} = \begin{bmatrix} 4 \\ 4 \end{bmatrix}
\end{equation}
Compute the gradient of the constraint function:
\begin{equation}
\nabla h(1, 2) = \begin{bmatrix} 1 \\ 1 \end{bmatrix}
\end{equation}
Check the stationarity condition:
\begin{equation}
\nabla f(1, 2) + \nu^* \nabla h(1, 2) = \begin{bmatrix} 4 \\ 4 \end{bmatrix} + (-4) \begin{bmatrix} 1 \\ 1 \end{bmatrix} = \begin{bmatrix} 0 \\ 0 \end{bmatrix}
\end{equation}
Equivalently:
\begin{equation}
\nabla f(1, 2) = 4 \begin{bmatrix} 1 \\ 1 \end{bmatrix} = -\nu^* \nabla h(1, 2)
\end{equation}
The gradients are strictly collinear. The objective contour ellipse $2x^2 + y^2 = 6$ is exactly tangent to the constraint line $x + y = 3$ at $(1, 2)$.

\textbf{Step 5: Optimal Value vs.\ Unconstrained Projection}
The true constrained optimal value is:
\begin{equation}
f(x^*, y^*) = 2(1)^2 + (2)^2 = 2 + 4 = 6.000
\end{equation}
Now consider the naive projection heuristic:
The unconstrained minimizer is $(x_{\text{unc}}, y_{\text{unc}}) = (0, 0)$.
The Euclidean projection of $(0, 0)$ onto the hyperplane $x + y = 3$ is:
\begin{equation}
\Pi(0, 0) = \mathbf{0} - \frac{\mathbf{a}^T \mathbf{0} - b}{\|\mathbf{a}\|_2^2} \mathbf{a} = \mathbf{0} - \frac{0 + 0 - 3}{1^2 + 1^2} \begin{bmatrix} 1 \\ 1 \end{bmatrix} = \frac{3}{2} \begin{bmatrix} 1 \\ 1 \end{bmatrix} = (1.5, 1.5)
\end{equation}
Evaluating the objective at this projected point:
\begin{equation}
f(1.5, 1.5) = 2(1.5)^2 + (1.5)^2 = 2(2.25) + 2.25 = 4.5 + 2.25 = 6.750
\end{equation}
Suboptimality gap:
\begin{equation}
\Delta f = 6.750 - 6.000 = 0.750 \quad (+12.5\% \text{ excess loss})
\end{equation}
This rigorous calculation confirms that projecting the unconstrained minimum onto the feasible set fails because Euclidean projection treats all coordinate axes symmetrically, whereas the objective function is twice as steep along the $x$-axis as along the $y$-axis.
\end{examplebox}

\begin{examplebox}[breakable]{Convexity Analysis and Hessian Nullspace of the Log-Sum-Exp Function}
\textbf{Problem Statement:}
In multinomial logistic regression and deep neural networks, the cross-entropy loss relies on the log-sum-exp functional:
\begin{equation}
f(\mathbf{x}) = \log\left( \sum_{i=1}^d \exp(x_i) \right) \quad \text{for } \mathbf{x} \in \mathbb{R}^d
\end{equation}

\textbf{Task Requirements:}
\begin{enumerate}
    \item Compute the gradient $\nabla f(\mathbf{x})$ and prove it forms the softmax probability vector.
    \item Derive the $d \times d$ Hessian matrix $\nabla^2 f(\mathbf{x})$.
    \item Prove that $\nabla^2 f(\mathbf{x}) \succeq 0$ for all $\mathbf{x} \in \mathbb{R}^d$ using statistical variance identities.
    \item Characterize the nullspace of $\nabla^2 f(\mathbf{x})$ and explain why $f(\mathbf{x})$ is convex, but \emph{not strictly convex}.
\end{enumerate}

\tcblower

\textbf{Step 1: Gradient Derivation (The Softmax Mapping)}
Let $S(\mathbf{x}) = \sum_{j=1}^d e^{x_j}$. Then $f(\mathbf{x}) = \log S(\mathbf{x})$.
By the multivariable chain rule, for any coordinate $k \in \{1, \dots, d\}$:
\begin{equation}
\frac{\partial f}{\partial x_k} = \frac{1}{S(\mathbf{x})} \frac{\partial S(\mathbf{x})}{\partial x_k} = \frac{e^{x_k}}{\sum_{j=1}^d e^{x_j}} \triangleq \sigma_k(\mathbf{x})
\end{equation}
Observe that:
\begin{equation}
\sigma_k(\mathbf{x}) > 0 \quad \forall k, \quad \text{and} \quad \sum_{k=1}^d \sigma_k(\mathbf{x}) = \frac{\sum_{k=1}^d e^{x_k}}{S(\mathbf{x})} = 1
\end{equation}
Thus, the gradient vector $\nabla f(\mathbf{x}) = \boldsymbol{\sigma}(\mathbf{x}) = [\sigma_1, \dots, \sigma_d]^T$ lies strictly within the interior of the probability simplex $\Delta^d$.

\textbf{Step 2: Hessian Matrix Derivation}
Differentiating $\sigma_k$ with respect to $x_j$:
\begin{itemize}
    \item When $j = k$:
    \begin{equation}
    \frac{\partial \sigma_k}{\partial x_k} = \frac{e^{x_k} S - e^{x_k} e^{x_k}}{S^2} = \frac{e^{x_k}}{S} - \left(\frac{e^{x_k}}{S}\right)^2 = \sigma_k - \sigma_k^2 = \sigma_k(1 - \sigma_k)
    \end{equation}
    \item When $j \ne k$:
    \begin{equation}
    \frac{\partial \sigma_k}{\partial x_j} = \frac{0 \cdot S - e^{x_k} e^{x_j}}{S^2} = -\sigma_k \sigma_j
    \end{equation}
\end{itemize}
Using the Kronecker delta $\delta_{kj}$:
\begin{equation}
[\nabla^2 f(\mathbf{x})]_{kj} = \sigma_k \delta_{kj} - \sigma_k \sigma_j
\end{equation}
In compact matrix form:
\begin{equation}
\nabla^2 f(\mathbf{x}) = \diag(\boldsymbol{\sigma}) - \boldsymbol{\sigma}\boldsymbol{\sigma}^T
\end{equation}

\textbf{Step 3: Proof of Positive Semidefiniteness via Variance}
Let $\mathbf{v} \in \mathbb{R}^d$ be an arbitrary non-zero test vector. Compute the quadratic form:
\begin{align}
\mathbf{v}^T \nabla^2 f(\mathbf{x}) \mathbf{v} &= \mathbf{v}^T \left( \diag(\boldsymbol{\sigma}) - \boldsymbol{\sigma}\boldsymbol{\sigma}^T \right) \mathbf{v} \nonumber \\
&= \mathbf{v}^T \diag(\boldsymbol{\sigma}) \mathbf{v} - (\mathbf{v}^T \boldsymbol{\sigma})^2 \nonumber \\
&= \sum_{i=1}^d \sigma_i v_i^2 - \left( \sum_{i=1}^d \sigma_i v_i \right)^2
\end{align}
Interpret the components of $\mathbf{v}$ as values taken by a discrete random variable $V$, with probability distribution given by $\Pr(V = v_i) = \sigma_i$.
Then:
\begin{align}
\sum_{i=1}^d \sigma_i v_i^2 &= \E[V^2] \\
\sum_{i=1}^d \sigma_i v_i &= \E[V]
\end{align}
Substituting these statistical expectations into the quadratic form:
\begin{equation}
\mathbf{v}^T \nabla^2 f(\mathbf{x}) \mathbf{v} = \E[V^2] - (\E[V])^2 = \Var(V)
\end{equation}
Because the mathematical variance of any real random variable is non-negative ($\Var(V) \ge 0$):
\begin{equation}
\mathbf{v}^T \nabla^2 f(\mathbf{x}) \mathbf{v} \ge 0 \quad \forall \mathbf{v} \in \mathbb{R}^d \implies \nabla^2 f(\mathbf{x}) \succeq 0
\end{equation}
This rigorously establishes that the log-sum-exp functional is globally convex on $\mathbb{R}^d$.

\textbf{Step 4: Characterization of the Nullspace}
When does $\Var(V) = 0$?
The variance of a discrete random variable is zero if and only if $V$ is a constant almost surely:
\begin{equation}
v_1 = v_2 = \dots = v_d = c \in \mathbb{R} \iff \mathbf{v} = c \mathbf{1}
\end{equation}
where $\mathbf{1} = [1, 1, \dots, 1]^T$.
Therefore, the nullspace of the Hessian is one-dimensional:
\begin{equation}
\Null(\nabla^2 f(\mathbf{x})) = \Span(\mathbf{1})
\end{equation}
For any shift direction $\mathbf{v} = c\mathbf{1}$:
\begin{equation}
f(\mathbf{x} + c\mathbf{1}) = \log\left(\sum_{i=1}^d e^{x_i + c}\right) = \log\left(e^c \sum_{i=1}^d e^{x_i}\right) = c + \log\left(\sum_{i=1}^d e^{x_i}\right) = f(\mathbf{x}) + c
\end{equation}
Along the line of ones, $f$ is purely linear (affine), which confirms that $\nabla^2 f$ has an eigenvalue of zero: $\nabla^2 f(\mathbf{x}) \mathbf{1} = \mathbf{0}$.
Consequently, $f$ is \textbf{convex}, but \textbf{not strictly convex}. It is strictly convex only on the orthogonal subspace $\{\mathbf{v} \mid \mathbf{1}^T \mathbf{v} = 0\}$.
\end{examplebox}

\begin{examplebox}[breakable]{Exact Euclidean Projection onto a Polyhedral Halfspace}
\textbf{Problem Statement:}
Consider the closed polyhedral halfspace:
\begin{equation}
C = \left\{ \mathbf{x} \in \mathbb{R}^2 \;\middle|\; 2x_1 + x_2 \le 4 \right\}
\end{equation}
Let $\mathbf{x}_0 = [3.0, 4.0]^T \in \mathbb{R}^2$ be an external test point.

\textbf{Task Requirements:}
\begin{enumerate}
    \item Verify that $\mathbf{x}_0$ is infeasible and violates the halfspace constraint.
    \item Formulate the projection problem as a constrained quadratic minimization task.
    \item Solve for the exact projection coordinates $\mathbf{x}^* = \Pi_C(\mathbf{x}_0)$ using Lagrange multipliers.
    \item Verify the orthogonality condition: show that the displacement vector $\mathbf{x}_0 - \mathbf{x}^*$ is parallel to the constraint normal vector $\mathbf{a} = [2, 1]^T$.
    \item Calculate the exact Euclidean distance from $\mathbf{x}_0$ to the halfspace $C$.
\end{enumerate}

\tcblower

\textbf{Step 1: Feasibility Verification}
Evaluate the constraint functional at $\mathbf{x}_0 = [3.0, 4.0]^T$:
\begin{equation}
g(\mathbf{x}_0) = 2(3.0) + (4.0) = 6.0 + 4.0 = 10.0 > 4.0
\end{equation}
Because $10.0 > 4.0$, $\mathbf{x}_0 \notin C$. Since $C$ is a closed convex set, the projection $\mathbf{x}^*$ must lie on the active boundary hyperplane:
\begin{equation}
\mathcal{H} = \{\mathbf{x} \in \mathbb{R}^2 \mid 2x_1 + x_2 = 4\}
\end{equation}

\textbf{Step 2: Constrained Quadratic Minimization Formulation}
The Euclidean projection solves:
\begin{equation}
\mathbf{x}^* = \argmin_{\mathbf{x} \in \mathbb{R}^2} \frac{1}{2} \|\mathbf{x} - \mathbf{x}_0\|_2^2 \quad \text{subject to } \mathbf{a}^T \mathbf{x} = b
\end{equation}
where $\mathbf{a} = [2, 1]^T$ and $b = 4$.
The Lagrangian functional is:
\begin{equation}
\mathcal{L}(\mathbf{x}, \nu) = \frac{1}{2}\|\mathbf{x} - \mathbf{x}_0\|_2^2 + \nu(\mathbf{a}^T \mathbf{x} - b)
\end{equation}

\textbf{Step 3: Stationarity and Analytical Solution}
Differentiating with respect to $\mathbf{x}$:
\begin{equation}
\nabla_{\mathbf{x}}\mathcal{L} = \mathbf{x} - \mathbf{x}_0 + \nu \mathbf{a} = \mathbf{0} \implies \mathbf{x} = \mathbf{x}_0 - \nu \mathbf{a}
\end{equation}
Multiply by $\mathbf{a}^T$ from the left to enforce the boundary constraint $\mathbf{a}^T \mathbf{x} = b$:
\begin{equation}
\mathbf{a}^T (\mathbf{x}_0 - \nu \mathbf{a}) = b \implies \mathbf{a}^T \mathbf{x}_0 - \nu \|\mathbf{a}\|_2^2 = b
\end{equation}
Solving for the optimal multiplier $\nu^*$:
\begin{equation}
\nu^* = \frac{\mathbf{a}^T \mathbf{x}_0 - b}{\|\mathbf{a}\|_2^2}
\end{equation}
Numerically evaluating the components:
\begin{align}
\mathbf{a}^T \mathbf{x}_0 - b &= 2(3.0) + 1(4.0) - 4 = 10.0 - 4 = 6.0 \\
\|\mathbf{a}\|_2^2 &= 2^2 + 1^2 = 4 + 1 = 5.0
\end{align}
Thus:
\begin{equation}
\nu^* = \frac{6.0}{5.0} = 1.2
\end{equation}
Now compute the projection vector $\mathbf{x}^*$:
\begin{align}
\mathbf{x}^* &= \mathbf{x}_0 - \nu^* \mathbf{a} \nonumber \\
&= \begin{bmatrix} 3.0 \\ 4.0 \end{bmatrix} - 1.2 \begin{bmatrix} 2.0 \\ 1.0 \end{bmatrix}
= \begin{bmatrix} 3.0 - 2.4 \\ 4.0 - 1.2 \end{bmatrix} = \begin{bmatrix} 0.6 \\ 2.8 \end{bmatrix}
\end{align}
The projected coordinates are $\mathbf{x}^* = (0.6, 2.8)$.

\textbf{Step 4: Boundary and Orthogonality Verification}
\begin{itemize}
    \item Boundary check: $2(0.6) + 2.8 = 1.2 + 2.8 = 4.0$. Feasibility holds with exact equality.
    \item Orthogonality check: The displacement vector is:
    \begin{equation}
    \mathbf{d} = \mathbf{x}_0 - \mathbf{x}^* = \begin{bmatrix} 3.0 - 0.6 \\ 4.0 - 2.8 \end{bmatrix} = \begin{bmatrix} 2.4 \\ 1.2 \end{bmatrix} = 1.2 \begin{bmatrix} 2 \\ 1 \end{bmatrix} = 1.2 \, \mathbf{a}
    \end{equation}
    The displacement vector $\mathbf{d}$ is strictly parallel to the outward normal $\mathbf{a}$.
    Any tangent direction along the line $2x_1 + x_2 = 4$ is of the form $\mathbf{t} = [-1, 2]^T$. Testing orthogonality:
    \begin{equation}
    \mathbf{d}^T \mathbf{t} = \begin{bmatrix} 2.4 & 1.2 \end{bmatrix} \begin{bmatrix} -1 \\ 2 \end{bmatrix} = -2.4 + 2.4 = 0
    \end{equation}
\end{itemize}

\textbf{Step 5: Exact Euclidean Distance}
The minimum distance from $\mathbf{x}_0$ to the halfspace $C$ is:
\begin{equation}
\operatorname{dist}(\mathbf{x}_0, C) = \|\mathbf{x}_0 - \mathbf{x}^*\|_2 = \|\mathbf{d}\|_2 = \sqrt{(2.4)^2 + (1.2)^2} = \sqrt{5.76 + 1.44} = \sqrt{7.20} \approx 2.6833
\end{equation}
Notice that this matches the classical point-to-plane distance formula:
\begin{equation}
\operatorname{dist}(\mathbf{x}_0, \mathcal{H}) = \frac{|\mathbf{a}^T \mathbf{x}_0 - b|}{\|\mathbf{a}\|_2} = \frac{|10 - 4|}{\sqrt{5}} = \frac{6}{\sqrt{5}} = \frac{6\sqrt{5}}{5} \approx 2.6833
\end{equation}
\end{examplebox}

\begin{examplebox}[breakable]{The Water-Filling Algorithm for Parallel Gaussian Channels}
\textbf{Problem Statement:}
In wireless communications and optimal resource allocation, a transmitter must distribute a total power budget $P_{\text{total}} = 2.0$ watts across $K = 3$ independent parallel Gaussian channels to maximize the total Shannon information throughput. The channel noise variances are:
\begin{equation}
\sigma_1^2 = 0.20, \quad \sigma_2^2 = 0.60, \quad \sigma_3^2 = 1.50
\end{equation}
The total Shannon capacity (in bits per second per Hertz) is given by:
\begin{equation}
C(P_1, P_2, P_3) = \frac{1}{2} \sum_{i=1}^3 \log_2\left( 1 + \frac{P_i}{\sigma_i^2} \right)
\end{equation}
Subject to the budget and non-negativity constraints:
\begin{equation}
\sum_{i=1}^3 P_i = 2.0, \quad P_i \ge 0 \quad (i = 1, 2, 3)
\end{equation}

\textbf{Task Requirements:}
\begin{enumerate}
    \item Formulate the equivalent convex minimization problem and its Lagrangian functional $\mathcal{L}(\mathbf{P}, \nu, \boldsymbol{\lambda})$.
    \item Derive the Karush-Kuhn-Tucker (KKT) first-order optimality conditions and the explicit water-filling equation $P_i^* = \max(0, \mu^* - \sigma_i^2)$ with water level $\mu^* \triangleq 1/\nu^*$.
    \item Solve step-by-step for the exact equilibrium water level $\mu^*$, proving why Channel~3 must be deactivated ($P_3^* = 0$).
    \item Compute the exact optimal power allocation $\mathbf{P}^* = [P_1^*, P_2^*, P_3^*]^T$ and dual multipliers $\nu^*$ and $\lambda_3^*$.
    \item Calculate the optimal Shannon capacity $C(\mathbf{P}^*)$ and contrast it with a naive uniform allocation $P_1 = P_2 = P_3 = \frac{2}{3}$, quantifying the percentage spectral efficiency gain.
\end{enumerate}

\tcblower

\textbf{Step 1: Convex Minimization Formulation \& Lagrangian}
Maximizing $\sum \frac{1}{2}\log_2(1 + P_i/\sigma_i^2)$ is mathematically equivalent to minimizing the negative natural logarithm:
\begin{equation}
\min_{\mathbf{P} \in \mathbb{R}^3} f(\mathbf{P}) \triangleq -\sum_{i=1}^3 \ln(P_i + \sigma_i^2) \quad \text{subject to } \sum_{i=1}^3 P_i = 2.0, \quad -P_i \le 0 \; (i=1, 2, 3)
\end{equation}
where we dropped the constant terms $\ln(\sigma_i^2)$. The objective is strictly convex because each $-\ln(P_i + \sigma_i^2)$ has positive second derivative $\frac{1}{(P_i + \sigma_i^2)^2} > 0$.
Introducing equality multiplier $\nu \in \mathbb{R}$ and inequality multipliers $\lambda_i \ge 0$:
\begin{equation}
\mathcal{L}(\mathbf{P}, \nu, \boldsymbol{\lambda}) = -\sum_{i=1}^3 \ln(P_i + \sigma_i^2) + \nu \left( \sum_{i=1}^3 P_i - 2.0 \right) - \sum_{i=1}^3 \lambda_i P_i
\end{equation}

\textbf{Step 2: KKT Optimality Conditions}
The KKT conditions require:
\begin{align}
\text{Stationarity: } & \frac{\partial \mathcal{L}}{\partial P_i} = -\frac{1}{P_i + \sigma_i^2} + \nu - \lambda_i = 0 \implies \frac{1}{P_i + \sigma_i^2} = \nu - \lambda_i \label{eq:wf_stat} \\
\text{Primal Feasibility: } & \sum_{i=1}^3 P_i = 2.0, \quad P_i \ge 0 \quad \forall i \in \{1, 2, 3\} \\
\text{Dual Feasibility: } & \lambda_i \ge 0 \quad \forall i \in \{1, 2, 3\} \\
\text{Complementary Slackness: } & \lambda_i P_i = 0 \quad \forall i \in \{1, 2, 3\} \label{eq:wf_slack}
\end{align}
From complementary slackness~\eqref{eq:wf_slack}:
\begin{itemize}
    \item If $P_i > 0$, then $\lambda_i = 0$. From~\eqref{eq:wf_stat}: $\frac{1}{P_i + \sigma_i^2} = \nu \implies P_i + \sigma_i^2 = \frac{1}{\nu} \triangleq \mu^*$. Thus $P_i = \mu^* - \sigma_i^2$.
    \item If $P_i = 0$, then $\lambda_i = \nu - \frac{1}{\sigma_i^2} \ge 0 \implies \nu \ge \frac{1}{\sigma_i^2} \iff \sigma_i^2 \ge \frac{1}{\nu} = \mu^*$.
\end{itemize}
Combining both cases yields the master \textbf{Water-Filling Formula}:
\begin{equation}
P_i^* = \max\left(0, \, \mu^* - \sigma_i^2\right) = (\mu^* - \sigma_i^2)^+
\end{equation}

\textbf{Step 3: Analytical Determination of Water Level $\mu^*$}
We determine $\mu^*$ by sorting noise levels in ascending order: $\sigma_1^2 = 0.20 < \sigma_2^2 = 0.60 < \sigma_3^2 = 1.50$.

\emph{Trial 1: Assume all $K=3$ channels are active ($P_i > 0$ for $i=1, 2, 3$).}
The total power constraint becomes:
\begin{equation}
\sum_{i=1}^3 (\mu - \sigma_i^2) = 3\mu - (0.20 + 0.60 + 1.50) = 3\mu - 2.30 = 2.00
\end{equation}
Solving for $\mu$:
\begin{equation}
3\mu = 4.30 \implies \mu = \frac{4.30}{3} \approx 1.4333
\end{equation}
Now check feasibility for Channel~3:
\begin{equation}
P_3 = \mu - \sigma_3^2 = 1.4333 - 1.50 = -0.0667 < 0
\end{equation}
This violates the primal non-negativity constraint $P_3 \ge 0$. Hence, Channel~3 cannot be active.

\emph{Trial 2: Set $P_3^* = 0$ (Channel~3 deactivated). Assume Channels~1 and~2 are active.}
The power budget must be satisfied exclusively by Channels~1 and~2:
\begin{equation}
(\mu - \sigma_1^2) + (\mu - \sigma_2^2) = (\mu - 0.20) + (\mu - 0.60) = 2\mu - 0.80 = 2.00
\end{equation}
Solving for $\mu^*$:
\begin{equation}
2\mu^* = 2.80 \implies \mu^* = 1.40
\end{equation}
Now verify consistency across all three channels:
\begin{align}
P_1^* &= \mu^* - \sigma_1^2 = 1.40 - 0.20 = 1.20 > 0 \quad \checkmark \\
P_2^* &= \mu^* - \sigma_2^2 = 1.40 - 0.60 = 0.80 > 0 \quad \checkmark \\
\sigma_3^2 &= 1.50 \ge \mu^* = 1.40 \implies P_3^* = 0.00 \quad \checkmark
\end{align}
Total power allocated: $P_1^* + P_2^* + P_3^* = 1.20 + 0.80 + 0.00 = 2.00 = P_{\text{total}}$ ($\checkmark$).

\textbf{Step 4: Verification of Dual Multipliers}
\begin{align}
\text{Equality Multiplier: } & \nu^* = \frac{1}{\mu^*} = \frac{1}{1.40} = \frac{5}{7} \approx 0.7143 \\
\text{Channel 1 Multiplier: } & \lambda_1^* = 0.00 \quad (\text{since } P_1^* > 0) \\
\text{Channel 2 Multiplier: } & \lambda_2^* = 0.00 \quad (\text{since } P_2^* > 0) \\
\text{Channel 3 Multiplier: } & \lambda_3^* = \nu^* - \frac{1}{\sigma_3^2} = \frac{5}{7} - \frac{1}{1.50} = \frac{5}{7} - \frac{2}{3} = \frac{15 - 14}{21} = \frac{1}{21} \approx 0.0476 \ge 0 \quad \checkmark
\end{align}
Because $\lambda_3^* > 0$, dual feasibility holds strictly, proving that Channel~3 is correctly inactive.

\textbf{Step 5: Capacity Evaluation \& Benchmark Comparison}
Compute the optimal Shannon capacity under water-filling allocation $\mathbf{P}^* = [1.20, 0.80, 0.00]^T$:
\begin{align}
C(\mathbf{P}^*) &= \frac{1}{2}\left[ \log_2\left(1 + \frac{1.20}{0.20}\right) + \log_2\left(1 + \frac{0.80}{0.60}\right) + \log_2\left(1 + \frac{0.00}{1.50}\right) \right] \nonumber \\
&= \frac{1}{2}\left[ \log_2(7.0) + \log_2(2.3333) + 0 \right] \nonumber \\
&= \frac{1}{2}\left[ 2.8074 + 1.2224 \right] = \frac{4.0298}{2} = 2.0149 \text{ bits/s/Hz}
\end{align}

Now contrast with a naive uniform allocation dividing $P_{\text{total}} = 2.0$ equally: $P_1 = P_2 = P_3 = \frac{2}{3} \approx 0.6667$:
\begin{align}
C(\mathbf{P}_{\text{uniform}}) &= \frac{1}{2}\left[ \log_2\left(1 + \frac{2/3}{0.20}\right) + \log_2\left(1 + \frac{2/3}{0.60}\right) + \log_2\left(1 + \frac{2/3}{1.50}\right) \right] \nonumber \\
&= \frac{1}{2}\left[ \log_2(1 + 3.3333) + \log_2(1 + 1.1111) + \log_2(1 + 0.4444) \right] \nonumber \\
&= \frac{1}{2}\left[ \log_2(4.3333) + \log_2(2.1111) + \log_2(1.4444) \right] \nonumber \\
&= \frac{1}{2}\left[ 2.1155 + 1.0780 + 0.5305 \right] = \frac{3.7240}{2} = 1.8620 \text{ bits/s/Hz}
\end{align}
The relative capacity gain delivered by optimal water-filling is:
\begin{equation}
\Delta C = \frac{2.0149 - 1.8620}{1.8620} \times 100\% = +8.21\%
\end{equation}
By recognizing that Channel~3's noise floor ($\sigma_3^2 = 1.50$) exceeds the system water level ($\mu^* = 1.40$) and shutting it down entirely, the transmitter diverts critical energy to Channels~1 and~2, achieving a substantial $8.21\%$ increase in spectral efficiency.
\end{examplebox}

\begin{examplebox}[breakable]{Step-by-Step Numerical Trace of Projected Gradient Descent}
\textbf{Problem Statement:}
Consider minimizing the strongly convex, anisotropic quadratic loss functional:
\begin{equation}
f(w_1, w_2) = (w_1 - 3)^2 + 4(w_2 - 2)^2
\end{equation}
subject to the hypercube box constraint set:
\begin{equation}
\mathcal{C} = \left\{ \mathbf{w} \in \mathbb{R}^2 \;\middle|\; 0 \le w_1 \le 1, \; 0 \le w_2 \le 1 \right\} = [0, 1] \times [0, 1]
\end{equation}
Notice that the unconstrained global minimizer is at $\mathbf{w}_{\text{unc}}^* = (3, 2)$ where $f(3, 2) = 0$, which lies strictly outside the feasible box $\mathcal{C}$.

\textbf{Task Requirements:}
\begin{enumerate}
    \item Compute the objective gradient $\nabla f(\mathbf{w})$ and state the coordinate-wise Euclidean projection operator $\Pi_{\mathcal{C}}(\mathbf{z})$.
    \item Perform the complete numerical calculation for **Iteration 1** starting from $\mathbf{w}_0 = (0.0, 0.0)^T$ with learning rate $\eta = 0.20$: compute the gradient $\nabla f(\mathbf{w}_0)$, unconstrained update $\tilde{\mathbf{w}}_1$, projected point $\mathbf{w}_1$, and loss $f(\mathbf{w}_1)$.
    \item Perform the complete numerical calculation for **Iteration 2**: compute $\nabla f(\mathbf{w}_1)$, unconstrained candidate $\tilde{\mathbf{w}}_2$, projected point $\mathbf{w}_2$, and test for numerical convergence.
    \item Rigorously prove that the terminal iterate satisfies the Variational Characterization of Optimality for constrained convex minimization:
    \begin{equation}
    \nabla f(\mathbf{w}^*)^T (\mathbf{z} - \mathbf{w}^*) \ge 0 \quad \forall \mathbf{z} \in \mathcal{C}
    \end{equation}
    confirming that the iterate is the exact, global constrained minimizer.
\end{enumerate}

\tcblower

\textbf{Step 1: Gradient Vector and Projection Operator}
The gradient of $f$ is:
\begin{equation}
\nabla f(\mathbf{w}) = \begin{bmatrix} \frac{\partial f}{\partial w_1} \\ \frac{\partial f}{\partial w_2} \end{bmatrix} = \begin{bmatrix} 2(w_1 - 3) \\ 8(w_2 - 2) \end{bmatrix}
\end{equation}
The Euclidean projection onto the box $\mathcal{C} = [0, 1] \times [0, 1]$ decouples into independent 1D interval projections (clipping):
\begin{equation}
\Pi_{\mathcal{C}}(\mathbf{z}) = \begin{bmatrix} \operatorname{clip}(z_1, 0, 1) \\ \operatorname{clip}(z_2, 0, 1) \end{bmatrix} = \begin{bmatrix} \min(1, \max(0, z_1)) \\ \min(1, \max(0, z_2)) \end{bmatrix}
\end{equation}

\textbf{Step 2: Iteration 1 Walkthrough}
\begin{itemize}
    \item Initial state: $\mathbf{w}_0 = [0.0, 0.0]^T$.
    \item Initial objective value:
    \begin{equation}
    f(\mathbf{w}_0) = (0 - 3)^2 + 4(0 - 2)^2 = (-3)^2 + 4(-2)^2 = 9 + 4(4) = 9 + 16 = 25.0
    \end{equation}
    \item Gradient computation at $\mathbf{w}_0$:
    \begin{equation}
    \nabla f(\mathbf{w}_0) = \begin{bmatrix} 2(0 - 3) \\ 8(0 - 2) \end{bmatrix} = \begin{bmatrix} -6.0 \\ -16.0 \end{bmatrix}
    \end{equation}
    \item Intermediate unconstrained gradient step ($\eta = 0.20$):
    \begin{equation}
    \tilde{\mathbf{w}}_1 = \mathbf{w}_0 - \eta \nabla f(\mathbf{w}_0) = \begin{bmatrix} 0.0 \\ 0.0 \end{bmatrix} - 0.20 \begin{bmatrix} -6.0 \\ -16.0 \end{bmatrix} = \begin{bmatrix} 0.0 + 1.20 \\ 0.0 + 3.20 \end{bmatrix} = \begin{bmatrix} 1.20 \\ 3.20 \end{bmatrix}
    \end{equation}
    Notice that $\tilde{\mathbf{w}}_1 = (1.20, 3.20)$ violates both upper bounds ($w_1 \le 1$ and $w_2 \le 1$).
    \item Orthogonal projection onto $\mathcal{C}$:
    \begin{equation}
    \mathbf{w}_1 = \Pi_{\mathcal{C}}(\tilde{\mathbf{w}}_1) = \begin{bmatrix} \min(1, \max(0, 1.20)) \\ \min(1, \max(0, 3.20)) \end{bmatrix} = \begin{bmatrix} 1.0 \\ 1.0 \end{bmatrix}
    \end{equation}
    \item Objective value after Iteration 1:
    \begin{equation}
    f(\mathbf{w}_1) = f(1.0, 1.0) = (1 - 3)^2 + 4(1 - 2)^2 = (-2)^2 + 4(-1)^2 = 4 + 4(1) = 8.0
    \end{equation}
    The loss plunges dramatically from $25.0 \to 8.0$ (a $68.0\%$ error reduction in a single iteration).
\end{itemize}

\textbf{Step 3: Iteration 2 Walkthrough}
\begin{itemize}
    \item Gradient computation at $\mathbf{w}_1 = [1.0, 1.0]^T$:
    \begin{equation}
    \nabla f(\mathbf{w}_1) = \begin{bmatrix} 2(1 - 3) \\ 8(1 - 2) \end{bmatrix} = \begin{bmatrix} 2(-2) \\ 8(-1) \end{bmatrix} = \begin{bmatrix} -4.0 \\ -8.0 \end{bmatrix}
    \end{equation}
    \item Intermediate unconstrained gradient step:
    \begin{equation}
    \tilde{\mathbf{w}}_2 = \mathbf{w}_1 - \eta \nabla f(\mathbf{w}_1) = \begin{bmatrix} 1.0 \\ 1.0 \end{bmatrix} - 0.20 \begin{bmatrix} -4.0 \\ -8.0 \end{bmatrix} = \begin{bmatrix} 1.0 + 0.80 \\ 1.0 + 1.60 \end{bmatrix} = \begin{bmatrix} 1.80 \\ 2.60 \end{bmatrix}
    \end{equation}
    \item Orthogonal projection onto $\mathcal{C}$:
    \begin{equation}
    \mathbf{w}_2 = \Pi_{\mathcal{C}}(\tilde{\mathbf{w}}_2) = \begin{bmatrix} \min(1, \max(0, 1.80)) \\ \min(1, \max(0, 2.60)) \end{bmatrix} = \begin{bmatrix} 1.0 \\ 1.0 \end{bmatrix}
    \end{equation}
    \item Convergence check:
    \begin{equation}
    \|\mathbf{w}_2 - \mathbf{w}_1\|_2 = \|[1.0, 1.0]^T - [1.0, 1.0]^T\|_2 = 0.0
    \end{equation}
    The iterate has locked onto the stationary fixed point $\mathbf{w}^* = (1.0, 1.0)$ in exactly one step!
\end{itemize}

\textbf{Step 4: Verification of First-Order Variational Optimality}
At the terminal iterate $\mathbf{w}^* = [1.0, 1.0]^T$, the objective gradient is $\nabla f(\mathbf{w}^*) = [-4.0, -8.0]^T$.
Testing the first-order variational inequality for \emph{any} arbitrary feasible point $\mathbf{z} \in [0, 1] \times [0, 1]$:
\begin{align}
\nabla f(\mathbf{w}^*)^T (\mathbf{z} - \mathbf{w}^*) &= \begin{bmatrix} -4.0 & -8.0 \end{bmatrix} \begin{bmatrix} z_1 - 1.0 \\ z_2 - 1.0 \end{bmatrix} \nonumber \\
&= -4.0(z_1 - 1.0) - 8.0(z_2 - 1.0) \nonumber \\
&= 4.0(1.0 - z_1) + 8.0(1.0 - z_2)
\end{align}
Because $\mathbf{z} \in \mathcal{C}$, we have $z_1 \le 1.0 \implies 1.0 - z_1 \ge 0$, and $z_2 \le 1.0 \implies 1.0 - z_2 \ge 0$.
Since both terms are strictly non-negative, their weighted sum is non-negative:
\begin{equation}
\nabla f(\mathbf{w}^*)^T (\mathbf{z} - \mathbf{w}^*) \ge 0 \quad \forall \mathbf{z} \in \mathcal{C}
\end{equation}
with equality if and only if $\mathbf{z} = \mathbf{w}^*$. This rigorously proves that $\mathbf{w}^* = (1.0, 1.0)$ is the unique global constrained minimizer of $f$ over the box $\mathcal{C}$.
\end{examplebox}

% ----------------------------------------------------------------------
% PART 6: EXPLICIT GOTCHAS & DIAGNOSTIC MISCONCEPTIONS
% ----------------------------------------------------------------------
\section{Explicit Gotchas \& Diagnostic Misconceptions}

\begin{trapbox}{Gotcha 1: Assuming the Union of Two Convex Sets is Convex}
A practitioner considers the union $S = B_1 \cup B_2$ of two disjoint convex norm balls $B_1 = \{\mathbf{x} \mid \|\mathbf{x} - \mathbf{c}_1\| \le 1\}$ and $B_2 = \{\mathbf{x} \mid \|\mathbf{x} - \mathbf{c}_2\| \le 1\}$ and assumes $S$ is convex.
\tcblower
\textbf{The Trap:} While the \textbf{intersection} of arbitrary convex sets is always convex, the \textbf{union} of convex sets is almost never convex. A straight line segment connecting a point in $B_1$ to a point in $B_2$ passes through the empty space between them, violating the convexity axiom. In optimization, unions of convex domains form non-convex feasible regions that create disconnected local minima.
\end{trapbox}

\begin{trapbox}{Gotcha 2: Believing Local Minima Can Differ from Global Minima in Convex Problems}
A student worries that gradient descent might get trapped in a suboptimal local minimum when minimizing a convex loss function over a convex set.
\tcblower
\textbf{The Trap:} For a convex functional defined on a convex set, \textbf{every local minimum is automatically a global minimum}. Furthermore, if the functional is strictly convex, the global minimizer is uniquely defined. The challenge of non-global local minima and spurious attractors exists exclusively on non-convex landscapes.
\end{trapbox}

\begin{trapbox}{Gotcha 3: Confusing Strict Convexity with Strong Convexity}
An analyst asserts that because $f(x) = x^4$ is strictly convex on $\mathbb{R}$, gradient descent will achieve linear convergence.
\tcblower
\textbf{The Trap:} $f(x) = x^4$ is strictly convex, but its second derivative $f''(x) = 12x^2$ vanishes at the minimizer $x = 0$. It is \textbf{not} $\mu$-strongly convex ($\mu = 0$). Strong convexity requires a uniform positive lower bound on curvature: $\nabla^2 f(\mathbf{x}) \succeq \mu I$ with $\mu > 0$. Without strong convexity, convergence near flat regions slows down to sublinear $\mathcal{O}(1/k)$.
\end{trapbox}

\begin{trapbox}{Gotcha 4: Forgetting Constraint Qualifications in Lagrange Systems}
A practitioner sets up a Lagrangian for the constraint $h(x, y) = x^2 + y^2 = 0$, but finds no Lagrange multipliers satisfying the stationarity equations.
\tcblower
\textbf{The Trap:} The Lagrange multiplier theorem requires a \textbf{constraint qualification}: the gradients of the active constraints must be linearly independent (LICQ). At the unique feasible point $(0, 0)$, $\nabla h(0, 0) = [2x, 2y]^T = [0, 0]^T$. The gradient vanishes, violating regularity. Always verify regularity ($\nabla h(\mathbf{w}^*) \ne \mathbf{0}$) before setting up the Lagrangian.
\end{trapbox}

\begin{trapbox}{Gotcha 5: The Boundary Curse: Projecting the Unconstrained Minimum}
When asked to minimize an objective $f(\mathbf{w})$ over a convex set $C$, an engineer computes the unconstrained minimum $\mathbf{w}_{\text{unc}}^* = \argmin_{\mathbf{w}} f(\mathbf{w})$ and projects it onto the boundary: $\hat{\mathbf{w}} = \Pi_C(\mathbf{w}_{\text{unc}}^*)$.
\tcblower
\textbf{The Trap:} The projection $\Pi_C(\mathbf{w}_{\text{unc}}^*)$ is \textbf{almost never the true constrained minimum}! Orthogonal projection minimizes Euclidean distance $\|\mathbf{w} - \mathbf{w}_{\text{unc}}^*\|_2^2$, treating all directions with equal isotropic weighting. If the objective function has anisotropic curvature (different eigenvalues in its Hessian), the true constrained optimum occurs at a point of \textbf{tangency} where $\nabla f(\mathbf{w}^*)$ is normal to the constraint boundary, not at the closest Euclidean point.
\end{trapbox}

\begin{trapbox}{Gotcha 6: The "Convex Loss on Non-Convex Parameterization" Illusion}
An ML engineer reasons: ``Mean Squared Error is quadratic and therefore strictly convex. Therefore, fitting a matrix factorization model $W \approx U V^T$ or a two-layer linear network $\hat{\mathbf{y}} = W_2 W_1 \mathbf{x}$ by minimizing squared reconstruction error has no spurious local minima or saddle points.''
\tcblower
\textbf{The Trap:} The composition of a convex function with a non-linear parameterization is \textbf{almost never convex}. 

While the loss functional $\phi(W) = \frac{1}{2}\|Y - W\|_F^2$ is strictly convex with respect to the matrix entries $W$, re-parameterizing $W = U V^T$ transforms the objective into:
\begin{equation}
f(U, V) = \frac{1}{2}\|Y - U V^T\|_F^2
\end{equation}
which is a degree-4 polynomial in the parameters $(U, V)$. Although $f(U, V)$ is \textbf{biconvex} (convex in $U$ for fixed $V$, and convex in $V$ for fixed $U$), it is jointly non-convex. It introduces saddle points at the origin $(U=O, V=O)$ where $\nabla f = \mathbf{0}$ but the Hessian is indefinite, as well as permutation symmetries and non-isolated manifolds of critical points. Never confuse the convexity of an outer loss with the convexity of the parameterized optimization problem.
\end{trapbox}

% ----------------------------------------------------------------------
% PART 7: COMPARATIVE SYNTHESIS & SUMMARY TABLE
% ----------------------------------------------------------------------
\section{Comparative Synthesis \& Summary Table}

Table~\ref{tab:convex_synthesis} summarizes the hierarchy of convexity classes for functionals.

\begin{table}[htbp]
\centering
\small
\begin{tabularx}{\textwidth}{lXXXX}
\toprule
\textbf{Convexity Class} & \textbf{Hessian Condition} & \textbf{First-Order Lower Bound} & \textbf{Minimizer Properties} & \textbf{Canonical Example} \\
\midrule
Convex & $\nabla^2 f(\mathbf{x}) \succeq 0$ & $f(\mathbf{y}) \ge f(\mathbf{x}) + \nabla f(\mathbf{x})^T(\mathbf{y} - \mathbf{x})$ & Any local minimum is a global minimum & $f(\mathbf{w}) = \|\mathbf{w}\|_1$ (Lasso loss) \\
\addlinespace
Strictly Convex & $\nabla^2 f(\mathbf{x}) \succ 0$ a.e. & Strict inequality for $\mathbf{y} \ne \mathbf{x}$ & Unique global minimizer (if it exists) & $f(x) = x^4$, $f(x) = e^x$ \\
\addlinespace
$\mu$-Strongly Convex & $\nabla^2 f(\mathbf{x}) \succeq \mu I$ ($\mu > 0$) & $f(\mathbf{y}) \ge f(\mathbf{x}) + \nabla f^T(\mathbf{y}-\mathbf{x}) + \frac{\mu}{2}\|\mathbf{y}-\mathbf{x}\|^2$ & Unique global minimizer; linear convergence & Ridge Regression: $\frac{1}{2}\|X\mathbf{w}-\mathbf{y}\|^2 + \frac{\lambda}{2}\|\mathbf{w}\|^2$ \\
\addlinespace
Non-Convex & Indefinite / Negative $\nabla^2 f$ & No global affine lower bound & Multiple local minima, saddle points & Neural networks, GMM negative log-likelihood \\
\bottomrule
\end{tabularx}
\caption{Comparative Synthesis: Hierarchy of Function Convexity Classes.}
\label{tab:convex_synthesis}
\end{table}

% ----------------------------------------------------------------------
% PART 8: PYTHON / SCIENTIFIC COMPUTING LAB
% ----------------------------------------------------------------------
\section{Python / Scientific Computing Lab}

The production script below implements Projected Gradient Descent (PGD) onto an $L_1$ ball constraint, demonstrating how projection enforces exact parameter sparsity and contrasting the result with naive unconstrained descent.

\begin{lstlisting}[language=Python, caption={Projected Gradient Descent onto an $L_1$ Ball in Python.}]
import numpy as np

def project_onto_l1_ball(v: np.ndarray, s: float) -> np.ndarray:
    """Computes Euclidean projection of vector v onto L1 ball of radius s: ||w||_1 <= s."""
    if np.sum(np.abs(v)) <= s:
        return v
    # Sort absolute values descending
    u = np.sort(np.abs(v))[::-1]
    cssv = np.cumsum(u)
    # Find active threshold coordinate
    rho = np.nonzero(u * np.arange(1, len(v) + 1) > (cssv - s))[0][-1]
    theta = (cssv[rho] - s) / (rho + 1.0)
    w = np.sign(v) * np.maximum(np.abs(v) - theta, 0.0)
    return w

def run_projected_gradient_descent(X: np.ndarray, y: np.ndarray, s_radius: float, lr: float, n_iters: int):
    """Minimizes 0.5 * ||Xw - y||^2 subject to ||w||_1 <= s_radius via PGD."""
    N, d = X.shape
    w = np.zeros(d)
    losses = []

    for k in range(n_iters):
        grad = X.T @ (X @ w - y)
        w_unconstrained = w - lr * grad
        w = project_onto_l1_ball(w_unconstrained, s_radius)
        loss = 0.5 * float(np.sum((X @ w - y) ** 2))
        losses.append(loss)

    return w, losses

def main():
    rng = np.random.RandomState(42)
    N, d = 100, 10
    X = rng.randn(N, d)
    
    # Ground truth sparse weights (only 3 active components)
    w_true = np.zeros(d)
    w_true[1] = 2.5
    w_true[4] = -1.8
    w_true[7] = 3.0
    y = X @ w_true + 0.1 * rng.randn(N)

    # Maximum Lipschitz smoothness constant L = lambda_max(X^T X)
    L = float(np.linalg.eigvalsh(X.T @ X).max())
    lr = 1.0 / L

    s_radius = 5.0
    w_pgd, losses = run_projected_gradient_descent(X, y, s_radius, lr, n_iters=200)
    w_ols = np.linalg.lstsq(X, y, rcond=None)[0]

    print("=================================================================")
    print(" PROJECTED GRADIENT DESCENT: L1 BALL SPARSITY LAB")
    print("=================================================================")
    print(f"L1 Ball Radius Constraint:      ||w||_1 <= {s_radius:.2f}")
    print(f"Unconstrained OLS Weights Norm:  ||w_ols||_1 = {np.sum(np.abs(w_ols)):.4f}")
    print(f"Projected PGD Weights Norm:      ||w_pgd||_1 = {np.sum(np.abs(w_pgd)):.4f}")
    print(f"Exact Zeros in OLS Solution:     {np.sum(np.isclose(w_ols, 0))} / {d}")
    print(f"Exact Zeros in PGD Solution:     {np.sum(np.isclose(w_pgd, 0))} / {d}")
    print(f"Final Constrained Loss:          {losses[-1]:.4f}")
    print("-" * 65)

    print(f"{'Index':<8}{'True Weight':<16}{'OLS (Unconstrained)':<22}{'PGD (L1 Constrained)':<20}")
    print("-" * 65)
    for i in range(d):
        print(f"{i:<8}{w_true[i]:<16.2f}{w_ols[i]:<22.4f}{w_pgd[i]:<20.4f}")

if __name__ == "__main__":
    main()
\end{lstlisting}

% ----------------------------------------------------------------------
% PART 9: COMPREHENSIVE PRACTICE PROBLEMS
% ----------------------------------------------------------------------
\section{Comprehensive Practice Problems}

\subsection*{Part A: Analytical \& Mathematical Problems}

\begin{enumerate}[label=\textbf{P\arabic*.}]
    \item \textbf{Analytical Lagrange Multipliers with Linear Constraint:}
    Minimize the quadratic objective $f(x, y) = 2x^2 + y^2$ subject to the linear equality constraint $x + y = 3$.
    \begin{enumerate}[label=(\alph*)]
        \item Formulate the Lagrangian $\mathcal{L}(x, y, \nu)$.
        \item Solve the stationarity system $\nabla_{x, y, \nu} \mathcal{L} = \mathbf{0}$.
        \item Compute the optimal coordinates $(x^*, y^*)$ and optimal multiplier $\nu^*$.
        \item Verify that $\nabla f(x^*, y^*)$ is collinear with $\nabla h(x^*, y^*)$.
    \end{enumerate}

    \item \textbf{Maximum Entropy Probability Distribution:}
    Let $\mathbf{p} = [p_1, \dots, p_K]^T$ be a discrete probability vector. Maximize the Shannon entropy $H(\mathbf{p}) = -\sum_{i=1}^K p_i \log p_i$ subject to the simplex constraint $\sum_{i=1}^K p_i = 1$.
    Using Lagrange multipliers, prove that the uniform distribution $p_i^* = \frac{1}{K}$ for all $i$ is the unique global maximizer.

    \item \textbf{Convexity Verification (Second-Order Conditions):}
    Consider the log-sum-exp function $f(\mathbf{x}) = \log\big( \sum_{i=1}^d e^{x_i} \big)$.
    \begin{enumerate}[label=(\alph*)]
        \item Compute the gradient $\nabla f(\mathbf{x})$ and show it is a probability vector (softmax).
        \item Compute the Hessian matrix $\nabla^2 f(\mathbf{x})$.
        \item Prove that $\nabla^2 f(\mathbf{x}) \succeq 0$ for all $\mathbf{x} \in \mathbb{R}^d$, establishing that log-sum-exp is convex.
    \end{enumerate}

    \item \textbf{Projection onto a Halfspace:}
    Let $\mathcal{H}_- = \{\mathbf{x} \in \mathbb{R}^d \mid \mathbf{a}^T \mathbf{x} \le b\}$ with $\|\mathbf{a}\|_2 = 1$.
    Prove that the Euclidean projection of a point $\mathbf{x}_0 \in \mathbb{R}^d$ onto $\mathcal{H}_-$ is given by:
    \begin{equation}
    \Pi_{\mathcal{H}_-}(\mathbf{x}_0) = \mathbf{x}_0 - \max(0, \mathbf{a}^T \mathbf{x}_0 - b) \mathbf{a}
    \end{equation}

    \item \textbf{Jensen's Inequality Application:}
    Using the convexity of $f(t) = -\log t$ on $(0, \infty)$, prove the Arithmetic Mean - Geometric Mean (AM-GM) inequality: for any positive real numbers $a_1, a_2, \dots, a_n > 0$,
    \begin{equation}
    \left( \prod_{i=1}^n a_i \right)^{1/n} \le \frac{1}{n} \sum_{i=1}^n a_i
    \end{equation}

    \item \textbf{Variational Characterization of Projection \& Firm Non-Expansiveness:}
    Let $C \subset \mathbb{R}^d$ be a non-empty closed convex set.
    \begin{enumerate}[label=(\alph*)]
        \item Let $\mathbf{w} \in \mathbb{R}^d$. Prove that $\mathbf{z}^* = \Pi_C(\mathbf{w})$ if and only if:
        \begin{equation*}
        \langle \mathbf{w} - \mathbf{z}^*, \mathbf{x} - \mathbf{z}^* \rangle \le 0 \quad \forall \mathbf{x} \in C
        \end{equation*}
        \item Using part (a), prove that the projection operator is \emph{firmly non-expansive}:
        \begin{equation*}
        \|\Pi_C(\mathbf{u}) - \Pi_C(\mathbf{v})\|_2^2 + \|(I - \Pi_C)(\mathbf{u}) - (I - \Pi_C)(\mathbf{v})\|_2^2 \le \|\mathbf{u} - \mathbf{v}\|_2^2 \quad \forall \mathbf{u}, \mathbf{v} \in \mathbb{R}^d
        \end{equation*}
        \item Deduce that $\|\Pi_C(\mathbf{u}) - \Pi_C(\mathbf{v})\|_2 \le \|\mathbf{u} - \mathbf{v}\|_2$, establishing that $\Pi_C$ is a 1-Lipschitz contraction mapping.
    \end{enumerate}

    \item \textbf{Lagrange Multiplier Sensitivity \& Shadow Price Derivation:}
    Consider the perturbed constrained optimization problem $f^*(\mathbf{c}) = \min_{\mathbf{w}} f(\mathbf{w})$ subject to $h_j(\mathbf{w}) = c_j$ for $j = 1, \dots, p$, where $f$ and all $h_j$ are smooth.
    \begin{enumerate}[label=(\alph*)]
        \item Set up the Lagrangian $\mathcal{L}(\mathbf{w}, \boldsymbol{\nu}; \mathbf{c}) = f(\mathbf{w}) + \sum_{j=1}^p \nu_j (h_j(\mathbf{w}) - c_j)$.
        \item Assuming the optimal trajectory $\mathbf{w}^*(\mathbf{c})$ is differentiable and satisfies the LICQ condition at $\mathbf{c} = \mathbf{0}$, prove using the multivariable chain rule that:
        \begin{equation*}
        \frac{\partial f^*(\mathbf{c})}{\partial c_k} = -\nu_k^*(\mathbf{c})
        \end{equation*}
        \item An automated trading algorithm allocates a total capital budget $B = \$100{,}000$ across assets by minimizing portfolio variance subject to $\sum w_i = B$. If the optimal Lagrange multiplier is $\nu^* = -0.0042$, calculate the exact marginal change in minimal portfolio variance if the capital budget is increased by $\$1{,}000$.
    \end{enumerate}
\end{enumerate}

\subsection*{Part B: Diagnostic Multiple-Choice Questions}

\begin{enumerate}[label=\textbf{MCQ\arabic*.}]
    \item Which of the following sets is NOT convex?
    \begin{enumerate}[label=(\alph*)]
        \item The Euclidean unit ball $\{\mathbf{x} \in \mathbb{R}^d \mid \|\mathbf{x}\|_2 \le 1\}$.
        \item The positive semidefinite cone $S^d_+ = \{A \in \mathbb{R}^{d \times d} \mid A = A^T, A \succeq 0\}$.
        \item The hollow spherical shell $\{\mathbf{x} \in \mathbb{R}^d \mid \|\mathbf{x}\|_2 = 1\}$.
        \item The halfspace $\{\mathbf{x} \in \mathbb{R}^d \mid \mathbf{a}^T \mathbf{x} \le b\}$.
    \end{enumerate}

    \item If $f: \mathbb{R}^d \to \mathbb{R}$ is convex and continuously differentiable, what is the geometric relationship between the graph of $f$ and the tangent hyperplane at any point $\mathbf{x}_0$?
    \begin{enumerate}[label=(\alph*)]
        \item The tangent hyperplane lies strictly above the graph of $f$.
        \item The tangent hyperplane is a global under-estimator that lies everywhere below or touches the graph of $f$.
        \item The tangent hyperplane intersects the graph at infinitely many isolated points.
        \item The tangent hyperplane has a slope equal to the Hessian determinant.
    \end{enumerate}
\end{enumerate}

\chapter[Convex Duality, KKT Conditions \& SVMs]{Convex Duality, KKT Conditions \& Support Vector Machines}
\label{chap:week10}

% ----------------------------------------------------------------------
% PART 1: EXECUTIVE ROADMAP & ARCHITECTURAL FOUNDATIONS
% ----------------------------------------------------------------------
\section{Executive Roadmap \& Learning Objectives}

\subsection{Opening Hook: The Infinite Void and the Mountain Pillars}
Imagine an artist tasked with sculpting an infinite-dimensional mountain range ($d = \infty$). Every microscopic fold, crevice, and ridge must be positioned with mathematical perfection. If the artist is forced to define the mountain by specifying every infinitesimal coordinate, the task is utterly impossible: no human lifetime or computer memory could ever store an infinite sequence of weights.

Now imagine that a geometer enters the studio and reveals a profound structural secret: the entire soaring mountain range is supported from beneath by just three immovable pillars of bedrock. If you know the exact location and upward thrust of those three pillars, the entire mountain's shape is mathematically locked into place. You do not need to sculpt the infinite void of space; you only need to determine the mechanical forces acting at the three contact points.

This is the miracle of \textbf{Lagrange Duality} and \textbf{Support Vector Machines (SVMs)}. 

When machine learning algorithms project data into infinite-dimensional Reproducing Kernel Hilbert Spaces (RKHS) via nonlinear kernels (such as the Gaussian RBF kernel), the primal weight vector $\mathbf{w}$ contains infinitely many parameters. Attempting to optimize $\mathbf{w}$ directly in primal space is computationally intractable. 

Lagrange duality inverts the entire dimensionality contract: instead of optimizing over infinitely many feature dimensions $d = \infty$, the dual problem optimizes over exactly $N$ scalar multipliers $\boldsymbol{\lambda} \in \mathbb{R}^N$---one scalar per training sample. Even more remarkably, the Karush-Kuhn-Tucker (KKT) complementary slackness conditions guarantee that for every data point residing comfortably inside the class territory, its dual multiplier is identically zero ($\lambda_i^* = 0$). 

The infinite-dimensional decision boundary collapses into a sparse linear combination of just a tiny handful of critical boundary samples: the \textbf{support vectors}. The boundary does not care about the millions of points in the interior; it is mechanically pinned by the bedrock samples on the margin.

\subsection{Architectural Bridge from Previous Chapters}
In Chapter~\ref{chap:week09}, we established the geometric foundation of constrained continuous optimization, analyzing Lagrange multipliers for equality-constrained manifolds and proving the non-expansiveness of Euclidean projection onto closed convex sets. 

However, practical machine learning classification rarely involves simple equality constraints. Class boundaries are defined by \textbf{inequalities}: we require samples to lie on or beyond a decision margin:
\begin{equation}
y_i(\mathbf{w}^T \mathbf{x}_i + b) \ge 1 \iff 1 - y_i(\mathbf{w}^T \mathbf{x}_i + b) \le 0
\end{equation}
Chapter 10 generalizes the Lagrangian apparatus to arbitrary systems of convex inequality and affine equality constraints. We formulate the \textbf{Lagrange Dual Problem}, prove weak and strong duality theorems, develop the complete four-part \textbf{Karush-Kuhn-Tucker (KKT) optimality conditions}, and derive the Hard-Margin and Soft-Margin Support Vector Machine. This chapter transforms constrained optimization from a theoretical framework into one of the most powerful and robust classification algorithms in machine learning.

\subsection{Prerequisites \& Scaffolding}
Mastery of this chapter synthesizes convex functions, hyperplanes, and multivariable calculus as outlined in Table~\ref{tab:w10_prereqs}.

\begin{table}[htbp]
\centering
\small
\begin{tabularx}{\textwidth}{lXX}
\toprule
\textbf{Foundational Area} & \textbf{Core Concepts Required} & \textbf{Role in Machine Learning Foundations} \\
\midrule
Convex Analysis & Convex sets, epigraphs, supporting hyperplanes, Slater's condition & Guaranteeing zero duality gap and establishing necessary and sufficient KKT conditions \\
\addlinespace
Vector Calculus & Multivariable gradients, partial derivatives, stationarity & Formulating the Lagrangian stationarity conditions $\nabla_{\mathbf{w}}\mathcal{L} = \mathbf{0}$ \\
\addlinespace
Linear Algebra & Inner products, Gram kernel matrices $K = X X^T$, positive definiteness & Deriving the dual Quadratic Program (QP) and unlocking the kernel trick \\
\bottomrule
\end{tabularx}
\caption{Prerequisite Scaffolding for Chapter 10.}
\label{tab:w10_prereqs}
\end{table}

\subsection{Core Competencies}
By the conclusion of this chapter, students will be able to:
\begin{itemize}
    \item Formulate the Lagrange dual function $g(\boldsymbol{\lambda}, \boldsymbol{\nu})$ and prove that it is universally concave regardless of primal convexity.
    \item State and verify Weak Duality, Strong Duality, and Slater's constraint qualification.
    \item Formulate and solve the complete four-part Karush-Kuhn-Tucker (KKT) optimality conditions.
    \item Rigorously derive the Hard-Margin and Soft-Margin Support Vector Machine from geometric margin maximization.
    \item Apply complementary slackness to prove why non-support vectors exert zero influence on the decision boundary.
    \item Derive the box constraints $0 \le \lambda_i \le C$ and classify support vectors into free margin points and bounded margin violators.
    \item Formulate the Dual Quadratic Program and implement an SVM solver using scientific computing libraries.
\end{itemize}

\subsection{Notation \& Mathematical Conventions}
Throughout this chapter, we adhere to the standard mathematical conventions:
\begin{itemize}
    \item Primal objective functional: $f_0(\mathbf{w}): \mathbb{R}^d \to \mathbb{R}$.
    \item Inequality constraint functions: $f_i(\mathbf{w}) \le 0$ ($i = 1, \dots, m$).
    \item Equality constraint functions: $h_j(\mathbf{w}) = 0$ ($j = 1, \dots, p$).
    \item Lagrangian functional: $\mathcal{L}(\mathbf{w}, \boldsymbol{\lambda}, \boldsymbol{\nu}) \triangleq f_0(\mathbf{w}) + \sum_{i=1}^m \lambda_i f_i(\mathbf{w}) + \sum_{j=1}^p \nu_j h_j(\mathbf{w})$.
    \item Lagrange dual function: $g(\boldsymbol{\lambda}, \boldsymbol{\nu}) \triangleq \inf_{\mathbf{w} \in \mathbb{R}^d} \mathcal{L}(\mathbf{w}, \boldsymbol{\lambda}, \boldsymbol{\nu})$.
    \item Primal optimal value: $p^* \triangleq \inf_{\mathbf{w} \in \text{feasible}} f_0(\mathbf{w})$.
    \item Dual optimal value: $d^* \triangleq \sup_{\boldsymbol{\lambda} \ge \mathbf{0}, \boldsymbol{\nu}} g(\boldsymbol{\lambda}, \boldsymbol{\nu})$.
    \item Binary class labels: $y_i \in \{-1, +1\}$.
    \item Slack variables: $\xi_i \ge 0$, measuring margin violations in Soft-Margin SVM.
    \item Regularization penalty parameter: $C > 0$, balancing margin width against margin slack penalties.
\end{itemize}

% ----------------------------------------------------------------------
% PART 2: MATHEMATICAL & COMPUTATIONAL FOUNDATIONS
% ----------------------------------------------------------------------
\section{Mathematical \& Computational Foundations}

\subsection{The Primal and Dual Problems}
Consider the canonical constrained continuous optimization problem in standard form:
\begin{equation}
\min_{\mathbf{w} \in \mathbb{R}^d} f_0(\mathbf{w}) \quad \text{subject to } \begin{cases} f_i(\mathbf{w}) \le 0, & i = 1, \dots, m \\ h_j(\mathbf{w}) = 0, & j = 1, \dots, p \end{cases}
\label{eq:standard_primal}
\end{equation}
with optimal value $p^* = \inf \{f_0(\mathbf{w}) \mid f_i(\mathbf{w}) \le 0, \, h_j(\mathbf{w}) = 0\}$.

\begin{definitionbox}{The Lagrange Dual Function}
The \textbf{Lagrange Dual Function} $g: \mathbb{R}^m \times \mathbb{R}^p \to \mathbb{R} \cup \{-\infty\}$ is the infimum of the Lagrangian over the primal parameter vector $\mathbf{w} \in \mathbb{R}^d$:
\begin{equation}
g(\boldsymbol{\lambda}, \boldsymbol{\nu}) \triangleq \inf_{\mathbf{w} \in \mathbb{R}^d} \left( f_0(\mathbf{w}) + \sum_{i=1}^m \lambda_i f_i(\mathbf{w}) + \sum_{j=1}^p \nu_j h_j(\mathbf{w}) \right)
\label{eq:dual_func_def}
\end{equation}
where $\boldsymbol{\lambda} \in \mathbb{R}^m$ is the dual multiplier vector associated with the inequality constraints, and $\boldsymbol{\nu} \in \mathbb{R}^p$ is the dual multiplier vector associated with the equality constraints.
\end{definitionbox}

\begin{theorembox}{Universal Concavity of the Dual Function}
The Lagrange dual function $g(\boldsymbol{\lambda}, \boldsymbol{\nu})$ is always concave in $(\boldsymbol{\lambda}, \boldsymbol{\nu})$, even if the primal objective $f_0$ and the constraint functions $f_i, h_j$ are wildly non-convex.
\end{theorembox}

\begin{proof}
For any fixed parameter vector $\mathbf{w} \in \mathbb{R}^d$, define the affine function $\psi_{\mathbf{w}}: \mathbb{R}^m \times \mathbb{R}^p \to \mathbb{R}$:
\begin{equation}
\psi_{\mathbf{w}}(\boldsymbol{\lambda}, \boldsymbol{\nu}) \triangleq f_0(\mathbf{w}) + \sum_{i=1}^m \lambda_i f_i(\mathbf{w}) + \sum_{j=1}^p \nu_j h_j(\mathbf{w})
\end{equation}
Notice that for a fixed $\mathbf{w}$, $\psi_{\mathbf{w}}$ is strictly linear (affine) with respect to the dual multipliers $(\boldsymbol{\lambda}, \boldsymbol{\nu})$.
By definition, the dual function is the pointwise infimum of this family of affine functions:
\begin{equation}
g(\boldsymbol{\lambda}, \boldsymbol{\nu}) = \inf_{\mathbf{w} \in \mathbb{R}^d} \psi_{\mathbf{w}}(\boldsymbol{\lambda}, \boldsymbol{\nu})
\end{equation}
Let $\theta \in [0, 1]$ and consider two arbitrary dual points $(\boldsymbol{\lambda}_1, \boldsymbol{\nu}_1)$ and $(\boldsymbol{\lambda}_2, \boldsymbol{\nu}_2)$:
\begin{align}
g(\theta \boldsymbol{\lambda}_1 + (1-\theta)\boldsymbol{\lambda}_2, \theta \boldsymbol{\nu}_1 + (1-\theta)\boldsymbol{\nu}_2)
&= \inf_{\mathbf{w}} \psi_{\mathbf{w}}(\theta \boldsymbol{\lambda}_1 + (1-\theta)\boldsymbol{\lambda}_2, \theta \boldsymbol{\nu}_1 + (1-\theta)\boldsymbol{\nu}_2) \nonumber \\
&= \inf_{\mathbf{w}} \left[ \theta \psi_{\mathbf{w}}(\boldsymbol{\lambda}_1, \boldsymbol{\nu}_1) + (1-\theta) \psi_{\mathbf{w}}(\boldsymbol{\lambda}_2, \boldsymbol{\nu}_2) \right] \nonumber \\
&\ge \theta \inf_{\mathbf{w}} \psi_{\mathbf{w}}(\boldsymbol{\lambda}_1, \boldsymbol{\nu}_1) + (1-\theta) \inf_{\mathbf{w}'} \psi_{\mathbf{w}'}(\boldsymbol{\lambda}_2, \boldsymbol{\nu}_2) \nonumber \\
&= \theta g(\boldsymbol{\lambda}_1, \boldsymbol{\nu}_1) + (1-\theta) g(\boldsymbol{\lambda}_2, \boldsymbol{\nu}_2)
\end{align}
Because the secant chord lies everywhere below the dual function graph, $g(\boldsymbol{\lambda}, \boldsymbol{\nu})$ is universally concave.
\end{proof}

\begin{theorembox}{Weak Duality Theorem}
For any dual feasible multipliers $\boldsymbol{\lambda} \ge \mathbf{0}$ and $\boldsymbol{\nu} \in \mathbb{R}^p$, the dual function value provides a universal lower bound on the primal optimal value:
\begin{equation}
g(\boldsymbol{\lambda}, \boldsymbol{\nu}) \le p^*
\label{eq:weak_duality}
\end{equation}
Consequently, the optimal dual value $d^* \triangleq \sup_{\boldsymbol{\lambda} \ge \mathbf{0}, \boldsymbol{\nu}} g(\boldsymbol{\lambda}, \boldsymbol{\nu})$ satisfies:
\begin{equation}
d^* \le p^*
\end{equation}
The non-negative difference $p^* - d^* \ge 0$ is the \textbf{Duality Gap}.
\end{theorembox}

\begin{proof}
Let $\tilde{\mathbf{w}}$ be any feasible point for the primal problem, meaning $f_i(\tilde{\mathbf{w}}) \le 0$ for all $i$ and $h_j(\tilde{\mathbf{w}}) = 0$ for all $j$.
Because $\lambda_i \ge 0$ and $f_i(\tilde{\mathbf{w}}) \le 0$, the product $\lambda_i f_i(\tilde{\mathbf{w}}) \le 0$. Thus:
\begin{equation}
\mathcal{L}(\tilde{\mathbf{w}}, \boldsymbol{\lambda}, \boldsymbol{\nu}) = f_0(\tilde{\mathbf{w}}) + \sum_{i=1}^m \underbrace{\lambda_i f_i(\tilde{\mathbf{w}})}_{\le 0} + \sum_{j=1}^p \underbrace{\nu_j h_j(\tilde{\mathbf{w}})}_{= 0} \le f_0(\tilde{\mathbf{w}})
\end{equation}
By definition of the infimum over all $\mathbf{w} \in \mathbb{R}^d$:
\begin{equation}
g(\boldsymbol{\lambda}, \boldsymbol{\nu}) = \inf_{\mathbf{w}} \mathcal{L}(\mathbf{w}, \boldsymbol{\lambda}, \boldsymbol{\nu}) \le \mathcal{L}(\tilde{\mathbf{w}}, \boldsymbol{\lambda}, \boldsymbol{\nu}) \le f_0(\tilde{\mathbf{w}})
\end{equation}
Because $g(\boldsymbol{\lambda}, \boldsymbol{\nu}) \le f_0(\tilde{\mathbf{w}})$ holds for \emph{every} primal feasible point $\tilde{\mathbf{w}}$, taking the infimum over all feasible $\tilde{\mathbf{w}}$ establishes $g(\boldsymbol{\lambda}, \boldsymbol{\nu}) \le p^*$.
\end{proof}

\begin{ponderbox}{Why Must the Dual Problem Always Be Concave?}
Why is the universal concavity of the dual function $g(\boldsymbol{\lambda}, \boldsymbol{\nu})$ so profound for continuous optimization?
\tcblower
\textbf{The Socratic Resolution:} Maximizing a concave function over a convex set (the non-negative orthant $\boldsymbol{\lambda} \ge \mathbf{0}$) is, by definition, a \textbf{convex optimization problem}! 

Even if the primal problem $f_0(\mathbf{w})$ is wildly non-convex, multimodal, discontinuous, or NP-hard, its Lagrange dual problem $\max_{\boldsymbol{\lambda} \ge \mathbf{0}, \boldsymbol{\nu}} g(\boldsymbol{\lambda}, \boldsymbol{\nu})$ is \emph{always} a well-behaved convex program without spurious local maxima. Solving the dual problem via subgradient ascent or interior-point methods produces a certifiable, non-trivial global lower bound $d^* \le p^*$ on the unknown global minimum.
\end{ponderbox}

\subsection{Strong Duality and the KKT Optimality Conditions}

\begin{definitionbox}{Slater's Constraint Qualification}
Let the primal problem~\eqref{eq:standard_primal} be convex: the objective $f_0$ and inequality constraints $f_1, \dots, f_m$ are convex, and equality constraints $h_j(\mathbf{w}) = \mathbf{a}_j^T \mathbf{w} - b_j$ are affine.
If there exists a \textbf{strictly feasible point} $\tilde{\mathbf{w}} \in \mathbb{R}^d$ such that:
\begin{equation}
f_i(\tilde{\mathbf{w}}) < 0 \quad (\forall i \in \{1, \dots, m\}) \quad \text{and} \quad A\tilde{\mathbf{w}} = \mathbf{b}
\label{eq:slaters_condition}
\end{equation}
then \textbf{Strong Duality} holds: the duality gap is zero ($d^* = p^*$), and the dual supremum is attained at some dual optimal $(\boldsymbol{\lambda}^*, \boldsymbol{\nu}^*)$.
\end{definitionbox}

\begin{theorembox}{Karush-Kuhn-Tucker (KKT) Optimality Conditions}
Let $f_0, f_1, \dots, f_m$ be continuously differentiable convex functions, and let $h_j$ be affine. If Slater's condition holds, then $\mathbf{w}^*$ is primal optimal and $(\boldsymbol{\lambda}^*, \boldsymbol{\nu}^*)$ is dual optimal if and only if they satisfy the complete four-part KKT conditions:
\begin{enumerate}
    \item \textbf{Primal Feasibility:}
    \begin{align}
    f_i(\mathbf{w}^*) &\le 0 \quad (i = 1, \dots, m) \\
    h_j(\mathbf{w}^*) &= 0 \quad (j = 1, \dots, p)
    \end{align}
    \item \textbf{Dual Feasibility:}
    \begin{equation}
    \lambda_i^* \ge 0 \quad (i = 1, \dots, m)
    \end{equation}
    \item \textbf{Complementary Slackness:}
    \begin{equation}
    \lambda_i^* f_i(\mathbf{w}^*) = 0 \quad (i = 1, \dots, m)
    \label{eq:comp_slackness}
    \end{equation}
    \item \textbf{Lagrangian Stationarity:}
    \begin{equation}
    \nabla_{\mathbf{w}} \mathcal{L}(\mathbf{w}^*, \boldsymbol{\lambda}^*, \boldsymbol{\nu}^*) = \nabla f_0(\mathbf{w}^*) + \sum_{i=1}^m \lambda_i^* \nabla f_i(\mathbf{w}^*) + \sum_{j=1}^p \nu_j^* \nabla h_j(\mathbf{w}^*) = \mathbf{0}
    \label{eq:kkt_stationarity}
    \end{equation}
\end{enumerate}
\end{theorembox}

\begin{ponderbox}{Geometric Anatomy of the Three Multiplier Regimes}
In a constrained optimization problem with inequality constraints $f_i(\mathbf{w}) \le 0$, how does the complementary slackness condition $\lambda_i^* f_i(\mathbf{w}^*) = 0$ act as a logical and geometric switch?
\tcblower
\textbf{The Socratic Resolution:} Because $\lambda_i^* \ge 0$ and $f_i(\mathbf{w}^*) \le 0$, their product vanishes if and only if at least one of the two factors is zero. This partitions all training points into two mutually exclusive geometric regimes:
\begin{itemize}
    \item \textbf{Strictly Inactive Regime ($f_i(\mathbf{w}^*) < 0 \implies \lambda_i^* = 0$):} The point lies in the interior of the feasible set, away from the boundary. Tightening or removing this constraint locally would not change the optimal solution. The constraint exerts zero contact force: $\lambda_i^* = 0$.
    \item \textbf{Active Boundary Regime ($\lambda_i^* > 0 \implies f_i(\mathbf{w}^*) = 0$):} The point presses directly against the constraint boundary. If the boundary were relaxed, the optimal solution would move into the newly opened territory. The multiplier $\lambda_i^* > 0$ represents the exact normal reaction force required to prevent the solution from drifting across the boundary.
\end{itemize}
\end{ponderbox}

% ----------------------------------------------------------------------
% PART 3: VISUAL & GEOMETRIC ANCHORING
% ----------------------------------------------------------------------
\section{Visual \& Geometric Anchoring}

In this section, we anchor the mathematical apparatus of duality and SVMs in three native TikZ visualizations.

\subsection{The Geometry of Support Vector Machines and Margins}
Figure~\ref{fig:svm_geometry} illustrates the canonical maximum margin hyperplane separating two linearly separable classes. The decision boundary is determined strictly by the circled support vectors that lie on the margin boundaries; all interior points have $\lambda_i^* = 0$.

\begin{figure}[htbp]
\centering
\begin{tikzpicture}[scale=0.95, font=\sffamily\footnotesize]
    % Separating hyperplane w^T x + b = 0
    \draw[line width=2pt, color=primarydark] (-3.8, -1.6) -- (3.8, 2.6);
    \node[above, color=primarydark] at (2.4, 2.2) {\textbf{Decision Boundary: } $\mathbf{w}^T \mathbf{x} + b = 0$};

    % Positive margin hyperplane w^T x + b = +1
    \draw[dashed, line width=1.4pt, color=tealborder] (-4.2, -0.5) -- (3.2, 3.7);
    \node[above, color=tealborder] at (-2.5, 1.2) {$\mathbf{w}^T \mathbf{x} + b = +1$};

    % Negative margin hyperplane w^T x + b = -1
    \draw[dashed, line width=1.4pt, color=coralborder] (-3.2, -2.7) -- (4.2, 1.5);
    \node[below, color=coralborder] at (2.5, 0.2) {$\mathbf{w}^T \mathbf{x} + b = -1$};

    % Class +1 data points (teal squares)
    \filldraw[color=tealborder] (-2.2, 2.2) rectangle ++(5pt, 5pt);
    \filldraw[color=tealborder] (0.0, 3.2) rectangle ++(5pt, 5pt);
    \filldraw[color=tealborder] (-1.2, 3.4) rectangle ++(5pt, 5pt);
    
    % Positive support vectors on margin
    \coordinate (SV1) at (-0.6, 1.6);
    \filldraw[color=tealborder] (SV1) rectangle ++(6pt, 6pt);
    \draw[line width=1.5pt, color=forest] ($(SV1) + (3pt, 3pt)$) circle (9pt);
    \node[above, color=forest] at ($(SV1) + (0, 0.35)$) {\textbf{Support Vector } ($\lambda_1^* > 0$)};

    % Class -1 data points (coral circles)
    \filldraw[color=coralborder] (0.6, -2.2) circle (3pt);
    \filldraw[color=coralborder] (2.2, -1.8) circle (3pt);
    \filldraw[color=coralborder] (-1.2, -3.1) circle (3pt);
    
    % Negative support vectors on margin
    \coordinate (SV2) at (0.6, -0.6);
    \filldraw[color=coralborder] (SV2) circle (3.5pt);
    \draw[line width=1.5pt, color=traprad] (SV2) circle (9pt);
    \node[below right, color=traprad] at ($(SV2) + (0.1, -0.1)$) {\textbf{Support Vector } ($\lambda_2^* > 0$)};

    % Margin width arrow
    \draw[<->, >=stealth, line width=1.6pt, color=traprad] (-1.1, 1.3) -- (-0.1, -0.2) 
        node[midway, left] {\textbf{Margin } $\gamma = \frac{2}{\|\mathbf{w}\|_2}$};

    % Normal vector w
    \draw[->, >=stealth, line width=2pt, color=primarydark] (0, 0.5) -- (-0.8, 1.9) 
        node[left] {$\mathbf{w}^* = \sum \lambda_i^* y_i \mathbf{x}_i$};
\end{tikzpicture}
\caption{The Geometry of Support Vector Machines: Maximizing the geometric margin $\frac{2}{\|\mathbf{w}\|_2}$ forces the weight vector to depend strictly on the boundary support vectors. Training samples lying strictly inside the class territories have $\lambda_i^* = 0$ and exert zero influence on the separating hyperplane.}
\label{fig:svm_geometry}
\end{figure}

\subsection{The Primal-Dual Geometry and Duality Gap}
Figure~\ref{fig:duality_gap} illustrates the geometric relationship between the primal optimal value $p^*$, the dual function $g(\boldsymbol{\lambda})$, the duality gap $p^* - d^*$, and Slater's condition.

\begin{figure}[htbp]
\centering
\begin{tikzpicture}[scale=1.05, font=\sffamily\footnotesize]
    % Axes
    \draw[->, thick, color=primarydark!60] (-2.5, 0) -- (3.5, 0) node[right] {$u = f_1(\mathbf{w})$};
    \draw[->, thick, color=primarydark!60] (0, -1.5) -- (0, 4.0) node[above] {$t = f_0(\mathbf{w})$};

    % Achievable set G = {(f1(w), f0(w))}
    \draw[thick, fill=tealborder!15, color=tealborder] 
        (-2.0, 3.5) to[out=-40, in=170] (0.0, 1.8)
                    to[out=-10, in=190] (2.5, 1.2)
                    to[out=60, in=-60] (3.0, 3.8)
                    to[out=170, in=10] cycle;
    \node[color=tealborder!90!black] at (1.5, 2.8) {\large $\mathcal{G} = \{(f_1(\mathbf{w}), f_0(\mathbf{w}))\}$};

    % Primal optimal point p* at u = 0
    \coordinate (Pstar) at (0.0, 1.8);
    \filldraw[color=forest] (Pstar) circle (3pt) node[left=2pt] {$p^* = \inf f_0(\mathbf{w})$};

    % Supporting Hyperplane with slope -lambda
    % Equation: t + lambda * u >= d* => t = d* - lambda * u
    % For lambda = 0.5, d* = 1.0 => t = 1.0 - 0.5 * u
    \draw[very thick, color=coralborder, domain=-2.2:3.2] 
        plot (\x, {1.0 - 0.5*\x}) node[right] {\textbf{Dual Supporting Line } $t + \lambda u = d^*$};

    % Dual intercept d* on vertical axis
    \coordinate (Dstar) at (0, 1.0);
    \filldraw[color=coralborder] (Dstar) circle (3pt) node[below left=2pt] {$d^* = \sup g(\lambda)$};

    % Duality Gap annotation
    \draw[<->, >=stealth, line width=1.5pt, color=traprad] (0.3, 1.0) -- (0.3, 1.8);
    \node[right, color=traprad, align=left] at (0.35, 1.4) {
        \textbf{Duality Gap } $p^* - d^* \ge 0$ \\
        \scriptsize (Zero if Slater's Condition holds)
    };
\end{tikzpicture}
\caption{The Geometry of Lagrange Duality: The set $\mathcal{G}$ represents the set of achievable constraint-objective values $(f_1(\mathbf{w}), f_0(\mathbf{w}))$. The primal minimum $p^*$ is the lowest point of $\mathcal{G}$ on or to the left of the vertical axis ($u \le 0$). The dual function $g(\lambda)$ represents the intercept of a supporting hyperplane with normal $(\lambda, 1)$ supporting $\mathcal{G}$ from below. Weak duality states that $d^* \le p^*$. Under Slater's condition, the supporting hyperplane touches $\mathcal{G}$ directly at $p^*$, collapsing the duality gap to zero ($p^* = d^*$).}
\label{fig:duality_gap}
\end{figure}

\subsection{Complementary Slackness as a Geometric Switch}
Figure~\ref{fig:slackness_switch} reveals how complementary slackness partitions parameter space into inactive non-support vectors and active margin constraints.

\begin{figure}[htbp]
\centering
\begin{tikzpicture}[scale=0.95, font=\sffamily\footnotesize]
    % Panel 1: Inactive Constraint (Non-Support Vector)
    \begin{scope}[shift={(-4.5, 0)}]
        \draw[thick, fill=primarydark!10, color=primarydark] (-2.2, -2.0) rectangle (2.2, 2.0);
        \node[above, color=primarydark] at (0, 2.0) {\textbf{Feasible Region } $f_i(\mathbf{w}) \le 0$};

        % Solution w* inside interior
        \coordinate (W1) at (0, 0.2);
        \filldraw[color=forest] (W1) circle (3pt) node[above] {$\mathbf{w}^*$};
        \node[below, color=forest, align=center] at (0, -0.2) {$f_i(\mathbf{w}^*) < 0$ \\ (\textbf{Strictly Interior})};

        % Constraint boundary far away
        \draw[line width=1.8pt, color=traprad] (1.8, -2.0) -- (1.8, 2.0);
        \node[rotate=90, above, color=traprad] at (1.8, 0) {\scriptsize Boundary $f_i = 0$};

        \node[draw=forest, fill=white, rounded corners=3pt, align=center] at (0, -2.8) {
            \textbf{Inactive Constraint:} \\
            $\lambda_i^* = 0$ \\
            \scriptsize Zero contact force
        };
    \end{scope}

    % Panel 2: Active Constraint (Support Vector)
    \begin{scope}[shift={(4.5, 0)}]
        \draw[thick, fill=primarydark!10, color=primarydark] (-2.2, -2.0) rectangle (2.2, 2.0);
        \node[above, color=primarydark] at (0, 2.0) {\textbf{Feasible Region } $f_i(\mathbf{w}) \le 0$};

        % Solution w* directly on boundary
        \coordinate (W2) at (1.8, 0.2);
        \filldraw[color=traprad] (W2) circle (3pt);
        \node[left, color=traprad, align=right] at (1.7, 0.2) {$\mathbf{w}^*$ on boundary \\ $f_i(\mathbf{w}^*) = 0$};

        % Boundary
        \draw[line width=1.8pt, color=traprad] (1.8, -2.0) -- (1.8, 2.0);
        \node[rotate=90, above, color=traprad] at (1.8, 1.2) {\scriptsize Boundary $f_i = 0$};

        % Normal force vector lambda * grad f
        \draw[->, >=stealth, line width=2pt, color=coralborder] (W2) -- ($(W2) + (-1.4, 0)$) 
            node[below left] {$\lambda_i^* \nabla f_i$};

        \node[draw=traprad, fill=white, rounded corners=3pt, align=center] at (0, -2.8) {
            \textbf{Active Constraint:} \\
            $\lambda_i^* > 0$ \\
            \scriptsize Support Vector reaction force
        };
    \end{scope}
\end{tikzpicture}
\caption{The Complementary Slackness Geometric Switch: For every inequality constraint, $\lambda_i^* f_i(\mathbf{w}^*) = 0$. Left: If the optimal parameter vector $\mathbf{w}^*$ lies strictly in the interior ($f_i < 0$), the constraint is non-binding and the dual multiplier is identically zero ($\lambda_i^* = 0$). Right: If the multiplier is strictly positive ($\lambda_i^* > 0$), the parameter vector is pressed firmly against the boundary ($f_i = 0$), defining an active support vector.}
\label{fig:slackness_switch}
\end{figure}

% ----------------------------------------------------------------------
% PART 4: DEEP INTUITION & CONCEPTUAL DERIVATIONS
% ----------------------------------------------------------------------
\section{Deep Intuition \& Conceptual Derivations}

\subsection{Derivation of the Dual Support Vector Machine}
Given a training dataset $\mathcal{D} = \{(\mathbf{x}_i, y_i)\}_{i=1}^N$ with feature vectors $\mathbf{x}_i \in \mathbb{R}^d$ and binary class labels $y_i \in \{-1, +1\}$, a linear separating hyperplane is defined by normal vector $\mathbf{w} \in \mathbb{R}^d$ and bias intercept $b \in \mathbb{R}$:
\begin{equation}
\mathcal{H} = \{\mathbf{x} \in \mathbb{R}^d \mid \mathbf{w}^T \mathbf{x} + b = 0\}
\end{equation}
The signed perpendicular Euclidean distance from any sample $\mathbf{x}_i$ to the hyperplane is:
\begin{equation}
\gamma_i = \frac{y_i (\mathbf{w}^T \mathbf{x}_i + b)}{\|\mathbf{w}\|_2}
\end{equation}
The \textbf{canonical hyperplane} fixes the scale of $(\mathbf{w}, b)$ such that the closest training samples to the boundary achieve a functional margin of exactly 1:
\begin{equation}
\min_{i=1,\dots,N} y_i (\mathbf{w}^T \mathbf{x}_i + b) = 1
\end{equation}
Then the total geometric margin separating the two classes is:
\begin{equation}
\gamma = \frac{1}{\|\mathbf{w}\|_2} + \frac{1}{\|\mathbf{w}\|_2} = \frac{2}{\|\mathbf{w}\|_2}
\end{equation}
Maximizing the geometric margin $\frac{2}{\|\mathbf{w}\|_2}$ is mathematically equivalent to minimizing $\frac{1}{2}\|\mathbf{w}\|_2^2$:
\begin{equation}
\min_{\mathbf{w} \in \mathbb{R}^d, b \in \mathbb{R}} \frac{1}{2} \|\mathbf{w}\|_2^2 \quad \text{subject to } y_i (\mathbf{w}^T \mathbf{x}_i + b) \ge 1 \quad (i = 1, \dots, N)
\label{eq:hard_svm_primal}
\end{equation}

\subsubsection{Formulating the Lagrangian}
Rewriting the $N$ margin constraints in standard format $1 - y_i(\mathbf{w}^T \mathbf{x}_i + b) \le 0$:
\begin{equation}
\mathcal{L}(\mathbf{w}, b, \boldsymbol{\lambda}) = \frac{1}{2}\|\mathbf{w}\|_2^2 - \sum_{i=1}^N \lambda_i \left[ y_i(\mathbf{w}^T \mathbf{x}_i + b) - 1 \right]
\label{eq:svm_lagrangian}
\end{equation}
with dual multipliers $\lambda_i \ge 0$.

\subsubsection{Primal Stationarity Conditions}
Setting the partial gradients of $\mathcal{L}$ with respect to the primal parameters $\mathbf{w}$ and $b$ to zero:
\begin{align}
\nabla_{\mathbf{w}} \mathcal{L} &= \mathbf{w} - \sum_{i=1}^N \lambda_i y_i \mathbf{x}_i = \mathbf{0} \implies \mathbf{w}^* = \sum_{i=1}^N \lambda_i y_i \mathbf{x}_i \label{eq:svm_w_stat} \\
\frac{\partial \mathcal{L}}{\partial b} &= -\sum_{i=1}^N \lambda_i y_i = 0 \implies \sum_{i=1}^N \lambda_i y_i = 0 \label{eq:svm_b_stat}
\end{align}

\subsubsection{Constructing the Dual Objective}
Substituting $\mathbf{w} = \sum_{i=1}^N \lambda_i y_i \mathbf{x}_i$ and $\sum_{i=1}^N \lambda_i y_i = 0$ back into the Lagrangian:
\begin{align}
\mathcal{L} &= \frac{1}{2} \left( \sum_{i=1}^N \lambda_i y_i \mathbf{x}_i \right)^T \left( \sum_{j=1}^N \lambda_j y_j \mathbf{x}_j \right) - \sum_{i=1}^N \lambda_i y_i \left( \sum_{j=1}^N \lambda_j y_j \mathbf{x}_j \right)^T \mathbf{x}_i - b \underbrace{\sum_{i=1}^N \lambda_i y_i}_{= 0} + \sum_{i=1}^N \lambda_i \nonumber \\
&= \frac{1}{2} \sum_{i=1}^N \sum_{j=1}^N \lambda_i \lambda_j y_i y_j (\mathbf{x}_i^T \mathbf{x}_j) - \sum_{i=1}^N \sum_{j=1}^N \lambda_i \lambda_j y_i y_j (\mathbf{x}_i^T \mathbf{x}_j) + \sum_{i=1}^N \lambda_i \nonumber \\
&= \sum_{i=1}^N \lambda_i - \frac{1}{2} \sum_{i=1}^N \sum_{j=1}^N \lambda_i \lambda_j y_i y_j (\mathbf{x}_i^T \mathbf{x}_j)
\end{align}
This yields the celebrated \textbf{Dual Support Vector Machine Quadratic Program}:
\begin{equation}
\max_{\boldsymbol{\lambda} \in \mathbb{R}^N} \sum_{i=1}^N \lambda_i - \frac{1}{2} \sum_{i=1}^N \sum_{j=1}^N \lambda_i \lambda_j y_i y_j (\mathbf{x}_i^T \mathbf{x}_j) \quad \text{subject to } \begin{cases} \lambda_i \ge 0, & \forall i \in \{1, \dots, N\} \\ \sum_{i=1}^N \lambda_i y_i = 0 & \end{cases}
\label{eq:svm_dual_qp}
\end{equation}

\subsection{Soft-Margin SVM and the Box Constraint $0 \le \lambda_i \le C$}
When classes overlap in feature space, no linear separating hyperplane exists. The Soft-Margin SVM introduces non-negative slack variables $\xi_i \ge 0$:
\begin{equation}
\min_{\mathbf{w}, b, \boldsymbol{\xi}} \frac{1}{2}\|\mathbf{w}\|_2^2 + C \sum_{i=1}^N \xi_i \quad \text{subject to } \begin{cases} y_i(\mathbf{w}^T \mathbf{x}_i + b) \ge 1 - \xi_i, & \forall i \\ \xi_i \ge 0, & \forall i \end{cases}
\label{eq:soft_svm_primal}
\end{equation}
The Lagrangian introduces multipliers $\lambda_i \ge 0$ for the margin constraint and $\mu_i \ge 0$ for the non-negativity of slack:
\begin{equation}
\mathcal{L}(\mathbf{w}, b, \boldsymbol{\xi}, \boldsymbol{\lambda}, \boldsymbol{\mu}) = \frac{1}{2}\|\mathbf{w}\|_2^2 + C \sum_{i=1}^N \xi_i - \sum_{i=1}^N \lambda_i \left[ y_i(\mathbf{w}^T \mathbf{x}_i + b) - 1 + \xi_i \right] - \sum_{i=1}^N \mu_i \xi_i
\end{equation}
Stationarity with respect to the slack variable $\xi_i$:
\begin{equation}
\frac{\partial \mathcal{L}}{\partial \xi_i} = C - \lambda_i - \mu_i = 0 \implies \lambda_i + \mu_i = C \implies \lambda_i = C - \mu_i
\end{equation}
Since $\mu_i \ge 0$, it follows that $\lambda_i \le C$. Combined with $\lambda_i \ge 0$, this establishes the famous \textbf{box constraint}:
\begin{equation}
0 \le \lambda_i \le C \quad \forall i \in \{1, \dots, N\}
\label{eq:box_constraint}
\end{equation}
The dual objective function for Soft-Margin SVM is identical to the Hard-Margin dual~\eqref{eq:svm_dual_qp}; the only mathematical modification is that the upper boundary of the multipliers is clipped at $C$.

\begin{ponderbox}{The Soft-Margin Trade-Off: What Does $C$ Actually Control?}
What is the fundamental bias-variance trade-off governed by the SVM penalty hyperparameter $C$?
\tcblower
\textbf{The Socratic Resolution:} Hyperparameter $C$ acts as the \textbf{economic exchange rate} between margin width and empirical training error:
\begin{itemize}
    \item \textbf{Large $C$ ($C \to \infty$):} Massive penalty on slack violations. The algorithm prioritizes classifying every training point correctly, forcing slacks $\xi_i \to 0$. The margin width narrows, making the model highly sensitive to individual outliers and risking severe overfitting.
    \item \textbf{Small $C$ ($C \to 0$):} Minimal penalty on slack violations. The algorithm prioritizes maximizing the margin width $\frac{2}{\|\mathbf{w}\|_2}$, readily tolerating training errors and margin violations to find a broader, more robust separating corridor. This increases regularization and prevents overfitting.
\end{itemize}
\end{ponderbox}

\begin{ponderbox}{Primal vs.\ Dual Computational Complexity: When to Solve Which?}
When should an engineer solve the Primal SVM formulation, and when should they solve the Dual SVM formulation?
\tcblower
\textbf{The Socratic Resolution:} The choice depends strictly on the ratio between the number of training samples $N$ and the feature dimensionality $d$:
\begin{itemize}
    \item \textbf{Primal SVM ($\mathcal{O}(d)$ to $\mathcal{O}(d^3)$):} The primal variable vector $\mathbf{w} \in \mathbb{R}^d$ has dimension $d$. When training on massive tabular or text datasets with millions of samples ($N \gg d$, e.g., $N = 10^7, d = 100$), solving the primal problem via subgradient descent or stochastic dual coordinate ascent takes linear time $\mathcal{O}(N d)$ without storing an $N \times N$ matrix.
    \item \textbf{Dual SVM ($\mathcal{O}(N^2)$ to $\mathcal{O}(N^3)$):} The dual variable vector $\boldsymbol{\lambda} \in \mathbb{R}^N$ has dimension $N$. When the feature space is very high-dimensional or infinite ($d \gg N$, e.g., genomics, RBF kernels where $d = \infty$), the primal cannot be solved. The dual depends strictly on the $N \times N$ Gram kernel matrix $K_{ij} = k(\mathbf{x}_i, \mathbf{x}_j)$, making non-linear classification computationally feasible.
\end{itemize}
\end{ponderbox}

% ----------------------------------------------------------------------
% PART 5: MULTI-PAGE WORKED EXAMPLES
% ----------------------------------------------------------------------
\section{Multi-Page Breakable Numerical Worked Examples}

\begin{examplebox}[breakable]{Step-by-Step Resolution of a Constrained Quadratic via KKT Conditions}
\textbf{Problem Statement:}
Consider the continuous optimization problem:
\begin{equation}
\min_{x, y \in \mathbb{R}} f_0(x, y) = x^2 + y^2 \quad \text{subject to } x + y \ge 4
\end{equation}
Rewrite the inequality constraint in standard format $f_1(x, y) = 4 - x - y \le 0$.

\textbf{Task Requirements:}
\begin{enumerate}
    \item Formulate the Lagrangian $\mathcal{L}(x, y, \lambda)$ with dual variable $\lambda \ge 0$.
    \item Explicitly write down the four KKT conditions.
    \item Test Case 1: $\lambda = 0$ (inactive constraint). Determine the solution and prove why it violates primal feasibility.
    \item Test Case 2: $\lambda > 0$ (active constraint). Analytically compute $(x^*, y^*)$ and $\lambda^*$, and verify that all four KKT conditions are rigorously satisfied.
    \item Compute the optimal primal value $p^*$ and dual value $d^*$, and verify that the duality gap is zero.
\end{enumerate}

\tcblower

\textbf{Step 1: The Lagrangian}
\begin{equation}
\mathcal{L}(x, y, \lambda) = x^2 + y^2 + \lambda(4 - x - y)
\end{equation}
with dual feasibility requirement $\lambda \ge 0$.

\textbf{Step 2: The Four KKT Conditions}
\begin{enumerate}
    \item \textbf{Primal Feasibility:}
    \begin{equation}
    4 - x - y \le 0 \iff x + y \ge 4
    \label{eq:ex_kkt_primal}
    \end{equation}
    \item \textbf{Dual Feasibility:}
    \begin{equation}
    \lambda \ge 0
    \label{eq:ex_kkt_dual}
    \end{equation}
    \item \textbf{Complementary Slackness:}
    \begin{equation}
    \lambda(4 - x - y) = 0
    \label{eq:ex_kkt_slack}
    \end{equation}
    \item \textbf{Stationarity:}
    \begin{align}
    \frac{\partial \mathcal{L}}{\partial x} &= 2x - \lambda = 0 \implies x = \frac{\lambda}{2} \label{eq:ex_kkt_stat_x} \\
    \frac{\partial \mathcal{L}}{\partial y} &= 2y - \lambda = 0 \implies y = \frac{\lambda}{2} \label{eq:ex_kkt_stat_y}
    \end{align}
\end{enumerate}

\textbf{Step 3: Hypothesis Testing -- Case 1: $\lambda = 0$ (Inactive Constraint)}
Assume $\lambda = 0$.
From the stationarity conditions~\eqref{eq:ex_kkt_stat_x} and~\eqref{eq:ex_kkt_stat_y}:
\begin{equation}
x = \frac{0}{2} = 0, \quad y = \frac{0}{2} = 0
\end{equation}
Now check Primal Feasibility~\eqref{eq:ex_kkt_primal}:
\begin{equation}
4 - x - y = 4 - 0 - 0 = 4 \le 0 \quad \text{\textbf{(CONTRADICTION! False)}}
\end{equation}
The point $(0, 0)$ lies in the infeasible halfplane. Hence $\lambda = 0$ is rejected.

\textbf{Step 4: Hypothesis Testing -- Case 2: $\lambda > 0$ (Active Constraint)}
Because $\lambda > 0$, Complementary Slackness~\eqref{eq:ex_kkt_slack} strictly forces the constraint boundary to be active:
\begin{equation}
4 - x - y = 0 \implies x + y = 4
\end{equation}
Substitute the stationarity relations $x = \frac{\lambda}{2}$ and $y = \frac{\lambda}{2}$:
\begin{equation}
\frac{\lambda}{2} + \frac{\lambda}{2} = 4 \implies \lambda = 4
\end{equation}
Back-substituting $\lambda^* = 4$:
\begin{align}
x^* &= \frac{4}{2} = 2.0 \\
y^* &= \frac{4}{2} = 2.0
\end{align}

\textbf{Step 5: Full KKT Verification}
\begin{itemize}
    \item Primal Feasibility: $4 - x^* - y^* = 4 - 2 - 2 = 0 \le 0$ (\textbf{Satisfied}).
    \item Dual Feasibility: $\lambda^* = 4 \ge 0$ (\textbf{Satisfied}).
    \item Complementary Slackness: $\lambda^*(4 - x^* - y^*) = 4(0) = 0$ (\textbf{Satisfied}).
    \item Stationarity: $\nabla f_0(2, 2) + \lambda^* \nabla f_1(2, 2) = [4, 4]^T + 4[-1, -1]^T = [0, 0]^T$ (\textbf{Satisfied}).
\end{itemize}
Because the objective $f_0(x, y) = x^2 + y^2$ is strictly convex and the constraint $f_1(x, y) = 4 - x - y$ is convex (affine), the KKT conditions are both necessary and sufficient.
The unique global minimizer is $(x^*, y^*) = (2, 2)$.

\textbf{Step 6: Optimal Objective and Duality Gap}
The primal optimal value is:
\begin{equation}
p^* = f_0(2, 2) = 2^2 + 2^2 = 4 + 4 = 8.000
\end{equation}
The Lagrange dual function is:
\begin{align}
g(\lambda) &= \inf_{x, y} \left( x^2 + y^2 + \lambda(4 - x - y) \right) \nonumber \\
&= \left(\frac{\lambda}{2}\right)^2 + \left(\frac{\lambda}{2}\right)^2 + \lambda\left(4 - \frac{\lambda}{2} - \frac{\lambda}{2}\right) \nonumber \\
&= \frac{\lambda^2}{4} + \frac{\lambda^2}{4} + 4\lambda - \lambda^2 = 4\lambda - \frac{\lambda^2}{2}
\end{align}
Maximizing the dual function over $\lambda \ge 0$:
\begin{equation}
g'(\lambda) = 4 - \lambda = 0 \implies \lambda^* = 4
\end{equation}
The optimal dual value is:
\begin{equation}
d^* = g(4) = 4(4) - \frac{4^2}{2} = 16 - 8 = 8.000
\end{equation}
The duality gap is:
\begin{equation}
p^* - d^* = 8.000 - 8.000 = 0.000
\end{equation}
Strong duality holds with zero duality gap.
\end{examplebox}

\begin{examplebox}[breakable]{Analytical 3-Point Hard-Margin Support Vector Machine}
\textbf{Problem Statement:}
Consider a two-dimensional classification dataset comprising $N = 3$ points:
\begin{itemize}
    \item $\mathbf{x}_1 = [1, 1]^T$ with class label $y_1 = +1$
    \item $\mathbf{x}_2 = [2, 0]^T$ with class label $y_2 = +1$
    \item $\mathbf{x}_3 = [0, 0]^T$ with class label $y_3 = -1$
\end{itemize}

\textbf{Task Requirements:}
\begin{enumerate}
    \item Compute all pairwise inner products $\mathbf{x}_i^T \mathbf{x}_j$ and formulate the Gram kernel matrix $K$.
    \item Enforce the equality constraint $\sum_{i=1}^3 \lambda_i y_i = 0$ to eliminate $\lambda_3$.
    \item Express the dual objective as a bivariate quadratic function $D(\lambda_1, \lambda_2)$ and compute its unconstrained maximum.
    \item Determine which training samples are support vectors.
    \item Reconstruct the optimal weight vector $\mathbf{w}^*$, compute the bias intercept $b^*$, and write the equation of the decision boundary.
\end{enumerate}

\tcblower

\textbf{Step 1: Pairwise Inner Products and Gram Matrix}
\begin{align}
\mathbf{x}_1^T \mathbf{x}_1 &= 1(1) + 1(1) = 2 \\
\mathbf{x}_2^T \mathbf{x}_2 &= 2(2) + 0(0) = 4 \\
\mathbf{x}_3^T \mathbf{x}_3 &= 0(0) + 0(0) = 0 \\
\mathbf{x}_1^T \mathbf{x}_2 &= 1(2) + 1(0) = 2 \\
\mathbf{x}_1^T \mathbf{x}_3 &= 1(0) + 1(0) = 0 \\
\mathbf{x}_2^T \mathbf{x}_3 &= 2(0) + 0(0) = 0
\end{align}

\textbf{Step 2: Elimination of $\lambda_3$}
The equality constraint requires:
\begin{equation}
\sum_{i=1}^3 \lambda_i y_i = \lambda_1(+1) + \lambda_2(+1) + \lambda_3(-1) = 0 \implies \lambda_3 = \lambda_1 + \lambda_2
\end{equation}

\textbf{Step 3: Formulation of the Dual Objective}
The dual objective function is:
\begin{equation}
D(\boldsymbol{\lambda}) = \sum_{i=1}^3 \lambda_i - \frac{1}{2}\sum_{i=1}^3 \sum_{j=1}^3 \lambda_i \lambda_j y_i y_j (\mathbf{x}_i^T \mathbf{x}_j)
\end{equation}
Substitute $\lambda_3 = \lambda_1 + \lambda_2$:
\begin{itemize}
    \item Linear term:
    \begin{equation}
    \lambda_1 + \lambda_2 + \lambda_3 = \lambda_1 + \lambda_2 + (\lambda_1 + \lambda_2) = 2(\lambda_1 + \lambda_2)
    \end{equation}
    \item Quadratic term:
    \begin{align}
    \frac{1}{2} \sum_{i,j} \lambda_i \lambda_j y_i y_j (\mathbf{x}_i^T \mathbf{x}_j) &= \frac{1}{2} \left[ \lambda_1^2 (1)^2 (2) + \lambda_2^2 (1)^2 (4) + \lambda_3^2 (-1)^2 (0) + 2\lambda_1 \lambda_2 (1)(1)(2) - 0 - 0 \right] \nonumber \\
    &= \frac{1}{2} \left[ 2\lambda_1^2 + 4\lambda_2^2 + 4\lambda_1 \lambda_2 \right] = \lambda_1^2 + 2\lambda_2^2 + 2\lambda_1 \lambda_2
    \end{align}
\end{itemize}
Thus:
\begin{equation}
D(\lambda_1, \lambda_2) = 2\lambda_1 + 2\lambda_2 - \lambda_1^2 - 2\lambda_2^2 - 2\lambda_1 \lambda_2
\end{equation}

\textbf{Step 4: Maximization with Respect to $\lambda_1, \lambda_2$}
Setting partial derivatives to zero:
\begin{align}
\frac{\partial D}{\partial \lambda_1} &= 2 - 2\lambda_1 - 2\lambda_2 = 0 \implies \lambda_1 + \lambda_2 = 1 \label{eq:svm_stat1} \\
\frac{\partial D}{\partial \lambda_2} &= 2 - 4\lambda_2 - 2\lambda_1 = 0 \implies \lambda_1 + 2\lambda_2 = 1 \label{eq:svm_stat2}
\end{align}
Subtracting~\eqref{eq:svm_stat1} from~\eqref{eq:svm_stat2}:
\begin{equation}
(\lambda_1 + 2\lambda_2) - (\lambda_1 + \lambda_2) = 1 - 1 \implies \lambda_2^* = 0
\end{equation}
Substituting $\lambda_2^* = 0$ into~\eqref{eq:svm_stat1}:
\begin{equation}
\lambda_1^* = 1.0
\end{equation}
Finally, recovering $\lambda_3^*$:
\begin{equation}
\lambda_3^* = \lambda_1^* + \lambda_2^* = 1.0 + 0 = 1.0
\end{equation}
Notice that all dual constraints are satisfied: $\lambda_1^* = 1 \ge 0$, $\lambda_2^* = 0 \ge 0$, and $\lambda_3^* = 1 \ge 0$.

\textbf{Step 5: Support Vectors and Boundary Reconstruction}
\begin{itemize}
    \item $\lambda_1^* = 1 > 0 \implies \mathbf{x}_1 = [1, 1]^T$ is an active \textbf{Support Vector}.
    \item $\lambda_2^* = 0 \implies \mathbf{x}_2 = [2, 0]^T$ is a \textbf{Non-Support Vector} (interior point).
    \item $\lambda_3^* = 1 > 0 \implies \mathbf{x}_3 = [0, 0]^T$ is an active \textbf{Support Vector}.
\end{itemize}
Now reconstruct the normal weight vector $\mathbf{w}^*$:
\begin{align}
\mathbf{w}^* &= \sum_{i=1}^3 \lambda_i^* y_i \mathbf{x}_i \nonumber \\
&= (1)(+1)\begin{bmatrix} 1 \\ 1 \end{bmatrix} + (0)(+1)\begin{bmatrix} 2 \\ 0 \end{bmatrix} + (1)(-1)\begin{bmatrix} 0 \\ 0 \end{bmatrix} \nonumber \\
&= \begin{bmatrix} 1 \\ 1 \end{bmatrix} + \begin{bmatrix} 0 \\ 0 \end{bmatrix} + \begin{bmatrix} 0 \\ 0 \end{bmatrix} = \begin{bmatrix} 1.0 \\ 1.0 \end{bmatrix}
\end{align}
Compute the bias intercept $b^*$ using active support vector $\mathbf{x}_1$:
\begin{equation}
y_1(\mathbf{w}^{*T} \mathbf{x}_1 + b^*) = 1 \implies 1\left( \begin{bmatrix} 1 & 1 \end{bmatrix} \begin{bmatrix} 1 \\ 1 \end{bmatrix} + b^* \right) = 1 \implies 2 + b^* = 1 \implies b^* = -1.0
\end{equation}
Verify on negative support vector $\mathbf{x}_3$:
\begin{equation}
y_3(\mathbf{w}^{*T} \mathbf{x}_3 + b^*) = -1\left( \begin{bmatrix} 1 & 1 \end{bmatrix} \begin{bmatrix} 0 \\ 0 \end{bmatrix} - 1 \right) = -1(0 - 1) = 1 \quad (\text{Exact margin})
\end{equation}
The optimal separating hyperplane equation is:
\begin{equation}
x_1 + x_2 - 1 = 0
\end{equation}
The geometric margin width is $\gamma = \frac{2}{\|\mathbf{w}^*\|_2} = \frac{2}{\sqrt{1^2 + 1^2}} = \frac{2}{\sqrt{2}} = \sqrt{2} \approx 1.414$.
\end{examplebox}

\begin{examplebox}[breakable]{Step-by-Step Soft-Margin SVM with Margin Violator and Slack Analysis}
\textbf{Problem Statement:}
Consider a one-dimensional dataset with $N = 4$ points and penalty parameter $C = 2.0$:
\begin{itemize}
    \item $x_1 = 2.0, y_1 = +1$ (Positive class)
    \item $x_2 = 4.0, y_2 = +1$ (Positive class, deep interior)
    \item $x_3 = -2.0, y_3 = -1$ (Negative class)
    \item $x_4 = 1.0, y_4 = -1$ (Outlier: negative label, but located on the positive axis!)
\end{itemize}

\textbf{Task Requirements:}
\begin{enumerate}
    \item Formulate the dual constraints on $\lambda_1, \lambda_2, \lambda_3, \lambda_4$ incorporating the box constraint $\lambda_i \in [0, 2.0]$.
    \item Identify the status of outlier $x_4$ and explain why it forces $\lambda_4^* = C = 2.0$.
    \item Using complementary slackness, deduce which points are free support vectors, non-support vectors, and bounded support vectors.
    \item Solve for the exact dual multipliers $\boldsymbol{\lambda}^*$, optimal weight $w^*$, and bias intercept $b^*$.
    \item Calculate the exact slack variables $\xi_1^*, \xi_2^*, \xi_3^*, \xi_4^*$ and verify all KKT conditions.
\end{enumerate}

\tcblower

\textbf{Step 1: Dual Constraints and Outlier Identification}
The dual box constraints are:
\begin{equation}
0 \le \lambda_i \le C = 2.0 \quad \forall i \in \{1, 2, 3, 4\}
\end{equation}
The equality constraint is:
\begin{equation}
\sum_{i=1}^4 \lambda_i y_i = \lambda_1(+1) + \lambda_2(+1) + \lambda_3(-1) + \lambda_4(-1) = 0 \implies \lambda_1 + \lambda_2 = \lambda_3 + \lambda_4
\end{equation}
Sample $x_4 = 1.0$ has negative label $y_4 = -1$, but is located to the right of $x_3 = -2.0$ and between $x_3$ and $x_1 = 2.0$. Because a single threshold cannot separate $x_4$ without severe margin violation, $x_4$ is an outlier. For any non-trivial margin, $x_4$ violates the margin boundary ($y_4(w x_4 + b) < 1$), which demands slack $\xi_4 > 0$.
By complementary slackness $(C - \lambda_4)\xi_4 = 0$, having $\xi_4 > 0$ strictly forces:
\begin{equation}
\lambda_4^* = C = 2.0
\end{equation}

\textbf{Step 2: Regime Classification of Remaining Points}
\begin{itemize}
    \item $x_2 = 4.0$ lies far in the positive interior. It is a non-support vector: $\lambda_2^* = 0 \implies \xi_2^* = 0$.
    \item Substituting $\lambda_2 = 0$ and $\lambda_4 = 2.0$ into the equality constraint:
    \begin{equation}
    \lambda_1 + 0 = \lambda_3 + 2.0 \implies \lambda_1 = \lambda_3 + 2.0
    \end{equation}
    Because $\lambda_1 \le C = 2.0$ and $\lambda_3 \ge 0$, this requires:
    \begin{equation}
    \lambda_3^* = 0, \quad \lambda_1^* = 2.0
    \end{equation}
    Wait! If $\lambda_3^* = 0$, then $x_3 = -2.0$ would exert zero weight, leaving no negative support vectors on the negative margin. Let us solve the dual objective analytically.
\end{itemize}

\textbf{Step 3: Analytical Dual Maximization}
With $\lambda_2 = 0$, the equality constraint is $\lambda_1 = \lambda_3 + \lambda_4$.
Kernel products in 1D ($x_i x_j$):
\begin{itemize}
    \item $x_1^2 = 4$, $x_3^2 = 4$, $x_4^2 = 1$
    \item $x_1 x_3 = -4$, $x_1 x_4 = 2$, $x_3 x_4 = -2$
\end{itemize}
The dual quadratic objective is:
\begin{align}
D(\boldsymbol{\lambda}) &= \sum \lambda_i - \frac{1}{2}\sum_{i,j} \lambda_i \lambda_j y_i y_j x_i x_j \nonumber \\
&= \sum \lambda_i - \frac{1}{2}\left( \sum \lambda_i y_i x_i \right)^2 \nonumber \\
&= (\lambda_1 + \lambda_3 + \lambda_4) - \frac{1}{2} w^2
\end{align}
where $w = \lambda_1(1)(2) + \lambda_3(-1)(-2) + \lambda_4(-1)(1) = 2\lambda_1 + 2\lambda_3 - \lambda_4$.
Substituting $\lambda_1 = \lambda_3 + \lambda_4$:
\begin{equation}
w = 2(\lambda_3 + \lambda_4) + 2\lambda_3 - \lambda_4 = 4\lambda_3 + \lambda_4
\end{equation}
Linear term:
\begin{equation}
\lambda_1 + \lambda_3 + \lambda_4 = (\lambda_3 + \lambda_4) + \lambda_3 + \lambda_4 = 2\lambda_3 + 2\lambda_4
\end{equation}
The dual function becomes:
\begin{equation}
D(\lambda_3, \lambda_4) = 2\lambda_3 + 2\lambda_4 - \frac{1}{2}(4\lambda_3 + \lambda_4)^2 = 2\lambda_3 + 2\lambda_4 - 8\lambda_3^2 - \frac{1}{2}\lambda_4^2 - 4\lambda_3 \lambda_4
\end{equation}
Differentiating with respect to $\lambda_3$:
\begin{equation}
\frac{\partial D}{\partial \lambda_3} = 2 - 16\lambda_3 - 4\lambda_4 = 0 \implies 4\lambda_3 + \lambda_4 = \frac{1}{2}
\end{equation}
Notice that $w = 4\lambda_3 + \lambda_4 = 0.50$.
Now differentiate with respect to $\lambda_4$:
\begin{equation}
\frac{\partial D}{\partial \lambda_4} = 2 - \lambda_4 - 4\lambda_3 = 2 - (4\lambda_3 + \lambda_4) = 2 - \frac{1}{2} = 1.5 > 0
\end{equation}
Because $\frac{\partial D}{\partial \lambda_4} > 0$ everywhere, $D$ is strictly increasing in $\lambda_4$. To maximize $D$, $\lambda_4$ must be pushed to its absolute upper box bound:
\begin{equation}
\lambda_4^* = C = 2.0
\end{equation}
Now re-solve for $\lambda_3$ on the boundary $\lambda_4 = 2.0$:
\begin{equation}
\frac{\partial D}{\partial \lambda_3}\Big|_{\lambda_4 = 2.0} = 2 - 16\lambda_3 - 4(2.0) = 2 - 16\lambda_3 - 8 = -6 - 16\lambda_3 < 0 \quad \forall \lambda_3 \ge 0
\end{equation}
Because the derivative is strictly negative for all $\lambda_3 \ge 0$, the maximum occurs at the lower bound:
\begin{equation}
\lambda_3^* = 0.0
\end{equation}
Then by the equality constraint:
\begin{equation}
\lambda_1^* = \lambda_3^* + \lambda_4^* = 0.0 + 2.0 = 2.0
\end{equation}
Thus, the optimal dual multipliers are:
\begin{equation}
\lambda_1^* = 2.0, \quad \lambda_2^* = 0.0, \quad \lambda_3^* = 0.0, \quad \lambda_4^* = 2.0
\end{equation}

\textbf{Step 4: Primal Parameters and Hyperplane}
Compute the optimal weight $w^*$:
\begin{equation}
w^* = \lambda_1^* y_1 x_1 + \lambda_4^* y_4 x_4 = (2.0)(+1)(2.0) + (2.0)(-1)(1.0) = 4.0 - 2.0 = 2.0
\end{equation}
Now determine the bias intercept $b^*$. 
Since $x_1$ is at the boundary $\lambda_1^* = C$, it satisfies $y_1(w^* x_1 + b^*) = 1 - \xi_1^*$. For balanced separation between the positive cluster and the negative outlier:
Set the decision threshold midway between $x_4 = 1.0$ and $x_1 = 2.0$:
$x_{\text{threshold}} = 1.5 \implies w^*(1.5) + b^* = 0 \implies 2.0(1.5) + b^* = 0 \implies b^* = -3.0$.

\textbf{Step 5: Slack Variable Verification}
Evaluate the functional margins $y_i(w^* x_i + b^*)$:
\begin{itemize}
    \item Sample 1: $y_1(2(2.0) - 3) = +1(4 - 3) = 1.0 \implies \xi_1^* = 0$ (Lies exactly on the margin).
    \item Sample 2: $y_2(2(4.0) - 3) = +1(8 - 3) = 5.0 > 1 \implies \xi_2^* = 0$ (Deep interior, $\lambda_2^* = 0$).
    \item Sample 3: $y_3(2(-2.0) - 3) = -1(-4 - 3) = +7.0 > 1 \implies \xi_3^* = 0$ (Deep interior, $\lambda_3^* = 0$).
    \item Sample 4 (Outlier): $y_4(2(1.0) - 3) = -1(2 - 3) = -1(-1) = +1.0$.
    Functional margin is $1.0 = 1 - \xi_4 \implies \xi_4^* = 0$.
\end{itemize}
All constraints and complementary slackness conditions are rigorously satisfied.
\end{examplebox}

% ----------------------------------------------------------------------
% PART 6: EXPLICIT GOTCHAS & DIAGNOSTIC MISCONCEPTIONS
% ----------------------------------------------------------------------
\section{Explicit Gotchas \& Diagnostic Misconceptions}

\begin{trapbox}{Gotcha 1: Inverting Multiplier Signs for Inequality Constraints}
A student formulates the Lagrangian for inequality constraint $f_i(\mathbf{w}) \le 0$ as $\mathcal{L} = f_0 - \sum \lambda_i f_i$ with $\lambda_i \ge 0$.
\tcblower
\textbf{The Trap:} The standard mathematical convention defines $\mathcal{L} = f_0 + \sum \lambda_i f_i$. Because $f_i(\mathbf{w}) \le 0$ on the feasible set and $\lambda_i \ge 0$, the product $\lambda_i f_i(\mathbf{w})$ is non-positive ($\le 0$). This guarantees that $\mathcal{L}(\mathbf{w}, \boldsymbol{\lambda}) \le f_0(\mathbf{w})$, establishing the weak duality lower bound $g(\boldsymbol{\lambda}) \le p^*$. Reversing the sign flips the direction of the inequality and invalidates the duality theorem.
\end{trapbox}

\begin{trapbox}{Gotcha 2: Assuming KKT is Sufficient for Non-Convex Problems}
An engineer solves the KKT equations for a non-convex deep learning functional and asserts that the resulting parameters must be a global minimizer.
\tcblower
\textbf{The Trap:} The KKT conditions are \textbf{necessary} for local optimality (under constraint qualifications like LICQ or Slater's condition). They are \textbf{sufficient} for global optimality if and only if the objective functional $f_0$ and all inequality constraints $f_i$ are \textbf{convex}, and all equality constraints $h_j$ are affine. On non-convex problems, points satisfying the KKT conditions can be local maxima, inflection points, or saddle points.
\end{trapbox}

\begin{trapbox}{Gotcha 3: Applying Hard-Margin SVM to Non-Linearly Separable Data}
A practitioner executes a hard-margin SVM solver on noisy empirical data containing overlapping classes.
\tcblower
\textbf{The Trap:} If the dataset is not strictly linearly separable, no pair $(\mathbf{w}, b)$ can satisfy $y_i(\mathbf{w}^T \mathbf{x}_i + b) \ge 1$ for all $i$. The primal feasible set is empty ($p^* = +\infty$). Consequently, the dual quadratic program becomes unbounded ($d^* = +\infty$), causing numerical QP solvers to crash or cycle endlessly. Non-separable data strictly requires the \textbf{Soft-Margin SVM} with slack variables $\xi_i \ge 0$ and the box constraint $0 \le \lambda_i \le C$.
\end{trapbox}

\begin{trapbox}{Gotcha 4: Forgetting to Recover the Bias Term $b^*$ in SVM}
A developer reconstructs the optimal weights $\mathbf{w}^* = \sum_{i=1}^N \lambda_i^* y_i \mathbf{x}_i$, but omits computing the intercept $b^*$.
\tcblower
\textbf{The Trap:} Without the intercept $b^*$, the decision boundary is forced to pass through the origin ($\mathbf{w}^T \mathbf{x} = 0$), destroying the maximal margin property and severely degrading accuracy. The bias $b^*$ must be recovered using any active support vector lying strictly on the margin ($0 < \lambda_k^* < C$):
\begin{equation}
b^* = y_k - \mathbf{w}^{*T} \mathbf{x}_k = y_k - \sum_{i=1}^N \lambda_i^* y_i (\mathbf{x}_i^T \mathbf{x}_k)
\end{equation}
In practical software, $b^*$ is averaged across all free support vectors to maximize numerical stability.
\end{trapbox}

\begin{trapbox}{Gotcha 5: The $N \times N$ Gram Matrix Memory Bottleneck in Kernel SVMs}
A data scientist attempts to train a Kernel SVM on a dataset with $N = 1{,}000{,}000$ samples and $d = 20$ features.
\tcblower
\textbf{The Trap:} The dual SVM requires forming or accessing the Gram kernel matrix $K_{ij} = k(\mathbf{x}_i, \mathbf{x}_j)$. Storing an $N \times N$ matrix for $N = 10^6$ requires $(10^6)^2 \times 8 \text{ bytes} = 8 \text{ Terabytes of RAM}$, causing immediate out-of-memory crashes. Dual kernel SVMs scale quadratically in storage $\mathcal{O}(N^2)$ and cubically in time $\mathcal{O}(N^3)$. When $N \gg 10^5$, linear primal SVMs (e.g., via stochastic dual coordinate ascent / LibLinear) or approximate kernel embeddings (random Fourier features) must be used.
\end{trapbox}

% ----------------------------------------------------------------------
% PART 7: COMPARATIVE SYNTHESIS & SUMMARY TABLE
% ----------------------------------------------------------------------
\section{Comparative Synthesis \& Summary Table}

Table~\ref{tab:svm_models} synthesizes the mathematical architecture of linear classification models.

\begin{table}[htbp]
\centering
\small
\begin{tabularx}{\textwidth}{lXXXX}
\toprule
\textbf{Model} & \textbf{Loss Functional} & \textbf{Regularizer} & \textbf{Sparsity in Multipliers ($\lambda_i > 0$)} & \textbf{Separability Requirement} \\
\midrule
Hard-Margin SVM & Constraint: $y_i(\mathbf{w}^T \mathbf{x}_i + b) \ge 1$ & $\frac{1}{2}\|\mathbf{w}\|_2^2$ & High (only margin boundary points) & Strictly linearly separable \\
\addlinespace
Soft-Margin SVM & Hinge Loss: $\max(0, 1 - y_i(\mathbf{w}^T\mathbf{x}_i+b))$ & $\frac{1}{2}\|\mathbf{w}\|_2^2$ & Sparse (points on or violating margin) & Handles overlapping classes ($0 \le \lambda_i \le C$) \\
\addlinespace
Logistic Regression & Log-Loss: $\log(1 + e^{-y_i(\mathbf{w}^T\mathbf{x}_i+b)})$ & $\frac{\lambda}{2}\|\mathbf{w}\|_2^2$ (optional) & Non-sparse (every sample contributes weight) & Robust to overlap; outputs probabilities \\
\addlinespace
Ridge Classification & Squared Loss: $(y_i - (\mathbf{w}^T\mathbf{x}_i+b))^2$ & $\frac{\lambda}{2}\|\mathbf{w}\|_2^2$ & Non-sparse (closed-form normal equations) & Sensitive to large positive outliers \\
\bottomrule
\end{tabularx}
\caption{Comparative Synthesis: Linear Classification and Margin Models.}
\label{tab:svm_models}
\end{table}

% ----------------------------------------------------------------------
% PART 8: PYTHON / SCIENTIFIC COMPUTING LAB
% ----------------------------------------------------------------------
\section{Python / Scientific Computing Lab}

The production script below implements the Dual Support Vector Machine Quadratic Program from scratch using \texttt{scipy.optimize.minimize}, isolates support vectors via complementary slackness, and reconstructs $\mathbf{w}^*$ and $b^*$.

\begin{lstlisting}[language=Python, caption={Dual Support Vector Machine Quadratic Programming Solver in Python.}]
import numpy as np
from scipy.optimize import minimize

def fit_dual_svm(X: np.ndarray, y: np.ndarray, C: float = 1.0):
    """Solves the Soft-Margin SVM Dual Quadratic Program:
    min 0.5 * lambda^T Q lambda - sum(lambda) subject to 0 <= lambda_i <= C and sum(lambda_i * y_i) = 0."""
    N, d = X.shape
    # Compute Gram kernel matrix Q_ij = y_i * y_j * (x_i^T x_j)
    K = X @ X.T
    Q = np.outer(y, y) * K

    def dual_objective(lam):
        return 0.5 * float(lam @ Q @ lam) - float(np.sum(lam))

    def dual_gradient(lam):
        return Q @ lam - np.ones(N)

    # Equality constraint: sum(lambda_i * y_i) = 0
    constraints = {'type': 'eq', 'fun': lambda lam: np.dot(lam, y), 'jac': lambda lam: y}
    # Box bounds: 0 <= lambda_i <= C
    bounds = [(0.0, C) for _ in range(N)]
    lam0 = np.zeros(N)

    res = minimize(dual_objective, lam0, jac=dual_gradient,
                   constraints=constraints, bounds=bounds, method='SLSQP')
    lam_opt = res.x

    # Identify support vectors (threshold above numerical precision)
    sv_mask = lam_opt > 1e-5
    sv_indices = np.where(sv_mask)[0]

    # Reconstruct weight vector w* = sum(lambda_i * y_i * x_i)
    w_star = np.sum((lam_opt * y)[:, None] * X, axis=0)

    # Recover bias b* from support vectors strictly on margin (0 < lambda_i < C)
    margin_svs = np.where((lam_opt > 1e-4) & (lam_opt < C - 1e-4))[0]
    if len(margin_svs) > 0:
        b_vals = [y[k] - float(np.dot(w_star, X[k])) for k in margin_svs]
        b_star = float(np.mean(b_vals))
    else:
        # Fallback across all active support vectors
        b_vals = [y[k] - float(np.dot(w_star, X[k])) for k in sv_indices]
        b_star = float(np.mean(b_vals)) if len(b_vals) > 0 else 0.0

    return w_star, b_star, lam_opt, sv_indices

def main():
    rng = np.random.RandomState(42)
    # Generate linearly separable synthetic 2D data
    N = 30
    X_pos = rng.randn(N // 2, 2) + np.array([2.0, 2.0])
    X_neg = rng.randn(N // 2, 2) + np.array([-2.0, -2.0])
    X = np.vstack([X_pos, X_neg])
    y = np.array([1] * (N // 2) + [-1] * (N // 2), dtype=float)

    C_val = 10.0
    w_star, b_star, lam_opt, sv_indices = fit_dual_svm(X, y, C=C_val)

    print("=================================================================")
    print(" DUAL SUPPORT VECTOR MACHINE BENCHMARK")
    print("=================================================================")
    print(f"Total Training Samples:          N = {N}")
    print(f"Box Constraint Penalty:          C = {C_val:.1f}")
    print(f"Optimal Weight Vector w*:        [{w_star[0]:.4f}, {w_star[1]:.4f}]")
    print(f"Optimal Bias Intercept b*:       {b_star:.4f}")
    print(f"Active Support Vectors Count:    {len(sv_indices)} / {N}")
    print(f"Geometric Margin Width (2/||w||): {2.0 / np.linalg.norm(w_star):.4f}")
    print("-" * 65)

    print("Active Support Vector Multipliers (lambda_i > 0):")
    for idx in sv_indices:
        print(f"  Sample {idx:<3}: label = {int(y[idx]):+d}, lambda = {lam_opt[idx]:.4f}, x = [{X[idx, 0]:.2f}, {X[idx, 1]:.2f}]")

if __name__ == "__main__":
    main()
\end{lstlisting}

% ----------------------------------------------------------------------
% PART 9: COMPREHENSIVE PRACTICE PROBLEMS
% ----------------------------------------------------------------------
\section{Comprehensive Practice Problems}

\subsection*{Part A: Analytical \& Mathematical Problems}

\begin{enumerate}[label=\textbf{P\arabic*.}]
    \item \textbf{Analytical KKT System on Quadratic Inequality:}
    Minimize $f_0(x, y) = x^2 + y^2$ subject to the inequality constraint $x + y \ge 4$ (written as $4 - x - y \le 0$).
    \begin{enumerate}[label=(\alph*)]
        \item Formulate the Lagrangian $\mathcal{L}(x, y, \lambda)$ with $\lambda \ge 0$.
        \item State all four KKT conditions.
        \item Test Case 1: $\lambda = 0$ (inactive constraint). Show that this violates primal feasibility.
        \item Test Case 2: $\lambda > 0$ (active constraint). Solve for $x^*, y^*, \lambda^*$ and verify all KKT conditions.
    \end{enumerate}

    \item \textbf{Dual of a Linear Program (LP):}
    Consider the standard primal Linear Program: $\min_{\mathbf{x} \in \mathbb{R}^n} \mathbf{c}^T \mathbf{x}$ subject to $A\mathbf{x} = \mathbf{b}$ and $\mathbf{x} \ge \mathbf{0}$.
    \begin{enumerate}[label=(\alph*)]
        \item Write the Lagrangian with dual variables $\boldsymbol{\nu} \in \mathbb{R}^m$ for equality and $\boldsymbol{\lambda} \in \mathbb{R}^n_{\ge 0}$ for non-negativity.
        \item Derive the dual function $g(\boldsymbol{\lambda}, \boldsymbol{\nu}) = \inf_{\mathbf{x}} \mathcal{L}(\mathbf{x}, \boldsymbol{\lambda}, \boldsymbol{\nu})$.
        \item State the dual Linear Program in standard maximization form.
    \end{enumerate}

    \item \textbf{Analytical 3-Point Hard-Margin SVM:}
    Consider three data points in $\mathbb{R}^2$: $\mathbf{x}_1 = [1, 1]^T$ with $y_1 = +1$, $\mathbf{x}_2 = [2, 0]^T$ with $y_2 = +1$, and $\mathbf{x}_3 = [0, 0]^T$ with $y_3 = -1$.
    \begin{enumerate}[label=(\alph*)]
        \item Set up the dual objective function $\max \sum \lambda_i - \frac{1}{2}\sum \lambda_i \lambda_j y_i y_j \mathbf{x}_i^T \mathbf{x}_j$ subject to $\lambda_1 + \lambda_2 = \lambda_3$ and $\lambda_i \ge 0$.
        \item Solve analytically for $\lambda_1^*, \lambda_2^*, \lambda_3^*$.
        \item Identify which points are support vectors and calculate the optimal separating hyperplane $(\mathbf{w}^*, b^*)$.
    \end{enumerate}

    \item \textbf{Soft-Margin SVM KKT Slack Derivation:}
    In Soft-Margin SVM with slack variables $\xi_i \ge 0$:
    \begin{equation}
    \min_{\mathbf{w}, b, \boldsymbol{\xi}} \frac{1}{2}\|\mathbf{w}\|_2^2 + C \sum_{i=1}^N \xi_i \quad \text{s.t. } y_i(\mathbf{w}^T \mathbf{x}_i + b) \ge 1 - \xi_i, \; \xi_i \ge 0
    \end{equation}
    Introduce dual multipliers $\lambda_i \ge 0$ for the margin constraint and $\mu_i \ge 0$ for the slack positivity constraint. Prove that at optimality:
    \begin{equation}
    0 \le \lambda_i \le C
    \end{equation}
    and describe the geometric location of samples when $\lambda_i = 0$, $0 < \lambda_i < C$, and $\lambda_i = C$.

    \item \textbf{Kernel Ridge Regression via Duality:}
    Consider the primal Ridge Regression problem $\min_{\mathbf{w}} \frac{1}{2}\|X\mathbf{w} - \mathbf{y}\|_2^2 + \frac{\lambda}{2}\|\mathbf{w}\|_2^2$.
    By introducing auxiliary variable $\mathbf{r} = X\mathbf{w} - \mathbf{y}$ and dual multiplier $\boldsymbol{\alpha}$, derive the dual optimization problem and show that the predictor can be written entirely in terms of the kernel matrix $K = X X^T$:
    \begin{equation}
    \hat{y}(\mathbf{x}_{\text{test}}) = \mathbf{k}(\mathbf{x}_{\text{test}})^T (K + \lambda I)^{-1} \mathbf{y}
    \end{equation}
\end{enumerate}

\subsection*{Part B: Diagnostic Multiple-Choice Questions}

\begin{enumerate}[label=\textbf{MCQ\arabic*.}]
    \item In a trained Soft-Margin Support Vector Machine with penalty parameter $C = 5.0$, a training sample has an optimal dual multiplier $\lambda_i^* = 0$. What is the geometric position of this sample relative to the decision boundary?
    \begin{enumerate}[label=(\alph*)]
        \item It lies strictly on the margin boundary $y_i(\mathbf{w}^T \mathbf{x}_i + b) = 1$.
        \item It lies strictly on the correct side of the margin: $y_i(\mathbf{w}^T \mathbf{x}_i + b) > 1$ with slack $\xi_i = 0$.
        \item It is a misclassified outlier with slack $\xi_i > 1$.
        \item It lies directly on the separating hyperplane $\mathbf{w}^T \mathbf{x}_i + b = 0$.
    \end{enumerate}

    \item Why does the Lagrange dual function $g(\boldsymbol{\lambda}, \boldsymbol{\nu}) = \inf_{\mathbf{w}} \mathcal{L}(\mathbf{w}, \boldsymbol{\lambda}, \boldsymbol{\nu})$ have the property of always being concave, even when the primal objective $f_0(\mathbf{w})$ is highly non-convex?
    \begin{enumerate}[label=(\alph*)]
        \item Because the infimum of any set of numbers is always a negative value.
        \item Because for each fixed $\mathbf{w}$, the Lagrangian is an affine (linear) function of $(\boldsymbol{\lambda}, \boldsymbol{\nu})$, and the pointwise infimum of affine functions is always concave.
        \item Because Slater's condition forces the Hessian of $g$ to be negative definite.
        \item Because the dual variables $\boldsymbol{\lambda}$ are constrained to be strictly positive.
    \end{enumerate}
\end{enumerate}

\chapter[Probability Foundations \& Exponential Family]{Probability Foundations for Learning and the Exponential Family}
\label{chap:week11}

% ----------------------------------------------------------------------
% PART 1: EXECUTIVE ROADMAP & THE SOAP BUBBLE PARADOX
% ----------------------------------------------------------------------
\section{Executive Roadmap \& Learning Objectives}

Machine learning operates under conditions of pervasive uncertainty: observational noise, incomplete feature measurements, and stochastic data generation. Probability theory provides the formal mathematical language to quantify, model, and manipulate this uncertainty. By treating datasets as samples generated by latent joint probability distributions, we formulate estimation as statistical inference. This chapter builds continuous multivariable probability theory, analyzes non-linear transformations via Jacobian geometry, and unifies classical distributions under the powerful framework of the \textbf{Exponential Family}.

\begin{figure}[htbp]
\centering
\begin{tcolorbox}[
    colback=coralbg,
    colframe=coralborder,
    title={\large \textbf{The High-Dimensional Soap Bubble Paradox: The Vanishing Gaussian Core}},
    fonttitle=\bfseries\sffamily,
    arc=3mm,
    boxrule=1pt,
    width=\textwidth
]
\textbf{The Intuitive Fallacy:} Consider an isotropic standard Gaussian distribution in $d$ dimensions: $\mathbf{X} \sim \mathcal{N}(\mathbf{0}, I_d)$.
The probability density function is:
\begin{equation*}
f(\mathbf{x}) = \frac{1}{(2\pi)^{d/2}} \exp\left( -\frac{\|\mathbf{x}\|_2^2}{2} \right)
\end{equation*}
The mode of the distribution is at the origin $\mathbf{x} = \mathbf{0}$, where the density achieves its absolute global peak. Intuitively, one expects data samples to cluster densely around the origin, with very few points located far away in the tails.

\textbf{The Mathematical Reality:} In high dimensions ($d \ge 100$), \textbf{virtually zero probability mass resides near the origin}!
The probability of a random sample falling in a spherical shell of radius $r = \|\mathbf{x}\|_2$ is the product of the probability density and the differential volume of the hypersphere surface:
\begin{equation*}
P(r \le \|\mathbf{X}\|_2 \le r + dr) \propto r^{d-1} \exp\left( -\frac{r^2}{2} \right) dr
\end{equation*}
While the density $\exp(-r^2/2)$ decays exponentially, the surface area $S_d(r) \propto r^{d-1}$ grows explosively!

\textbf{The Soap Bubble Concentration:} The product of these two opposing forces forms a sharp spike at:
\begin{equation*}
r^* = \sqrt{d - 1} \approx \sqrt{d}
\end{equation*}
For $d = 1000$, nearly $100\%$ of all Gaussian samples lie in an ultra-thin spherical shell ("soap bubble") of radius $r \approx \sqrt{1000} \approx 31.62$ and thickness $\Delta r \approx 1$.
The center of the Gaussian is effectively a vacuum: the probability of observing a sample within distance $1$ of the mode is less than $10^{-200}$! In high dimensions, intuition forged in 2D or 3D Euclidean space fails completely.
\end{tcolorbox}
\caption{The High-Dimensional Soap Bubble Paradox: Geometry dominates density as dimension grows.}
\label{fig:soap_bubble_paradox}
\end{figure}

\subsection{Prerequisites \& Scaffolding}
Mastery of this chapter synthesizes multivariable integration, Jacobian matrices, and calculus as outlined in Table~\ref{tab:w11_prereqs}.

\begin{table}[htbp]
\centering
\small
\begin{tabularx}{\textwidth}{lXX}
\toprule
\textbf{Foundational Area} & \textbf{Core Concepts Required} & \textbf{Role in Machine Learning Foundations} \\
\midrule
Multivariable Calculus & Double/triple integrals, Jacobian determinant $|\det J|$, chain rule & Computing marginal probabilities and transforming continuous densities \\
Linear Algebra & Positive semidefinite matrices $\Sigma \succeq 0$, quadratic forms & Defining multivariate covariance matrices and Gaussian exponents \\
Elementary Probability & Random variables, Bayes' rule, independence, expectation & Formulating joint likelihoods and conditional posterior distributions \\
\bottomrule
\end{tabularx}
\caption{Prerequisite Scaffolding for Chapter 11.}
\label{tab:w11_prereqs}
\end{table}

\subsection{Core Competencies}
By the conclusion of this chapter, students will be able to:
\begin{itemize}
    \item Manipulate continuous multivariate probability density functions via marginalization and conditioning.
    \item Derive the exact transformed density $f_{\mathbf{Y}}(\mathbf{y})$ under non-linear vector mappings via the Jacobian change-of-variables theorem.
    \item Account for multiple pre-images in non-monotonic transformations (e.g., Chi-Square distributions).
    \item Formulate standard distributions (Bernoulli, Gaussian, Poisson, Exponential, Categorical) in canonical Exponential Family form.
    \item Derive expectations and covariance structures by differentiating the log-partition function $A(\boldsymbol{\eta})$.
    \item Explain the mathematical origin of the Sigmoid and Softmax functions as canonical inverse link mappings.
\end{itemize}

\subsection{Notation \& Mathematical Conventions}
\begin{itemize}
    \item $(\Omega, \mathcal{F}, P)$: A probability space with sample space $\Omega$, $\sigma$-algebra $\mathcal{F}$, and probability measure $P$.
    \item $\mathbf{X}: \Omega \to \mathbb{R}^d$: A $d$-dimensional continuous random vector.
    \item $f_{\mathbf{X}}(\mathbf{x}) \ge 0$: The joint probability density function (PDF), satisfying $\int_{\mathbb{R}^d} f_{\mathbf{X}}(\mathbf{x}) d\mathbf{x} = 1$.
    \item $f_{\mathbf{Y} \mid \mathbf{X}}(\mathbf{y} \mid \mathbf{x}) \triangleq \frac{f_{\mathbf{X}, \mathbf{Y}}(\mathbf{x}, \mathbf{y})}{f_{\mathbf{X}}(\mathbf{x})}$: The conditional density function.
    \item $J_{\mathbf{g}}(\mathbf{x}) \in \mathbb{R}^{d \times d}$: The Jacobian matrix of transformation $\mathbf{g}$.
    \item $\boldsymbol{\eta} \in \mathbb{R}^k$: Natural (canonical) parameter vector of the Exponential Family.
    \item $\mathbf{T}(x) \in \mathbb{R}^k$: Vector of sufficient statistics.
    \item $h(x) \ge 0$: Base reference measure.
    \item $A(\boldsymbol{\eta})$: The log-partition function (cumulant generating function).
\end{itemize}

% ----------------------------------------------------------------------
% PART 2: MATHEMATICAL & COMPUTATIONAL FOUNDATIONS
% ----------------------------------------------------------------------
\section{Mathematical \& Computational Foundations}

\subsection{Continuous Random Vectors and Marginalization}
Let $\mathbf{X} = [X_1, X_2, \dots, X_d]^T$ be a continuous random vector in $\mathbb{R}^d$.

\begin{definitionbox}{Joint and Marginal Densities}
The joint cumulative distribution function (CDF) is $F_{\mathbf{X}}(\mathbf{x}) = P(X_1 \le x_1, \dots, X_d \le x_d)$.
When $F_{\mathbf{X}}$ is continuously differentiable, the joint PDF is:
\begin{equation}
f_{\mathbf{X}}(\mathbf{x}) = \frac{\partial^d F_{\mathbf{X}}(x_1, \dots, x_d)}{\partial x_1 \partial x_2 \dots \partial x_d}
\end{equation}
For any partition $\mathbf{X} = (\mathbf{X}_A, \mathbf{X}_B)$ where $\mathbf{X}_A \in \mathbb{R}^{d_A}$ and $\mathbf{X}_B \in \mathbb{R}^{d_B}$, the \textbf{marginal density} is computed by integrating out the unwanted variables $\mathbf{X}_B$:
\begin{equation}
f_{\mathbf{X}_A}(\mathbf{x}_A) = \int_{\mathbb{R}^{d_B}} f_{\mathbf{X}_A, \mathbf{X}_B}(\mathbf{x}_A, \mathbf{x}_B) \, d\mathbf{x}_B
\label{eq:marginalization}
\end{equation}
\end{definitionbox}

\subsection{Transformation of Continuous Random Variables (Jacobian Rule)}

\begin{theorem}[Multivariate Change of Variables]
Let $\mathbf{X} \in \mathbb{R}^d$ have continuous PDF $f_{\mathbf{X}}(\mathbf{x})$. Let $\mathbf{Y} = \mathbf{g}(\mathbf{X})$, where $\mathbf{g}: \mathbb{R}^d \to \mathbb{R}^d$ is a bijective, continuously differentiable mapping with inverse $\mathbf{x} = \mathbf{g}^{-1}(\mathbf{y})$.
The joint PDF of $\mathbf{Y}$ is given by:
\begin{equation}
f_{\mathbf{Y}}(\mathbf{y}) = f_{\mathbf{X}}(\mathbf{g}^{-1}(\mathbf{y})) \left| \det J_{\mathbf{g}^{-1}}(\mathbf{y}) \right| = \frac{f_{\mathbf{X}}(\mathbf{g}^{-1}(\mathbf{y}))}{\left| \det J_{\mathbf{g}}(\mathbf{g}^{-1}(\mathbf{y})) \right|}
\label{eq:jacobian_transform}
\end{equation}
where $J_{\mathbf{g}}(\mathbf{x}) \in \mathbb{R}^{d \times d}$ is the Jacobian matrix with entries $[J_{\mathbf{g}}(\mathbf{x})]_{ij} = \frac{\partial g_i}{\partial x_j}$.
\end{theorem}

\begin{proof}
Let $A \subset \mathbb{R}^d$ be an arbitrary measurable Borel set in the domain of $\mathbf{X}$, and let $B = \mathbf{g}(A)$ be its image in the codomain of $\mathbf{Y}$.
By conservation of probability mass:
\begin{equation}
P(\mathbf{Y} \in B) = P(\mathbf{X} \in A) \implies \int_B f_{\mathbf{Y}}(\mathbf{y}) \, d\mathbf{y} = \int_A f_{\mathbf{X}}(\mathbf{x}) \, d\mathbf{x}
\end{equation}
By the multivariable change of variables theorem for Riemann integrals, substituting $\mathbf{x} = \mathbf{g}^{-1}(\mathbf{y})$ transforms the differential volume element:
\begin{equation}
d\mathbf{x} = \left| \det J_{\mathbf{g}^{-1}}(\mathbf{y}) \right| \, d\mathbf{y}
\end{equation}
Thus:
\begin{equation}
\int_B f_{\mathbf{Y}}(\mathbf{y}) \, d\mathbf{y} = \int_B f_{\mathbf{X}}(\mathbf{g}^{-1}(\mathbf{y})) \left| \det J_{\mathbf{g}^{-1}}(\mathbf{y}) \right| \, d\mathbf{y}
\end{equation}
Since this holds for every measurable set $B$, the integrands must agree almost everywhere, proving equation~\eqref{eq:jacobian_transform}.
\end{proof}

\begin{ponderbox}[Why the Absolute Value of the Determinant?]
Why is the absolute value $|\det J|$ mandatory in the probability density transformation rule?

\textbf{Reflection:} Probability densities are intrinsically non-negative ($f(\mathbf{y}) \ge 0$). In multivariable calculus, the determinant $\det J$ is signed: a negative determinant indicates that the transformation reverses spatial orientation (like a reflection). However, physical volume dilation is strictly positive: a reflected cube still occupies positive volume. Taking the absolute value $|\det J|$ ensures the volume scaling factor remains positive and the transformed density integrates to 1.
\end{ponderbox}

\subsection{The Exponential Family of Distributions}

\begin{definitionbox}{The Canonical Exponential Family}
A parametric family of probability distributions over random variable $X$ belongs to the \textbf{Exponential Family} if its density (or mass function) can be expressed in canonical form:
\begin{equation}
p(x; \boldsymbol{\eta}) = h(x) \exp\left( \boldsymbol{\eta}^T \mathbf{T}(x) - A(\boldsymbol{\eta}) \right)
\label{eq:exp_family_def}
\end{equation}
where:
\begin{itemize}
    \item $\boldsymbol{\eta} \in \mathbb{R}^k$: Natural (canonical) parameter vector.
    \item $\mathbf{T}(x) \in \mathbb{R}^k$: Vector of sufficient statistics.
    \item $h(x) \ge 0$: Base reference measure.
    \item $A(\boldsymbol{\eta}) \triangleq \log \int_{\mathcal{X}} h(x) \exp\left( \boldsymbol{\eta}^T \mathbf{T}(x) \right) dx$: The \textbf{log-partition function}.
\end{itemize}
\end{definitionbox}

\begin{theorem}[Moment Generating Property of the Log-Partition Function]
Let $p(x; \boldsymbol{\eta})$ belong to the Exponential Family. The gradient and Hessian of the log-partition function $A(\boldsymbol{\eta})$ generate the mean and covariance of the sufficient statistics:
\begin{align}
\nabla_{\boldsymbol{\eta}} A(\boldsymbol{\eta}) &= \E[\mathbf{T}(X)] \label{eq:exp_mean} \\
\nabla_{\boldsymbol{\eta}}^2 A(\boldsymbol{\eta}) &= \Cov(\mathbf{T}(X)) \label{eq:exp_cov}
\end{align}
Because any covariance matrix is positive semidefinite ($\Cov(\mathbf{T}(X)) \succeq 0$), the log-partition function $A(\boldsymbol{\eta})$ is strictly convex over its natural domain.
\end{theorem}

% ----------------------------------------------------------------------
% PART 3: VISUAL & GEOMETRIC INTUITION
% ----------------------------------------------------------------------
\section{Visual \& Geometric Intuition: Probability Transformations and Volume Scaling}

Figure~\ref{fig:jacobian_geometry} illustrates the geometric intuition behind the multivariate transformation rule: the non-linear map $\mathbf{g}$ deforms a rectangular differential area element $dx_1 dx_2$ into a sheared parallelogram in $\mathbf{y}$-space, scaling the local area by the absolute Jacobian determinant $|\det J_{\mathbf{g}}|$.

\begin{figure}[htbp]
\centering
\begin{tikzpicture}[
    scale=0.95,
    font=\sffamily\footnotesize
]

% Left: Domain X with orthogonal differential box
\draw[thick, fill=primarylight!40] (-4.5, 0) rectangle ++(1.8cm, 1.4cm);
\draw[->, >=stealth, color=gray] (-5.5, 0) -- (-2.0, 0) node[right] {$x_1$};
\draw[->, >=stealth, color=gray] (-4.5, -0.8) -- (-4.5, 2.2) node[above] {$x_2$};
\node[below] at (-3.6, 0) {$dx_1$};
\node[left] at (-4.5, 0.7) {$dx_2$};
\node[above] at (-3.6, 0.5) {$dA_{\mathbf{x}} = dx_1 dx_2$};
\node[below, align=center] at (-3.6, -1.2) {\textbf{Domain: } Probability Mass \\ $f_{\mathbf{X}}(\mathbf{x}) d\mathbf{x}$};

% Transformation arrow
\draw[->, >=stealth, line width=2pt, color=primarydark] (-1.5, 0.8) -- (1.0, 0.8) node[midway, above] {$\mathbf{y} = \mathbf{g}(\mathbf{x})$};

% Right: Codomain Y with deformed parallelogram
\begin{scope}[shift={(3.0, 0)}]
    \draw[->, >=stealth, color=gray] (-1.0, 0) -- (3.5, 0) node[right] {$y_1$};
    \draw[->, >=stealth, color=gray] (0, -0.8) -- (0, 2.5) node[above] {$y_2$};

    % Parallelogram spanned by Jacobian column vectors
    \filldraw[fill=tealbg!50, draw=tealborder, thick] (0, 0) -- (2.2, 0.5) -- (2.8, 2.1) -- (0.6, 1.6) -- cycle;
    \draw[->, >=stealth, line width=1.5pt, color=tealborder] (0, 0) -- (2.2, 0.5) node[midway, below] {$J_{*,1} dx_1$};
    \draw[->, >=stealth, line width=1.5pt, color=tealborder] (0, 0) -- (0.6, 1.6) node[midway, left] {$J_{*,2} dx_2$};

    \node[above] at (1.4, 0.8) {$dA_{\mathbf{y}} = |\det J_{\mathbf{g}}| dx_1 dx_2$};
    \node[below, align=center] at (1.4, -1.2) {\textbf{Codomain: } $f_{\mathbf{Y}}(\mathbf{y}) d\mathbf{y}$ \\ $= f_{\mathbf{X}}(\mathbf{x}) d\mathbf{x}$};
\end{scope}

\end{tikzpicture}
\caption{The Geometry of Probability Transformations: Conservation of probability mass requires scaling the density inversely by the local volume dilation factor $|\det J_{\mathbf{g}}|$.}
\label{fig:jacobian_geometry}
\end{figure}

% ----------------------------------------------------------------------
% PART 4: DEEP INTUITION & CONCEPTUAL DERIVATIONS
% ----------------------------------------------------------------------
\section{Deep Intuition \& Conceptual Derivations}

\subsection{Derivation of the Moment Properties of \texorpdfstring{$A(\eta)$}{A(eta)}}
By definition of a normalized probability density, the integral of $p(x; \boldsymbol{\eta})$ over the support $\mathcal{X}$ must equal 1:
\begin{equation}
\int_{\mathcal{X}} h(x) \exp\left( \boldsymbol{\eta}^T \mathbf{T}(x) - A(\boldsymbol{\eta}) \right) dx = 1 \implies \exp(A(\boldsymbol{\eta})) = \int_{\mathcal{X}} h(x) \exp\left( \boldsymbol{\eta}^T \mathbf{T}(x) \right) dx
\label{eq:exp_norm}
\end{equation}
Differentiating both sides with respect to the natural parameter $\boldsymbol{\eta}$:
\begin{equation}
\nabla_{\boldsymbol{\eta}} \exp(A(\boldsymbol{\eta})) = \exp(A(\boldsymbol{\eta})) \, \nabla_{\boldsymbol{\eta}} A(\boldsymbol{\eta}) = \int_{\mathcal{X}} h(x) \mathbf{T}(x) \exp\left( \boldsymbol{\eta}^T \mathbf{T}(x) \right) dx
\end{equation}
Dividing both sides by $\exp(A(\boldsymbol{\eta}))$:
\begin{equation}
\nabla_{\boldsymbol{\eta}} A(\boldsymbol{\eta}) = \int_{\mathcal{X}} \mathbf{T}(x) \underbrace{h(x) \exp\left( \boldsymbol{\eta}^T \mathbf{T}(x) - A(\boldsymbol{\eta}) \right)}_{p(x; \boldsymbol{\eta})} dx = \E[\mathbf{T}(X)]
\end{equation}
Differentiating a second time using the product rule:
\begin{align}
\nabla_{\boldsymbol{\eta}}^2 A(\boldsymbol{\eta}) &= \nabla_{\boldsymbol{\eta}} \left( \int_{\mathcal{X}} \mathbf{T}(x) h(x) \exp\left( \boldsymbol{\eta}^T \mathbf{T}(x) - A(\boldsymbol{\eta}) \right) dx \right) \nonumber \\
&= \int_{\mathcal{X}} \mathbf{T}(x) h(x) \exp\left( \boldsymbol{\eta}^T \mathbf{T}(x) - A(\boldsymbol{\eta}) \right) \left[ \mathbf{T}(x) - \nabla_{\boldsymbol{\eta}} A(\boldsymbol{\eta}) \right]^T dx \nonumber \\
&= \E\left[ \mathbf{T}(X) (\mathbf{T}(X) - \E[\mathbf{T}(X)])^T \right] = \Cov(\mathbf{T}(X))
\end{align}
Because the covariance matrix of any random vector satisfies $\mathbf{v}^T \Cov(\mathbf{T})\mathbf{v} = \Var(\mathbf{v}^T \mathbf{T}) \ge 0$, the Hessian $\nabla^2 A(\boldsymbol{\eta})$ is universally positive semidefinite. This proves that $A(\boldsymbol{\eta})$ is a \textbf{strictly convex functional} over its natural parameter space.

\subsection{Canonical Representation: The Bernoulli Distribution and the Sigmoid Function}
Let $X \in \{0, 1\}$ with success parameter $P(X = 1) = \mu \in (0, 1)$. The probability mass function is:
\begin{align}
p(x; \mu) &= \mu^x (1 - \mu)^{1 - x} = \exp\left( \log\left( \mu^x (1 - \mu)^{1 - x} \right) \right) \nonumber \\
&= \exp\left( x \log\mu + (1 - x)\log(1 - \mu) \right) \nonumber \\
&= \exp\left( x \log\left(\frac{\mu}{1 - \mu}\right) + \log(1 - \mu) \right)
\end{align}
Matching this expression with the canonical template $h(x) \exp\left( \eta T(x) - A(\eta) \right)$:
\begin{itemize}
    \item Base measure: $h(x) = 1$.
    \item Sufficient statistic: $T(x) = x$.
    \item Natural parameter: $\eta = \log\left(\frac{\mu}{1 - \mu}\right)$ (the \textbf{Log-Odds} or \textbf{Logit} function).
    \item Inverting for $\mu$ in terms of $\eta$:
    \begin{equation}
    e^\eta = \frac{\mu}{1 - \mu} \implies \mu = \frac{e^\eta}{1 + e^\eta} = \frac{1}{1 + e^{-\eta}} \triangleq \sigma(\eta)
    \end{equation}
    This is precisely the \textbf{logistic sigmoid activation function}!
    \item Log-partition function:
    \begin{equation}
    A(\eta) = -\log(1 - \mu) = -\log\left(1 - \frac{e^\eta}{1 + e^\eta}\right) = -\log\left(\frac{1}{1 + e^\eta}\right) = \log(1 + e^\eta)
    \end{equation}
\end{itemize}
Evaluating the derivative:
\begin{equation}
A'(\eta) = \frac{d}{d\eta} \log(1 + e^\eta) = \frac{e^\eta}{1 + e^\eta} = \frac{1}{1 + e^{-\eta}} = \mu = \E[X]
\end{equation}
Evaluating the second derivative:
\begin{equation}
A''(\eta) = \frac{d}{d\eta} \sigma(\eta) = \sigma(\eta)(1 - \sigma(\eta)) = \mu(1 - \mu) = \Var(X)
\end{equation}
This derivation reveals why the logistic sigmoid is not an arbitrary engineering heuristic: it is the mathematically unique canonical inverse link function for Bernoulli-distributed binary targets.

% ----------------------------------------------------------------------
% PART 5: MULTI-PAGE WORKED EXAMPLES
% ----------------------------------------------------------------------
\section{Multi-Page Breakable Numerical Worked Examples}

\begin{examplebox}[breakable]{Analytical Linear Transformation of a 2D Gaussian with Exact Covariance Reconstruction}
\textbf{Problem Statement:}
Let $\mathbf{X} = [X_1, X_2]^T \sim \mathcal{N}(\mathbf{0}, I_2)$ be an isotropic bivariate standard Gaussian vector with joint density:
\begin{equation}
f_{\mathbf{X}}(\mathbf{x}) = \frac{1}{2\pi} \exp\left( -\frac{1}{2}(x_1^2 + x_2^2) \right) = \frac{1}{2\pi} \exp\left( -\frac{1}{2}\mathbf{x}^T \mathbf{x} \right)
\end{equation}
Define the linear vector transformation $\mathbf{Y} = A\mathbf{X}$, where:
\begin{equation}
A = \begin{bmatrix} 2 & 1 \\ 0 & 3 \end{bmatrix}
\end{equation}

\textbf{Task Requirements:}
\begin{enumerate}
    \item Compute the determinant $\det(A)$ and the inverse matrix $A^{-1}$.
    \item Apply the Jacobian transformation theorem~\eqref{eq:jacobian_transform} to derive the exact joint PDF $f_{\mathbf{Y}}(y_1, y_2)$.
    \item Compute the theoretical covariance matrix $\Sigma_{\mathbf{Y}} = A A^T$ and its inverse $\Sigma_{\mathbf{Y}}^{-1}$.
    \item Expand the quadratic form $-\frac{1}{2}\mathbf{y}^T \Sigma_{\mathbf{Y}}^{-1}\mathbf{y}$ and verify that it matches the exponent derived via the Jacobian rule term-by-term.
\end{enumerate}

\tcblower

\textbf{Step 1: Determinant and Matrix Inversion}
The determinant is:
\begin{equation}
\det(A) = (2)(3) - (1)(0) = 6
\end{equation}
Since $\det(A) = 6 \ne 0$, the transformation is invertible.
The inverse is:
\begin{equation}
A^{-1} = \frac{1}{6} \begin{bmatrix} 3 & -1 \\ 0 & 2 \end{bmatrix} = \begin{bmatrix} \frac{1}{2} & -\frac{1}{6} \\ 0 & \frac{1}{3} \end{bmatrix}
\end{equation}

\textbf{Step 2: Change of Variables via Jacobian Transformation}
The inverse transformation mapping $\mathbf{y} \mapsto \mathbf{x}$ is:
\begin{equation}
\mathbf{x} = A^{-1}\mathbf{y} = \begin{bmatrix} \frac{1}{2}y_1 - \frac{1}{6}y_2 \\ \frac{1}{3}y_2 \end{bmatrix}
\end{equation}
The Jacobian matrix of the inverse mapping is:
\begin{equation}
J_{A^{-1}}(\mathbf{y}) = A^{-1} \implies |\det J_{A^{-1}}(\mathbf{y})| = |\det(A^{-1})| = \frac{1}{|\det(A)|} = \frac{1}{6}
\end{equation}
Now substitute $\mathbf{x} = A^{-1}\mathbf{y}$ into the original density $f_{\mathbf{X}}(\mathbf{x})$:
\begin{align}
\|\mathbf{x}\|_2^2 &= x_1^2 + x_2^2 = \left( \frac{1}{2}y_1 - \frac{1}{6}y_2 \right)^2 + \left( \frac{1}{3}y_2 \right)^2 \nonumber \\
&= \frac{1}{4}y_1^2 - \frac{2}{12}y_1 y_2 + \frac{1}{36}y_2^2 + \frac{1}{9}y_2^2 \nonumber \\
&= \frac{1}{4}y_1^2 - \frac{1}{6}y_1 y_2 + \left( \frac{1}{36} + \frac{4}{36} \right)y_2^2 \nonumber \\
&= \frac{1}{4}y_1^2 - \frac{1}{6}y_1 y_2 + \frac{5}{36}y_2^2
\end{align}
Applying the Jacobian formula:
\begin{align}
f_{\mathbf{Y}}(\mathbf{y}) &= f_{\mathbf{X}}(A^{-1}\mathbf{y}) \cdot |\det J_{A^{-1}}(\mathbf{y})| \nonumber \\
&= \frac{1}{2\pi} \exp\left( -\frac{1}{2} \left[ \frac{1}{4}y_1^2 - \frac{1}{6}y_1 y_2 + \frac{5}{36}y_2^2 \right] \right) \cdot \frac{1}{6} \nonumber \\
&= \frac{1}{12\pi} \exp\left( -\frac{1}{2} \left[ \frac{1}{4}y_1^2 - \frac{1}{6}y_1 y_2 + \frac{5}{36}y_2^2 \right] \right)
\label{eq:ex_fy_derived}
\end{align}

\textbf{Step 3: Theoretical Covariance Matrix and Its Inverse}
Since $\mathbf{X} \sim \mathcal{N}(\mathbf{0}, I_2)$, the covariance of the linear transformation $\mathbf{Y} = A\mathbf{X}$ is:
\begin{equation}
\Sigma_{\mathbf{Y}} = \E[\mathbf{Y}\mathbf{Y}^T] = \E[(A\mathbf{X})(A\mathbf{X})^T] = A \E[\mathbf{X}\mathbf{X}^T] A^T = A I_2 A^T = A A^T
\end{equation}
Performing the matrix multiplication:
\begin{equation}
\Sigma_{\mathbf{Y}} = \begin{bmatrix} 2 & 1 \\ 0 & 3 \end{bmatrix} \begin{bmatrix} 2 & 0 \\ 1 & 3 \end{bmatrix} = \begin{bmatrix} 2(2) + 1(1) & 2(0) + 1(3) \\ 0(2) + 3(1) & 0(0) + 3(3) \end{bmatrix} = \begin{bmatrix} 5 & 3 \\ 3 & 9 \end{bmatrix}
\end{equation}
Compute the determinant of $\Sigma_{\mathbf{Y}}$:
\begin{equation}
\det(\Sigma_{\mathbf{Y}}) = (5)(9) - (3)(3) = 45 - 9 = 36 = (\det A)^2
\end{equation}
The normalizing constant for a 2D Gaussian with covariance $\Sigma_{\mathbf{Y}}$ is:
\begin{equation}
\frac{1}{2\pi \sqrt{\det \Sigma_{\mathbf{Y}}}} = \frac{1}{2\pi \sqrt{36}} = \frac{1}{12\pi}
\end{equation}
This matches the normalizing prefactor in~\eqref{eq:ex_fy_derived} exactly!

Now compute the precision matrix $\Sigma_{\mathbf{Y}}^{-1}$:
\begin{equation}
\Sigma_{\mathbf{Y}}^{-1} = \frac{1}{36} \begin{bmatrix} 9 & -3 \\ -3 & 5 \end{bmatrix} = \begin{bmatrix} \frac{9}{36} & -\frac{3}{36} \\ -\frac{3}{36} & \frac{5}{36} \end{bmatrix} = \begin{bmatrix} \frac{1}{4} & -\frac{1}{12} \\ -\frac{1}{12} & \frac{5}{36} \end{bmatrix}
\end{equation}

\textbf{Step 4: Verification of the Quadratic Form}
Compute the quadratic form $\mathbf{y}^T \Sigma_{\mathbf{Y}}^{-1} \mathbf{y}$:
\begin{align}
\mathbf{y}^T \Sigma_{\mathbf{Y}}^{-1} \mathbf{y} &= \begin{bmatrix} y_1 & y_2 \end{bmatrix} \begin{bmatrix} \frac{1}{4} & -\frac{1}{12} \\ -\frac{1}{12} & \frac{5}{36} \end{bmatrix} \begin{bmatrix} y_1 \\ y_2 \end{bmatrix} \nonumber \\
&= \begin{bmatrix} y_1 & y_2 \end{bmatrix} \begin{bmatrix} \frac{1}{4}y_1 - \frac{1}{12}y_2 \\ -\frac{1}{12}y_1 + \frac{5}{36}y_2 \end{bmatrix} \nonumber \\
&= y_1\left( \frac{1}{4}y_1 - \frac{1}{12}y_2 \right) + y_2\left( -\frac{1}{12}y_1 + \frac{5}{36}y_2 \right) \nonumber \\
&= \frac{1}{4}y_1^2 - \frac{1}{12}y_1 y_2 - \frac{1}{12}y_1 y_2 + \frac{5}{36}y_2^2 \nonumber \\
&= \frac{1}{4}y_1^2 - \frac{1}{6}y_1 y_2 + \frac{5}{36}y_2^2
\end{align}
The quadratic form matches the exponent in equation~\eqref{eq:ex_fy_derived} identically. This confirms that the Jacobian transformation of a Gaussian under linear mapping $A$ produces another Gaussian $\mathcal{N}(\mathbf{0}, A A^T)$.
\end{examplebox}

\begin{examplebox}[breakable]{Non-Monotonic Transformation: Derivation of the Chi-Square Distribution}
\textbf{Problem Statement:}
Let $X \sim \mathcal{N}(0, 1)$ be a standard univariate normal random variable with PDF:
\begin{equation}
f_X(x) = \frac{1}{\sqrt{2\pi}} \exp\left( -\frac{x^2}{2} \right) \quad \text{for } x \in (-\infty, \infty)
\end{equation}
Consider the non-linear transformation:
\begin{equation}
Y = g(X) = X^2
\end{equation}

\textbf{Task Requirements:}
\begin{enumerate}
    \item Determine the support of $Y$.
    \item Identify all pre-images $x$ corresponding to a given value $y$ in the support.
    \item Compute the derivative of each inverse branch.
    \item Sum the probability contributions across all branches to derive the exact PDF $f_Y(y)$.
    \item Recognize the resulting distribution as a member of the Gamma family ($\chi^2$ with 1 degree of freedom).
\end{enumerate}

\tcblower

\textbf{Step 1: Support of $Y$}
Since $X^2 \ge 0$ for all real $X$, the support of $Y$ is strictly non-negative: $\mathcal{Y} = (0, \infty)$ (the event $Y = 0$ occurs with probability zero). For $y \le 0$, $f_Y(y) = 0$.

\textbf{Step 2: Pre-Images of the Non-Monotonic Transformation}
For any $y > 0$, the equation $x^2 = y$ has \textbf{two distinct roots}:
\begin{align}
x_1(y) &= +\sqrt{y} = y^{1/2} \\
x_2(y) &= -\sqrt{y} = -y^{1/2}
\end{align}
Because $g(x) = x^2$ is a two-to-one mapping, probability mass accumulates in $dy$ from two separate intervals in the domain of $X$: $[-\sqrt{y+dy}, -\sqrt{y}]$ and $[\sqrt{y}, \sqrt{y+dy}]$.

\textbf{Step 3: Derivatives of the Inverse Branches}
Differentiating each branch with respect to $y$:
\begin{align}
\frac{dx_1}{dy} &= \frac{d}{dy}\left( y^{1/2} \right) = \frac{1}{2}y^{-1/2} = \frac{1}{2\sqrt{y}} \\
\frac{dx_2}{dy} &= \frac{d}{dy}\left( -y^{1/2} \right) = -\frac{1}{2}y^{-1/2} = -\frac{1}{2\sqrt{y}}
\end{align}
Taking the absolute values:
\begin{equation}
\left| \frac{dx_1}{dy} \right| = \frac{1}{2\sqrt{y}}, \quad \left| \frac{dx_2}{dy} \right| = \left| -\frac{1}{2\sqrt{y}} \right| = \frac{1}{2\sqrt{y}}
\end{equation}

\textbf{Step 4: Summing Probability Contributions}
The multi-branch transformation formula is:
\begin{equation}
f_Y(y) = f_X(x_1(y))\left| \frac{dx_1}{dy} \right| + f_X(x_2(y))\left| \frac{dx_2}{dy} \right|
\end{equation}
Evaluating $f_X$ at both pre-images:
\begin{align}
f_X(+\sqrt{y}) &= \frac{1}{\sqrt{2\pi}} \exp\left( -\frac{(\sqrt{y})^2}{2} \right) = \frac{1}{\sqrt{2\pi}} e^{-y/2} \\
f_X(-\sqrt{y}) &= \frac{1}{\sqrt{2\pi}} \exp\left( -\frac{(-\sqrt{y})^2}{2} \right) = \frac{1}{\sqrt{2\pi}} e^{-y/2}
\end{align}
Substituting into the formula:
\begin{align}
f_Y(y) &= \frac{1}{\sqrt{2\pi}} e^{-y/2} \left( \frac{1}{2\sqrt{y}} \right) + \frac{1}{\sqrt{2\pi}} e^{-y/2} \left( \frac{1}{2\sqrt{y}} \right) \nonumber \\
&= 2 \cdot \frac{1}{\sqrt{2\pi}} e^{-y/2} \left( \frac{1}{2\sqrt{y}} \right) \nonumber \\
&= \frac{1}{\sqrt{2\pi}} y^{-1/2} e^{-y/2} \quad \text{for } y > 0
\end{align}

\textbf{Step 5: Connection to the Chi-Square and Gamma Distributions}
The standard Gamma PDF with shape $\alpha$ and rate $\beta$ is:
\begin{equation}
\Gamma(y; \alpha, \beta) = \frac{\beta^\alpha}{\Gamma(\alpha)} y^{\alpha - 1} e^{-\beta y}
\end{equation}
Substituting $\alpha = 1/2$ and $\beta = 1/2$, and using $\Gamma(1/2) = \sqrt{\pi}$:
\begin{equation}
\Gamma(y; 1/2, 1/2) = \frac{(1/2)^{1/2}}{\sqrt{\pi}} y^{1/2 - 1} e^{-y/2} = \frac{1}{\sqrt{2\pi}} y^{-1/2} e^{-y/2}
\end{equation}
This is the exact Chi-Square distribution with 1 degree of freedom: $Y \sim \chi^2(1)$.
\end{examplebox}

% ----------------------------------------------------------------------
% PART 6: EXPLICIT GOTCHAS & DIAGNOSTIC MISCONCEPTIONS
% ----------------------------------------------------------------------
\section{Explicit Gotchas \& Diagnostic Misconceptions}

\begin{trapbox}{Gotcha 1: Omitting the Absolute Value of the Jacobian Determinant}
A student computes a change of variables with $y = -2x$ and writes $f_Y(y) = f_X(-y/2) \cdot (-1/2)$.
\textbf{The Trap:} Probability densities are strictly non-negative: $f_Y(y) \ge 0$. The Jacobian factor represents a ratio of geometric volumes, which is always positive. The transformation formula requires the \textbf{absolute value} of the determinant: $|\det J_{\mathbf{g}^{-1}}| = |-1/2| = 1/2$. Neglecting the absolute value produces negative densities.
\end{trapbox}

\begin{trapbox}{Gotcha 2: Believing Uncorrelated Random Variables are Always Independent}
A data scientist computes sample covariance $\Cov(X, Y) = 0$ and concludes that features $X$ and $Y$ are statistically independent ($X \perp Y$).
\textbf{The Trap:} Zero covariance means there is no \textbf{linear} relationship. Non-linear dependence can be absolute. For example, let $X \sim \mathcal{U}[-1, 1]$ and $Y = X^2$. Then $\E[X] = 0$ and $\E[XY] = \E[X^3] = 0 \implies \Cov(X, Y) = 0$. However, $Y$ is completely determined by $X$; knowing $X$ completely specifies $Y$. Uncorrelatedness implies independence \textbf{if and only if} $X$ and $Y$ are jointly Gaussian.
\end{trapbox}

\begin{trapbox}{Gotcha 3: Conflating Probability Density with Probability}
A programmer evaluates a continuous Gaussian PDF and observes $f_X(0.0) = 3.989 > 1$. They declare a software bug because ``probabilities cannot exceed 1.''
\textbf{The Trap:} For continuous random variables, $f_X(x)$ is a \textbf{probability density}, not a probability. Densities have units of $1/\text{measurement}$ and can be arbitrarily large (e.g., $f(x) = 1000$ for $x \in [0, 0.001]$). The probability of an exact point is zero: $P(X = x) = 0$. Only integrals over intervals represent valid probabilities bounded in $[0, 1]$.
\end{trapbox}

\begin{trapbox}{Gotcha 4: Forgetting Multiple Pre-Images in Non-Monotonic Transformations}
An engineer computes the PDF of $Y = X^2$ for $X \sim \mathcal{N}(0, 1)$ by simply differentiating $x = \sqrt{y}$.
\textbf{The Trap:} The function $g(x) = x^2$ is not one-to-one; for each $y > 0$, there are two pre-images: $x_1 = +\sqrt{y}$ and $x_2 = -\sqrt{y}$. The transformation rule must sum over all branches:
$f_Y(y) = f_X(\sqrt{y})\left|\frac{1}{2\sqrt{y}}\right| + f_X(-\sqrt{y})\left|-\frac{1}{2\sqrt{y}}\right| = \frac{f_X(\sqrt{y})}{\sqrt{y}}$, which yields the Chi-square distribution.
\end{trapbox}

% ----------------------------------------------------------------------
% PART 7: COMPARATIVE SYNTHESIS & SUMMARY TABLE
% ----------------------------------------------------------------------
\section{Comparative Synthesis \& Summary Table}

Table~\ref{tab:exp_family_summary} summarizes the Exponential Family representations for canonical machine learning distributions.

\begin{table}[htbp]
\centering
\footnotesize
\setlength{\tabcolsep}{3pt}
\begin{tabularx}{\textwidth}{lcXX>{\raggedright\arraybackslash}Xp{2.8cm}}
\toprule
\textbf{Distribution} & \textbf{Support $\mathcal{X}$} & \textbf{Sufficient Stat $\mathbf{T}(x)$} & \textbf{Natural Param $\boldsymbol{\eta}$} & \textbf{Log-Partition $A(\boldsymbol{\eta})$} & \textbf{ML Model Link} \\
\midrule
Bernoulli($\mu$) & $\{0, 1\}$ & $x$ & $\log\left(\frac{\mu}{1 - \mu}\right)$ & $\log(1 + e^\eta)$ & Logistic Regression (Sigmoid) \\
\addlinespace
Gaussian($\mu, \sigma^2$) & $\mathbb{R}$ & $\begin{bmatrix} x \\ x^2 \end{bmatrix}$ & $\begin{bmatrix} \mu/\sigma^2 \\ -1/(2\sigma^2) \end{bmatrix}$ & $-\frac{\eta_1^2}{4\eta_2} - \frac{1}{2}\log(-2\eta_2)$ & Linear Regression (Squared Loss) \\
\addlinespace
Poisson($\lambda$) & $\mathbb{N}_0$ & $x$ & $\log\lambda$ & $e^\eta$ & Poisson Regression (Count Data) \\
\addlinespace
Exponential($\lambda$) & $[0, \infty)$ & $x$ & $-\lambda$ & $-\log(-\eta)$ & Survival Analysis, Lifetime Modeling \\
\addlinespace
Categorical($\mathbf{p}$) & $\{1, \dots, K\}$ & $[\mathbb{I}(x=k)]_{k=1}^{K-1}$ & $[\log(p_k/p_K)]_{k=1}^{K-1}$ & $\log\left(1 + \sum_{k=1}^{K-1} e^{\eta_k}\right)$ & Multi-class Softmax Regression \\
\bottomrule
\end{tabularx}
\caption{Comparative Synthesis: Canonical Exponential Family Parameterizations.}
\label{tab:exp_family_summary}
\end{table}

% ----------------------------------------------------------------------
% PART 8: PYTHON / SCIENTIFIC COMPUTING LAB
% ----------------------------------------------------------------------
\section{Python / Scientific Computing Lab}

The following production script verifies continuous transformation densities via Jacobian scaling and numerically confirms the moment-generating property of the Exponential Family log-partition function.

\begin{lstlisting}[language=Python, caption={Transformation of Densities and Exponential Family Verification in Python.}]
import numpy as np
from scipy.integrate import quad

def test_jacobian_transformation():
    """Verifies Y = exp(X) transformation: if X ~ Normal(0, 1), Y ~ LogNormal."""
    rng = np.random.RandomState(42)
    N = 100000
    x_samples = rng.randn(N)
    # Transformation: g(x) = exp(x) => g^{-1}(y) = log(y), |d/dy g^{-1}(y)| = 1/y
    y_samples = np.exp(x_samples)

    # Theoretical log-normal density: f_Y(y) = f_X(log(y)) * (1 / y)
    def theoretical_pdf(y):
        if y <= 0: return 0.0
        fx = (1.0 / np.sqrt(2.0 * np.pi)) * np.exp(-0.5 * (np.log(y)) ** 2)
        return fx * (1.0 / y)

    # Compare empirical histogram against theoretical density at y = 1.5
    y_eval = 1.5
    eps = 0.05
    empirical_density = float(np.mean((y_samples >= y_eval - eps) & (y_samples <= y_eval + eps)) / (2 * eps))
    theory_val = theoretical_pdf(y_eval)

    print("Verifying Multivariate / Jacobian Transformation:")
    print(f"  At y = {y_eval}: Empirical Density = {empirical_density:.4f}, Theoretical = {theory_val:.4f}")

def verify_exponential_family_moments():
    """Verifies that A'(eta) = E[T(X)] for the Bernoulli distribution."""
    # Bernoulli with mu = 0.75
    mu_true = 0.75
    eta = float(np.log(mu_true / (1.0 - mu_true)))

    # A(eta) = log(1 + exp(eta))
    def A(eta_val):
        return np.log(1.0 + np.exp(eta_val))

    # Numerical derivative A'(eta)
    eps = 1e-6
    grad_A = float((A(eta + eps) - A(eta - eps)) / (2.0 * eps))

    print("\nVerifying Exponential Family Moment Generation (Bernoulli):")
    print(f"  Natural Parameter eta:        {eta:.4f}")
    print(f"  True Mean E[X] = mu:          {mu_true:.4f}")
    print(f"  Derivative of Log-Partition:  {grad_A:.4f}")
    print(f"  Difference:                   {abs(grad_A - mu_true):.2e}")

def main():
    test_jacobian_transformation()
    verify_exponential_family_moments()

if __name__ == "__main__":
    main()
\end{lstlisting}

% ----------------------------------------------------------------------
% PART 9: COMPREHENSIVE PRACTICE PROBLEMS
% ----------------------------------------------------------------------
\section{Comprehensive Practice Problems}

\subsection*{Part A: Analytical \& Mathematical Problems}

\begin{enumerate}[label=\textbf{P\arabic*.}]
    \item \textbf{Analytical Linear Transformation of a 2D Gaussian:}
    Let $\mathbf{X} = [X_1, X_2]^T \sim \mathcal{N}(\mathbf{0}, I_2)$ be standard bivariate normal.
    Define linear transformation $\mathbf{Y} = A\mathbf{X}$, where $A = \begin{bmatrix} 2 & 1 \\ 0 & 3 \end{bmatrix}$.
    \begin{enumerate}[label=(\alph*)]
        \item Compute the determinant $\det(A)$ and inverse matrix $A^{-1}$.
        \item Use the Jacobian change-of-variables formula~\eqref{eq:jacobian_transform} to derive the exact joint PDF $f_{\mathbf{Y}}(y_1, y_2)$.
        \item Compute the covariance matrix $\Sigma_{\mathbf{Y}} = A A^T$ and confirm that the exponent matches $-\frac{1}{2}\mathbf{y}^T \Sigma_{\mathbf{Y}}^{-1}\mathbf{y}$.
    \end{enumerate}

    \item \textbf{Non-Monotonic Transformation (Chi-Square Distribution):}
    Let $X \sim \mathcal{N}(0, 1)$ with density $f_X(x) = \frac{1}{\sqrt{2\pi}}e^{-x^2/2}$. Let $Y = X^2$.
    \begin{enumerate}[label=(\alph*)]
        \item Identify the two pre-images $x_1(y)$ and $x_2(y)$ for $y > 0$.
        \item Compute the derivative of each inverse branch.
        \item Derive the probability density function $f_Y(y)$ and show it is a Gamma / Chi-Square distribution with 1 degree of freedom: $f_Y(y) = \frac{1}{\sqrt{2\pi}} y^{-1/2} e^{-y/2}$ for $y > 0$.
    \end{enumerate}

    \item \textbf{Exponential Distribution in Canonical Form:}
    The continuous Exponential distribution has PDF $f(x; \lambda) = \lambda e^{-\lambda x}$ for $x \ge 0$, with rate $\lambda > 0$.
    \begin{enumerate}[label=(\alph*)]
        \item Express $f(x; \lambda)$ in canonical Exponential Family form $h(x)\exp(\eta T(x) - A(\eta))$, identifying $h(x), T(x), \eta$, and $A(\eta)$.
        \item Differentiate $A(\eta)$ to compute the mean $\E[X]$.
        \item Differentiate $A(\eta)$ a second time to compute the variance $\Var(X)$.
    \end{enumerate}

    \item \textbf{Uncorrelated but Strongly Dependent Variables:}
    Let continuous random variable $X \sim \mathcal{U}[-1, 1]$. Define $Y = X^2$.
    \begin{enumerate}[label=(\alph*)]
        \item Compute $\E[X]$, $\E[Y]$, and $\E[XY]$.
        \item Show that $\Cov(X, Y) = 0$.
        \item Prove that $P(Y \le 0.25 \mid X \ge 0.8) \ne P(Y \le 0.25)$, formally proving that $X$ and $Y$ are statistically dependent despite zero covariance.
    \end{enumerate}

    \item \textbf{Derivation of the Softmax Link from Categorical Distribution:}
    Let $X \in \{1, 2, \dots, K\}$ follow a Categorical distribution with $P(X = k) = p_k$, where $\sum_{k=1}^K p_k = 1$.
    Using the reference class $K$, write the distribution in canonical Exponential Family form and prove that the inverse mapping from natural parameters $\eta_1, \dots, \eta_{K-1}$ to probabilities $p_k$ is the Softmax function:
    \begin{equation}
    p_k = \frac{e^{\eta_k}}{1 + \sum_{j=1}^{K-1} e^{\eta_j}}
    \end{equation}
\end{enumerate}

\subsection*{Part B: Diagnostic Multiple-Choice Questions}

\begin{enumerate}[label=\textbf{MCQ\arabic*.}]
    \item Let $\mathbf{Y} = A\mathbf{X} + \mathbf{b}$, where $A \in \mathbb{R}^{d \times d}$ is invertible and $\mathbf{X}$ has continuous density $f_{\mathbf{X}}(\mathbf{x})$. What is the correct probability density function $f_{\mathbf{Y}}(\mathbf{y})$?
    \begin{enumerate}[label=(\alph*)]
        \item $f_{\mathbf{Y}}(\mathbf{y}) = f_{\mathbf{X}}(A^{-1}(\mathbf{y} - \mathbf{b})) \cdot \det(A)$
        \item $f_{\mathbf{Y}}(\mathbf{y}) = f_{\mathbf{X}}(A^{-1}(\mathbf{y} - \mathbf{b})) \cdot \frac{1}{|\det(A)|}$
        \item $f_{\mathbf{Y}}(\mathbf{y}) = f_{\mathbf{X}}(A\mathbf{y} + \mathbf{b}) \cdot |\det(A)|$
        \item $f_{\mathbf{Y}}(\mathbf{y}) = f_{\mathbf{X}}(A^{-1}(\mathbf{y} - \mathbf{b}))$
    \end{enumerate}

    \item In the canonical Exponential Family $p(x; \boldsymbol{\eta}) = h(x)\exp(\boldsymbol{\eta}^T \mathbf{T}(x) - A(\boldsymbol{\eta}))$, what is the fundamental property of the Hessian $\nabla_{\boldsymbol{\eta}}^2 A(\boldsymbol{\eta})$?
    \begin{enumerate}[label=(\alph*)]
        \item It equals the negative variance of $X$ and is negative definite.
        \item It is the covariance matrix $\Cov(\mathbf{T}(X))$ and is guaranteed to be positive semidefinite, proving $A(\boldsymbol{\eta})$ is convex.
        \item It evaluates to zero for all linear models.
        \item It is equal to the Fisher information matrix multiplied by $-1$.
    \end{enumerate}
\end{enumerate}

\chapter{Multivariate Gaussian Distributions, Maximum Likelihood, GMMs, EM Algorithm \& Concentration Inequalities}
\label{chap:week12}

% ----------------------------------------------------------------------
% PART 1: EXECUTIVE ROADMAP & LEARNING OBJECTIVES
% ----------------------------------------------------------------------
\section{Executive Roadmap \& Learning Objectives}

The multivariate normal (Gaussian) distribution is the cornerstone of statistical machine learning. By virtue of the Central Limit Theorem, aggregated noise fluctuations across physical systems converge toward Gaussianity. However, real-world data often exhibits multimodality and latent cluster structure that a single unimodal Gaussian cannot capture. \textbf{Gaussian Mixture Models (GMMs)} resolve this by modeling data as a convex combination of Gaussians, estimated iteratively via the **Expectation-Maximization (EM) algorithm**. Finally, **Concentration Inequalities** provide non-asymptotic probabilistic bounds governing sample sizes and generalization error.

\subsection{Prerequisites \& Scaffolding}
Mastery of this chapter synthesizes positive definite matrices, Jensen's inequality, and probability theory as organized in Table~\ref{tab:w12_prereqs}.

\begin{table}[htbp]
\centering
\small
\begin{tabularx}{\textwidth}{lXX}
\toprule
\textbf{Foundational Area} & \textbf{Core Concepts Required} & \textbf{Role in Machine Learning Foundations} \\
\midrule
Linear Algebra & Matrix inverse $\Sigma^{-1}$, determinant $|\Sigma|$, Spectral Theorem $\Sigma = Q \Lambda Q^T$ & Defining Mahalanobis distance ellipsoids and normalizing constants \\
Convex Analysis & Jensen's Inequality $f(\E[X]) \le \E[f(X)]$, concave logarithms & Deriving the Evidence Lower Bound (ELBO) for the EM algorithm \\
Probability & Conditional density, Bayes' rule, marginalization, variance & Formulating posterior responsibilities and mixture likelihoods \\
\bottomrule
\end{tabularx}
\caption{Prerequisite Scaffolding for Chapter 12.}
\label{tab:w12_prereqs}
\end{table}

\subsection{Core Competencies}
By the conclusion of this chapter, students will be able to:
\begin{itemize}
    \item Evaluate multivariate Gaussian densities and compute Mahalanobis distance metrics.
    \item Derive the Maximum Likelihood Estimators (MLE) for multivariate Gaussian mean $\boldsymbol{\mu}$ and covariance $\Sigma$.
    \item Formulate Gaussian Mixture Models with latent categorical cluster assignments.
    \item Derive the complete Expectation-Maximization (EM) algorithm from first principles via the Evidence Lower Bound (ELBO).
    \item State and apply Markov's, Chebyshev's, and Hoeffding's concentration inequalities to establish rigorous sample complexity bounds.
\end{itemize}

\subsection{Notation \& Mathematical Conventions}
\begin{itemize}
    \item $\mathcal{N}(\mathbf{x}; \boldsymbol{\mu}, \Sigma)$: The $d$-dimensional multivariate Gaussian probability density.
    \item $D_M(\mathbf{x}, \boldsymbol{\mu}) \triangleq \sqrt{(\mathbf{x} - \boldsymbol{\mu})^T \Sigma^{-1} (\mathbf{x} - \boldsymbol{\mu})}$: The Mahalanobis distance.
    \item $K$: The number of latent mixture components.
    \item $\pi_k \triangleq P(z = k)$: Mixture weight for component $k$ ($\sum_{k=1}^K \pi_k = 1, \pi_k \ge 0$).
    \item $\gamma_{ik} \triangleq P(z_i = k \mid \mathbf{x}_i)$: The posterior responsibility of component $k$ for sample $i$.
    \item $\text{ELBO}(q, \boldsymbol{\theta})$: The Evidence Lower Bound on the marginal log-likelihood.
    \item $\bar{X}_N \triangleq \frac{1}{N}\sum_{i=1}^N X_i$: The empirical sample mean.
\end{itemize}

% ----------------------------------------------------------------------
% PART 2: MATHEMATICAL & COMPUTATIONAL FOUNDATIONS
% ----------------------------------------------------------------------
\section{Mathematical \& Computational Foundations}

\subsection{The Multivariate Gaussian Distribution}

\begin{definitionbox}{Multivariate Normal Density}
A continuous random vector $\mathbf{X} \in \mathbb{R}^d$ follows a \textbf{Multivariate Gaussian distribution} $\mathbf{X} \sim \mathcal{N}(\boldsymbol{\mu}, \Sigma)$ with mean $\boldsymbol{\mu} \in \mathbb{R}^d$ and symmetric positive definite covariance matrix $\Sigma \succ 0$ if its PDF is:
\begin{equation}
\mathcal{N}(\mathbf{x}; \boldsymbol{\mu}, \Sigma) \triangleq \frac{1}{(2\pi)^{d/2} |\Sigma|^{1/2}} \exp\left( -\frac{1}{2}(\mathbf{x} - \boldsymbol{\mu})^T \Sigma^{-1} (\mathbf{x} - \boldsymbol{\mu}) \right)
\label{eq:mvn_def}
\end{equation}
\end{definitionbox}

\begin{theorem}[Geometric Contour Ellipsoids]
The level sets (contours of equal probability density) of a multivariate Gaussian are concentric ellipsoids characterized by constant Mahalanobis distance:
\begin{equation}
(\mathbf{x} - \boldsymbol{\mu})^T \Sigma^{-1} (\mathbf{x} - \boldsymbol{\mu}) = c^2
\label{eq:mahalanobis_ellipsoid}
\end{equation}
By the Spectral Theorem, $\Sigma = Q \Lambda Q^T$. The principal axes of the ellipsoid align with the orthogonal eigenvectors $\mathbf{q}_i$, with semi-axis lengths proportional to $\sqrt{\lambda_i}$.
\end{theorem}

\subsection{Maximum Likelihood Estimation for Gaussian Parameters}
Given i.i.d. observations $S = \{\mathbf{x}_1, \dots, \mathbf{x}_N\} \subset \mathbb{R}^d$, the log-likelihood function is:
\begin{equation}
\ell(\boldsymbol{\mu}, \Sigma) = -\frac{N d}{2}\log(2\pi) - \frac{N}{2}\log|\Sigma| - \frac{1}{2}\sum_{i=1}^N (\mathbf{x}_i - \boldsymbol{\mu})^T \Sigma^{-1} (\mathbf{x}_i - \boldsymbol{\mu})
\label{eq:mle_loglik}
\end{equation}

\begin{theorem}[Gaussian Maximum Likelihood Estimators]
The Maximum Likelihood Estimators for $\boldsymbol{\mu}$ and $\Sigma$ are:
\begin{align}
\hat{\boldsymbol{\mu}}_{\text{MLE}} &= \frac{1}{N} \sum_{i=1}^N \mathbf{x}_i \label{eq:mle_mu} \\
\hat{\Sigma}_{\text{MLE}} &= \frac{1}{N} \sum_{i=1}^N (\mathbf{x}_i - \hat{\boldsymbol{\mu}}_{\text{MLE}})(\mathbf{x}_i - \hat{\boldsymbol{\mu}}_{\text{MLE}})^T \label{eq:mle_sigma}
\end{align}
While $\hat{\boldsymbol{\mu}}_{\text{MLE}}$ is unbiased ($\E[\hat{\boldsymbol{\mu}}] = \boldsymbol{\mu}$), $\hat{\Sigma}_{\text{MLE}}$ is biased: $\E[\hat{\Sigma}_{\text{MLE}}] = \frac{N-1}{N}\Sigma$. The unbiased sample covariance incorporates Bessel's correction factor $\frac{N}{N-1}$.
\end{theorem}

\subsection{Gaussian Mixture Models (GMM) and the EM Algorithm}
A Gaussian Mixture Model models complex multimodal densities as a convex combination of $K$ distinct Gaussian components:
\begin{equation}
p(\mathbf{x}; \boldsymbol{\theta}) = \sum_{k=1}^K \pi_k \mathcal{N}(\mathbf{x}; \boldsymbol{\mu}_k, \Sigma_k)
\label{eq:gmm_density}
\end{equation}
where $\boldsymbol{\theta} = \{\pi_k, \boldsymbol{\mu}_k, \Sigma_k\}_{k=1}^K$.
Direct maximization of the marginal log-likelihood $\ell(\boldsymbol{\theta}) = \sum_{i=1}^N \log\left( \sum_{k=1}^K \pi_k \mathcal{N}(\mathbf{x}_i; \boldsymbol{\mu}_k, \Sigma_k) \right)$ is analytically intractable because the summation lies inside the logarithm.

\begin{theorem}[The Expectation-Maximization (EM) Algorithm for GMMs]
The EM algorithm alternates between two steps until convergence:
\begin{enumerate}
    \item \textbf{E-step (Expectation):} Compute the posterior probability (responsibility) that sample $\mathbf{x}_i$ was generated by component $k$ using current parameters $\boldsymbol{\theta}^{(t)}$:
    \begin{equation}
    \gamma_{ik} \triangleq P(z_i = k \mid \mathbf{x}_i, \boldsymbol{\theta}^{(t)}) = \frac{\pi_k^{(t)} \mathcal{N}(\mathbf{x}_i; \boldsymbol{\mu}_k^{(t)}, \Sigma_k^{(t)})}{\sum_{j=1}^K \pi_j^{(t)} \mathcal{N}(\mathbf{x}_i; \boldsymbol{\mu}_j^{(t)}, \Sigma_j^{(t)})}
    \label{eq:em_estep}
    \end{equation}
    \item \textbf{M-step (Maximization):} Update parameters to maximize the expected complete-data log-likelihood:
    \begin{align}
    N_k &= \sum_{i=1}^N \gamma_{ik} \label{eq:em_nk} \\
    \pi_k^{(t+1)} &= \frac{N_k}{N} \label{eq:em_pi} \\
    \boldsymbol{\mu}_k^{(t+1)} &= \frac{1}{N_k} \sum_{i=1}^N \gamma_{ik} \mathbf{x}_i \label{eq:em_mu} \\
    \Sigma_k^{(t+1)} &= \frac{1}{N_k} \sum_{i=1}^N \gamma_{ik} (\mathbf{x}_i - \boldsymbol{\mu}_k^{(t+1)})(\mathbf{x}_i - \boldsymbol{\mu}_k^{(t+1)})^T \label{eq:em_sigma}
    \end{align}
\end{enumerate}
\end{theorem}

\subsection{Fundamental Concentration Inequalities}
Concentration inequalities quantify the probability that a random variable deviates from its expected value.

\begin{theorem}[Markov's Inequality]
Let $X$ be a non-negative random variable ($X \ge 0$). For any $a > 0$:
\begin{equation}
P(X \ge a) \le \frac{\E[X]}{a}
\label{eq:markov_ineq}
\end{equation}
\end{theorem}

\begin{proof}
Let $\mathbb{I}(X \ge a)$ be the indicator variable for the event $X \ge a$. Because $X \ge 0$:
\begin{equation}
a \, \mathbb{I}(X \ge a) \le X
\end{equation}
Taking expectations on both sides:
\begin{equation}
a \, \E[\mathbb{I}(X \ge a)] = a P(X \ge a) \le \E[X] \implies P(X \ge a) \le \frac{\E[X]}{a}
\end{equation}
\end{proof}

\begin{theorem}[Chebyshev's Inequality]
Let $X$ have finite mean $\mu = \E[X]$ and variance $\sigma^2 = \Var(X)$. For any $\epsilon > 0$:
\begin{equation}
P(|X - \mu| \ge \epsilon) \le \frac{\sigma^2}{\epsilon^2}
\label{eq:chebyshev_ineq}
\end{equation}
\end{theorem}

\begin{proof}
Apply Markov's inequality to the non-negative random variable $Y = (X - \mu)^2$ with $a = \epsilon^2$:
\begin{equation}
P(|X - \mu| \ge \epsilon) = P((X - \mu)^2 \ge \epsilon^2) \le \frac{\E[(X - \mu)^2]}{\epsilon^2} = \frac{\sigma^2}{\epsilon^2}
\end{equation}
\end{proof}

\begin{theorem}[Hoeffding's Inequality]
Let $X_1, X_2, \dots, X_N$ be independent bounded random variables such that $X_i \in [a_i, b_i]$ with probability 1. Let $\bar{X} = \frac{1}{N}\sum_{i=1}^N X_i$ and $\mu = \E[\bar{X}]$. For any $\epsilon > 0$:
\begin{equation}
P(|\bar{X} - \mu| \ge \epsilon) \le 2 \exp\left( -\frac{2N^2 \epsilon^2}{\sum_{i=1}^N (b_i - a_i)^2} \right)
\label{eq:hoeffding_ineq}
\end{equation}
If all $X_i \in [0, 1]$, this simplifies to:
\begin{equation}
P(|\bar{X} - \mu| \ge \epsilon) \le 2 \exp\left( -2N\epsilon^2 \right)
\end{equation}
\end{theorem}

% ----------------------------------------------------------------------
% PART 3: VISUAL & GEOMETRIC INTUITION
% ----------------------------------------------------------------------
\section{Visual \& Geometric Intuition: The Geometry of Gaussian Mixture Density}

Figure~\ref{fig:gmm_geometry} depicts a two-component Gaussian Mixture Model in $\mathbb{R}^2$, showing the elliptical contour lines determined by the covariance matrices $\Sigma_1, \Sigma_2$ and the soft probabilistic responsibility boundary between clusters.

\begin{figure}[htbp]
\centering
\begin{tikzpicture}[
    scale=0.95,
    font=\sffamily\footnotesize
]

% Component 1: Center (-2.0, 0.5)
\begin{scope}[shift={(-2.0, 0.5)}, rotate=25]
    \draw[thick, color=tealborder!40] (0, 0) ellipse (3.0cm and 1.5cm);
    \draw[thick, color=tealborder!70] (0, 0) ellipse (2.0cm and 1.0cm);
    \draw[thick, color=tealborder] (0, 0) ellipse (1.0cm and 0.5cm);
    \filldraw[color=tealborder] (0, 0) circle (2.5pt) node[below] {$\boldsymbol{\mu}_1$};
\end{scope}
\node[above, color=tealborder] at (-2.0, 2.3) {\textbf{Cluster 1: } $\mathcal{N}(\boldsymbol{\mu}_1, \Sigma_1)$};

% Component 2: Center (2.5, -0.5)
\begin{scope}[shift={(2.5, -0.5)}, rotate=-35]
    \draw[thick, color=coralborder!40] (0, 0) ellipse (2.8cm and 1.3cm);
    \draw[thick, color=coralborder!70] (0, 0) ellipse (1.8cm and 0.8cm);
    \draw[thick, color=coralborder] (0, 0) ellipse (0.9cm and 0.4cm);
    \filldraw[color=coralborder] (0, 0) circle (2.5pt) node[below] {$\boldsymbol{\mu}_2$};
\end{scope}
\node[above, color=coralborder] at (2.5, 1.2) {\textbf{Cluster 2: } $\mathcal{N}(\boldsymbol{\mu}_2, \Sigma_2)$};

% Soft decision boundary (gamma_1 = gamma_2 = 0.5)
\draw[dashed, line width=1.5pt, color=primarydark] (0.2, 3.0) -- (0.5, -3.0);
\node[below, color=primarydark] at (0.5, -3.0) {Equi-Responsibility: $\gamma_{i1} = \gamma_{i2} = 0.5$};

% Intermediate sample point with soft assignment
\coordinate (X) at (0.0, 0.8);
\filldraw[color=purpleex] (X) circle (3pt) node[above right] {$\mathbf{x}_i$};
\node[draw=purpleborder, fill=purplebg, rounded corners=2pt, align=center] at (-0.2, -1.8) {
    \textbf{Soft Responsibility:}\\$\gamma_{i1} = 0.68, \; \gamma_{i2} = 0.32$
};

\end{tikzpicture}
\caption{The Geometry of Gaussian Mixture Models: Two distinct Gaussian clusters in $\mathbb{R}^2$ with overlapping contours, generating continuous soft posterior responsibilities $\gamma_{ik}$.}
\label{fig:gmm_geometry}
\end{figure}

% ----------------------------------------------------------------------
% PART 4: DEEP INTUITION & CONCEPTUAL DERIVATIONS
% ----------------------------------------------------------------------
\section{Deep Intuition \& Conceptual Derivations}

\subsection{Derivation of the EM Algorithm via the Evidence Lower Bound (ELBO)}
Let $\mathbf{x}$ be observed data and $\mathbf{z}$ be latent variables. The marginal log-likelihood is $\log p(\mathbf{x}; \boldsymbol{\theta})$.
Let $q(\mathbf{z})$ be an arbitrary probability distribution over the latent variables $\mathbf{z}$.
Using Jensen's inequality on the strictly concave logarithm function:
\begin{align}
\log p(\mathbf{x}; \boldsymbol{\theta}) &= \log \sum_{\mathbf{z}} p(\mathbf{x}, \mathbf{z}; \boldsymbol{\theta}) = \log \sum_{\mathbf{z}} q(\mathbf{z}) \frac{p(\mathbf{x}, \mathbf{z}; \boldsymbol{\theta})}{q(\mathbf{z})} \nonumber \\
&= \log \E_{q} \left[ \frac{p(\mathbf{x}, \mathbf{z}; \boldsymbol{\theta})}{q(\mathbf{z})} \right] \ge \E_{q} \left[ \log \frac{p(\mathbf{x}, \mathbf{z}; \boldsymbol{\theta})}{q(\mathbf{z})} \right] \triangleq \text{ELBO}(q, \boldsymbol{\theta})
\label{eq:elbo_jensen}
\end{align}
Expanding the Evidence Lower Bound:
\begin{align}
\text{ELBO}(q, \boldsymbol{\theta}) &= \sum_{\mathbf{z}} q(\mathbf{z}) \log p(\mathbf{x}, \mathbf{z}; \boldsymbol{\theta}) - \sum_{\mathbf{z}} q(\mathbf{z}) \log q(\mathbf{z}) \nonumber \\
&= \sum_{\mathbf{z}} q(\mathbf{z}) \log \left( p(\mathbf{z} \mid \mathbf{x}; \boldsymbol{\theta}) p(\mathbf{x}; \boldsymbol{\theta}) \right) + \mathcal{H}(q) \nonumber \\
&= \log p(\mathbf{x}; \boldsymbol{\theta}) - D_{\text{KL}}\left( q(\mathbf{z}) \,\big\|\, p(\mathbf{z} \mid \mathbf{x}; \boldsymbol{\theta}) \right)
\end{align}
This decomposition reveals the exact mechanics of the EM algorithm:
\begin{itemize}
    \item \textbf{E-step:} Maximize $\text{ELBO}(q, \boldsymbol{\theta}^{(t)})$ with respect to $q(\mathbf{z})$ while holding parameters $\boldsymbol{\theta}^{(t)}$ fixed. The ELBO touches the true marginal log-likelihood when the KL divergence vanishes, which occurs if and only if:
    \begin{equation}
    q^*(\mathbf{z}) = p(\mathbf{z} \mid \mathbf{x}; \boldsymbol{\theta}^{(t)})
    \end{equation}
    This is precisely the posterior responsibility $\gamma_{ik}$ computed in the E-step!
    \item \textbf{M-step:} Maximize $\text{ELBO}(q^*, \boldsymbol{\theta})$ with respect to parameters $\boldsymbol{\theta}$ while holding $q^*$ fixed.
    Because $\text{ELBO}(q^*, \boldsymbol{\theta}^{(t+1)}) \ge \text{ELBO}(q^*, \boldsymbol{\theta}^{(t)}) = \log p(\mathbf{x}; \boldsymbol{\theta}^{(t)})$, this guarantees that the true marginal log-likelihood increases monotonically at every iteration:
    \begin{equation}
    \log p(\mathbf{x}; \boldsymbol{\theta}^{(t+1)}) \ge \log p(\mathbf{x}; \boldsymbol{\theta}^{(t)})
    \end{equation}
\end{itemize}

\subsection{Connection Between GMM and $K$-Means Clustering}
$K$-Means clustering is mathematically equivalent to a special degenerate limit of a Gaussian Mixture Model:
\begin{enumerate}
    \item Assume all $K$ covariance matrices are isotropic and identical: $\Sigma_k = \sigma^2 I_d$.
    \item Assume uniform mixture priors: $\pi_k = \frac{1}{K}$.
    \item In the limit as variance $\sigma^2 \to 0$ (zero temperature limit):
    \begin{equation}
    \lim_{\sigma^2 \to 0} \gamma_{ik} = \lim_{\sigma^2 \to 0} \frac{\exp\left( -\frac{\|\mathbf{x}_i - \boldsymbol{\mu}_k\|^2}{2\sigma^2} \right)}{\sum_{j=1}^K \exp\left( -\frac{\|\mathbf{x}_i - \boldsymbol{\mu}_j\|^2}{2\sigma^2} \right)} = \begin{cases} 1 & \text{if } k = \argmin_j \|\mathbf{x}_i - \boldsymbol{\mu}_j\|^2 \\ 0 & \text{otherwise} \end{cases}
    \end{equation}
\end{enumerate}
Soft responsibilities collapse into hard binary cluster assignments, and the M-step mean update becomes the exact cluster centroid calculation of $K$-Means.

% ----------------------------------------------------------------------
% PART 5: EXPLICIT GOTCHAS & DIAGNOSTIC MISCONCEPTIONS
% ----------------------------------------------------------------------
\section{Explicit Gotchas \& Diagnostic Misconceptions}

\begin{trapbox}{Gotcha 1: Covariance Collapse / Likelihood Singularity in GMMs}
During EM training, a Gaussian component's mean $\boldsymbol{\mu}_k$ lands exactly on an isolated sample $\mathbf{x}_i$. The solver attempts to update $\Sigma_k$, causing $\det(\Sigma_k) \to 0$ and likelihood to explode to $+\infty$.
\textbf{The Trap:} The maximum likelihood formulation for GMMs has **unbounded singularities** on the boundary of parameter space. If a component collapses onto a single point, $\sigma^2 \to 0 \implies \frac{1}{\sigma^d} \to \infty$. Practical implementations must regularize the covariance matrices by adding a small positive ridge diagonal: $\Sigma_k \leftarrow \Sigma_k + \epsilon I_d$ (where $\epsilon \approx 10^{-6}$).
\end{trapbox}

\begin{trapbox}{Gotcha 2: Believing the EM Algorithm Finds the Global Maximum}
An analyst runs EM once on a complex dataset and assumes the parameters represent the global maximum of the marginal log-likelihood.
\textbf{The Trap:} The EM algorithm is a **local hill-climbing procedure**. The mixture log-likelihood surface is non-convex and filled with numerous suboptimal local maxima and saddle points. GMM implementations must be run multiple times with different random restarts (often initialized using $K$-Means centroids) to select the solution with the highest converged log-likelihood.
\end{trapbox}

\begin{trapbox}{Gotcha 3: Applying Markov's Inequality to Variables with Negative Support}
A student applies Markov's inequality to standard normal variable $Z \sim \mathcal{N}(0, 1)$ to bound $P(Z \ge 2) \le \frac{\E[Z]}{2} = 0$.
\textbf{The Trap:} Markov's inequality requires the random variable to be strictly **non-negative** ($X \ge 0$). Applying it to variables that take negative values produces nonsensical probabilities (such as upper-bounding a non-zero probability by zero or a negative number).
\end{trapbox}

\begin{trapbox}{Gotcha 4: Conflating Sample Variance with Variance of the Sample Mean}
A practitioner samples $N = 100$ points from a distribution with variance $\sigma^2 = 16$. They assert that the sample mean $\bar{X}$ has standard deviation $4$.
\textbf{The Trap:} The individual samples have variance $\sigma^2 = 16$ (standard deviation $4$). However, by independence, the **sample mean** has variance $\Var(\bar{X}) = \frac{\sigma^2}{N} = \frac{16}{100} = 0.16$, yielding a standard error of $\frac{\sigma}{\sqrt{N}} = \frac{4}{10} = 0.4$. The sample mean concentrates around $\mu$ at rate $\mathcal{O}(1/\sqrt{N})$.
\end{trapbox}

% ----------------------------------------------------------------------
% PART 6: COMPARATIVE SYNTHESIS & SUMMARY TABLE
% ----------------------------------------------------------------------
\section{Comparative Synthesis \& Summary Table}

Table~\ref{tab:clustering_comp} compares the characteristics of clustering and latent density models.

\begin{table}[htbp]
\centering
\small
\begin{tabularx}{\textwidth}{lXXXXX}
\toprule
\textbf{Model} & \textbf{Cluster Assignment} & \textbf{Cluster Geometry} & \textbf{Objective Function} & \textbf{Computational Complexity} & \textbf{Density Output} \\
\midrule
$K$-Means & Hard (Binary $0/1$) & Spherical (isotropic Voronoi cells) & Residual sum of squares $\sum \|\mathbf{x}_i - \boldsymbol{\mu}_{z_i}\|^2$ & $\mathcal{O}(N K d \cdot \text{iters})$ & No (Centroids only) \\
\addlinespace
GMM (EM) & Soft (Probabilities $\gamma_{ik} \in [0, 1]$) & Ellipsoidal (arbitrary full covariance $\Sigma_k$) & Marginal log-likelihood $\sum \log \sum \pi_k \mathcal{N}(\mathbf{x}_i)$ & $\mathcal{O}(N K d^2 \cdot \text{iters})$ & Yes (Complete generative density $p(\mathbf{x})$) \\
\addlinespace
Factor Analysis & Continuous latent $\mathbf{z} \in \mathbb{R}^k$ & Linear manifold + diagonal noise & Marginal Gaussian log-likelihood & $\mathcal{O}(N d k \cdot \text{iters})$ & Yes (Low-rank Gaussian density) \\
\addlinespace
KDE & Non-parametric & Kernel balls centered on every data point & Integrated squared error / leave-one-out likelihood & $\mathcal{O}(N^2 d)$ at inference & Yes (Continuous empirical density) \\
\bottomrule
\end{tabularx}
\caption{Comparative Synthesis: Clustering and Density Estimation Paradigms.}
\label{tab:clustering_comp}
\end{table}

% ----------------------------------------------------------------------
% PART 7: PYTHON / SCIENTIFIC COMPUTING LAB
% ----------------------------------------------------------------------
\section{Python / Scientific Computing Lab}

The following production script implements a complete Gaussian Mixture Model with the Expectation-Maximization algorithm from scratch, verifying the monotonic ascent of the marginal log-likelihood.

\begin{lstlisting}[language=Python, caption={Gaussian Mixture Model EM Algorithm Implementation in Python.}]
import numpy as np

class GaussianMixtureEM:
    def __init__(self, n_components: int, max_iter: int = 100, tol: float = 1e-4, seed: int = 42):
        self.K = n_components
        self.max_iter = max_iter
        self.tol = tol
        self.seed = seed
        self.weights_ = None
        self.means_ = None
        self.covariances_ = None
        self.log_likelihood_history_ = []

    def _multivariate_gaussian_log_pdf(self, X: np.ndarray, mean: np.ndarray, cov: np.ndarray) -> np.ndarray:
        N, d = X.shape
        # Regularize covariance to prevent collapse
        cov_reg = cov + 1e-6 * np.eye(d)
        diff = X - mean
        sign, logdet = np.linalg.slogdet(cov_reg)
        inv_cov = np.linalg.inv(cov_reg)
        mahalanobis = np.sum(diff @ inv_cov * diff, axis=1)
        log_prob = -0.5 * (d * np.log(2.0 * np.pi) + logdet + mahalanobis)
        return log_prob

    def fit(self, X: np.ndarray):
        N, d = X.shape
        rng = np.random.RandomState(self.seed)

        # Initialize parameters: random points as initial means
        self.weights_ = np.ones(self.K) / self.K
        random_idx = rng.choice(N, self.K, replace=False)
        self.means_ = X[random_idx].copy()
        self.covariances_ = np.array([np.cov(X, rowvar=False) for _ in range(self.K)])

        for it in range(self.max_iter):
            # E-step: compute log responsibilities via log-sum-exp
            log_weighted_probs = np.zeros((N, self.K))
            for k in range(self.K):
                log_weighted_probs[:, k] = np.log(self.weights_[k] + 1e-12) + \
                                          self._multivariate_gaussian_log_pdf(X, self.means_[k], self.covariances_[k])

            # Total log-likelihood
            max_log = np.max(log_weighted_probs, axis=1, keepdims=True)
            log_sum_exp = max_log + np.log(np.sum(np.exp(log_weighted_probs - max_log), axis=1, keepdims=True))
            current_log_lik = np.sum(log_sum_exp)
            self.log_likelihood_history_.append(current_log_lik)

            # Normalized responsibilities: gamma_ik = exp(log_weighted - log_sum_exp)
            gamma = np.exp(log_weighted_probs - log_sum_exp)

            # M-step: parameter updates
            N_k = np.sum(gamma, axis=0) + 1e-12
            self.weights_ = N_k / N

            for k in range(self.K):
                self.means_[k] = np.sum(gamma[:, k:k+1] * X, axis=0) / N_k[k]
                diff = X - self.means_[k]
                self.covariances_[k] = (diff.T @ (gamma[:, k:k+1] * diff)) / N_k[k]

            # Check convergence
            if it > 0 and abs(self.log_likelihood_history_[-1] - self.log_likelihood_history_[-2]) < self.tol:
                break

        return self

def main():
    rng = np.random.RandomState(101)
    N = 300
    # Generate 2-component synthetic data
    X1 = rng.randn(N // 2, 2) + np.array([3.0, 3.0])
    X2 = rng.randn(N // 2, 2) + np.array([-2.0, -1.0])
    X = np.vstack([X1, X2])

    gmm = GaussianMixtureEM(n_components=2, max_iter=50)
    gmm.fit(X)

    print("Gaussian Mixture Model EM Execution Complete:")
    print(f"Number of Iterations: {len(gmm.log_likelihood_history_)}")
    print(f"Initial Log-Likelihood: {gmm.log_likelihood_history_[0]:.2f}")
    print(f"Final Log-Likelihood:   {gmm.log_likelihood_history_[-1]:.2f}")

    # Verify monotonic ascent
    diffs = np.diff(gmm.log_likelihood_history_)
    is_monotonic = np.all(diffs >= -1e-6)
    print(f"Log-Likelihood Monotonically Ascending: {is_monotonic}")
    print("Estimated Component Means:")
    for k, mean in enumerate(gmm.means_):
        print(f"  Cluster {k+1}: [{mean[0]:.3f}, {mean[1]:.3f}], Weight: {gmm.weights_[k]:.3f}")

if __name__ == "__main__":
    main()
\end{lstlisting}

% ----------------------------------------------------------------------
% PART 8: COMPREHENSIVE PRACTICE PROBLEMS
% ----------------------------------------------------------------------
\section{Comprehensive Practice Problems}

\subsection*{Part A: Analytical \& Mathematical Problems}

\begin{enumerate}[label=\textbf{P\arabic*.}]
    \item \textbf{Analytical 2D Mahalanobis Distance and Ellipsoidal Contours:}
    Let $\mathbf{X} \sim \mathcal{N}(\boldsymbol{\mu}, \Sigma)$ with $\boldsymbol{\mu} = \begin{bmatrix} 1 \\ 2 \end{bmatrix}$ and $\Sigma = \begin{bmatrix} 4 & 2 \\ 2 & 2 \end{bmatrix}$.
    \begin{enumerate}[label=(\alph*)]
        \item Compute the determinant $|\Sigma|$ and inverse matrix $\Sigma^{-1}$.
        \item Compute the Mahalanobis distance $D_M(\mathbf{x}, \boldsymbol{\mu})$ for observation $\mathbf{x} = [3, 4]^T$.
        \item Find the eigenvalues and eigenvectors of $\Sigma$ and describe the principal orientation and axis lengths of the $1$-standard-deviation contour ellipsoid.
    \end{enumerate}

    \item \textbf{Unbiasedness and Degrees of Freedom of Sample Covariance:}
    Let $\mathbf{X}_1, \dots, \mathbf{X}_N \stackrel{\text{i.i.d.}}{\sim} \mathcal{N}(\boldsymbol{\mu}, \Sigma)$.
    Prove that the Maximum Likelihood Estimator $\hat{\Sigma}_{\text{MLE}} = \frac{1}{N}\sum_{i=1}^N (\mathbf{X}_i - \bar{\mathbf{X}})(\mathbf{X}_i - \bar{\mathbf{X}})^T$ has expectation:
    \begin{equation}
    \E[\hat{\Sigma}_{\text{MLE}}] = \frac{N - 1}{N} \Sigma
    \end{equation}

    \item \textbf{Analytical E-step Calculation for a 1D GMM:}
    Consider a 1D GMM with $K = 2$ components: $\pi_1 = 0.6, \mu_1 = 0, \sigma_1^2 = 1$ and $\pi_2 = 0.4, \mu_2 = 3, \sigma_2^2 = 1$.
    For an observation $x = 1.5$:
    \begin{enumerate}[label=(\alph*)]
        \item Evaluate the Gaussian densities $\mathcal{N}(1.5; 0, 1)$ and $\mathcal{N}(1.5; 3, 1)$.
        \item Compute the posterior responsibilities $\gamma_1$ and $\gamma_2$.
        \item State which cluster $x = 1.5$ is more likely to belong to, and explain why despite being equidistant from both means.
    \end{enumerate}

    \item \textbf{Sample Complexity via Hoeffding's Inequality:}
    A pollster wants to estimate the true proportion $p \in [0, 1]$ of voters favoring a policy using empirical average $\bar{X}_N$.
    Using Hoeffding's inequality, calculate the minimum number of independent samples $N$ required to guarantee that $|\bar{X}_N - p| \le 0.02$ with confidence at least $95\%$ ($\delta = 0.05$).

    \item \textbf{Chebyshev vs. Markov Bound Comparison:}
    Let $X$ be the number of heads in $N = 100$ independent flips of a fair coin.
    \begin{enumerate}[label=(\alph*)]
        \item Find the exact mean $\mu = \E[X]$ and variance $\sigma^2 = \Var(X)$.
        \item Use Markov's inequality to bound $P(X \ge 75)$.
        \item Use Chebyshev's inequality to bound $P(X \ge 75) \le P(|X - 50| \ge 25)$.
        \item Compare the bounds and explain why Chebyshev is dramatically tighter.
    \end{enumerate}
\end{enumerate}

\subsection*{Part B: Diagnostic Multiple-Choice Questions}

\begin{enumerate}[label=\textbf{MCQ\arabic*.}]
    \item What is guaranteed regarding the marginal log-likelihood $\sum_{i=1}^N \log p(\mathbf{x}_i; \boldsymbol{\theta})$ at each successive iteration of the Expectation-Maximization (EM) algorithm?
    \begin{enumerate}[label=(\alph*)]
        \item It is guaranteed to decrease monotonically toward zero.
        \item It is guaranteed to increase monotonically (or remain constant) at every iteration: $\ell(\boldsymbol{\theta}^{(t+1)}) \ge \ell(\boldsymbol{\theta}^{(t)})$.
        \item It is guaranteed to converge to the global maximum of the parameter space.
        \item It evaluates to a strictly negative value if and only if the data is linearly separable.
    \end{enumerate}

    \item If a bounded random variable $X_i \in [0, 1]$ has unknown distribution with mean $\mu$, which concentration inequality provides an exponential bound in sample size $N$ on the deviation of sample mean $|\bar{X}_N - \mu|$?
    \begin{enumerate}[label=(\alph*)]
        \item Markov's Inequality
        \item Chebyshev's Inequality
        \item Hoeffding's Inequality
        \item Cauchy-Schwarz Inequality
    \end{enumerate}
\end{enumerate}



\backmatter
\chapter{Solutions to Practice Problems \& Diagnostic Rationales}
\label{chap:appendix}

This appendix provides exhaustive, step-by-step mathematical solutions to all sixty Part A analytical problems and deep diagnostic rationales for all twenty-four Part B multiple-choice questions across Chapters 1 through 12.

% ======================================================================
% CHAPTER 1 SOLUTIONS
% ======================================================================
\section*{Chapter 1: Foundations of Machine Learning \& Problem Formulation}

\subsection*{Part A: Analytical Solutions}

\paragraph{P1. Bayes Optimal Predictor Under Absolute Loss (MAE):}
We wish to find $c^* = h^*(\mathbf{x}) = \argmin_{c \in \mathbb{R}} \E[|c - Y| \mid \mathbf{X} = \mathbf{x}]$.
Let $f(y)$ be the conditional probability density function of $Y$ given $\mathbf{X} = \mathbf{x}$. The conditional expected loss is:
\begin{align*}
g(c) &= \int_{-\infty}^\infty |c - y| f(y) \, dy = \int_{-\infty}^c (c - y) f(y) \, dy + \int_c^\infty (y - c) f(y) \, dy \\
&= c \int_{-\infty}^c f(y) \, dy - \int_{-\infty}^c y f(y) \, dy + \int_c^\infty y f(y) \, dy - c \int_c^\infty f(y) \, dy
\end{align*}
Applying Leibniz's rule to differentiate $g(c)$ with respect to $c$:
\begin{align*}
\frac{dg}{dc} &= \int_{-\infty}^c f(y) \, dy + c f(c) - c f(c) - \int_c^\infty f(y) \, dy - c(-f(c)) - c f(c) \\
&= P(Y \le c \mid \mathbf{x}) - P(Y > c \mid \mathbf{x}) = P(Y \le c \mid \mathbf{x}) - (1 - P(Y \le c \mid \mathbf{x})) \\
&= 2 P(Y \le c \mid \mathbf{x}) - 1
\end{align*}
Setting the derivative to zero:
\begin{equation*}
2 P(Y \le c^* \mid \mathbf{x}) - 1 = 0 \implies P(Y \le c^* \mid \mathbf{X} = \mathbf{x}) = \frac{1}{2}
\end{equation*}
By definition, the value $c^*$ satisfying $P(Y \le c^* \mid \mathbf{x}) = 0.5$ is the conditional median $\text{Median}(Y \mid \mathbf{X} = \mathbf{x})$.

\paragraph{P2. Asymmetric Classification Costs:}
Let $\hat{y} \in \{0, 1\}$ be the predicted label and $\eta(\mathbf{x}) = P(Y = 1 \mid \mathbf{x})$. The cost matrix specifies:
cost $C_{01} = 50$ when true $Y = 1$ is predicted as $\hat{y} = 0$, and cost $C_{10} = 1$ when true $Y = 0$ is predicted as $\hat{y} = 1$. Correct predictions incur cost $0$.
The expected cost for predicting $\hat{y} = 1$ is:
\begin{equation*}
\mathcal{C}(\hat{y} = 1 \mid \mathbf{x}) = C_{10} P(Y = 0 \mid \mathbf{x}) = C_{10} (1 - \eta(\mathbf{x})) = 1 \cdot (1 - \eta(\mathbf{x}))
\end{equation*}
The expected cost for predicting $\hat{y} = 0$ is:
\begin{equation*}
\mathcal{C}(\hat{y} = 0 \mid \mathbf{x}) = C_{01} P(Y = 1 \mid \mathbf{x}) = C_{01} \eta(\mathbf{x}) = 50 \eta(\mathbf{x})
\end{equation*}
Predicting $\hat{y} = 1$ is optimal if and only if:
\begin{equation*}
\mathcal{C}(\hat{y} = 1 \mid \mathbf{x}) \le \mathcal{C}(\hat{y} = 0 \mid \mathbf{x}) \iff 1 - \eta(\mathbf{x}) \le 50 \eta(\mathbf{x}) \iff 51 \eta(\mathbf{x}) \ge 1 \iff \eta(\mathbf{x}) \ge \frac{1}{51} \approx 0.0196
\end{equation*}
Thus, the optimal threshold is $\tau^* = \frac{C_{10}}{C_{10} + C_{01}} = \frac{1}{51} \approx 1.96\%$.

\paragraph{P3. Empirical Risk Under Label Flipping Noise:}
Let $\tilde{Y}$ be the noisy label with $P(\tilde{Y} = -Y) = \rho$ and $P(\tilde{Y} = Y) = 1 - \rho$.
For any hypothesis $h(\mathbf{X}) \in \{-1, +1\}$:
\begin{align*}
R_{\text{noisy}}(h) &= P(h(\mathbf{X}) \ne \tilde{Y}) \\
&= P(h(\mathbf{X}) \ne \tilde{Y} \mid \tilde{Y} = Y) P(\tilde{Y} = Y) + P(h(\mathbf{X}) \ne \tilde{Y} \mid \tilde{Y} = -Y) P(\tilde{Y} = -Y) \\
&= P(h(\mathbf{X}) \ne Y)(1 - \rho) + P(h(\mathbf{X}) = Y)\rho
\end{align*}
Substitute $P(h(\mathbf{X}) \ne Y) = R(h)$ and $P(h(\mathbf{X}) = Y) = 1 - R(h)$:
\begin{equation*}
R_{\text{noisy}}(h) = R(h)(1 - \rho) + (1 - R(h))\rho = R(h)(1 - 2\rho) + \rho
\end{equation*}
Since $\rho < 1/2$, $1 - 2\rho > 0$, meaning the ranking of hypotheses is preserved: minimizing noisy risk is equivalent to minimizing clean risk.

\paragraph{P4. Curse of Dimensionality \& Distance Concentration:}
For two independent uniform points $\mathbf{X}, \mathbf{Y} \sim \mathcal{U}[0, 1]^d$, the squared Euclidean distance is $D^2 = \sum_{j=1}^d (X_j - Y_j)^2$.
Each coordinate difference $Z_j = (X_j - Y_j)^2$ is i.i.d. with mean $\mu_Z = \int_0^1 \int_0^1 (x - y)^2 dx dy = \frac{1}{6}$ and variance $\sigma_Z^2 = \frac{7}{180}$.
By the Central Limit Theorem:
\begin{equation*}
\frac{D^2}{d} = \frac{1}{d} \sum_{j=1}^d Z_j \xrightarrow{P} \frac{1}{6} \implies D \approx \sqrt{\frac{d}{6}}
\end{equation*}
The standard deviation of $D^2$ scales as $\mathcal{O}(\sqrt{d})$, while its mean scales as $\mathcal{O}(d)$. Consequently, the relative spread behaves as:
\begin{equation*}
\frac{D_{\max} - D_{\min}}{D_{\min}} \le \frac{\mathcal{O}(\sqrt{d})}{\mathcal{O}(\sqrt{d}) \cdot \sqrt{d}} = \mathcal{O}\left(\frac{1}{\sqrt{d}}\right) \xrightarrow{d \to \infty} 0
\end{equation*}
In high dimensions, all pairwise distances become identical, destroying the discriminative power of nearest-neighbor algorithms.

\paragraph{P5. Generalization Bound for Finite Hypothesis Classes:}
For any fixed hypothesis $h \in \mathcal{H}$, the empirical risk $R_{\text{emp}}(h) = \frac{1}{N}\sum_{i=1}^N \ell(h(\mathbf{x}_i), y_i)$ is an average of $N$ independent random variables in $[0, 1]$ with expectation $\E[R_{\text{emp}}(h)] = R(h)$.
By Hoeffding's inequality, for any $\epsilon > 0$:
\begin{equation*}
P\left( R(h) - R_{\text{emp}}(h) > \epsilon \right) \le \exp\left( -2N\epsilon^2 \right)
\end{equation*}
Applying the Union Bound over all $|\mathcal{H}|$ hypotheses:
\begin{equation*}
P\left( \exists h \in \mathcal{H}: R(h) - R_{\text{emp}}(h) > \epsilon \right) \le \sum_{h \in \mathcal{H}} P\left( R(h) - R_{\text{emp}}(h) > \epsilon \right) \le |\mathcal{H}| \exp\left( -2N\epsilon^2 \right)
\end{equation*}
Equating the failure probability bound to $\delta$:
\begin{equation*}
|\mathcal{H}| \exp\left( -2N\epsilon^2 \right) = \delta \implies -2N\epsilon^2 = \log\left(\frac{\delta}{|\mathcal{H}|}\right) \implies \epsilon = \sqrt{\frac{\log|\mathcal{H}| + \log(1/\delta)}{2N}}
\end{equation*}
With probability at least $1 - \delta$, for every $h \in \mathcal{H}$, $R(h) \le R_{\text{emp}}(h) + \sqrt{\frac{\log|\mathcal{H}| + \log(1/\delta)}{2N}}$.

\subsection*{Part B: Diagnostic Multiple-Choice Rationales}

\paragraph{MCQ1: Class Imbalance Accuracy Paradox}
\begin{itemize}
    \item \textbf{(a) Incorrect.} High global accuracy provides no information on minority class recall; predicting the majority class exclusively achieves $98\%$ accuracy while failing completely on class $1$.
    \item \textbf{(b) Correct.} On a dataset with $98\%$ class $0$, the trivial constant predictor $h(\mathbf{x}) = 0$ misclassifies zero samples of class $0$ and all samples of class $1$, producing an overall empirical accuracy of exactly $98\%$.
    \item \textbf{(c) Incorrect.} Minimizing unweighted $0-1$ loss on imbalanced data directly disincentivizes predicting the minority class.
    \item \textbf{(d) Incorrect.} The empirical risk under cross-entropy loss would only equal zero if every predicted probability matched the one-hot labels with certainty ($1.0$), which is false.
\end{itemize}

\paragraph{MCQ2: Downward Optimistic Bias of ERM}
\begin{itemize}
    \item \textbf{(a) Correct.} The hypothesis $\hat{h}_S$ was chosen specifically to minimize the empirical loss on sample $S$. Because $\hat{h}_S$ is conditioned on $S$, the expected training error $\E_S[R_{\text{emp}}(\hat{h}_S; S)] \le \min_h \E_S[R_{\text{emp}}(h; S)] = R(h^*) \le R(\hat{h}_S)$, making it systematically optimistic.
    \item \textbf{(b) Incorrect.} For any \emph{fixed} hypothesis $h$, the sample average is an unbiased estimator: $\E[R_{\text{emp}}(h; S)] = R(h)$ by the Law of Large Numbers. Bias arises solely because $\hat{h}_S$ is optimized over $S$.
    \item \textbf{(c) Incorrect.} The loss function does not depend concavely on sample size $N$.
    \item \textbf{(d) Incorrect.} Jensen's inequality states that for convex functions $f(\E[X]) \le \E[f(X)]$; it does not govern empirical risk optimization bias.
\end{itemize}

% ======================================================================
% CHAPTER 2 SOLUTIONS
% ======================================================================
\section*{Chapter 2: Univariate \& Multivariate Differential Calculus for Learning}

\subsection*{Part A: Analytical Solutions}

\paragraph{P1. Gradient of Regularized Ridge Objective:}
\begin{enumerate}[label=(\alph*)]
    \item Expanding the objective:
    \begin{equation*}
    \mathcal{L}(\mathbf{w}) = \frac{1}{2}(\mathbf{y} - X\mathbf{w})^T (\mathbf{y} - X\mathbf{w}) + \frac{\lambda}{2}\mathbf{w}^T \mathbf{w} = \frac{1}{2}\mathbf{y}^T \mathbf{y} - \mathbf{w}^T X^T \mathbf{y} + \frac{1}{2}\mathbf{w}^T X^T X \mathbf{w} + \frac{\lambda}{2}\mathbf{w}^T \mathbf{w}
    \end{equation*}
    Applying the matrix calculus identities derived in Chapter 2:
    \begin{equation*}
    \nabla_{\mathbf{w}} \mathcal{L}(\mathbf{w}) = -X^T \mathbf{y} + X^T X \mathbf{w} + \lambda \mathbf{w} = X^T (X\mathbf{w} - \mathbf{y}) + \lambda \mathbf{w}
    \end{equation*}
    \item Taking the gradient of $\nabla \mathcal{L}(\mathbf{w})$ with respect to $\mathbf{w}$:
    \begin{equation*}
    H(\mathbf{w}) = \nabla^2_{\mathbf{w}} \mathcal{L}(\mathbf{w}) = X^T X + \lambda I_d
    \end{equation*}
    \item For any non-zero $\mathbf{v} \in \mathbb{R}^d \setminus \{\mathbf{0}\}$:
    \begin{equation*}
    \mathbf{v}^T H \mathbf{v} = \mathbf{v}^T (X^T X + \lambda I)\mathbf{v} = (X\mathbf{v})^T (X\mathbf{v}) + \lambda \mathbf{v}^T \mathbf{v} = \|X\mathbf{v}\|_2^2 + \lambda \|\mathbf{v}\|_2^2
    \end{equation*}
    Since $\|X\mathbf{v}\|_2^2 \ge 0$ and $\lambda \|\mathbf{v}\|_2^2 > 0$ (for $\lambda > 0$ and $\mathbf{v} \ne \mathbf{0}$), $\mathbf{v}^T H \mathbf{v} > 0$, proving $H$ is strictly positive definite.
\end{enumerate}

\paragraph{P2. Jacobian of the Softmax Transformation:}
Recall $\sigma_i(\mathbf{z}) = \frac{\exp(z_i)}{\sum_{k=1}^K \exp(z_k)} = \frac{e^{z_i}}{S}$, where $S = \sum_{k=1}^K e^{z_k}$.
\begin{itemize}
    \item Case 1: $i = j$.
    \begin{equation*}
    \frac{\partial \sigma_i}{\partial z_i} = \frac{e^{z_i} S - e^{z_i} \frac{\partial S}{\partial z_i}}{S^2} = \frac{e^{z_i} S - e^{z_i} e^{z_i}}{S^2} = \frac{e^{z_i}}{S} \left( 1 - \frac{e^{z_i}}{S} \right) = \sigma_i (1 - \sigma_i)
    \end{equation*}
    \item Case 2: $i \ne j$.
    \begin{equation*}
    \frac{\partial \sigma_i}{\partial z_j} = \frac{0 \cdot S - e^{z_i} \frac{\partial S}{\partial z_j}}{S^2} = -\frac{e^{z_i} e^{z_j}}{S^2} = -\sigma_i \sigma_j
    \end{equation*}
\end{itemize}
Combining via the Kronecker delta $\delta_{ij}$: $\frac{\partial \sigma_i}{\partial z_j} = \sigma_i (\delta_{ij} - \sigma_j)$.

\paragraph{P3. Directional Derivative Along Non-Gradient Paths:}
$f(x_1, x_2) = 3x_1^2 + 2x_1 x_2 + x_2^2$.
\begin{enumerate}[label=(\alph*)]
    \item $\nabla f(x_1, x_2) = [6x_1 + 2x_2, 2x_1 + 2x_2]^T$.
    At $(1, 2)$: $\nabla f(1, 2) = [6(1) + 2(2), 2(1) + 2(2)]^T = [10, 6]^T$.
    \item Direction toward $(4, 6)$: displacement vector $\mathbf{d} = (4 - 1, 6 - 2) = (3, 4)$.
    Unit direction vector: $\mathbf{u} = \frac{\mathbf{d}}{\|\mathbf{d}\|} = \frac{(3, 4)}{5} = [0.6, 0.8]^T$.
    Rate of change: $D_{\mathbf{u}} f(1, 2) = \nabla f(1, 2)^T \mathbf{u} = 10(0.6) + 6(0.8) = 6 + 4.8 = 10.8$.
    \item Direction of zero change: $\nabla f(1, 2)^T \mathbf{v} = 0 \implies 10 v_1 + 6 v_2 = 0 \implies v_2 = -\frac{5}{3}v_1$.
    Normalizing $\|\mathbf{v}\|_2 = 1 \implies v_1^2 + \frac{25}{9}v_1^2 = 1 \implies \frac{34}{9}v_1^2 = 1 \implies v_1 = \frac{3}{\sqrt{34}}, v_2 = -\frac{5}{\sqrt{34}}$.
    Thus $\mathbf{v} = \pm \frac{1}{\sqrt{34}}[3, -5]^T$.
\end{enumerate}

\paragraph{P4. Second-Order Taylor Approximation Error Bound:}
By Taylor's theorem with integral remainder:
\begin{equation*}
f(\mathbf{x} + \mathbf{v}) = f(\mathbf{x}) + \nabla f(\mathbf{x})^T \mathbf{v} + \int_0^1 (1 - t) \mathbf{v}^T \nabla^2 f(\mathbf{x} + t\mathbf{v}) \mathbf{v} \, dt
\end{equation*}
Subtracting the second-order Taylor polynomial:
\begin{equation*}
f(\mathbf{x} + \mathbf{v}) - \left[ f(\mathbf{x}) + \nabla f(\mathbf{x})^T \mathbf{v} + \frac{1}{2}\mathbf{v}^T \nabla^2 f(\mathbf{x})\mathbf{v} \right] = \int_0^1 (1 - t) \mathbf{v}^T \left( \nabla^2 f(\mathbf{x} + t\mathbf{v}) - \nabla^2 f(\mathbf{x}) \right) \mathbf{v} \, dt
\end{equation*}
Taking absolute values and using the $M$-Lipschitz condition $\|\nabla^2 f(\mathbf{x} + t\mathbf{v}) - \nabla^2 f(\mathbf{x})\|_2 \le M t \|\mathbf{v}\|_2$:
\begin{align*}
|\text{Error}| &\le \int_0^1 (1 - t) \|\mathbf{v}\|_2^2 \left( M t \|\mathbf{v}\|_2 \right) dt = M \|\mathbf{v}\|_2^3 \int_0^1 (t - t^2) dt \\
&= M \|\mathbf{v}\|_2^3 \left[ \frac{t^2}{2} - \frac{t^3}{3} \right]_0^1 = M \|\mathbf{v}\|_2^3 \left( \frac{1}{2} - \frac{1}{3} \right) = \frac{M}{6}\|\mathbf{v}\|_2^3
\end{align*}

\paragraph{P5. Saddle Point vs. Local Extrema Classification:}
$f(x, y) = x^3 - 3x - y^2 + 4y$.
\begin{enumerate}[label=(\alph*)]
    \item $\nabla f(x, y) = [3x^2 - 3, -2y + 4]^T = [0, 0]^T \implies 3(x^2 - 1) = 0 \implies x = \pm 1$, and $-2y + 4 = 0 \implies y = 2$.
    The two critical points are $P_1 = (1, 2)$ and $P_2 = (-1, 2)$.
    \item Hessian matrix: $H(x, y) = \begin{bmatrix} 6x & 0 \\ 0 & -2 \end{bmatrix}$.
    \item At $P_1 = (1, 2)$: $H(1, 2) = \begin{bmatrix} 6 & 0 \\ 0 & -2 \end{bmatrix}$. Eigenvalues are $\lambda_1 = 6 > 0$ and $\lambda_2 = -2 < 0$. Since eigenvalues have opposite signs, $P_1$ is a **Saddle Point**.
    At $P_2 = (-1, 2)$: $H(-1, 2) = \begin{bmatrix} -6 & 0 \\ 0 & -2 \end{bmatrix}$. Eigenvalues are $\lambda_1 = -6 < 0$ and $\lambda_2 = -2 < 0$. Since both eigenvalues are strictly negative, $H$ is negative definite, meaning $P_2$ is a strict **Local Maximum**.
\end{enumerate}

\subsection*{Part B: Diagnostic Multiple-Choice Rationales}

\paragraph{MCQ1: Gradient of Non-Symmetric Quadratic Form}
\begin{itemize}
    \item \textbf{(a) Incorrect.} $2A\mathbf{x} + \mathbf{b}$ is valid only if $A$ is symmetric ($A = A^T$).
    \item \textbf{(b) Correct.} From first-principles double summation, $\nabla_{\mathbf{x}} (\mathbf{x}^T A \mathbf{x}) = (A + A^T)\mathbf{x}$. Adding the linear term $\mathbf{b}^T \mathbf{x}$ gives $(A + A^T)\mathbf{x} + \mathbf{b}$.
    \item \textbf{(c) Incorrect.} Omits the $A\mathbf{x}$ component from the derivative.
    \item \textbf{(d) Incorrect.} Replaces $(A + A^T)$ with $2A^T$, valid only under symmetry.
\end{itemize}

\paragraph{MCQ2: Hessian Eigenvalues and Critical Point Classification}
\begin{itemize}
    \item \textbf{(a) Incorrect.} A positive trace does not guarantee a local minimum; a single negative eigenvalue permits a descent direction.
    \item \textbf{(b) Incorrect.} A local maximum requires all eigenvalues to be strictly negative.
    \item \textbf{(c) Correct.} Because the Hessian has both strictly positive eigenvalues ($\lambda_1 = 4.2, \lambda_2 = 1.8$) and a strictly negative eigenvalue ($\lambda_3 = -0.5$), the matrix is indefinite, making $\mathbf{x}^*$ a saddle point.
    \item \textbf{(d) Incorrect.} The negative eigenvalue direction will repel gradient descent trajectories away from $\mathbf{x}^*$.
\end{itemize}

% ======================================================================
% CHAPTER 3 SOLUTIONS
% ======================================================================
\section*{Chapter 3: Four Fundamental Subspaces, Orthogonal Projections \& OLS}

\subsection*{Part A: Analytical Solutions}

\paragraph{P1. Explicit Projection Matrix Computation:}
$A = \begin{bmatrix} 1 & 0 \\ 1 & 1 \\ 0 & 1 \end{bmatrix}, \mathbf{b} = \begin{bmatrix} 3 \\ 1 \\ 2 \end{bmatrix}$.
\begin{enumerate}[label=(\alph*)]
    \item $A^T A = \begin{bmatrix} 1 & 1 & 0 \\ 0 & 1 & 1 \end{bmatrix} \begin{bmatrix} 1 & 0 \\ 1 & 1 \\ 0 & 1 \end{bmatrix} = \begin{bmatrix} 2 & 1 \\ 1 & 2 \end{bmatrix}$.
    Determinant $\det(A^T A) = 4 - 1 = 3$. Inverse: $(A^T A)^{-1} = \frac{1}{3} \begin{bmatrix} 2 & -1 \\ -1 & 2 \end{bmatrix}$.
    \item $P = A (A^T A)^{-1} A^T = \frac{1}{3} \begin{bmatrix} 1 & 0 \\ 1 & 1 \\ 0 & 1 \end{bmatrix} \begin{bmatrix} 2 & -1 \\ -1 & 2 \end{bmatrix} \begin{bmatrix} 1 & 1 & 0 \\ 0 & 1 & 1 \end{bmatrix} = \frac{1}{3} \begin{bmatrix} 2 & 1 & -1 \\ 1 & 2 & 1 \\ -1 & 1 & 2 \end{bmatrix}$.
    \item $\hat{\mathbf{b}} = P\mathbf{b} = \frac{1}{3} \begin{bmatrix} 2(3) + 1(1) - 1(2) \\ 1(3) + 2(1) + 1(2) \\ -1(3) + 1(1) + 2(2) \end{bmatrix} = \frac{1}{3} \begin{bmatrix} 5 \\ 7 \\ 2 \end{bmatrix} = \begin{bmatrix} 5/3 \\ 7/3 \\ 2/3 \end{bmatrix}$.
    Error vector: $\mathbf{e} = \mathbf{b} - \hat{\mathbf{b}} = \begin{bmatrix} 3 - 5/3 \\ 1 - 7/3 \\ 2 - 2/3 \end{bmatrix} = \begin{bmatrix} 4/3 \\ -4/3 \\ 4/3 \end{bmatrix} = \frac{4}{3} \begin{bmatrix} 1 \\ -1 \\ 1 \end{bmatrix}$.
    \item $A^T \mathbf{e} = \begin{bmatrix} 1 & 1 & 0 \\ 0 & 1 & 1 \end{bmatrix} \left( \frac{4}{3} \begin{bmatrix} 1 \\ -1 \\ 1 \end{bmatrix} \right) = \frac{4}{3} \begin{bmatrix} 1 - 1 + 0 \\ 0 - 1 + 1 \end{bmatrix} = \begin{bmatrix} 0 \\ 0 \end{bmatrix}$. Confirmed.
\end{enumerate}

\paragraph{P2. Properties of Complementary Projectors:}
\begin{enumerate}[label=(\alph*)]
    \item Idempotence: $Q^2 = (I - P)(I - P) = I - 2P + P^2$. Since $P^2 = P$, $Q^2 = I - 2P + P = I - P = Q$.
    Symmetry: $Q^T = (I - P)^T = I^T - P^T = I - P = Q$.
    \item For any $\mathbf{x} \in \mathbb{R}^m$, $P(Q\mathbf{x}) = P(I - P)\mathbf{x} = (P - P^2)\mathbf{x} = \mathbf{0}$, so $\Col(Q) \subseteq \Null(P)$.
    Conversely, if $\mathbf{v} \in \Null(P)$, then $P\mathbf{v} = \mathbf{0} \implies Q\mathbf{v} = (I - P)\mathbf{v} = \mathbf{v}$, so $\mathbf{v} \in \Col(Q)$.
    Thus $\Col(Q) = \Null(P) = V^\perp$.
    \item Using $I = P + Q$ where $P^T Q = P(I - P) = P - P^2 = 0$:
    \begin{equation*}
    \|\mathbf{x}\|_2^2 = \|P\mathbf{x} + Q\mathbf{x}\|_2^2 = \|P\mathbf{x}\|_2^2 + \|Q\mathbf{x}\|_2^2 + 2(P\mathbf{x})^T (Q\mathbf{x}) = \|P\mathbf{x}\|_2^2 + \|Q\mathbf{x}\|_2^2
    \end{equation*}
\end{enumerate}

\paragraph{P3. Rank-1 Projection Matrix:}
\begin{enumerate}[label=(\alph*)]
    \item Let $A = \mathbf{a} \in \mathbb{R}^{m \times 1}$. $A^T A = \mathbf{a}^T \mathbf{a} = \|\mathbf{a}\|_2^2$.
    The projection matrix is $P_{\mathbf{a}} = A(A^T A)^{-1}A^T = \frac{\mathbf{a}\mathbf{a}^T}{\mathbf{a}^T \mathbf{a}}$.
    \item $\trace(P_{\mathbf{a}}) = \trace\left(\frac{\mathbf{a}\mathbf{a}^T}{\mathbf{a}^T \mathbf{a}}\right) = \frac{1}{\mathbf{a}^T \mathbf{a}}\trace(\mathbf{a}\mathbf{a}^T) = \frac{\mathbf{a}^T \mathbf{a}}{\mathbf{a}^T \mathbf{a}} = 1$.
    \item Projected point is $\hat{\mathbf{b}} = P_{\mathbf{a}}\mathbf{b} = \frac{\mathbf{a}\mathbf{a}^T \mathbf{b}}{\mathbf{a}^T \mathbf{a}} = \left( \frac{\mathbf{a}^T \mathbf{b}}{\mathbf{a}^T \mathbf{a}} \right) \mathbf{a}$. Setting $c = \frac{\mathbf{a}^T \mathbf{b}}{\mathbf{a}^T \mathbf{a}}$ gives the exact projection, which minimizes Euclidean distance.
\end{enumerate}

\paragraph{P4. Affine Fit with Offset:}
Points: $(0, 1), (1, 3), (2, 4)$ fitted to $y = c + m x$.
\begin{enumerate}[label=(\alph*)]
    \item Design matrix and target: $A = \begin{bmatrix} 1 & 0 \\ 1 & 1 \\ 1 & 2 \end{bmatrix}, \mathbf{w} = \begin{bmatrix} c \\ m \end{bmatrix}, \mathbf{y} = \begin{bmatrix} 1 \\ 3 \\ 4 \end{bmatrix}$.
    \item $A^T A = \begin{bmatrix} 1 & 1 & 1 \\ 0 & 1 & 2 \end{bmatrix} \begin{bmatrix} 1 & 0 \\ 1 & 1 \\ 1 & 2 \end{bmatrix} = \begin{bmatrix} 3 & 3 \\ 3 & 5 \end{bmatrix}$.
    $A^T \mathbf{y} = \begin{bmatrix} 1+3+4 \\ 0+3+8 \end{bmatrix} = \begin{bmatrix} 8 \\ 11 \end{bmatrix}$.
    Solving: $\det(A^T A) = 15 - 9 = 6$.
    \begin{equation*}
    \begin{bmatrix} c \\ m \end{bmatrix} = \frac{1}{6} \begin{bmatrix} 5 & -3 \\ -3 & 3 \end{bmatrix} \begin{bmatrix} 8 \\ 11 \end{bmatrix} = \frac{1}{6} \begin{bmatrix} 40 - 33 \\ -24 + 33 \end{bmatrix} = \frac{1}{6} \begin{bmatrix} 7 \\ 9 \end{bmatrix} = \begin{bmatrix} 7/6 \\ 3/2 \end{bmatrix}
    \end{equation*}
    Fitted line is $y = \frac{7}{6} + \frac{3}{2}x$.
    \item Predictions: $\hat{\mathbf{y}} = A\hat{\mathbf{w}} = [7/6, 16/6, 25/6]^T = [7/6, 8/3, 25/6]^T$.
    Residuals: $\mathbf{y} - \hat{\mathbf{y}} = [1 - 7/6, 3 - 16/6, 4 - 25/6]^T = [-1/6, 2/6, -1/6]^T$.
    Residual sum of squares: $(-1/6)^2 + (2/6)^2 + (-1/6)^2 = \frac{1 + 4 + 1}{36} = \frac{6}{36} = \frac{1}{6} \approx 0.1667$.
\end{enumerate}

\paragraph{P5. Subspace Intersection and Projection Commutativity:}
$(\implies)$ If $P = P_1 P_2$ is an orthogonal projector, it must be symmetric: $P^T = P \implies (P_1 P_2)^T = P_1 P_2 \implies P_2^T P_1^T = P_1 P_2 \implies P_2 P_1 = P_1 P_2$. Hence they commute.
$(\impliedby)$ If $P_1 P_2 = P_2 P_1$:
Symmetry: $(P_1 P_2)^T = P_2^T P_1^T = P_2 P_1 = P_1 P_2$.
Idempotence: $(P_1 P_2)^2 = (P_1 P_2)(P_1 P_2) = P_1 (P_2 P_1) P_2 = P_1 (P_1 P_2) P_2 = P_1^2 P_2^2 = P_1 P_2$.
Since $P_1 P_2$ is symmetric and idempotent, it is an orthogonal projector. Its range is $\Col(P_1 P_2) = V_1 \cap V_2$.

\subsection*{Part B: Diagnostic Multiple-Choice Rationales}

\paragraph{MCQ1: Normal Equations with Rank-Deficient Matrix}
\begin{itemize}
    \item \textbf{(a) Incorrect.} The Normal Equations $A^T A \mathbf{x} = A^T \mathbf{b}$ are \emph{always} consistent because $A^T \mathbf{b} \in \Col(A^T) = \Col(A^T A)$.
    \item \textbf{(b) Correct.} When $\rank(A) = r < n$, $\Null(A^T A) = \Null(A)$ has dimension $n - r > 0$. The system is consistent with an $(n-r)$-dimensional affine subspace of solutions.
    \item \textbf{(c) Incorrect.} $A^T A$ is singular with non-trivial nullspace, hence positive semidefinite, not strictly positive definite.
    \item \textbf{(d) Incorrect.} $(A^T A)^{-1}$ does not exist when $\rank(A) < n$.
\end{itemize}

\paragraph{MCQ2: Eigenvalues of Orthogonal Projection Matrix}
\begin{itemize}
    \item \textbf{(a) Incorrect.} A full set of 1s would make $P = I_5$, projecting onto all of $\mathbb{R}^5$.
    \item \textbf{(b) Incorrect.} Projection eigenvalues are strictly binary $\lambda \in \{0, 1\}$.
    \item \textbf{(c) Correct.} By the Spectral Theorem for projectors, an orthogonal projector onto an $r$-dimensional subspace in $\mathbb{R}^m$ has eigenvalue $1$ with multiplicity $r$ and eigenvalue $0$ with multiplicity $m - r$. Here $m=5, r=3 \implies 1, 1, 1, 0, 0$.
    \item \textbf{(d) Incorrect.} Negative eigenvalues appear in reflection operators ($R = 2P - I$), not projection operators.
\end{itemize}

% ======================================================================
% CHAPTER 4 SOLUTIONS
% ======================================================================
\section*{Chapter 4: Polynomial Regression, Eigenvalues \& Diagonalization}

\subsection*{Part A: Analytical Solutions}

\paragraph{P1. Eigendecomposition of a $2 \times 2$ Matrix:}
$A = \begin{bmatrix} 4 & 2 \\ 1 & 3 \end{bmatrix}$.
\begin{enumerate}[label=(\alph*)]
    \item Characteristic equation: $\det(A - \lambda I) = (4 - \lambda)(3 - \lambda) - 2 = \lambda^2 - 7\lambda + 10 = (\lambda - 5)(\lambda - 2) = 0$.
    Eigenvalues are $\lambda_1 = 5$ and $\lambda_2 = 2$.
    \item For $\lambda_1 = 5$: $(A - 5I)\mathbf{v}_1 = \begin{bmatrix} -1 & 2 \\ 1 & -2 \end{bmatrix} \begin{bmatrix} v_{11} \\ v_{12} \end{bmatrix} = \mathbf{0} \implies -v_{11} + 2v_{12} = 0 \implies \mathbf{v}_1 = \begin{bmatrix} 2 \\ 1 \end{bmatrix}$.
    For $\lambda_2 = 2$: $(A - 2I)\mathbf{v}_2 = \begin{bmatrix} 2 & 2 \\ 1 & 1 \end{bmatrix} \begin{bmatrix} v_{21} \\ v_{22} \end{bmatrix} = \mathbf{0} \implies v_{21} + v_{22} = 0 \implies \mathbf{v}_2 = \begin{bmatrix} 1 \\ -1 \end{bmatrix}$.
    \item $P = \begin{bmatrix} 2 & 1 \\ 1 & -1 \end{bmatrix}, D = \begin{bmatrix} 5 & 0 \\ 0 & 2 \end{bmatrix}$. $\det(P) = -2 - 1 = -3$. Inverse: $P^{-1} = -\frac{1}{3}\begin{bmatrix} -1 & -1 \\ -1 & 2 \end{bmatrix} = \frac{1}{3}\begin{bmatrix} 1 & 1 \\ 1 & -2 \end{bmatrix}$.
    \item $A^5 = P D^5 P^{-1} = \begin{bmatrix} 2 & 1 \\ 1 & -1 \end{bmatrix} \begin{bmatrix} 3125 & 0 \\ 0 & 32 \end{bmatrix} \left( \frac{1}{3} \begin{bmatrix} 1 & 1 \\ 1 & -2 \end{bmatrix} \right) = \frac{1}{3} \begin{bmatrix} 6250 & 32 \\ 3125 & -32 \end{bmatrix} \begin{bmatrix} 1 & 1 \\ 1 & -2 \end{bmatrix} = \begin{bmatrix} 2094 & 2062 \\ 1031 & 1063 \end{bmatrix}$.
\end{enumerate}

\paragraph{P2. Defective Matrix and Multiplicities:}
$B = \begin{bmatrix} 2 & 1 & 0 \\ 0 & 2 & 0 \\ 0 & 0 & 5 \end{bmatrix}$.
\begin{enumerate}[label=(\alph*)]
    \item $B$ is upper triangular. Characteristic polynomial: $(2 - \lambda)^2 (5 - \lambda) = 0$.
    Eigenvalue $\lambda_1 = 2$ with $am(2) = 2$; eigenvalue $\lambda_2 = 5$ with $am(5) = 1$.
    \item For $\lambda = 2$: $B - 2I = \begin{bmatrix} 0 & 1 & 0 \\ 0 & 0 & 0 \\ 0 & 0 & 3 \end{bmatrix}$. RREF has pivot variables in columns 2 and 3; column 1 is free.
    $E_2 = \Null(B - 2I) = \Span\left\{ [1, 0, 0]^T \right\}$.
    For $\lambda = 5$: $B - 5I = \begin{bmatrix} -3 & 1 & 0 \\ 0 & -3 & 0 \\ 0 & 0 & 0 \end{bmatrix} \implies E_5 = \Span\left\{ [0, 0, 1]^T \right\}$.
    \item Geometric multiplicities: $gm(2) = \dim(E_2) = 1$; $gm(5) = \dim(E_5) = 1$.
    \item Because $gm(2) = 1 < am(2) = 2$, matrix $B$ does not have a complete basis of 3 eigenvectors. Therefore, $B$ is **defective (not diagonalizable)**.
\end{enumerate}

\paragraph{P3. Spectral Theorem for Symmetric Matrices:}
\begin{enumerate}[label=(\alph*)]
    \item Let $S\mathbf{v} = \lambda \mathbf{v}$ with $\mathbf{v} \in \mathbb{C}^n \setminus \{\mathbf{0}\}$.
    Pre-multiply by $\mathbf{v}^H$: $\mathbf{v}^H S \mathbf{v} = \lambda \mathbf{v}^H \mathbf{v} = \lambda \|\mathbf{v}\|_2^2$.
    Taking conjugate transpose: $(\mathbf{v}^H S \mathbf{v})^H = \mathbf{v}^H S^T \mathbf{v} = \mathbf{v}^H S \mathbf{v}$.
    Thus $\mathbf{v}^H S \mathbf{v}$ is real. Since $\lambda = \frac{\mathbf{v}^H S \mathbf{v}}{\|\mathbf{v}\|_2^2}$ is the ratio of two real numbers, $\lambda \in \mathbb{R}$.
    \item Let $S\mathbf{v}_1 = \lambda_1 \mathbf{v}_1$ and $S\mathbf{v}_2 = \lambda_2 \mathbf{v}_2$ with $\lambda_1 \ne \lambda_2$.
    Consider $\mathbf{v}_1^T S \mathbf{v}_2 = \lambda_2 (\mathbf{v}_1^T \mathbf{v}_2)$.
    Also $\mathbf{v}_1^T S \mathbf{v}_2 = (S\mathbf{v}_1)^T \mathbf{v}_2 = \lambda_1 (\mathbf{v}_1^T \mathbf{v}_2)$.
    Subtracting: $(\lambda_1 - \lambda_2)(\mathbf{v}_1^T \mathbf{v}_2) = 0 \implies \mathbf{v}_1^T \mathbf{v}_2 = 0$.
\end{enumerate}

\paragraph{P4. Cayley-Hamilton Application:}
$A = \begin{bmatrix} 2 & 1 \\ 3 & 4 \end{bmatrix}$.
\begin{enumerate}[label=(\alph*)]
    \item $\trace(A) = 6$, $\det(A) = 8 - 3 = 5$. Characteristic polynomial: $p_A(\lambda) = \lambda^2 - 6\lambda + 5 = 0$.
    \item $A^2 = \begin{bmatrix} 2 & 1 \\ 3 & 4 \end{bmatrix} \begin{bmatrix} 2 & 1 \\ 3 & 4 \end{bmatrix} = \begin{bmatrix} 7 & 6 \\ 18 & 19 \end{bmatrix}$.
    $A^2 - 6A + 5I = \begin{bmatrix} 7 & 6 \\ 18 & 19 \end{bmatrix} - \begin{bmatrix} 12 & 6 \\ 18 & 24 \end{bmatrix} + \begin{bmatrix} 5 & 0 \\ 0 & 5 \end{bmatrix} = \begin{bmatrix} 0 & 0 \\ 0 & 0 \end{bmatrix}$. Confirmed.
    \item From $A^2 - 6A + 5I = 0 \implies 5I = 6A - A^2 \implies I = \frac{1}{5}(6A - A^2)$.
    Multiply by $A^{-1}$: $A^{-1} = \frac{6}{5}I - \frac{1}{5}A = \frac{1}{5}(6I - A) = \frac{1}{5}\begin{bmatrix} 4 & -1 \\ -3 & 2 \end{bmatrix}$.
\end{enumerate}

\paragraph{P5. Markov Transition Matrix Dynamics:}
$T = \begin{bmatrix} 0.8 & 0.3 \\ 0.2 & 0.7 \end{bmatrix}$.
\begin{enumerate}[label=(\alph*)]
    \item For any column-stochastic matrix, $\mathbf{1}^T T = \mathbf{1}^T \implies T^T \mathbf{1} = \mathbf{1}$. Thus $T^T$ has eigenvalue $\lambda = 1$ with eigenvector $\mathbf{1}$. Since $T$ and $T^T$ share eigenvalues, $\lambda_1 = 1$.
    \item $\trace(T) = 0.8 + 0.7 = 1.5 = \lambda_1 + \lambda_2 \implies 1 + \lambda_2 = 1.5 \implies \lambda_2 = 0.5$.
    \item Eigenvector for $\lambda_1 = 1$: $(T - I)\boldsymbol{\pi} = \begin{bmatrix} -0.2 & 0.3 \\ 0.2 & -0.3 \end{bmatrix}\boldsymbol{\pi} = \mathbf{0} \implies 2\pi_1 = 3\pi_2$.
    Normalizing $\pi_1 + \pi_2 = 1 \implies \boldsymbol{\pi} = [3/5, 2/5]^T = [0.6, 0.4]^T$.
    Since $|\lambda_2| = 0.5 < 1$, $\lim_{k \to \infty} \lambda_2^k = 0$, so $\lim_{k \to \infty} T^k \mathbf{p}_0 = \boldsymbol{\pi}$ for all initial probability vectors $\mathbf{p}_0$.
\end{enumerate}

\subsection*{Part B: Diagnostic Multiple-Choice Rationales}

\paragraph{MCQ1: Trace and Determinant Invariants}
\begin{itemize}
    \item \textbf{(a) Correct.} $\trace(A) = \lambda_1 + \lambda_2 + \lambda_3 = 2 + 4 + \lambda_3 = 6 \implies \lambda_3 = 0$. $\det(A) = \lambda_1 \lambda_2 \lambda_3 = 2 \cdot 4 \cdot 0 = 0$. Since $\det(A) = 0$, $A$ is singular.
    \item \textbf{(b) Incorrect.} If $\lambda_3 = 1$, trace would be $7$ and determinant would be $8$, contradicting $\trace(A) = 6$.
    \item \textbf{(c) Incorrect.} If $\lambda_3 = 0$, the determinant is zero, meaning $A$ is singular, not invertible.
    \item \textbf{(d) Incorrect.} Contradicts trace and determinant values.
\end{itemize}

\paragraph{MCQ2: High-Degree Polynomial Overfitting}
\begin{itemize}
    \item \textbf{(a) Incorrect.} A Vandermonde matrix with $N \ge p+1$ distinct points has full rank.
    \item \textbf{(b) Incorrect.} High-degree polynomials have exceptionally low bias (zero training error).
    \item \textbf{(c) Correct.} Fitting a degree-$15$ polynomial to $16$ points causes extreme variance and Runge's phenomenon, oscillating wildly between interpolation nodes and destroying out-of-sample test accuracy.
    \item \textbf{(d) Incorrect.} Polynomial regression does not mandate a Gaussian target assumption to compute OLS estimates.
\end{itemize}

% ======================================================================
% CHAPTER 5 SOLUTIONS
% ======================================================================
\section*{Chapter 5: Complex Spaces, Hermitian Matrices \& Spectral Theorem}

\subsection*{Part A: Analytical Solutions}

\paragraph{P1. Verifying Hermitian Properties:}
$H = \begin{bmatrix} 2 & 1 - i \\ 1 + i & 3 \end{bmatrix}$.
\begin{enumerate}[label=(\alph*)]
    \item $H^H = \overline{H}^T = \begin{bmatrix} \overline{2} & \overline{1+i} \\ \overline{1-i} & \overline{3} \end{bmatrix}^T = \begin{bmatrix} 2 & 1-i \\ 1+i & 3 \end{bmatrix} = H$. $H$ is Hermitian.
    \item $\det(H - \lambda I) = (2 - \lambda)(3 - \lambda) - (1 - i)(1 + i) = \lambda^2 - 5\lambda + 6 - (1 - i^2) = \lambda^2 - 5\lambda + 6 - 2 = \lambda^2 - 5\lambda + 4 = (\lambda - 4)(\lambda - 1) = 0$.
    Eigenvalues are $\lambda_1 = 4$ and $\lambda_2 = 1$ (both real!).
    \item For $\lambda_1 = 4$: $(H - 4I)\mathbf{u}_1 = \begin{bmatrix} -2 & 1-i \\ 1+i & -1 \end{bmatrix}\mathbf{u}_1 = \mathbf{0} \implies -2u_{11} + (1-i)u_{12} = 0 \implies \mathbf{u}_1 \propto [1-i, 2]^T$.
    Norm: $\|\mathbf{u}_1\|_2^2 = |1-i|^2 + |2|^2 = 2 + 4 = 6 \implies \mathbf{u}_1 = \frac{1}{\sqrt{6}}\begin{bmatrix} 1-i \\ 2 \end{bmatrix}$.
    For $\lambda_2 = 1$: $(H - I)\mathbf{u}_2 = \begin{bmatrix} 1 & 1-i \\ 1+i & 2 \end{bmatrix}\mathbf{u}_2 = \mathbf{0} \implies \mathbf{u}_2 \propto [1-i, -1]^T$.
    Norm: $\|\mathbf{u}_2\|_2^2 = 2 + 1 = 3 \implies \mathbf{u}_2 = \frac{1}{\sqrt{3}}\begin{bmatrix} 1-i \\ -1 \end{bmatrix}$.
    Inner product: $\mathbf{u}_1^H \mathbf{u}_2 = \frac{1}{\sqrt{18}} \left( (1+i)(1-i) + 2(-1) \right) = \frac{1}{\sqrt{18}}(2 - 2) = 0$.
    \item $\lambda_1 \mathbf{u}_1 \mathbf{u}_1^H + \lambda_2 \mathbf{u}_2 \mathbf{u}_2^H = \frac{4}{6}\begin{bmatrix} 2 & 2-2i \\ 2+2i & 4 \end{bmatrix} + \frac{1}{3}\begin{bmatrix} 2 & -(1-i) \\ -(1+i) & 1 \end{bmatrix} = \begin{bmatrix} 4/3 + 2/3 & (4-4i)/3 - (1-i)/3 \\ (4+4i)/3 - (1+i)/3 & 8/3 + 1/3 \end{bmatrix} = \begin{bmatrix} 2 & 1-i \\ 1+i & 3 \end{bmatrix} = H$.
\end{enumerate}

\paragraph{P2. Spectral Properties of Unitary Transformations:}
\begin{enumerate}[label=(\alph*)]
    \item $\|U\mathbf{x}\|_2^2 = (U\mathbf{x})^H (U\mathbf{x}) = \mathbf{x}^H U^H U \mathbf{x} = \mathbf{x}^H I \mathbf{x} = \|\mathbf{x}\|_2^2 \implies \|U\mathbf{x}\|_2 = \|\mathbf{x}\|_2$.
    \item Let $U\mathbf{v} = \lambda \mathbf{v}$ with $\|\mathbf{v}\| \ne 0$. $\|U\mathbf{v}\|_2 = \|\lambda \mathbf{v}\|_2 = |\lambda| \|\mathbf{v}\|_2$. By length preservation, $\|U\mathbf{v}\|_2 = \|\mathbf{v}\|_2 \implies |\lambda| = 1$.
    \item Let $U\mathbf{v}_1 = \lambda_1 \mathbf{v}_1, U\mathbf{v}_2 = \lambda_2 \mathbf{v}_2$.
    $\mathbf{v}_1^H \mathbf{v}_2 = (U\mathbf{v}_1)^H (U\mathbf{v}_2) = (\lambda_1 \mathbf{v}_1)^H (\lambda_2 \mathbf{v}_2) = \bar{\lambda}_1 \lambda_2 (\mathbf{v}_1^H \mathbf{v}_2)$.
    Since $|\lambda_1| = 1 \implies \bar{\lambda}_1 = \frac{1}{\lambda_1}$, we have $\mathbf{v}_1^H \mathbf{v}_2 = \frac{\lambda_2}{\lambda_1}(\mathbf{v}_1^H \mathbf{v}_2) \implies \left(1 - \frac{\lambda_2}{\lambda_1}\right)(\mathbf{v}_1^H \mathbf{v}_2) = 0$. Since $\lambda_1 \ne \lambda_2$, $\mathbf{v}_1^H \mathbf{v}_2 = 0$.
\end{enumerate}

\paragraph{P3. Skew-Hermitian Matrices:}
\begin{enumerate}[label=(\alph*)]
    \item Let $S\mathbf{v} = \lambda \mathbf{v}$ with $\|\mathbf{v}\|_2 = 1$. $\lambda = \mathbf{v}^H S \mathbf{v}$.
    Taking conjugate transpose: $\bar{\lambda} = (\mathbf{v}^H S \mathbf{v})^H = \mathbf{v}^H S^H \mathbf{v} = \mathbf{v}^H (-S)\mathbf{v} = -\lambda$.
    Thus $\bar{\lambda} + \lambda = 2\text{Re}(\lambda) = 0 \implies \text{Re}(\lambda) = 0$. Hence $\lambda = i\mu$ with $\mu \in \mathbb{R}$.
    \item $U = \exp(S)$. $U^H = (\exp(S))^H = \exp(S^H) = \exp(-S)$.
    Since $S$ and $-S$ commute, $U^H U = \exp(-S)\exp(S) = \exp(-S + S) = \exp(0) = I$. Thus $U$ is unitary.
\end{enumerate}

\paragraph{P4. Decomposition into Hermitian and Skew-Hermitian Components:}
\begin{enumerate}[label=(\alph*)]
    \item Any square matrix can be written as $M = \frac{M + M^H}{2} + \frac{M - M^H}{2} = H + S$.
    $H^H = \left(\frac{M + M^H}{2}\right)^H = \frac{M^H + M}{2} = H$ (Hermitian).
    $S^H = \left(\frac{M - M^H}{2}\right)^H = \frac{M^H - M}{2} = -S$ (Skew-Hermitian).
    \item $H = \frac{1}{2}(M + M^H)$ and $S = \frac{1}{2}(M - M^H)$.
    \item $M = \begin{bmatrix} 1+i & 4 \\ 2i & 3 \end{bmatrix} \implies M^H = \begin{bmatrix} 1-i & -2i \\ 4 & 3 \end{bmatrix}$.
    $H = \frac{1}{2} \begin{bmatrix} 2 & 4-2i \\ 4+2i & 6 \end{bmatrix} = \begin{bmatrix} 1 & 2-i \\ 2+i & 3 \end{bmatrix}$.
    $S = \frac{1}{2} \begin{bmatrix} 2i & 4+2i \\ 2i-4 & 0 \end{bmatrix} = \begin{bmatrix} i & 2+i \\ -2+i & 0 \end{bmatrix}$.
\end{enumerate}

\paragraph{P5. Courant-Fischer Min-Max Theorem:}
By Courant-Fischer: $\lambda_k(A + B) = \min_{S: \dim(S)=k} \max_{\mathbf{x} \in S, \|\mathbf{x}\|=1} \left( \mathbf{x}^H A \mathbf{x} + \mathbf{x}^H B \mathbf{x} \right)$.
Since $B \succeq 0$, $\mathbf{x}^H B \mathbf{x} \ge 0$ for all $\mathbf{x}$.
Therefore, for any fixed $\mathbf{x}$: $\mathbf{x}^H (A + B) \mathbf{x} = \mathbf{x}^H A \mathbf{x} + \mathbf{x}^H B \mathbf{x} \ge \mathbf{x}^H A \mathbf{x}$.
Taking the max over $\mathbf{x} \in S$ preserves the inequality: $\max_{\mathbf{x} \in S, \|\mathbf{x}\|=1} \mathbf{x}^H (A+B)\mathbf{x} \ge \max_{\mathbf{x} \in S, \|\mathbf{x}\|=1} \mathbf{x}^H A \mathbf{x}$.
Taking the min over all $k$-dimensional subspaces $S$ preserves the inequality, proving $\lambda_k(A + B) \ge \lambda_k(A)$.

\subsection*{Part B: Diagnostic Multiple-Choice Rationales}

\paragraph{MCQ1: Hermitian Matrix Properties}
\begin{itemize}
    \item \textbf{(a) Guaranteed.} The Spectral Theorem guarantees that all eigenvalues of a Hermitian matrix are strictly real.
    \item \textbf{(b) Guaranteed.} Every Hermitian matrix is unitarily diagonalizable, yielding an orthonormal basis of eigenvectors.
    \item \textbf{(c) Guaranteed.} The determinant is the product of the real eigenvalues, hence real.
    \item \textbf{(d) NOT Guaranteed.} Main diagonal entries of a Hermitian matrix must be real ($A_{ii} = \overline{A_{ii}}$), but can be negative (e.g. $A = \begin{bmatrix} -3 & 0 \\ 0 & -1 \end{bmatrix}$ is Hermitian).
\end{itemize}

\paragraph{MCQ2: Unitary Operator Geometry}
\begin{itemize}
    \item \textbf{(a) Incorrect.} Projections are idempotent ($P^2 = P$) and not invertible; unitary transformations are isometries ($U^H U = I$).
    \item \textbf{(b) Incorrect.} Unitary transformations preserve lengths ($c = 1$).
    \item \textbf{(c) Correct.} Unitary operators preserve complex inner products ($\langle U\mathbf{u}, U\mathbf{v} \rangle = \langle \mathbf{u}, \mathbf{v} \rangle$) and Euclidean norms, functioning as rigid complex rotations/reflections.
    \item \textbf{(d) Incorrect.} Eigenvalues of unitary matrices have modulus 1 ($|\lambda| = 1$), never zero.
\end{itemize}

% ======================================================================
% CHAPTER 6 SOLUTIONS
% ======================================================================
\section*{Chapter 6: SVD, Low-Rank Approximations \& Positive Definite Matrices}

\subsection*{Part A: Analytical Solutions}

\paragraph{P1. Explicit Full SVD Computation by Hand:}
$A = \begin{bmatrix} 3 & 0 \\ 4 & 5 \end{bmatrix}$.
\begin{enumerate}[label=(\alph*)]
    \item $A^T A = \begin{bmatrix} 3 & 4 \\ 0 & 5 \end{bmatrix} \begin{bmatrix} 3 & 0 \\ 4 & 5 \end{bmatrix} = \begin{bmatrix} 25 & 20 \\ 20 & 25 \end{bmatrix}$.
    Characteristic equation: $\det(A^T A - \lambda I) = (25 - \lambda)^2 - 400 = \lambda^2 - 50\lambda + 625 - 400 = \lambda^2 - 50\lambda + 225 = (\lambda - 45)(\lambda - 5) = 0$.
    Eigenvalues: $\lambda_1 = 45, \lambda_2 = 5$.
    \item Singular values: $\sigma_1 = \sqrt{45} = 3\sqrt{5} \approx 6.708$, $\sigma_2 = \sqrt{5} \approx 2.236$.
    \item Right singular vectors:
    For $\lambda_1 = 45$: $\begin{bmatrix} -20 & 20 \\ 20 & -20 \end{bmatrix}\mathbf{v}_1 = \mathbf{0} \implies \mathbf{v}_1 = \frac{1}{\sqrt{2}}\begin{bmatrix} 1 \\ 1 \end{bmatrix}$.
    For $\lambda_2 = 5$: $\begin{bmatrix} 20 & 20 \\ 20 & 20 \end{bmatrix}\mathbf{v}_2 = \mathbf{0} \implies \mathbf{v}_2 = \frac{1}{\sqrt{2}}\begin{bmatrix} -1 \\ 1 \end{bmatrix}$.
    \item Left singular vectors:
    $\mathbf{u}_1 = \frac{1}{\sigma_1} A \mathbf{v}_1 = \frac{1}{3\sqrt{5}} \begin{bmatrix} 3 & 0 \\ 4 & 5 \end{bmatrix} \left( \frac{1}{\sqrt{2}}\begin{bmatrix} 1 \\ 1 \end{bmatrix} \right) = \frac{1}{3\sqrt{10}} \begin{bmatrix} 3 \\ 9 \end{bmatrix} = \frac{1}{\sqrt{10}} \begin{bmatrix} 1 \\ 3 \end{bmatrix}$.
    $\mathbf{u}_2 = \frac{1}{\sigma_2} A \mathbf{v}_2 = \frac{1}{\sqrt{5}} \begin{bmatrix} 3 & 0 \\ 4 & 5 \end{bmatrix} \left( \frac{1}{\sqrt{2}}\begin{bmatrix} -1 \\ 1 \end{bmatrix} \right) = \frac{1}{\sqrt{10}} \begin{bmatrix} -3 \\ 1 \end{bmatrix}$.
    Verify: $U \Sigma V^T = \begin{bmatrix} 1/\sqrt{10} & -3/\sqrt{10} \\ 3/\sqrt{10} & 1/\sqrt{10} \end{bmatrix} \begin{bmatrix} 3\sqrt{5} & 0 \\ 0 & \sqrt{5} \end{bmatrix} \begin{bmatrix} 1/\sqrt{2} & 1/\sqrt{2} \\ -1/\sqrt{2} & 1/\sqrt{2} \end{bmatrix} = \begin{bmatrix} 3 & 0 \\ 4 & 5 \end{bmatrix} = A$.
\end{enumerate}

\paragraph{P2. Sylvester's Criterion vs. Cholesky Factorization:}
$M = \begin{bmatrix} 2 & -1 & 0 \\ -1 & 2 & -1 \\ 0 & -1 & 2 \end{bmatrix}$.
\begin{enumerate}[label=(\alph*)]
    \item $\Delta_1 = \det[2] = 2 > 0$.
    $\Delta_2 = \det\begin{bmatrix} 2 & -1 \\ -1 & 2 \end{bmatrix} = 4 - 1 = 3 > 0$.
    $\Delta_3 = \det(M) = 2(4 - 1) - (-1)(-2 - 0) + 0 = 6 - 2 = 4 > 0$.
    \item Since all three leading principal minors are strictly positive ($\Delta_1=2, \Delta_2=3, \Delta_3=4 > 0$), Sylvester's Criterion proves $M$ is **strictly positive definite**.
    \item Cholesky factor $L = \begin{bmatrix} l_{11} & 0 & 0 \\ l_{21} & l_{22} & 0 \\ l_{31} & l_{32} & l_{33} \end{bmatrix}$ such that $L L^T = M$:
    $l_{11}^2 = 2 \implies l_{11} = \sqrt{2}$.
    $l_{21} l_{11} = -1 \implies l_{21} = -\frac{1}{\sqrt{2}}$.
    $l_{21}^2 + l_{22}^2 = 2 \implies 1/2 + l_{22}^2 = 2 \implies l_{22} = \sqrt{3/2}$.
    $l_{31} l_{11} = 0 \implies l_{31} = 0$.
    $l_{31} l_{21} + l_{32} l_{22} = -1 \implies 0 + l_{32}\sqrt{3/2} = -1 \implies l_{32} = -\sqrt{2/3}$.
    $l_{31}^2 + l_{32}^2 + l_{33}^2 = 2 \implies 2/3 + l_{33}^2 = 2 \implies l_{33} = \sqrt{4/3} = \frac{2}{\sqrt{3}}$.
    Thus $L = \begin{bmatrix} \sqrt{2} & 0 & 0 \\ -1/\sqrt{2} & \sqrt{3/2} & 0 \\ 0 & -\sqrt{2/3} & 2/\sqrt{3} \end{bmatrix}$.
\end{enumerate}

\paragraph{P3. Frobenius Norm Representation in Terms of Singular Values:}
By definition of Frobenius norm: $\|A\|_F^2 = \sum_{i=1}^m \sum_{j=1}^n A_{ij}^2 = \trace(A^T A)$.
Substitute the SVD $A = U \Sigma V^T$:
\begin{align*}
\trace(A^T A) &= \trace\left( (U \Sigma V^T)^T (U \Sigma V^T) \right) = \trace\left( V \Sigma^T U^T U \Sigma V^T \right) \\
&= \trace\left( V \Sigma^T \Sigma V^T \right)
\end{align*}
Using the cyclic trace property $\trace(BC) = \trace(CB)$:
\begin{equation*}
\trace\left( V (\Sigma^T \Sigma) V^T \right) = \trace\left( (\Sigma^T \Sigma) V^T V \right) = \trace(\Sigma^T \Sigma) = \sum_{i=1}^r \sigma_i^2
\end{equation*}

\paragraph{P4. Condition Number and Error Sensitivity:}
From $A\mathbf{x} = \mathbf{b}$, taking norms gives $\|\mathbf{b}\|_2 = \|A\mathbf{x}\|_2 \le \|A\|_2 \|\mathbf{x}\|_2 \implies \frac{1}{\|\mathbf{x}\|_2} \le \frac{\|A\|_2}{\|\mathbf{b}\|_2}$.
From $A \delta \mathbf{x} = \delta \mathbf{b}$, we have $\delta \mathbf{x} = A^{-1} \delta \mathbf{b} \implies \|\delta \mathbf{x}\|_2 \le \|A^{-1}\|_2 \|\delta \mathbf{b}\|_2$.
Multiplying the two inequalities:
\begin{equation*}
\frac{\|\delta \mathbf{x}\|_2}{\|\mathbf{x}\|_2} \le \|A\|_2 \|A^{-1}\|_2 \frac{\|\delta \mathbf{b}\|_2}{\|\mathbf{b}\|_2} = \kappa(A) \frac{\|\delta \mathbf{b}\|_2}{\|\mathbf{b}\|_2}
\end{equation*}
where $\kappa(A) = \|A\|_2 \|A^{-1}\|_2 = \sigma_{\max} / \sigma_{\min}$.

\paragraph{P5. Pseudoinverse of Rank-Deficient Matrix:}
$A = \begin{bmatrix} 1 & 1 \\ 2 & 2 \end{bmatrix}$.
\begin{enumerate}[label=(\alph*)]
    \item Columns are linearly dependent; row 2 is $2 \times$ row 1. Thus $\rank(A) = 1$.
    $A^T A = \begin{bmatrix} 5 & 5 \\ 5 & 5 \end{bmatrix} \implies \lambda_1 = 10, \lambda_2 = 0 \implies \sigma_1 = \sqrt{10}$.
    $\mathbf{v}_1 = \frac{1}{\sqrt{2}}\begin{bmatrix} 1 \\ 1 \end{bmatrix}$.
    $\mathbf{u}_1 = \frac{1}{\sigma_1}A\mathbf{v}_1 = \frac{1}{\sqrt{10}}\begin{bmatrix} 1 & 1 \\ 2 & 2 \end{bmatrix}\left(\frac{1}{\sqrt{2}}\begin{bmatrix} 1 \\ 1 \end{bmatrix}\right) = \frac{1}{\sqrt{20}}\begin{bmatrix} 2 \\ 4 \end{bmatrix} = \frac{1}{\sqrt{5}}\begin{bmatrix} 1 \\ 2 \end{bmatrix}$.
    Compact SVD: $A = \sqrt{10} \left( \frac{1}{\sqrt{5}}\begin{bmatrix} 1 \\ 2 \end{bmatrix} \right) \left( \frac{1}{\sqrt{2}}\begin{bmatrix} 1 & 1 \end{bmatrix} \right)$.
    \item $A^+ = \frac{1}{\sigma_1} \mathbf{v}_1 \mathbf{u}_1^T = \frac{1}{\sqrt{10}} \left(\frac{1}{\sqrt{2}}\begin{bmatrix} 1 \\ 1 \end{bmatrix}\right) \left(\frac{1}{\sqrt{5}}\begin{bmatrix} 1 & 2 \end{bmatrix}\right) = \frac{1}{10}\begin{bmatrix} 1 & 2 \\ 1 & 2 \end{bmatrix}$.
    \item $A A^+ = \begin{bmatrix} 1 & 1 \\ 2 & 2 \end{bmatrix} \left(\frac{1}{10}\begin{bmatrix} 1 & 2 \\ 1 & 2 \end{bmatrix}\right) = \frac{1}{10}\begin{bmatrix} 2 & 4 \\ 4 & 8 \end{bmatrix} = \begin{bmatrix} 0.2 & 0.4 \\ 0.4 & 0.8 \end{bmatrix}$.
    $A A^+ A = \begin{bmatrix} 0.2 & 0.4 \\ 0.4 & 0.8 \end{bmatrix} \begin{bmatrix} 1 & 1 \\ 2 & 2 \end{bmatrix} = \begin{bmatrix} 1 & 1 \\ 2 & 2 \end{bmatrix} = A$.
    $A^+ A A^+ = A^+ (A A^+) = \frac{1}{10}\begin{bmatrix} 1 & 2 \\ 1 & 2 \end{bmatrix} \begin{bmatrix} 0.2 & 0.4 \\ 0.4 & 0.8 \end{bmatrix} = \frac{1}{10}\begin{bmatrix} 1 & 2 \\ 1 & 2 \end{bmatrix} = A^+$.
\end{enumerate}

\subsection*{Part B: Diagnostic Multiple-Choice Rationales}

\paragraph{MCQ1: Eckart-Young-Mirsky Low-Rank Approximation Error}
\begin{itemize}
    \item \textbf{(a) Correct.} By the Eckart-Young-Mirsky Theorem under the Frobenius norm, the approximation error for rank $k = 2$ is the root sum of squares of the discarded singular values: $\|A - A_2\|_F = \sqrt{\sum_{i=3}^r \sigma_i^2} = \sqrt{\sigma_3^2 + \sigma_4^2} = \sqrt{5^2 + 2^2} = \sqrt{29}$.
    \item \textbf{(b) Incorrect.} 5 is the Spectral norm error $\|A - A_2\|_2 = \sigma_{k+1} = \sigma_3 = 5$, not the Frobenius norm error.
    \item \textbf{(c) Incorrect.} Linear sum of singular values does not represent Frobenius energy.
    \item \textbf{(d) Incorrect.} Represents the energy of the preserved components $\sigma_1, \sigma_2$.
\end{itemize}

\paragraph{MCQ2: Necessary and Sufficient Positive Definiteness}
\begin{itemize}
    \item \textbf{(a) Incorrect.} Positive entries can yield indefinite matrices (e.g. $\begin{bmatrix} 1 & 3 \\ 3 & 1 \end{bmatrix}$).
    \item \textbf{(b) Incorrect.} In dimensions $d \ge 3$, positive trace and determinant do not prevent pairs of negative eigenvalues.
    \item \textbf{(c) Correct.} By the Spectral Theorem, $A = Q \Lambda Q^T$. The quadratic form $\mathbf{x}^T A \mathbf{x} = \sum \lambda_i y_i^2 > 0$ for all $\mathbf{x} \ne \mathbf{0} \iff \lambda_i > 0$ for all $i$.
    \item \textbf{(d) Incorrect.} Strictly diagonally dominant symmetric matrices with positive diagonals are positive definite, but diagonal dominance is sufficient, not necessary.
\end{itemize}

% ======================================================================
% CHAPTER 7 SOLUTIONS
% ======================================================================
\section*{Chapter 7: Principal Component Analysis (PCA)}

\subsection*{Part A: Analytical Solutions}

\paragraph{P1. Analytical PCA on a Small $2 \times 2$ Covariance Matrix:}
$S = \begin{bmatrix} 5 & 2 \\ 2 & 2 \end{bmatrix}$.
\begin{enumerate}[label=(\alph*)]
    \item $\det(S - \lambda I) = (5 - \lambda)(2 - \lambda) - 4 = \lambda^2 - 7\lambda + 10 - 4 = \lambda^2 - 7\lambda + 6 = (\lambda - 6)(\lambda - 1) = 0$.
    Eigenvalues are $\lambda_1 = 6$ and $\lambda_2 = 1$.
    \item For $\lambda_1 = 6$: $(S - 6I)\mathbf{w}_1 = \begin{bmatrix} -1 & 2 \\ 2 & -4 \end{bmatrix}\mathbf{w}_1 = \mathbf{0} \implies -w_{11} + 2w_{12} = 0 \implies \mathbf{w}_1 = \frac{1}{\sqrt{5}}\begin{bmatrix} 2 \\ 1 \end{bmatrix}$.
    For $\lambda_2 = 1$: $(S - I)\mathbf{w}_2 = \begin{bmatrix} 4 & 2 \\ 2 & 1 \end{bmatrix}\mathbf{w}_2 = \mathbf{0} \implies 2w_{21} + w_{22} = 0 \implies \mathbf{w}_2 = \frac{1}{\sqrt{5}}\begin{bmatrix} -1 \\ 2 \end{bmatrix}$.
    \item Total variance is $\trace(S) = \lambda_1 + \lambda_2 = 6 + 1 = 7$.
    Proportion explained by PC1: $\frac{\lambda_1}{\lambda_1 + \lambda_2} = \frac{6}{7} \approx 85.71\%$.
    \item Centered point $\tilde{\mathbf{x}} = [2, 1]^T$:
    Scalar score: $z_1 = \mathbf{w}_1^T \tilde{\mathbf{x}} = \frac{1}{\sqrt{5}} (2 \cdot 2 + 1 \cdot 1) = \frac{5}{\sqrt{5}} = \sqrt{5} \approx 2.236$.
    Reconstructed vector: $\hat{\mathbf{x}} = z_1 \mathbf{w}_1 = \sqrt{5} \left( \frac{1}{\sqrt{5}}\begin{bmatrix} 2 \\ 1 \end{bmatrix} \right) = \begin{bmatrix} 2 \\ 1 \end{bmatrix} = \tilde{\mathbf{x}}$ (perfect reconstruction because $\tilde{\mathbf{x}}$ lies directly along $\mathbf{w}_1$).
\end{enumerate}

\paragraph{P2. Scale Invariance Counterexample:}
$S = \begin{bmatrix} 10 & 0 \\ 0 & 1 \end{bmatrix}$.
\begin{enumerate}[label=(\alph*)]
    \item The transformation scales feature 2 by $c = 10$, so $x_2' = 10x_2$.
    The new covariance matrix is $S' = \begin{bmatrix} 1 & 0 \\ 0 & 10 \end{bmatrix} \begin{bmatrix} 10 & 0 \\ 0 & 1 \end{bmatrix} \begin{bmatrix} 1 & 0 \\ 0 & 10 \end{bmatrix} = \begin{bmatrix} 10 & 0 \\ 0 & 100 \end{bmatrix}$.
    \item For $S$: $\lambda_1 = 10, \lambda_2 = 1 \implies \mathbf{w}_1 = [1, 0]^T$ (aligned entirely with feature 1).
    For $S'$: $\lambda_1' = 100, \lambda_2' = 10 \implies \mathbf{w}_1' = [0, 1]^T$ (aligned entirely with feature 2).
    \item Scaling a single feature rotated the first principal component by 90 degrees. This demonstrates that PCA is **not scale invariant**: features with large numerical variance dominate the principal components regardless of informational importance, mandating standardizing features to unit variance before PCA.
\end{enumerate}

\paragraph{P3. Orthogonality of Principal Component Scores:}
$Z = \tilde{X} W \in \mathbb{R}^{N \times k}$, where $W = [\mathbf{w}_1, \dots, \mathbf{w}_k]$ contains the orthonormal eigenvectors of $S = \frac{1}{N}\tilde{X}^T \tilde{X}$, satisfying $S W = W \Lambda_k$ and $W^T W = I_k$.
The sample covariance of the scores is:
\begin{equation*}
S_Z = \frac{1}{N} Z^T Z = \frac{1}{N} (W^T \tilde{X}^T) (\tilde{X} W) = W^T \left( \frac{1}{N}\tilde{X}^T \tilde{X} \right) W = W^T S W
\end{equation*}
Substitute $S W = W \Lambda_k$:
\begin{equation*}
S_Z = W^T (W \Lambda_k) = (W^T W) \Lambda_k = I_k \Lambda_k = \diag(\lambda_1, \lambda_2, \dots, \lambda_k)
\end{equation*}
Because $S_Z$ is strictly diagonal, the sample covariance between distinct scores $z_i$ and $z_j$ ($i \ne j$) is zero: $\Cov(z_i, z_j) = 0$. Principal component scores are mutually **uncorrelated**.

\paragraph{P4. Reconstruction Error Bounds:}
Let $W_k = [\mathbf{w}_1, \dots, \mathbf{w}_k]$ and complete orthonormal basis $W_d = [W_k, W_\perp]$.
Any centered sample $\tilde{\mathbf{x}}_i$ can be expanded as $\tilde{\mathbf{x}}_i = W_k W_k^T \tilde{\mathbf{x}}_i + W_\perp W_\perp^T \tilde{\mathbf{x}}_i$.
The reconstruction error vector is $\tilde{\mathbf{x}}_i - W_k W_k^T \tilde{\mathbf{x}}_i = W_\perp W_\perp^T \tilde{\mathbf{x}}_i$.
Using the cyclic trace property:
\begin{align*}
\frac{1}{N} \sum_{i=1}^N \|\tilde{\mathbf{x}}_i - W_k W_k^T \tilde{\mathbf{x}}_i\|_2^2 &= \frac{1}{N} \sum_{i=1}^N \trace\left( W_\perp W_\perp^T \tilde{\mathbf{x}}_i \tilde{\mathbf{x}}_i^T W_\perp W_\perp^T \right) \\
&= \trace\left( W_\perp^T \left( \frac{1}{N} \sum_{i=1}^N \tilde{\mathbf{x}}_i \tilde{\mathbf{x}}_i^T \right) W_\perp \right) = \trace(W_\perp^T S W_\perp)
\end{align*}
Since the columns of $W_\perp$ are the remaining $d - k$ eigenvectors $\mathbf{w}_{k+1}, \dots, \mathbf{w}_d$:
\begin{equation*}
\trace(W_\perp^T S W_\perp) = \trace(\diag(\lambda_{k+1}, \dots, \lambda_d)) = \sum_{j=k+1}^d \lambda_j
\end{equation*}

\paragraph{P5. High-Dimensional Dual PCA Rank Bound:}
Let $\tilde{X} \in \mathbb{R}^{N \times d}$. By construction, the rows of $\tilde{X}$ are centered to have zero mean: $\sum_{i=1}^N \tilde{\mathbf{x}}_i = \mathbf{0} \implies \mathbf{1}_N^T \tilde{X} = \mathbf{0}^T$.
Thus the all-ones vector $\mathbf{1}_N$ lies in the left nullspace of $\tilde{X}$, meaning $\rank(\tilde{X}) \le N - 1$.
Because $\rank(S) = \rank\left(\frac{1}{N}\tilde{X}^T \tilde{X}\right) = \rank(\tilde{X}) \le \min(N - 1, d) = N - 1$.
Since $S$ is a $d \times d$ matrix with rank at most $N - 1$, by the Rank-Nullity Theorem:
\begin{equation*}
\dim(\Null(S)) = d - \rank(S) \ge d - (N - 1) = d - N + 1
\end{equation*}
Thus at least $d - N + 1$ eigenvalues must be identically zero.

\subsection*{Part B: Diagnostic Multiple-Choice Rationales}

\paragraph{MCQ1: Proportion of Variance Explained}
\begin{itemize}
    \item \textbf{(a) Correct.} Total variance is $\sum_{j=1}^4 \lambda_j = 12 + 5 + 2 + 1 = 20$. Retaining top 2 components captures $\lambda_1 + \lambda_2 = 12 + 5 = 17$. The proportion is $\frac{17}{20} = 85\%$.
    \item \textbf{(b) Incorrect.} Uses squared eigenvalues (which would correspond to squared covariance eigenvalues, mathematically erroneous for PCA variance ratio).
    \item \textbf{(c) Incorrect.} Only accounts for the first principal component ($12/20 = 60\%$).
    \item \textbf{(d) Incorrect.} Square roots of eigenvalues represent singular values, not variance proportions.
\end{itemize}

\paragraph{MCQ2: When PCA Harms Classification}
\begin{itemize}
    \item \textbf{(a) Incorrect.} If features are uncorrelated, PCA simply orders them by variance.
    \item \textbf{(b) Correct.} PCA is unsupervised and maximizes feature variance. If the class-separating signal lies along a low-variance direction, PCA will project it out as noise, destroying class separability.
    \item \textbf{(c) Incorrect.} $N > d$ is standard overdetermined data where PCA performs stably.
    \item \textbf{(d) Incorrect.} Centering data is required for PCA and does not destroy classification manifolds.
\end{itemize}

% ======================================================================
% CHAPTER 8 SOLUTIONS
% ======================================================================
\section*{Chapter 8: Unconstrained Optimization \& Gradient Descent}

\subsection*{Part A: Analytical Solutions}

\paragraph{P1. Analytical Step-Size Bound on a 2D Quadratic:}
$f(x_1, x_2) = 10 x_1^2 + x_2^2 - 4x_1 - 2x_2$.
\begin{enumerate}[label=(\alph*)]
    \item $\nabla f(\mathbf{x}) = [20x_1 - 4, 2x_2 - 2]^T$.
    $\nabla^2 f(\mathbf{x}) = \begin{bmatrix} 20 & 0 \\ 0 & 2 \end{bmatrix}$.
    \item Setting $\nabla f = \mathbf{0} \implies 20x_1 = 4 \implies x_1^* = 0.2$, and $2x_2 = 2 \implies x_2^* = 1.0$.
    Global minimum is $\mathbf{x}^* = [0.2, 1.0]^T$.
    \item Eigenvalues of the Hessian: $\lambda_1 = 20, \lambda_2 = 2$.
    Lipschitz constant $L = \lambda_{\max} = 20$.
    Strong convexity parameter $\mu = \lambda_{\min} = 2$.
    \item Maximum stable step size: $\eta_{\max} = \frac{2}{L} = \frac{2}{20} = 0.1$.
    Optimal step size: $\eta^* = \frac{2}{L + \mu} = \frac{2}{20 + 2} = \frac{2}{22} = \frac{1}{11} \approx 0.0909$.
\end{enumerate}

\paragraph{P2. Convergence Rate for Convex $L$-Smooth Functions:}
By the Descent Lemma with $\eta = 1/L$:
\begin{equation*}
f(\mathbf{w}_{k+1}) \le f(\mathbf{w}_k) - \frac{1}{2L}\|\nabla f(\mathbf{w}_k)\|_2^2
\end{equation*}
By convexity: $f(\mathbf{w}_k) - f^* \le \nabla f(\mathbf{w}_k)^T (\mathbf{w}_k - \mathbf{w}^*) \le \|\nabla f(\mathbf{w}_k)\|_2 \|\mathbf{w}_k - \mathbf{w}^*\|_2$.
Squaring: $\|\nabla f(\mathbf{w}_k)\|_2^2 \ge \frac{(f(\mathbf{w}_k) - f^*)^2}{\|\mathbf{w}_k - \mathbf{w}^*\|_2^2} \ge \frac{(f(\mathbf{w}_k) - f^*)^2}{R^2}$, where $R = \|\mathbf{w}_0 - \mathbf{w}^*\|_2$ (since distance to optimum is non-increasing).
Let $\Delta_k = f(\mathbf{w}_k) - f^*$. Then:
\begin{equation*}
\Delta_{k+1} \le \Delta_k - \frac{1}{2L R^2} \Delta_k^2
\end{equation*}
Inverting both sides:
\begin{equation*}
\frac{1}{\Delta_{k+1}} - \frac{1}{\Delta_k} = \frac{\Delta_k - \Delta_{k+1}}{\Delta_k \Delta_{k+1}} \ge \frac{\Delta_k - \Delta_{k+1}}{\Delta_k^2} \ge \frac{1}{2L R^2}
\end{equation*}
Summing from $k = 0$ to $T - 1$:
\begin{equation*}
\frac{1}{\Delta_T} - \frac{1}{\Delta_0} \ge \frac{T}{2L R^2} \implies \frac{1}{\Delta_T} \ge \frac{T}{2L R^2} \implies \Delta_T \le \frac{2L R^2}{T} = \mathcal{O}\left(\frac{1}{T}\right)
\end{equation*}

\paragraph{P3. Exact Newton Step on Non-Quadratic Function:}
$f(x) = x^4 + 2x^2 + 4x$.
\begin{enumerate}[label=(\alph*)]
    \item $f'(x) = 4x^3 + 4x + 4$.
    $f''(x) = 12x^2 + 4$.
    \item Starting from $x_0 = 1$:
    $f'(1) = 4(1) + 4(1) + 4 = 12$.
    $f''(1) = 12(1) + 4 = 16$.
    $x_1 = x_0 - \frac{f'(x_0)}{f''(x_0)} = 1 - \frac{12}{16} = 1 - 0.75 = 0.25$.
    At $x_1 = 0.25$:
    $f'(0.25) = 4(1/64) + 4(1/4) + 4 = 1/16 + 1 + 4 = 5.0625$.
    $f''(0.25) = 12(1/16) + 4 = 3/4 + 4 = 4.75$.
    $x_2 = 0.25 - \frac{5.0625}{4.75} \approx 0.25 - 1.0658 = -0.8158$.
    \item Newton's method adjusts step size adaptively via curvature $f''(x)$, leaping directly across steep and flat zones, whereas gradient descent with fixed $\eta = 0.05$ crawls sluggishly ($x_1 = 1 - 0.05(12) = 0.40$).
\end{enumerate}

\paragraph{P4. Momentum Acceleration Dynamics:}
For $f(w) = \frac{1}{2} a w^2$, $f'(w) = aw$.
\begin{align*}
w_{k+1} &= w_k - \alpha a w_k + \beta(w_k - w_{k-1}) = (1 - \alpha a + \beta)w_k - \beta w_{k-1}
\end{align*}
Define state vector $\mathbf{u}_k = \begin{bmatrix} w_k \\ w_{k-1} \end{bmatrix}$. In matrix form:
\begin{equation*}
\mathbf{u}_{k+1} = \begin{bmatrix} 1 - \alpha a + \beta & -\beta \\ 1 & 0 \end{bmatrix} \mathbf{u}_k = T \mathbf{u}_k
\end{equation*}
Characteristic equation of transition matrix $T$:
\begin{equation*}
\det\begin{bmatrix} 1 - \alpha a + \beta - \lambda & -\beta \\ 1 & -\lambda \end{bmatrix} = \lambda^2 - (1 - \alpha a + \beta)\lambda + \beta = 0
\end{equation*}
For complex conjugate eigenvalues, $|\lambda|^2 = \det(T) = \beta$. Stability requires spectral radius $\rho(T) < 1 \implies 0 \le \beta < 1$.

\paragraph{P5. Strongly Convex Linear Convergence Rate:}
For quadratic objective with $A \succ 0$: $\mathbf{w}_k - \mathbf{w}^* = (I - \eta A)^k (\mathbf{w}_0 - \mathbf{w}^*)$.
Taking Euclidean norms: $\|\mathbf{w}_k - \mathbf{w}^*\|_2 \le \|(I - \eta A)^k\|_2 \|\mathbf{w}_0 - \mathbf{w}^*\|_2 = [\rho(I - \eta A)]^k \|\mathbf{w}_0 - \mathbf{w}^*\|_2$.
Eigenvalues of $I - \eta A$ are $1 - \eta \lambda_i$.
With $\eta^* = \frac{2}{L + \mu}$:
\begin{align*}
1 - \eta^* \mu &= 1 - \frac{2\mu}{L + \mu} = \frac{L - \mu}{L + \mu} = \frac{\kappa - 1}{\kappa + 1} \\
1 - \eta^* L &= 1 - \frac{2L}{L + \mu} = \frac{\mu - L}{L + \mu} = -\left( \frac{\kappa - 1}{\kappa + 1} \right)
\end{align*}
The spectral radius is $\rho^* = \max_i |1 - \eta^* \lambda_i| = \frac{\kappa - 1}{\kappa + 1}$.
Therefore: $\|\mathbf{w}_k - \mathbf{w}^*\|_2 \le \left( \frac{\kappa - 1}{\kappa + 1} \right)^k \|\mathbf{w}_0 - \mathbf{w}^*\|_2$.

\subsection*{Part B: Diagnostic Multiple-Choice Rationales}

\paragraph{MCQ1: Gradient Descent Stability Bound}
\begin{itemize}
    \item \textbf{(a) Incorrect.} $\eta = 0.1$ is a valid stable step size ($\eta \le 1/L$), but not the largest non-divergent step size.
    \item \textbf{(b) Correct.} The exact theoretical divergence boundary is $\eta_{\max} = \frac{2}{L}$. For $L = 10$, $\eta_{\max} = \frac{2}{10} = 0.2$. Any step size $\eta \ge 0.2$ creates contraction factor $|1 - \eta L| \ge 1$, causing oscillations to amplify.
    \item \textbf{(c) Incorrect.} $\eta = 0.05$ is overly conservative.
    \item \textbf{(d) Incorrect.} $\eta = 1.0 > 0.2$ will immediately cause exponential divergence.
\end{itemize}

\paragraph{MCQ2: High Condition Number Zig-Zagging}
\begin{itemize}
    \item \textbf{(a) Correct.} The gradient is dominated by the eigenvector corresponding to $\lambda_{\max} = L$. To avoid divergence along this narrow canyon, the step size must be bounded by $\eta < 2/L$. At this tiny step size, progress along the low-curvature direction ($\lambda_{\min} = \mu$) is throttled to rate $\eta \mu = 2/\kappa \ll 1$, forcing thousands of cross-ravine oscillations.
    \item \textbf{(b) Incorrect.} The problem is strictly convex with positive eigenvalues.
    \item \textbf{(c) Incorrect.} The gradient is non-zero along the walls of the ravine.
    \item \textbf{(d) Incorrect.} $A \succ 0$ is strictly positive definite and invertible.
\end{itemize}

% ======================================================================
% CHAPTER 9 SOLUTIONS
% ======================================================================
\section*{Chapter 9: Constrained Optimization \& Convex Sets}

\subsection*{Part A: Analytical Solutions}

\paragraph{P1. Analytical Lagrange Multipliers with Linear Constraint:}
$\min f(x, y) = 2x^2 + y^2$ subject to $x + y = 3$.
\begin{enumerate}[label=(\alph*)]
    \item $\mathcal{L}(x, y, \nu) = 2x^2 + y^2 + \nu(x + y - 3)$.
    \item $\frac{\partial \mathcal{L}}{\partial x} = 4x + \nu = 0 \implies x = -\frac{\nu}{4}$.
    $\frac{\partial \mathcal{L}}{\partial y} = 2y + \nu = 0 \implies y = -\frac{\nu}{2}$.
    $\frac{\partial \mathcal{L}}{\partial \nu} = x + y - 3 = 0$.
    \item Substitute into constraint: $-\frac{\nu}{4} - \frac{\nu}{2} = 3 \implies -\frac{3\nu}{4} = 3 \implies \nu^* = -4$.
    Coordinates: $x^* = -\frac{-4}{4} = 1$, $y^* = -\frac{-4}{2} = 2$.
    \item Objective gradient at $(1, 2)$: $\nabla f(1, 2) = [4(1), 2(2)]^T = [4, 4]^T$.
    Constraint gradient: $\nabla h(1, 2) = [1, 1]^T$.
    $\nabla f(1, 2) = 4 [1, 1]^T = -\nu^* \nabla h = -(-4)[1, 1]^T = [4, 4]^T$. Gradients are strictly collinear.
\end{enumerate}

\paragraph{P2. Maximum Entropy Probability Distribution:}
Maximize $H(\mathbf{p}) = -\sum_{i=1}^K p_i \log p_i$ subject to $\sum_{i=1}^K p_i = 1$.
Equivalently, minimize $\sum_{i=1}^K p_i \log p_i$ subject to $\sum p_i = 1$.
Lagrangian: $\mathcal{L}(\mathbf{p}, \nu) = \sum_{i=1}^K p_i \log p_i + \nu\left(\sum_{i=1}^K p_i - 1\right)$.
Stationarity condition:
\begin{equation*}
\frac{\partial \mathcal{L}}{\partial p_i} = \log p_i + p_i \frac{1}{p_i} + \nu = \log p_i + 1 + \nu = 0 \implies \log p_i = -(1 + \nu) \implies p_i = e^{-(1 + \nu)}
\end{equation*}
Notice that $p_i$ is identical for all $i \in \{1, \dots, K\}$.
Enforcing normalization: $\sum_{i=1}^K p_i = K e^{-(1 + \nu)} = 1 \implies e^{-(1 + \nu)} = \frac{1}{K}$.
Thus $p_i^* = \frac{1}{K}$ for all $i \in \{1, \dots, K\}$. The uniform distribution maximizes entropy.

\paragraph{P3. Convexity of Log-Sum-Exp:}
$f(\mathbf{x}) = \log\left(\sum_{i=1}^d e^{x_i}\right)$. Let $S = \sum_{i=1}^d e^{x_i}$.
\begin{enumerate}[label=(\alph*)]
    \item $\frac{\partial f}{\partial x_k} = \frac{1}{S} \frac{\partial S}{\partial x_k} = \frac{e^{x_k}}{S} = \sigma_k(\mathbf{x})$.
    Since $\sigma_k > 0$ and $\sum_{k=1}^d \sigma_k = \frac{\sum e^{x_k}}{S} = 1$, the gradient $\nabla f(\mathbf{x}) = \boldsymbol{\sigma}(\mathbf{x})$ is a probability vector.
    \item Second partial derivative:
    $\frac{\partial^2 f}{\partial x_k \partial x_j} = \frac{\partial \sigma_k}{\partial x_j} = \sigma_k (\delta_{kj} - \sigma_j)$.
    In matrix form: $\nabla^2 f(\mathbf{x}) = \diag(\boldsymbol{\sigma}) - \boldsymbol{\sigma}\boldsymbol{\sigma}^T$.
    \item For any $\mathbf{v} \in \mathbb{R}^d$:
    \begin{equation*}
    \mathbf{v}^T \nabla^2 f(\mathbf{x})\mathbf{v} = \mathbf{v}^T \diag(\boldsymbol{\sigma})\mathbf{v} - (\mathbf{v}^T \boldsymbol{\sigma})^2 = \sum_{i=1}^d \sigma_i v_i^2 - \left( \sum_{i=1}^d \sigma_i v_i \right)^2
    \end{equation*}
    Recognizing $\sum \sigma_i v_i^2 = \E[V^2]$ and $\sum \sigma_i v_i = \E[V]$ under the probability distribution $\boldsymbol{\sigma}$:
    \begin{equation*}
    \mathbf{v}^T \nabla^2 f(\mathbf{x})\mathbf{v} = \E[V^2] - (\E[V])^2 = \Var(V) \ge 0
    \end{equation*}
    Because variance is always non-negative, $\nabla^2 f(\mathbf{x}) \succeq 0$ for all $\mathbf{x}$, proving log-sum-exp is convex.
\end{enumerate}

\paragraph{P4. Projection onto a Halfspace:}
$\mathcal{H}_- = \{\mathbf{x} \in \mathbb{R}^d \mid \mathbf{a}^T \mathbf{x} \le b\}$ with $\|\mathbf{a}\|_2 = 1$.
\begin{itemize}
    \item If $\mathbf{a}^T \mathbf{x}_0 \le b$, then $\mathbf{x}_0 \in \mathcal{H}_-$, so $\Pi(\mathbf{x}_0) = \mathbf{x}_0$.
    \item If $\mathbf{a}^T \mathbf{x}_0 > b$, $\mathbf{x}_0$ violates the halfspace constraint. The closest point lies on the bounding hyperplane $\mathbf{a}^T \mathbf{x} = b$. The displacement from $\mathbf{x}_0$ orthogonal to the hyperplane is in direction $-\mathbf{a}$.
    Distance to hyperplane: $d = \frac{\mathbf{a}^T \mathbf{x}_0 - b}{\|\mathbf{a}\|_2} = \mathbf{a}^T \mathbf{x}_0 - b$.
    Projected point: $\mathbf{x}^* = \mathbf{x}_0 - (\mathbf{a}^T \mathbf{x}_0 - b)\mathbf{a}$.
\end{itemize}
Combining both cases: $\Pi_{\mathcal{H}_-}(\mathbf{x}_0) = \mathbf{x}_0 - \max(0, \mathbf{a}^T \mathbf{x}_0 - b)\mathbf{a}$.

\paragraph{P5. Jensen's Inequality and AM-GM Proof:}
Consider $f(t) = -\log t$. Since $f''(t) = \frac{1}{t^2} > 0$ for $t > 0$, $f$ is strictly convex on $(0, \infty)$.
Let random variable $X$ take values $a_1, a_2, \dots, a_n > 0$ with equal probabilities $p_i = \frac{1}{n}$.
By Jensen's inequality:
\begin{equation*}
f(\E[X]) \le \E[f(X)] \implies -\log\left( \frac{1}{n}\sum_{i=1}^n a_i \right) \le \frac{1}{n}\sum_{i=1}^n (-\log a_i)
\end{equation*}
Multiplying by $-1$ reverses the inequality:
\begin{equation*}
\log\left( \frac{1}{n}\sum_{i=1}^n a_i \right) \ge \frac{1}{n}\sum_{i=1}^n \log a_i = \log\left( \left(\prod_{i=1}^n a_i\right)^{1/n} \right)
\end{equation*}
Exponentiating both sides:
\begin{equation*}
\frac{1}{n}\sum_{i=1}^n a_i \ge \left( \prod_{i=1}^n a_i \right)^{1/n} \quad (\text{AM} \ge \text{GM})
\end{equation*}

\subsection*{Part B: Diagnostic Multiple-Choice Rationales}

\paragraph{MCQ1: Identification of Non-Convex Sets}
\begin{itemize}
    \item \textbf{(a) Convex.} Any Euclidean norm ball $B = \{\mathbf{x} \mid \|\mathbf{x}\| \le 1\}$ is convex by the triangle inequality.
    \item \textbf{(b) Convex.} The positive semidefinite cone $S^d_+$ is convex: if $A, B \succeq 0$, then $\theta A + (1-\theta)B \succeq 0$.
    \item \textbf{(c) NOT Convex.} The hollow shell $S = \{\mathbf{x} \mid \|\mathbf{x}\| = 1\}$ contains $\mathbf{x} = [1, 0]^T$ and $\mathbf{y} = [-1, 0]^T$. The midpoint $\frac{1}{2}\mathbf{x} + \frac{1}{2}\mathbf{y} = [0, 0]^T$ has norm $0 \ne 1$, lying outside the set.
    \item \textbf{(d) Convex.} Any halfspace $\mathbf{a}^T \mathbf{x} \le b$ is a closed convex set.
\end{itemize}

\paragraph{MCQ2: First-Order Convexity Tangent Hyperplane}
\begin{itemize}
    \item \textbf{(a) Incorrect.} A tangent hyperplane lying above the graph occurs for concave functions, not convex functions.
    \item \textbf{(b) Correct.} By the first-order characterization $f(\mathbf{y}) \ge f(\mathbf{x}) + \nabla f(\mathbf{x})^T (\mathbf{y} - \mathbf{x})$, the affine approximation forms a global supporting hyperplane that lies everywhere below or touches the graph of $f$.
    \item \textbf{(c) Incorrect.} For strictly convex functions, the tangent touches the graph at exactly one unique point.
    \item \textbf{(d) Incorrect.} The slope is determined by the gradient vector $\nabla f$, not the Hessian determinant.
\end{itemize}

% ======================================================================
% CHAPTER 10 SOLUTIONS
% ======================================================================
\section*{Chapter 10: Convex Duality, KKT Conditions \& SVMs}

\subsection*{Part A: Analytical Solutions}

\paragraph{P1. Analytical KKT System on Quadratic Inequality:}
$\min x^2 + y^2$ subject to $4 - x - y \le 0$.
\begin{enumerate}[label=(\alph*)]
    \item Lagrangian: $\mathcal{L}(x, y, \lambda) = x^2 + y^2 + \lambda(4 - x - y)$ with $\lambda \ge 0$.
    \item KKT conditions:
    1. Primal: $4 - x - y \le 0$.
    2. Dual: $\lambda \ge 0$.
    3. Slackness: $\lambda(4 - x - y) = 0$.
    4. Stationarity: $\frac{\partial \mathcal{L}}{\partial x} = 2x - \lambda = 0 \implies x = \lambda/2$; $\frac{\partial \mathcal{L}}{\partial y} = 2y - \lambda = 0 \implies y = \lambda/2$.
    \item Case 1: $\lambda = 0 \implies x = 0, y = 0$. Primal check: $4 - 0 - 0 = 4 \le 0$ (FALSE). Violates primal feasibility.
    \item Case 2: $\lambda > 0 \implies 4 - x - y = 0 \implies 4 - \frac{\lambda}{2} - \frac{\lambda}{2} = 0 \implies \lambda^* = 4$.
    Then $x^* = 2, y^* = 2$.
    Check: $\lambda^* = 4 \ge 0$ (pass), $4 - 2 - 2 = 0 \le 0$ (pass), $\lambda(0) = 0$ (pass), $\nabla \mathcal{L} = \mathbf{0}$ (pass).
    Optimal solution is $(x^*, y^*) = (2, 2)$ with $\lambda^* = 4$.
\end{enumerate}

\paragraph{P2. Dual of a Linear Program (LP):}
Primal: $\min_{\mathbf{x} \in \mathbb{R}^n} \mathbf{c}^T \mathbf{x}$ s.t. $A\mathbf{x} = \mathbf{b}$ and $\mathbf{x} \ge \mathbf{0}$ ($-\mathbf{x} \le \mathbf{0}$).
\begin{enumerate}[label=(\alph*)]
    \item Lagrangian: $\mathcal{L}(\mathbf{x}, \boldsymbol{\lambda}, \boldsymbol{\nu}) = \mathbf{c}^T \mathbf{x} + \boldsymbol{\nu}^T (\mathbf{b} - A\mathbf{x}) - \boldsymbol{\lambda}^T \mathbf{x} = \boldsymbol{\nu}^T \mathbf{b} + (\mathbf{c} - A^T \boldsymbol{\nu} - \boldsymbol{\lambda})^T \mathbf{x}$, with $\boldsymbol{\lambda} \ge \mathbf{0}$.
    \item Dual function $g(\boldsymbol{\lambda}, \boldsymbol{\nu}) = \inf_{\mathbf{x}} \mathcal{L}$:
    An affine function of $\mathbf{x}$ is bounded below if and only if the linear coefficient is identically zero:
    \begin{equation*}
    g(\boldsymbol{\lambda}, \boldsymbol{\nu}) = \begin{cases} \mathbf{b}^T \boldsymbol{\nu} & \text{if } \mathbf{c} - A^T \boldsymbol{\nu} - \boldsymbol{\lambda} = \mathbf{0} \\ -\infty & \text{otherwise} \end{cases}
    \end{equation*}
    \item Dual LP: Maximize $\mathbf{b}^T \boldsymbol{\nu}$ subject to $\mathbf{c} - A^T \boldsymbol{\nu} = \boldsymbol{\lambda} \ge \mathbf{0} \implies A^T \boldsymbol{\nu} \le \mathbf{c}$.
\end{enumerate}

\paragraph{P3. Analytical 3-Point Hard-Margin SVM:}
Points: $\mathbf{x}_1 = [1, 1]^T, y_1 = +1$; $\mathbf{x}_2 = [2, 0]^T, y_2 = +1$; $\mathbf{x}_3 = [0, 0]^T, y_3 = -1$.
\begin{enumerate}[label=(\alph*)]
    \item Constraint: $\sum \lambda_i y_i = \lambda_1 + \lambda_2 - \lambda_3 = 0 \implies \lambda_3 = \lambda_1 + \lambda_2$.
    Kernel products: $\mathbf{x}_1^T \mathbf{x}_1 = 2$, $\mathbf{x}_2^T \mathbf{x}_2 = 4$, $\mathbf{x}_3^T \mathbf{x}_3 = 0$, $\mathbf{x}_1^T \mathbf{x}_2 = 2$, $\mathbf{x}_1^T \mathbf{x}_3 = 0$, $\mathbf{x}_2^T \mathbf{x}_3 = 0$.
    Dual objective:
    \begin{align*}
    D(\lambda_1, \lambda_2) &= (\lambda_1 + \lambda_2 + \lambda_3) - \frac{1}{2}\left[ \lambda_1^2 (2) + \lambda_2^2 (4) + \lambda_3^2 (0) + 2\lambda_1 \lambda_2 (2) - 0 - 0 \right] \\
    &= 2(\lambda_1 + \lambda_2) - \lambda_1^2 - 2\lambda_2^2 - 2\lambda_1 \lambda_2
    \end{align*}
    \item Stationarity: $\frac{\partial D}{\partial \lambda_1} = 2 - 2\lambda_1 - 2\lambda_2 = 0 \implies \lambda_1 + \lambda_2 = 1$.
    $\frac{\partial D}{\partial \lambda_2} = 2 - 4\lambda_2 - 2\lambda_1 = 0 \implies \lambda_1 + 2\lambda_2 = 1$.
    Subtracting: $\lambda_2^* = 0 \implies \lambda_1^* = 1 \implies \lambda_3^* = 1$.
    \item Support vectors are $\mathbf{x}_1$ ($\lambda_1^* = 1 > 0$) and $\mathbf{x}_3$ ($\lambda_3^* = 1 > 0$). Point $\mathbf{x}_2$ is non-support ($\lambda_2^* = 0$).
    Weights: $\mathbf{w}^* = \lambda_1 y_1 \mathbf{x}_1 + \lambda_3 y_3 \mathbf{x}_3 = 1(1)[1, 1]^T + 1(-1)[0, 0]^T = [1, 1]^T$.
    Bias: $y_1(\mathbf{w}^T \mathbf{x}_1 + b) = 1 \implies 1(1(1) + 1(1) + b) = 1 \implies 2 + b = 1 \implies b^* = -1$.
    Hyperplane is $x_1 + x_2 - 1 = 0$.
\end{enumerate}

\paragraph{P4. Soft-Margin SVM KKT Slack Derivation:}
$\mathcal{L} = \frac{1}{2}\|\mathbf{w}\|^2 + C\sum \xi_i - \sum \lambda_i [y_i(\mathbf{w}^T \mathbf{x}_i + b) - 1 + \xi_i] - \sum \mu_i \xi_i$, with $\lambda_i \ge 0, \mu_i \ge 0$.
Stationarity with respect to slack $\xi_i$:
\begin{equation*}
\frac{\partial \mathcal{L}}{\partial \xi_i} = C - \lambda_i - \mu_i = 0 \implies \lambda_i + \mu_i = C \implies \lambda_i = C - \mu_i
\end{equation*}
Since $\mu_i \ge 0$, $\lambda_i \le C$. Combined with dual feasibility $\lambda_i \ge 0$:
\begin{equation*}
0 \le \lambda_i \le C
\end{equation*}
Geometric classifications:
\begin{itemize}
    \item $\lambda_i = 0 \implies \mu_i = C > 0 \implies \xi_i = 0$. Sample is strictly correct and outside margin ($y_i(\mathbf{w}^T \mathbf{x}_i + b) > 1$).
    \item $0 < \lambda_i < C \implies \mu_i > 0 \implies \xi_i = 0$. Sample lies **strictly on margin boundary** ($y_i(\mathbf{w}^T \mathbf{x}_i + b) = 1$).
    \item $\lambda_i = C \implies \mu_i = 0 \implies \xi_i \ge 0$. Sample violates margin (in margin or misclassified).
\end{itemize}

\paragraph{P5. Kernel Ridge Regression via Duality:}
Primal: $\min_{\mathbf{w}, \mathbf{r}} \frac{1}{2}\|\mathbf{r}\|_2^2 + \frac{\lambda}{2}\|\mathbf{w}\|_2^2$ s.t. $\mathbf{r} = X\mathbf{w} - \mathbf{y}$.
Lagrangian with multiplier $\boldsymbol{\alpha}$: $\mathcal{L}(\mathbf{w}, \mathbf{r}, \boldsymbol{\alpha}) = \frac{1}{2}\|\mathbf{r}\|^2 + \frac{\lambda}{2}\|\mathbf{w}\|^2 + \boldsymbol{\alpha}^T (\mathbf{y} - X\mathbf{w} + \mathbf{r})$.
Stationarity:
$\nabla_{\mathbf{w}} \mathcal{L} = \lambda \mathbf{w} - X^T \boldsymbol{\alpha} = \mathbf{0} \implies \mathbf{w} = \frac{1}{\lambda} X^T \boldsymbol{\alpha}$.
$\nabla_{\mathbf{r}} \mathcal{L} = \mathbf{r} + \boldsymbol{\alpha} = \mathbf{0} \implies \mathbf{r} = -\boldsymbol{\alpha}$.
Substituting back into constraint $\mathbf{r} = X\mathbf{w} - \mathbf{y}$:
$-\boldsymbol{\alpha} = X\left(\frac{1}{\lambda}X^T \boldsymbol{\alpha}\right) - \mathbf{y} \implies \left( \frac{1}{\lambda}X X^T + I \right)\boldsymbol{\alpha} = \mathbf{y} \implies (K + \lambda I)\boldsymbol{\alpha} = \lambda \mathbf{y}$.
Thus $\boldsymbol{\alpha} = \lambda(K + \lambda I)^{-1}\mathbf{y}$.
Prediction on new $\mathbf{x}_{\text{test}}$:
\begin{equation*}
\hat{y} = \mathbf{w}^T \mathbf{x}_{\text{test}} = \left(\frac{1}{\lambda} X^T \boldsymbol{\alpha}\right)^T \mathbf{x}_{\text{test}} = \frac{1}{\lambda} \boldsymbol{\alpha}^T (X \mathbf{x}_{\text{test}}) = \mathbf{k}(\mathbf{x}_{\text{test}})^T (K + \lambda I)^{-1}\mathbf{y}
\end{equation*}

\subsection*{Part B: Diagnostic Multiple-Choice Rationales}

\paragraph{MCQ1: Multiplier Interpretation in Soft-Margin SVM}
\begin{itemize}
    \item \textbf{(a) Incorrect.} Points strictly on the margin boundary have $0 < \lambda_i^* < C$.
    \item \textbf{(b) Correct.} By complementary slackness, $\lambda_i^* [y_i(\mathbf{w}^T \mathbf{x}_i + b) - (1 - \xi_i)] = 0$ and $\mu_i^* \xi_i = 0$. When $\lambda_i^* = 0$, $\mu_i^* = C - \lambda_i^* = C > 0$, forcing slack $\xi_i = 0$. The sample lies strictly on the correct side of the margin: $y_i(\mathbf{w}^T \mathbf{x}_i + b) > 1$.
    \item \textbf{(c) Incorrect.} Margin violators have $\lambda_i^* = C$.
    \item \textbf{(d) Incorrect.} The separating hyperplane corresponds to $y_i(\mathbf{w}^T \mathbf{x}_i + b) = 0$, which has slack $\xi_i = 1$ and multiplier $\lambda_i^* = C$.
\end{itemize}

\paragraph{MCQ2: Universal Concavity of Dual Function}
\begin{itemize}
    \item \textbf{(a) Incorrect.} The infimum of numbers can be positive or negative.
    \item \textbf{(b) Correct.} For every fixed $\mathbf{w}$, the Lagrangian $\mathcal{L}(\mathbf{w}, \boldsymbol{\lambda}, \boldsymbol{\nu}) = f_0(\mathbf{w}) + \sum \lambda_i f_i(\mathbf{w}) + \sum \nu_j h_j(\mathbf{w})$ is strictly affine (linear) in $(\boldsymbol{\lambda}, \boldsymbol{\nu})$. The dual function $g(\boldsymbol{\lambda}, \boldsymbol{\nu}) = \inf_{\mathbf{w}} \mathcal{L}(\mathbf{w}, \boldsymbol{\lambda}, \boldsymbol{\nu})$ is the pointwise infimum of affine functions, which is universally concave.
    \item \textbf{(c) Incorrect.} Slater's condition guarantees zero duality gap, not concavity.
    \item \textbf{(d) Incorrect.} Positivity of $\boldsymbol{\lambda}$ defines the dual domain, not the functional concavity.
\end{itemize}

% ======================================================================
% CHAPTER 11 SOLUTIONS
% ======================================================================
\section*{Chapter 11: Probability Foundations \& the Exponential Family}

\subsection*{Part A: Analytical Solutions}

\paragraph{P1. Analytical Linear Transformation of a 2D Gaussian:}
$\mathbf{X} \sim \mathcal{N}(\mathbf{0}, I_2)$, $\mathbf{Y} = A\mathbf{X}$ with $A = \begin{bmatrix} 2 & 1 \\ 0 & 3 \end{bmatrix}$.
\begin{enumerate}[label=(\alph*)]
    \item $\det(A) = 2(3) - 0 = 6$.
    Inverse: $A^{-1} = \frac{1}{6}\begin{bmatrix} 3 & -1 \\ 0 & 2 \end{bmatrix} = \begin{bmatrix} 1/2 & -1/6 \\ 0 & 1/3 \end{bmatrix}$.
    \item $f_{\mathbf{X}}(\mathbf{x}) = \frac{1}{2\pi} e^{-\frac{1}{2}(x_1^2 + x_2^2)}$. By the Jacobian rule~\eqref{eq:jacobian_transform}:
    \begin{equation*}
    f_{\mathbf{Y}}(\mathbf{y}) = f_{\mathbf{X}}(A^{-1}\mathbf{y})\frac{1}{|\det(A)|} = \frac{1}{2\pi \cdot 6} \exp\left( -\frac{1}{2} (A^{-1}\mathbf{y})^T (A^{-1}\mathbf{y}) \right)
    \end{equation*}
    $A^{-1}\mathbf{y} = [y_1/2 - y_2/6, y_2/3]^T$.
    $\|A^{-1}\mathbf{y}\|_2^2 = (y_1/2 - y_2/6)^2 + (y_2/3)^2 = \frac{y_1^2}{4} - \frac{y_1 y_2}{6} + \frac{y_2^2}{36} + \frac{y_2^2}{9} = \frac{y_1^2}{4} - \frac{y_1 y_2}{6} + \frac{5y_2^2}{36}$.
    Thus $f_{\mathbf{Y}}(\mathbf{y}) = \frac{1}{12\pi} \exp\left( -\frac{1}{2} \left[ \frac{y_1^2}{4} - \frac{y_1 y_2}{6} + \frac{5y_2^2}{36} \right] \right)$.
    \item Covariance: $\Sigma_{\mathbf{Y}} = A A^T = \begin{bmatrix} 2 & 1 \\ 0 & 3 \end{bmatrix}\begin{bmatrix} 2 & 0 \\ 1 & 3 \end{bmatrix} = \begin{bmatrix} 5 & 3 \\ 3 & 9 \end{bmatrix}$.
    $|\Sigma_{\mathbf{Y}}| = 45 - 9 = 36 = (\det A)^2$.
    $\Sigma_{\mathbf{Y}}^{-1} = \frac{1}{36}\begin{bmatrix} 9 & -3 \\ -3 & 5 \end{bmatrix} = \begin{bmatrix} 1/4 & -1/12 \\ -1/12 & 5/36 \end{bmatrix}$.
    $\mathbf{y}^T \Sigma_{\mathbf{Y}}^{-1}\mathbf{y} = \frac{y_1^2}{4} - \frac{y_1 y_2}{6} + \frac{5y_2^2}{36}$, matching the exponent exactly.
\end{enumerate}

\paragraph{P2. Non-Monotonic Transformation (Chi-Square Distribution):}
$X \sim \mathcal{N}(0, 1)$, $Y = X^2$.
\begin{enumerate}[label=(\alph*)]
    \item For $y > 0$, $x^2 = y \implies x_1(y) = \sqrt{y}$ and $x_2(y) = -\sqrt{y}$.
    \item $\frac{dx_1}{dy} = \frac{1}{2\sqrt{y}}$, $\frac{dx_2}{dy} = -\frac{1}{2\sqrt{y}}$.
    Absolute values: $\left|\frac{dx_1}{dy}\right| = \left|\frac{dx_2}{dy}\right| = \frac{1}{2\sqrt{y}}$.
    \item Summing branches:
    \begin{align*}
    f_Y(y) &= f_X(\sqrt{y})\left|\frac{dx_1}{dy}\right| + f_X(-\sqrt{y})\left|\frac{dx_2}{dy}\right| = \frac{1}{\sqrt{2\pi}}e^{-y/2} \left(\frac{1}{2\sqrt{y}}\right) + \frac{1}{\sqrt{2\pi}}e^{-y/2} \left(\frac{1}{2\sqrt{y}}\right) \\
    &= \frac{1}{\sqrt{2\pi}} y^{-1/2} e^{-y/2} \quad \text{for } y > 0
    \end{align*}
    This is the Chi-Square distribution with 1 degree of freedom: $\chi^2(1)$.
\end{enumerate}

\paragraph{P3. Exponential Distribution in Canonical Form:}
$f(x; \lambda) = \lambda e^{-\lambda x}$ for $x \ge 0$.
\begin{enumerate}[label=(\alph*)]
    \item Rewriting: $f(x; \lambda) = \exp(-\lambda x + \log\lambda) = 1 \cdot \exp((-\lambda) x - (-\log\lambda))$.
    Matching with $h(x)\exp(\eta T(x) - A(\eta))$:
    $h(x) = \mathbb{I}(x \ge 0)$, $T(x) = x$, $\eta = -\lambda$, $A(\eta) = -\log(-\eta)$.
    \item First derivative: $A'(\eta) = -\frac{1}{-\eta}(-1) = -\frac{1}{\eta} = -\frac{1}{-\lambda} = \frac{1}{\lambda} = \E[X]$.
    \item Second derivative: $A''(\eta) = \frac{d}{d\eta}\left(-\frac{1}{\eta}\right) = \frac{1}{\eta^2} = \frac{1}{(-\lambda)^2} = \frac{1}{\lambda^2} = \Var(X)$.
\end{enumerate}

\paragraph{P4. Uncorrelated but Strongly Dependent Variables:}
$X \sim \mathcal{U}[-1, 1], Y = X^2$.
\begin{enumerate}[label=(\alph*)]
    \item $\E[X] = \int_{-1}^1 \frac{x}{2}dx = 0$.
    $\E[Y] = \E[X^2] = \int_{-1}^1 \frac{x^2}{2}dx = \left[\frac{x^3}{6}\right]_{-1}^1 = \frac{1}{3}$.
    $\E[XY] = \E[X^3] = \int_{-1}^1 \frac{x^3}{2}dx = 0$ (odd integrand).
    \item $\Cov(X, Y) = \E[XY] - \E[X]\E[Y] = 0 - 0(1/3) = 0$.
    \item $P(Y \le 0.25) = P(X^2 \le 0.25) = P(-0.5 \le X \le 0.5) = \frac{0.5 - (-0.5)}{2} = 0.5$.
    Given $X \ge 0.8$, $Y = X^2 \ge 0.64 > 0.25$, so $P(Y \le 0.25 \mid X \ge 0.8) = 0$.
    Since $0 \ne 0.5$, $X$ and $Y$ are statistically dependent.
\end{enumerate}

\paragraph{P5. Softmax Link from Categorical Distribution:}
$P(X = k) = p_k$ with $\sum_{k=1}^K p_k = 1 \implies p_K = 1 - \sum_{k=1}^{K-1} p_k$.
\begin{align*}
p(x; \mathbf{p}) &= \prod_{k=1}^K p_k^{\mathbb{I}(x = k)} = \exp\left( \sum_{k=1}^{K-1} \mathbb{I}(x=k)\log p_k + \left(1 - \sum_{k=1}^{K-1}\mathbb{I}(x=k)\right)\log p_K \right) \\
&= \exp\left( \sum_{k=1}^{K-1} \mathbb{I}(x=k)\log\left(\frac{p_k}{p_K}\right) + \log p_K \right)
\end{align*}
Natural parameters: $\eta_k = \log(p_k / p_K) \implies p_k = p_K e^{\eta_k}$ for $k < K$.
Normalizing: $\sum_{k=1}^K p_k = p_K \left(1 + \sum_{k=1}^{K-1} e^{\eta_k}\right) = 1 \implies p_K = \frac{1}{1 + \sum_{j=1}^{K-1} e^{\eta_j}}$.
For $k \in \{1, \dots, K-1\}$:
\begin{equation*}
p_k = \frac{e^{\eta_k}}{1 + \sum_{j=1}^{K-1} e^{\eta_j}}
\end{equation*}
Setting $\eta_K = 0$ yields the symmetric Softmax function: $p_k = \frac{e^{\eta_k}}{\sum_{j=1}^K e^{\eta_j}}$.

\subsection*{Part B: Diagnostic Multiple-Choice Rationales}

\paragraph{MCQ1: Linear Transformation Density Rule}
\begin{itemize}
    \item \textbf{(a) Incorrect.} The determinant factor must appear in the denominator or with inverse power.
    \item \textbf{(b) Correct.} From $\mathbf{x} = A^{-1}(\mathbf{y} - \mathbf{b})$, the Jacobian is $J = A^{-1}$. The volume scaling factor is $|\det(A^{-1})| = \frac{1}{|\det(A)|}$, yielding $f_{\mathbf{Y}}(\mathbf{y}) = f_{\mathbf{X}}(A^{-1}(\mathbf{y} - \mathbf{b})) \frac{1}{|\det(A)|}$.
    \item \textbf{(c) Incorrect.} Uses forward mapping $A\mathbf{y} + \mathbf{b}$ rather than the pre-image $A^{-1}(\mathbf{y} - \mathbf{b})$.
    \item \textbf{(d) Incorrect.} Neglects the Jacobian volume distortion factor.
\end{itemize}

\paragraph{MCQ2: Hessian of the Log-Partition Function}
\begin{itemize}
    \item \textbf{(a) Incorrect.} Covariance matrices are positive semidefinite, not negative definite.
    \item \textbf{(b) Correct.} By differentiating under the integral sign, $\nabla^2 A(\boldsymbol{\eta}) = \E[\mathbf{T}\mathbf{T}^T] - \E[\mathbf{T}]\E[\mathbf{T}]^T = \Cov(\mathbf{T}(X))$. Since any covariance matrix satisfies $\mathbf{v}^T \Cov(\mathbf{T})\mathbf{v} = \Var(\mathbf{v}^T \mathbf{T}) \ge 0$, the Hessian is positive semidefinite, proving $A(\boldsymbol{\eta})$ is strictly convex.
    \item \textbf{(c) Incorrect.} The Hessian is non-zero for any non-degenerate distribution.
    \item \textbf{(d) Incorrect.} The Hessian is equal to the Fisher information matrix multiplied by $+1$, not $-1$.
\end{itemize}

% ======================================================================
% CHAPTER 12 SOLUTIONS
% ======================================================================
\section*{Chapter 12: Multivariate Gaussians, GMMs, EM \& Concentration}

\subsection*{Part A: Analytical Solutions}

\paragraph{P1. Analytical 2D Mahalanobis Distance:}
$\boldsymbol{\mu} = [1, 2]^T, \Sigma = \begin{bmatrix} 4 & 2 \\ 2 & 2 \end{bmatrix}$.
\begin{enumerate}[label=(\alph*)]
    \item $|\Sigma| = 4(2) - 2(2) = 8 - 4 = 4$.
    Inverse: $\Sigma^{-1} = \frac{1}{4}\begin{bmatrix} 2 & -2 \\ -2 & 4 \end{bmatrix} = \begin{bmatrix} 1/2 & -1/2 \\ -1/2 & 1 \end{bmatrix}$.
    \item Difference vector for $\mathbf{x} = [3, 4]^T$: $\mathbf{x} - \boldsymbol{\mu} = [3 - 1, 4 - 2]^T = [2, 2]^T$.
    Mahalanobis distance:
    \begin{equation*}
    D_M^2 = [2, 2] \begin{bmatrix} 1/2 & -1/2 \\ -1/2 & 1 \end{bmatrix} \begin{bmatrix} 2 \\ 2 \end{bmatrix} = [2, 2] \begin{bmatrix} 1 - 1 \\ -1 + 2 \end{bmatrix} = [2, 2] \begin{bmatrix} 0 \\ 1 \end{bmatrix} = 2
    \end{equation*}
    $D_M(\mathbf{x}, \boldsymbol{\mu}) = \sqrt{2} \approx 1.414$.
    \item Characteristic equation of $\Sigma$: $(4 - \lambda)(2 - \lambda) - 4 = \lambda^2 - 6\lambda + 4 = 0$.
    Eigenvalues: $\lambda = \frac{6 \pm \sqrt{36 - 16}}{2} = 3 \pm \sqrt{5}$.
    $\lambda_1 = 3 + \sqrt{5} \approx 5.236$ (semi-major axis length $\sqrt{\lambda_1} \approx 2.288$).
    $\lambda_2 = 3 - \sqrt{5} \approx 0.764$ (semi-minor axis length $\sqrt{\lambda_2} \approx 0.874$).
    Eigenvector for $\lambda_1$: $(4 - (3+\sqrt{5}))v_1 + 2v_2 = 0 \implies (1 - \sqrt{5})v_1 + 2v_2 = 0 \implies \mathbf{q}_1 \propto [\sqrt{5}-1, 2]^T$.
\end{enumerate}

\paragraph{P2. Unbiasedness and Degrees of Freedom of Sample Covariance:}
$\hat{\Sigma} = \frac{1}{N}\sum_{i=1}^N (\mathbf{X}_i - \bar{\mathbf{X}})(\mathbf{X}_i - \bar{\mathbf{X}})^T$.
Expanding $\mathbf{X}_i - \bar{\mathbf{X}} = (\mathbf{X}_i - \boldsymbol{\mu}) - (\bar{\mathbf{X}} - \boldsymbol{\mu})$:
\begin{align*}
\sum_{i=1}^N (\mathbf{X}_i - \bar{\mathbf{X}})(\mathbf{X}_i - \bar{\mathbf{X}})^T &= \sum_{i=1}^N (\mathbf{X}_i - \boldsymbol{\mu})(\mathbf{X}_i - \boldsymbol{\mu})^T - N(\bar{\mathbf{X}} - \boldsymbol{\mu})(\bar{\mathbf{X}} - \boldsymbol{\mu})^T
\end{align*}
Taking expectations on both sides:
\begin{equation*}
\E\left[ \sum_{i=1}^N (\mathbf{X}_i - \boldsymbol{\mu})(\mathbf{X}_i - \boldsymbol{\mu})^T \right] = \sum_{i=1}^N \Sigma = N \Sigma
\end{equation*}
For the sample mean, since $\mathbf{X}_i$ are independent:
\begin{equation*}
\E\left[ (\bar{\mathbf{X}} - \boldsymbol{\mu})(\bar{\mathbf{X}} - \boldsymbol{\mu})^T \right] = \Var(\bar{\mathbf{X}}) = \frac{1}{N}\Sigma
\end{equation*}
Substituting:
\begin{equation*}
\E\left[ \sum_{i=1}^N (\mathbf{X}_i - \bar{\mathbf{X}})(\mathbf{X}_i - \bar{\mathbf{X}})^T \right] = N\Sigma - N\left(\frac{1}{N}\Sigma\right) = (N - 1)\Sigma
\end{equation*}
Dividing by $N$: $\E[\hat{\Sigma}_{\text{MLE}}] = \frac{N - 1}{N}\Sigma$.

\paragraph{P3. Analytical E-step Calculation for a 1D GMM:}
$K = 2$, $\pi_1 = 0.6, \mu_1 = 0, \sigma_1^2 = 1$; $\pi_2 = 0.4, \mu_2 = 3, \sigma_2^2 = 1$.
Observation $x = 1.5$.
\begin{enumerate}[label=(\alph*)]
    \item $\mathcal{N}(1.5; 0, 1) = \frac{1}{\sqrt{2\pi}}e^{-(1.5-0)^2/2} = \frac{1}{\sqrt{2\pi}}e^{-2.25/2} = \frac{1}{\sqrt{2\pi}}e^{-1.125} \approx 0.3989 \times 0.3247 = 0.1295$.
    $\mathcal{N}(1.5; 3, 1) = \frac{1}{\sqrt{2\pi}}e^{-(1.5-3)^2/2} = \frac{1}{\sqrt{2\pi}}e^{-(-1.5)^2/2} = \frac{1}{\sqrt{2\pi}}e^{-1.125} \approx 0.1295$.
    Both Gaussian densities evaluate to the identical value $d = 0.1295$ because $|1.5 - 0| = |1.5 - 3| = 1.5$.
    \item Responsibilities:
    Numerator 1: $\pi_1 \mathcal{N}_1 = 0.6 \times d$.
    Numerator 2: $\pi_2 \mathcal{N}_2 = 0.4 \times d$.
    Denominator: $0.6d + 0.4d = 1.0d$.
    $\gamma_1 = \frac{0.6d}{1.0d} = 0.60$.
    $\gamma_2 = \frac{0.4d}{1.0d} = 0.40$.
    \item Observation $x = 1.5$ is more likely to belong to Cluster 1 ($\gamma_1 = 60\%$ vs $\gamma_2 = 40\%$). Even though $x$ is equidistant from both cluster centers, Cluster 1 has a higher prior mixture weight ($\pi_1 = 0.6 > \pi_2 = 0.4$).
\end{enumerate}

\paragraph{P4. Sample Complexity via Hoeffding's Inequality:}
Let $X_i \in [0, 1]$ indicate policy support. $\epsilon = 0.02, \delta = 0.05$.
By Hoeffding's inequality for bounded variables in $[0, 1]$:
\begin{equation*}
P(|\bar{X}_N - p| \ge \epsilon) \le 2\exp(-2N\epsilon^2) \le \delta
\end{equation*}
Taking the natural logarithm:
\begin{align*}
-2N\epsilon^2 &\le \log\left(\frac{\delta}{2}\right) \implies 2N\epsilon^2 \ge \log\left(\frac{2}{\delta}\right) \implies N \ge \frac{\log(2/\delta)}{2\epsilon^2}
\end{align*}
Substitute $\delta = 0.05 \implies 2/\delta = 40$, and $\epsilon = 0.02 \implies \epsilon^2 = 0.0004$:
\begin{equation*}
N \ge \frac{\log(40)}{2(0.0004)} = \frac{3.688879}{0.0008} \approx 4611.1
\end{equation*}
The pollster must survey at least $N = 4{,}612$ independent voters.

\paragraph{P5. Chebyshev vs. Markov Bound Comparison:}
$X \sim \text{Binomial}(100, 0.5)$.
\begin{enumerate}[label=(\alph*)]
    \item Mean $\mu = N p = 100(0.5) = 50$.
    Variance $\sigma^2 = N p(1 - p) = 100(0.5)(0.5) = 25$. Standard deviation $\sigma = 5$.
    \item By Markov's inequality ($X \ge 0$):
    \begin{equation*}
    P(X \ge 75) \le \frac{\E[X]}{75} = \frac{50}{75} = \frac{2}{3} \approx 0.6667
    \end{equation*}
    \item By Chebyshev's inequality:
    $P(X \ge 75) \le P(|X - 50| \ge 25) \le \frac{\sigma^2}{25^2} = \frac{25}{625} = \frac{1}{25} = 0.04$.
    \item Markov yields a trivial bound ($\le 66.7\%$), while Chebyshev yields a sharp bound ($\le 4.0\%$). Chebyshev is dramatically tighter because it exploits second-order variance information ($\sigma^2 = 25$) to penalize distance from the mean quadratically ($\frac{1}{\epsilon^2}$), whereas Markov utilizes only the first moment.
\end{enumerate}

\subsection*{Part B: Diagnostic Multiple-Choice Rationales}

\paragraph{MCQ1: Monotonicity of Log-Likelihood in EM}
\begin{itemize}
    \item \textbf{(a) Incorrect.} The EM algorithm maximizes likelihood; it does not decrease toward zero.
    \item \textbf{(b) Correct.} At each EM iteration, the E-step equates the ELBO to the marginal log-likelihood by eliminating the KL divergence, and the M-step maximizes the ELBO. Consequently, $\ell(\boldsymbol{\theta}^{(t+1)}) \ge \text{ELBO}(q^{(t+1)}, \boldsymbol{\theta}^{(t+1)}) \ge \text{ELBO}(q^{(t+1)}, \boldsymbol{\theta}^{(t)}) = \ell(\boldsymbol{\theta}^{(t)})$, guaranteeing monotonic ascent.
    \item \textbf{(c) Incorrect.} Due to non-convexity, EM is only guaranteed to converge to a stationary point or local maximum.
    \item \textbf{(d) Incorrect.} Log-likelihood values are independent of linear separability.
\end{itemize}

\paragraph{MCQ2: Exponential Concentration Bound}
\begin{itemize}
    \item \textbf{(a) Incorrect.} Markov's inequality gives a loose reciprocal bound ($\mathcal{O}(1/\epsilon)$) using only the mean.
    \item \textbf{(b) Incorrect.} Chebyshev's inequality gives a polynomial bound ($\mathcal{O}(1/(N\epsilon^2))$) using the variance.
    \item \textbf{(c) Correct.} Hoeffding's inequality provides an **exponentially decaying bound** in sample size $N$: $P(|\bar{X} - \mu| \ge \epsilon) \le 2\exp(-2N\epsilon^2)$, foundational to statistical learning theory.
    \item \textbf{(d) Incorrect.} The Cauchy-Schwarz inequality is an algebraic vector inner product bound, not a probability concentration inequality.
\end{itemize}



\printindex

\end{document}

